% ===========================================================================
%  E-Trike System Documentation v0.0.3-alpha
%  Generated: 2026-06-24
%  Compile:  pdflatex -interaction=nonstopmode etrike-0.0.2-alpha.tex
%            pdflatex -interaction=nonstopmode etrike-0.0.2-alpha.tex  (twice for TOC)
% ===========================================================================

\documentclass[12pt,a4paper,oneside]{report}

% ── Packages ──────────────────────────────────────────────────────────
\usepackage[utf8]{inputenc}
\usepackage[T1]{fontenc}
\usepackage{XCharter}
\usepackage[inner=30mm,outer=20mm,top=22mm,bottom=25mm,footskip=14mm]{geometry}
\usepackage{amsmath,amssymb}
\usepackage{booktabs}
\usepackage{longtable}
\usepackage{array}

% ── Table spacing and behavior ──────────────────────────────────────
\setlength{\LTpre}{6pt plus 2pt minus 2pt}
\setlength{\LTpost}{6pt plus 2pt minus 2pt}
\setlength{\tabcolsep}{5pt}
\renewcommand{\arraystretch}{1.15}
% Allow longtables to break more gracefully
\raggedbottom
\usepackage{listings}
\usepackage[table]{xcolor}
\usepackage{soul}
\usepackage{hyperref}
\usepackage{fancyhdr}
\usepackage{titlesec}
\usepackage{enumitem}
\usepackage{alltt}
\usepackage{makecell}
\usepackage{caption}
\usepackage{float}
\usepackage{parskip}
\usepackage{pmboxdraw}
\usepackage{textcomp}
\usepackage{calc}

% ── Pandoc compatibility ────────────────────────────────────────────
\providecommand{\tightlist}{\setlength{\itemsep}{0pt}\setlength{\parskip}{0pt}}

% ── Prevent overfull hboxes in narrow table columns ─────────────────
\setlength{\emergencystretch}{3em}
\tolerance=3000
\hbadness=10000

% ── Unicode character support for source docs ─────────────────────────
\DeclareUnicodeCharacter{03A9}{\ensuremath{\Omega}}
\DeclareUnicodeCharacter{03C0}{\ensuremath{\pi}}
\DeclareUnicodeCharacter{03B4}{\ensuremath{\delta}}
\DeclareUnicodeCharacter{03C9}{\ensuremath{\omega}}
\DeclareUnicodeCharacter{03B1}{\ensuremath{\alpha}}
\DeclareUnicodeCharacter{03B2}{\ensuremath{\beta}}
\DeclareUnicodeCharacter{03B8}{\ensuremath{\theta}}
\DeclareUnicodeCharacter{2264}{\ensuremath{\leq}}
\DeclareUnicodeCharacter{2265}{\ensuremath{\geq}}
\DeclareUnicodeCharacter{00B0}{\ensuremath{{}^\circ}}
\DeclareUnicodeCharacter{00B1}{\ensuremath{\pm}}
\DeclareUnicodeCharacter{00D7}{\ensuremath{\times}}
\DeclareUnicodeCharacter{00B5}{\ensuremath{\mu}}
\DeclareUnicodeCharacter{2192}{\ensuremath{\rightarrow}}
\DeclareUnicodeCharacter{2190}{\ensuremath{\leftarrow}}
\DeclareUnicodeCharacter{2194}{\ensuremath{\leftrightarrow}}
\DeclareUnicodeCharacter{25BA}{$\triangleright$}
\DeclareUnicodeCharacter{25C0}{$\triangleleft$}
\DeclareUnicodeCharacter{25BC}{$\triangledown$}
\DeclareUnicodeCharacter{25B2}{$\triangle$}

% ── Title formatting (14pt chapter/section headings, 12pt body) ───────
\titleformat{\chapter}[display]
  {\normalfont\bfseries\fontsize{14}{17}\selectfont}
  {\chaptertitlename\ \thechapter}{12pt}{}
\titlespacing*{\chapter}{0pt}{0pt}{20pt}
\titleformat{\section}
  {\normalfont\bfseries\fontsize{14}{17}\selectfont}
  {\thesection}{1em}{}
\titleformat{\subsection}
  {\normalfont\bfseries\fontsize{13}{16}\selectfont}
  {\thesubsection}{1em}{}
\titleformat{\subsubsection}
  {\normalfont\bfseries\fontsize{12}{15}\selectfont}
  {\thesubsubsection}{1em}{}

% ── Hyperref setup ────────────────────────────────────────────────────
\hypersetup{
  colorlinks=true,
  linkcolor=black,
  urlcolor=blue!60!black,
  citecolor=black,
  pdfauthor={E-Trike Project},
  pdftitle={E-Trike System Documentation v0.0.3-alpha},
  pdfsubject={Three-node distributed control vehicle architecture},
  pdfkeywords={CAN bus, embedded systems, distributed control, steer-by-wire, brake-by-wire, tricycle, vehicle architecture},
  bookmarksnumbered=true,
  bookmarksopen=true,
  bookmarksopenlevel=1
}

% ── Code listing style ────────────────────────────────────────────────
\lstdefinestyle{etrike}{
  basicstyle=\ttfamily\scriptsize,
  breaklines=true,
  breakatwhitespace=true,
  postbreak=\mbox{$\hookrightarrow$\space},
  frame=single,
  framesep=4pt,
  xleftmargin=6pt,
  xrightmargin=6pt,
  aboveskip=8pt,
  belowskip=8pt,
  showstringspaces=false,
  columns=flexible,
  keepspaces=true,
  tabsize=2,
  captionpos=b
}
\lstset{style=etrike}

% ── Header / footer ───────────────────────────────────────────────────
\pagestyle{fancy}
\fancyhf{}
\fancyhead[L]{\small\textit{E-Trike System Documentation}}
\fancyhead[R]{\small\textit{v0.0.3-alpha}}
\fancyfoot[C]{\thepage}
\renewcommand{\headrulewidth}{0.4pt}
\renewcommand{\footrulewidth}{0pt}
\setlength{\headheight}{14pt}

% ── Part page style ───────────────────────────────────────────────────
\makeatletter
\renewcommand{\@endpart}{\vfill\newpage}
\makeatother

% ── Convenience commands ──────────────────────────────────────────────
\newcommand{\canid}[1]{\texttt{0x#1}}
\newcommand{\canname}[1]{\textsc{#1}}
\newcommand{\gpio}[1]{\texttt{GPIO#1}}
\newcommand{\voltage}[1]{\ensuremath{#1\,\text{V}}}
\newcommand{\bitrate}[1]{\ensuremath{#1\,\text{kbit/s}}}
\newcommand{\frequency}[1]{\ensuremath{#1\,\text{Hz}}}
\newcommand{\mcu}[1]{\texttt{#1}}
\newcommand{\gap}[1]{\textbf{Gap \##1}}
\newcommand{\gapresolved}[1]{\textbf{RESOLVED \##1}}
\newcommand{\syntree}[1]{\textsf{#1}}

\newcolumntype{P}[1]{>{\centering\arraybackslash}p{#1}}
\newcolumntype{L}[1]{>{\raggedright\arraybackslash}p{#1}}

% ── Page style for title page ─────────────────────────────────────────
\fancypagestyle{plain}{
  \fancyhf{}
  \renewcommand{\headrulewidth}{0pt}
}

\begin{document}

% ═══════════════════════════════════════════════════════════════════════
%  TITLE PAGE
% ═══════════════════════════════════════════════════════════════════════

\thispagestyle{empty}

\vspace*{60mm}

\begin{center}
  {\fontsize{28}{34}\selectfont\textbf{E-Trike}}\\[8mm]
  {\fontsize{22}{28}\selectfont\textbf{System Documentation}}\\[20mm]
  {\fontsize{14}{18}\selectfont\textit{Three-Node Distributed Control Architecture}}\\[30mm]
  {\fontsize{12}{15}\selectfont Version 0.0.2-alpha}\\[4mm]
  {\fontsize{11}{14}\selectfont 2026-06-24}
\end{center}

\vfill

\begin{center}
  {\footnotesize
  Three-Node Distributed Control Architecture\\[2mm]
  Dual CAN Bus (500 kbit/s) --- Real-Time Operating System --- Hard Realtime Safety\\[2mm]
  Steer-by-Wire \textbullet\ Electro-Hydraulic Brake-by-Wire
  }
\end{center}

\newpage

% ═══════════════════════════════════════════════════════════════════════
%  TABLE OF CONTENTS
% ═══════════════════════════════════════════════════════════════════════

\setcounter{tocdepth}{1}
\tableofcontents
\newpage

% ═══════════════════════════════════════════════════════════════════════
%  PART I -- SYSTEM ARCHITECTURE & CAN DICTIONARY
% ═══════════════════════════════════════════════════════════════════════

\part{System Architecture \& CAN Dictionary}
\label{part:architecture}

\chapter{E-Trike System Architecture}\label{e-trike-system-architecture}

Four-node distributed control: \textbf{Jetson Orin} (ROS 2 perception/planning),
\textbf{RT ESP32-S3} (realtime physics, steering, brake \& CAN gateway),
\textbf{SYS ESP32-S3} (safety \& body control), \textbf{MTR STM32} (motor
actuation).

\begin{quote}
\textbf{Variant:} A consolidated single-controller version exists at
\href{rt-aurix-lite/rt-aurix-lite-architecture.md}{\texttt{rt-aurix-lite/rt-aurix-lite-architecture.md}}
--- combines RT+SYS into one AURIX TC3xx on a single CAN bus. This document
describes the distributed reference architecture. Both variants share the same
CAN protocol and actuator interfaces.
\end{quote}

Two physical CAN buses at 500 kbit/s. RT bridges selected messages between
buses. Actuators are \textbf{SYNTREE} CAN modules: EPS-C (steer-by-wire) and SEB
(electro-hydraulic brake). \textbf{Mode-gated dual control (Option D):} RT
directly commands both EPS-C and SEB in AUTO mode; SYS commands SEB in MANUAL
mode (lever pass-through) and on ESTOP. This ensures no single MCU failure takes
both actuators, and AUTO-mode brake+steer are 1-hop from the kinematics engine.
Motor control is on a dedicated STM32 board (MTR) for safety isolation per ISO
26262 EGAS 3-level concept.

\begin{quote}
\textbf{Autoware.Auto v0.0.4-alpha:} A new ROS 2 node on the Jetson
(\texttt{autoware\_vehicle\_bridge}) provides Autoware.Auto topic compatibility
(\texttt{AckermannControlCommand} in,
\texttt{VelocityReport}/\texttt{SteeringReport} out). The CAN protocol across
both buses is unchanged. See §5.1.
\end{quote}

\begin{center}\rule{0.5\linewidth}{0.5pt}\end{center}

\section{1. Topology}\label{topology}

\begin{verbatim}
  ┌────────────────── High-Level CAN (500 kbit/s) ──────────────────┐
  │                                                                  │
  │  ┌──────────┐            ┌──────────────┐                       │
  │  │  Jetson  │            │  RT ESP32-S3 │                       │
  │  │  Orin │            │              │                       │
  │  │          │            │ Physics      │                       │
  │  │ ROS 2    │            │ Steering     │                       │
  │  │ Planning │            │ Gateway      │                       │
  │  └────┬─────┘            │              │                       │
  │       │                  └──────┬───────┘                       │
  │  TX:  0x300,0x301,    TX: 0x011,0x120, │                        │
  │       0x302,0x001           0x210,0x220,│                        │
  │                              0x400,0x600,│                       │
  │  RX:  0x011,0x120,          0x001,0x7FC │                        │
  │       0x210,0x220,      RX: 0x300,0x301,│                        │
  │       0x400,0x600,          0x302,0x001 │                        │
  │       0x001,0x7FD                       │                        │
  └─────────────────────────────────────────┘                        │
                                            │                        │
                 ┌──────────────────────────┘                        │
                 │                                                    │
  ┌──────────────▼─────── Low-Level CAN (500 kbit/s) ───────────────┐│
  │                                                                  ││
  │  ┌──────────────┐      ┌──────────────┐                         ││
  │  │  RT ESP32-S3 │      │ SYS ESP32-S3 │                         ││
  │  │  (gateway)   │      │              │                         ││
  │  └──────┬───────┘      │ Safety       │                         ││
  │         │              │ Motor        │                         ││
  │    TX:  0x169,0x204,   │ Brake        │                         ││
  │         0x302,0x001,   └──────┬───────┘                         ││
  │         0x7FD                 │                                  ││
  │                          TX:  0x7B9,0x011,                      ││
  │    RX:  0x001,0x011,          0x012,0x110,                      ││
  │         0x110,0x120,          0x120,0x600,                      ││
  │         0x201,0x600,          0x001,0x7FE                       ││
  │         0x7FD,0x7FE              RX:  0x001,0x204,                    ││
  │                               0x302,0x721,                      ││
  │                               0x7FD                             ││
  └─────────────────────────────────────────────────────────────────┘│
                                        │                            │
                ┌───────────────────────┼────────────────────┐       │
                │                       │                    │       │
          ┌─────▼─────┐          ┌─────▼─────┐        ┌─────▼─────┐ │
          │  SYNTREE  │          │  SYNTREE  │        │   DC-DC   │ │
          │  SEB      │          │  EPS-C    │        │ Converter │ │
          │  (Brake)  │          │ (Steering)│        │ 72V→12V   │ │
          │ 0x7B9 cmd │          │ 0x169 cmd │        │ (0x012)   │ │
          │ 0x721 stat│          │ 0x201 stat│        └───────────┘ │
          └───────────┘          └───────────┘                      │
                                                                    │
  ┌───────────┐                                                    │
  │  Motor    │  (analog: 0–5 V throttle via MCP4725,              │
  │Controller │   72 V gear lines via relay module, from SYS)      │
  └───────────┘                                                    │
\end{verbatim}

\begin{center}\rule{0.5\linewidth}{0.5pt}\end{center}

\section{2. CAN message catalog}\label{can-message-catalog}

\subsection{2.1 Low-level CAN}\label{low-level-can}

\begin{longtable}[]{@{}
  >{\raggedright\arraybackslash}p{(\columnwidth - 14\tabcolsep) * \real{0.0678}}
  >{\raggedright\arraybackslash}p{(\columnwidth - 14\tabcolsep) * \real{0.1017}}
  >{\raggedright\arraybackslash}p{(\columnwidth - 14\tabcolsep) * \real{0.1356}}
  >{\raggedright\arraybackslash}p{(\columnwidth - 14\tabcolsep) * \real{0.2203}}
  >{\raggedright\arraybackslash}p{(\columnwidth - 14\tabcolsep) * \real{0.0847}}
  >{\raggedright\arraybackslash}p{(\columnwidth - 14\tabcolsep) * \real{0.1525}}
  >{\raggedright\arraybackslash}p{(\columnwidth - 14\tabcolsep) * \real{0.1356}}
  >{\raggedright\arraybackslash}p{(\columnwidth - 14\tabcolsep) * \real{0.1017}}@{}}
\toprule\noalign{}
\begin{minipage}[b]{\linewidth}\raggedright
ID
\end{minipage} & \begin{minipage}[b]{\linewidth}\raggedright
Name
\end{minipage} & \begin{minipage}[b]{\linewidth}\raggedright
Sender
\end{minipage} & \begin{minipage}[b]{\linewidth}\raggedright
Receiver(s)
\end{minipage} & \begin{minipage}[b]{\linewidth}\raggedright
DLC
\end{minipage} & \begin{minipage}[b]{\linewidth}\raggedright
Payload
\end{minipage} & \begin{minipage}[b]{\linewidth}\raggedright
Period
\end{minipage} & \begin{minipage}[b]{\linewidth}\raggedright
Prio
\end{minipage} \\
\midrule\noalign{}
\endhead
\bottomrule\noalign{}
\endlastfoot
\texttt{0x001} & SAFETY\_ESTOP & Any & All (bridged to high) & 0 & (none) &
Event & Highest \\
\texttt{0x011} & SYS\_SAFETY\_STS & SYS & RT (→ Jetson) & 2 & u8 estop, u8
hb\_ok & 5 Hz & V.High \\
\texttt{0x012} & SYS\_DCDC\_CMD & SYS & DC-DC converter & 1 & u8 enable & Change
& V.High \\
\texttt{0x110} & SYS\_MODE\_CMD & SYS & RT & 1 & u8 mode (0=M, 1=A, 2=ESTOP) &
Change & High \\
\texttt{0x120} & SYS\_THROTTLE\_STS & MTR & RT (→ Jetson), SYS & 2 & i16
speed\_mmps & 100 Hz & Medium \\
\texttt{0x169} & VCU\_SES\_REQ & RT & EPS-C (steering) & 8 & Angle cmd +
security bytes & 50 Hz & Medium \\
\texttt{0x201} & SES\_STATUS & EPS-C & RT & 8 & Steering angle + status feedback
& 100 Hz & Medium \\
\texttt{0x202} & SES\_ErrInfo & EPS-C & RT & 8 & 25 fault flags (8× L3) & 10 Hz
& Medium \\
\texttt{0x203} & SES\_Version & EPS-C & RT & 8 & SW + HW version & 1 Hz &
Lowest \\
\texttt{0x204} & RT\_DRIVE\_CMD & RT & \textbf{MTR (STM32), SYS} & 5 & i32
speed\_mmps, u8 gear & 100 Hz & Medium \\
\texttt{0x205} & RT\_BRAKE\_CMD & RT & SYS & 4 & i32 brake\_pressure\_kpa & 50
Hz & Low \\
\texttt{0x206} & MTR\_MOTOR\_FBK & MTR & SYS, RT & 4 & i16 actual\_speed, u8
gear\_state, u8 fault\_flags & 50 Hz & Low \\
\texttt{0x302} & HOST\_LIGHT\_CMD & RT (fwd) & SYS & 1 & u8 lights bitfield &
Change & Medium \\
\texttt{0x600} & SYS\_DIAG\_RPT & SYS & RT (→ Jetson) & 8 & diag struct & 1 Hz &
Lowest \\
\texttt{0x6FA} & SES\_Test & EPS-C & RT & 8 & Motor current, ECU temp, supply
voltage & 100 Hz & Lowest \\
\texttt{0x6FB} & SEB\_Test & SEB & SYS & 8 & Motor current, ECU temp, supply
voltage & 100 Hz & Lowest \\
\texttt{0x721} & SEB\_STATUS & SEB & SYS & 8 & Brake stroke + status + error
level feedback & 100 Hz & Lowest \\
\texttt{0x731} & SEB\_ErrInfo & SEB & SYS & 8 & 23 fault flags (16× L3) & 10 Hz
& Lowest \\
\texttt{0x741} & SEB\_Version & SEB & SYS & 8 & SW + HW version & 1 Hz &
Lowest \\
\texttt{0x7B9} & VCU\_SEB\_REQ & RT (AUTO) / SYS (MANUAL, ESTOP) † & SEB (brake)
& 8 & Stroke/pressure cmd + security & 50 Hz & Lowest \\
\texttt{0x7FD} & RT\_HEARTBEAT & RT & SYS & 1 & u8 alive\_ctr & 2 Hz & Lowest \\
\texttt{0x7FE} & SYS\_HEARTBEAT & SYS & RT & 1 & u8 alive\_ctr & 10 Hz &
Lowest \\
\end{longtable}

\begin{quote}
\textbf{ID note}: SYNTREE units are preprogrammed and cannot be reconfigured.
EPS-C uses factory command \texttt{0x169} and status \texttt{0x201}, plus
diagnostic frames \texttt{0x202} (err info), \texttt{0x203} (version),
\texttt{0x6FA} (telemetry). SEB uses factory command \texttt{0x7B9} and status
\texttt{0x721}, plus diagnostic frames \texttt{0x731} (err info), \texttt{0x741}
(version), \texttt{0x6FB} (telemetry). \texttt{RT\_DRIVE\_CMD} is placed at
\texttt{0x204} to avoid collision with EPS-C \texttt{0x169}. \texttt{0x205}
RT\_BRAKE\_CMD avoids collision with \texttt{0x202} and \texttt{0x203}.

\textbf{† \texttt{0x7B9} dual-sender note:} Architecture §6.2 Option D specifies
RT sends \texttt{0x7B9} directly in AUTO (1-hop from kinematics), SYS in
MANUAL/ESTOP. \textbf{Current implementation (v0.0.4):} SYS is the sole
\texttt{0x7B9} sender in all modes --- RT sends \texttt{0x205\ RT\_BRAKE\_CMD}
to SYS, which converts to SEB protocol. Direct RT→SEB transmission is planned
(gap \#13).
\end{quote}

\subsection{2.2 High-level CAN}\label{high-level-can}

\begin{longtable}[]{@{}
  >{\raggedright\arraybackslash}p{(\columnwidth - 14\tabcolsep) * \real{0.0678}}
  >{\raggedright\arraybackslash}p{(\columnwidth - 14\tabcolsep) * \real{0.1017}}
  >{\raggedright\arraybackslash}p{(\columnwidth - 14\tabcolsep) * \real{0.1356}}
  >{\raggedright\arraybackslash}p{(\columnwidth - 14\tabcolsep) * \real{0.2203}}
  >{\raggedright\arraybackslash}p{(\columnwidth - 14\tabcolsep) * \real{0.0847}}
  >{\raggedright\arraybackslash}p{(\columnwidth - 14\tabcolsep) * \real{0.1525}}
  >{\raggedright\arraybackslash}p{(\columnwidth - 14\tabcolsep) * \real{0.1356}}
  >{\raggedright\arraybackslash}p{(\columnwidth - 14\tabcolsep) * \real{0.1017}}@{}}
\toprule\noalign{}
\begin{minipage}[b]{\linewidth}\raggedright
ID
\end{minipage} & \begin{minipage}[b]{\linewidth}\raggedright
Name
\end{minipage} & \begin{minipage}[b]{\linewidth}\raggedright
Sender
\end{minipage} & \begin{minipage}[b]{\linewidth}\raggedright
Receiver(s)
\end{minipage} & \begin{minipage}[b]{\linewidth}\raggedright
DLC
\end{minipage} & \begin{minipage}[b]{\linewidth}\raggedright
Payload
\end{minipage} & \begin{minipage}[b]{\linewidth}\raggedright
Period
\end{minipage} & \begin{minipage}[b]{\linewidth}\raggedright
Prio
\end{minipage} \\
\midrule\noalign{}
\endhead
\bottomrule\noalign{}
\endlastfoot
\texttt{0x001} & SAFETY\_ESTOP & RT (fwd), Jetson & Jetson, RT & 0 & (none) &
Event & Highest \\
\texttt{0x011} & SYS\_SAFETY\_STS & RT (fwd) & Jetson & 2 & u8 estop, u8 hb\_ok
& 5 Hz & V.High \\
\texttt{0x120} & SYS\_THROTTLE\_STS & RT (fwd) & Jetson & 2 & i16 speed\_mmps &
100 Hz & Medium \\
\texttt{0x210} & RT\_STATE\_RPT & RT & Jetson & 3 & u8 mode, u8 steer\_valid, u8
reversing & 10 Hz & Low \\
\texttt{0x220} & RT\_PID\_RPT & RT & Jetson & 6 & RESERVED --- future
closed-loop PID telemetry & --- (inactive) & Low \\
\texttt{0x300} & HOST\_DRIVE\_CMD & Jetson & RT & 8 & i32 speed\_mmps, i24
yaw\_rate\_mrad\_s, u8 gear & ≤100 Hz & Medium \\
\texttt{0x301} & HOST\_BRAKE\_REQ & Jetson & RT & 4 & i32 brake\_pressure\_kpa &
Demand & Medium \\
\texttt{0x302} & HOST\_LIGHT\_CMD & Jetson & RT (→ SYS) & 1 & u8 lights bitfield
& Change & Medium \\
\texttt{0x400} & HOST\_OBSTACLE\_DIST & Jetson & RT & 4 & u32 distance\_mm & 10
Hz & Low \\
\texttt{0x600} & SYS\_DIAG\_RPT & RT (fwd) & Jetson & 8 & diag struct & 1 Hz &
Lowest \\
\texttt{0x7FD} & RT\_HEARTBEAT & RT & Jetson & 1 & u8 alive\_ctr & 2 Hz &
Lowest \\
\texttt{0x7FC} & JETSON\_HEARTBEAT & Jetson & RT & 1 & u8 alive\_ctr & 2 Hz &
Lowest \\
\end{longtable}

\begin{quote}
Bit-level signal layouts in
\href{can-dictionary.md}{\texttt{can-dictionary.md}}.
\end{quote}

\subsection{2.3 RT CAN gateway --- message handling by
category}\label{rt-can-gateway-message-handling-by-category}

RT is the only dual-bus node. Every CAN message falls into exactly one of three
categories:

\subsubsection{Category 1: Transparent
forward}\label{category-1-transparent-forward}

RT copies the frame to the other bus unchanged --- same CAN ID, same payload,
same DLC. The receiving node cannot tell whether RT or the original sender
transmitted it.

\begin{longtable}[]{@{}
  >{\raggedright\arraybackslash}p{(\columnwidth - 4\tabcolsep) * \real{0.3235}}
  >{\raggedright\arraybackslash}p{(\columnwidth - 4\tabcolsep) * \real{0.4118}}
  >{\raggedright\arraybackslash}p{(\columnwidth - 4\tabcolsep) * \real{0.2647}}@{}}
\toprule\noalign{}
\begin{minipage}[b]{\linewidth}\raggedright
Direction
\end{minipage} & \begin{minipage}[b]{\linewidth}\raggedright
IDs forwarded
\end{minipage} & \begin{minipage}[b]{\linewidth}\raggedright
Example
\end{minipage} \\
\midrule\noalign{}
\endhead
\bottomrule\noalign{}
\endlastfoot
Low → High & \texttt{0x001}, \texttt{0x011}, \texttt{0x120}, \texttt{0x206},
\texttt{0x600} & SYS sends \texttt{0x011} on low → RT transmits \texttt{0x011}
on high → Jetson receives \\
High → Low & \texttt{0x001}, \texttt{0x302} & Jetson sends \texttt{0x302} on
high → RT transmits \texttt{0x302} on low → SYS receives \\
\end{longtable}

\begin{quote}
\texttt{0x001} is the only bidirectional forward --- ESTOP originates on either
bus and must reach all nodes on both buses.
\end{quote}

\subsubsection{Category 2: Consumed by RT → different message generated on other
bus}\label{category-2-consumed-by-rt-different-message-generated-on-other-bus}

RT receives a command on one bus, processes it internally, and transmits a
\textbf{different} CAN ID with \textbf{different} payload on the other bus.
These are not forwards --- they are translations.

\begin{longtable}[]{@{}
  >{\raggedright\arraybackslash}p{(\columnwidth - 8\tabcolsep) * \real{0.3175}}
  >{\raggedright\arraybackslash}p{(\columnwidth - 8\tabcolsep) * \real{0.0794}}
  >{\raggedright\arraybackslash}p{(\columnwidth - 8\tabcolsep) * \real{0.1746}}
  >{\raggedright\arraybackslash}p{(\columnwidth - 8\tabcolsep) * \real{0.3492}}
  >{\raggedright\arraybackslash}p{(\columnwidth - 8\tabcolsep) * \real{0.0794}}@{}}
\toprule\noalign{}
\begin{minipage}[b]{\linewidth}\raggedright
Inbound (consumed)
\end{minipage} & \begin{minipage}[b]{\linewidth}\raggedright
Bus
\end{minipage} & \begin{minipage}[b]{\linewidth}\raggedright
Processing
\end{minipage} & \begin{minipage}[b]{\linewidth}\raggedright
Outbound (generated)
\end{minipage} & \begin{minipage}[b]{\linewidth}\raggedright
Bus
\end{minipage} \\
\midrule\noalign{}
\endhead
\bottomrule\noalign{}
\endlastfoot
\texttt{0x300} HOST\_DRIVE\_CMD & High & Kinematics → \texttt{ResolvedSetpoint}
& \texttt{0x204} RT\_DRIVE\_CMD + \texttt{0x169} VCU\_SES\_REQ & Low \\
\texttt{0x301} HOST\_BRAKE\_REQ & High & Max-select arbitration → \texttt{0x205}
RT\_BRAKE\_CMD & \texttt{0x205} (RT→SYS, AUTO only); SYS converts to
\texttt{0x7B9} → SEB & Low \\
\end{longtable}

\subsubsection{Category 3: Bus-local (never forwarded, never
regenerated)}\label{category-3-bus-local-never-forwarded-never-regenerated}

These messages serve only nodes on a single bus. RT neither forwards nor
translates them.

\begin{longtable}[]{@{}
  >{\raggedright\arraybackslash}p{(\columnwidth - 4\tabcolsep) * \real{0.2778}}
  >{\raggedright\arraybackslash}p{(\columnwidth - 4\tabcolsep) * \real{0.2778}}
  >{\raggedright\arraybackslash}p{(\columnwidth - 4\tabcolsep) * \real{0.4444}}@{}}
\toprule\noalign{}
\begin{minipage}[b]{\linewidth}\raggedright
Bus
\end{minipage} & \begin{minipage}[b]{\linewidth}\raggedright
IDs
\end{minipage} & \begin{minipage}[b]{\linewidth}\raggedright
Reason
\end{minipage} \\
\midrule\noalign{}
\endhead
\bottomrule\noalign{}
\endlastfoot
Low only & \texttt{0x012}, \texttt{0x110}, \texttt{0x169}, \texttt{0x204},
\texttt{0x205}, \texttt{0x7B9} & DC-DC, mode, steering, drive, brake, and SEB
commands. RT-generated or SYS-generated --- never leave low bus. \\
Low only & \texttt{0x201}, \texttt{0x721} & SYNTREE status feedback --- RT/SYS
consume locally for boot sync and monitoring \\
High only & \texttt{0x210}, \texttt{0x220}, \texttt{0x400} & RT telemetry ---
generated by RT for Jetson consumption only \\
Both (independent) & \texttt{0x7FD}, \texttt{0x7FE}, \texttt{0x7FC} & Per-node
heartbeat with alive counter --- MUST NOT be bridged (see §8.6). Each bus has
its own liveness domain. \\
\end{longtable}

\begin{center}\rule{0.5\linewidth}{0.5pt}\end{center}

\section{3. Mode state machine}\label{mode-state-machine}

\begin{verbatim}
         ┌──────────┐
    ┌───▶│  MANUAL  │◀───┐
    │    └─────┬────┘    │
    │     push btn=AUTO  push btn=MANUAL
    │          │          │
    │    ┌─────▼────┐    │
    │    │   AUTO   │    │
    │    └─────┬────┘    │
    │          │          │
    │  ESTOP button / CAN 0x001 / HB timeout
    │          │          │
    │    ┌─────▼────┐    │
    │    │  ESTOP   │────┘
    │    └─────┬────┘
    │         │ START button (GPIO32)
    │         ▼
    └──────── MANUAL
\end{verbatim}

\begin{longtable}[]{@{}
  >{\raggedright\arraybackslash}p{(\columnwidth - 2\tabcolsep) * \real{0.3750}}
  >{\raggedright\arraybackslash}p{(\columnwidth - 2\tabcolsep) * \real{0.6250}}@{}}
\toprule\noalign{}
\begin{minipage}[b]{\linewidth}\raggedright
Mode
\end{minipage} & \begin{minipage}[b]{\linewidth}\raggedright
Behavior
\end{minipage} \\
\midrule\noalign{}
\endhead
\bottomrule\noalign{}
\endlastfoot
\textbf{MANUAL} & Rider steers / rides throttle. \textbf{MTR} reads throttle ADC
+ gear sense → pass-through via MCP4725 + relays. Brake lever → SYS GPIO → CAN
\texttt{0x7B9} → SEB. EPS-C standalone (RT idle). DC-DC on. Mode gated by SYS
\texttt{0x110}. \\
\textbf{AUTO} & Jetson \texttt{AckermannControlCommand} → high CAN
\texttt{0x300} → RT kinematics → low CAN \texttt{0x204} (MTR: speed+gear) +
\texttt{0x169} (EPS-C: angle). \textbf{MTR} drives MCP4725 + gear relays
following \texttt{0x204}. Lights from Jetson via \texttt{0x302} (RT fwd). Brake
via \texttt{0x7B9}. Mode gated by SYS \texttt{0x110}. \\
\textbf{ESTOP} & MCP4725 = 0 V, all gear outputs OFF, \texttt{0x7B9} stroke=max
(full brake), steering ramps to 0° at 20°/s via active \texttt{0x169} (unless
obstacle-triggered → hold then silent-stop), DC-DC ON (\texttt{0x012\ enable=1}
--- maintains 12V for MCUs, CAN transceivers, brake light), 12V accessory relay
OFF (kills headlight, turn signals, mode bulbs). Brake light ON (powered from
always-on DC-DC rail, not through accessory relay). Exit: \textbf{START button}
or \textbf{MODE long-press (3s)} → MANUAL. Steering ramp completes before 0x169
handoff (brake/motor/lights transition immediately; steering defers until
centered). Or power-cycle. \\
\end{longtable}

\begin{center}\rule{0.5\linewidth}{0.5pt}\end{center}

\section{4. Signal flow}\label{signal-flow}

\subsection{4.1 Manual mode}\label{manual-mode}

\begin{verbatim}
Throttle grip (0–5V) ──► MTR ADC ──► MTR MCP4725 (0–5V) ──► Motor controller
Gear selector (72V)  ──► TLP281 opto → MTR GPIO ──► relay module → 72V → ECU
Brake lever           ──► SYS GPIO ──► CAN 0x7B9 → SEB (stroke=MAX if pressed)
Steering wheel        ──► EPS-C standalone (RT idle, monitors 0x201)
Signal lights         ──► Turn: handlebar switches (SYS GPIO3/6). Head: toggle (SYS GPIO7). Brake: OR logic → SYS GPIO21
DC-DC converter       ──► SYS CAN 0x012 enable=1 → 12V rail on
SYS → CAN 0x110       ──► MTR receives mode=Manual → ADC pass-through
\end{verbatim}

\subsection{4.2 Auto mode}\label{auto-mode}

\begin{verbatim}
Jetson AckermannControlCommand ──► High CAN 0x300 ──► RT kinematics
                                          │
               ┌──────────────────────────┤
               ▼ (low CAN)                ▼ (low CAN)
   0x204 {speed, gear} → MTR       0x169 {angle} → EPS-C
               │                          │
               ├──► MCP4725 → Motor       │  RT listens 0x201 for feedback
               ├──► Relays → ECU gear     │  Dynamic angle clamp + following error
               └──► CAN 0x120, 0x206 →    │
                     RT (fwd to Host),     │
                     SYS (EGAS L2 monitor) │

SYS → CAN 0x110 ──► MTR receives mode=Auto → follows 0x204
Jetson ──► High CAN 0x301 ──► RT brake arbitration → Low CAN 0x7B9 → SEB (RT sends directly in AUTO)
Jetson ──► High CAN 0x302 ──► RT fwd → Low CAN 0x302 → SYS → light relays
SYS ────► Low CAN 0x7B9 → SEB (MANUAL/ESTOP only; RT sends in AUTO)
\end{verbatim}

\begin{center}\rule{0.5\linewidth}{0.5pt}\end{center}

\section{5. Responsibility split}\label{responsibility-split}

\begin{longtable}[]{@{}lllll@{}}
\toprule\noalign{}
Concern & Jetson & RT & SYS & MTR \\
\midrule\noalign{}
\endhead
\bottomrule\noalign{}
\endlastfoot
Perception / planning & ✓ & & & \\
ROS 2 → CAN bridge (\texttt{AckermannControlCommand}) & ✓ & & & \\
Autoware.Auto bridge (\texttt{autoware\_vehicle\_bridge}) & ✓ & & & \\
CAN gateway (low ↔ high) & & ✓ & & \\
Tricycle kinematics & & ✓ & & \\
Steering angle compute + CAN TX (\texttt{0x169}) & & ✓ & & \\
Steering boot sync (Listen-Before-Speaking) & & ✓ & & \\
Steering safety: dynamic angle clamp, hard-stops, following error & & ✓ & & \\
Obstacle speed limit & & ✓ & & \\
Command staleness watchdog & & ✓ & & \\
E-stop GPIO + button (wired to both) & & & ✓ & ✓ \\
Brake lever → CAN (\texttt{0x7B9}, 50 Hz continuous) & & & ✓ & \\
Brake boot sync (Listen-Before-Speaking) & & & ✓ & \\
Brake rolling counter + checksum & & & ✓ & \\
DC-DC converter CAN control (\texttt{0x012}) & & & ✓ & \\
Heartbeat monitoring & & ✓ (Jetson high + SYS low) & ✓ (RT, low) & \\
Mode switch reading & & & ✓ & \\
Throttle ADC read (0--5V) & & & & ✓ \\
Throttle MCP4725 DAC output (0--5V) & & & & ✓ \\
Gear 72V read (TLP281 opto) & & & & ✓ \\
Gear 72V output (relay module) & & & & ✓ \\
Motor feedback CAN TX (\texttt{0x206}) & & & & ✓ \\
12V accessory power relay & & & ✓ & \\
Mode indicator lights & & & ✓ & \\
Signal lights (turn, brake, head) & & & ✓ & \\
System diagnostics & & & ✓ & \\
\end{longtable}

\begin{center}\rule{0.5\linewidth}{0.5pt}\end{center}

\subsection{5.1 Jetson Orin --- Autoware.Auto ROS 2 Bridge
(v0.0.4-alpha)}\label{jetson-orin-autoware.auto-ros-2-bridge-v0.0.4-alpha}

A new ROS 2 lifecycle node (\texttt{autoware\_vehicle\_bridge}) at
\texttt{jetson/src/autoware\_vehicle\_bridge/} provides Autoware.Auto topic
compatibility. The CAN protocol on both buses is unchanged.

\textbf{Subscriptions:}
\texttt{autoware\_auto\_control\_msgs/msg/AckermannControlCommand} → CAN
\texttt{0x300}/\texttt{0x301}. \texttt{GearCommand},
\texttt{TurnIndicatorsCommand}, \texttt{HazardLightsCommand} → CAN
\texttt{0x302}. \texttt{Engage} → internal state (false suppresses commands).

\textbf{Publications:} \texttt{VelocityReport} ← CAN \texttt{0x120}.
\texttt{SteeringReport} ← CAN \texttt{0x210}. \texttt{GearReport} ← CAN
\texttt{0x210}. \texttt{ControlModeReport} ← CAN \texttt{0x210} +
\texttt{0x011}.

\textbf{Heartbeats:} Sends \texttt{0x7FC} at 2 Hz. Monitors \texttt{0x7FD}
(1500ms timeout → RCLCPP\_WARN\_THROTTLE). RT's own staleness watchdog (500ms)
stops the vehicle independently.

\textbf{Node parameters:} \texttt{wheel\_base:\ 1.5m},
\texttt{max\_speed\_forward:\ 3.0\ m/s},
\texttt{max\_steering\_angle:\ 0.698\ rad},
\texttt{max\_brake\_pressure\_kpa:\ 5000}, \texttt{loop\_rate:\ 100\ Hz},
\texttt{command\_timeout\_ms:\ 500}.

\begin{center}\rule{0.5\linewidth}{0.5pt}\end{center}

\section{6. Design principles}\label{design-principles}

\begin{enumerate}
\def\labelenumi{\arabic{enumi}.}
\tightlist
\item
  \textbf{Queues over shared state.} No mutexes, no semaphores. Thread-safe
  queue pipes.
\item
  \textbf{ESTOP bypasses queues.} Safety task preempts and writes directly to
  actuators.
\item
  \textbf{One CAN ID = one sender per bus.} Every CAN ID has exactly one
  originator on a given bus. Each node's heartbeat uses its own ID
  (\texttt{0x7FD} RT, \texttt{0x7FE} SYS, \texttt{0x7FC} Jetson).
\item
  \textbf{Lower CAN ID = higher bus priority.} Safety IDs (\texttt{0x00X}) win
  arbitration.
\item
  \textbf{All multi-byte CAN fields are big-endian (MSB first)} --- unless
  SYNTREE protocol specifies otherwise (see \texttt{0x169}, \texttt{0x7B9} in
  can-dictionary).
\item
  \textbf{Manual mode is pass-through, not dead.} SYS mirrors physical inputs to
  outputs.
\item
  \textbf{Actuators are standalone CAN modules.} EPS-C, SEB, and DC-DC are
  commanded via CAN.
\item
  \textbf{RT is the only dual-bus node.} No direct Jetson ↔ SYS path.
\item
  \textbf{Listen Before Speaking.} SYNTREE units require receiving status
  feedback before any command is sent. Boot state machines enforce this.
\item
  \textbf{EGAS 3-level safety separation for motor actuation.} The motor
  controller takes raw analog signals (0--5V throttle, 72V gear) with no
  internal intelligence. This is the only actuator without built-in CAN
  monitoring. Per ISO 26262, it is isolated on a dedicated STM32 (MTR) with
  three independent safety levels.
\item
  \textbf{Mode-gated dual control of SYNTREE actuators (Option D).} In AUTO, RT
  commands both EPS-C and SEB directly --- 1-hop from the kinematics engine. In
  MANUAL, SYS commands SEB based on lever position; EPS-C runs standalone. The
  \texttt{0x7B9} SEB command is dual-sender but mode-gated --- only one node
  transmits at a time, no collision. Per ISO 26262-5:2018 §7.4.4, this is
  redundant-controller practice. No single MCU failure can take both actuators
  offline.
\end{enumerate}

\begin{center}\rule{0.5\linewidth}{0.5pt}\end{center}

\subsection{6.1 EGAS 3-Level Motor Safety
Architecture}\label{egas-3-level-motor-safety-architecture}

The separation of \textbf{MTR (STM32)} from \textbf{SYS (ESP32-S3)} follows the
EGAS 3-level electronic throttle monitoring concept (ISO 26262 ASIL-C):

\begin{verbatim}
Level 3: Hardware — ESTOP button wired direct to both MTR and SYS
         TPS3850 external watchdog on each MCU. No software, no CAN.
         ESTOP press → MTR cuts throttle + gear instantly (local).
         
Level 2: Function Monitor — SYS ESP32-S3
         Monitors MTR via CAN: compares 0x204 setpoint vs 0x206 feedback.
         Mismatch > threshold → CAN 0x001 ESTOP.
         Also handles QM body functions (lights, DCDC, indicators).
         
Level 1: Function Controller — MTR STM32
         Normal actuation: reads sensors, drives MCP4725 DAC + gear relays.
         MANUAL: pass-through from grip/gear. AUTO: follows CAN 0x204.
         No wireless, no OS, minimal attack surface.
\end{verbatim}

\begin{longtable}[]{@{}
  >{\raggedright\arraybackslash}p{(\columnwidth - 2\tabcolsep) * \real{0.4231}}
  >{\raggedright\arraybackslash}p{(\columnwidth - 2\tabcolsep) * \real{0.5769}}@{}}
\toprule\noalign{}
\begin{minipage}[b]{\linewidth}\raggedright
Principle
\end{minipage} & \begin{minipage}[b]{\linewidth}\raggedright
Implementation
\end{minipage} \\
\midrule\noalign{}
\endhead
\bottomrule\noalign{}
\endlastfoot
\textbf{Freedom from interference} & Separate MCUs --- a SYS crash/hang cannot
block motor kill \\
\textbf{ASIL decomposition} & MTR (QM level sensor reads) + SYS monitor (ASIL-B
comparison) → combined ASIL-C \\
\textbf{Independent safe state path} & ESTOP button wired to both MCUs --- MTR
cuts throttle locally, zero CAN delay \\
\textbf{Diverse monitoring} & SYS compares commanded speed vs actual from 0x206
--- mismatch → ESTOP \\
\textbf{Why only motor needed MTR} & SYNTREE EPS-C and SEB already do EGAS Level
1 internally (CAN, PID, angle sensor, fault detection). Motor controller is a
dumb analog device --- it has no intelligence, no CAN, no internal monitoring.
That gap is filled by the dedicated MTR board. \\
\textbf{Body control is QM} & Lights, indicators, DCDC, mode bulbs on SYS are
non-safety --- safe to share MCU with Level 2 monitoring per ISO 26262 QM
classification \\
\end{longtable}

\subsection{6.2 Mode-Gated Dual Control (Option D) --- SYNTREE
Actuators}\label{mode-gated-dual-control-option-d-syntree-actuators}

The EPS-C and SEB are commanded by different nodes depending on mode. Four
options were evaluated (see §6.3). Option D was selected:

\textbf{Why Option D over A (distributed fixed):} In AUTO, RT computes steering
angle AND brake pressure from the same kinematics model. Option D sends both
from RT --- a single control tick produces both CAN frames. No cross-node sync
needed for obstacle response. In A, brake required RT→0x205→SYS→0x7B9 (3 hops);
in D it's RT→0x7B9 (1 hop).

\textbf{Why Option D over C (SYS owns both):} C makes SYS a single point of
failure for both actuators. D preserves independent failure modes --- RT failure
loses AUTO steering/brake but SYS can still brake in MANUAL; SYS failure loses
MANUAL brake but RT can still brake in AUTO. The EPS-C's internal timeout (20ms
comm loss → lock) and SEB's internal timeout provide hardware backup regardless.

\textbf{Why not B (private CAN):} Adds hardware cost (MCP2515 + transceiver +
bus wiring) and SYS bridging latency without safety improvement. The shared
low-level CAN with SYNTREE's own rolling counter + checksum validation already
prevents accidental actuation.

\begin{longtable}[]{@{}
  >{\raggedright\arraybackslash}p{(\columnwidth - 4\tabcolsep) * \real{0.2381}}
  >{\raggedright\arraybackslash}p{(\columnwidth - 4\tabcolsep) * \real{0.2619}}
  >{\raggedright\arraybackslash}p{(\columnwidth - 4\tabcolsep) * \real{0.5000}}@{}}
\toprule\noalign{}
\begin{minipage}[b]{\linewidth}\raggedright
Property
\end{minipage} & \begin{minipage}[b]{\linewidth}\raggedright
RT (AUTO)
\end{minipage} & \begin{minipage}[b]{\linewidth}\raggedright
SYS (MANUAL/ESTOP)
\end{minipage} \\
\midrule\noalign{}
\endhead
\bottomrule\noalign{}
\endlastfoot
EPS-C (0x169) & Sends angle from kinematics & Does NOT send --- EPS-C
standalone \\
SEB (0x7B9) & Sends arbitrated brake kPa→pressure & Sends lever stroke / ESTOP
max \\
0x201 (EPS-C status) & Monitors angle for following error & Does not monitor \\
0x721 (SEB status) & Does not monitor & Monitors stroke for boot sync \\
Mode switch & Receives 0x110 from SYS & Sends 0x110 on change \\
\end{longtable}

\begin{quote}
\textbf{Dual-sender exception}: \texttt{0x7B9} has two senders on the low-level
bus. But only one is active at a time --- mode-gated. In AUTO, RT sends
continuously at 50 Hz and SYS is silent. In MANUAL/ESTOP, SYS sends and RT is
silent. No simultaneous transmission in normal operation. The SEB accepts
whichever frame has a valid checksum and rolling counter. This is standard
automotive redundant-controller practice (ISO 26262-5:2018 §7.4.4).
\end{quote}

\subsection{6.3 Options Comparison
(Archived)}\label{options-comparison-archived}

\begin{longtable}[]{@{}
  >{\raggedright\arraybackslash}p{(\columnwidth - 8\tabcolsep) * \real{0.3793}}
  >{\raggedright\arraybackslash}p{(\columnwidth - 8\tabcolsep) * \real{0.1379}}
  >{\raggedright\arraybackslash}p{(\columnwidth - 8\tabcolsep) * \real{0.1724}}
  >{\raggedright\arraybackslash}p{(\columnwidth - 8\tabcolsep) * \real{0.1379}}
  >{\raggedright\arraybackslash}p{(\columnwidth - 8\tabcolsep) * \real{0.1724}}@{}}
\toprule\noalign{}
\begin{minipage}[b]{\linewidth}\raggedright
Criterion
\end{minipage} & \begin{minipage}[b]{\linewidth}\raggedright
A (distributed)
\end{minipage} & \begin{minipage}[b]{\linewidth}\raggedright
B (private CAN)
\end{minipage} & \begin{minipage}[b]{\linewidth}\raggedright
C (SYS only)
\end{minipage} & \begin{minipage}[b]{\linewidth}\raggedright
\textbf{D (mode-gated)}
\end{minipage} \\
\midrule\noalign{}
\endhead
\bottomrule\noalign{}
\endlastfoot
Single point takes both actuators? & No & Yes & Yes & \textbf{No} \\
AUTO brake latency & \textasciitilde30ms & \textasciitilde40ms &
\textasciitilde20ms & \textbf{\textasciitilde20ms} \\
AUTO steer latency & \textasciitilde20ms & \textasciitilde30ms &
\textasciitilde30ms & \textbf{\textasciitilde20ms} \\
SYS failure → brake lost? & Yes & Yes & Yes & \textbf{No (RT takes over)} \\
RT failure → steer lost? & Yes & Same & No (SYS bridges) & Yes (EPS-C timeout
backup) \\
Code complexity (1-5) & 3 & 4 & 4 & \textbf{3} \\
Hardware cost (1-5) & 4 & 2 & 4 & \textbf{4} \\
Safety score (1-5) & 4 & 2 & 2 & \textbf{5} \\
\textbf{Total (higher=better)} & 23 & 13 & 15 & \textbf{26} \\
\end{longtable}

\begin{center}\rule{0.5\linewidth}{0.5pt}\end{center}

\section{7. RT ESP32-S3 --- Realtime Physics, Steering \& CAN
Gateway}\label{rt-esp32-s3-realtime-physics-steering-can-gateway}

\subsection{7.1 Role}\label{role}

Converts ROS 2 motion commands (high CAN \texttt{0x300}) into: - \textbf{Speed +
gear} → low CAN \texttt{0x204} → SYS - \textbf{Steering angle} → low CAN
\texttt{0x169} → SYNTREE EPS-C

Bridges selected CAN messages (§2.3). Listens to \texttt{0x201\ SES\_STATUS} for
steering feedback and safety monitoring.

\textbf{8 FreeRTOS tasks} on ESP32-S3 @ 240 MHz, 1000 Hz tick.

\subsection{7.2 Dual CAN hardware}\label{dual-can-hardware}

\begin{longtable}[]{@{}
  >{\raggedright\arraybackslash}p{(\columnwidth - 8\tabcolsep) * \real{0.1087}}
  >{\raggedright\arraybackslash}p{(\columnwidth - 8\tabcolsep) * \real{0.2391}}
  >{\raggedright\arraybackslash}p{(\columnwidth - 8\tabcolsep) * \real{0.2391}}
  >{\raggedright\arraybackslash}p{(\columnwidth - 8\tabcolsep) * \real{0.1304}}
  >{\raggedright\arraybackslash}p{(\columnwidth - 8\tabcolsep) * \real{0.2826}}@{}}
\toprule\noalign{}
\begin{minipage}[b]{\linewidth}\raggedright
Bus
\end{minipage} & \begin{minipage}[b]{\linewidth}\raggedright
Controller
\end{minipage} & \begin{minipage}[b]{\linewidth}\raggedright
Interface
\end{minipage} & \begin{minipage}[b]{\linewidth}\raggedright
GPIO
\end{minipage} & \begin{minipage}[b]{\linewidth}\raggedright
Transceiver
\end{minipage} \\
\midrule\noalign{}
\endhead
\bottomrule\noalign{}
\endlastfoot
Low-level & Built-in TWAI & Direct & TX=5, RX=4 & SN65HVD230 \\
High-level & MCP2515 & SPI & SCK=36, MOSI=37, MISO=38, CS=39, INT=40 &
SN65HVD230 \\
\end{longtable}

\subsection{7.3 CAN messages received}\label{can-messages-received}

\begin{longtable}[]{@{}
  >{\raggedright\arraybackslash}p{(\columnwidth - 8\tabcolsep) * \real{0.1515}}
  >{\raggedright\arraybackslash}p{(\columnwidth - 8\tabcolsep) * \real{0.1515}}
  >{\raggedright\arraybackslash}p{(\columnwidth - 8\tabcolsep) * \real{0.1818}}
  >{\raggedright\arraybackslash}p{(\columnwidth - 8\tabcolsep) * \real{0.2727}}
  >{\raggedright\arraybackslash}p{(\columnwidth - 8\tabcolsep) * \real{0.2424}}@{}}
\toprule\noalign{}
\begin{minipage}[b]{\linewidth}\raggedright
Bus
\end{minipage} & \begin{minipage}[b]{\linewidth}\raggedright
ID
\end{minipage} & \begin{minipage}[b]{\linewidth}\raggedright
Name
\end{minipage} & \begin{minipage}[b]{\linewidth}\raggedright
Payload
\end{minipage} & \begin{minipage}[b]{\linewidth}\raggedright
Action
\end{minipage} \\
\midrule\noalign{}
\endhead
\bottomrule\noalign{}
\endlastfoot
Low & \texttt{0x001} & SAFETY\_ESTOP & --- & \texttt{mode\_set(Estop)}, forward
to high \\
Low & \texttt{0x011} & SYS\_SAFETY\_STS & \texttt{\{u8\ estop,\ u8\ hb\_ok\}} &
Forward to high \\
Low & \texttt{0x110} & SYS\_MODE\_CMD & \texttt{u8\ mode} &
\texttt{mode\_set(Manual/Auto)} \\
Low & \texttt{0x120} & SYS\_THROTTLE\_STS & \texttt{i16\ speed\_mmps} & Forward
to high \\
Low & \texttt{0x201} & SES\_STATUS &
\texttt{\{u8\ status,\ i16\ angle,\ u8\ torq,\ …\}} (8 bytes) & Steering
feedback: sync boot angle, following error check \\
Low & \texttt{0x202} & SES\_ErrInfo & 25 fault flags (8 bytes) & Log faults;
escalate any L3 flag to ESTOP via \texttt{0x001} \\
Low & \texttt{0x203} & SES\_Version & \texttt{\{u8\ sw\_ver,\ u8\ hw\_ver\}} &
Log on boot for compatibility check \\
Low & \texttt{0x206} & MTR\_MOTOR\_FBK &
\texttt{\{i16\ actual\_speed,\ u8\ gear\_state,\ u8\ fault\_flags\}} & Forward
to high; monitor motor feedback \\
Low & \texttt{0x600} & SYS\_DIAG\_RPT & 8 bytes & Forward to high \\
Low & \texttt{0x6FA} & SES\_Test &
\texttt{\{i16\ mtr\_curr,\ u16\ ecu\_temp,\ u16\ pow\_volt\}} & Monitor motor
current / ECU temp for degradation \\
Low & \texttt{0x7FE} & SYS\_HEARTBEAT & \texttt{u8\ alive\_ctr} & Feed SYS alive
counter (10 Hz); if no frame for \textgreater200ms → zero setpoints, RT takes
over brake \\
High & \texttt{0x001} & SAFETY\_ESTOP & --- & \texttt{mode\_set(Estop)}, forward
to low \\
High & \texttt{0x300} & HOST\_DRIVE\_CMD & \texttt{\{i32\ speed,\ i32\ yaw\}} &
→ \texttt{cmd\_queue} \\
High & \texttt{0x301} & HOST\_BRAKE\_REQ & \texttt{i32\ pressure\_kpa} & →
atomic store \\
High & \texttt{0x302} & HOST\_LIGHT\_CMD & \texttt{u8} bitfield & Forward to
low \\
High & \texttt{0x7FC} & JETSON\_HEARTBEAT & \texttt{u8\ alive\_ctr} & Feed
Jetson alive counter; frozen \textgreater1500ms → stale command \\
\end{longtable}

\subsection{7.4 CAN messages sent}\label{can-messages-sent}

\begin{longtable}[]{@{}
  >{\raggedright\arraybackslash}p{(\columnwidth - 8\tabcolsep) * \real{0.1613}}
  >{\raggedright\arraybackslash}p{(\columnwidth - 8\tabcolsep) * \real{0.1613}}
  >{\raggedright\arraybackslash}p{(\columnwidth - 8\tabcolsep) * \real{0.1935}}
  >{\raggedright\arraybackslash}p{(\columnwidth - 8\tabcolsep) * \real{0.2903}}
  >{\raggedright\arraybackslash}p{(\columnwidth - 8\tabcolsep) * \real{0.1935}}@{}}
\toprule\noalign{}
\begin{minipage}[b]{\linewidth}\raggedright
Bus
\end{minipage} & \begin{minipage}[b]{\linewidth}\raggedright
ID
\end{minipage} & \begin{minipage}[b]{\linewidth}\raggedright
Name
\end{minipage} & \begin{minipage}[b]{\linewidth}\raggedright
Payload
\end{minipage} & \begin{minipage}[b]{\linewidth}\raggedright
Rate
\end{minipage} \\
\midrule\noalign{}
\endhead
\bottomrule\noalign{}
\endlastfoot
Low & \texttt{0x001} & SAFETY\_ESTOP & --- & Event \\
Low & \texttt{0x204} & RT\_DRIVE\_CMD & \texttt{\{i32\ speed,\ u8\ gear\}} & 100
Hz \\
Low & \texttt{0x205} & RT\_BRAKE\_CMD & \texttt{i32\ brake\_pressure\_kpa} &
\textbf{50 Hz} \\
Low & \texttt{0x169} & VCU\_SES\_REQ &
\texttt{\{u8\ ctrl,\ i16\ angle,\ u8\ speed,\ u8\ sec,\ u8\ cnt+cksum,\ u8\ cksum\}}
(8 bytes) & \textbf{50 Hz} \\
Low & \texttt{0x302} & HOST\_LIGHT\_CMD (fwd) & \texttt{u8} bitfield & Change \\
Low & \texttt{0x7FD} & RT\_HEARTBEAT & \texttt{u8\ alive\_ctr} & 2 Hz \\
High & \texttt{0x001} & SAFETY\_ESTOP (fwd) & --- & Event \\
High & \texttt{0x011} & SYS\_SAFETY\_STS (fwd) &
\texttt{\{u8\ estop,\ u8\ hb\_ok\}} & 5 Hz \\
High & \texttt{0x120} & SYS\_THROTTLE\_STS (fwd) & \texttt{i16\ speed\_mmps} &
100 Hz \\
High & \texttt{0x210} & RT\_STATE\_RPT &
\texttt{\{u8\ mode,\ u8\ steer\_valid,\ u8\ reversing\}} & 10 Hz \\
High & \texttt{0x220} & RT\_PID\_RPT & RESERVED (inactive until encoders fitted)
& --- \\
High & \texttt{0x400} & RT\_OBSTACLE\_RPT & \texttt{u32\ distance\_mm} & 10
Hz \\
High & \texttt{0x600} & SYS\_DIAG\_RPT (fwd) & 8 bytes & 1 Hz \\
High & \texttt{0x7FD} & RT\_HEARTBEAT & \texttt{u8\ alive\_ctr} & 2 Hz \\
\end{longtable}

\subsection{7.5 Internal data types}\label{internal-data-types}

\begin{verbatim}
enum class Mode : uint8_t { Manual = 0, Auto = 1, Estop = 2 };
enum class Gear : uint8_t { N = 0, D = 1, S = 2, R = 3 };

enum class SteerState : uint8_t {
    STEER_BOOT_WAIT,          // 500ms power-on delay — do NOT transmit
    STEER_LISTEN_SYNC,        // Waiting for 0x201 SES_STATUS (angle + alignment), 5s timeout → FAULT
    STEER_ACTIVE,             // Normal operation — transmit 0x169 at 50 Hz
    ESTOP_RAMP_TO_ZERO,       // Non-obstacle ESTOP: ramp to 0° at 20°/s, continue transmitting
    ESTOP_HOLD_THEN_SILENT,   // Obstacle ESTOP: hold current angle 500ms, then silent-stop
    STEER_FAULT               // Timeout or silent-stop — stop transmitting
};

struct DriveCmd {
    int32_t speed_mmps      = 0;   // [-500, 3000]
    int32_t yaw_rate_mrad_s = 0;   // [-3000, 3000]
};

struct ResolvedSetpoint {
    int32_t motor_speed_mmps = 0;
    int32_t steer_angle_mdeg = 0;  // ±45000, +right (internal; convert to decideg for SYNTREE)
    uint8_t gear             = 0;
    bool    steer_valid      = false;
    bool    reversing        = false;
};
\end{verbatim}

\subsection{7.6 Control mechanisms}\label{control-mechanisms}

\subsubsection{Tricycle kinematics}\label{tricycle-kinematics}

\[\delta = \arctan\left(\frac{L \cdot \omega}{|v|}\right) \quad L = 1500\text{ mm}\]

\begin{verbatim}
physics_resolve(cmd):
  1. Convert mm/s→m/s, mrad/s→rad/s
  2. If |v| > 50 mm/s: δ = atan2(L·ω, |v|), steer_valid = true
     Else: δ = steer_hold · 0.8 (decay), steer_valid = false
  3. Clamp δ to ±steer_limit (dynamic, see below)
  4. Clamp v to [-500, 3000] mm/s
  5. reversing = v < 0
  6. gear: v > 0 → D, v == 0 → N, v < 0 → R
\end{verbatim}

\subsubsection{\texorpdfstring{Steering --- SYNTREE EPS-C via CAN
(\texttt{0x169})}{Steering --- SYNTREE EPS-C via CAN (0x169)}}\label{steering-syntree-eps-c-via-can-0x169}

\textbf{Boot sequence --- ``Listen Before Speaking'':}

\begin{verbatim}
State machine (steer_state_machine_loop):

STEER_BOOT_WAIT:
  - 500ms delay after power-on
  - DO NOT transmit any 0x169 frames
  - → STEER_LISTEN_SYNC

STEER_LISTEN_SYNC:
  - Wait for 0x201 SES_STATUS frame (timeout: 5s)
      - On timeout → STEER_FAULT (rider can retry: short-press START → STEER_LISTEN_SYNC)
  - Extract SES_StrAngle (int16, scale 0.1°/bit → convert to internal mdeg)
  - CRITICAL: Set active_target_angle = current_physical_angle
  - Wait for SES_INF_Angle_Status == 1 (aligned)
  - → STEER_ACTIVE

STEER_ACTIVE:
  - Transmit 0x169 at 50 Hz
  - First frame commands wheels to stay exactly where they are
  - Then follow Jetson targets with dynamic clamp
  - Monitor following error: if |cmd - actual| > threshold for > timeout → ESTOP
      - ESTOP behavior depends on trigger:
          Obstacle-triggered: hold current angle (clamped to dynamic limit for speed),
            silent-stop after 500ms. If current angle exceeds dynamic limit, ramp to
            limit at 20°/s first, then hold. Prevents rollover during cornering + braking.
          Non-obstacle (heartbeat loss, cmd stale, manual button): ramp to 0° at 20°/s,
            continue transmitting 0x169 during ramp, then hold at 0°
      - If following error persists during ESTOP centering ramp (>5° for >1s):
          fall back to silent-stop (linkage likely mechanically jammed)

STEER_FAULT:
  - Stop transmitting 0x169
  - EPS-C will timeout-fault internally
      - Short-press START button → reset to STEER_LISTEN_SYNC (manual retry)
      - Long-press START (3s) + throttle at zero → force-activate with target=0°
        (MANUAL mode only; AUTO bulb blinks 2 Hz = degraded steering, AUTO mode locked out)
\end{verbatim}

\textbf{Unit conversion} (internal mdeg ↔ SYNTREE decideg):

\begin{verbatim}
SYNTREE raw = internal_angle_mdeg / 100   (45500 mdeg → 455 raw → 45.5°)
internal_mdeg = SYNTREE raw * 100         (455 raw → 45500 mdeg)
\end{verbatim}

\textbf{SYNTREE protocol specifics:}

\begin{longtable}[]{@{}
  >{\raggedright\arraybackslash}p{(\columnwidth - 2\tabcolsep) * \real{0.6111}}
  >{\raggedright\arraybackslash}p{(\columnwidth - 2\tabcolsep) * \real{0.3889}}@{}}
\toprule\noalign{}
\begin{minipage}[b]{\linewidth}\raggedright
Parameter
\end{minipage} & \begin{minipage}[b]{\linewidth}\raggedright
Value
\end{minipage} \\
\midrule\noalign{}
\endhead
\bottomrule\noalign{}
\endlastfoot
Command ID & \texttt{0x169} (factory default --- SYNTREE preprogrammed, not
reconfigurable. \texttt{0x7B9} is sent by RT in AUTO, SYS in MANUAL/ESTOP ---
one sender active at a time, no collision. \texttt{0x205} removed (RT now sends
SEB directly)) \\
Rate & 50 Hz (20 ms period) --- continuous transmission required \\
Control mode & 1 = Angle Mode \\
Angle range & ±700 raw (±70.0°, unit limit; software clamp tighter) \\
Angle resolution & 0.1°/bit (int16) \\
Slew rate & \texttt{VCU\_SES\_Tgt\_StrAngleSpd} {[}°/s{]} --- speed-dependent \\
Rolling counter & 4-bit, increment 0→15 every frame \\
Checksum & XOR of bytes 0--6, then \texttt{\^{}\ 0xFF} (verify against spec) \\
Security enables & Byte 5: \texttt{roll\_cnt\_enable=1},
\texttt{checksum\_enable=1} \\
\end{longtable}

\textbf{Safety mechanisms:}

\begin{longtable}[]{@{}
  >{\raggedright\arraybackslash}p{(\columnwidth - 2\tabcolsep) * \real{0.4583}}
  >{\raggedright\arraybackslash}p{(\columnwidth - 2\tabcolsep) * \real{0.5417}}@{}}
\toprule\noalign{}
\begin{minipage}[b]{\linewidth}\raggedright
Mechanism
\end{minipage} & \begin{minipage}[b]{\linewidth}\raggedright
Description
\end{minipage} \\
\midrule\noalign{}
\endhead
\bottomrule\noalign{}
\endlastfoot
\textbf{Software hard-stops} & Clamp commanded angle to ±40° (inside physical
end-stops). Reject any Jetson command exceeding this regardless of unit's ±78°
capability. \\
\textbf{Dynamic angle clamp} & Max allowable angle inversely proportional to
\texttt{RT\_PidMeasured} speed. At 25 km/h → max \textasciitilde5°. At 2 km/h →
max \textasciitilde40°. Prevents rollover. \\
\textbf{Following error} & Compare commanded angle (\texttt{0x169}) vs actual
(\texttt{SES\_StrAngle} from \texttt{0x201}). If abs(error) \textgreater{} 5°
for \textgreater{} 300 ms → trigger ESTOP (stuck linkage / rock jam). \\
\textbf{Timeout fault} & If \texttt{0x169} stops for \textgreater20 ms, EPS-C
triggers internal comm fault. RT must maintain 50 Hz in AUTO. \\
\textbf{Alignment check} & \texttt{SES\_INF\_Angle\_Status} must be 1 before
AUTO mode engages. Drive motor locked out until aligned. \\
\end{longtable}

\textbf{Mode behavior:}

\begin{longtable}[]{@{}
  >{\raggedright\arraybackslash}p{(\columnwidth - 2\tabcolsep) * \real{0.2500}}
  >{\raggedright\arraybackslash}p{(\columnwidth - 2\tabcolsep) * \real{0.7500}}@{}}
\toprule\noalign{}
\begin{minipage}[b]{\linewidth}\raggedright
Mode
\end{minipage} & \begin{minipage}[b]{\linewidth}\raggedright
Steering behavior
\end{minipage} \\
\midrule\noalign{}
\endhead
\bottomrule\noalign{}
\endlastfoot
MANUAL & RT does NOT send \texttt{0x169}. EPS-C standalone. RT still listens
\texttt{0x201} for telemetry. \\
AUTO & RT sends \texttt{0x169} at 50 Hz with resolved angle, dynamic clamp +
slew rate applied. \\
ESTOP & Obstacle: hold current angle clamped to dynamic limit (ramp to limit if
exceeding), silent-stop after 500ms. Non-obstacle: ramp to 0° at 20°/s via
active \texttt{0x169}. Fall back to silent-stop if following error persists
\textgreater1s. Both paths non-interruptible by START button (steering
transition defers until ramp/hold complete). \\
\end{longtable}

\subsubsection{Speed control --- open-loop today,
PID-ready}\label{speed-control-open-loop-today-pid-ready}

Motor speed control is \textbf{open-loop}: \texttt{0x204} desired speed → MTR
MCP4725 DAC = \texttt{speed/3000\ ×\ 4095} (fixed voltage mapping). No feedback
compensation. Acceptable for flat-ground steady-state operation but will not
compensate for hills, headwinds, or load changes.

\textbf{Why open-loop?} PID gains are defined in config.h (Kp=1.0, Ki=0.1,
Kd=0.05) but the implementation is deferred. The rear motor encoder hasn't been
physically fitted to the trike yet. The PCNT encoder inputs (GPIO1/2) currently
read zero. Running PID against \texttt{measured=0} would command full throttle
--- it MUST NOT be active until the encoder is wired.

\textbf{Integration plan (gap \#5):} Once the rear motor encoder is fitted: 1.
\texttt{control\_task} on RT feeds
\texttt{speed\_pid\_update(desired,\ measured,\ dt)} with real encoder data 2.
PID output (effort correction) is sent to SYS via an added field in
\texttt{0x204} or a new CAN ID 3. MTR adds the PID trim to the base MCP4725
voltage --- closing the loop 4. Until then, the PID runs as a \textbf{shadow
controller} in RT, outputting to \texttt{0x220\ RT\_PID\_RPT} (telemetry only)
for validation against the open-loop command --- the two should track closely on
flat ground and diverge under load, telling us the PID gains are reasonable
before wiring it in.

\subsubsection{Obstacle speed limit}\label{obstacle-speed-limit}

Jetson perception (LiDAR/camera/stereo) provides minimum obstacle distance via
CAN \texttt{0x400} at 10 Hz. RT applies linear clamp: 300 mm→0, 3000 mm→full
speed. Jetson's own MPPI planner produces safe speed targets; RT's limiter is a
safety backstop, not the primary avoidance mechanism.

\subsubsection{Command staleness watchdog}\label{command-staleness-watchdog}

500 ms, checked @ 10 Hz. Zero \texttt{0x204} + stop \texttt{0x169}.

\subsubsection{Brake arbitration
(max-select)}\label{brake-arbitration-max-select}

\begin{verbatim}
brake_kpa = max(rt_computed_obstacle, jetson_request_0x301)
send 0x205 {brake_kpa} at 50 Hz on low bus → SYS brake_task
\end{verbatim}

RT computes obstacle-emergency brake pressure from distance sensor. Jetson sends
deceleration requests via \texttt{0x301}. The worse (higher) pressure wins.
Result goes to SYS via \texttt{0x205\ RT\_BRAKE\_CMD}.

SYS forwards to SEB in Pressure Mode when \texttt{0x205\ \textgreater{}\ 0},
falling back to Stroke Mode for lever/ESTOP binary triggers.

\subsection{7.7 RTOS task layout}\label{rtos-task-layout}

\begin{verbatim}
Pri 5  can_rx_low   ── TWAI → can_rx_low_queue (16)
      can_rx_high  ── MCP2515 SPI → can_rx_high_queue (16)

Pri 4  dispatch     ◀── both RX queues
           Routes: high 0x300→cmd_queue, 0x301→atomic, 0x302→gw_tx_low
                   low 0x011→gw_tx_high, 0x120→gw_tx_high, 0x600→gw_tx_high
                   low 0x201→steer_feedback (sync angle, following error)
                   any 0x001→mode_set(Estop)+gateway, low 0x110→mode_set

Pri 4  control      ◀── cmd_queue (4, overwrite)
           100 Hz: kinematics  + dynamic angle clamp + obstacle + brake → setpoint_queue

Pri 3  can_tx_low   ◀── setpoint_queue + gw_tx_low_queue
           → 0x204 (100 Hz), 0x205 (50 Hz, drive setpoint), 0x169 (50 Hz, steer state machine), 0x302 (change)
      can_tx_high  ◀── telemetry + gw_tx_high_queue → 0x011,0x120,0x210,0x220,0x400,0x600

Pri 1  watchdog     ── 10 Hz staleness check
      heartbeat    ── 2 Hz 0x7FD on both buses
\end{verbatim}

\begin{longtable}[]{@{}
  >{\raggedright\arraybackslash}p{(\columnwidth - 8\tabcolsep) * \real{0.1622}}
  >{\raggedright\arraybackslash}p{(\columnwidth - 8\tabcolsep) * \real{0.1622}}
  >{\raggedright\arraybackslash}p{(\columnwidth - 8\tabcolsep) * \real{0.1892}}
  >{\raggedright\arraybackslash}p{(\columnwidth - 8\tabcolsep) * \real{0.2162}}
  >{\raggedright\arraybackslash}p{(\columnwidth - 8\tabcolsep) * \real{0.2703}}@{}}
\toprule\noalign{}
\begin{minipage}[b]{\linewidth}\raggedright
Task
\end{minipage} & \begin{minipage}[b]{\linewidth}\raggedright
Prio
\end{minipage} & \begin{minipage}[b]{\linewidth}\raggedright
Stack
\end{minipage} & \begin{minipage}[b]{\linewidth}\raggedright
Period
\end{minipage} & \begin{minipage}[b]{\linewidth}\raggedright
Behavior
\end{minipage} \\
\midrule\noalign{}
\endhead
\bottomrule\noalign{}
\endlastfoot
\texttt{can\_rx\_low} & 5 & 4096 B & Event & \texttt{twai\_receive()} → queue \\
\texttt{can\_rx\_high} & 5 & 4096 B & Event & MCP2515 SPI → queue \\
\texttt{dispatch} & 4 & 4096 B & Event & Route both RX queues + gateway + steer
feedback \\
\texttt{control} & 4 & 4096 B & 100 Hz & Kinematics, dynamic angle clamp,
obstacle, brake arbitration, gear derivation \\
\texttt{can\_tx\_low} & 3 & 3072 B & Event & 0x204@100Hz, 0x169@50Hz (steer SM
gated), 0x302 \\
\texttt{can\_tx\_high} & 3 & 3072 B & Event & Telemetry → MCP2515 SPI \\
\texttt{watchdog} & 1 & 2048 B & 10 Hz & Staleness → zero setpoints + stop
steer \\
\texttt{heartbeat} & 1 & 2048 B & 2 Hz & \texttt{0x7FD} on both buses (per-bus,
not bridged): \texttt{alive\_ctr++\ \&\ 0xFF}, DLC=1 \\
\end{longtable}

\subsection{7.8 Hardware pin assignments --- Board: RT
ESP32-S3}\label{hardware-pin-assignments-board-rt-esp32-s3}

\begin{quote}
⚠️ \textbf{Board identity:} This is the \textbf{RT} ESP32-S3 (realtime physics,
steering, CAN gateway). Do not confuse with SYS ESP32-S3 pin assignments in
§8.8. Same GPIO numbers on different boards are different physical pins ---
connect to the board labeled ``RT,'' not the one labeled ``SYS.''
\end{quote}

\begin{longtable}[]{@{}
  >{\raggedright\arraybackslash}p{(\columnwidth - 6\tabcolsep) * \real{0.2500}}
  >{\raggedright\arraybackslash}p{(\columnwidth - 6\tabcolsep) * \real{0.1875}}
  >{\raggedright\arraybackslash}p{(\columnwidth - 6\tabcolsep) * \real{0.3438}}
  >{\raggedright\arraybackslash}p{(\columnwidth - 6\tabcolsep) * \real{0.2188}}@{}}
\toprule\noalign{}
\begin{minipage}[b]{\linewidth}\raggedright
Signal
\end{minipage} & \begin{minipage}[b]{\linewidth}\raggedright
GPIO
\end{minipage} & \begin{minipage}[b]{\linewidth}\raggedright
Direction
\end{minipage} & \begin{minipage}[b]{\linewidth}\raggedright
Notes
\end{minipage} \\
\midrule\noalign{}
\endhead
\bottomrule\noalign{}
\endlastfoot
CAN TX (low) & 5 & Out & SN65HVD230 \\
CAN RX (low) & 4 & In & SN65HVD230 \\
SPI SCK & 36 & Out & MCP2515 (high CAN) \\
SPI MOSI & 37 & Out & MCP2515 \\
SPI MISO & 38 & In & MCP2515 \\
SPI CS & 39 & Out & MCP2515 \\
MCP INT & 40 & In & MCP2515 interrupt \\
Encoder A (rear motor) & 1 & In & Speed feedback (PCNT), quadrature \\
Encoder B (rear motor) & 2 & In & \\
Encoder A (front wheel) & 3 & In & Speed/angle feedback (PCNT), quadrature ---
\textbf{sensor TBD} \\
Encoder B (front wheel) & 6 & In & \\
Encoder A (rear left wheel) & 9 & In & Differential speed feedback (PCNT),
quadrature --- \textbf{sensor TBD} \\
Encoder B (rear left wheel) & 12 & In & \\
Encoder A (rear right wheel) & 13 & In & Differential speed feedback (PCNT),
quadrature --- \textbf{sensor TBD} \\
Encoder B (rear right wheel) & 14 & In & \\
I2C SDA & 10 & I/O & IMU (optional) \\
I2C SCL & 11 & Out & IMU (optional) \\
WDT toggle & 21 & Out & External watchdog IC (TPS3850). Toggled by
\texttt{control\_task} at 100 Hz. \\
\end{longtable}

\subsection{7.9 Configuration constants}\label{configuration-constants}

\begin{verbatim}
namespace rt {
// Vehicle
constexpr float kWheelbaseMM = 1500.0f;
// Steering (SYNTREE EPS-C)
constexpr float kSteerHardLimitDeg = 40.0f;       // software hard-stop inside mechanical limit
constexpr float kSteerFollowingErrMinDeg = 2.0f;   // floor threshold (was fixed 5.0)
constexpr float kSteerFollowingErrFactor = 0.25f;  // × dynamic_limit → threshold
constexpr int   kSteerFollowingErrMs = 300;        // duration before ESTOP
constexpr int   kSteerCmdRateHz = 50;              // SYNTREE requires 20ms period
constexpr int   kSteerBootWaitMs = 500;            // power-on delay before listening
// Dynamic angle clamp (0.0.4): limit_deg = 40.0 − (speed_kmh − 2.0) × (35.0/23.0), clamped [5.0, 40.0]
constexpr float kAngleClampBaseDeg    = 40.0f;   // max at 2 km/h
constexpr float kAngleClampMinDeg     =  5.0f;   // min at ≥25 km/h
constexpr float kAngleClampRangeDeg   = 35.0f;   // base − min
constexpr float kAngleClampSpeedRange = 23.0f;   // 25 − 2 km/h
// Speed limits
constexpr int kMaxSpeedFwdMmps = 3000, kMaxSpeedRevMmps = 500;
constexpr int kLowSpeedThreshMmps = 50;
// PID (deferred — see gap #5; gains TBD once encoders fitted)
// Obstacle
constexpr unsigned kObstacleStopMM = 300, kObstacleClearMM = 3000;
// Timing
constexpr int kControlLoopHz = 100, kCmdStaleTimeoutMs = 500;
constexpr int kHeartbeatIntervalMs = 500;
constexpr int kHeartbeatTimeoutMsSys = 200;       // SYS heartbeat loss (2 missed at 10 Hz). RT takes over 0x7B9.
constexpr int kHeartbeatTimeoutMsJetson = 1500;  // Jetson heartbeat loss (3 missed frames) → assisted stop
// CAN IDs are in shared/can/can_protocol.h (namespace can):
//   kIdRtHeartbeatLow (0x7FD), kIdSysHeartbeat (0x7FE), kIdJetsonHeartbeat (0x7FC),
//   kIdRtBrakeCmd (0x205), kIdSyntreeEpsCmd (0x169), kIdSyntreeSebCmd (0x7B9), etc.
constexpr int kWdtToggleGpio = 21;
// CAN
constexpr int kCanLowBitrateHz = 500000, kCanHighBitrateHz = 500000;
// Encoders (quadrature, PCNT)
constexpr int kEncRearMotorA = 1, kEncRearMotorB = 2;
constexpr int kEncFrontWheelA = 3, kEncFrontWheelB = 6;     // sensor TBD
constexpr int kEncRearLeftA = 9, kEncRearLeftB = 12;        // sensor TBD
constexpr int kEncRearRightA = 13, kEncRearRightB = 14;     // sensor TBD
} // namespace rt
\end{verbatim}

\subsection{7.10 Error handling}\label{error-handling}

\begin{longtable}[]{@{}
  >{\raggedright\arraybackslash}p{(\columnwidth - 4\tabcolsep) * \real{0.3000}}
  >{\raggedright\arraybackslash}p{(\columnwidth - 4\tabcolsep) * \real{0.3667}}
  >{\raggedright\arraybackslash}p{(\columnwidth - 4\tabcolsep) * \real{0.3333}}@{}}
\toprule\noalign{}
\begin{minipage}[b]{\linewidth}\raggedright
Failure
\end{minipage} & \begin{minipage}[b]{\linewidth}\raggedright
Detection
\end{minipage} & \begin{minipage}[b]{\linewidth}\raggedright
Response
\end{minipage} \\
\midrule\noalign{}
\endhead
\bottomrule\noalign{}
\endlastfoot
Low CAN bus-off & TWAI TEC \textgreater{} 255 & Log, auto-recover; ESTOP if
persistent \\
High CAN bus-off & MCP2515 error flags & Log, auto-recover; zero setpoints until
restored \\
Command stale & Watchdog 500 ms & Zero \texttt{0x204} + stop \texttt{0x169} \\
Obstacle timeout & Echo \textgreater{} 30 ms & Distance = UINT32\_MAX \\
Rear motor encoder missing & Speed = 0 & open-loop - no encoder feedback (I-term
saturates) \\
Wheel encoder missing (any) & No pulses for \textgreater1s at known speed & Log
warning; differential odometry degraded. Does NOT trigger ESTOP. \\
Steering CAN TX fail & TWAI TX errors & Log, EPS-C will timeout-fault \\
Steering following error & abs(cmd − actual) \textgreater{} 5° for 300 ms &
\texttt{mode\_set(Estop)} \\
Steering sync timeout & No \texttt{0x201} within 5s of boot & → STEER\_FAULT;
rider can retry via START button short-press \\
Gateway queue full & \texttt{xQueueSend} fail & Drop (except 0x001 --- direct
TX) \\
\end{longtable}

\subsection{7.11 Startup}\label{startup}

\begin{verbatim}
1. can_low_init() → TWAI, low CAN
2. can_high_init() → SPI + MCP2515, high CAN
3. steer SM → STEER_BOOT_WAIT (500ms) → STEER_LISTEN_SYNC (await 0x201) → STEER_ACTIVE
4. watchdog_init() → timestamp
5. heartbeat_init() → alive counter
6. Create queues (6), Create 8 tasks
7. ESP_LOGI("Ready")
\end{verbatim}

\begin{center}\rule{0.5\linewidth}{0.5pt}\end{center}

\section{8. SYS ESP32-S3 --- Safety, Motor Actuation \& Body
Control}\label{sys-esp32-s3-safety-motor-actuation-body-control}

\subsection{8.1 Role}\label{role-1}

Safety (E-stop, brake lever, RT heartbeat), motor actuation (0--5V via MCP4725,
72V gear via relays), brake control (SYNTREE SEB via CAN \texttt{0x7B9}), DC-DC
converter, signal lights, mode indicators, 12V power, diagnostics.

Low-level CAN only. Jetson communication via RT.

\textbf{15 FreeRTOS tasks} on ESP32-S3 @ 240 MHz, 1000 Hz tick.

\subsection{8.2 CAN interface}\label{can-interface}

Built-in TWAI, GPIO 4/5, 500 kbit/s, SN65HVD230.

\subsection{8.3 CAN messages received}\label{can-messages-received-1}

\begin{longtable}[]{@{}
  >{\raggedright\arraybackslash}p{(\columnwidth - 8\tabcolsep) * \real{0.1143}}
  >{\raggedright\arraybackslash}p{(\columnwidth - 8\tabcolsep) * \real{0.1714}}
  >{\raggedright\arraybackslash}p{(\columnwidth - 8\tabcolsep) * \real{0.2571}}
  >{\raggedright\arraybackslash}p{(\columnwidth - 8\tabcolsep) * \real{0.2286}}
  >{\raggedright\arraybackslash}p{(\columnwidth - 8\tabcolsep) * \real{0.2286}}@{}}
\toprule\noalign{}
\begin{minipage}[b]{\linewidth}\raggedright
ID
\end{minipage} & \begin{minipage}[b]{\linewidth}\raggedright
Name
\end{minipage} & \begin{minipage}[b]{\linewidth}\raggedright
Payload
\end{minipage} & \begin{minipage}[b]{\linewidth}\raggedright
Source
\end{minipage} & \begin{minipage}[b]{\linewidth}\raggedright
Action
\end{minipage} \\
\midrule\noalign{}
\endhead
\bottomrule\noalign{}
\endlastfoot
\texttt{0x001} & SAFETY\_ESTOP & --- & RT or any & \texttt{mode\_set(Estop)} \\
\texttt{0x204} & RT\_DRIVE\_CMD & \texttt{\{i32\ speed,\ u8\ gear\}} & RT & →
\texttt{setpoint\_queue}; stale \textgreater200ms → zero speed + N \\
\texttt{0x205} & RT\_BRAKE\_CMD & \texttt{i32\ brake\_pressure\_kpa} & RT & →
\texttt{g\_brake\_pressure\_kpa} atomic; \textgreater0 → SEB Pressure Mode \\
\texttt{0x206} & MTR\_MOTOR\_FBK &
\texttt{\{i16\ actual\_speed,\ u8\ gear\_state,\ u8\ fault\_flags\}} & MTR &
EGAS L2: compare speed setpoint vs actual; mismatch → ESTOP \\
\texttt{0x302} & HOST\_LIGHT\_CMD (fwd) & \texttt{u8} bitfield & RT & →
\texttt{g\_light\_state} \\
\texttt{0x6FB} & SEB\_Test &
\texttt{\{i16\ mtr\_curr,\ u16\ ecu\_temp,\ u16\ pow\_volt\}} & SEB & Monitor
motor current / ECU temp trends for degradation early warning \\
\texttt{0x721} & SEB\_STATUS &
\texttt{\{u8\ status,\ u16\ stroke,\ u16\ angle,\ u8\ press,\ …\}} (8 bytes) &
SEB & Sync boot stroke, brake feedback, error level \\
\texttt{0x731} & SEB\_ErrInfo & 23 fault flags (8 bytes) & SEB & Log faults;
escalate any L3 flag to ESTOP via \texttt{0x001} \\
\texttt{0x741} & SEB\_Version & \texttt{\{u8\ sw\_ver,\ u8\ hw\_ver\}} & SEB &
Log on boot for compatibility check \\
\texttt{0x7FD} & RT\_HEARTBEAT & \texttt{u8\ alive\_ctr} & RT & Feed RT alive
counter; timeout \textgreater1000ms → ESTOP. Faster detection: \texttt{0x204}
staleness at 200ms → zero speed. \\
\end{longtable}

\subsection{8.4 CAN messages sent}\label{can-messages-sent-1}

\begin{longtable}[]{@{}
  >{\raggedright\arraybackslash}p{(\columnwidth - 8\tabcolsep) * \real{0.1250}}
  >{\raggedright\arraybackslash}p{(\columnwidth - 8\tabcolsep) * \real{0.1875}}
  >{\raggedright\arraybackslash}p{(\columnwidth - 8\tabcolsep) * \real{0.2812}}
  >{\raggedright\arraybackslash}p{(\columnwidth - 8\tabcolsep) * \real{0.1875}}
  >{\raggedright\arraybackslash}p{(\columnwidth - 8\tabcolsep) * \real{0.2188}}@{}}
\toprule\noalign{}
\begin{minipage}[b]{\linewidth}\raggedright
ID
\end{minipage} & \begin{minipage}[b]{\linewidth}\raggedright
Name
\end{minipage} & \begin{minipage}[b]{\linewidth}\raggedright
Payload
\end{minipage} & \begin{minipage}[b]{\linewidth}\raggedright
Rate
\end{minipage} & \begin{minipage}[b]{\linewidth}\raggedright
Notes
\end{minipage} \\
\midrule\noalign{}
\endhead
\bottomrule\noalign{}
\endlastfoot
\texttt{0x001} & SAFETY\_ESTOP & --- & Event & → All nodes (ESTOP on L3 fault,
RT heartbeat loss) \\
\texttt{0x011} & SYS\_SAFETY\_STS & \texttt{\{u8\ estop,\ u8\ hb\_ok\}} & 5 Hz &
→ RT (fwd to Jetson) \\
\texttt{0x012} & SYS\_DCDC\_CMD & \texttt{u8\ enable} & Change & → DC-DC
converter \\
\texttt{0x110} & SYS\_MODE\_CMD & \texttt{u8\ mode} & Change & → RT \\
\texttt{0x600} & SYS\_DIAG\_RPT & 8 bytes & 1 Hz & → RT (fwd to Jetson) \\
\texttt{0x7B9} & VCU\_SEB\_REQ &
\texttt{\{u8\ ctrl{[}2{]},\ u16\ stroke,\ u16\ press,\ u8\ sec,\ u8\ cksum\}} (8
bytes) & \textbf{50 Hz} & → SYNTREE SEB \\
\texttt{0x7FE} & SYS\_HEARTBEAT & \texttt{u8\ alive\_ctr} & 10 Hz & → RT \\
\end{longtable}

\subsection{8.5 Internal data types}\label{internal-data-types-1}

\begin{verbatim}
enum class SysMode : uint8_t { Manual = 0, Auto = 1, Estop = 2 };
enum class Gear : uint8_t { N = 0, D = 1, S = 2, R = 3 };

enum class BrakeState : uint8_t {
    BRAKE_BOOT_WAIT,     // 500ms — do NOT transmit
    BRAKE_LISTEN_SYNC,   // Wait for 0x721 SEB_STATUS, read current stroke
    BRAKE_ACTIVE,        // Transmit 0x7B9 at 50 Hz, synced to SEB
    BRAKE_DEGRADED       // LISTEN_SYNC timeout — transmit with lever defaults, no sync
};

struct ActuatorSetpoint {
    int32_t motor_speed_mmps = 0;
    Gear    gear             = Gear::N;
};
\end{verbatim}

\subsection{8.6 Control mechanisms}\label{control-mechanisms-1}

\subsubsection{Throttle --- MCP4725 I2C DAC
(0--5V)}\label{throttle-mcp4725-i2c-dac-05v}

\begin{longtable}[]{@{}ll@{}}
\toprule\noalign{}
Parameter & Value \\
\midrule\noalign{}
\endhead
\bottomrule\noalign{}
\endlastfoot
DAC device & MCP4725, I2C addr 0x60, SDA=GPIO15, SCL=GPIO16 \\
Resolution & 12-bit (0--4095), VCC=5V → 0--5V output (no op-amp) \\
ADC read & ADC1\_CH5, GPIO10, 12-bit, voltage divider 5V→3.3V \\
Dead zone & 200 (raw ADC) \\
Max speed & 3000 mm/s \\
Update rate & 100 Hz \\
\end{longtable}

\begin{longtable}[]{@{}ll@{}}
\toprule\noalign{}
Mode & Behavior \\
\midrule\noalign{}
\endhead
\bottomrule\noalign{}
\endlastfoot
MANUAL & ADC read → MCP4725 write (pass-through) \\
AUTO & \texttt{setpoint.speed} → \texttt{abs(speed)/3000\ ×\ 4095} → MCP4725 \\
ESTOP & MCP4725 = 0 \\
\end{longtable}

Direction via gear lines --- MCP4725 outputs 0--5V proportional to speed
magnitude only.

\subsubsection{Gear --- TLP281 input + relay output
(72V)}\label{gear-tlp281-input-relay-output-72v}

\textbf{Input} (manual sense, galvanic isolation):

\begin{longtable}[]{@{}lll@{}}
\toprule\noalign{}
Signal & GPIO & Conditioning \\
\midrule\noalign{}
\endhead
\bottomrule\noalign{}
\endlastfoot
D sense & 12 & TLP281 optoisolator ch1 (72V→3.3V) \\
S sense & 13 & TLP281 optoisolator ch2 \\
R sense & 14 & TLP281 optoisolator ch3 \\
\end{longtable}

\textbf{Output} (mimic 72V to ECU):

\begin{longtable}[]{@{}lll@{}}
\toprule\noalign{}
Signal & GPIO & Path \\
\midrule\noalign{}
\endhead
\bottomrule\noalign{}
\endlastfoot
D out & 33 & Relay ch1: GPIO→IN, 72V→1A fuse→COM→NO→ECU, TVS SMCJ90CA→GND \\
S out & 34 & Relay ch2: same path \\
R out & 35 & Relay ch3: same path \\
\end{longtable}

\begin{longtable}[]{@{}ll@{}}
\toprule\noalign{}
Mode & Behavior \\
\midrule\noalign{}
\endhead
\bottomrule\noalign{}
\endlastfoot
MANUAL & Read TLP281 → mirror to relays \\
AUTO & Gear from CAN \texttt{0x204} → energize relay \\
ESTOP & All OFF (N) \\
\end{longtable}

\subsubsection{\texorpdfstring{DC-DC converter --- CAN
\texttt{0x012}}{DC-DC converter --- CAN 0x012}}\label{dc-dc-converter-can-0x012}

\begin{longtable}[]{@{}
  >{\raggedright\arraybackslash}p{(\columnwidth - 4\tabcolsep) * \real{0.2075}}
  >{\raggedright\arraybackslash}p{(\columnwidth - 4\tabcolsep) * \real{0.2264}}
  >{\raggedright\arraybackslash}p{(\columnwidth - 4\tabcolsep) * \real{0.5660}}@{}}
\toprule\noalign{}
\begin{minipage}[b]{\linewidth}\raggedright
Condition
\end{minipage} & \begin{minipage}[b]{\linewidth}\raggedright
CAN \texttt{0x012}
\end{minipage} & \begin{minipage}[b]{\linewidth}\raggedright
12V accessory relay (GPIO27)
\end{minipage} \\
\midrule\noalign{}
\endhead
\bottomrule\noalign{}
\endlastfoot
MANUAL or AUTO & \texttt{enable\ =\ 1} & ON (all accessories powered) \\
ESTOP & \textbf{\texttt{enable\ =\ 1}} (maintains 12V for MCUs, CAN
transceivers, and brake light) & OFF (cuts headlight, turn signals, mode
bulbs) \\
\end{longtable}

Sent on state change. \textbf{DC-DC stays ON during ESTOP} --- MCUs need
12V→3.3V to run, CAN transceivers need 5V, and the brake light must illuminate
during ESTOP (fail-visible). The 12V accessory relay (GPIO27) provides the
secondary cut for non-safety loads. The brake light is wired to the always-on
DC-DC output, not through the accessory relay, so it remains powered during
ESTOP.

\begin{quote}
\textbf{Note:} \href{can-dictionary.md}{\texttt{can-dictionary.md}} §0x012 still
shows DC-DC → OFF on ESTOP (line 81). This paragraph documents the deliberate
design change: DC-DC stays ON during ESTOP to maintain MCU power and brake
light. The can-dictionary entry is outdated and should be updated to match this
section.
\end{quote}

\subsubsection{Signal lights}\label{signal-lights}

\begin{longtable}[]{@{}lll@{}}
\toprule\noalign{}
Signal & GPIO & Notes \\
\midrule\noalign{}
\endhead
\bottomrule\noalign{}
\endlastfoot
Left turn & 18 & Blink 500ms on/off while active \\
Right turn & 19 & Blink 500ms on/off while active \\
Brake light & 21 & \textbf{OR of all braking sources} (see below) \\
Headlight & 22 & On/off \\
\end{longtable}

\textbf{Brake light logic --- OR of all braking sources:}

\begin{verbatim}
brake_light_on = safety_brake_lever_pressed()   // GPIO2 — physical lever
              OR (mode == Estop)                // ESTOP — full brake
              OR g_light_state.brake_light;     // Jetson CAN 0x302 — predictive / hazard
              // Future: OR (brake_stroke_mm > 0) — SEB feedback confirms actual braking
\end{verbatim}

All four sources are local to SYS. \texttt{g\_light\_state.brake\_light} from
Jetson is a \textbf{supplemental} trigger --- useful for predictive illumination
(Jetson sees obstacle before pressure builds) or hazard flashing --- but can
never be the \emph{only} trigger. The physical braking state always wins.

\textbf{Manual mode --- handlebar switches:}

\begin{longtable}[]{@{}
  >{\raggedright\arraybackslash}p{(\columnwidth - 6\tabcolsep) * \real{0.2857}}
  >{\raggedright\arraybackslash}p{(\columnwidth - 6\tabcolsep) * \real{0.2143}}
  >{\raggedright\arraybackslash}p{(\columnwidth - 6\tabcolsep) * \real{0.2143}}
  >{\raggedright\arraybackslash}p{(\columnwidth - 6\tabcolsep) * \real{0.2857}}@{}}
\toprule\noalign{}
\begin{minipage}[b]{\linewidth}\raggedright
Switch
\end{minipage} & \begin{minipage}[b]{\linewidth}\raggedright
GPIO
\end{minipage} & \begin{minipage}[b]{\linewidth}\raggedright
Type
\end{minipage} & \begin{minipage}[b]{\linewidth}\raggedright
Action
\end{minipage} \\
\midrule\noalign{}
\endhead
\bottomrule\noalign{}
\endlastfoot
Left turn & 3 & Momentary, active-low, pull-up & Press → left turn blinks (500ms
on/off). Press again → cancel. \\
Right turn & 6 & Momentary, active-low, pull-up & Press → right turn blinks.
Press again → cancel. \\
Headlight & 7 & Toggle, active-low, pull-up & Each press toggles headlight
on/off. \\
\end{longtable}

\textbf{Mode-dependent behavior:}

\begin{longtable}[]{@{}
  >{\raggedright\arraybackslash}p{(\columnwidth - 6\tabcolsep) * \real{0.1395}}
  >{\raggedright\arraybackslash}p{(\columnwidth - 6\tabcolsep) * \real{0.3023}}
  >{\raggedright\arraybackslash}p{(\columnwidth - 6\tabcolsep) * \real{0.2558}}
  >{\raggedright\arraybackslash}p{(\columnwidth - 6\tabcolsep) * \real{0.3023}}@{}}
\toprule\noalign{}
\begin{minipage}[b]{\linewidth}\raggedright
Mode
\end{minipage} & \begin{minipage}[b]{\linewidth}\raggedright
Turn signals
\end{minipage} & \begin{minipage}[b]{\linewidth}\raggedright
Headlight
\end{minipage} & \begin{minipage}[b]{\linewidth}\raggedright
Brake light
\end{minipage} \\
\midrule\noalign{}
\endhead
\bottomrule\noalign{}
\endlastfoot
MANUAL & Handlebar switches → \texttt{lights\_task} & Handlebar switch → GPIO22
& \textbf{OR logic} --- lever + Jetson bit \\
AUTO & \texttt{g\_light\_state} from CAN \texttt{0x302} &
\texttt{g\_light\_state.headlight} & \textbf{OR logic} --- lever + ESTOP +
Jetson bit \\
ESTOP & OFF & OFF & \textbf{ON} (forced, overrides all) \\
\end{longtable}

\textbf{Turn signal blink pattern} (\texttt{lights\_task}): - 500 ms ON, 500 ms
OFF, repeating while \texttt{left\_turn} or \texttt{right\_turn} is active -
Pressing the same switch again cancels (toggle behavior) - Pressing the opposite
switch cancels the current side and starts the new side - Both pressed
simultaneously → hazard flashers (both blink in sync)

\textbf{Headlight toggle} (\texttt{lights\_task}): - Each press of GPIO7 toggles
\texttt{headlight\_on\ =\ !headlight\_on} - AUTO mode: override from
\texttt{g\_light\_state.headlight} via CAN \texttt{0x302}

\subsubsection{Mode switch --- push button
toggle}\label{mode-switch-push-button-toggle}

A momentary push button on GPIO11 (active-low, internal pull-up, debounced).
Each press toggles the mode: MANUAL → AUTO → MANUAL. ESTOP cannot be exited via
the MODE button --- use the START button (GPIO32) or power-cycle.

\begin{verbatim}
mode_task @ 10 Hz:
  // Mode toggle button (GPIO11)
  read GPIO11
  if falling edge (prev_mode==HIGH, now==LOW) and debounce == 0:
      if current mode == MANUAL → mode_set(Auto)
      elif current mode == AUTO  → mode_set(Manual)
      debounce = kDebounceMs / 100

  // Start button (GPIO32) — exit ESTOP
  read GPIO32
  if falling edge (prev_start==HIGH, now==LOW) and debounce == 0:
      if current mode == ESTOP → mode_set(Manual)
      debounce = kDebounceMs / 100

  if debounce > 0: debounce--
  prev_mode = GPIO11; prev_start = GPIO32
\end{verbatim}

\begin{quote}
A toggle switch would require the rider to physically change switch position. A
push button is simpler to operate while riding --- one tap to switch modes.
\end{quote}

\subsubsection{Mode indicator bulbs}\label{mode-indicator-bulbs}

Two relay-driven bulbs (visible in sunlight, not just PCB LEDs):

\begin{longtable}[]{@{}lll@{}}
\toprule\noalign{}
Signal & GPIO & Active for \\
\midrule\noalign{}
\endhead
\bottomrule\noalign{}
\endlastfoot
AUTO bulb & 25 & AUTO mode \\
MANUAL bulb & 26 & MANUAL mode \\
(both OFF) & --- & ESTOP \\
\end{longtable}

\begin{quote}
GPIO → relay coil → bulb. Bulbs are powered from the 12V accessory rail --- they
go dark on ESTOP regardless of MCU state.
\end{quote}

\subsubsection{12V relay --- unchanged}\label{v-relay-unchanged}

GPIO27 (HIGH=ON, ESTOP→OFF).

\subsubsection{Heartbeat --- automotive liveness
supervision}\label{heartbeat-automotive-liveness-supervision}

Each node sends a \textbf{1-byte alive counter} on its own CAN ID at 2 Hz. A
frozen counter = a frozen node, even if the CAN controller is still DMA-ing from
a hardware buffer.

\textbf{Heartbeat IDs (one sender per ID):}

\begin{longtable}[]{@{}llll@{}}
\toprule\noalign{}
ID & Sender & Bus & Receiver \\
\midrule\noalign{}
\endhead
\bottomrule\noalign{}
\endlastfoot
\texttt{0x7FD} & RT & Both (per-bus, NOT bridged) & SYS (low), Jetson (high) \\
\texttt{0x7FE} & SYS & Low only & RT \\
\texttt{0x7FC} & Jetson & High only & RT \\
\end{longtable}

\textbf{Why heartbeats are NOT bridged:}

Each CAN bus is an independent liveness domain. RT sends \texttt{0x7FD}
independently on each bus (same ID, separate alive counters per bus). SYS and
Jetson heartbeats never leave their respective buses. This means every receiver
sees exactly one sender per heartbeat ID --- no ambiguity, no software demuxing
needed.

\textbf{Liveness matrix (who monitors whom):}

\begin{verbatim}
       SYS ──── low CAN ──── RT ──── high CAN ──── Jetson
              0x7FE (SYS)          0x7FC (Jetson)
              0x7FD (RT)           0x7FD (RT)

SYS watches:   RT (0x7FD on low, 1000ms timeout)
RT watches:    SYS (0x7FE on low, 200ms timeout — 2 missed frames at 10 Hz) AND Jetson (0x7FC on high, 1500ms timeout)
Jetson watches: RT (0x7FD on high, 1500ms timeout)

SYS does NOT watch Jetson — RT handles Jetson failure (zero setpoints)
Jetson does NOT watch SYS — RT handles SYS failure (0x001 ESTOP)
\end{verbatim}

\begin{longtable}[]{@{}
  >{\raggedright\arraybackslash}p{(\columnwidth - 6\tabcolsep) * \real{0.1642}}
  >{\raggedright\arraybackslash}p{(\columnwidth - 6\tabcolsep) * \real{0.2687}}
  >{\raggedright\arraybackslash}p{(\columnwidth - 6\tabcolsep) * \real{0.2836}}
  >{\raggedright\arraybackslash}p{(\columnwidth - 6\tabcolsep) * \real{0.2836}}@{}}
\toprule\noalign{}
\begin{minipage}[b]{\linewidth}\raggedright
Parameter
\end{minipage} & \begin{minipage}[b]{\linewidth}\raggedright
RT↔SYS (low bus)
\end{minipage} & \begin{minipage}[b]{\linewidth}\raggedright
RT→Jetson (high)
\end{minipage} & \begin{minipage}[b]{\linewidth}\raggedright
Jetson→RT (high)
\end{minipage} \\
\midrule\noalign{}
\endhead
\bottomrule\noalign{}
\endlastfoot
RT heartbeat ID & \texttt{0x7FD} & \texttt{0x7FD} & --- \\
Peer heartbeat ID & \texttt{0x7FE} (SYS) & --- & \texttt{0x7FC} (Jetson) \\
DLC & 1 & 1 & 1 \\
Payload & \texttt{u8\ alive\_ctr} (0--255, wraps) & \texttt{u8\ alive\_ctr} &
\texttt{u8\ alive\_ctr} \\
Period & RT: 500 ms (2 Hz) / SYS: 100 ms (10 Hz) & 500 ms (2 Hz) & 500 ms (2
Hz) \\
Timeout & \textbf{200ms} (2 missed frames) & 1500ms (3 missed frames) &
\textbf{1500ms} (3 missed frames) \\
Action on loss & \textbf{RT takes over 0x7B9 (stroke=max) + CAN 0x001.} MTR
kills motor/gear locally. Brake gap ≤220ms. & Jetson: \textbf{assisted stop} ---
zero 0x204 + stop 0x169 + 0x205 brake=2000 kPa + SYS→MANUAL. Rider can override.
& RT: zero \texttt{0x204} + stop \texttt{0x169} \\
\end{longtable}

\textbf{Why 200ms for SYS heartbeat (10 Hz):}

SYS controls the SEB brake actuator --- there is no independent fast-path check
for brake command loss (unlike steering which has the 300ms following-error
check, and motor which has the 200ms 0x204 staleness check). The heartbeat IS
the brake fast-path check. At 25 km/h, 200ms is \textasciitilde1.4m of travel
--- within FTTI for brake loss. At 10 Hz (100ms period), 2 missed frames = 200ms
worst case. This is tight enough for brake safety while providing a 100ms
debounce against CAN jitter.

\textbf{RT brake takeover on SYS heartbeat loss:} Architecture §6.2 Option D
explicitly requires RT to take over brake on SYS failure: ``SYS failure → brake
lost? No (RT takes over).'' When RT detects SYS heartbeat loss, RT immediately
begins transmitting 0x7B9 with stroke=max (full brake) at 50 Hz, regardless of
current mode. RT already knows the full SYNTREE SEB protocol --- this is the
same transmission it performs in AUTO mode, with stroke=max substituted for the
computed pressure target. The mode gate opens on emergency: both RT and SYS may
briefly transmit 0x7B9 during the takeover transition, which is within the
dual-sender exception documented in §6.2. Total brake gap: 200ms detection +
20ms first frame = 220ms worst case.

\textbf{Jetson→RT at 1500ms --- assisted stop:}

Jetson runs the perception stack (obstacle detection). When it dies, the vehicle
must decelerate actively --- pure coast is insufficient in traffic. Three missed
frames at 2 Hz = 1500ms protects against false triggers from Jetson's
non-realtime Linux CAN stack (a Linux machine that can't send a single CAN frame
in 1.5s is genuinely dead). On timeout, RT commands: zero 0x204 speed, stop
0x169 steering, 0x205 brake = 2000 kPa (\textasciitilde2 MPa, moderate
deceleration without wheel lockup), and transitions SYS to MANUAL mode. Brake
light illuminates. Rider can override with lever. This is ``assisted stop'' ---
between pure coast and full ESTOP --- providing active deceleration while
preserving lights and rider agency.

\textbf{MTR-side \texttt{0x204} staleness check (fast path):}

MTR's control loop checks data freshness independently of heartbeat. If
\texttt{0x204\ RT\_DRIVE\_CMD} stops arriving in AUTO mode (RT control loop
crashed, CAN TX failed):

\begin{verbatim}
motor_task @ 100Hz:
  if (time_since_last_0x204 > 200ms):  // 2 missed frames at 100 Hz
      setpoint.speed = 0
      setpoint.gear = N
\end{verbatim}

This is a data-quality check, not a node-liveness check. It triggers in 200ms
(vs 1000ms for heartbeat) because stale data is immediately dangerous. Combined
with the heartbeat, SYS now has two independent checks:

\begin{longtable}[]{@{}llll@{}}
\toprule\noalign{}
Check & Detects & Timeout & Action \\
\midrule\noalign{}
\endhead
\bottomrule\noalign{}
\endlastfoot
\texttt{0x204} staleness & RT control loop stopped & 200ms & Zero speed +
neutral \\
RT heartbeat loss (\texttt{0x7FD}) & RT node dead & 1000ms & ESTOP (full
stop) \\
\end{longtable}

\textbf{Startup grace period:}

At boot, heartbeats haven't been established. \texttt{safety\_heartbeat\_ok()}
returns \texttt{true} for the first \textbf{3 seconds} if
\texttt{last\_hb\_timestamp\ ==\ 0}. After the grace period, real heartbeat
checking begins. This prevents false ESTOP during boot.

\begin{verbatim}
bool safety_heartbeat_ok() {
    int64_t now = esp_timer_get_time();
    if (last_hb_rt_us == 0) {
        return (now < kStartupGracePeriodUs);  // 3_000_000
    }
    return ((now - last_hb_rt_us) / 1000) < kHeartbeatTimeoutMs;  // 1000
}
\end{verbatim}

\textbf{Alive counter validation:}

A heartbeat frame with the same counter value as the previous frame = stuck CAN
controller (MCU hung, controller still DMA-ing from buffer). Treated as a missed
heartbeat.

\begin{verbatim}
bool heartbeat_is_fresh(uint8_t new_ctr) {
    if (new_ctr != last_alive_ctr) { last_alive_ctr = new_ctr; return true; }
    return false;  // frozen — same counter twice
}
\end{verbatim}

\subsubsection{External watchdog}\label{external-watchdog}

Each ESP32 toggles a dedicated \textbf{external watchdog GPIO} every iteration
of its highest-priority task. A hardware window watchdog IC (e.g., TPS3850)
resets the MCU if toggling stops for \textgreater100 ms. On reset, all outputs
default to safe state.

\begin{longtable}[]{@{}lllll@{}}
\toprule\noalign{}
Node & GPIO & Toggled by & Period & Watchdog IC \\
\midrule\noalign{}
\endhead
\bottomrule\noalign{}
\endlastfoot
RT & \textbf{21} & \texttt{control\_task} & 100 Hz & TPS3850 or equiv, 100ms
window \\
SYS & \textbf{23} & \texttt{safety\_task} & 20 Hz & TPS3850 or equiv, 100ms
window \\
\end{longtable}

\begin{quote}
This is independent of CAN heartbeat. A hung MCU with a frozen CAN controller is
invisible to heartbeat --- but the external watchdog catches it.
\end{quote}

\subsubsection{Physical controls}\label{physical-controls}

Three buttons on the dashboard:

\begin{longtable}[]{@{}
  >{\raggedright\arraybackslash}p{(\columnwidth - 6\tabcolsep) * \real{0.2857}}
  >{\raggedright\arraybackslash}p{(\columnwidth - 6\tabcolsep) * \real{0.2143}}
  >{\raggedright\arraybackslash}p{(\columnwidth - 6\tabcolsep) * \real{0.2143}}
  >{\raggedright\arraybackslash}p{(\columnwidth - 6\tabcolsep) * \real{0.2857}}@{}}
\toprule\noalign{}
\begin{minipage}[b]{\linewidth}\raggedright
Button
\end{minipage} & \begin{minipage}[b]{\linewidth}\raggedright
GPIO
\end{minipage} & \begin{minipage}[b]{\linewidth}\raggedright
Type
\end{minipage} & \begin{minipage}[b]{\linewidth}\raggedright
Action
\end{minipage} \\
\midrule\noalign{}
\endhead
\bottomrule\noalign{}
\endlastfoot
\textbf{ESTOP} & 1 & Big red mushroom, NC, active-low, hardware ISR & Instant
ESTOP --- motor kill, brake engage, DCDC off \\
\textbf{START} & 32 & Green momentary, active-low, pull-up, debounced & Exit
ESTOP → MANUAL. No effect in AUTO/MANUAL. \\
\textbf{MODE} & 11 & Momentary, active-low, pull-up, debounced & Toggle MANUAL ↔
AUTO. Ignored in ESTOP. \\
\end{longtable}

Plus brake lever on GPIO2 (active-low, pull-up). Safety task polls ESTOP + brake
lever at 20 Hz. Mode task handles MODE + START buttons at 10 Hz.

\begin{quote}
Industrial safety pattern: separate STOP (red mushroom) and START (green)
buttons. STOP is NC (normally-closed) --- a cut wire triggers ESTOP, not a
failure-silent state.
\end{quote}

\subsubsection{\texorpdfstring{Brake --- SYNTREE SEB via CAN
(\texttt{0x7B9})}{Brake --- SYNTREE SEB via CAN (0x7B9)}}\label{brake-syntree-seb-via-can-0x7b9}

\textbf{Boot sequence --- ``Listen Before Speaking'':}

\begin{verbatim}
BRAKE_BOOT_WAIT:
  - 500ms delay after power-on
  - DO NOT transmit any 0x7B9 frames
  - → BRAKE_LISTEN_SYNC

BRAKE_LISTEN_SYNC:
  - Wait for 0x721 SEB_STATUS frame
  - Extract SEB_Stroke_Value (u16, scale 0.05, offset -30)
  - Set initial command target = current stroke (hold position)
  - Wait for SEB_Alignment_Status == 1
  - → BRAKE_ACTIVE

BRAKE_ACTIVE:
  - Transmit 0x7B9 at 50 Hz continuously
  - Rolling counter increments 0→15 every frame
  - Checksum = XOR(bytes 0–6) ^ 0xFF (verify against spec)

BRAKE_DEGRADED:
  - Entered when LISTEN_SYNC times out (no 0x721 within 2s)
  - Transmit 0x7B9 at 50 Hz with conservative defaults:
      Lever pressed (GPIO2 LOW) → stroke = kBrakeManualStroke (~15 mm)
      Lever released → stroke = 0 mm
  - Rolling counter still increments, checksum still computed
  - When first valid 0x721 arrives: sync current stroke, → BRAKE_ACTIVE
  - Diagnostic flag set in 0x600 SYS_DIAG_RPT
  - Brake lever remains functional — this is not a terminal fault state
\end{verbatim}

\textbf{SYNTREE SEB protocol:}

\begin{longtable}[]{@{}
  >{\raggedright\arraybackslash}p{(\columnwidth - 2\tabcolsep) * \real{0.6111}}
  >{\raggedright\arraybackslash}p{(\columnwidth - 2\tabcolsep) * \real{0.3889}}@{}}
\toprule\noalign{}
\begin{minipage}[b]{\linewidth}\raggedright
Parameter
\end{minipage} & \begin{minipage}[b]{\linewidth}\raggedright
Value
\end{minipage} \\
\midrule\noalign{}
\endhead
\bottomrule\noalign{}
\endlastfoot
Command ID & \texttt{0x7B9} \\
Rate & 50 Hz (20 ms) --- continuous transmission required \\
Control mode & 0 = Stroke Mode, 1 = Pressure Mode (1-bit per CSV) \\
Stroke range & -5 to 27 mm (raw: 500--1140, scale 0.05, offset -30) \\
Pressure range & 0 to 5 MPa (raw: 0--100, u8, scale 0.05 MPa/bit, offset 0) \\
Pressure conversion & \texttt{seb\_raw\ =\ kPa\ ×\ 0.02} (1 MPa = 1000 kPa; 1
bit = 0.05 MPa) \\
Rolling counter & 4-bit, increment every frame \\
Checksum & XOR of bytes 0--6, then \texttt{\^{}\ 0xFF} (verify against spec) \\
\end{longtable}

\begin{quote}
\textbf{Pressure Mode (1) --- verified kPa→raw conversion:}

RT sends \texttt{0x205\ \{i32\ brake\_pressure\_kpa\}}. SYS converts using
verified SYNTREE SEB spec: \texttt{seb\_raw\ =\ (uint8\_t)(kpa\ *\ 0.02f)}.
Scale: 0.05 MPa/bit, range 0--5 MPa (raw 0--100). At 5000 kPa (5 MPa) → raw=100.
Clamp to \texttt{kSebMaxPressureRaw} (100).

\textbf{Mode-switching:} 0x205 transitions 0→positive: hold current stroke,
switch mode bit to 1 (Pressure), ramp pressure from current measured to target.
0x205 drops to 0: switch mode to 1 (Stroke), stroke=0. Prevents pressure
transients.
\end{quote}

\textbf{Mode-dependent behavior:}

\begin{longtable}[]{@{}
  >{\raggedright\arraybackslash}p{(\columnwidth - 2\tabcolsep) * \real{0.2857}}
  >{\raggedright\arraybackslash}p{(\columnwidth - 2\tabcolsep) * \real{0.7143}}@{}}
\toprule\noalign{}
\begin{minipage}[b]{\linewidth}\raggedright
Mode
\end{minipage} & \begin{minipage}[b]{\linewidth}\raggedright
Brake behavior
\end{minipage} \\
\midrule\noalign{}
\endhead
\bottomrule\noalign{}
\endlastfoot
MANUAL & Brake lever GPIO2 LOW → stroke = \texttt{kBrakeManualStroke}
(\textasciitilde15 mm). Released → stroke = 0 mm. Transmit at 50 Hz in Stroke
Mode. \\
AUTO & Lever pressed → \texttt{kBrakeManualStroke} (driver override always
wins). No lever + \texttt{0x205\ \textgreater{}\ 0} → Pressure Mode with kPa→MPa
mapping. No lever + \texttt{0x205\ ==\ 0} → stroke = 0. ESTOP → max. \\
ESTOP & Stroke = \texttt{kBrakeMaxStroke} (full brake, \textasciitilde27 mm).
Transmit at 50 Hz in Stroke Mode. \\
\end{longtable}

\textbf{Stroke value calculation:}

\begin{verbatim}
Physical stroke [mm] → raw = (physical + 30.0) / 0.05
Example: 0 mm → (0+30)/0.05 = 600
         15 mm → (15+30)/0.05 = 900
         27 mm → (27+30)/0.05 = 1140
\end{verbatim}

\subsection{8.7 RTOS task layout}\label{rtos-task-layout-1}

\begin{verbatim}
Pri 5  can_rx      ── TWAI → can_rx_queue (16)
       safety      ── GPIO poll @ 20 Hz → ESTOP / HB check

Pri 4  dispatch    ◀── can_rx_queue: 0x204→setpoint, 0x302→light, 0x001→ESTOP, 0x721→brake_feedback
       mode        ── Push button (GPIO11) @ 10 Hz → toggle MANUAL↔AUTO, CAN 0x110
       motor       ◀── setpoint_queue (4, overwrite)
             100 Hz: AUTO→MCP4725+gear, MANUAL→pass-through, ESTOP→all off

Pri 3  throttle    ── ADC @ 100 Hz → CAN 0x120
       gear        ── Gear FSM @ 50 Hz
       brake       ── Brake FSM @ 50 Hz → CAN 0x7B9 (continuous, rolling ctr + cksum)
       lights      ── Light FSM @ 20 Hz (blink timing, ESTOP=brake ON)
       dcdc        ── DCDC FSM @ 5 Hz → CAN 0x012

Pri 2  indicator   ── Mode bulbs @ 5 Hz
       power       ── 12V relay @ 5 Hz
       can_tx      ── Safety status @ 5 Hz → CAN 0x011

Pri 1  diag        ── System health @ 1 Hz → CAN 0x600
       hb          ── 0x7FE @ 10 Hz
\end{verbatim}

\begin{longtable}[]{@{}
  >{\raggedright\arraybackslash}p{(\columnwidth - 8\tabcolsep) * \real{0.1622}}
  >{\raggedright\arraybackslash}p{(\columnwidth - 8\tabcolsep) * \real{0.1622}}
  >{\raggedright\arraybackslash}p{(\columnwidth - 8\tabcolsep) * \real{0.1892}}
  >{\raggedright\arraybackslash}p{(\columnwidth - 8\tabcolsep) * \real{0.2162}}
  >{\raggedright\arraybackslash}p{(\columnwidth - 8\tabcolsep) * \real{0.2703}}@{}}
\toprule\noalign{}
\begin{minipage}[b]{\linewidth}\raggedright
Task
\end{minipage} & \begin{minipage}[b]{\linewidth}\raggedright
Prio
\end{minipage} & \begin{minipage}[b]{\linewidth}\raggedright
Stack
\end{minipage} & \begin{minipage}[b]{\linewidth}\raggedright
Period
\end{minipage} & \begin{minipage}[b]{\linewidth}\raggedright
Behavior
\end{minipage} \\
\midrule\noalign{}
\endhead
\bottomrule\noalign{}
\endlastfoot
\texttt{can\_rx} & 5 & 4096 B & Event & \texttt{twai\_receive()}, copy to
queue \\
\texttt{safety} & 5 & 2048 B & 20 Hz & ESTOP GPIO, RT HB timeout, EGAS L2:
compare 0x204 vs 0x206 \\
\texttt{dispatch} & 4 & 3072 B & Event & Route 0x206, 0x302, 0x001, 0x721 \\
\texttt{mode} & 4 & 2048 B & 10 Hz & MODE btn toggle + START btn (ESTOP→MANUAL),
CAN 0x110 → RT + MTR \\
\texttt{brake} & 3 & 2048 B & \textbf{50 Hz} & Brake SM (boot sync) + CAN 0x7B9
with rolling ctr + checksum \\
\texttt{lights} & 3 & 1536 B & 20 Hz & Light GPIOs + blink \\
\texttt{dcdc} & 3 & 1024 B & 5 Hz & DCDC FSM, CAN 0x012 \\
\texttt{indicator} & 2 & 1024 B & 5 Hz & Mode bulbs (AUTO/MANUAL) \\
\texttt{power} & 2 & 1024 B & 5 Hz & 12V relay \\
\texttt{can\_tx} & 2 & 3072 B & 5 Hz & CAN \texttt{0x011} (mode task sends
\texttt{0x110}) \\
\texttt{diag} & 1 & 2048 B & 1 Hz & CAN 0x600 \\
\texttt{hb} & 1 & 2048 B & 10 Hz & CAN \texttt{0x7FE} (SYS heartbeat to RT, fast
path for brake loss detection) \\
\end{longtable}

\begin{quote}
\textbf{Note:} Motor actuation (MCP4725 DAC), throttle ADC, and gear relay
control run on SYS in MANUAL/AUTO modes. MTR STM32 handles only the bare-metal
motor controller interface (EGAS L1). In ESTOP, SYS zeros the DAC and disengages
all relays.
\end{quote}

\subsection{8.8 Hardware pin assignments --- Board: SYS
ESP32-S3}\label{hardware-pin-assignments-board-sys-esp32-s3}

\begin{quote}
⚠️ \textbf{Board identity:} This is the \textbf{SYS} ESP32-S3 (safety, brake
control, body control). Motor actuation is on the dedicated \textbf{MTR} STM32
board (§5.0). Do not confuse with RT ESP32-S3 pin assignments in §7.8. Same GPIO
numbers on different boards are different physical pins --- connect to the board
labeled ``SYS,'' not the one labeled ``RT.''

\textbf{Shared GPIO numbers (safe --- different chips):} ESTOP button is shared
between SYS GPIO1 and MTR (hardwired to both MCUs for Level 3 kill). No other
signals are split. GPIO 3,6,7 appear in both RT and SYS tables but are
physically separate pins on separate ESP32-S3 packages --- no electrical
conflict. Throttle ADC, gear sense, MCP4725 I2C, and gear outputs are on MTR
only (§5.0).
\end{quote}

\begin{longtable}[]{@{}
  >{\raggedright\arraybackslash}p{(\columnwidth - 6\tabcolsep) * \real{0.2105}}
  >{\raggedright\arraybackslash}p{(\columnwidth - 6\tabcolsep) * \real{0.1579}}
  >{\raggedright\arraybackslash}p{(\columnwidth - 6\tabcolsep) * \real{0.2895}}
  >{\raggedright\arraybackslash}p{(\columnwidth - 6\tabcolsep) * \real{0.3421}}@{}}
\toprule\noalign{}
\begin{minipage}[b]{\linewidth}\raggedright
Signal
\end{minipage} & \begin{minipage}[b]{\linewidth}\raggedright
GPIO
\end{minipage} & \begin{minipage}[b]{\linewidth}\raggedright
Direction
\end{minipage} & \begin{minipage}[b]{\linewidth}\raggedright
Conditioning
\end{minipage} \\
\midrule\noalign{}
\endhead
\bottomrule\noalign{}
\endlastfoot
CAN TX (low) & 5 & Out & SN65HVD230 \\
CAN RX (low) & 4 & In & SN65HVD230 \\
E-stop button & 1 & In & Big red mushroom, NC (active-low), pull-up, hardware
ISR. Shared with MTR. \\
Brake lever & 2 & In & Active-low, pull-up → CAN 0x7B9 → SEB \\
Start button & \textbf{32} & In & Momentary, active-low, pull-up, debounced.
Exits ESTOP → MANUAL. \\
Mode button & 11 & In & Push button, active-low, pull-up, debounced (momentary
toggle MANUAL↔AUTO). Publishes CAN 0x110 to RT + MTR. \\
Left turn switch & \textbf{3} & In & Handlebar switch, active-low, pull-up \\
Right turn switch & \textbf{6} & In & Handlebar switch, active-low, pull-up \\
Headlight switch & \textbf{7} & In & Handlebar toggle, active-low, pull-up \\
Left turn & 18 & Out & Relay → lamp \\
Right turn & 19 & Out & Relay → lamp \\
Brake light & 21 & Out & Relay → lamp \\
Headlight & 22 & Out & Relay → lamp \\
AUTO bulb & 25 & Out & Relay → bulb (12V rail) \\
MANUAL bulb & 26 & Out & Relay → bulb (12V rail) \\
12V relay & 27 & Out & Secondary cut on ESTOP \\
WDT toggle & \textbf{23} & Out & External watchdog IC (TPS3850). Toggled by
\texttt{safety\_task} every 50 ms. \\
\end{longtable}

\begin{quote}
\textbf{Motor I/O (throttle, gear, MCP4725) is on MTR STM32 only} --- see §5.0
for MTR pin assignments.
\end{quote}

\subsection{8.9 Configuration constants}\label{configuration-constants-1}

\begin{verbatim}
namespace sys {
// CAN
constexpr int kCanBitrateHz = 500000, kCanTxGpio = 5, kCanRxGpio = 4;
// Safety
constexpr int kEstopGpio = 1, kBrakeLeverGpio = 2, kModeSwitchGpio = 11;
constexpr int kWdtToggleGpio = 23;
// Throttle/gear I/O is on MTR (mtr-stm32/src/config.h), not SYS
// Light inputs (handlebar switches, MANUAL mode)
constexpr int kSwitchLeftTurn = 3, kSwitchRightTurn = 6, kSwitchHeadlight = 7;
// Light outputs
constexpr int kLightLeftTurn = 18, kLightRightTurn = 19;
constexpr int kLightBrake = 21, kLightHead = 22;
// Indicators & power
constexpr int kBulbAuto = 25, kBulbManual = 26, kPower12vRelay = 27;
constexpr int kModeBtnGpio = 11, kStartBtnGpio = 32;
constexpr int kDebounceMs = 500;         // push button debounce period
// Turn blink
constexpr int kTurnBlinkOnMs = 500, kTurnBlinkOffMs = 500;
// Timing
constexpr int kControlLoopHz = 100;
constexpr int kHeartbeatIntervalMs    = 100;   // SYS sends 0x7FE at 10 Hz (fast path for brake loss detection)
constexpr int kHeartbeatTimeoutMsRt   = 1000;  // RT heartbeat loss (0x7FD at 2 Hz, 2 missed frames = 1000ms). Faster 0x204 staleness (200ms) catches RT crash first.
// CAN IDs are in shared/can/can_protocol.h (namespace can):
//   kIdRtHeartbeatLow (0x7FD), kIdSysHeartbeat (0x7FE), kIdRtBrakeCmd (0x205),
//   kIdSyntreeSebCmd (0x7B9), etc.
// Shared constants are in shared/shared_config.h (namespace shared):
//   kSebMaxPressureRaw (100), kStartupGracePeriodMs (3000),
//   kBrakeStrokeScale (0.05), kBrakeStrokeOffset (-30.0)
constexpr int kSetpointStaleMs = 200;           // 0x204 staleness → zero speed
constexpr int kSafetyCheckHz = 20, kGearCheckHz = 50;
// Brake (SYNTREE SEB)
constexpr int kBrakeCmdRateHz = 50, kBrakeBootWaitMs = 500;
constexpr float kBrakeManualStroke = 15.0f, kBrakeMaxStroke = 27.0f;
} // namespace sys
\end{verbatim}

\subsection{8.10 Error handling}\label{error-handling-1}

\begin{longtable}[]{@{}
  >{\raggedright\arraybackslash}p{(\columnwidth - 4\tabcolsep) * \real{0.3000}}
  >{\raggedright\arraybackslash}p{(\columnwidth - 4\tabcolsep) * \real{0.3667}}
  >{\raggedright\arraybackslash}p{(\columnwidth - 4\tabcolsep) * \real{0.3333}}@{}}
\toprule\noalign{}
\begin{minipage}[b]{\linewidth}\raggedright
Failure
\end{minipage} & \begin{minipage}[b]{\linewidth}\raggedright
Detection
\end{minipage} & \begin{minipage}[b]{\linewidth}\raggedright
Response
\end{minipage} \\
\midrule\noalign{}
\endhead
\bottomrule\noalign{}
\endlastfoot
E-stop pressed & GPIO1 LOW & ESTOP → DAC=0, gears off, brake=max, \textbf{DCDC
on} (maintains 12V for MCUs, CAN transceivers, brake light), 12V accessory relay
off \\
CAN bus-off & TWAI TEC \textgreater{} 255 & Log, auto-recover \\
RT HB timeout & \texttt{0x7FD} alive ctr frozen \textgreater1000ms & ESTOP (AUTO
only) \\
\texttt{0x204} stale & No \texttt{0x204} for \textgreater200ms & Zero speed + N
gear (controlled stop) \\
Brake lever & GPIO2 LOW & Engage brake \\
ADC fail & \texttt{adc1\_get\_raw()==0} & Throttle = 0 \\
Gear sense conflict & Multiple lines HIGH & Treat as N (fail-safe) \\
DCDC CAN TX fail & TWAI TX errors & 12V relay backup cut \\
SEB sync timeout & No \texttt{0x721} within 2s of boot &
\texttt{BRAKE\_DEGRADED} --- transmits 0x7B9 with lever-based defaults (0mm
released, 15mm pressed). Recovers when 0x721 arrives. Lever always
functional. \\
SEB checksum fail & SEB rejects frame & Frame dropped; counter still
increments \\
SEB stroke following error & abs(0x7B9 cmd − 0x721 actual) \textgreater{} 3mm
for \textgreater100ms & Log brake fault to 0x600 diag. Set persistent fault
flag. Cannot escalate (ESTOP is already max brake) but critical for
post-incident analysis. \\
SEB Level 3 fault & \texttt{SEB\_Error\_Status\ ≥\ 3} in 0x721 or 0x731 & Log
SEB severe fault to 0x600 diag. Alert rider via brake light flash pattern (if
12V available). \\
0x721 staleness & No 0x721 for \textgreater100ms (10 missed at 100 Hz) & SEB
comms lost --- log fault. If sustained, treat as brake system offline. \\
MTR ESTOP ACK timeout & \texttt{ESTOP\_ACTIVE} bit not set in 0x206 fault\_flags
within 100ms of ESTOP & Log MTR failure to acknowledge. If 0x206 also stale, MTR
comms lost. \\
External WDT timeout & TPS3850 MR pin & MCU hardware reset → all outputs safe
state \\
Queue full & \texttt{xQueueSend} fail & Frame dropped \\
\end{longtable}

\subsection{8.11 Startup}\label{startup-1}

\begin{verbatim}
 1. can_driver_init()     → TWAI, low-level CAN
 2. mode_manager_init()   → GPIO11 (MODE), GPIO32 (START)
 3. safety_monitor_init() → GPIO1 (ESTOP), GPIO2 (brake lever), WDT GPIO23
 4. throttle_init()       → ADC1_CH5 + I2C + MCP4725 (output=0)
 5. gear_init()           → GPIO12-14 (IN), GPIO33-35 (OUT, LOW)
 6. lights_init()         → GPIO3,6,7 (IN, switches) + GPIO18-22,25-26 (OUT)
 7. power_init()          → GPIO27 (OUT, LOW)
 8. brake_init()          → BRAKE_BOOT_WAIT (500ms) → LISTEN_SYNC (await 0x721) → ACTIVE
 9. dcdc_init()           → CAN 0x012 enable=0
10. Create queues         → can_rx(16), setpoint(4)
11. Create 15 tasks
12. power_task → 12V relay ON (if not ESTOP)
13. dcdc_task → CAN 0x012 enable=1 (if not ESTOP)
14. safety_task starts WDT toggle → external watchdog armed
15. ESP_LOGI("Ready")
\end{verbatim}

\begin{center}\rule{0.5\linewidth}{0.5pt}\end{center}

\section{9. CAN bus device maps}\label{can-bus-device-maps}

\subsection{Low-level}\label{low-level}

\begin{verbatim}
 Low-Level CAN (500 kbit/s)
  ├── RT ESP32-S3 (TWAI)        TX: 0x169,0x204,0x205,0x302,0x001,0x7FD
  │                              RX: 0x001,0x011,0x110,0x120,0x201,0x202,0x203,0x206,0x600,0x6FA,0x7FD,0x7FE
  ├── SYS ESP32-S3               TX: 0x011,0x012,0x110,0x600,0x7B9,0x001,0x7FE
  │                              RX: 0x001,0x204,0x205,0x206,0x302,0x6FB,0x721,0x731,0x741,0x7FD
  ├── SYNTREE EPS-C (steering)   TX: 0x201 | RX: 0x169
  ├── SYNTREE SEB (brake)        TX: 0x721 | RX: 0x7B9
  └── DC-DC converter (72→12V)  RX: 0x012
\end{verbatim}

\subsection{High-level}\label{high-level}

\begin{verbatim}
 High-Level CAN (500 kbit/s)
  ├── Jetson Orin             TX: 0x300,0x301,0x302,0x001,0x7FC
  │                              RX: 0x001,0x011,0x120,0x210,0x220,0x400,0x600,0x7FD
  └── RT ESP32-S3 (MCP2515)      TX: 0x011,0x120,0x210,0x220,0x400,0x600,0x001,0x7FD
                                  RX: 0x001,0x300,0x301,0x302,0x7FC
\end{verbatim}

\begin{center}\rule{0.5\linewidth}{0.5pt}\end{center}

\section{10. Hardware summary}\label{hardware-summary}

\begin{longtable}[]{@{}
  >{\raggedright\arraybackslash}p{(\columnwidth - 8\tabcolsep) * \real{0.1250}}
  >{\raggedright\arraybackslash}p{(\columnwidth - 8\tabcolsep) * \real{0.2292}}
  >{\raggedright\arraybackslash}p{(\columnwidth - 8\tabcolsep) * \real{0.1042}}
  >{\raggedright\arraybackslash}p{(\columnwidth - 8\tabcolsep) * \real{0.2292}}
  >{\raggedright\arraybackslash}p{(\columnwidth - 8\tabcolsep) * \real{0.3125}}@{}}
\toprule\noalign{}
\begin{minipage}[b]{\linewidth}\raggedright
Node
\end{minipage} & \begin{minipage}[b]{\linewidth}\raggedright
Controller
\end{minipage} & \begin{minipage}[b]{\linewidth}\raggedright
MCU
\end{minipage} & \begin{minipage}[b]{\linewidth}\raggedright
Framework
\end{minipage} & \begin{minipage}[b]{\linewidth}\raggedright
CAN interfaces
\end{minipage} \\
\midrule\noalign{}
\endhead
\bottomrule\noalign{}
\endlastfoot
Jetson & Orin & --- & ROS 2 & 1× CAN (high) \\
RT & ESP32-S3 @ 240 MHz & Xtensa LX7 & ESP-IDF + FreeRTOS & TWAI (low) + MCP2515
SPI (high) \\
SYS & ESP32-S3 @ 240 MHz & Xtensa LX7 & ESP-IDF + FreeRTOS & TWAI (low only) \\
\end{longtable}

\begin{longtable}[]{@{}ll@{}}
\toprule\noalign{}
Parameter & Value \\
\midrule\noalign{}
\endhead
\bottomrule\noalign{}
\endlastfoot
CAN bitrate (both) & 500 kbit/s \\
CAN transceiver & SN65HVD230 \\
FreeRTOS tick & 1000 Hz \\
RT tasks & 8 \\
SYS tasks & 15 \\
\end{longtable}

\begin{center}\rule{0.5\linewidth}{0.5pt}\end{center}

\section{11. Build}\label{build}

\begin{verbatim}
cd rt-esp32 && pio run && pio run -t upload && pio device monitor
cd sys-esp32 && pio run && pio run -t upload && pio device monitor
\end{verbatim}

\begin{center}\rule{0.5\linewidth}{0.5pt}\end{center}

\section{12. Known design gaps}\label{known-design-gaps}

\begin{longtable}[]{@{}
  >{\raggedright\arraybackslash}p{(\columnwidth - 6\tabcolsep) * \real{0.1071}}
  >{\raggedright\arraybackslash}p{(\columnwidth - 6\tabcolsep) * \real{0.1786}}
  >{\raggedright\arraybackslash}p{(\columnwidth - 6\tabcolsep) * \real{0.2857}}
  >{\raggedright\arraybackslash}p{(\columnwidth - 6\tabcolsep) * \real{0.4286}}@{}}
\toprule\noalign{}
\begin{minipage}[b]{\linewidth}\raggedright
\#
\end{minipage} & \begin{minipage}[b]{\linewidth}\raggedright
Gap
\end{minipage} & \begin{minipage}[b]{\linewidth}\raggedright
Impact
\end{minipage} & \begin{minipage}[b]{\linewidth}\raggedright
Resolution
\end{minipage} \\
\midrule\noalign{}
\endhead
\bottomrule\noalign{}
\endlastfoot
1 & \st{Brake arbitration has no CAN path to SYS} & \st{Jetson
\mbox{\texttt{0x301}} + RT obstacle braking never actuated} & \textbf{RESOLVED:}
\texttt{0x205\ RT\_BRAKE\_CMD} (RT→SYS, DLC=4, i32 kPa, 50 Hz). SYS maps kPa→SEB
Pressure Mode. Mode-switching protocol defined in §8.6. \\
2 & \st{No CAN message for Jetson to request S (Sport) gear} & \st{AUTO can only
select D/N/R} & \textbf{RESOLVED:} \texttt{0x300\ HOST\_DRIVE\_CMD} repacked:
i32 speed + i24 yaw + u8 gear (N=0,D=1,S=2,R=3). RT passes gear through to
\texttt{0x204}. \\
3 & \st{EPS-C timeout-fault behavior unknown} & \st{On ESTOP or comm loss,
steering may lock, center, or freewheel} & \textbf{RESOLVED:} Two-tier ESTOP
steering (§7.6). Active-zero centering for non-obstacle ESTOP with silent-stop
fallback on mechanical jam. Obstacle-triggered ESTOP holds current angle. EPS-C
timeout-fault is now a last-resort fallback (CAN bus dead), not the primary
ESTOP mechanism. \\
4 & \st{SEB pressure control mode not defined} & \st{SYS currently uses stroke
mode only} & \textbf{RESOLVED:} Verified SYNTREE SEB spec:
\texttt{VCU\_SEB\_Pre\_Value\_Req} is u8 at bit 32, scale 0.05 MPa/bit, range
0--5 MPa (raw 0--100). Conversion: \texttt{seb\_raw\ =\ kPa\ ×\ 0.02}. Clamp to
100. \\
5 & Speed control open-loop --- PID exists but encoder not fitted &
\texttt{speed\_pid.cpp} is correct but runs with \texttt{measured=0} (encoder
not physically installed). PID must NOT be in the active control path until
GPIO1/2 PCNT is wired. Currently runs as shadow controller → \texttt{0x220}
telemetry for validation. & 1) Fit rear motor encoder to GPIO1/2. 2) Verify
encoder pulses on PCNT. 3) Enable \texttt{pid\_update(desired,\ measured,\ dt)}
with real data. 4) Route PID effort trim to SYS via new field in \texttt{0x204}
or dedicated ID. 5) Validate closed-loop response against open-loop baseline. \\
6 & \textbf{ESTOP exit race condition --- START button during steering ramp} &
Pressing START before non-obstacle centering ramp completes stops 0x169
mid-ramp, leaving EPS-C to comm-fault at an off-center angle. A 30° steering
lock in MANUAL mode (where EPS-C is standalone with no active centering) makes
the vehicle unrideable. & \textbf{Software fix:} RT defers ESTOP→MANUAL steering
transition until ramp completes. Brake/motor/lights transition immediately. No
new CAN messages needed. See \texttt{docs/emergency-safety-analysis.md} §1. \\
7 & \textbf{SEB BRAKE\_FAULT makes physical brake lever inoperative} & Original
BRAKE\_FAULT state permanently stops 0x7B9 transmission on boot sync timeout.
Since SEB is by-wire (electro-hydraulic, no mechanical linkage), this disables
the brake lever --- contradicting the claim ``brake lever always works.'' &
\textbf{Software fix:} Replaced BRAKE\_FAULT with BRAKE\_DEGRADED (§8.6). On
sync timeout, transmit 0x7B9 with lever-based defaults (0mm/15mm) without
waiting for sync. Recovers when 0x721 arrives. Lever always functional. See
\texttt{docs/emergency-safety-analysis.md} §2. \\
8 & \textbf{Watchdog reset unbraked window --- SEB comm-fault behavior
unverified} & When SYS watchdog resets, SEB enters comm-fault after 20ms. If SEB
releases on timeout, the vehicle coasts without brake for \textasciitilde2.5s
(SYS reboot + brake LBS). If SEB holds, the window is only 20ms. Behavior is
empirically unverified. & \textbf{Test SEB comm-fault behavior} (stroke=27mm,
stop CAN, measure pressure over 5s). If release: add NC brake-hold relay gated
by TPS3850 RST line. If hold: document as verified safety property. \textbf{Also
mitigated by gap \#11 --- RT brake takeover closes most of the window.} See
\texttt{docs/emergency-safety-analysis.md} §3. \\
9 & \textbf{Obstacle ESTOP ``hold angle'' can cause rollover during cornering} &
Obstacle-triggered ESTOP holds current steering angle regardless of vehicle
speed. A cornering vehicle under hard braking experiences lateral load transfer
that the dynamic angle clamp was designed to prevent --- but the clamp is only
applied to commanded angles, not to the ESTOP hold angle. Physics model §8
defines the rollover threshold: a\_y = v²/L·tan(δ) \textgreater{} g·w/(2h). &
\textbf{Software fix:} During obstacle-triggered ESTOP, clamp the hold angle to
the dynamic angle clamp limit for the current speed. If current angle exceeds
the limit, ramp down to the limit at 20°/s. Straight-line cases unchanged. See
\texttt{docs/emergency-safety-analysis.md} §4. \\
10 & \textbf{Jetson heartbeat loss → pure coast is insufficient for perception
failure} & Jetson runs obstacle detection. When it dies, the vehicle coasts with
no active brake, no perception, and no steering. At 25 km/h, \textgreater50m
coast-to-stop. The rider must recognize the failure and manually brake. The
heartbeat timeout (1500ms) is long enough that false positives from Linux CAN
jitter are unlikely. & \textbf{Software fix:} On Jetson heartbeat loss, RT
commands 0x205 = 2000 kPa (\textasciitilde2 MPa moderate brake) + transitions
SYS to MANUAL. Brake light ON. Rider can override with lever. DC-DC stays on
(lights work). This is ``assisted stop'' --- between coast and full ESTOP. See
\texttt{docs/emergency-safety-analysis.md} §5. \\
11 & \textbf{No secondary ESTOP exit path --- START button is a single point of
failure} & GPIO32 is the only ESTOP exit. Stuck HIGH (broken wire) = can never
exit ESTOP. The only backup is power-cycle, which restarts all nodes and
requires brake/steering LBS at roadside. & \textbf{Software fix:} MODE button
(GPIO11) long-press (3s) exits ESTOP → MANUAL as secondary path. Two independent
GPIOs on separate physical buttons. Add START button health monitoring in
diag\_task. See \texttt{docs/emergency-safety-analysis.md} §6. \\
12 & \textbf{SYS crash --- 1000ms brake gap before RT responds, and RT doesn't
take over 0x7B9} & Architecture §6.2 Option D explicitly claims ``SYS failure →
brake lost? No (RT takes over).'' But RT only broadcasts CAN 0x001 on SYS
heartbeat loss --- which SEB ignores (it only responds to 0x7B9). The 1000ms
heartbeat timeout creates a \textasciitilde980ms brake gap. The heartbeat
justification (1000ms acceptable because fast-path checks cover actuators) omits
the brake --- there is no fast-path check for SEB command loss. &
\textbf{Software fix:} (1) Increase SYS heartbeat rate from 2 Hz to 10 Hz; RT
timeout shortened to 200ms (2 missed frames at 100ms period). (2) On SYS
heartbeat loss, RT immediately takes over 0x7B9 transmission with stroke=max,
regardless of current mode (mode-gate opens on emergency). Brake gap: 220ms
worst case. \textbf{Also mitigates gap \#8} by reducing watchdog brake window
from \textasciitilde2.5s to \textasciitilde220ms. See
\texttt{docs/emergency-safety-analysis.md} §7. \\
13 & \textbf{No independent brake monitor --- ESTOP could silently have no
brakes} & All 8 safety layers protect motor+steering; none verify SEB actually
applied braking force. If SEB's CAN receiver is faulted, ESTOP
\texttt{0x7B9\ stroke=max} is never received and the system never knows.
SEB\_STATUS (\texttt{0x721}, 100 Hz) already provides SEB\_Stroke\_Value,
SEB\_Pressure\_Value, and SEB\_Error\_Status (Level 3 = severe/shutdown). This
is sufficient for a brake following-error monitor with zero new sensors. &
\textbf{Software fix:} Add brake following-error check in SYS \texttt{dispatch}:
compare 0x7B9 cmd stroke vs 0x721 actual stroke (threshold 3mm, debounce 100ms).
Monitor SEB\_Error\_Status for Level 3 faults. Add 0x721 staleness check (100ms
timeout). Faults logged via 0x600 diagnostic --- cannot escalate beyond ESTOP
(already max brake) but essential for incident analysis. See
\texttt{docs/emergency-safety-analysis.md} §8. \\
14 & \textbf{CAN 0x001 spoofable --- no authentication, DoS-vulnerable} & 0x001
is DLC=0 with no sender ID, checksum, or rolling counter. A corrupted node can
flood 0x001, triggering persistent ESTOP and saturating the CAN bus. CAN error
containment (bus-off at TEC\textgreater255) eventually isolates the node, but
the flood window is disruptive. & \textbf{Software fix:} (1) Rate-limit 0x001
processing: max 2 frames per 500ms window per node (the first one already
triggered ESTOP). (2) Change 0x001 DLC from 0 to 1 with sender-ID byte (0=SYS,
1=RT, 2=Jetson, 3=MTR) for diagnostics and per-sender rate limiting. See
\texttt{docs/emergency-safety-analysis.md} §9. \\
15 & \textbf{No ESTOP acknowledgment from MTR STM32} & SYS/RT have no
confirmation that MTR received and acted on ESTOP. MTR has no dedicated
heartbeat. If MTR freezes at zero output, the SYS monitor sees 0x204=0, 0x206=0
and assumes MTR responded correctly. No 0x206 staleness check exists. &
\textbf{Software fix:} MTR sets \texttt{ESTOP\_ACTIVE} bit in 0x206 fault\_flags
when it has locally cut throttle+gear. SYS checks for ACK within 100ms of ESTOP.
Add 0x206 staleness check (\textgreater200ms → MTR comms lost). See
\texttt{docs/emergency-safety-analysis.md} §10. \\
16 & \textbf{Startup grace period masks heartbeat but not 0x204 staleness} & On
cold boot, 0x204 is absent for \textgreater200ms before RT comes online. SYS
0x204 staleness check triggers (harmlessly --- forces speed=0 which is already
default). But the startup grace concept is applied inconsistently: heartbeat
checks are masked, staleness is not. & \textbf{Trivial fix:} Gate 0x204
staleness check with same 3-second startup grace period as heartbeat. See
\texttt{docs/emergency-safety-analysis.md} §11. \\
17 & \textbf{ESTOP HMI ambiguous --- ``both bulbs OFF'' identical to powered-off
vehicle} & During ESTOP, DC-DC is OFF → 12V rail dead → brake light, mode bulbs,
and indicators all dark. A rider returning to the vehicle cannot distinguish
ESTOP from power-off. The OR logic claim ``brake light ON during ESTOP'' is
physically impossible with DC-DC off. & \textbf{Software + wiring fix:} Keep
DC-DC ON during ESTOP (needed for MCU power anyway). Cut only the 12V accessory
relay (GPIO27). Rewire brake light to always-on DC-DC rail (not accessory relay
output). Result: ESTOP = brake light ON + mode bulbs OFF. Power-off = everything
OFF. See \texttt{docs/emergency-safety-analysis.md} §12. \\
18 & \textbf{EPS-C mechanical jam silent-stop recovery path unclear} & When
mechanical jam triggers silent-stop during ESTOP centering, the architecture
says ``fall back to silent-stop'' without specifying steer SM state transition.
It's unclear whether the jam is recoverable via START short-press (STEER\_FAULT
path) or requires power-cycle. & \textbf{Documentation fix:} Explicitly
transition to STEER\_FAULT on mechanical jam during ESTOP centering. Existing
STEER\_FAULT recovery paths apply: START short-press → LISTEN\_SYNC retry; START
long-press 3s + throttle=0 → force-activate at 0° (MANUAL only). See
\texttt{docs/emergency-safety-analysis.md} §13. \\
19 & \textbf{Fixed 5° steering following error threshold wrong at high speed} &
At 25 km/h, dynamic clamp limits steering to \textasciitilde5°. A fixed 5°
threshold represents a 100\% error --- the EPS-C must have ZERO response to
trigger. A 4° error at speed (80\% authority loss) would NOT trigger ESTOP. At 2
km/h (40° limit), 5° is only 12.5\% --- potentially too sensitive for parking
maneuvers. & \textbf{Software fix:} Speed-scaled threshold:
\texttt{max(2°,\ 0.25\ ×\ dynamic\_limit)}. Result: 2° at 25 km/h (tight), 4.5°
at 10 km/h, 10° at 2 km/h (tolerant). One-line change in RT following-error
check. See \texttt{docs/emergency-safety-analysis.md} §14. \\
\end{longtable}

\begin{center}\rule{0.5\linewidth}{0.5pt}\end{center}

\section{13. Reference documents}\label{reference-documents}

\begin{longtable}[]{@{}
  >{\raggedright\arraybackslash}p{(\columnwidth - 2\tabcolsep) * \real{0.4000}}
  >{\raggedright\arraybackslash}p{(\columnwidth - 2\tabcolsep) * \real{0.6000}}@{}}
\toprule\noalign{}
\begin{minipage}[b]{\linewidth}\raggedright
File
\end{minipage} & \begin{minipage}[b]{\linewidth}\raggedright
Content
\end{minipage} \\
\midrule\noalign{}
\endhead
\bottomrule\noalign{}
\endlastfoot
\href{can-dictionary.md}{\texttt{can-dictionary.md}} & Bit-level CAN signal
layouts for all IDs on both buses \\
\href{docs/steering-unit.md}{\texttt{docs/steering-unit.md}} & SYNTREE EPS-C
protocol reference \\
\href{docs/brake-unit.md}{\texttt{docs/brake-unit.md}} & SYNTREE SEB protocol
reference \\
\href{rt-esp32/README.md}{\texttt{rt-esp32/README.md}} & RT build \& test \\
\href{sys-esp32/README.md}{\texttt{sys-esp32/README.md}} & SYS build \& test \\
\href{notes/can-protocol.md}{\texttt{notes/can-protocol.md}} & CAN protocol
theory --- arbitration, frame types, standards \\
\href{notes/can-hardware-basics.md}{\texttt{notes/can-hardware-basics.md}} & CAN
physical layer --- termination, topology, transceivers \\
\href{legacy/notes/can-addressing-for-etrike.md}{\texttt{legacy/notes/can-addressing-for-etrike.md}}
& DEPRECATED --- CAN addressing scheme (single-ESP32 design, historical
reference) \\
\href{notes/can-troubleshooting.md}{\texttt{notes/can-troubleshooting.md}} & CAN
debugging --- common mistakes, error states, tools \\
\href{notes/rtos-task-design.md}{\texttt{notes/rtos-task-design.md}} & RTOS task
design --- priorities, preemptive scheduling, queues, tick rate \\
\href{notes/pid-control.md}{\texttt{notes/pid-control.md}} & PID control
fundamentals --- P/I/D terms, integral windup, tuning procedure \\
\href{notes/state-machine-design.md}{\texttt{notes/state-machine-design.md}} &
State machine design --- FSMs, hierarchical states, C++ implementation
patterns \\
\href{notes/heartbeat-monitoring.md}{\texttt{notes/heartbeat-monitoring.md}} &
Heartbeat \& liveness --- alive counters, FTTI, timeout selection, startup
grace \\
\href{notes/distributed-safety-patterns.md}{\texttt{notes/distributed-safety-patterns.md}}
& Distributed safety --- defense in depth, ESTOP, fail-safe, NC wiring,
debounce \\
\href{notes/endianness-binary-protocols.md}{\texttt{notes/endianness-binary-protocols.md}}
& Endianness \& binary protocols --- big/little endian, bit numbering, struct
packing \\
\href{notes/can-gateway-design.md}{\texttt{notes/can-gateway-design.md}} & CAN
gateway design --- multi-bus architecture, forwarding categories, isolation \\
\href{notes/analog-interfacing.md}{\texttt{notes/analog-interfacing.md}} &
Analog interfacing --- ADC, DAC, optoisolators, relays, protection circuits \\
\href{docs/physics-model.md}{\texttt{docs/physics-model.md}} & Tricycle
kinematics --- forward/inverse, rollover, slip angles \\
\href{docs/listen-before-speaking.md}{\texttt{docs/listen-before-speaking.md}} &
CAN actuator safe bootstrapping pattern \\
\href{docs/can-gateway-bridging.md}{\texttt{docs/can-gateway-bridging.md}} & CAN
gateway forwarding rules and implementation \\
\href{docs/emergency-system.md}{\texttt{docs/emergency-system.md}} &
\textbf{Primary:} ESTOP triggers, 8-layer defense, emergency response matrix,
rider's guide, EGAS 3-level architecture, testing procedures \\
\href{docs/emergency-safety-analysis.md}{\texttt{docs/emergency-safety-analysis.md}}
& \textbf{Safety analysis:} ESTOP exit race condition, SEB brake lever
contradiction, watchdog unbraked window --- causal traces, risk quantification,
recommended fixes \\
\href{docs/defense-in-depth-safety.md}{\texttt{docs/defense-in-depth-safety.md}}
& Layered safety --- ESTOP, following error, dynamic clamp, OR logic \\
\href{docs/syntree-security-protocol.md}{\texttt{docs/syntree-security-protocol.md}}
& Rolling counter + XOR checksum for SYNTREE actuators \\
\href{docs/high-voltage-isolation.md}{\texttt{docs/high-voltage-isolation.md}} &
72V galvanic isolation --- TLP281 optos, relays, fuses, TVS \\
\href{docs/distributed-architecture.md}{\texttt{docs/distributed-architecture.md}}
& Three-node rationale --- Jetson/RT/SYS split, dual-CAN on RT \\
\href{docs/actuator-interfacing.md}{\texttt{docs/actuator-interfacing.md}} &
MCP4725 DAC throttle, gear pass-through, relay logic \\
\href{docs/external-watchdog.md}{\texttt{docs/external-watchdog.md}} & External
watchdog IC --- timeout, safe state, testing \\
\href{docs/pid-speed-control.md}{\texttt{docs/pid-speed-control.md}} & PID speed
control theory, tuning, anti-windup \\
\end{longtable}


\chapter{CAN Signal Dictionary --- E-Trike}\label{can-signal-dictionary-e-trike}

Two physical CAN buses at 500 kbit/s. All fields big-endian (MSB first) unless
noted (SYNTREE protocol uses Motorola LSB).

\subsection{Type Notation}\label{type-notation}

\begin{longtable}[]{@{}lll@{}}
\toprule\noalign{}
Notation & C Type & Meaning \\
\midrule\noalign{}
\endhead
\bottomrule\noalign{}
\endlastfoot
\texttt{i8\ /\ i16\ /\ i32} & \texttt{int8\_t\ /\ int16\_t\ /\ int32\_t} &
Signed integer \\
\texttt{i24} & 24-bit signed (packed) & Non-standard width, CAN frame only \\
\texttt{u8\ /\ u16\ /\ u32} & \texttt{uint8\_t\ /\ uint16\_t\ /\ uint32\_t} &
Unsigned integer \\
\texttt{u8\ bool} & \texttt{uint8\_t} & Boolean in a byte (0 or 1) \\
\texttt{u8\ enum} & \texttt{uint8\_t} & Enumeration in a byte \\
\texttt{u8\ bitmask} & \texttt{uint8\_t} & Bitfield, each bit is a flag \\
\texttt{DLC=0} & --- & Zero-length CAN frame (event signal, no payload) \\
\end{longtable}

\subsection{Physical Units}\label{physical-units}

\begin{longtable}[]{@{}lll@{}}
\toprule\noalign{}
Unit & Meaning & Example \\
\midrule\noalign{}
\endhead
\bottomrule\noalign{}
\endlastfoot
\texttt{mm/s} & Millimeters per second & Speed: 3000 = 3.0 m/s \\
\texttt{mrad/s} & Milliradians per second & Yaw rate \\
\texttt{kPa} & Kilopascals & Brake pressure \\
\texttt{mm} & Millimeters & Distance \\
\texttt{0.1°/bit} & Tenths of a degree per LSB & Steer angle: 455 = 45.5° \\
\texttt{0.05\ mm/bit} & 0.05 mm per LSB & Brake stroke \\
\texttt{0.05\ MPa/bit} & 0.05 MPa per LSB & Brake pressure \\
\texttt{°/s} & Degrees per second & Steer angle speed \\
\texttt{Nm} & Newton-meters & Torque \\
\end{longtable}

\begin{center}\rule{0.5\linewidth}{0.5pt}\end{center}

\section{1. Low-Level CAN Bus}\label{low-level-can-bus}

Nodes: RT ESP32-S3, SYS ESP32-S3, SYNTREE EPS-C (steering), SYNTREE SEB (brake),
DC-DC converter.

\begin{center}\rule{0.5\linewidth}{0.5pt}\end{center}

\subsection{0x001 --- SAFETY\_ESTOP}\label{x001-safety_estop}

\begin{longtable}[]{@{}ll@{}}
\toprule\noalign{}
Property & Value \\
\midrule\noalign{}
\endhead
\bottomrule\noalign{}
\endlastfoot
\textbf{Sender} & Any (RT, SYS) \\
\textbf{Receiver(s)} & All nodes \\
\textbf{DLC} & 0 \\
\textbf{Period} & On event \\
\end{longtable}

Presence of this frame = emergency stop. Motor stop, brake engage, steering
disable, DCDC off.

\begin{center}\rule{0.5\linewidth}{0.5pt}\end{center}

\subsection{0x011 --- SYS\_SAFETY\_STS}\label{x011-sys_safety_sts}

\begin{longtable}[]{@{}ll@{}}
\toprule\noalign{}
Property & Value \\
\midrule\noalign{}
\endhead
\bottomrule\noalign{}
\endlastfoot
\textbf{Sender} & SYS \\
\textbf{Receiver(s)} & RT (→ Jetson) \\
\textbf{DLC} & 2 \\
\textbf{Period} & 5 Hz \\
\end{longtable}

\begin{longtable}[]{@{}lllllllll@{}}
\toprule\noalign{}
Signal & Start bit & Len & Type & Scale & Offset & Min & Max & Unit \\
\midrule\noalign{}
\endhead
\bottomrule\noalign{}
\endlastfoot
\texttt{SYS\_EstopActive} & 0 & 8 & u8 & 1 & 0 & 0 & 1 & --- \\
\texttt{SYS\_HeartbeatOk} & 8 & 8 & u8 & 1 & 0 & 0 & 1 & --- \\
\end{longtable}

\begin{center}\rule{0.5\linewidth}{0.5pt}\end{center}

\subsection{0x012 --- SYS\_DCDC\_CMD}\label{x012-sys_dcdc_cmd}

\begin{longtable}[]{@{}ll@{}}
\toprule\noalign{}
Property & Value \\
\midrule\noalign{}
\endhead
\bottomrule\noalign{}
\endlastfoot
\textbf{Sender} & SYS \\
\textbf{Receiver(s)} & DC-DC converter (72V→12V) \\
\textbf{DLC} & 1 \\
\textbf{Period} & On change \\
\end{longtable}

\begin{longtable}[]{@{}lllllllll@{}}
\toprule\noalign{}
Signal & Start bit & Len & Type & Scale & Offset & Min & Max & Unit \\
\midrule\noalign{}
\endhead
\bottomrule\noalign{}
\endlastfoot
\texttt{SYS\_DcdcEnable} & 0 & 8 & u8 & 1 & 0 & 0 & 1 & --- \\
\end{longtable}

ESTOP → \textbf{1 (on)} --- maintains 12V for MCUs, CAN transceivers, and brake
light. The 12V accessory relay (GPIO27) provides the secondary cut for
non-safety loads (see architecture §8.6). All other modes → 1 (on).

\begin{center}\rule{0.5\linewidth}{0.5pt}\end{center}

\subsection{0x110 --- SYS\_MODE\_CMD}\label{x110-sys_mode_cmd}

\begin{longtable}[]{@{}ll@{}}
\toprule\noalign{}
Property & Value \\
\midrule\noalign{}
\endhead
\bottomrule\noalign{}
\endlastfoot
\textbf{Sender} & SYS \\
\textbf{Receiver(s)} & RT \\
\textbf{DLC} & 1 \\
\textbf{Period} & On change \\
\end{longtable}

\begin{longtable}[]{@{}lllllllll@{}}
\toprule\noalign{}
Signal & Start bit & Len & Type & Scale & Offset & Min & Max & Unit \\
\midrule\noalign{}
\endhead
\bottomrule\noalign{}
\endlastfoot
\texttt{SYS\_Mode} & 0 & 8 & u8 & 1 & 0 & 0 & 2 & enum (0=M, 1=A, 2=ESTOP) \\
\end{longtable}

\begin{center}\rule{0.5\linewidth}{0.5pt}\end{center}

\subsection{0x120 --- SYS\_THROTTLE\_STS}\label{x120-sys_throttle_sts}

\begin{longtable}[]{@{}ll@{}}
\toprule\noalign{}
Property & Value \\
\midrule\noalign{}
\endhead
\bottomrule\noalign{}
\endlastfoot
\textbf{Sender} & MTR \\
\textbf{Receiver(s)} & RT (→ Jetson) \\
\textbf{DLC} & 2 \\
\textbf{Period} & 100 Hz \\
\end{longtable}

\begin{longtable}[]{@{}lllllllll@{}}
\toprule\noalign{}
Signal & Start bit & Len & Type & Scale & Offset & Min & Max & Unit \\
\midrule\noalign{}
\endhead
\bottomrule\noalign{}
\endlastfoot
\texttt{SYS\_ThrottleSpeed} & 0 & 16 & i16 & 1 & 0 & -500 & 3000 & mm/s \\
\end{longtable}

\begin{center}\rule{0.5\linewidth}{0.5pt}\end{center}

\subsection{0x204 --- RT\_DRIVE\_CMD}\label{x204-rt_drive_cmd}

\begin{longtable}[]{@{}ll@{}}
\toprule\noalign{}
Property & Value \\
\midrule\noalign{}
\endhead
\bottomrule\noalign{}
\endlastfoot
\textbf{Sender} & RT \\
\textbf{Receiver(s)} & SYS, \textbf{MTR (STM32)} \\
\textbf{DLC} & 5 \\
\textbf{Period} & 100 Hz \\
\end{longtable}

\begin{longtable}[]{@{}lllllllll@{}}
\toprule\noalign{}
Signal & Start bit & Len & Type & Scale & Offset & Min & Max & Unit \\
\midrule\noalign{}
\endhead
\bottomrule\noalign{}
\endlastfoot
\texttt{RT\_MotorSpeed} & 0 & 32 & i32 & 1 & 0 & -500 & 3000 & mm/s \\
\texttt{RT\_Gear} & 32 & 8 & u8 & 1 & 0 & 0 & 3 & enum (0=N,1=D,2=S,3=R) \\
\end{longtable}

Byte layout (big-endian): Byte 0-3 = speed {[}31:0{]}, Byte 4 = gear.

MTR receives 0x204 directly for motor actuation (speed → MCP4725 DAC, gear →
relay module). SYS also receives 0x204 for EGAS Level 2 monitoring (compares
setpoint vs 0x206 feedback).

\begin{quote}
Placed at \texttt{0x204} to avoid collision with EPS-C error info frame at
\texttt{0x202}. SYNTREE units are preprogrammed and cannot be reconfigured.
\end{quote}

\begin{center}\rule{0.5\linewidth}{0.5pt}\end{center}

\subsection{0x205 --- RT\_BRAKE\_CMD}\label{x205-rt_brake_cmd}

\begin{longtable}[]{@{}ll@{}}
\toprule\noalign{}
Property & Value \\
\midrule\noalign{}
\endhead
\bottomrule\noalign{}
\endlastfoot
\textbf{Sender} & RT \\
\textbf{Receiver(s)} & SYS \\
\textbf{DLC} & 4 \\
\textbf{Period} & 50 Hz (20 ms) \\
\end{longtable}

\begin{longtable}[]{@{}lllllllll@{}}
\toprule\noalign{}
Signal & Start bit & Len & Type & Scale & Offset & Min & Max & Unit \\
\midrule\noalign{}
\endhead
\bottomrule\noalign{}
\endlastfoot
\texttt{RT\_BrakePressure} & 0 & 32 & i32 & 1 & 0 & 0 & 20000 & kPa \\
\end{longtable}

Byte layout (big-endian): Bytes 0-3 = brake pressure {[}kPa{]}.

RT max-select: \texttt{brake\_kpa\ =\ max(rt\_obstacle,\ jetson\_0x301)}. SYS
converts: \texttt{seb\_raw\ =\ (uint8\_t)(kpa\ *\ 0.02f)} (verified SYNTREE
spec: \texttt{VCU\_SEB\_Pre\_Value\_Req} is u8, scale 0.05 MPa/bit, range 0--5
MPa). When \texttt{0x205\ \textgreater{}\ 0}, SYS switches SEB to Pressure Mode
(mode=2). When \texttt{0x205\ ==\ 0}, falls back to Stroke Mode for lever/ESTOP
triggers.

\subsection{0x206 --- MTR\_MOTOR\_FBK}\label{x206-mtr_motor_fbk}

\begin{longtable}[]{@{}ll@{}}
\toprule\noalign{}
Property & Value \\
\midrule\noalign{}
\endhead
\bottomrule\noalign{}
\endlastfoot
\textbf{Sender} & MTR STM32 \\
\textbf{Receiver(s)} & SYS, RT \\
\textbf{DLC} & 4 \\
\textbf{Period} & 50 Hz (20 ms) \\
\end{longtable}

\begin{longtable}[]{@{}
  >{\raggedright\arraybackslash}p{(\columnwidth - 16\tabcolsep) * \real{0.1311}}
  >{\raggedright\arraybackslash}p{(\columnwidth - 16\tabcolsep) * \real{0.1803}}
  >{\raggedright\arraybackslash}p{(\columnwidth - 16\tabcolsep) * \real{0.0820}}
  >{\raggedright\arraybackslash}p{(\columnwidth - 16\tabcolsep) * \real{0.0984}}
  >{\raggedright\arraybackslash}p{(\columnwidth - 16\tabcolsep) * \real{0.1148}}
  >{\raggedright\arraybackslash}p{(\columnwidth - 16\tabcolsep) * \real{0.1311}}
  >{\raggedright\arraybackslash}p{(\columnwidth - 16\tabcolsep) * \real{0.0820}}
  >{\raggedright\arraybackslash}p{(\columnwidth - 16\tabcolsep) * \real{0.0820}}
  >{\raggedright\arraybackslash}p{(\columnwidth - 16\tabcolsep) * \real{0.0984}}@{}}
\toprule\noalign{}
\begin{minipage}[b]{\linewidth}\raggedright
Signal
\end{minipage} & \begin{minipage}[b]{\linewidth}\raggedright
Start bit
\end{minipage} & \begin{minipage}[b]{\linewidth}\raggedright
Len
\end{minipage} & \begin{minipage}[b]{\linewidth}\raggedright
Type
\end{minipage} & \begin{minipage}[b]{\linewidth}\raggedright
Scale
\end{minipage} & \begin{minipage}[b]{\linewidth}\raggedright
Offset
\end{minipage} & \begin{minipage}[b]{\linewidth}\raggedright
Min
\end{minipage} & \begin{minipage}[b]{\linewidth}\raggedright
Max
\end{minipage} & \begin{minipage}[b]{\linewidth}\raggedright
Unit
\end{minipage} \\
\midrule\noalign{}
\endhead
\bottomrule\noalign{}
\endlastfoot
\texttt{MTR\_ActualSpeed} & 0 & 16 & i16 & 1 & 0 & -500 & 3000 & mm/s \\
\texttt{MTR\_GearState} & 16 & 8 & u8 enum & 1 & 0 & 0 & 3 &
\{0=N,1=D,2=S,3=R\} \\
\texttt{MTR\_FaultFlags} & 24 & 8 & u8 bitmask & --- & --- & --- & --- &
bit0=ESTOP active, bit1=CMD timeout (0x204 stale), bit2=ADC fault, bit3=gear
conflict \\
\end{longtable}

Byte layout (big-endian): Bytes 0-1=speed, Byte 2=gear, Byte 3=faults.

\subsection{0x201 --- SES\_STATUS (SYNTREE EPS-C
Feedback)}\label{x201-ses_status-syntree-eps-c-feedback}

\begin{longtable}[]{@{}ll@{}}
\toprule\noalign{}
Property & Value \\
\midrule\noalign{}
\endhead
\bottomrule\noalign{}
\endlastfoot
\textbf{Sender} & SYNTREE EPS-C (steering module) \\
\textbf{Receiver(s)} & RT \\
\textbf{DLC} & 8 \\
\textbf{Period} & 10 ms (100 Hz) \\
\textbf{Endianness} & Motorola LSB (little-endian) \\
\end{longtable}

\begin{longtable}[]{@{}
  >{\raggedright\arraybackslash}p{(\columnwidth - 18\tabcolsep) * \real{0.1081}}
  >{\raggedright\arraybackslash}p{(\columnwidth - 18\tabcolsep) * \real{0.1486}}
  >{\raggedright\arraybackslash}p{(\columnwidth - 18\tabcolsep) * \real{0.0676}}
  >{\raggedright\arraybackslash}p{(\columnwidth - 18\tabcolsep) * \real{0.0811}}
  >{\raggedright\arraybackslash}p{(\columnwidth - 18\tabcolsep) * \real{0.0946}}
  >{\raggedright\arraybackslash}p{(\columnwidth - 18\tabcolsep) * \real{0.1081}}
  >{\raggedright\arraybackslash}p{(\columnwidth - 18\tabcolsep) * \real{0.0676}}
  >{\raggedright\arraybackslash}p{(\columnwidth - 18\tabcolsep) * \real{0.0676}}
  >{\raggedright\arraybackslash}p{(\columnwidth - 18\tabcolsep) * \real{0.0811}}
  >{\raggedright\arraybackslash}p{(\columnwidth - 18\tabcolsep) * \real{0.1757}}@{}}
\toprule\noalign{}
\begin{minipage}[b]{\linewidth}\raggedright
Signal
\end{minipage} & \begin{minipage}[b]{\linewidth}\raggedright
Start bit
\end{minipage} & \begin{minipage}[b]{\linewidth}\raggedright
Len
\end{minipage} & \begin{minipage}[b]{\linewidth}\raggedright
Type
\end{minipage} & \begin{minipage}[b]{\linewidth}\raggedright
Scale
\end{minipage} & \begin{minipage}[b]{\linewidth}\raggedright
Offset
\end{minipage} & \begin{minipage}[b]{\linewidth}\raggedright
Min
\end{minipage} & \begin{minipage}[b]{\linewidth}\raggedright
Max
\end{minipage} & \begin{minipage}[b]{\linewidth}\raggedright
Unit
\end{minipage} & \begin{minipage}[b]{\linewidth}\raggedright
Description
\end{minipage} \\
\midrule\noalign{}
\endhead
\bottomrule\noalign{}
\endlastfoot
\texttt{SES\_INF\_Angle\_Status} & 0 & 1 & bool & 1 & 0 & 0 & 1 & --- & Center
Finding Status. 0=Center Finding, 1=Found. (CSV Row 11) \\
\texttt{SES\_Control\_Mode\_Status} & 1 & 2 & u8 & 1 & 0 & 0 & 3 & enum &
Control Mode Feedback. 0=Manual, 1=Automatic. (CSV Row 12) \\
(unaccounted) & 3 & 3 & --- & --- & --- & --- & --- & --- & Byte 0 bits 3--5 ---
not enumerated in CSV \\
\texttt{SES\_Error\_Status} & 6 & 2 & u8 & 1 & 0 & 0 & 3 & enum & Error Status.
0=Normal, 1=L1 Warning, 2=L2, 3=L3. (CSV Row 13) \\
(unaccounted) & --- & 8 & --- & --- & --- & --- & --- & --- & Byte 1 --- not
enumerated in CSV \\
\texttt{SES\_StrAngle} & 16 & 16 & u16 & 0.1 & -3000 & -700 & 700 & ° & Steering
Angle. Unsigned per CSV. Raw 0→-3000°, raw 30000→0°, raw 23000→-700°, raw
37000→700°. (CSV Row 14) \\
\texttt{SES\_Tgt\_StrAngleSpd} & 32 & 16 & i16 & 0.5 & 0 & 0 & 1480 & °/s &
Target Angle Speed. 16-bit signed per CSV. Overlaps Torq at byte 5. (CSV Row
15) \\
\texttt{EPS\_SteeringWheel\_Torq} & 40 & 8 & u8 & 0.1 & -12.1 & -12 & 12 & Nm &
Steering Wheel Torque Feedback. Overlaps StrAngleSpd{[}15:8{]} at byte 5. Init
0x79 (121 raw = 0 Nm). (CSV Row 16) \\
\texttt{SES\_RollCnt\_Enable\_Status} & 48 & 1 & bool & 1 & 0 & 0 & 1 & --- &
Life Signal Enable Feedback. 0=Invalid, 1=Valid. (CSV Row 17) \\
\texttt{SES\_CheckSum\_Enable\_Status} & 49 & 1 & bool & 1 & 0 & 0 & 1 & --- &
Checksum Enable Feedback. 0=Invalid, 1=Valid. (CSV Row 18) \\
(unaccounted) & 50 & 2 & --- & --- & --- & --- & --- & --- & Byte 6 bits 2--3
--- not enumerated in CSV \\
\texttt{SES\_RollCnt\_Status} & 52 & 4 & u8 & 1 & 0 & 0 & 15 & --- & Life Signal
Feedback. Rolling counter 0--15. (CSV Row 19) \\
\texttt{SES\_CheckSum\_Status} & 56 & 8 & u8 & 1 & 0 & 0 & 255 & --- & Checksum
Feedback = XOR(bytes 0--6) \^{} 0xFF. (CSV Row 20) \\
\end{longtable}

\textbf{Byte layout} (little-endian, CSV as source of truth):

\begin{longtable}[]{@{}
  >{\raggedright\arraybackslash}p{(\columnwidth - 16\tabcolsep) * \real{0.2000}}
  >{\raggedright\arraybackslash}p{(\columnwidth - 16\tabcolsep) * \real{0.1000}}
  >{\raggedright\arraybackslash}p{(\columnwidth - 16\tabcolsep) * \real{0.1000}}
  >{\raggedright\arraybackslash}p{(\columnwidth - 16\tabcolsep) * \real{0.1000}}
  >{\raggedright\arraybackslash}p{(\columnwidth - 16\tabcolsep) * \real{0.1000}}
  >{\raggedright\arraybackslash}p{(\columnwidth - 16\tabcolsep) * \real{0.1000}}
  >{\raggedright\arraybackslash}p{(\columnwidth - 16\tabcolsep) * \real{0.1000}}
  >{\raggedright\arraybackslash}p{(\columnwidth - 16\tabcolsep) * \real{0.1000}}
  >{\raggedright\arraybackslash}p{(\columnwidth - 16\tabcolsep) * \real{0.1000}}@{}}
\toprule\noalign{}
\begin{minipage}[b]{\linewidth}\raggedright
Byte
\end{minipage} & \begin{minipage}[b]{\linewidth}\raggedright
0
\end{minipage} & \begin{minipage}[b]{\linewidth}\raggedright
1
\end{minipage} & \begin{minipage}[b]{\linewidth}\raggedright
2
\end{minipage} & \begin{minipage}[b]{\linewidth}\raggedright
3
\end{minipage} & \begin{minipage}[b]{\linewidth}\raggedright
4
\end{minipage} & \begin{minipage}[b]{\linewidth}\raggedright
5
\end{minipage} & \begin{minipage}[b]{\linewidth}\raggedright
6
\end{minipage} & \begin{minipage}[b]{\linewidth}\raggedright
7
\end{minipage} \\
\midrule\noalign{}
\endhead
\bottomrule\noalign{}
\endlastfoot
Content & AngleSts{[}0{]}+ModeSts{[}1:2{]}+(gap)+Error{[}6:7{]} & (unacc.) &
StrAngle {[}7:0{]} & StrAngle {[}15:8{]} & Speed {[}7:0{]} & Speed{[}15:8{]} /
Torq {[}7:0{]} (overlap) & RollCntEn{[}0{]}+CksEn{[}1{]}+(gap)+RollCnt{[}4:7{]}
& CksSum\_Stat \\
\end{longtable}

\begin{quote}
\textbf{StrAngle conversion (CSV Unsigned, offset=-3000):}
\texttt{physical\_deg\ =\ raw\ ×\ 0.1\ −\ 3000}. The EPS-C encodes steering
angle as an unsigned 16-bit value with -3000 offset. Raw 30000 → 0° (straight).
Raw 23000 → -700° (full left). Raw 37000 → 700° (full right). In practice, RT
uses \texttt{internal\_angle\_mdeg\ /\ 100} to produce the raw value; adjust per
actual calibration.

\textbf{Torq encoding (CSV scale=0.1, offset=-12.1):}
\texttt{physical\_Nm\ =\ raw\ ×\ 0.1\ −\ 12.1}. Raw 121 (0x79) → 0.0 Nm (no
torque). Raw 0 → -12.1 Nm. Raw 241 → 12.0 Nm. The EPS-C biases torque readings
so zero torque is at raw 121.

\textbf{Byte 5 overlap:} CSV declares \texttt{SES\_Tgt\_StrAngleSpd} as 16-bit
(bytes 4--5, Signed) AND \texttt{EPS\_SteeringWheel\_Torq} as 8-bit (byte 5).
Both are listed --- the EPS-C may report these in alternate frames or the CSV
represents signals available across firmware versions.
\end{quote}

\begin{center}\rule{0.5\linewidth}{0.5pt}\end{center}

\subsection{0x169 --- VCU\_SES\_REQ (SYNTREE EPS-C
Command)}\label{x169-vcu_ses_req-syntree-eps-c-command}

\begin{longtable}[]{@{}
  >{\raggedright\arraybackslash}p{(\columnwidth - 2\tabcolsep) * \real{0.5882}}
  >{\raggedright\arraybackslash}p{(\columnwidth - 2\tabcolsep) * \real{0.4118}}@{}}
\toprule\noalign{}
\begin{minipage}[b]{\linewidth}\raggedright
Property
\end{minipage} & \begin{minipage}[b]{\linewidth}\raggedright
Value
\end{minipage} \\
\midrule\noalign{}
\endhead
\bottomrule\noalign{}
\endlastfoot
\textbf{Sender} & RT ESP32-S3 \\
\textbf{Receiver(s)} & SYNTREE EPS-C (steering module) \\
\textbf{DLC} & 8 \\
\textbf{Period} & 20 ms (50 Hz) --- \textbf{continuous, every frame} \\
\textbf{Endianness} & Motorola LSB (little-endian) \\
\textbf{Note} & Factory default \texttt{0x169}. SYNTREE unit is preprogrammed
and not reconfigurable. \texttt{RT\_DRIVE\_CMD} placed at \texttt{0x204} to
avoid collision. \\
\end{longtable}

\begin{longtable}[]{@{}
  >{\raggedright\arraybackslash}p{(\columnwidth - 18\tabcolsep) * \real{0.1081}}
  >{\raggedright\arraybackslash}p{(\columnwidth - 18\tabcolsep) * \real{0.1486}}
  >{\raggedright\arraybackslash}p{(\columnwidth - 18\tabcolsep) * \real{0.0676}}
  >{\raggedright\arraybackslash}p{(\columnwidth - 18\tabcolsep) * \real{0.0811}}
  >{\raggedright\arraybackslash}p{(\columnwidth - 18\tabcolsep) * \real{0.0946}}
  >{\raggedright\arraybackslash}p{(\columnwidth - 18\tabcolsep) * \real{0.1081}}
  >{\raggedright\arraybackslash}p{(\columnwidth - 18\tabcolsep) * \real{0.0676}}
  >{\raggedright\arraybackslash}p{(\columnwidth - 18\tabcolsep) * \real{0.0676}}
  >{\raggedright\arraybackslash}p{(\columnwidth - 18\tabcolsep) * \real{0.0811}}
  >{\raggedright\arraybackslash}p{(\columnwidth - 18\tabcolsep) * \real{0.1757}}@{}}
\toprule\noalign{}
\begin{minipage}[b]{\linewidth}\raggedright
Signal
\end{minipage} & \begin{minipage}[b]{\linewidth}\raggedright
Start bit
\end{minipage} & \begin{minipage}[b]{\linewidth}\raggedright
Len
\end{minipage} & \begin{minipage}[b]{\linewidth}\raggedright
Type
\end{minipage} & \begin{minipage}[b]{\linewidth}\raggedright
Scale
\end{minipage} & \begin{minipage}[b]{\linewidth}\raggedright
Offset
\end{minipage} & \begin{minipage}[b]{\linewidth}\raggedright
Min
\end{minipage} & \begin{minipage}[b]{\linewidth}\raggedright
Max
\end{minipage} & \begin{minipage}[b]{\linewidth}\raggedright
Unit
\end{minipage} & \begin{minipage}[b]{\linewidth}\raggedright
Description
\end{minipage} \\
\midrule\noalign{}
\endhead
\bottomrule\noalign{}
\endlastfoot
\texttt{VCU\_SES\_Alignment\_Enable} & 0 & 1 & bool & 1 & 0 & 0 & 1 & --- & SES
Angle Initial Alignment Enable. 0=disabled, 1=centering. (CSV Row 2) \\
\texttt{VCU\_SES\_Control\_Enable} & 1 & 1 & bool & 1 & 0 & 0 & 1 & --- & VCU
Direction Control Enable. 0=Disabled (Default Assist), 1=Rising Edge Enable
(Angle Control Mode). (CSV Row 3) \\
(unaccounted) & 2 & 6 & --- & --- & --- & --- & --- & --- & Byte 0 bits 2--7 ---
not enumerated in CSV \\
(unaccounted) & --- & 8 & --- & --- & --- & --- & --- & --- & Byte 1 --- not
enumerated in CSV \\
\texttt{VCU\_SES\_Tgt\_StrAngle} & 16 & 16 & i16 & 0.1 & -3000 & -700 & 700 & °
& Target Steering Angle. Negative = left. (CSV Row 4). Note: CSV offset=-3000
(see conversion note below). \\
\texttt{VCU\_SES\_Tgt\_StrAngleSpd} & 32 & 16 & u16 & 1 & 0 & 125 & 525 & °/s &
Target Steering Angle Speed. 16-bit per CSV. Overlaps security signals at byte
5. (CSV Row 5) \\
\texttt{VCU\_SES\_RollCnt\_Enable} & 40 & 1 & bool & 1 & 0 & 0 & 1 & --- & Life
Signal Enable --- \textbf{Must be 1}. Overlaps StrAngleSpd{[}15:8{]}. (CSV Row
6) \\
\texttt{VCU\_SES\_CheckSum\_Enable} & 41 & 1 & bool & 1 & 0 & 0 & 1 & --- &
Checksum Enable --- \textbf{Must be 1}. Overlaps StrAngleSpd{[}15:8{]}. (CSV Row
7) \\
(unaccounted) & 42 & 2 & --- & --- & --- & --- & --- & --- & Byte 5 bits 2--3
--- not enumerated in CSV \\
\texttt{VCU\_SES\_RollCnt} & 44 & 4 & u8 & 1 & 0 & 0 & 15 & --- & Life Signal
rolling counter. Increment every frame. Overlaps StrAngleSpd{[}15:8{]}. (CSV Row
8) \\
\texttt{VCU\_Veh\_Spd\_Value} & 48 & 8 & u8 & 1 & 0 & 0 & 255 & --- & Vehicle
Speed. RT must populate with current speed. (CSV Row 9) \\
\texttt{VCU\_SES\_CheckSum} & 56 & 8 & u8 & 1 & 0 & 0 & 255 & --- & Checksum =
XOR(bytes 0--6) \^{} 0xFF (CSV Row 10) \\
\end{longtable}

\textbf{Byte layout} (little-endian, CSV as source of truth):

\begin{longtable}[]{@{}
  >{\raggedright\arraybackslash}p{(\columnwidth - 16\tabcolsep) * \real{0.2000}}
  >{\raggedright\arraybackslash}p{(\columnwidth - 16\tabcolsep) * \real{0.1000}}
  >{\raggedright\arraybackslash}p{(\columnwidth - 16\tabcolsep) * \real{0.1000}}
  >{\raggedright\arraybackslash}p{(\columnwidth - 16\tabcolsep) * \real{0.1000}}
  >{\raggedright\arraybackslash}p{(\columnwidth - 16\tabcolsep) * \real{0.1000}}
  >{\raggedright\arraybackslash}p{(\columnwidth - 16\tabcolsep) * \real{0.1000}}
  >{\raggedright\arraybackslash}p{(\columnwidth - 16\tabcolsep) * \real{0.1000}}
  >{\raggedright\arraybackslash}p{(\columnwidth - 16\tabcolsep) * \real{0.1000}}
  >{\raggedright\arraybackslash}p{(\columnwidth - 16\tabcolsep) * \real{0.1000}}@{}}
\toprule\noalign{}
\begin{minipage}[b]{\linewidth}\raggedright
Byte
\end{minipage} & \begin{minipage}[b]{\linewidth}\raggedright
0
\end{minipage} & \begin{minipage}[b]{\linewidth}\raggedright
1
\end{minipage} & \begin{minipage}[b]{\linewidth}\raggedright
2
\end{minipage} & \begin{minipage}[b]{\linewidth}\raggedright
3
\end{minipage} & \begin{minipage}[b]{\linewidth}\raggedright
4
\end{minipage} & \begin{minipage}[b]{\linewidth}\raggedright
5
\end{minipage} & \begin{minipage}[b]{\linewidth}\raggedright
6
\end{minipage} & \begin{minipage}[b]{\linewidth}\raggedright
7
\end{minipage} \\
\midrule\noalign{}
\endhead
\bottomrule\noalign{}
\endlastfoot
Content & Align{[}0{]}+CtrlEn{[}1{]} & (unacc.) & Angle {[}7:0{]} & Angle
{[}15:8{]} & Speed {[}7:0{]} & Speed{[}15:8{]} /
RollCntEn{[}0{]}+CksEn{[}1{]}+(gap)+RollCnt{[}4:7{]} (overlap) & Veh\_Spd
{[}7:0{]} & CheckSum \\
\end{longtable}

\begin{quote}
\textbf{Angle conversion (CSV offset=-3000):}
\texttt{physical\_deg\ =\ raw\ ×\ 0.1\ +\ (-3000)}. This encoding places the
±700° physical range at raw values approximately 23000--37000. Raw 0 → -3000°
(outside normal range). This offset is unusual for a signed integer --- 0° does
not map to raw 0. The CSV offset may be a tool artifact; verify against observed
CAN bus values.

\textbf{Byte 5 overlap:} CSV declares \texttt{VCU\_SES\_Tgt\_StrAngleSpd} as
16-bit (bytes 4--5) AND security signals at byte 5 (RollCnt\_Enable,
CheckSum\_Enable, RollCnt). Both are listed in the manufacturer's DBC export.
The EPS-C may internally separate these --- the upper nibble of byte 5 carries
security data while the lower bits carry speed. RT firmware should write the
speed value to bytes 4--5 AND set security fields at byte 5 bits 0--3 + upper
nibble; the EPS-C validates the security portion independently of the speed
portion.

\textbf{Slew rate:} Speed-dependent. RT computes
\texttt{VCU\_SES\_Tgt\_StrAngleSpd} based on speed to ensure smooth steering.
CSV lists range 125--525 °/s. The EPS-C may reject speed commands below 125 °/s.
\end{quote}

\textbf{Internal conversion (architecture, offset=0)}:
\texttt{VCU\_SES\_Tgt\_StrAngle\_raw\ =\ internal\_angle\_mdeg\ /\ 100} (45500
mdeg → 455 raw → 45.5°). CSV declares offset=-3000; if that encoding is used,
the formula would be \texttt{raw\ =\ internal\_angle\_mdeg\ /\ 100\ +\ 30000}
(45500 mdeg → 30455 raw → 45.5°). Verify which encoding the EPS-C actually
expects by observing CAN bus traffic.

\textbf{Security}: If \texttt{roll\_cnt\_enable=0} or
\texttt{checksum\_enable=0}, unit may reject frames. Both must be 1. Checksum
algorithm: \texttt{XOR(bytes{[}0..6{]})\ \^{}\ 0xFF} (verify exact formula
against SYNTREE spec).

\textbf{Slew rate}: Speed-dependent. RT computes
\texttt{VCU\_SES\_Tgt\_StrAngleSpd} based on speed to ensure smooth steering.
Lower speed → lower slew rate for comfort; higher speed → higher slew rate for
responsiveness (within dynamic clamp).

\begin{center}\rule{0.5\linewidth}{0.5pt}\end{center}

\subsection{0x202 --- SES\_ErrInfo (SYNTREE EPS-C Error
Detail)}\label{x202-ses_errinfo-syntree-eps-c-error-detail}

\begin{longtable}[]{@{}ll@{}}
\toprule\noalign{}
Property & Value \\
\midrule\noalign{}
\endhead
\bottomrule\noalign{}
\endlastfoot
\textbf{Sender} & SYNTREE EPS-C (steering module) \\
\textbf{Receiver(s)} & RT ESP32-S3 \\
\textbf{DLC} & 8 \\
\textbf{Period} & 100 ms (10 Hz) \\
\textbf{Endianness} & Motorola LSB (little-endian) \\
\end{longtable}

Detailed fault flags. Each bit is an independent fault indicator. 1 = fault
active.

\begin{longtable}[]{@{}
  >{\raggedright\arraybackslash}p{(\columnwidth - 10\tabcolsep) * \real{0.1429}}
  >{\raggedright\arraybackslash}p{(\columnwidth - 10\tabcolsep) * \real{0.1964}}
  >{\raggedright\arraybackslash}p{(\columnwidth - 10\tabcolsep) * \real{0.0893}}
  >{\raggedright\arraybackslash}p{(\columnwidth - 10\tabcolsep) * \real{0.1071}}
  >{\raggedright\arraybackslash}p{(\columnwidth - 10\tabcolsep) * \real{0.2321}}
  >{\raggedright\arraybackslash}p{(\columnwidth - 10\tabcolsep) * \real{0.2321}}@{}}
\toprule\noalign{}
\begin{minipage}[b]{\linewidth}\raggedright
Signal
\end{minipage} & \begin{minipage}[b]{\linewidth}\raggedright
Start bit
\end{minipage} & \begin{minipage}[b]{\linewidth}\raggedright
Len
\end{minipage} & \begin{minipage}[b]{\linewidth}\raggedright
Type
\end{minipage} & \begin{minipage}[b]{\linewidth}\raggedright
Fault Level
\end{minipage} & \begin{minipage}[b]{\linewidth}\raggedright
Description
\end{minipage} \\
\midrule\noalign{}
\endhead
\bottomrule\noalign{}
\endlastfoot
\texttt{SES\_ECUUnderVolt\_Err} & 0 & 1 & bool & L2 & Controller Under
Voltage \\
\texttt{SES\_ECUOverVolt\_Err} & 1 & 1 & bool & L2 & Controller Over Voltage \\
\texttt{SES\_CanCom\_Err} & 2 & 1 & bool & L1 & CAN Communication Fault \\
\texttt{SES\_ECUTemp\_Err} & 3 & 1 & bool & L1 & Controller Temp Fault \\
\texttt{SES\_Domain\_drive\_SC\_Err} & 4 & 1 & bool & L2 & Domain Drive Short
Circuit \\
\texttt{SES\_Domain\_drive\_V\_Err} & 5 & 1 & bool & L2 & Domain Drive Voltage
Fault \\
\texttt{SES\_Domain\_drive\_T\_Err} & 6 & 1 & bool & L2 & Domain Drive
Temperature Fault \\
\texttt{SES\_TempSensor\_Err} & 7 & 1 & bool & --- & Temperature Sensor Fault \\
\texttt{SES\_AngleSensor\_P\_OC\_Err} & 8 & 1 & bool & \textbf{L3} & Angle
Sensor Pri. Open Circuit \\
\texttt{SES\_AngleSensor\_P\_AF\_Err} & 9 & 1 & bool & \textbf{L3} & Angle
Sensor Pri. Out of Range \\
\texttt{SES\_AngleSensor\_S\_OC\_Err} & 10 & 1 & bool & \textbf{L3} & Angle
Sensor Sec. Open Circuit \\
\texttt{SES\_AngleSensor\_S\_AF\_Err} & 11 & 1 & bool & \textbf{L3} & Angle
Sensor Sec. Out of Range \\
\texttt{SES\_SensorPow\_Err} & 12 & 1 & bool & L2 & Sensor Power Fault \\
\texttt{SES\_Alignment\_Err} & 13 & 1 & bool & L1 & Centering Fault \\
\texttt{SES\_OverAngle\_Err} & 14 & 1 & bool & L2 & Over Angle Fault \\
\texttt{SES\_StrMtr\_Stall\_Err} & 15 & 1 & bool & L1 & Motor Stall Fault \\
\texttt{SES\_MtrCurt\_Err} & 16 & 1 & bool & L2 & Motor Current Fault \\
\texttt{SES\_SensorCL\_Err} & 17 & 1 & bool & L2 & Sensor 5V Power 1/2 Fault \\
\texttt{SES\_TorqSensor\_T1\_OC\_Err} & 18 & 1 & bool & \textbf{L3} & Torque
Sensor T1 Open Circuit \\
\texttt{SES\_TorqSensor\_T1\_AF\_Err} & 19 & 1 & bool & \textbf{L3} & Torque
Sensor T1 Out of Range \\
\texttt{SES\_TorqSensor\_T2\_OC\_Err} & 20 & 1 & bool & \textbf{L3} & Torque
Sensor T2 Open Circuit \\
\texttt{SES\_TorqSensor\_T2\_AF\_Err} & 21 & 1 & bool & \textbf{L3} & Torque
Sensor T2 Out of Range \\
\texttt{SentAngle\_Err} & 22 & 1 & bool & L1 & Angle Error \\
\texttt{SES\_StrMtr\_Idling\_Err} & 23 & 1 & bool & L2 & Motor Idling Fault \\
\texttt{SES\_EPROM\_Err} & 24 & 1 & bool & L2 & EEPROM Fault \\
(reserved) & 25 & 31 & --- & --- & Bits 25--55 (byte 3 bits 1--7 + bytes 4--6).
Not enumerated in CSV. \\
\texttt{SES\_Veh\_Spd\_Value} & 56 & 8 & u8 & --- & Vehicle speed at fault
snapshot (CSV Row 46) \\
\end{longtable}

\textbf{Byte layout} (little-endian):

\begin{longtable}[]{@{}
  >{\raggedright\arraybackslash}p{(\columnwidth - 16\tabcolsep) * \real{0.2000}}
  >{\raggedright\arraybackslash}p{(\columnwidth - 16\tabcolsep) * \real{0.1000}}
  >{\raggedright\arraybackslash}p{(\columnwidth - 16\tabcolsep) * \real{0.1000}}
  >{\raggedright\arraybackslash}p{(\columnwidth - 16\tabcolsep) * \real{0.1000}}
  >{\raggedright\arraybackslash}p{(\columnwidth - 16\tabcolsep) * \real{0.1000}}
  >{\raggedright\arraybackslash}p{(\columnwidth - 16\tabcolsep) * \real{0.1000}}
  >{\raggedright\arraybackslash}p{(\columnwidth - 16\tabcolsep) * \real{0.1000}}
  >{\raggedright\arraybackslash}p{(\columnwidth - 16\tabcolsep) * \real{0.1000}}
  >{\raggedright\arraybackslash}p{(\columnwidth - 16\tabcolsep) * \real{0.1000}}@{}}
\toprule\noalign{}
\begin{minipage}[b]{\linewidth}\raggedright
Byte
\end{minipage} & \begin{minipage}[b]{\linewidth}\raggedright
0
\end{minipage} & \begin{minipage}[b]{\linewidth}\raggedright
1
\end{minipage} & \begin{minipage}[b]{\linewidth}\raggedright
2
\end{minipage} & \begin{minipage}[b]{\linewidth}\raggedright
3
\end{minipage} & \begin{minipage}[b]{\linewidth}\raggedright
4
\end{minipage} & \begin{minipage}[b]{\linewidth}\raggedright
5
\end{minipage} & \begin{minipage}[b]{\linewidth}\raggedright
6
\end{minipage} & \begin{minipage}[b]{\linewidth}\raggedright
7
\end{minipage} \\
\midrule\noalign{}
\endhead
\bottomrule\noalign{}
\endlastfoot
Content & faults {[}7:0{]} & faults {[}15:8{]} & faults {[}23:16{]} & faults
{[}31:24{]} & rsvd & rsvd & rsvd & \texttt{SES\_Veh\_Spd\_Value} \\
\end{longtable}

\begin{quote}
\textbf{Safety note:} 8 faults are L3 (angle sensor primary/secondary
open-circuit/out-of-range × 4 + torque sensor T1/T2 open-circuit/out-of-range ×
4). RT must subscribe and escalate L3 faults to ESTOP. The EPS-C has
dual-redundant angle sensors (P/S) and dual torque sensors (T1/T2) --- L3 on
either channel of a redundant pair indicates critical sensor failure. Note fault
levels differ from SEB: steering CAN comms is L1 (minor) vs brake CAN comms L3
(severe) --- steering comm loss is less immediately dangerous than brake comm
loss.
\end{quote}

\begin{quote}
⚠️ \textbf{CSV Row 40--41 description swap:} Row 40 signal
\texttt{SES\_TorqSensor\_T2\_OC\_Err} is labeled ``Torque Sensor T1 Out of
Range'' --- descriptions appear swapped. Signal names used above are
authoritative; descriptions are corrected.
\end{quote}

\begin{center}\rule{0.5\linewidth}{0.5pt}\end{center}

\subsection{0x203 --- SES\_Version (SYNTREE EPS-C Firmware
Version)}\label{x203-ses_version-syntree-eps-c-firmware-version}

\begin{longtable}[]{@{}ll@{}}
\toprule\noalign{}
Property & Value \\
\midrule\noalign{}
\endhead
\bottomrule\noalign{}
\endlastfoot
\textbf{Sender} & SYNTREE EPS-C (steering module) \\
\textbf{Receiver(s)} & RT ESP32-S3 \\
\textbf{DLC} & 8 \\
\textbf{Period} & 1000 ms (1 Hz) \\
\textbf{Endianness} & Motorola LSB (little-endian) \\
\end{longtable}

\begin{longtable}[]{@{}
  >{\raggedright\arraybackslash}p{(\columnwidth - 18\tabcolsep) * \real{0.1081}}
  >{\raggedright\arraybackslash}p{(\columnwidth - 18\tabcolsep) * \real{0.1486}}
  >{\raggedright\arraybackslash}p{(\columnwidth - 18\tabcolsep) * \real{0.0676}}
  >{\raggedright\arraybackslash}p{(\columnwidth - 18\tabcolsep) * \real{0.0811}}
  >{\raggedright\arraybackslash}p{(\columnwidth - 18\tabcolsep) * \real{0.0946}}
  >{\raggedright\arraybackslash}p{(\columnwidth - 18\tabcolsep) * \real{0.1081}}
  >{\raggedright\arraybackslash}p{(\columnwidth - 18\tabcolsep) * \real{0.0676}}
  >{\raggedright\arraybackslash}p{(\columnwidth - 18\tabcolsep) * \real{0.0676}}
  >{\raggedright\arraybackslash}p{(\columnwidth - 18\tabcolsep) * \real{0.0811}}
  >{\raggedright\arraybackslash}p{(\columnwidth - 18\tabcolsep) * \real{0.1757}}@{}}
\toprule\noalign{}
\begin{minipage}[b]{\linewidth}\raggedright
Signal
\end{minipage} & \begin{minipage}[b]{\linewidth}\raggedright
Start bit
\end{minipage} & \begin{minipage}[b]{\linewidth}\raggedright
Len
\end{minipage} & \begin{minipage}[b]{\linewidth}\raggedright
Type
\end{minipage} & \begin{minipage}[b]{\linewidth}\raggedright
Scale
\end{minipage} & \begin{minipage}[b]{\linewidth}\raggedright
Offset
\end{minipage} & \begin{minipage}[b]{\linewidth}\raggedright
Min
\end{minipage} & \begin{minipage}[b]{\linewidth}\raggedright
Max
\end{minipage} & \begin{minipage}[b]{\linewidth}\raggedright
Unit
\end{minipage} & \begin{minipage}[b]{\linewidth}\raggedright
Description
\end{minipage} \\
\midrule\noalign{}
\endhead
\bottomrule\noalign{}
\endlastfoot
\texttt{SES\_SW\_Version} & 0 & 8 & u8 & 0.01 & 0 & 0 & 255 & --- & Software
version (e.g., 0x64 = 1.00) \\
\texttt{SES\_HW\_Version} & 8 & 8 & u8 & 0.1 & 0 & 0 & 25.5 & --- & Hardware
version (e.g., 0x0D = 1.3) \\
(reserved) & 16 & 48 & --- & --- & --- & --- & --- & --- & Bytes 2--7 (code also
reads bytes 2-3 as extended HW version info --- minor logging discrepancy) \\
\end{longtable}

\textbf{RT usage}: Log on boot for compatibility check
(\texttt{SW=\%02X.\%02X\ HW=\%02X.\%02X}). Report via telemetry to Jetson.

\begin{center}\rule{0.5\linewidth}{0.5pt}\end{center}

\subsection{0x6FA --- SES\_Test (SYNTREE EPS-C
Telemetry)}\label{x6fa-ses_test-syntree-eps-c-telemetry}

\begin{longtable}[]{@{}ll@{}}
\toprule\noalign{}
Property & Value \\
\midrule\noalign{}
\endhead
\bottomrule\noalign{}
\endlastfoot
\textbf{Sender} & SYNTREE EPS-C (steering module) \\
\textbf{Receiver(s)} & RT ESP32-S3 \\
\textbf{DLC} & 8 \\
\textbf{Period} & 10 ms (100 Hz) \\
\textbf{Endianness} & Motorola LSB (little-endian) \\
\end{longtable}

\begin{longtable}[]{@{}
  >{\raggedright\arraybackslash}p{(\columnwidth - 18\tabcolsep) * \real{0.1081}}
  >{\raggedright\arraybackslash}p{(\columnwidth - 18\tabcolsep) * \real{0.1486}}
  >{\raggedright\arraybackslash}p{(\columnwidth - 18\tabcolsep) * \real{0.0676}}
  >{\raggedright\arraybackslash}p{(\columnwidth - 18\tabcolsep) * \real{0.0811}}
  >{\raggedright\arraybackslash}p{(\columnwidth - 18\tabcolsep) * \real{0.0946}}
  >{\raggedright\arraybackslash}p{(\columnwidth - 18\tabcolsep) * \real{0.1081}}
  >{\raggedright\arraybackslash}p{(\columnwidth - 18\tabcolsep) * \real{0.0676}}
  >{\raggedright\arraybackslash}p{(\columnwidth - 18\tabcolsep) * \real{0.0676}}
  >{\raggedright\arraybackslash}p{(\columnwidth - 18\tabcolsep) * \real{0.0811}}
  >{\raggedright\arraybackslash}p{(\columnwidth - 18\tabcolsep) * \real{0.1757}}@{}}
\toprule\noalign{}
\begin{minipage}[b]{\linewidth}\raggedright
Signal
\end{minipage} & \begin{minipage}[b]{\linewidth}\raggedright
Start bit
\end{minipage} & \begin{minipage}[b]{\linewidth}\raggedright
Len
\end{minipage} & \begin{minipage}[b]{\linewidth}\raggedright
Type
\end{minipage} & \begin{minipage}[b]{\linewidth}\raggedright
Scale
\end{minipage} & \begin{minipage}[b]{\linewidth}\raggedright
Offset
\end{minipage} & \begin{minipage}[b]{\linewidth}\raggedright
Min
\end{minipage} & \begin{minipage}[b]{\linewidth}\raggedright
Max
\end{minipage} & \begin{minipage}[b]{\linewidth}\raggedright
Unit
\end{minipage} & \begin{minipage}[b]{\linewidth}\raggedright
Description
\end{minipage} \\
\midrule\noalign{}
\endhead
\bottomrule\noalign{}
\endlastfoot
(reserved) & 0 & 8 & --- & --- & --- & --- & --- & --- & Byte 0 \\
\texttt{SES\_MtrCurt} & 8 & 16 & i16 & 0.0078125 & 0 & 0 & 60 & A & Motor
current. Bytes 1--2. Narrower range than brake SEB\_Test (±255A). \\
\texttt{SES\_ECUTemp} & 24 & 16 & u16 & 0.5 & 0 & 0 & 255 & °C & ECU
temperature. Bytes 3--4. \\
\texttt{SES\_PowVolt} & 40 & 16 & u16 & 0.00390625 & 0 & 0 & 18 & V & Supply
voltage. Bytes 5--6. Narrower range than brake (0--18V vs 0--32V). \\
(reserved) & 56 & 8 & --- & --- & --- & --- & --- & --- & Byte 7 \\
\end{longtable}

\begin{quote}
\textbf{Note:} CSV uses byte-local start-bit numbering for this message.
Converted to absolute Motorola LSB above. Range limits differ from brake
SEB\_Test (0x6FB): steering motor current 0--60A vs brake ±255A; steering supply
voltage 0--18V vs brake 0--32V. \textbf{RT usage:} Monitor \texttt{SES\_MtrCurt}
for mechanical binding / rack damage. \texttt{SES\_ECUTemp} for thermal
throttling.
\end{quote}

\begin{center}\rule{0.5\linewidth}{0.5pt}\end{center}

\subsection{0x302 --- HOST\_LIGHT\_CMD
(forwarded)}\label{x302-host_light_cmd-forwarded}

\begin{longtable}[]{@{}ll@{}}
\toprule\noalign{}
Property & Value \\
\midrule\noalign{}
\endhead
\bottomrule\noalign{}
\endlastfoot
\textbf{Sender} & RT (fwd from Jetson) \\
\textbf{Receiver(s)} & SYS \\
\textbf{DLC} & 1 \\
\textbf{Period} & On change \\
\end{longtable}

\begin{longtable}[]{@{}llll@{}}
\toprule\noalign{}
Signal & Start bit & Len & Type \\
\midrule\noalign{}
\endhead
\bottomrule\noalign{}
\endlastfoot
\texttt{HOST\_LeftTurn} & 0 & 1 & bool \\
\texttt{HOST\_RightTurn} & 1 & 1 & bool \\
\texttt{HOST\_BrakeLight} & 2 & 1 & bool \\
\texttt{HOST\_Headlight} & 3 & 1 & bool \\
(reserved) & 4 & 4 & --- \\
\end{longtable}

\begin{center}\rule{0.5\linewidth}{0.5pt}\end{center}

\subsection{0x600 --- SYS\_DIAG\_RPT}\label{x600-sys_diag_rpt}

\begin{longtable}[]{@{}ll@{}}
\toprule\noalign{}
Property & Value \\
\midrule\noalign{}
\endhead
\bottomrule\noalign{}
\endlastfoot
\textbf{Sender} & SYS \\
\textbf{Receiver(s)} & RT (→ Jetson) \\
\textbf{DLC} & 8 \\
\textbf{Period} & 1 Hz \\
\end{longtable}

\begin{longtable}[]{@{}llll@{}}
\toprule\noalign{}
Signal & Start bit & Len & Type \\
\midrule\noalign{}
\endhead
\bottomrule\noalign{}
\endlastfoot
\texttt{SYS\_DiagMode} & 0 & 8 & u8 \\
\texttt{SYS\_DiagBrakeEngaged} & 8 & 8 & u8 \\
\texttt{SYS\_DiagHeartbeatOk} & 16 & 8 & u8 \\
\texttt{SYS\_DiagEstopActive} & 24 & 8 & u8 \\
\texttt{SYS\_DiagFreeHeapKb} & 32 & 16 & u16 \\
\texttt{SYS\_DiagTec} & 48 & 8 & u8 \\
\texttt{SYS\_DiagRec} & 56 & 8 & u8 \\
\end{longtable}

Byte layout (big-endian): Byte 0=mode, 1=brake, 2=hb\_ok, 3=estop, 4-5=heap,
6=tec, 7=rec.

\begin{center}\rule{0.5\linewidth}{0.5pt}\end{center}

\subsection{0x7B9 --- VCU\_SEB\_REQ (SYNTREE SEB Brake
Command)}\label{x7b9-vcu_seb_req-syntree-seb-brake-command}

\begin{longtable}[]{@{}ll@{}}
\toprule\noalign{}
Property & Value \\
\midrule\noalign{}
\endhead
\bottomrule\noalign{}
\endlastfoot
\textbf{Sender} & SYS ESP32-S3 \\
\textbf{Receiver(s)} & SYNTREE SEB (brake module) \\
\textbf{DLC} & 8 \\
\textbf{Period} & 20 ms (50 Hz) --- \textbf{continuous, every frame} \\
\textbf{Endianness} & Motorola LSB (little-endian) \\
\end{longtable}

\begin{longtable}[]{@{}
  >{\raggedright\arraybackslash}p{(\columnwidth - 18\tabcolsep) * \real{0.1081}}
  >{\raggedright\arraybackslash}p{(\columnwidth - 18\tabcolsep) * \real{0.1486}}
  >{\raggedright\arraybackslash}p{(\columnwidth - 18\tabcolsep) * \real{0.0676}}
  >{\raggedright\arraybackslash}p{(\columnwidth - 18\tabcolsep) * \real{0.0811}}
  >{\raggedright\arraybackslash}p{(\columnwidth - 18\tabcolsep) * \real{0.0946}}
  >{\raggedright\arraybackslash}p{(\columnwidth - 18\tabcolsep) * \real{0.1081}}
  >{\raggedright\arraybackslash}p{(\columnwidth - 18\tabcolsep) * \real{0.0676}}
  >{\raggedright\arraybackslash}p{(\columnwidth - 18\tabcolsep) * \real{0.0676}}
  >{\raggedright\arraybackslash}p{(\columnwidth - 18\tabcolsep) * \real{0.0811}}
  >{\raggedright\arraybackslash}p{(\columnwidth - 18\tabcolsep) * \real{0.1757}}@{}}
\toprule\noalign{}
\begin{minipage}[b]{\linewidth}\raggedright
Signal
\end{minipage} & \begin{minipage}[b]{\linewidth}\raggedright
Start bit
\end{minipage} & \begin{minipage}[b]{\linewidth}\raggedright
Len
\end{minipage} & \begin{minipage}[b]{\linewidth}\raggedright
Type
\end{minipage} & \begin{minipage}[b]{\linewidth}\raggedright
Scale
\end{minipage} & \begin{minipage}[b]{\linewidth}\raggedright
Offset
\end{minipage} & \begin{minipage}[b]{\linewidth}\raggedright
Min
\end{minipage} & \begin{minipage}[b]{\linewidth}\raggedright
Max
\end{minipage} & \begin{minipage}[b]{\linewidth}\raggedright
Unit
\end{minipage} & \begin{minipage}[b]{\linewidth}\raggedright
Description
\end{minipage} \\
\midrule\noalign{}
\endhead
\bottomrule\noalign{}
\endlastfoot
\texttt{VCU\_SEB\_Alignment\_Enable} & 0 & 1 & bool & 1 & 0 & 0 & 1 & --- &
Calibration enable (CSV Row 2) \\
\texttt{VCU\_SEB\_Control\_Enable} & 1 & 1 & bool & 1 & 0 & 0 & 1 & --- & Active
control enable (CSV Row 3) \\
\texttt{VCU\_SEB\_Control\_Mode} & 2 & 1 & bool & 1 & 0 & 0 & 1 & enum &
0=Stroke, 1=Pressure (CSV Row 4) \\
\texttt{VCU\_SEB\_AutoBrake} & 3 & 1 & bool & 1 & 0 & 0 & 1 & --- & Auto-brake /
emergency trigger (CSV Row 5) \\
(unaccounted) & 4 & 4 & --- & --- & --- & --- & --- & --- & Byte 0 bits 4--7 ---
not enumerated in CSV \\
(unaccounted) & --- & 8 & --- & --- & --- & --- & --- & --- & Byte 1 --- not
enumerated in CSV \\
\texttt{VCU\_SEB\_Stroke\_Value\_Req} & 16 & 16 & u16 & 0.05 & -30 & -5 & 27 &
mm & Requested stroke position (CSV Row 6) \\
\texttt{VCU\_SEB\_Pre\_Value\_Req} & 24 & 8 & u8 & 0.05 & 0 & 0 & 5 & MPa &
Requested pressure (CSV Row 7). Raw = kPa × 0.02. Overlaps Stroke{[}15:8{]} at
byte 3 --- mode-dependent: Stroke uses full 16-bit in Mode 0, byte 3 carries
pressure in Mode 1. \\
(unaccounted) & 32 & 16 & --- & --- & --- & --- & --- & --- & Bytes 4--5 --- not
enumerated in CSV \\
\texttt{VCU\_SEB\_RollCnt\_Enable} & 48 & 1 & bool & 1 & 0 & 0 & 1 & --- & Life
Signal Validity --- \textbf{Must be 1} (CSV Row 8) \\
\texttt{VCU\_SEB\_CheckSum\_Enable} & 49 & 1 & bool & 1 & 0 & 0 & 1 & --- &
Checksum Validity --- \textbf{Must be 1} (CSV Row 9) \\
(unaccounted) & 50 & 2 & --- & --- & --- & --- & --- & --- & Byte 6 bits 2--3
--- not enumerated in CSV \\
\texttt{VCU\_SEB\_RollCnt} & 52 & 4 & u8 & 1 & 0 & 0 & 15 & --- & Life Signal
rolling counter. Increment every frame. (CSV Row 10) \\
\texttt{VCU\_SEB\_CheckSum} & 56 & 8 & u8 & 1 & 0 & 0 & 255 & --- & Checksum =
XOR(bytes 0--6) \^{} 0xFF (CSV Row 11) \\
\end{longtable}

\textbf{Byte layout} (little-endian, CSV as source of truth):

\begin{longtable}[]{@{}
  >{\raggedright\arraybackslash}p{(\columnwidth - 16\tabcolsep) * \real{0.2000}}
  >{\raggedright\arraybackslash}p{(\columnwidth - 16\tabcolsep) * \real{0.1000}}
  >{\raggedright\arraybackslash}p{(\columnwidth - 16\tabcolsep) * \real{0.1000}}
  >{\raggedright\arraybackslash}p{(\columnwidth - 16\tabcolsep) * \real{0.1000}}
  >{\raggedright\arraybackslash}p{(\columnwidth - 16\tabcolsep) * \real{0.1000}}
  >{\raggedright\arraybackslash}p{(\columnwidth - 16\tabcolsep) * \real{0.1000}}
  >{\raggedright\arraybackslash}p{(\columnwidth - 16\tabcolsep) * \real{0.1000}}
  >{\raggedright\arraybackslash}p{(\columnwidth - 16\tabcolsep) * \real{0.1000}}
  >{\raggedright\arraybackslash}p{(\columnwidth - 16\tabcolsep) * \real{0.1000}}@{}}
\toprule\noalign{}
\begin{minipage}[b]{\linewidth}\raggedright
Byte
\end{minipage} & \begin{minipage}[b]{\linewidth}\raggedright
0
\end{minipage} & \begin{minipage}[b]{\linewidth}\raggedright
1
\end{minipage} & \begin{minipage}[b]{\linewidth}\raggedright
2
\end{minipage} & \begin{minipage}[b]{\linewidth}\raggedright
3
\end{minipage} & \begin{minipage}[b]{\linewidth}\raggedright
4
\end{minipage} & \begin{minipage}[b]{\linewidth}\raggedright
5
\end{minipage} & \begin{minipage}[b]{\linewidth}\raggedright
6
\end{minipage} & \begin{minipage}[b]{\linewidth}\raggedright
7
\end{minipage} \\
\midrule\noalign{}
\endhead
\bottomrule\noalign{}
\endlastfoot
Content & Align{[}0{]}+CtrlEn{[}1{]}+Mode{[}2{]}+AutoBrk{[}3{]} & (unaccounted)
& Stroke\_Req {[}7:0{]} & Stroke\_Req {[}15:8{]} / Pre\_Req {[}7:0{]}
(mode-muxed) & (unaccounted) & (unaccounted) &
RollCntEn{[}0{]}+CksEn{[}1{]}+(gap)+RollCnt{[}4:7{]} & CheckSum \\
\end{longtable}

\begin{quote}
\textbf{Byte 3 multiplexing:} In Stroke Mode (Mode=0), bytes 2--3 carry the
16-bit stroke value (\texttt{VCU\_SEB\_Stroke\_Value\_Req}). In Pressure Mode
(Mode=1), byte 3 carries the 8-bit pressure value
(\texttt{VCU\_SEB\_Pre\_Value\_Req}). Both signals are declared in the CSV at
overlapping positions --- the SEB interprets byte 3 based on the active mode
bit. Bytes 1, 4, and 5 are not enumerated in the CSV; the SEB may ignore them or
use them for undocumented functions.
\end{quote}

\textbf{Stroke conversion}: \texttt{raw\ =\ (physical\_mm\ +\ 30.0)\ /\ 0.05}

\begin{longtable}[]{@{}lll@{}}
\toprule\noalign{}
Physical & Raw & Use case \\
\midrule\noalign{}
\endhead
\bottomrule\noalign{}
\endlastfoot
-5 mm & 500 & Min \\
0 mm & 600 & Released \\
15 mm & 900 & Manual lever pressed \\
27 mm & 1140 & ESTOP full brake \\
\end{longtable}

\textbf{Security}: Rolling counter must increment 0→15 every frame. Same value
twice → SEB rejects (assumes frozen controller). Checksum =
\texttt{XOR(bytes{[}0..6{]})\ \^{}\ 0xFF} (verify against SYNTREE spec).

\textbf{Mode 0 (Stroke)}: Command a specific pushrod position in mm. Best for
mimicking pedal travel / ESTOP full brake / manual lever. \textbf{Mode 1
(Pressure)}: Command hydraulic pressure in MPa. SEB's internal PID maintains
target. Best for autonomous deceleration control (compensates for pad wear,
temperature).

\begin{center}\rule{0.5\linewidth}{0.5pt}\end{center}

\subsection{0x721 --- SEB\_STATUS (SYNTREE SEB Brake
Feedback)}\label{x721-seb_status-syntree-seb-brake-feedback}

\begin{longtable}[]{@{}ll@{}}
\toprule\noalign{}
Property & Value \\
\midrule\noalign{}
\endhead
\bottomrule\noalign{}
\endlastfoot
\textbf{Sender} & SYNTREE SEB (brake module) \\
\textbf{Receiver(s)} & SYS ESP32-S3 \\
\textbf{DLC} & 8 \\
\textbf{Period} & 10 ms (100 Hz) \\
\textbf{Endianness} & Motorola LSB (little-endian) \\
\end{longtable}

\begin{longtable}[]{@{}
  >{\raggedright\arraybackslash}p{(\columnwidth - 18\tabcolsep) * \real{0.1081}}
  >{\raggedright\arraybackslash}p{(\columnwidth - 18\tabcolsep) * \real{0.1486}}
  >{\raggedright\arraybackslash}p{(\columnwidth - 18\tabcolsep) * \real{0.0676}}
  >{\raggedright\arraybackslash}p{(\columnwidth - 18\tabcolsep) * \real{0.0811}}
  >{\raggedright\arraybackslash}p{(\columnwidth - 18\tabcolsep) * \real{0.0946}}
  >{\raggedright\arraybackslash}p{(\columnwidth - 18\tabcolsep) * \real{0.1081}}
  >{\raggedright\arraybackslash}p{(\columnwidth - 18\tabcolsep) * \real{0.0676}}
  >{\raggedright\arraybackslash}p{(\columnwidth - 18\tabcolsep) * \real{0.0676}}
  >{\raggedright\arraybackslash}p{(\columnwidth - 18\tabcolsep) * \real{0.0811}}
  >{\raggedright\arraybackslash}p{(\columnwidth - 18\tabcolsep) * \real{0.1757}}@{}}
\toprule\noalign{}
\begin{minipage}[b]{\linewidth}\raggedright
Signal
\end{minipage} & \begin{minipage}[b]{\linewidth}\raggedright
Start bit
\end{minipage} & \begin{minipage}[b]{\linewidth}\raggedright
Len
\end{minipage} & \begin{minipage}[b]{\linewidth}\raggedright
Type
\end{minipage} & \begin{minipage}[b]{\linewidth}\raggedright
Scale
\end{minipage} & \begin{minipage}[b]{\linewidth}\raggedright
Offset
\end{minipage} & \begin{minipage}[b]{\linewidth}\raggedright
Min
\end{minipage} & \begin{minipage}[b]{\linewidth}\raggedright
Max
\end{minipage} & \begin{minipage}[b]{\linewidth}\raggedright
Unit
\end{minipage} & \begin{minipage}[b]{\linewidth}\raggedright
Description
\end{minipage} \\
\midrule\noalign{}
\endhead
\bottomrule\noalign{}
\endlastfoot
\texttt{SEB\_Alignment\_Status} & 0 & 1 & bool & 1 & 0 & 0 & 1 & --- & Alignment
Info Feedback. 1 = aligned. (CSV Row 12) \\
\texttt{SEB\_Control\_Enable\_Status} & 1 & 1 & bool & 1 & 0 & 0 & 1 & --- &
Control Enable Feedback (CSV Row 13) \\
\texttt{SEB\_Control\_Mode\_Status} & 2 & 2 & u8 & 1 & 0 & 0 & 3 & enum &
Control Mode Feedback: 0=?, 1=Stroke?, 2=Pressure?, 3=? (CSV Row 14) \\
\texttt{SEB\_AutoBrake\_Status} & 4 & 1 & bool & 1 & 0 & 0 & 1 & --- & Auto
Brake Status Feedback (CSV Row 15) \\
(unaccounted) & 5 & 1 & --- & --- & --- & --- & --- & --- & Byte 0 bit 5 --- not
enumerated in CSV \\
\texttt{SEB\_Error\_Status} & 6 & 2 & u8 & 1 & 0 & 0 & 3 & enum & 0=No fault,
1=L1 minor, 2=L2 general, 3=L3 severe (CSV Row 16) \\
(unaccounted) & --- & 8 & --- & --- & --- & --- & --- & --- & Byte 1 --- not
enumerated in CSV \\
\texttt{SEB\_Stroke\_Value} & 16 & 16 & u16 & 0.05 & -30 & -5 & 27 & mm & Stroke
Value Feedback (CSV Row 17) \\
\texttt{SEB\_Pressure\_Value} & 24 & 8 & u8 & 0.05 & 0 & 0 & 5 & MPa & Pressure
Value Feedback (CSV Row 18). Overlaps Stroke{[}15:8{]} at byte 3 ---
mode-dependent. \\
\texttt{SEB\_Angle\_Value} & 40 & 16 & i16 & 0.5 & 0 & -150 & 840 & --- & Angle
Feedback (CSV Row 19). Overlaps security echo bits at byte 6 --- see note
below. \\
\texttt{SEB\_RollCnt\_Enable\_Status} & 48 & 1 & bool & 1 & 0 & 0 & 1 & --- &
Life Signal Status Feedback (CSV Row 20). Overlaps Angle\_Value{[}15:8{]}. \\
\texttt{SEB\_CheckSum\_Enable\_Status} & 49 & 1 & bool & 1 & 0 & 0 & 1 & --- &
Checksum Status Feedback (CSV Row 21). Overlaps Angle\_Value{[}15:8{]}. \\
(unaccounted) & 50 & 2 & --- & --- & --- & --- & --- & --- & Byte 6 bits 2--3
--- not enumerated in CSV \\
\texttt{SEB\_RollCnt\_Status} & 52 & 4 & u8 & 1 & 0 & 0 & 15 & --- & Life Signal
Feedback --- echoes received rolling counter (CSV Row 22) \\
\texttt{SEB\_CheckSum\_Status} & 56 & 8 & u8 & 1 & 0 & 0 & 255 & --- & Checksum
Feedback (CSV Row 23) \\
\end{longtable}

\textbf{Byte layout} (little-endian, CSV as source of truth):

\begin{longtable}[]{@{}
  >{\raggedright\arraybackslash}p{(\columnwidth - 16\tabcolsep) * \real{0.2000}}
  >{\raggedright\arraybackslash}p{(\columnwidth - 16\tabcolsep) * \real{0.1000}}
  >{\raggedright\arraybackslash}p{(\columnwidth - 16\tabcolsep) * \real{0.1000}}
  >{\raggedright\arraybackslash}p{(\columnwidth - 16\tabcolsep) * \real{0.1000}}
  >{\raggedright\arraybackslash}p{(\columnwidth - 16\tabcolsep) * \real{0.1000}}
  >{\raggedright\arraybackslash}p{(\columnwidth - 16\tabcolsep) * \real{0.1000}}
  >{\raggedright\arraybackslash}p{(\columnwidth - 16\tabcolsep) * \real{0.1000}}
  >{\raggedright\arraybackslash}p{(\columnwidth - 16\tabcolsep) * \real{0.1000}}
  >{\raggedright\arraybackslash}p{(\columnwidth - 16\tabcolsep) * \real{0.1000}}@{}}
\toprule\noalign{}
\begin{minipage}[b]{\linewidth}\raggedright
Byte
\end{minipage} & \begin{minipage}[b]{\linewidth}\raggedright
0
\end{minipage} & \begin{minipage}[b]{\linewidth}\raggedright
1
\end{minipage} & \begin{minipage}[b]{\linewidth}\raggedright
2
\end{minipage} & \begin{minipage}[b]{\linewidth}\raggedright
3
\end{minipage} & \begin{minipage}[b]{\linewidth}\raggedright
4
\end{minipage} & \begin{minipage}[b]{\linewidth}\raggedright
5
\end{minipage} & \begin{minipage}[b]{\linewidth}\raggedright
6
\end{minipage} & \begin{minipage}[b]{\linewidth}\raggedright
7
\end{minipage} \\
\midrule\noalign{}
\endhead
\bottomrule\noalign{}
\endlastfoot
Content &
Align{[}0{]}+CtrlEn{[}1{]}+Mode{[}2:3{]}+AutoBrk{[}4{]}+(gap)+Error{[}6:7{]} &
(unacc.) & Stroke {[}7:0{]} & Stroke {[}15:8{]} / Pressure {[}7:0{]}
(mode-muxed) & (unacc.) & Angle {[}7:0{]} & Angle{[}15:8{]} /
RollCntEn{[}0{]}+CksEn{[}1{]}+(gap)+RollCnt{[}4:7{]} (overlap) & CksSum\_Stat \\
\end{longtable}

\begin{quote}
\textbf{Byte 3 multiplexing:} Same pattern as command frame ---
\texttt{SEB\_Stroke\_Value} uses full 16-bit at bytes 2--3 in Stroke Mode;
\texttt{SEB\_Pressure\_Value} uses byte 3 in Pressure Mode. The SEB reports
whichever is active.

\textbf{Byte 6 overlap:} CSV lists \texttt{SEB\_Angle\_Value} as 16-bit (bytes
5--6) AND security echo bits at byte 6 (bits 48--49, 52--55). These overlap. The
CSV (manufacturer DBC export) declares both --- they may represent different
firmware versions or the Angle\_Value may be 8-bit in practice (byte 5 only).
Trust the CSV's declaration and handle in firmware by reading Angle as 16-bit
from bytes 5--6, understanding that the upper byte may carry security echo data
in some SEB firmware revisions.
\end{quote}

\textbf{SYS usage}: Boot sync --- read \texttt{SEB\_Stroke\_Value} as initial
command target. Active --- confirm \texttt{SEB\_Alignment\_Status\ ==\ 1}.
\texttt{SEB\_Error\_Status\ \textgreater{}\ 0} → log and report via
\texttt{0x011}. Subscribe to \texttt{0x731\ SEB\_ErrInfo} for detailed fault
flags --- escalate L3 faults to ESTOP.

\begin{center}\rule{0.5\linewidth}{0.5pt}\end{center}

\subsection{0x731 --- SEB\_ErrInfo (SYNTREE SEB Error
Detail)}\label{x731-seb_errinfo-syntree-seb-error-detail}

\begin{longtable}[]{@{}ll@{}}
\toprule\noalign{}
Property & Value \\
\midrule\noalign{}
\endhead
\bottomrule\noalign{}
\endlastfoot
\textbf{Sender} & SYNTREE SEB (brake module) \\
\textbf{Receiver(s)} & SYS ESP32-S3 \\
\textbf{DLC} & 8 \\
\textbf{Period} & 100 ms (10 Hz) \\
\textbf{Endianness} & Motorola LSB (little-endian) \\
\end{longtable}

Detailed fault flags. Each bit is an independent fault indicator. 1 = fault
active.

\begin{longtable}[]{@{}
  >{\raggedright\arraybackslash}p{(\columnwidth - 10\tabcolsep) * \real{0.1429}}
  >{\raggedright\arraybackslash}p{(\columnwidth - 10\tabcolsep) * \real{0.1964}}
  >{\raggedright\arraybackslash}p{(\columnwidth - 10\tabcolsep) * \real{0.0893}}
  >{\raggedright\arraybackslash}p{(\columnwidth - 10\tabcolsep) * \real{0.1071}}
  >{\raggedright\arraybackslash}p{(\columnwidth - 10\tabcolsep) * \real{0.2321}}
  >{\raggedright\arraybackslash}p{(\columnwidth - 10\tabcolsep) * \real{0.2321}}@{}}
\toprule\noalign{}
\begin{minipage}[b]{\linewidth}\raggedright
Signal
\end{minipage} & \begin{minipage}[b]{\linewidth}\raggedright
Start bit
\end{minipage} & \begin{minipage}[b]{\linewidth}\raggedright
Len
\end{minipage} & \begin{minipage}[b]{\linewidth}\raggedright
Type
\end{minipage} & \begin{minipage}[b]{\linewidth}\raggedright
Fault Level
\end{minipage} & \begin{minipage}[b]{\linewidth}\raggedright
Description
\end{minipage} \\
\midrule\noalign{}
\endhead
\bottomrule\noalign{}
\endlastfoot
\texttt{SEB\_ECUUnderVolt\_Err} & 0 & 1 & bool & L2 & Controller Undervoltage \\
\texttt{SEB\_ECUOverVolt\_Err} & 1 & 1 & bool & L2 & Controller Overvoltage \\
\texttt{SEB\_CanCom\_Err} & 2 & 1 & bool & \textbf{L3} & CAN Communication
Fault \\
\texttt{SEB\_ECUTemp\_Err} & 3 & 1 & bool & \textbf{L3} & Controller Temperature
Fault \\
\texttt{SEB\_Domain\_drive\_SC\_Err} & 4 & 1 & bool & \textbf{L3} & Domain Drive
Short Circuit \\
\texttt{SEB\_Domain\_drive\_V\_Err} & 5 & 1 & bool & \textbf{L3} & Domain Drive
Voltage Fault \\
\texttt{SEB\_Domain\_drive\_T\_Err} & 6 & 1 & bool & \textbf{L3} & Domain Drive
Temperature Fault \\
\texttt{SEB\_AngleSensor\_P\_OC\_Err} & 7 & 1 & bool & \textbf{L3} & Angle
Sensor P Open Circuit \\
\texttt{SEB\_AngleSensor\_P\_AF\_Err} & 8 & 1 & bool & \textbf{L3} & Angle
Sensor P Mainboard Abnormal \\
\texttt{SEB\_AngleSensor\_S\_OC\_Err} & 9 & 1 & bool & \textbf{L3} & Angle
Sensor S Open Circuit \\
\texttt{SEB\_AngleSensor\_S\_AF\_Err} & 10 & 1 & bool & \textbf{L3} & Angle
Sensor S Sub-board Abnormal \\
\texttt{SEB\_NOPreSensor\_Err} & 11 & 1 & bool & \textbf{L3} & Unconnected Oil
Pressure Sensor \\
(reserved) & 12 & 1 & --- & --- & \\
\texttt{SEB\_SensorUCL\_Err} & 13 & 1 & bool & \textbf{L3} & Sensor Plausibility
Fault \\
\texttt{SEB\_Alignment\_Err} & 14 & 1 & bool & L2 & Alignment Fault \\
\texttt{SEB\_AngleOver\_Err} & 15 & 1 & bool & L2 & Angle Out of Bounds \\
(reserved) & 16 & 1 & --- & --- & \\
\texttt{SEB\_Mtr\_Stall\_Err} & 17 & 1 & bool & \textbf{L3} & Motor Stall
Fault \\
\texttt{SEB\_MtrDC\_Err} & 18 & 1 & bool & \textbf{L3} & Motor Disconnect
Fault \\
\texttt{SEB\_Oil\_Err} & 19 & 1 & bool & L2 & Oil Pressure Error \\
\texttt{SEB\_InitOil\_Err} & 20 & 1 & bool & \textbf{L3} & Initial Oil Pressure
Fault \\
\texttt{SEB\_SentValue\_Err} & 21 & 1 & bool & \textbf{L3} & Send Value Error \\
\texttt{SEB\_Mtr\_NoLoad\_Err} & 22 & 1 & bool & \textbf{L3} & Motor No-load
Fault \\
(reserved) & 23 & 1 & --- & --- & \\
\texttt{SEB\_PreSensorOver\_Err} & 24 & 1 & bool & L2 & Oil Pressure Sensor
Overvoltage \\
\texttt{SEB\_LowVolt\_Charging\_Err} & 25 & 1 & bool & L2 & Low Voltage Charging
Failure \\
(reserved) & 26 & 38 & --- & --- & Bits 26--63 (byte 3 bits 2--7 + bytes 4--7).
Not enumerated in CSV. \\
\end{longtable}

\textbf{Byte layout} (little-endian):

\begin{longtable}[]{@{}
  >{\raggedright\arraybackslash}p{(\columnwidth - 16\tabcolsep) * \real{0.2000}}
  >{\raggedright\arraybackslash}p{(\columnwidth - 16\tabcolsep) * \real{0.1000}}
  >{\raggedright\arraybackslash}p{(\columnwidth - 16\tabcolsep) * \real{0.1000}}
  >{\raggedright\arraybackslash}p{(\columnwidth - 16\tabcolsep) * \real{0.1000}}
  >{\raggedright\arraybackslash}p{(\columnwidth - 16\tabcolsep) * \real{0.1000}}
  >{\raggedright\arraybackslash}p{(\columnwidth - 16\tabcolsep) * \real{0.1000}}
  >{\raggedright\arraybackslash}p{(\columnwidth - 16\tabcolsep) * \real{0.1000}}
  >{\raggedright\arraybackslash}p{(\columnwidth - 16\tabcolsep) * \real{0.1000}}
  >{\raggedright\arraybackslash}p{(\columnwidth - 16\tabcolsep) * \real{0.1000}}@{}}
\toprule\noalign{}
\begin{minipage}[b]{\linewidth}\raggedright
Byte
\end{minipage} & \begin{minipage}[b]{\linewidth}\raggedright
0
\end{minipage} & \begin{minipage}[b]{\linewidth}\raggedright
1
\end{minipage} & \begin{minipage}[b]{\linewidth}\raggedright
2
\end{minipage} & \begin{minipage}[b]{\linewidth}\raggedright
3
\end{minipage} & \begin{minipage}[b]{\linewidth}\raggedright
4
\end{minipage} & \begin{minipage}[b]{\linewidth}\raggedright
5
\end{minipage} & \begin{minipage}[b]{\linewidth}\raggedright
6
\end{minipage} & \begin{minipage}[b]{\linewidth}\raggedright
7
\end{minipage} \\
\midrule\noalign{}
\endhead
\bottomrule\noalign{}
\endlastfoot
Content & faults {[}7:0{]} & faults {[}15:8{]} & faults {[}23:16{]} & faults
{[}31:24{]} & rsvd & rsvd & rsvd & rsvd \\
\end{longtable}

\begin{quote}
\textbf{Safety note:} 14 of 23 documented faults are L3 (severe, request
shutdown). SYS must subscribe to this message and escalate any L3 fault to ESTOP
via \texttt{0x001}. The summary \texttt{SEB\_Error\_Status} in \texttt{0x721}
only provides a 2-bit aggregate level; this message reveals \emph{which} sensor
or subsystem failed and at what severity.
\end{quote}

\begin{center}\rule{0.5\linewidth}{0.5pt}\end{center}

\subsection{0x741 --- SEB\_Version (SYNTREE SEB Firmware
Version)}\label{x741-seb_version-syntree-seb-firmware-version}

\begin{longtable}[]{@{}ll@{}}
\toprule\noalign{}
Property & Value \\
\midrule\noalign{}
\endhead
\bottomrule\noalign{}
\endlastfoot
\textbf{Sender} & SYNTREE SEB (brake module) \\
\textbf{Receiver(s)} & SYS ESP32-S3 \\
\textbf{DLC} & 8 \\
\textbf{Period} & 1000 ms (1 Hz) \\
\textbf{Endianness} & Motorola LSB (little-endian) \\
\end{longtable}

\begin{longtable}[]{@{}
  >{\raggedright\arraybackslash}p{(\columnwidth - 18\tabcolsep) * \real{0.1081}}
  >{\raggedright\arraybackslash}p{(\columnwidth - 18\tabcolsep) * \real{0.1486}}
  >{\raggedright\arraybackslash}p{(\columnwidth - 18\tabcolsep) * \real{0.0676}}
  >{\raggedright\arraybackslash}p{(\columnwidth - 18\tabcolsep) * \real{0.0811}}
  >{\raggedright\arraybackslash}p{(\columnwidth - 18\tabcolsep) * \real{0.0946}}
  >{\raggedright\arraybackslash}p{(\columnwidth - 18\tabcolsep) * \real{0.1081}}
  >{\raggedright\arraybackslash}p{(\columnwidth - 18\tabcolsep) * \real{0.0676}}
  >{\raggedright\arraybackslash}p{(\columnwidth - 18\tabcolsep) * \real{0.0676}}
  >{\raggedright\arraybackslash}p{(\columnwidth - 18\tabcolsep) * \real{0.0811}}
  >{\raggedright\arraybackslash}p{(\columnwidth - 18\tabcolsep) * \real{0.1757}}@{}}
\toprule\noalign{}
\begin{minipage}[b]{\linewidth}\raggedright
Signal
\end{minipage} & \begin{minipage}[b]{\linewidth}\raggedright
Start bit
\end{minipage} & \begin{minipage}[b]{\linewidth}\raggedright
Len
\end{minipage} & \begin{minipage}[b]{\linewidth}\raggedright
Type
\end{minipage} & \begin{minipage}[b]{\linewidth}\raggedright
Scale
\end{minipage} & \begin{minipage}[b]{\linewidth}\raggedright
Offset
\end{minipage} & \begin{minipage}[b]{\linewidth}\raggedright
Min
\end{minipage} & \begin{minipage}[b]{\linewidth}\raggedright
Max
\end{minipage} & \begin{minipage}[b]{\linewidth}\raggedright
Unit
\end{minipage} & \begin{minipage}[b]{\linewidth}\raggedright
Description
\end{minipage} \\
\midrule\noalign{}
\endhead
\bottomrule\noalign{}
\endlastfoot
\texttt{SEB\_SW\_Version} & 0 & 8 & u8 & 0.01 & 0 & 0 & 25.5 & --- & Software
version (e.g., 0xC8 = 2.00) \\
\texttt{SEB\_HW\_Version} & 8 & 8 & u8 & 0.1 & 0 & 0 & 25.5 & --- & Hardware
version (e.g., 0x0D = 1.3) \\
(reserved) & 16 & 48 & --- & --- & --- & --- & --- & --- & Bytes 2--7 \\
\end{longtable}

\textbf{SYS usage}: Log on boot for compatibility check and field diagnostics.
Report via \texttt{0x600\ SYS\_DIAG\_RPT}.

\begin{center}\rule{0.5\linewidth}{0.5pt}\end{center}

\subsection{0x6FB --- SEB\_Test (SYNTREE SEB
Telemetry)}\label{x6fb-seb_test-syntree-seb-telemetry}

\begin{longtable}[]{@{}ll@{}}
\toprule\noalign{}
Property & Value \\
\midrule\noalign{}
\endhead
\bottomrule\noalign{}
\endlastfoot
\textbf{Sender} & SYNTREE SEB (brake module) \\
\textbf{Receiver(s)} & SYS ESP32-S3 \\
\textbf{DLC} & 8 \\
\textbf{Period} & 10 ms (100 Hz) \\
\textbf{Endianness} & Motorola LSB (little-endian) \\
\end{longtable}

\begin{longtable}[]{@{}
  >{\raggedright\arraybackslash}p{(\columnwidth - 18\tabcolsep) * \real{0.1081}}
  >{\raggedright\arraybackslash}p{(\columnwidth - 18\tabcolsep) * \real{0.1486}}
  >{\raggedright\arraybackslash}p{(\columnwidth - 18\tabcolsep) * \real{0.0676}}
  >{\raggedright\arraybackslash}p{(\columnwidth - 18\tabcolsep) * \real{0.0811}}
  >{\raggedright\arraybackslash}p{(\columnwidth - 18\tabcolsep) * \real{0.0946}}
  >{\raggedright\arraybackslash}p{(\columnwidth - 18\tabcolsep) * \real{0.1081}}
  >{\raggedright\arraybackslash}p{(\columnwidth - 18\tabcolsep) * \real{0.0676}}
  >{\raggedright\arraybackslash}p{(\columnwidth - 18\tabcolsep) * \real{0.0676}}
  >{\raggedright\arraybackslash}p{(\columnwidth - 18\tabcolsep) * \real{0.0811}}
  >{\raggedright\arraybackslash}p{(\columnwidth - 18\tabcolsep) * \real{0.1757}}@{}}
\toprule\noalign{}
\begin{minipage}[b]{\linewidth}\raggedright
Signal
\end{minipage} & \begin{minipage}[b]{\linewidth}\raggedright
Start bit
\end{minipage} & \begin{minipage}[b]{\linewidth}\raggedright
Len
\end{minipage} & \begin{minipage}[b]{\linewidth}\raggedright
Type
\end{minipage} & \begin{minipage}[b]{\linewidth}\raggedright
Scale
\end{minipage} & \begin{minipage}[b]{\linewidth}\raggedright
Offset
\end{minipage} & \begin{minipage}[b]{\linewidth}\raggedright
Min
\end{minipage} & \begin{minipage}[b]{\linewidth}\raggedright
Max
\end{minipage} & \begin{minipage}[b]{\linewidth}\raggedright
Unit
\end{minipage} & \begin{minipage}[b]{\linewidth}\raggedright
Description
\end{minipage} \\
\midrule\noalign{}
\endhead
\bottomrule\noalign{}
\endlastfoot
(reserved) & 0 & 8 & --- & --- & --- & --- & --- & --- & Byte 0 \\
\texttt{SEB\_MtrCurr} & 8 & 16 & i16 & 0.0078125 & 0 & -255 & 255 & A & Motor
current. Bytes 1--2. ±1.99 A range. \\
\texttt{SEB\_ECUTemp} & 24 & 16 & u16 & 0.5 & 0 & -40 & 215 & °C & ECU
temperature. Bytes 3--4. \\
\texttt{SEB\_PowVolt} & 40 & 16 & u16 & 0.00390625 & 0 & 0 & 32 & V & Power
supply voltage. Bytes 5--6. \\
(reserved) & 56 & 8 & --- & --- & --- & --- & --- & --- & Byte 7 \\
\end{longtable}

\begin{quote}
\textbf{Note:} CSV uses byte-local start-bit numbering for this message (start
bit = offset within start byte). Converted to absolute Motorola LSB above.
\textbf{SYS usage:} Monitor \texttt{SEB\_MtrCurr} for mechanical binding / motor
degradation trends. \texttt{SEB\_ECUTemp} for over-temperature early warning.
\end{quote}

\begin{center}\rule{0.5\linewidth}{0.5pt}\end{center}

\subsection{0x7FD --- RT\_HEARTBEAT
(low-level)}\label{x7fd-rt_heartbeat-low-level}

\begin{longtable}[]{@{}ll@{}}
\toprule\noalign{}
Property & Value \\
\midrule\noalign{}
\endhead
\bottomrule\noalign{}
\endlastfoot
\textbf{Sender} & RT \\
\textbf{Receiver(s)} & SYS \\
\textbf{DLC} & 1 \\
\textbf{Period} & 2 Hz (500 ms) \\
\textbf{Timeout} & 1000ms (2 missed frames) → SYS triggers ESTOP (AUTO only) \\
\end{longtable}

\begin{longtable}[]{@{}
  >{\raggedright\arraybackslash}p{(\columnwidth - 8\tabcolsep) * \real{0.1860}}
  >{\raggedright\arraybackslash}p{(\columnwidth - 8\tabcolsep) * \real{0.2558}}
  >{\raggedright\arraybackslash}p{(\columnwidth - 8\tabcolsep) * \real{0.1163}}
  >{\raggedright\arraybackslash}p{(\columnwidth - 8\tabcolsep) * \real{0.1395}}
  >{\raggedright\arraybackslash}p{(\columnwidth - 8\tabcolsep) * \real{0.3023}}@{}}
\toprule\noalign{}
\begin{minipage}[b]{\linewidth}\raggedright
Signal
\end{minipage} & \begin{minipage}[b]{\linewidth}\raggedright
Start bit
\end{minipage} & \begin{minipage}[b]{\linewidth}\raggedright
Len
\end{minipage} & \begin{minipage}[b]{\linewidth}\raggedright
Type
\end{minipage} & \begin{minipage}[b]{\linewidth}\raggedright
Description
\end{minipage} \\
\midrule\noalign{}
\endhead
\bottomrule\noalign{}
\endlastfoot
\texttt{alive\_ctr} & 0 & 8 & u8 & Increments every frame (wraps at 255). Frozen
= hung RT. \\
\end{longtable}

\texttt{0x204} staleness check at 200ms provides faster detection. RT sends
\texttt{0x7FD} independently on both buses (per-bus, NOT bridged; separate
counters per bus).

\begin{center}\rule{0.5\linewidth}{0.5pt}\end{center}

\subsection{0x7FE --- SYS\_HEARTBEAT
(low-level)}\label{x7fe-sys_heartbeat-low-level}

\begin{longtable}[]{@{}
  >{\raggedright\arraybackslash}p{(\columnwidth - 2\tabcolsep) * \real{0.5882}}
  >{\raggedright\arraybackslash}p{(\columnwidth - 2\tabcolsep) * \real{0.4118}}@{}}
\toprule\noalign{}
\begin{minipage}[b]{\linewidth}\raggedright
Property
\end{minipage} & \begin{minipage}[b]{\linewidth}\raggedright
Value
\end{minipage} \\
\midrule\noalign{}
\endhead
\bottomrule\noalign{}
\endlastfoot
\textbf{Sender} & SYS \\
\textbf{Receiver(s)} & RT \\
\textbf{DLC} & 1 \\
\textbf{Period} & 10 Hz (100 ms) \\
\textbf{Timeout} & 200ms (2 missed frames at 10 Hz) → RT takes over
\texttt{0x7B9} (stroke=max) + sends CAN \texttt{0x001}. Total brake gap
≤220ms. \\
\end{longtable}

\begin{longtable}[]{@{}
  >{\raggedright\arraybackslash}p{(\columnwidth - 8\tabcolsep) * \real{0.1860}}
  >{\raggedright\arraybackslash}p{(\columnwidth - 8\tabcolsep) * \real{0.2558}}
  >{\raggedright\arraybackslash}p{(\columnwidth - 8\tabcolsep) * \real{0.1163}}
  >{\raggedright\arraybackslash}p{(\columnwidth - 8\tabcolsep) * \real{0.1395}}
  >{\raggedright\arraybackslash}p{(\columnwidth - 8\tabcolsep) * \real{0.3023}}@{}}
\toprule\noalign{}
\begin{minipage}[b]{\linewidth}\raggedright
Signal
\end{minipage} & \begin{minipage}[b]{\linewidth}\raggedright
Start bit
\end{minipage} & \begin{minipage}[b]{\linewidth}\raggedright
Len
\end{minipage} & \begin{minipage}[b]{\linewidth}\raggedright
Type
\end{minipage} & \begin{minipage}[b]{\linewidth}\raggedright
Description
\end{minipage} \\
\midrule\noalign{}
\endhead
\bottomrule\noalign{}
\endlastfoot
\texttt{alive\_ctr} & 0 & 8 & u8 & Increments every frame (wraps at 255). Frozen
= hung SYS. \\
\end{longtable}

SYS heartbeat never leaves low bus. Startup grace period: 3 seconds (both
heartbeat monitors).

\begin{center}\rule{0.5\linewidth}{0.5pt}\end{center}

\section{2. High-Level CAN Bus}\label{high-level-can-bus}

Nodes: Jetson Orin, RT ESP32-S3 (MCP2515 SPI).

\begin{center}\rule{0.5\linewidth}{0.5pt}\end{center}

\subsection{0x001 --- SAFETY\_ESTOP
(high-level)}\label{x001-safety_estop-high-level}

\begin{longtable}[]{@{}ll@{}}
\toprule\noalign{}
Property & Value \\
\midrule\noalign{}
\endhead
\bottomrule\noalign{}
\endlastfoot
\textbf{Sender} & Jetson or RT (fwd from low) \\
\textbf{Receiver(s)} & Jetson, RT \\
\textbf{DLC} & 0 \\
\textbf{Period} & On event \\
\end{longtable}

Bridged by RT between buses.

\begin{center}\rule{0.5\linewidth}{0.5pt}\end{center}

\subsection{0x011 --- SYS\_SAFETY\_STS
(forwarded)}\label{x011-sys_safety_sts-forwarded}

Forwarded from low-level by RT. Same payload layout as §1 \texttt{0x011}.

\begin{center}\rule{0.5\linewidth}{0.5pt}\end{center}

\subsection{0x120 --- SYS\_THROTTLE\_STS
(forwarded)}\label{x120-sys_throttle_sts-forwarded}

Forwarded from low-level by RT. Same payload layout as §1 \texttt{0x120}.

\begin{center}\rule{0.5\linewidth}{0.5pt}\end{center}

\subsection{0x210 --- RT\_STATE\_RPT}\label{x210-rt_state_rpt}

\begin{longtable}[]{@{}ll@{}}
\toprule\noalign{}
Property & Value \\
\midrule\noalign{}
\endhead
\bottomrule\noalign{}
\endlastfoot
\textbf{Sender} & RT \\
\textbf{Receiver(s)} & Jetson \\
\textbf{DLC} & 3 \\
\textbf{Period} & 10 Hz \\
\end{longtable}

\begin{longtable}[]{@{}lllllll@{}}
\toprule\noalign{}
Signal & Start bit & Len & Type & Min & Max & Unit \\
\midrule\noalign{}
\endhead
\bottomrule\noalign{}
\endlastfoot
\texttt{RT\_Mode} & 0 & 8 & u8 (enum) & 0 & 2 & --- \\
\texttt{RT\_SteerValid} & 8 & 8 & u8 (bool) & 0 & 1 & --- \\
\texttt{RT\_Reversing} & 16 & 8 & u8 (bool) & 0 & 1 & --- \\
\end{longtable}

Byte layout (big-endian): Byte 0=mode, 1=steer\_valid, 2=reversing.

\begin{center}\rule{0.5\linewidth}{0.5pt}\end{center}

\subsection{0x220 --- RT\_PID\_RPT}\label{x220-rt_pid_rpt}

\begin{longtable}[]{@{}ll@{}}
\toprule\noalign{}
Property & Value \\
\midrule\noalign{}
\endhead
\bottomrule\noalign{}
\endlastfoot
\textbf{Sender} & RT \\
\textbf{Receiver(s)} & Jetson \\
\textbf{DLC} & 6 \\
\textbf{Period} & 10 Hz \\
\end{longtable}

\begin{longtable}[]{@{}lllll@{}}
\toprule\noalign{}
Signal & Start bit & Len & Type & Unit \\
\midrule\noalign{}
\endhead
\bottomrule\noalign{}
\endlastfoot
\texttt{RT\_PidSetpoint} & 0 & 16 & i16 & mm/s \\
\texttt{RT\_PidMeasured} & 16 & 16 & i16 & mm/s \\
\texttt{RT\_PidOutput} & 32 & 16 & i16 & --- \\
\end{longtable}

Byte layout (big-endian): Bytes 0-1=sp, 2-3=meas, 4-5=out.

\begin{center}\rule{0.5\linewidth}{0.5pt}\end{center}

\subsection{0x300 --- HOST\_DRIVE\_CMD}\label{x300-host_drive_cmd}

\begin{longtable}[]{@{}ll@{}}
\toprule\noalign{}
Property & Value \\
\midrule\noalign{}
\endhead
\bottomrule\noalign{}
\endlastfoot
\textbf{Sender} & Jetson \\
\textbf{Receiver(s)} & RT \\
\textbf{DLC} & 8 \\
\textbf{Period} & ≤100 Hz \\
\end{longtable}

\begin{longtable}[]{@{}lllllllll@{}}
\toprule\noalign{}
Signal & Start bit & Len & Type & Scale & Offset & Min & Max & Unit \\
\midrule\noalign{}
\endhead
\bottomrule\noalign{}
\endlastfoot
\texttt{HOST\_DriveSpeed} & 0 & 32 & i32 & 1 & 0 & -500 & 3000 & mm/s \\
\texttt{HOST\_YawRate} & 32 & 24 & i24 & 1 & 0 & -3000 & 3000 & mrad/s \\
\texttt{HOST\_Gear} & 56 & 8 & u8 & 1 & 0 & 0 & 3 & enum (0=N,1=D,2=S,3=R) \\
\end{longtable}

Byte layout (big-endian): Bytes 0-3=speed, 4-6=yaw{[}23:0{]}, 7=gear.

ROS 2 conversion: \texttt{speed\_mmps\ =\ linear.x\ ×\ 1000},
\texttt{yaw\_rate\_mrad\_s\ =\ angular.z\ ×\ 1000}.

\begin{center}\rule{0.5\linewidth}{0.5pt}\end{center}

\subsection{0x301 --- HOST\_BRAKE\_REQ}\label{x301-host_brake_req}

\begin{longtable}[]{@{}ll@{}}
\toprule\noalign{}
Property & Value \\
\midrule\noalign{}
\endhead
\bottomrule\noalign{}
\endlastfoot
\textbf{Sender} & Jetson \\
\textbf{Receiver(s)} & RT \\
\textbf{DLC} & 4 \\
\textbf{Period} & On demand \\
\end{longtable}

\begin{longtable}[]{@{}lllll@{}}
\toprule\noalign{}
Signal & Start bit & Len & Type & Unit \\
\midrule\noalign{}
\endhead
\bottomrule\noalign{}
\endlastfoot
\texttt{HOST\_BrakePressure} & 0 & 32 & i32 & kPa \\
\end{longtable}

Byte layout (big-endian): Bytes 0-3. RT arbitrates: max(RT\_computed,
HOST\_request). Result forwarded to SYS via \texttt{0x205} RT\_BRAKE\_CMD (i32
kPa, 50 Hz).

\begin{center}\rule{0.5\linewidth}{0.5pt}\end{center}

\subsection{0x302 --- HOST\_LIGHT\_CMD}\label{x302-host_light_cmd}

\begin{longtable}[]{@{}ll@{}}
\toprule\noalign{}
Property & Value \\
\midrule\noalign{}
\endhead
\bottomrule\noalign{}
\endlastfoot
\textbf{Sender} & Jetson \\
\textbf{Receiver(s)} & RT (→ SYS) \\
\textbf{DLC} & 1 \\
\textbf{Period} & On change \\
\end{longtable}

Layout identical to low-level \texttt{0x302}. RT forwards transparently.

\begin{center}\rule{0.5\linewidth}{0.5pt}\end{center}

\subsection{0x400 --- RT\_OBSTACLE\_RPT}\label{x400-rt_obstacle_rpt}

\begin{longtable}[]{@{}ll@{}}
\toprule\noalign{}
Property & Value \\
\midrule\noalign{}
\endhead
\bottomrule\noalign{}
\endlastfoot
\textbf{Sender} & RT \\
\textbf{Receiver(s)} & Jetson \\
\textbf{DLC} & 4 \\
\textbf{Period} & 10 Hz \\
\end{longtable}

\begin{longtable}[]{@{}lllllll@{}}
\toprule\noalign{}
Signal & Start bit & Len & Type & Min & Max & Unit \\
\midrule\noalign{}
\endhead
\bottomrule\noalign{}
\endlastfoot
\texttt{RT\_ObstacleDistance} & 0 & 32 & u32 & 0 & 2³²−1 & mm \\
\end{longtable}

UINT32\_MAX = no reading / timeout.

RT reports min obstacle distance to Jetson at 10 Hz. Jetson perception
(LiDAR/camera) sends distance via internal path to RT, which rebroadcasts on the
high bus for logging/monitoring.

\begin{center}\rule{0.5\linewidth}{0.5pt}\end{center}

\subsection{0x600 --- SYS\_DIAG\_RPT
(forwarded)}\label{x600-sys_diag_rpt-forwarded}

Forwarded from low-level by RT. Same layout as §1 \texttt{0x600}.

\begin{center}\rule{0.5\linewidth}{0.5pt}\end{center}

\subsection{0x7FD --- RT\_HEARTBEAT
(high-level)}\label{x7fd-rt_heartbeat-high-level}

\begin{longtable}[]{@{}ll@{}}
\toprule\noalign{}
Property & Value \\
\midrule\noalign{}
\endhead
\bottomrule\noalign{}
\endlastfoot
\textbf{Sender} & RT \\
\textbf{Receiver(s)} & Jetson \\
\textbf{DLC} & 1 \\
\textbf{Period} & 2 Hz (500 ms) \\
\textbf{Timeout} & 1500ms (3 missed frames) → Jetson stops publishing
\texttt{/cmd\_vel} \\
\end{longtable}

\begin{longtable}[]{@{}
  >{\raggedright\arraybackslash}p{(\columnwidth - 8\tabcolsep) * \real{0.1860}}
  >{\raggedright\arraybackslash}p{(\columnwidth - 8\tabcolsep) * \real{0.2558}}
  >{\raggedright\arraybackslash}p{(\columnwidth - 8\tabcolsep) * \real{0.1163}}
  >{\raggedright\arraybackslash}p{(\columnwidth - 8\tabcolsep) * \real{0.1395}}
  >{\raggedright\arraybackslash}p{(\columnwidth - 8\tabcolsep) * \real{0.3023}}@{}}
\toprule\noalign{}
\begin{minipage}[b]{\linewidth}\raggedright
Signal
\end{minipage} & \begin{minipage}[b]{\linewidth}\raggedright
Start bit
\end{minipage} & \begin{minipage}[b]{\linewidth}\raggedright
Len
\end{minipage} & \begin{minipage}[b]{\linewidth}\raggedright
Type
\end{minipage} & \begin{minipage}[b]{\linewidth}\raggedright
Description
\end{minipage} \\
\midrule\noalign{}
\endhead
\bottomrule\noalign{}
\endlastfoot
\texttt{alive\_ctr} & 0 & 8 & u8 & Increments every frame (wraps at 255). Frozen
= hung RT. \\
\end{longtable}

RT sends \texttt{0x7FD} independently on both buses (per-bus, NOT bridged).

\begin{center}\rule{0.5\linewidth}{0.5pt}\end{center}

\subsection{0x7FC --- JETSON\_HEARTBEAT
(high-level)}\label{x7fc-jetson_heartbeat-high-level}

\begin{longtable}[]{@{}
  >{\raggedright\arraybackslash}p{(\columnwidth - 2\tabcolsep) * \real{0.5882}}
  >{\raggedright\arraybackslash}p{(\columnwidth - 2\tabcolsep) * \real{0.4118}}@{}}
\toprule\noalign{}
\begin{minipage}[b]{\linewidth}\raggedright
Property
\end{minipage} & \begin{minipage}[b]{\linewidth}\raggedright
Value
\end{minipage} \\
\midrule\noalign{}
\endhead
\bottomrule\noalign{}
\endlastfoot
\textbf{Sender} & Jetson \\
\textbf{Receiver(s)} & RT \\
\textbf{DLC} & 1 \\
\textbf{Period} & 2 Hz (500 ms) \\
\textbf{Timeout} & 1500ms (3 missed frames) → RT zeroes \texttt{0x204} + stops
\texttt{0x169} (controlled stop) \\
\end{longtable}

\begin{longtable}[]{@{}
  >{\raggedright\arraybackslash}p{(\columnwidth - 8\tabcolsep) * \real{0.1860}}
  >{\raggedright\arraybackslash}p{(\columnwidth - 8\tabcolsep) * \real{0.2558}}
  >{\raggedright\arraybackslash}p{(\columnwidth - 8\tabcolsep) * \real{0.1163}}
  >{\raggedright\arraybackslash}p{(\columnwidth - 8\tabcolsep) * \real{0.1395}}
  >{\raggedright\arraybackslash}p{(\columnwidth - 8\tabcolsep) * \real{0.3023}}@{}}
\toprule\noalign{}
\begin{minipage}[b]{\linewidth}\raggedright
Signal
\end{minipage} & \begin{minipage}[b]{\linewidth}\raggedright
Start bit
\end{minipage} & \begin{minipage}[b]{\linewidth}\raggedright
Len
\end{minipage} & \begin{minipage}[b]{\linewidth}\raggedright
Type
\end{minipage} & \begin{minipage}[b]{\linewidth}\raggedright
Description
\end{minipage} \\
\midrule\noalign{}
\endhead
\bottomrule\noalign{}
\endlastfoot
\texttt{alive\_ctr} & 0 & 8 & u8 & Increments every frame (wraps at 255). Frozen
= hung Jetson. \\
\end{longtable}

Jetson is QM, not safety-critical. Heartbeat loss triggers controlled stop, not
ESTOP.

\begin{center}\rule{0.5\linewidth}{0.5pt}\end{center}

\section{3. CAN ID Summary}\label{can-id-summary}

\subsection{Low-level bus}\label{low-level-bus}

\begin{longtable}[]{@{}llllll@{}}
\toprule\noalign{}
ID & Name & Sender & Receiver & DLC & Rate \\
\midrule\noalign{}
\endhead
\bottomrule\noalign{}
\endlastfoot
\texttt{0x001} & SAFETY\_ESTOP & RT, SYS & All & 0 & Event \\
\texttt{0x011} & SYS\_SAFETY\_STS & SYS & RT (→Jetson) & 2 & 5 Hz \\
\texttt{0x012} & SYS\_DCDC\_CMD & SYS & DC-DC & 1 & Change \\
\texttt{0x110} & SYS\_MODE\_CMD & SYS & RT & 1 & Change \\
\texttt{0x120} & SYS\_THROTTLE\_STS & MTR & RT (→Jetson) & 2 & 100 Hz \\
\texttt{0x169} & VCU\_SES\_REQ & RT & EPS-C & 8 & \textbf{50 Hz} \\
\texttt{0x201} & SES\_STATUS & EPS-C & RT & 8 & 100 Hz \\
\texttt{0x202} & SES\_ErrInfo & EPS-C & RT & 8 & 10 Hz \\
\texttt{0x203} & SES\_Version & EPS-C & RT & 8 & 1 Hz \\
\texttt{0x204} & RT\_DRIVE\_CMD & RT & SYS, MTR & 5 & 100 Hz \\
\texttt{0x205} & RT\_BRAKE\_CMD & RT & SYS & 4 & \textbf{50 Hz} \\
\texttt{0x302} & HOST\_LIGHT\_CMD & RT (fwd) & SYS & 1 & Change \\
\texttt{0x600} & SYS\_DIAG\_RPT & SYS & RT (→Jetson) & 8 & 1 Hz \\
\texttt{0x6FA} & SES\_Test & EPS-C & RT & 8 & 100 Hz \\
\texttt{0x6FB} & SEB\_Test & SEB & SYS & 8 & 100 Hz \\
\texttt{0x721} & SEB\_STATUS & SEB & SYS & 8 & 100 Hz \\
\texttt{0x731} & SEB\_ErrInfo & SEB & SYS & 8 & 10 Hz \\
\texttt{0x741} & SEB\_Version & SEB & SYS & 8 & 1 Hz \\
\texttt{0x7B9} & VCU\_SEB\_REQ & SYS & SEB & 8 & \textbf{50 Hz} \\
\texttt{0x7FD} & RT\_HEARTBEAT & RT & SYS & 1 & 2 Hz \\
\texttt{0x7FE} & SYS\_HEARTBEAT & SYS & RT & 1 & 10 Hz \\
\end{longtable}

\subsection{High-level bus}\label{high-level-bus}

\begin{longtable}[]{@{}llllll@{}}
\toprule\noalign{}
ID & Name & Sender & Receiver & DLC & Rate \\
\midrule\noalign{}
\endhead
\bottomrule\noalign{}
\endlastfoot
\texttt{0x001} & SAFETY\_ESTOP & Jetson, RT & Jetson, RT & 0 & Event \\
\texttt{0x011} & SYS\_SAFETY\_STS & RT (fwd) & Jetson & 2 & 5 Hz \\
\texttt{0x120} & SYS\_THROTTLE\_STS & RT (fwd) & Jetson & 2 & 100 Hz \\
\texttt{0x210} & RT\_STATE\_RPT & RT & Jetson & 3 & 10 Hz \\
\texttt{0x220} & RT\_PID\_RPT & RT & Jetson & 6 & --- (RESERVED, inactive) \\
\texttt{0x300} & HOST\_DRIVE\_CMD & Jetson & RT & 8 & ≤100 Hz \\
\texttt{0x301} & HOST\_BRAKE\_REQ & Jetson & RT & 4 & Demand \\
\texttt{0x302} & HOST\_LIGHT\_CMD & Jetson & RT (→SYS) & 1 & Change \\
\texttt{0x400} & RT\_OBSTACLE\_RPT & RT & Jetson & 4 & 10 Hz \\
\texttt{0x600} & SYS\_DIAG\_RPT & RT (fwd) & Jetson & 8 & 1 Hz \\
\texttt{0x7FD} & RT\_HEARTBEAT & RT & Jetson & 1 & 2 Hz \\
\texttt{0x7FC} & JETSON\_HEARTBEAT & Jetson & RT & 1 & 2 Hz \\
\end{longtable}

\begin{center}\rule{0.5\linewidth}{0.5pt}\end{center}

\section{4. Forwarding Rules (RT Gateway)}\label{forwarding-rules-rt-gateway}

RT is the only dual-bus node. Every CAN message falls into exactly one of three
categories:

\subsection{Category 1: Transparent forward (same ID, same
payload)}\label{category-1-transparent-forward-same-id-same-payload}

\begin{longtable}[]{@{}ll@{}}
\toprule\noalign{}
Direction & IDs \\
\midrule\noalign{}
\endhead
\bottomrule\noalign{}
\endlastfoot
Low → High & \texttt{0x001}, \texttt{0x011}, \texttt{0x120}, \texttt{0x206},
\texttt{0x600} \\
High → Low & \texttt{0x001}, \texttt{0x302} \\
\end{longtable}

\subsection{Category 2: Consumed by RT → different message
generated}\label{category-2-consumed-by-rt-different-message-generated}

\begin{longtable}[]{@{}
  >{\raggedright\arraybackslash}p{(\columnwidth - 6\tabcolsep) * \real{0.3103}}
  >{\raggedright\arraybackslash}p{(\columnwidth - 6\tabcolsep) * \real{0.1724}}
  >{\raggedright\arraybackslash}p{(\columnwidth - 6\tabcolsep) * \real{0.3448}}
  >{\raggedright\arraybackslash}p{(\columnwidth - 6\tabcolsep) * \real{0.1724}}@{}}
\toprule\noalign{}
\begin{minipage}[b]{\linewidth}\raggedright
Inbound
\end{minipage} & \begin{minipage}[b]{\linewidth}\raggedright
Bus
\end{minipage} & \begin{minipage}[b]{\linewidth}\raggedright
Outbound
\end{minipage} & \begin{minipage}[b]{\linewidth}\raggedright
Bus
\end{minipage} \\
\midrule\noalign{}
\endhead
\bottomrule\noalign{}
\endlastfoot
\texttt{0x300} HOST\_DRIVE\_CMD & High & \texttt{0x204} RT\_DRIVE\_CMD +
\texttt{0x169} VCU\_SES\_REQ & Low \\
\texttt{0x301} HOST\_BRAKE\_REQ & High & \texttt{0x205} RT\_BRAKE\_CMD & Low \\
\end{longtable}

\subsection{Category 3: Bus-local (never forwarded, never
regenerated)}\label{category-3-bus-local-never-forwarded-never-regenerated}

\begin{longtable}[]{@{}
  >{\raggedright\arraybackslash}p{(\columnwidth - 2\tabcolsep) * \real{0.5000}}
  >{\raggedright\arraybackslash}p{(\columnwidth - 2\tabcolsep) * \real{0.5000}}@{}}
\toprule\noalign{}
\begin{minipage}[b]{\linewidth}\raggedright
Bus
\end{minipage} & \begin{minipage}[b]{\linewidth}\raggedright
IDs
\end{minipage} \\
\midrule\noalign{}
\endhead
\bottomrule\noalign{}
\endlastfoot
Low only & \texttt{0x012}, \texttt{0x110}, \texttt{0x169}, \texttt{0x202},
\texttt{0x203}, \texttt{0x204}, \texttt{0x205}, \texttt{0x6FA}, \texttt{0x6FB},
\texttt{0x721}, \texttt{0x731}, \texttt{0x741}, \texttt{0x7B9} \\
Low only & \texttt{0x201} (SYNTREE EPS-C feedback) \\
High only & \texttt{0x210}, \texttt{0x220}, \texttt{0x400} (RT telemetry) \\
Both independent & \texttt{0x7FD}, \texttt{0x7FE}, \texttt{0x7FC} (per-node
heartbeat --- NOT bridged) \\
\end{longtable}

\begin{center}\rule{0.5\linewidth}{0.5pt}\end{center}

\section{5. Priority Groups}\label{priority-groups}

\begin{longtable}[]{@{}
  >{\raggedright\arraybackslash}p{(\columnwidth - 4\tabcolsep) * \real{0.4000}}
  >{\raggedright\arraybackslash}p{(\columnwidth - 4\tabcolsep) * \real{0.4000}}
  >{\raggedright\arraybackslash}p{(\columnwidth - 4\tabcolsep) * \real{0.2000}}@{}}
\toprule\noalign{}
\begin{minipage}[b]{\linewidth}\raggedright
Priority
\end{minipage} & \begin{minipage}[b]{\linewidth}\raggedright
ID Range
\end{minipage} & \begin{minipage}[b]{\linewidth}\raggedright
IDs
\end{minipage} \\
\midrule\noalign{}
\endhead
\bottomrule\noalign{}
\endlastfoot
Highest & \texttt{0x001} & ESTOP \\
Very High & \texttt{0x010}--\texttt{0x01F} & SAFETY\_STATUS, DCDC\_CMD \\
High & \texttt{0x100}--\texttt{0x11F} & MODE\_CMD \\
Medium & \texttt{0x120}--\texttt{0x3FF} & THROTTLE, DRIVE, MTR\_MOTOR\_FBK,
SES\_STATUS/REQ/ErrInfo/Version, DRIVE\_CMD, BRAKE\_REQ, LIGHT\_CMD \\
Low & \texttt{0x400}--\texttt{0x5FF} & OBSTACLE, STATE\_REPORT, PID\_FEEDBACK \\
Lowest & \texttt{0x600}--\texttt{0x7FF} & DIAG, SES\_Test (\texttt{0x6FA}),
SEB\_Test (\texttt{0x6FB}), SEB\_STATUS/ErrInfo/Version
(\texttt{0x721}/\texttt{0x731}/\texttt{0x741}), SEB\_REQ (\texttt{0x7B9}),
HEARTBEAT (\texttt{0x7FC}--\texttt{0x7FE}) \\
\end{longtable}

\begin{quote}
Lower CAN ID = higher bus arbitration priority. Safety-critical frames occupy
\texttt{0x00X}. SYNTREE IDs (\texttt{0x2XX}, \texttt{0x7XX}) are in
medium/lowest ranges per manufacturer assignment.
\end{quote}


% ═══════════════════════════════════════════════════════════════════════
%  PART II -- COMPONENT DOCUMENTATION (docs/)
% ═══════════════════════════════════════════════════════════════════════

\part{Component Documentation}
\label{part:docs}

\chapter{Actuator Interfacing --- Throttle DAC \& Gear
Pass-Through}\label{actuator-interfacing-throttle-dac-gear-pass-through}

The SYS ESP32-S3 interfaces with the motor controller through two channels:

\begin{enumerate}
\def\labelenumi{\arabic{enumi}.}
\tightlist
\item
  \textbf{Throttle:} A 0--5 V analog signal (standard E-bike/e-scooter throttle
  input). Generated by an MCP4725 I2C DAC.
\item
  \textbf{Gear:} Three 72 V lines (D, S, R) sensed via TLP281 optoisolators and
  driven via mechanical relays.
\end{enumerate}

In MANUAL mode, SYS acts as a \textbf{transparent pass-through}: physical inputs
are mirrored directly to outputs. In AUTO mode, SYS drives outputs from CAN
commands (via RT).

\begin{center}\rule{0.5\linewidth}{0.5pt}\end{center}

\section{Throttle --- MCP4725 I2C DAC}\label{throttle-mcp4725-i2c-dac}

\subsection{Why a DAC instead of PWM?}\label{why-a-dac-instead-of-pwm}

The motor controller expects a smooth 0--5 V analog voltage. The ESP32 has no
true analog output --- only PWM (LEDC) with 8--16 bit resolution.

\begin{longtable}[]{@{}
  >{\raggedright\arraybackslash}p{(\columnwidth - 4\tabcolsep) * \real{0.4545}}
  >{\raggedright\arraybackslash}p{(\columnwidth - 4\tabcolsep) * \real{0.2727}}
  >{\raggedright\arraybackslash}p{(\columnwidth - 4\tabcolsep) * \real{0.2727}}@{}}
\toprule\noalign{}
\begin{minipage}[b]{\linewidth}\raggedright
Approach
\end{minipage} & \begin{minipage}[b]{\linewidth}\raggedright
Pros
\end{minipage} & \begin{minipage}[b]{\linewidth}\raggedright
Cons
\end{minipage} \\
\midrule\noalign{}
\endhead
\bottomrule\noalign{}
\endlastfoot
PWM + RC low-pass filter & Free (uses built-in LEDC), fast update & Ripple, slow
settling, sensitive to filter component tolerance. Motor controller may misread
ripple as a throttle command. \\
MCP4725 I2C DAC & Clean 12-bit output, \textless10 µs settling, no ripple &
Additional BOM item (\textasciitilde\$1), I2C bus overhead \\
\end{longtable}

For a safety-critical throttle signal, the clean output of a DAC is worth the
cost. PWM ripple on a throttle line can cause the motor controller to oscillate
between drive and coast --- producing surging that's both uncomfortable and
potentially dangerous.

\subsection{MCP4725 characteristics}\label{mcp4725-characteristics}

\begin{longtable}[]{@{}
  >{\raggedright\arraybackslash}p{(\columnwidth - 2\tabcolsep) * \real{0.6111}}
  >{\raggedright\arraybackslash}p{(\columnwidth - 2\tabcolsep) * \real{0.3889}}@{}}
\toprule\noalign{}
\begin{minipage}[b]{\linewidth}\raggedright
Parameter
\end{minipage} & \begin{minipage}[b]{\linewidth}\raggedright
Value
\end{minipage} \\
\midrule\noalign{}
\endhead
\bottomrule\noalign{}
\endlastfoot
Resolution & 12-bit (4096 steps) \\
Output range & 0 V to VCC \\
VCC & 5 V (powered by 12V→5V LDO, not ESP32 3.3V) \\
Interface & I2C, address 0x60 (A0 pin to GND) \\
Settling time & 6 µs typical \\
I2C clock & 100 kHz (standard mode) or 400 kHz (fast mode) \\
Write time & \textasciitilde30 µs per value at 400 kHz (3 bytes: address, data
MSB, data LSB) \\
\end{longtable}

\subsection{Voltage mapping}\label{voltage-mapping}

\begin{verbatim}
Speed [mm/s]          MCP4725 output [V]      12-bit code
─────────────         ──────────────────       ──────────
   3000 (max fwd)     5.0 V                    4095
   1500               2.5 V                    2047
      0 (neutral)     0.0 V (or offset?)       0
   -500 (max rev)     0.0 V (?)               0
\end{verbatim}

\begin{quote}
\textbf{Note:} The motor controller's reverse behavior is TBD --- some
controllers use a separate reverse input, others use a 0--2.5V = reverse,
2.5--5V = forward mapping. The current implementation treats 0 V as stop and
uses the gear relay for direction. Verify against the specific motor controller
datasheet.
\end{quote}

\subsection{I2C write in C++ (ESP-IDF)}\label{i2c-write-in-c-esp-idf}

\begin{verbatim}
#include <driver/i2c.h>

#define MCP4725_ADDR  0x60
#define I2C_MASTER_NUM I2C_NUM_0

void throttle_write(uint16_t value_12bit) {
    // MCP4725 expects: [upper 4 bits in first byte, lower 8 bits in second byte]
    uint8_t data[2];
    data[0] = (value_12bit >> 4) & 0xFF;   // upper 8 bits (includes 4 data + 4 zero)
    data[1] = (value_12bit & 0x0F) << 4;    // lower 4 bits shifted to upper nibble

    i2c_cmd_handle_t cmd = i2c_cmd_link_create();
    i2c_master_start(cmd);
    i2c_master_write_byte(cmd, (MCP4725_ADDR << 1) | I2C_MASTER_WRITE, true);
    i2c_master_write_byte(cmd, data[0], true);
    i2c_master_write_byte(cmd, data[1], true);
    i2c_master_stop(cmd);
    i2c_master_cmd_begin(I2C_MASTER_NUM, cmd, pdMS_TO_TICKS(10));
    i2c_cmd_link_delete(cmd);
}
\end{verbatim}

\subsection{Throttle in MANUAL mode}\label{throttle-in-manual-mode}

\begin{verbatim}
Throttle grip (0–5V) ──► SYS ADC1 ──► map to 12-bit ──► MCP4725 ──► Motor controller
\end{verbatim}

The ADC value is read at 100 Hz by the \texttt{throttle} task and used directly:

\begin{verbatim}
void throttle_task() {
    // Read ADC, convert to 12-bit MCP4725 value
    int adc_raw = adc1_get_raw(ADC1_CHANNEL_5);  // 0–4095 (12-bit ESP32 ADC)
    // Optional: deadband, smoothing, rate limiting
    uint16_t dac_value = adc_raw;  // direct pass-through
    throttle_write(dac_value);
    // Also send to CAN 0x120 for Jetson telemetry
    int16_t speed_mmps = adc_to_speed(adc_raw);
    can_send_throttle_position(speed_mmps);
}
\end{verbatim}

\subsection{Throttle in AUTO mode}\label{throttle-in-auto-mode}

\begin{verbatim}
CAN 0x204 (speed_mmps) ──► SYS dispatch ──► setpoint_queue ──► motor_task ──► MCP4725
\end{verbatim}

The \texttt{motor\_task} reads \texttt{setpoint\_queue} and writes the DAC:

\begin{verbatim}
void motor_task() {
    ActuatorSetpoint sp;
    if (xQueueReceive(setpoint_queue, &sp, 0) == pdTRUE) {
        latest_setpoint = sp;
    }
    // In AUTO: use setpoint
    // In MANUAL: throttle_task writes DAC directly (motor_task skips)
    // In ESTOP: write 0 to DAC regardless
    if (g_mode == Mode::Estop) {
        throttle_write(0);
    } else if (g_mode == Mode::Auto) {
        uint16_t dac = speed_to_dac(latest_setpoint.motor_speed_mmps);
        throttle_write(dac);
    }
    // MANUAL: throttle_task handles it
}
\end{verbatim}

\begin{center}\rule{0.5\linewidth}{0.5pt}\end{center}

\section{Gear --- TLP281 input + relay
output}\label{gear-tlp281-input-relay-output}

The motor controller uses 72 V logic on its gear select inputs: - Apply 72 V to
the \textbf{D} (Drive) pin → forward. - Apply 72 V to the \textbf{S} (Sport) pin
→ sport mode. - Apply 72 V to the \textbf{R} (Reverse) pin → reverse. - No pin
energized → Neutral.

\subsection{Input: TLP281 optoisolator
reading}\label{input-tlp281-optoisolator-reading}

The physical gear selector is a handlebar switch wired to 72 V. SYS senses which
gear is selected via TLP281 optoisolators (see
{[}{[}high-voltage-isolation{]}{]}):

\begin{verbatim}
Physical switch              SYS GPIO (via TLP281)
──────────────               ─────────────────────
D selected → 72V on D line → GPIO_X reads LOW
S selected → 72V on S line → GPIO_Y reads LOW
R selected → 72V on R line → GPIO_Z reads LOW
Neutral     → no 72V        → all GPIOs read HIGH
\end{verbatim}

A GPIO reading LOW means ``this gear is active'' (active-low).

\subsection{Output: Relay module}\label{output-relay-module}

The 72 V gear lines to the motor controller are switched by mechanical relays:

\begin{verbatim}
ESP32 GPIO ── transistor driver ── relay coil ── 12V
Relay NO contact: 72V bus ── fuse ── motor controller gear input
\end{verbatim}

\subsection{Gear FSM}\label{gear-fsm}

\begin{verbatim}
State: current_gear ∈ {N, D, S, R}

MANUAL mode:
  - Read TLP281 GPIOs at 50 Hz
  - Mirror to relay outputs: gear_in == gear_out
  - The rider's physical switch directly controls the motor controller

AUTO mode:
  - Read gear from 0x204 RT_DRIVE_CMD.gear field
  - Drive relay outputs accordingly
  - Ignore physical switch position (rider can't override AUTO gear)

ESTOP:
  - All relays OFF → all gear lines de-energized → motor controller sees Neutral
\end{verbatim}

\begin{verbatim}
void gear_task() {
    if (g_mode == Mode::Estop) {
        gear_relay_drive(Gear::N);  // force neutral
        return;
    }

    if (g_mode == Mode::Manual) {
        // Pass-through: read physical switch, mirror to relays
        Gear physical = gear_read_tlp281();
        gear_relay_drive(physical);
    } else {
        // AUTO: follow setpoint
        gear_relay_drive(latest_setpoint.gear);
    }
}
\end{verbatim}

\subsection{Why pass-through matters}\label{why-pass-through-matters}

In MANUAL mode, SYS is a transparent conduit. The rider twists the throttle →
SYS reads ADC → SYS writes DAC. The rider selects a gear → SYS reads optos → SYS
drives relays. This means:

\begin{enumerate}
\def\labelenumi{\arabic{enumi}.}
\tightlist
\item
  \textbf{SYS firmware is not in the critical path for manual operation.} If the
  AUTO path has a bug, the rider can switch to MANUAL and still control the
  vehicle.
\item
  \textbf{No latency added.} The pass-through is just ADC→DAC and GPIO→relay,
  with no CAN round-trip.
\item
  \textbf{SYS can still monitor.} Even in pass-through, SYS reads throttle
  position and sends it on CAN \texttt{0x120} for Jetson telemetry.
\end{enumerate}

\begin{center}\rule{0.5\linewidth}{0.5pt}\end{center}

\section{Combined throttle + gear behavior by
mode}\label{combined-throttle-gear-behavior-by-mode}

\begin{longtable}[]{@{}llll@{}}
\toprule\noalign{}
Mode & Throttle DAC & Gear relays & Source \\
\midrule\noalign{}
\endhead
\bottomrule\noalign{}
\endlastfoot
MANUAL & ADC pass-through & TLP281 mirror & Physical controls \\
AUTO & \texttt{0x204} setpoint → DAC & \texttt{0x204} gear field → relays &
Jetson via RT \\
ESTOP & 0 V (forced) & All OFF (forced) & Safety override \\
\end{longtable}

\begin{center}\rule{0.5\linewidth}{0.5pt}\end{center}

\emph{Primary reference: {[}{[}emergency-system{]}{]} for ESTOP behavior on
throttle DAC and gear relays, and the emergency response matrix.}

\emph{See also: {[}{[}high-voltage-isolation{]}{]} for TLP281 and relay
isolation details, {[}{[}defense-in-depth-safety{]}{]} for ESTOP behavior,
{[}{[}architecture{]}{]} §4 for signal flow diagrams, §8.6 for SYS control
mechanisms.}


\chapter{Autoware.Auto Integration --- E-Trike
v0.0.4-alpha}\label{autoware.auto-integration-e-trike-v0.0.4-alpha}

\textbf{Supersedes:} All previous 0.0.4 documents \textbf{Base document:}
\texttt{docs/autoware-auto-communication-architecture.md} \textbf{Audit basis:}
12 agents across 4 rounds \textbf{Date:} 2026-06-23

\begin{center}\rule{0.5\linewidth}{0.5pt}\end{center}

\section{Optimal Architecture: Minimum Change, Maximum
Compatibility}\label{optimal-architecture-minimum-change-maximum-compatibility}

The current E-Trike CAN protocol works. RT does kinematics. Safety is verified.
The only thing missing for Autoware.Auto compatibility is the ROS 2 topic
interface on the Jetson.

\textbf{What changes:} Jetson gets a new ROS node
(\texttt{autoware\_vehicle\_bridge}). The CAN bus stays exactly as it is.

\begin{verbatim}
AckermannControlCommand (ROS)
        │
        ▼
┌─────────────────────────┐
│ autoware_vehicle_bridge │  ← NEW (Jetson)
│  ROS → CAN: 0x300       │
│  CAN → ROS: reports     │
└───────────┬─────────────┘
            │ 0x300 {speed, yaw, gear}  ← UNCHANGED
    High CAN│ 500 kbit/s
            ▼
┌─────── RT ESP32-S3 ─────┐  ← UNCHANGED
│  kinematics, gateway,    │
│  steering, safety        │
└───────────┬─────────────┘
    Low CAN │ 500 kbit/s
            ▼
    EPS-C + SEB + SYS + MTR  ← UNCHANGED
\end{verbatim}

Everything below the red line stays identical to today. The only new code is the
Jetson ROS node.

\begin{center}\rule{0.5\linewidth}{0.5pt}\end{center}

\section{1. ROS 2 Topic Interface (Jetson
Node)}\label{ros-2-topic-interface-jetson-node}

\subsection{1.1 Subscriptions (Inputs)}\label{subscriptions-inputs}

\begin{longtable}[]{@{}
  >{\raggedright\arraybackslash}p{(\columnwidth - 6\tabcolsep) * \real{0.2500}}
  >{\raggedright\arraybackslash}p{(\columnwidth - 6\tabcolsep) * \real{0.2500}}
  >{\raggedright\arraybackslash}p{(\columnwidth - 6\tabcolsep) * \real{0.2500}}
  >{\raggedright\arraybackslash}p{(\columnwidth - 6\tabcolsep) * \real{0.2500}}@{}}
\toprule\noalign{}
\begin{minipage}[b]{\linewidth}\raggedright
Topic
\end{minipage} & \begin{minipage}[b]{\linewidth}\raggedright
Message Type
\end{minipage} & \begin{minipage}[b]{\linewidth}\raggedright
Maps To
\end{minipage} & \begin{minipage}[b]{\linewidth}\raggedright
Notes
\end{minipage} \\
\midrule\noalign{}
\endhead
\bottomrule\noalign{}
\endlastfoot
\texttt{\textasciitilde{}/input/control\_cmd} &
\texttt{autoware\_auto\_control\_msgs/msg/AckermannControlCommand} & High CAN
\texttt{0x300} & \texttt{longitudinal.speed} → speed\_mmps.
\texttt{lateral.steering\_tire\_angle} → yaw via inverse tricycle kinematics.
\texttt{longitudinal.acceleration} → brake kPa via deceleration model. \\
\texttt{\textasciitilde{}/input/gear\_cmd} &
\texttt{autoware\_auto\_vehicle\_msgs/msg/GearCommand} & \texttt{0x300} gear
field & Override auto-derived gear \\
\texttt{\textasciitilde{}/input/turn\_indicators\_cmd} &
\texttt{autoware\_auto\_vehicle\_msgs/msg/TurnIndicatorsCommand} &
\texttt{0x302} turn bits & Forwarded by RT to SYS \\
\texttt{\textasciitilde{}/input/hazard\_lights\_cmd} &
\texttt{autoware\_auto\_vehicle\_msgs/msg/HazardLightsCommand} & \texttt{0x302}
both turn bits & \\
\texttt{\textasciitilde{}/input/engage} &
\texttt{autoware\_auto\_vehicle\_msgs/msg/Engage} & Internal state &
\texttt{false} → stop publishing commands \\
\end{longtable}

\subsection{1.2 Publications (Outputs)}\label{publications-outputs}

\begin{longtable}[]{@{}
  >{\raggedright\arraybackslash}p{(\columnwidth - 6\tabcolsep) * \real{0.2500}}
  >{\raggedright\arraybackslash}p{(\columnwidth - 6\tabcolsep) * \real{0.2500}}
  >{\raggedright\arraybackslash}p{(\columnwidth - 6\tabcolsep) * \real{0.2500}}
  >{\raggedright\arraybackslash}p{(\columnwidth - 6\tabcolsep) * \real{0.2500}}@{}}
\toprule\noalign{}
\begin{minipage}[b]{\linewidth}\raggedright
Topic
\end{minipage} & \begin{minipage}[b]{\linewidth}\raggedright
Message Type
\end{minipage} & \begin{minipage}[b]{\linewidth}\raggedright
Source CAN
\end{minipage} & \begin{minipage}[b]{\linewidth}\raggedright
Notes
\end{minipage} \\
\midrule\noalign{}
\endhead
\bottomrule\noalign{}
\endlastfoot
\texttt{\textasciitilde{}/output/velocity\_status} &
\texttt{autoware\_auto\_vehicle\_msgs/msg/VelocityReport} & \texttt{0x120}
(speed) & \texttt{longitudinal\_velocity} = mmps/1000. \texttt{heading\_rate} =
0 (no sensor). \texttt{lateral\_velocity} = 0. \\
\texttt{\textasciitilde{}/output/steering\_status} &
\texttt{autoware\_auto\_vehicle\_msgs/msg/SteeringReport} & \emph{(via
\texttt{0x201} SES\_STATUS on low bus --- RT forwards as \texttt{0x210} steering
validity)} & \texttt{steering\_tire\_angle} from RT state report or derived from
\texttt{0x300} echo \\
\texttt{\textasciitilde{}/output/gear\_status} &
\texttt{autoware\_auto\_vehicle\_msgs/msg/GearReport} & \texttt{0x210}
(mode+steer+reversing) & RT derives gear from speed direction \\
\texttt{\textasciitilde{}/output/control\_mode} &
\texttt{autoware\_auto\_vehicle\_msgs/msg/ControlModeReport} & \texttt{0x210}
mode field + \texttt{0x011} estop & Manual→MANUAL(4),
Auto+engaged→AUTONOMOUS(1), Estop→DISENGAGED(5) \\
\end{longtable}

\subsection{1.3 What Gets Fixed in the Original
Document}\label{what-gets-fixed-in-the-original-document}

\begin{longtable}[]{@{}
  >{\raggedright\arraybackslash}p{(\columnwidth - 4\tabcolsep) * \real{0.3333}}
  >{\raggedright\arraybackslash}p{(\columnwidth - 4\tabcolsep) * \real{0.3333}}
  >{\raggedright\arraybackslash}p{(\columnwidth - 4\tabcolsep) * \real{0.3333}}@{}}
\toprule\noalign{}
\begin{minipage}[b]{\linewidth}\raggedright
Original
\end{minipage} & \begin{minipage}[b]{\linewidth}\raggedright
Corrected
\end{minipage} & \begin{minipage}[b]{\linewidth}\raggedright
Reason
\end{minipage} \\
\midrule\noalign{}
\endhead
\bottomrule\noalign{}
\endlastfoot
\texttt{autoware\_control\_msgs} & \texttt{autoware\_auto\_control\_msgs} &
Missing \texttt{.auto} prefix \\
\texttt{autoware\_vehicle\_msgs} & \texttt{autoware\_auto\_vehicle\_msgs} &
Missing \texttt{.auto} prefix \\
\texttt{Control} message & \texttt{AckermannControlCommand} & Current canonical
type \\
\texttt{tier4\_vehicle\_msgs/*} & \emph{(removed)} & Not Autoware.Auto. Replaced
with \texttt{autoware\_auto\_*\_msgs} equivalents or removed. \\
\texttt{ActuationCommandStamped} & \emph{(removed)} & Nav2 pattern. Control →
CAN directly. \\
\texttt{VehicleEmergencyStamped} & \texttt{Engage} message & Original type
doesn't exist in Autoware.Auto \\
Gear \texttt{10=REVERSE,\ 20=PARK} & \texttt{REVERSE=20,\ PARK=22,\ LOW=23} &
Match \texttt{autoware\_auto\_vehicle\_msgs/GearCommand} constants \\
\texttt{\textasciitilde{}/input/from\_can\_bus} /
\texttt{\textasciitilde{}/output/to\_can\_bus} & \emph{(removed)} & Raw CAN
pass-through is a security concern \\
\texttt{longitudinal.speed} & \emph{(keep --- IS correct)} & Field is
\texttt{speed} in \texttt{LongitudinalCommand.msg} \\
\texttt{1=Autonomous,\ 4=Manual} & \emph{(keep --- matches Autoware constants)}
& \texttt{ControlModeReport}: AUTONOMOUS=1, MANUAL=4 \\
\end{longtable}

\begin{center}\rule{0.5\linewidth}{0.5pt}\end{center}

\section{2. CAN Protocol (Unchanged from Current
Architecture)}\label{can-protocol-unchanged-from-current-architecture}

The high-bus CAN protocol stays exactly as defined in \texttt{architecture.md}
§2.2 and \texttt{can-dictionary.md} §2. No CAN IDs change. No encoding changes.
No new frames.

\subsection{2.1 Commands (Jetson → RT) ---
UNCHANGED}\label{commands-jetson-rt-unchanged}

\begin{longtable}[]{@{}
  >{\raggedright\arraybackslash}p{(\columnwidth - 6\tabcolsep) * \real{0.2500}}
  >{\raggedright\arraybackslash}p{(\columnwidth - 6\tabcolsep) * \real{0.2500}}
  >{\raggedright\arraybackslash}p{(\columnwidth - 6\tabcolsep) * \real{0.2500}}
  >{\raggedright\arraybackslash}p{(\columnwidth - 6\tabcolsep) * \real{0.2500}}@{}}
\toprule\noalign{}
\begin{minipage}[b]{\linewidth}\raggedright
CAN ID
\end{minipage} & \begin{minipage}[b]{\linewidth}\raggedright
Name
\end{minipage} & \begin{minipage}[b]{\linewidth}\raggedright
Content
\end{minipage} & \begin{minipage}[b]{\linewidth}\raggedright
Rate
\end{minipage} \\
\midrule\noalign{}
\endhead
\bottomrule\noalign{}
\endlastfoot
\texttt{0x001} & SAFETY\_ESTOP & DLC=0 (event) & Event \\
\texttt{0x300} & HOST\_DRIVE\_CMD &
\texttt{\{i32\ speed\_mmps,\ i24\ yaw\_rate\_mrad\_s,\ u8\ gear\}} & ≤100 Hz \\
\texttt{0x301} & HOST\_BRAKE\_REQ & \texttt{\{i32\ brake\_pressure\_kpa\}} & On
demand \\
\texttt{0x302} & HOST\_LIGHT\_CMD &
\texttt{\{u8\ bits:\ LT{[}0{]}\ RT{[}1{]}\ BRK{[}2{]}\ HEAD{[}3{]}\}} & On
change \\
\texttt{0x7FC} & JETSON\_HEARTBEAT & \texttt{\{u8\ alive\_ctr\}} & 2 Hz \\
\end{longtable}

\subsection{2.2 Feedback (RT → Jetson) ---
UNCHANGED}\label{feedback-rt-jetson-unchanged}

\begin{longtable}[]{@{}
  >{\raggedright\arraybackslash}p{(\columnwidth - 6\tabcolsep) * \real{0.2500}}
  >{\raggedright\arraybackslash}p{(\columnwidth - 6\tabcolsep) * \real{0.2500}}
  >{\raggedright\arraybackslash}p{(\columnwidth - 6\tabcolsep) * \real{0.2500}}
  >{\raggedright\arraybackslash}p{(\columnwidth - 6\tabcolsep) * \real{0.2500}}@{}}
\toprule\noalign{}
\begin{minipage}[b]{\linewidth}\raggedright
CAN ID
\end{minipage} & \begin{minipage}[b]{\linewidth}\raggedright
Name
\end{minipage} & \begin{minipage}[b]{\linewidth}\raggedright
Content
\end{minipage} & \begin{minipage}[b]{\linewidth}\raggedright
Rate
\end{minipage} \\
\midrule\noalign{}
\endhead
\bottomrule\noalign{}
\endlastfoot
\texttt{0x001} & SAFETY\_ESTOP & DLC=0 (forwarded from low bus) & Event \\
\texttt{0x011} & SYS\_SAFETY\_STS & \texttt{\{u8\ estop,\ u8\ hb\_ok\}}
(forwarded from SYS) & 5 Hz \\
\texttt{0x120} & SYS\_THROTTLE\_STS & \texttt{\{i16\ speed\_mmps\}} (forwarded
from MTR) & 100 Hz \\
\texttt{0x210} & RT\_STATE\_RPT &
\texttt{\{u8\ mode,\ u8\ steer\_valid,\ u8\ reversing\}} & 10 Hz \\
\texttt{0x400} & HOST\_OBSTACLE\_DIST & \texttt{\{u32\ distance\_mm\}} & 10
Hz \\
\texttt{0x600} & SYS\_DIAG\_RPT &
\texttt{\{u8\ mode,\ u8\ brake,\ u8\ hb,\ u8\ estop,\ u16\ heap,\ u8\ tec,\ u8\ rec\}}
& 1 Hz \\
\texttt{0x7FD} & RT\_HEARTBEAT & \texttt{\{u8\ alive\_ctr\}} & 2 Hz \\
\end{longtable}

\subsection{2.3 Why No Changes}\label{why-no-changes}

\begin{itemize}
\tightlist
\item
  \textbf{CAN IDs:} The current IDs work. Changing them requires coordinated RT
  + Jetson firmware updates with zero functional benefit.
\item
  \textbf{Encoding:} Physical units (mm/s, kPa) are directly readable on
  \texttt{candump}. Ratio encoding adds conversion steps without improving
  anything.
\item
  \textbf{Split frames:} The current combined \texttt{0x300} is atomic per
  planning cycle. Splitting into steer/brake/throttle CAN IDs creates
  correlation problems that don't exist today.
\item
  \textbf{Autoware.Auto doesn't standardize CAN.} The ROS topics are the
  interface contract. The CAN bus is an internal implementation detail.
\end{itemize}

\begin{center}\rule{0.5\linewidth}{0.5pt}\end{center}

\section{3. RT ESP32-S3 (Unchanged)}\label{rt-esp32-s3-unchanged}

The RT firmware stays as described in \texttt{architecture.md} §7. No changes
to: - CAN gateway forwarding rules (§2.3) - Tricycle kinematics (§7.6) -
Steering LBS state machine + dynamic clamp + following error - Brake max-select
arbitration - Obstacle speed limit - Command staleness watchdog (500ms) - 8
FreeRTOS tasks + priorities (§7.7) - Any CAN TX/RX behavior

\textbf{One addition (optional, not required for Autoware.Auto):} Add
\texttt{0x721} SEB\_STATUS to RT's low-bus RX list for SEB error monitoring.
This is a pre-existing gap, not an upgrade requirement.

\begin{center}\rule{0.5\linewidth}{0.5pt}\end{center}

\section{4. Safety (Unchanged)}\label{safety-unchanged}

All safety mechanisms in \texttt{architecture.md} §§6-8 remain identical: - 8
ESTOP trigger paths (physical button, CAN 0x001, heartbeat loss, following
error, external watchdog, 0x204 staleness) - EGAS 3-level motor safety (MTR L1,
SYS L2, hardware ESTOP L3) - Mode-gated dual control (Option D) -
Heartbeat/liveness matrix (unchanged rates and timeouts) - SYS heartbeat: 10 Hz
/ 200ms (the post-gap-\#12 rate --- can-dictionary.md to be updated)

\begin{center}\rule{0.5\linewidth}{0.5pt}\end{center}

\section{5. Jetson Node Implementation}\label{jetson-node-implementation}

\subsection{5.1 Package}\label{package}

\begin{verbatim}
jetson/src/autoware_vehicle_bridge/
  ├── include/autoware_vehicle_bridge/
  │   └── vehicle_bridge_node.hpp
  ├── src/
  │   ├── vehicle_bridge_node.cpp    # Lifecycle node
  │   ├── ros_to_can.cpp             # AckermannControlCommand → 0x300/0x301/0x302
  │   └── can_to_ros.cpp             # 0x120/0x011/0x210/0x600 → ROS messages
  ├── config/
  │   └── etrike.param.yaml          # Trike-specific parameters
  └── CMakeLists.txt
\end{verbatim}

\subsection{5.2 Encoding (ROS → CAN)}\label{encoding-ros-can}

\begin{verbatim}
// AckermannControlCommand → 0x300 HOST_DRIVE_CMD
void encode_drive_cmd(const AckermannControlCommand& cmd, can::Frame& fr) {
    fr.id = 0x300;
    fr.dlc = 8;
    
    // Speed: m/s → mm/s
    int32_t speed_mmps = (int32_t)(cmd.longitudinal.speed * 1000.0f);
    speed_mmps = std::clamp(speed_mmps, -500, 3000);
    
    // Steering angle → yaw rate via inverse tricycle kinematics
    // δ = arctan(L·ω/|v|) → ω = |v|·tan(δ)/L  where L=1.5m
    float v_ms = std::abs(cmd.longitudinal.speed);
    float delta_rad = cmd.lateral.steering_tire_angle;
    int32_t yaw_mrad_s = 0;
    if (v_ms > 0.05f) {
        float omega_rad_s = v_ms * std::tan(delta_rad) / 1.5f;
        yaw_mrad_s = (int32_t)(omega_rad_s * 1000.0f);
    }
    yaw_mrad_s = std::clamp(yaw_mrad_s, -3000, 3000);
    
    // Gear: auto-derived from speed direction (CAN override via GearCommand)
    uint8_t gear = (speed_mmps > 50) ? 1 : (speed_mmps < -50) ? 3 : 0;  // D=1, R=3, N=0
    
    // Big-endian encode
    fr.put_i32(0, speed_mmps);
    fr.data[4] = (yaw_mrad_s >> 16) & 0xFF;
    fr.data[5] = (yaw_mrad_s >> 8) & 0xFF;
    fr.data[6] = yaw_mrad_s & 0xFF;
    fr.data[7] = gear;
}

// Brake: longitudinal.acceleration (negative) → 0x301 HOST_BRAKE_REQ
void encode_brake_req(const AckermannControlCommand& cmd, can::Frame& fr) {
    if (!cmd.longitudinal.is_defined_acceleration) return;
    float decel = std::max(0.0f, -cmd.longitudinal.acceleration);
    int32_t kpa = (int32_t)(decel / 5.0f * 5000.0f);  // 5 m/s² → 5000 kPa (max)
    kpa = std::clamp(kpa, 0, 5000);
    
    fr.id = 0x301;
    fr.dlc = 4;
    fr.put_i32(0, kpa);
}
\end{verbatim}

\subsection{5.3 Decoding (CAN → ROS)}\label{decoding-can-ros}

\begin{verbatim}
// 0x120 SYS_THROTTLE_STS → VelocityReport
void decode_velocity(const can::Frame& fr, VelocityReport& msg) {
    int16_t speed_mmps = fr.get_i16(0);
    msg.longitudinal_velocity = speed_mmps / 1000.0f;  // mm/s → m/s
    msg.lateral_velocity = 0.0f;   // No sensor
    msg.heading_rate = 0.0f;        // No sensor (TBD: IMU or inverse kinematics)
}

// 0x210 RT_STATE_RPT → ControlModeReport + GearReport
void decode_state(const can::Frame& fr, ControlModeReport& mode_msg, GearReport& gear_msg) {
    uint8_t mode = fr.data[0];
    bool reversing = fr.data[2];
    
    // Map trike mode → Autoware constants
    switch (mode) {
        case 0: mode_msg.mode = ControlModeReport::MANUAL; break;      // 4
        case 1: mode_msg.mode = ControlModeReport::AUTONOMOUS; break;   // 1
        case 2: mode_msg.mode = ControlModeReport::DISENGAGED; break;   // 5
    }
    
    uint8_t gear = reversing ? 3 : 1;  // R=3, D=1 (simplified — full mapping TBD)
    gear_msg.report = gear;
}
\end{verbatim}

\subsection{5.4 Parameters
(etrike.param.yaml)}\label{parameters-etrike.param.yaml}

\begin{verbatim}
/**:
  ros__parameters:
    wheel_base: 1.5                # m (tricycle)
    max_speed_forward: 3.0         # m/s
    max_speed_reverse: 0.5         # m/s
    max_steering_angle: 0.698      # rad (40°)
    max_brake_pressure_kpa: 5000   # 5 MPa
    loop_rate: 100.0               # Hz
    command_timeout_ms: 500
    can_interface: "can0"
    rt_heartbeat_timeout_ms: 1500
\end{verbatim}

\textbf{Removed from original document's §3} (not applicable to trike): VGR
coefficients, wheel\_radius, margin\_time\_for\_gear\_change,
use\_actuation\_cmd, convert\_steer\_cmd, max\_throttle (pedal ratio),
max\_brake (pedal ratio), emergency\_brake (pedal ratio)

\begin{center}\rule{0.5\linewidth}{0.5pt}\end{center}

\section{\texorpdfstring{6. Document Updates to
\texttt{autoware-auto-communication-architecture.md}}{6. Document Updates to autoware-auto-communication-architecture.md}}\label{document-updates-to-autoware-auto-communication-architecture.md}

\subsection{6.1 §1 ROS 2 Topics: Fix compliance
errors}\label{ros-2-topics-fix-compliance-errors}

\begin{itemize}
\tightlist
\item
  \texttt{.auto} prefix on all package names
\item
  \texttt{Control} → \texttt{AckermannControlCommand}
\item
  Remove \texttt{ActuationCommandStamped}, \texttt{VehicleEmergencyStamped},
  \texttt{ActuationStatusStamped} (tier4)
\item
  Remove raw CAN pass-through topics
\item
  Add \texttt{Engage} subscription
\item
  Keep correct field names (\texttt{longitudinal.speed},
  \texttt{1=Autonomous/4=Manual})
\end{itemize}

\subsection{6.2 §2 CAN Protocol: Replace with actual trike CAN
IDs}\label{can-protocol-replace-with-actual-trike-can-ids}

\begin{itemize}
\tightlist
\item
  Use \texttt{architecture.md} §2.2 / \texttt{can-dictionary.md} §2 as source of
  truth
\item
  Commands: \texttt{0x300} drive, \texttt{0x301} brake, \texttt{0x302} lights,
  \texttt{0x001} ESTOP
\item
  Feedback: \texttt{0x011}, \texttt{0x120}, \texttt{0x210}, \texttt{0x400},
  \texttt{0x600}
\item
  Heartbeats: \texttt{0x7FC}, \texttt{0x7FD}
\item
  Physical units throughout (no ratio encoding)
\item
  Remove the original's \texttt{0x200}--\texttt{0x206},
  \texttt{0x300}--\texttt{0x340} CAN IDs (reference design for different
  vehicle)
\end{itemize}

\subsection{6.3 §3 Parameters: Set trike
values}\label{parameters-set-trike-values}

\begin{itemize}
\tightlist
\item
  wheel\_base: 1.5m, max\_steering\_angle: 0.698rad (40°), etc.
\item
  Remove passenger-car params (VGR, wheel\_radius, margin\_time)
\end{itemize}

\subsection{6.4 §4 micro-ROS: Extract}\label{micro-ros-extract}

\begin{itemize}
\tightlist
\item
  Move to separate reference document or remove
\item
  Not applicable to E-Trike (ESP32-S3 + raw CAN)
\end{itemize}

\subsection{6.5 §5 HIL Serial: Extract}\label{hil-serial-extract}

\begin{itemize}
\tightlist
\item
  Move to \texttt{docs/hil-simulation-protocol.md}
\item
  Not production hardware
\end{itemize}

\subsection{6.6 §6 Safety Limits: Adapt}\label{safety-limits-adapt}

\begin{itemize}
\tightlist
\item
  Keep safety concept, update references to E-Trike architecture
\item
  Reference \texttt{architecture.md} §§7.10, 8.10 for actual enforcement
\end{itemize}

\begin{center}\rule{0.5\linewidth}{0.5pt}\end{center}

\section{7. Pre-Existing Issues (Separate from This
Upgrade)}\label{pre-existing-issues-separate-from-this-upgrade}

These issues exist in the current system regardless of Autoware.Auto
integration. They should be tracked separately:

\begin{longtable}[]{@{}
  >{\raggedright\arraybackslash}p{(\columnwidth - 4\tabcolsep) * \real{0.3333}}
  >{\raggedright\arraybackslash}p{(\columnwidth - 4\tabcolsep) * \real{0.3333}}
  >{\raggedright\arraybackslash}p{(\columnwidth - 4\tabcolsep) * \real{0.3333}}@{}}
\toprule\noalign{}
\begin{minipage}[b]{\linewidth}\raggedright
\#
\end{minipage} & \begin{minipage}[b]{\linewidth}\raggedright
Issue
\end{minipage} & \begin{minipage}[b]{\linewidth}\raggedright
Severity
\end{minipage} \\
\midrule\noalign{}
\endhead
\bottomrule\noalign{}
\endlastfoot
B1 & Steering angle offset (EPS-C CSV offset=-3000 vs architecture offset=0) &
CRITICAL \\
B2 & SEB comm-fault behavior (hold vs release on CAN loss) & CRITICAL \\
H1 & Dynamic angle clamp interpolation formula undefined & HIGH \\
H2 & SYS heartbeat rate: 10Hz in architecture §8.6/§8.9 but 2Hz in
can-dictionary & HIGH \\
H3 & 0x721 not in RT RX list (needed for SEB error monitoring) & HIGH \\
H4 & 0x400 direction contradictory (Jetson→RT vs RT→Jetson) & MEDIUM \\
M1 & Following error threshold should be speed-scaled, not fixed 5° & MEDIUM \\
M2 & Obstacle→kPa formula undefined & MEDIUM \\
M3 & 0x206 MTR\_MOTOR\_FBK missing from can-dictionary & MEDIUM \\
M4 & VCU\_Veh\_Spd\_Value source (commanded vs measured speed) undocumented &
MEDIUM \\
M5 & 0x120 sender: MTR vs SYS documentation conflict & LOW \\
\end{longtable}

\textbf{These are not caused by the Autoware.Auto upgrade. They exist today.}
They should be fixed independently.

\begin{center}\rule{0.5\linewidth}{0.5pt}\end{center}

\section{8. What This Approach Avoids}\label{what-this-approach-avoids}

\begin{longtable}[]{@{}
  >{\raggedright\arraybackslash}p{(\columnwidth - 2\tabcolsep) * \real{0.5000}}
  >{\raggedright\arraybackslash}p{(\columnwidth - 2\tabcolsep) * \real{0.5000}}@{}}
\toprule\noalign{}
\begin{minipage}[b]{\linewidth}\raggedright
Avoided Complexity
\end{minipage} & \begin{minipage}[b]{\linewidth}\raggedright
Why
\end{minipage} \\
\midrule\noalign{}
\endhead
\bottomrule\noalign{}
\endlastfoot
CAN ID changes & Current IDs work. Changing them requires coordinated RT+Jetson
firmware updates. \\
Ratio encoding & Two conversion steps (physical→ratio→CAN→ratio→physical) that
cancel out. \\
Split command frames & PendingSetpoint assembly, 3 queues, per-field staleness,
cross-cycle correlation. \\
Dual-protocol transition & Gated deployment, RX for both old and new, phased ID
enable/disable. \\
\texttt{0x300} direction reversal & Old: Jetson→RT command. New: RT→Jetson
feedback. Same ID, same bus, opposite direction. \\
Extra feedback CAN IDs & \texttt{0x301} BRAKE\_FBK and \texttt{0x302}
THROTTLE\_FBK not needed for any ROS message. \\
\end{longtable}

\begin{center}\rule{0.5\linewidth}{0.5pt}\end{center}

\section{9. Implementation}\label{implementation}

\subsection{Phase 1: Document Update}\label{phase-1-document-update}

\begin{itemize}
\tightlist
\item[$\square$]
  Apply §6 changes to \texttt{autoware-auto-communication-architecture.md}
\item[$\square$]
  Fix can-dictionary.md: \texttt{0x7FE}→10Hz, add \texttt{0x206}, fix
  \texttt{0x120} sender
\end{itemize}

\subsection{Phase 2: Jetson Node}\label{phase-2-jetson-node}

\begin{itemize}
\tightlist
\item[$\square$]
  Create \texttt{jetson/src/autoware\_vehicle\_bridge/} package
\item[$\square$]
  Lifecycle node (configure/activate/deactivate/cleanup)
\item[$\square$]
  ROS→CAN encoding (AckermannControlCommand → 0x300/0x301/0x302)
\item[$\square$]
  CAN→ROS decoding (0x120/0x011/0x210 → reports)
\item[$\square$]
  Heartbeat (send \texttt{0x7FC}, monitor \texttt{0x7FD})
\end{itemize}

\subsection{Phase 3: Integration Test}\label{phase-3-integration-test}

\begin{itemize}
\tightlist
\item[$\square$]
  End-to-end: AckermannControlCommand → 0x300 → RT → actuator → feedback → ROS
  report
\item[$\square$]
  Engage=false → stop commands
\item[$\square$]
  CAN 0x001 ESTOP path from Jetson
\item[$\square$]
  4-hour endurance soak
\end{itemize}

\begin{center}\rule{0.5\linewidth}{0.5pt}\end{center}

\section{10. CAN ID Quick Reference}\label{can-id-quick-reference}

\subsection{High Bus (Jetson ↔ RT) ---
UNCHANGED}\label{high-bus-jetson-rt-unchanged}

\begin{longtable}[]{@{}llll@{}}
\toprule\noalign{}
ID & Name & Direction & Content \\
\midrule\noalign{}
\endhead
\bottomrule\noalign{}
\endlastfoot
\texttt{0x001} & SAFETY\_ESTOP & Bidir (bridged) & DLC=0 \\
\texttt{0x011} & SYS\_SAFETY\_STS & ← (fwd) &
\texttt{\{u8\ estop,\ u8\ hb\_ok\}} \\
\texttt{0x120} & SYS\_THROTTLE\_STS & ← (fwd) & \texttt{\{i16\ speed\_mmps\}} \\
\texttt{0x210} & RT\_STATE\_RPT & ← &
\texttt{\{u8\ mode,\ u8\ steer\_valid,\ u8\ reversing\}} \\
\texttt{0x300} & HOST\_DRIVE\_CMD & → &
\texttt{\{i32\ speed\_mmps,\ i24\ yaw\_mrad\_s,\ u8\ gear\}} \\
\texttt{0x301} & HOST\_BRAKE\_REQ & → &
\texttt{\{i32\ brake\_pressure\_kpa\}} \\
\texttt{0x302} & HOST\_LIGHT\_CMD & → & \texttt{\{u8\ bits\}} \\
\texttt{0x400} & HOST\_OBSTACLE\_DIST & → & \texttt{\{u32\ mm\}} \\
\texttt{0x600} & SYS\_DIAG\_RPT & ← (fwd) & \texttt{\{8\ bytes\}} \\
\texttt{0x7FC} & JETSON\_HEARTBEAT & → & \texttt{\{u8\ ctr\}} \\
\texttt{0x7FD} & RT\_HEARTBEAT & ← & \texttt{\{u8\ ctr\}} \\
\end{longtable}

\subsection{Low Bus (RT/SYS/Actuators) ---
UNCHANGED}\label{low-bus-rtsysactuators-unchanged}

\begin{longtable}[]{@{}llll@{}}
\toprule\noalign{}
ID & Name & Sender & Locked? \\
\midrule\noalign{}
\endhead
\bottomrule\noalign{}
\endlastfoot
\texttt{0x001} & SAFETY\_ESTOP & Any & --- \\
\texttt{0x011} & SYS\_SAFETY\_STS & SYS & --- \\
\texttt{0x110} & SYS\_MODE\_CMD & SYS & --- \\
\texttt{0x120} & SYS\_THROTTLE\_STS & MTR & --- \\
\texttt{0x169} & VCU\_SES\_REQ & RT & SYNTREE \\
\texttt{0x201} & SES\_STATUS & EPS-C & SYNTREE \\
\texttt{0x204} & RT\_DRIVE\_CMD & RT & --- \\
\texttt{0x205} & RT\_BRAKE\_CMD & RT & --- \\
\texttt{0x302} & HOST\_LIGHT\_CMD & RT (fwd) & --- \\
\texttt{0x721} & SEB\_STATUS & SEB & SYNTREE \\
\texttt{0x7B9} & VCU\_SEB\_REQ & RT/SYS (mode-gated) & SYNTREE \\
\texttt{0x7FD} & RT\_HEARTBEAT & RT & --- \\
\texttt{0x7FE} & SYS\_HEARTBEAT & SYS & --- \\
\end{longtable}

\begin{center}\rule{0.5\linewidth}{0.5pt}\end{center}

\emph{Version: 0.0.4-alpha. Optimal architecture --- minimum change for maximum
Autoware.Auto compatibility.}


\chapter{Autoware Communication Architecture \& Interfaces (Tabular
Reference)}\label{autoware-communication-architecture-interfaces-tabular-reference}

This document contains a structured reference of the communication interfaces,
CAN commands, internal states, and serialization protocols for the Autoware
vehicle interface stack.

\begin{center}\rule{0.5\linewidth}{0.5pt}\end{center}

\section{1. ROS 2 Topic Interfaces (C++ Lifecycle
Node)}\label{ros-2-topic-interfaces-c-lifecycle-node}

\begin{quote}
\textbf{Source Location:}
\texttt{vehicle\_interface/include/vehicle\_interface/vehicle\_interface\_node.hpp}
\end{quote}

\subsection{\texorpdfstring{A. Subscriptions (Inputs to
\texttt{vehicle\_interface})}{A. Subscriptions (Inputs to vehicle\_interface)}}\label{a.-subscriptions-inputs-to-vehicle_interface}

\begin{longtable}[]{@{}
  >{\raggedright\arraybackslash}p{(\columnwidth - 8\tabcolsep) * \real{0.2000}}
  >{\raggedright\arraybackslash}p{(\columnwidth - 8\tabcolsep) * \real{0.2000}}
  >{\raggedright\arraybackslash}p{(\columnwidth - 8\tabcolsep) * \real{0.2000}}
  >{\raggedright\arraybackslash}p{(\columnwidth - 8\tabcolsep) * \real{0.2000}}
  >{\raggedright\arraybackslash}p{(\columnwidth - 8\tabcolsep) * \real{0.2000}}@{}}
\toprule\noalign{}
\begin{minipage}[b]{\linewidth}\raggedright
Topic Name
\end{minipage} & \begin{minipage}[b]{\linewidth}\raggedright
ROS 2 Message Type
\end{minipage} & \begin{minipage}[b]{\linewidth}\raggedright
Key Member Fields
\end{minipage} & \begin{minipage}[b]{\linewidth}\raggedright
Internal Target Variable
\end{minipage} & \begin{minipage}[b]{\linewidth}\raggedright
Description
\end{minipage} \\
\midrule\noalign{}
\endhead
\bottomrule\noalign{}
\endlastfoot
\texttt{\textasciitilde{}/input/control\_cmd} &
\texttt{autoware\_control\_msgs/msg/Control} &
\texttt{lateral.steering\_tire\_angle}
(float)\texttt{lateral.steering\_tire\_rotation\_rate}
(float)\texttt{longitudinal.speed} (float)\texttt{longitudinal.acceleration}
(float) & \texttt{control\_cmd\_ptr\_} & Planning targets from the trajectory
follower pipeline \\
\texttt{\textasciitilde{}/input/actuation\_cmd} &
\texttt{tier4\_vehicle\_msgs/msg/ActuationCommandStamped} &
\texttt{actuation.accel\_cmd} (float)\texttt{actuation.brake\_cmd}
(float)\texttt{actuation.steer\_cmd} (float) & \texttt{actuation\_cmd\_ptr\_} &
Normalized pedal commands from \texttt{raw\_vehicle\_cmd\_converter} \\
\texttt{\textasciitilde{}/input/gear\_cmd} &
\texttt{autoware\_vehicle\_msgs/msg/GearCommand} & \texttt{command} (uint8) &
\texttt{gear\_cmd\_ptr\_} & Request gear shifts (e.g.~Drive, Reverse, Park) \\
\texttt{\textasciitilde{}/input/turn\_indicators\_cmd} &
\texttt{autoware\_vehicle\_msgs/msg/TurnIndicatorsCommand} & \texttt{command}
(uint8) & \texttt{turn\_indicators\_cmd\_ptr\_} & Request turn signal indicators
activation \\
\texttt{\textasciitilde{}/input/hazard\_lights\_cmd} &
\texttt{autoware\_vehicle\_msgs/msg/HazardLightsCommand} & \texttt{command}
(uint8) & \texttt{hazard\_lights\_cmd\_ptr\_} & Request hazard lights
activation \\
\texttt{\textasciitilde{}/input/emergency\_cmd} &
\texttt{tier4\_vehicle\_msgs/msg/VehicleEmergencyStamped} & \texttt{emergency}
(bool) & \texttt{emergency\_cmd\_ptr\_} & Triggers immediate safety stops and
system disengagement \\
\texttt{\textasciitilde{}/input/from\_can\_bus} & \texttt{can\_msgs::msg::Frame}
& \texttt{id} (uint32)\texttt{data} (uint8{[}8{]})\texttt{dlc} (uint8) & N/A
(Callback trigger) & Raw incoming frames received via SocketCAN \\
\end{longtable}

\begin{center}\rule{0.5\linewidth}{0.5pt}\end{center}

\subsection{\texorpdfstring{B. Publications (Outputs from
\texttt{vehicle\_interface})}{B. Publications (Outputs from vehicle\_interface)}}\label{b.-publications-outputs-from-vehicle_interface}

\begin{longtable}[]{@{}
  >{\raggedright\arraybackslash}p{(\columnwidth - 8\tabcolsep) * \real{0.2000}}
  >{\raggedright\arraybackslash}p{(\columnwidth - 8\tabcolsep) * \real{0.2000}}
  >{\raggedright\arraybackslash}p{(\columnwidth - 8\tabcolsep) * \real{0.2000}}
  >{\raggedright\arraybackslash}p{(\columnwidth - 8\tabcolsep) * \real{0.2000}}
  >{\raggedright\arraybackslash}p{(\columnwidth - 8\tabcolsep) * \real{0.2000}}@{}}
\toprule\noalign{}
\begin{minipage}[b]{\linewidth}\raggedright
Topic Name
\end{minipage} & \begin{minipage}[b]{\linewidth}\raggedright
ROS 2 Message Type
\end{minipage} & \begin{minipage}[b]{\linewidth}\raggedright
Key Member Fields
\end{minipage} & \begin{minipage}[b]{\linewidth}\raggedright
Internal Source Variable
\end{minipage} & \begin{minipage}[b]{\linewidth}\raggedright
Description
\end{minipage} \\
\midrule\noalign{}
\endhead
\bottomrule\noalign{}
\endlastfoot
\texttt{\textasciitilde{}/output/to\_can\_bus} & \texttt{can\_msgs::msg::Frame}
& \texttt{id} (uint32)\texttt{data} (uint8{[}8{]})\texttt{dlc} (uint8) & N/A
(Encoded inline) & Outgoing CAN frames dispatched to SocketCAN \\
\texttt{\textasciitilde{}/output/velocity\_status} &
\texttt{autoware\_vehicle\_msgs/msg/VelocityReport} &
\texttt{longitudinal\_velocity} (float)\texttt{lateral\_velocity}
(float)\texttt{heading\_rate} (float) & \texttt{speed\_status\_} & Speed and yaw
feedback for localized Odometry \\
\texttt{\textasciitilde{}/output/steering\_status} &
\texttt{autoware\_vehicle\_msgs/msg/SteeringReport} &
\texttt{steering\_tire\_angle} (float) &
\texttt{steering\_status\_.actual\_angle} & Confirms current tire direction
angle \\
\texttt{\textasciitilde{}/output/gear\_status} &
\texttt{autoware\_vehicle\_msgs/msg/GearReport} & \texttt{report} (uint8) &
\texttt{current\_gear\_} & Active transmission gear status \\
\texttt{\textasciitilde{}/output/control\_mode} &
\texttt{autoware\_vehicle\_msgs/msg/ControlModeReport} & \texttt{mode} (uint8) &
\texttt{is\_engaged\_} & Reports active mode (\texttt{1\ =\ Autonomous},
\texttt{4\ =\ Manual}) \\
\texttt{\textasciitilde{}/output/turn\_indicators\_status} &
\texttt{autoware\_vehicle\_msgs/msg/TurnIndicatorsReport} & \texttt{report}
(uint8) & N/A (Stubbed/Read from CAN) & Confirms physical turn signal status \\
\texttt{\textasciitilde{}/output/hazard\_lights\_status} &
\texttt{autoware\_vehicle\_msgs/msg/HazardLightsReport} & \texttt{report}
(uint8) & N/A (Stubbed/Read from CAN) & Confirms physical hazard lights
status \\
\texttt{\textasciitilde{}/output/actuation\_status} &
\texttt{tier4\_vehicle\_msgs/msg/ActuationStatusStamped} &
\texttt{status.accel\_status} (float)\texttt{status.brake\_status}
(float)\texttt{status.steer\_status} (float) &
\texttt{throttle\_status\_}\texttt{brake\_status\_}\texttt{steering\_status\_} &
Feedback of actual pedal strokes \\
\end{longtable}

\begin{center}\rule{0.5\linewidth}{0.5pt}\end{center}

\section{2. Default CAN Bus Protocol
Specifications}\label{default-can-bus-protocol-specifications}

\begin{quote}
\textbf{Source Location:}
\texttt{vehicle\_interface/include/vehicle\_interface/can\_protocol.hpp}
\end{quote}

\subsection{\texorpdfstring{A. Commands (Tier 2 Node \(\rightarrow\) Tier 3
Actuation
ECUs)}{A. Commands (Tier 2 Node \textbackslash rightarrow Tier 3 Actuation ECUs)}}\label{a.-commands-tier-2-node-rightarrow-tier-3-actuation-ecus}

\begin{longtable}[]{@{}
  >{\raggedright\arraybackslash}p{(\columnwidth - 8\tabcolsep) * \real{0.2000}}
  >{\raggedright\arraybackslash}p{(\columnwidth - 8\tabcolsep) * \real{0.2000}}
  >{\raggedright\arraybackslash}p{(\columnwidth - 8\tabcolsep) * \real{0.2000}}
  >{\raggedright\arraybackslash}p{(\columnwidth - 8\tabcolsep) * \real{0.2000}}
  >{\raggedright\arraybackslash}p{(\columnwidth - 8\tabcolsep) * \real{0.2000}}@{}}
\toprule\noalign{}
\begin{minipage}[b]{\linewidth}\raggedright
CAN ID
\end{minipage} & \begin{minipage}[b]{\linewidth}\raggedright
Signal Name
\end{minipage} & \begin{minipage}[b]{\linewidth}\raggedright
Bytes Used
\end{minipage} & \begin{minipage}[b]{\linewidth}\raggedright
Data Type
\end{minipage} & \begin{minipage}[b]{\linewidth}\raggedright
Encoding Formula / Values
\end{minipage} \\
\midrule\noalign{}
\endhead
\bottomrule\noalign{}
\endlastfoot
\textbf{\texttt{0x200}} & Steering target angle & \texttt{0..1} &
\texttt{int16\_t} &
\(\text{Raw} = (\delta_{\text{tire\_rad}} + 3.14159) / 0.001\) \\
& Steering target rate & \texttt{2..3} & \texttt{int16\_t} &
\(\text{Raw} = (\dot{\delta}_{\text{rad\_s}} + 10.0) / 0.01\) \\
& Steering enable flag & \texttt{4} & \texttt{uint8\_t} & \texttt{0x01} = Active
control, \texttt{0x00} = Manual bypass \\
& Unused & \texttt{5..7} & \texttt{uint8\_t} & \texttt{0x00} (Reserved) \\
\textbf{\texttt{0x201}} & Brake target pressure & \texttt{0..1} &
\texttt{uint16\_t} & \(\text{Raw} = P_{\text{brake\_ratio}} \times 10000.0\)
(Range: \(0.0 - 1.0\)) \\
& Brake enable flag & \texttt{2} & \texttt{uint8\_t} & \texttt{0x01} = Active
control, \texttt{0x00} = Manual bypass \\
& Unused & \texttt{3..7} & \texttt{uint8\_t} & \texttt{0x00} \\
\textbf{\texttt{0x202}} & Throttle target position & \texttt{0..1} &
\texttt{uint16\_t} & \(\text{Raw} = T_{\text{accel\_ratio}} \times 10000.0\)
(Range: \(0.0 - 1.0\)) \\
& Throttle enable flag & \texttt{2} & \texttt{uint8\_t} & \texttt{0x01} = Active
control, \texttt{0x00} = Manual bypass \\
& Unused & \texttt{3..7} & \texttt{uint8\_t} & \texttt{0x00} \\
\textbf{\texttt{0x203}} & Gear selection & \texttt{0} & \texttt{uint8\_t} &
\texttt{0} = NONE, \texttt{1} = NEUTRAL, \texttt{2} = DRIVE, \texttt{3} = LOW,
\texttt{10} = REVERSE, \texttt{20} = PARK \\
& Unused & \texttt{1..7} & \texttt{uint8\_t} & \texttt{0x00} \\
\textbf{\texttt{0x204}} & Turn Signal command & \texttt{0} & \texttt{uint8\_t} &
\texttt{0} = NONE, \texttt{1} = LEFT, \texttt{2} = RIGHT, \texttt{3} = HAZARD \\
& Unused & \texttt{1..7} & \texttt{uint8\_t} & \texttt{0x00} \\
\textbf{\texttt{0x206}} & System autonomous command & \texttt{0} &
\texttt{uint8\_t} & \texttt{0x01} = Enable override authority, \texttt{0x00} =
Disable \\
& Clear override flag & \texttt{1} & \texttt{uint8\_t} & \texttt{0x01} = Reset
override blockages, \texttt{0x00} = Neutral \\
& Unused & \texttt{2..7} & \texttt{uint8\_t} & \texttt{0x00} \\
\end{longtable}

\begin{center}\rule{0.5\linewidth}{0.5pt}\end{center}

\subsection{\texorpdfstring{B. Feedback Status (Tier 3 Actuation ECUs
\(\rightarrow\) Tier 2
Node)}{B. Feedback Status (Tier 3 Actuation ECUs \textbackslash rightarrow Tier 2 Node)}}\label{b.-feedback-status-tier-3-actuation-ecus-rightarrow-tier-2-node}

\begin{longtable}[]{@{}
  >{\raggedright\arraybackslash}p{(\columnwidth - 8\tabcolsep) * \real{0.2000}}
  >{\raggedright\arraybackslash}p{(\columnwidth - 8\tabcolsep) * \real{0.2000}}
  >{\raggedright\arraybackslash}p{(\columnwidth - 8\tabcolsep) * \real{0.2000}}
  >{\raggedright\arraybackslash}p{(\columnwidth - 8\tabcolsep) * \real{0.2000}}
  >{\raggedright\arraybackslash}p{(\columnwidth - 8\tabcolsep) * \real{0.2000}}@{}}
\toprule\noalign{}
\begin{minipage}[b]{\linewidth}\raggedright
CAN ID
\end{minipage} & \begin{minipage}[b]{\linewidth}\raggedright
Signal Name
\end{minipage} & \begin{minipage}[b]{\linewidth}\raggedright
Bytes Used
\end{minipage} & \begin{minipage}[b]{\linewidth}\raggedright
Data Type
\end{minipage} & \begin{minipage}[b]{\linewidth}\raggedright
Decoding Formula / Values
\end{minipage} \\
\midrule\noalign{}
\endhead
\bottomrule\noalign{}
\endlastfoot
\textbf{\texttt{0x300}} & Actual steer angle & \texttt{0..1} & \texttt{int16\_t}
& \(\theta_{\text{tire\_rad}} = (\text{Raw} \times 0.001) - 3.14159\) \\
& Actual steer rate & \texttt{2..3} & \texttt{int16\_t} &
\(\dot{\theta}_{\text{rad\_s}} = (\text{Raw} \times 0.01) - 10.0\) \\
& Steering system fault & \texttt{4} & \texttt{uint8\_t} & \texttt{0} = Nominal
(OK), \texttt{1} = Fault, \texttt{2} = Auto-Initializing \\
& Unused & \texttt{5..7} & \texttt{uint8\_t} & N/A \\
\textbf{\texttt{0x301}} & Actual brake pressure & \texttt{0..1} &
\texttt{uint16\_t} & \(P_{\text{actual}} = \text{Raw} / 10000.0\) \\
& Brake system fault & \texttt{2} & \texttt{uint8\_t} & \texttt{0} = Nominal,
\texttt{1} = Hardware fault \\
& Unused & \texttt{3..7} & \texttt{uint8\_t} & N/A \\
\textbf{\texttt{0x302}} & Actual throttle position & \texttt{0..1} &
\texttt{uint16\_t} & \(T_{\text{actual}} = \text{Raw} / 10000.0\) \\
& Throttle system fault & \texttt{2} & \texttt{uint8\_t} & \texttt{0} = Nominal,
\texttt{1} = Hardware fault \\
& Unused & \texttt{3..7} & \texttt{uint8\_t} & N/A \\
\textbf{\texttt{0x303}} & Actual gear & \texttt{0} & \texttt{uint8\_t} &
\texttt{0} = NONE, \texttt{1} = NEUTRAL, \texttt{2} = DRIVE, \texttt{3} = LOW,
\texttt{10} = REVERSE, \texttt{20} = PARK \\
& Unused & \texttt{1..7} & \texttt{uint8\_t} & N/A \\
\textbf{\texttt{0x320}} & Longitudinal Speed & \texttt{0..1} & \texttt{int16\_t}
& \(v = \text{Raw} \times 0.01\) {[}m/s{]} \\
& Yaw rate (Heading rate) & \texttt{2..3} & \texttt{int16\_t} &
\(\omega = \text{Raw} \times 0.001\) {[}rad/s{]} \\
& Unused & \texttt{4..7} & \texttt{uint8\_t} & N/A \\
\textbf{\texttt{0x340}} & Core System Status flags & \texttt{0} &
\texttt{uint8\_t} & Bit \texttt{0} (0x01): \texttt{autonomous\_enabled}Bit
\texttt{1} (0x02): \texttt{override\_active}Bit \texttt{2} (0x04):
\texttt{system\_fault}Bit \texttt{3} (0x08): \texttt{estop\_active} \\
& Diagnostic Fault code & \texttt{1} & \texttt{uint8\_t} & Actuator or
microcontroller fault code value \\
& Unused & \texttt{2..7} & \texttt{uint8\_t} & N/A \\
\end{longtable}

\begin{center}\rule{0.5\linewidth}{0.5pt}\end{center}

\section{3. Node Configuration Parameters (YAML
Config)}\label{node-configuration-parameters-yaml-config}

\begin{quote}
\textbf{Source Location:}
\texttt{vehicle\_interface/config/vehicle\_interface.param.yaml}
\end{quote}

Parameters configured inside the YAML map to the C++ struct \texttt{Params}:

\begin{longtable}[]{@{}
  >{\raggedright\arraybackslash}p{(\columnwidth - 8\tabcolsep) * \real{0.2000}}
  >{\raggedright\arraybackslash}p{(\columnwidth - 8\tabcolsep) * \real{0.2000}}
  >{\raggedright\arraybackslash}p{(\columnwidth - 8\tabcolsep) * \real{0.2000}}
  >{\raggedright\arraybackslash}p{(\columnwidth - 8\tabcolsep) * \real{0.2000}}
  >{\raggedright\arraybackslash}p{(\columnwidth - 8\tabcolsep) * \real{0.2000}}@{}}
\toprule\noalign{}
\begin{minipage}[b]{\linewidth}\raggedright
Parameter Name
\end{minipage} & \begin{minipage}[b]{\linewidth}\raggedright
Data Type
\end{minipage} & \begin{minipage}[b]{\linewidth}\raggedright
Default Value
\end{minipage} & \begin{minipage}[b]{\linewidth}\raggedright
Units / Range
\end{minipage} & \begin{minipage}[b]{\linewidth}\raggedright
Description
\end{minipage} \\
\midrule\noalign{}
\endhead
\bottomrule\noalign{}
\endlastfoot
\texttt{loop\_rate} & \texttt{double} & \texttt{30.0} & Hz & Publish command
frame execution frequency \\
\texttt{command\_timeout\_ms} & \texttt{int} & \texttt{1000} & Milliseconds &
Triggers emergency stop on command loss \\
\texttt{max\_throttle} & \texttt{double} & \texttt{0.4} & \(0.0 - 1.0\) (Pedal
ratio) & Upper throttle command clamp \\
\texttt{max\_brake} & \texttt{double} & \texttt{0.8} & \(0.0 - 1.0\) (Pedal
ratio) & Upper brake pressure command clamp \\
\texttt{emergency\_brake} & \texttt{double} & \texttt{0.7} & \(0.0 - 1.0\)
(Pedal ratio) & Brake pressure applied during failsafe events \\
\texttt{max\_steering\_angle} & \texttt{double} & \texttt{1.0} & Radians &
Maximum allowable physical tire angle \\
\texttt{max\_steering\_rate} & \texttt{double} & \texttt{5.0} & Rad/sec &
Steering rotation velocity limit \\
\texttt{steering\_offset} & \texttt{double} & \texttt{0.0} & Radians &
Calibration alignment adjustment \\
\texttt{vgr\_coef\_a} & \texttt{double} & \texttt{15.713} & Constant & Variable
Gear Ratio: base offset coefficient \\
\texttt{vgr\_coef\_b} & \texttt{double} & \texttt{0.053} & Constant & Variable
Gear Ratio: velocity scaling coefficient \\
\texttt{vgr\_coef\_c} & \texttt{double} & \texttt{0.042} & Constant & Variable
Gear Ratio: angle scaling coefficient \\
\texttt{wheel\_base} & \texttt{double} & \texttt{2.79} & Meters & Axle
center-to-center wheelbase span \\
\texttt{wheel\_radius} & \texttt{double} & \texttt{0.383} & Meters & Active
rolling wheel radius \\
\texttt{margin\_time\_for\_gear\_change} & \texttt{double} & \texttt{2.0} &
Seconds & Anti-gear-chattering shift lockout window \\
\texttt{base\_frame\_id} & \texttt{std::string} & \texttt{"base\_link"} & N/A &
TF kinematic frame ID of the vehicle \\
\texttt{can\_interface} & \texttt{std::string} & \texttt{"can0"} & N/A &
Destination physical socket name \\
\texttt{use\_actuation\_cmd} & \texttt{bool} & \texttt{true} & \texttt{true} /
\texttt{false} & True: Use pedal mappings. False: Bypass to Control \\
\texttt{convert\_steer\_cmd} & \texttt{bool} & \texttt{true} & \texttt{true} /
\texttt{false} & Enable VGR steering wheel angle conversion \\
\end{longtable}

\begin{center}\rule{0.5\linewidth}{0.5pt}\end{center}

\section{\texorpdfstring{4. micro-ROS Internal Events
(\texttt{EventsMicroAutowareHandle})}{4. micro-ROS Internal Events (EventsMicroAutowareHandle)}}\label{micro-ros-internal-events-eventsmicroautowarehandle}

\begin{quote}
\textbf{Source Locations:} *
\texttt{vehicle\_interface/microAutoware/src/microAutoware/microAutoware.h} *
\texttt{vehicle\_interface/microAutoware/src/nucleo-H753ZI\_microAutoware\_Serial\_HIL/Core/Inc/microAutoware.h}
(STM32 Firmware Project Copy)
\end{quote}

Event flags used within the FreeRTOS middleware tasks
(\texttt{StartMicroAutoware} and \texttt{StartTaskControl}):

\begin{longtable}[]{@{}
  >{\raggedright\arraybackslash}p{(\columnwidth - 8\tabcolsep) * \real{0.2000}}
  >{\raggedright\arraybackslash}p{(\columnwidth - 8\tabcolsep) * \real{0.2000}}
  >{\raggedright\arraybackslash}p{(\columnwidth - 8\tabcolsep) * \real{0.2000}}
  >{\raggedright\arraybackslash}p{(\columnwidth - 8\tabcolsep) * \real{0.2000}}
  >{\raggedright\arraybackslash}p{(\columnwidth - 8\tabcolsep) * \real{0.2000}}@{}}
\toprule\noalign{}
\begin{minipage}[b]{\linewidth}\raggedright
Constant Name
\end{minipage} & \begin{minipage}[b]{\linewidth}\raggedright
Value (Binary)
\end{minipage} & \begin{minipage}[b]{\linewidth}\raggedright
Value (Hex)
\end{minipage} & \begin{minipage}[b]{\linewidth}\raggedright
Action Trigger / Target
\end{minipage} & \begin{minipage}[b]{\linewidth}\raggedright
Description
\end{minipage} \\
\midrule\noalign{}
\endhead
\bottomrule\noalign{}
\endlastfoot
\texttt{MA\_TO\_AUTOWARE\_MODE\_FLAG} & \texttt{0b0000000010} & \texttt{0x002} &
Control Task \(\rightarrow\) micro-ROS & Requests shift of current control
execution to Autoware \\
\texttt{MA\_TO\_MANUAL\_MODE\_FLAG} & \texttt{0b0000000100} & \texttt{0x004} &
Control Task \(\rightarrow\) micro-ROS & Requests fallback shift to manual
control (joystick) \\
\texttt{MA\_TO\_EMERGENCY\_MODE\_FLAG} & \texttt{0b0000001000} & \texttt{0x008}
& Control Task \(\rightarrow\) micro-ROS & Triggers fallback shift to Emergency
Stop control status \\
\texttt{SYS\_TO\_AUTOWARE\_MODE\_FLAG} & \texttt{0b0000010000} & \texttt{0x010}
& micro-ROS \(\rightarrow\) Control Task & Acknowledges system state mode is now
Autoware \\
\texttt{SYS\_TO\_MANUAL\_MODE\_FLAG} & \texttt{0b0000100000} & \texttt{0x020} &
micro-ROS \(\rightarrow\) Control Task & Acknowledges system state mode is now
Manual \\
\texttt{SYS\_TO\_EMERGENCY\_MODE\_FLAG} & \texttt{0b0001000000} & \texttt{0x040}
& micro-ROS \(\rightarrow\) Control Task & Acknowledges system state mode is now
Emergency \\
\texttt{VEHICLE\_NEW\_DATA\_FLAG} & \texttt{0b0010000000} & \texttt{0x080} &
Control Task \(\rightarrow\) micro-ROS & Signifies new HIL/CARLA feedback parsed
from UART \\
\texttt{AUTOWARE\_NEW\_DATA\_FLAG} & \texttt{0b0100000000} & \texttt{0x100} &
micro-ROS \(\rightarrow\) Control Task & Signifies new control command message
received \\
\texttt{MICRO\_ROS\_AGENT\_ONLINE\_FLAG} & \texttt{0b1000000000} &
\texttt{0x200} & micro-ROS \(\rightarrow\) Control Task & Set when handshake
ping checks successfully with agent \\
\end{longtable}

\begin{center}\rule{0.5\linewidth}{0.5pt}\end{center}

\section{\texorpdfstring{5. HIL Serial Communication Protocol (STM32
\(\leftrightarrow\)
CARLA)}{5. HIL Serial Communication Protocol (STM32 \textbackslash leftrightarrow CARLA)}}\label{hil-serial-communication-protocol-stm32-leftrightarrow-carla}

\begin{quote}
\textbf{Source Locations:} *
\texttt{vehicle\_interface/microAutoware/src/HIL/taksControl.c} /
\texttt{utils.c} *
\texttt{vehicle\_interface/microAutoware/src/nucleo-H753ZI\_microAutoware\_Serial\_HIL/Core/Src/taksControl.c}
/ \texttt{utils.c} (STM32 Firmware Project Copy)
\end{quote}

\subsection{\texorpdfstring{A. Command Message Packet Structure (STM32
\(\rightarrow\)
CARLA)}{A. Command Message Packet Structure (STM32 \textbackslash rightarrow CARLA)}}\label{a.-command-message-packet-structure-stm32-rightarrow-carla}

\emph{Total Packet Frame: 30 Bytes}

\begin{longtable}[]{@{}
  >{\raggedright\arraybackslash}p{(\columnwidth - 8\tabcolsep) * \real{0.2000}}
  >{\raggedright\arraybackslash}p{(\columnwidth - 8\tabcolsep) * \real{0.2000}}
  >{\raggedright\arraybackslash}p{(\columnwidth - 8\tabcolsep) * \real{0.2000}}
  >{\raggedright\arraybackslash}p{(\columnwidth - 8\tabcolsep) * \real{0.2000}}
  >{\raggedright\arraybackslash}p{(\columnwidth - 8\tabcolsep) * \real{0.2000}}@{}}
\toprule\noalign{}
\begin{minipage}[b]{\linewidth}\raggedright
Byte Index
\end{minipage} & \begin{minipage}[b]{\linewidth}\raggedright
ASCII Tag Prefix
\end{minipage} & \begin{minipage}[b]{\linewidth}\raggedright
Data Payload Field
\end{minipage} & \begin{minipage}[b]{\linewidth}\raggedright
Byte Type
\end{minipage} & \begin{minipage}[b]{\linewidth}\raggedright
Description
\end{minipage} \\
\midrule\noalign{}
\endhead
\bottomrule\noalign{}
\endlastfoot
\texttt{0} & \texttt{\#} (0x23) & Frame Start Tag & \texttt{char} & Marks start
of transmission frame \\
\texttt{1} & \texttt{S} (0x53) & Steering Angle prefix & \texttt{char} &
Identifies following float payload \\
\texttt{2..5} & N/A & \texttt{xSteeringAngle} & \texttt{float} (4 Bytes) &
Target steering tire angle command {[}rad{]} \\
\texttt{6} & \texttt{W} (0x57) & Steering Velocity prefix & \texttt{char} &
Identifies following float payload \\
\texttt{7..10} & N/A & \texttt{xSteeringVelocity} & \texttt{float} (4 Bytes) &
Target steering angular rate {[}rad/s{]} \\
\texttt{11} & \texttt{V} (0x56) & Target Speed prefix & \texttt{char} &
Identifies following float payload \\
\texttt{12..15} & N/A & \texttt{xSpeed} & \texttt{float} (4 Bytes) &
Longitudinal target velocity {[}m/s{]} \\
\texttt{16} & \texttt{A} (0x41) & Target Accel prefix & \texttt{char} &
Identifies following float payload \\
\texttt{17..20} & N/A & \texttt{xAcceleration} & \texttt{float} (4 Bytes) &
Target acceleration command {[}m/s²{]} \\
\texttt{21} & \texttt{J} (0x4A) & Target Jerk prefix & \texttt{char} &
Identifies following float payload \\
\texttt{22..25} & N/A & \texttt{xJerk} & \texttt{float} (4 Bytes) & Target jerk
command {[}m/s³{]} \\
\texttt{26} & \texttt{M} (0x4D) & Control Mode prefix & \texttt{char} &
Identifies following control mode byte \\
\texttt{27} & N/A & \texttt{ucControlMode} & \texttt{uint8\_t} (1 Byte) &
Current mode: \texttt{0\ =\ Emergency}, \texttt{1\ =\ Autoware},
\texttt{4\ =\ Manual} \\
\texttt{28} & \texttt{\$} (0x24) & Frame Stop Tag & \texttt{char} & Marks end of
transmission frame \\
\texttt{29} & \texttt{\textbackslash{}0} (0x00) & String Terminator &
\texttt{char} & Null terminator character \\
\end{longtable}

\begin{center}\rule{0.5\linewidth}{0.5pt}\end{center}

\subsection{\texorpdfstring{B. Telemetry Message Packet Structure (CARLA
\(\rightarrow\)
STM32)}{B. Telemetry Message Packet Structure (CARLA \textbackslash rightarrow STM32)}}\label{b.-telemetry-message-packet-structure-carla-rightarrow-stm32}

\emph{Total Packet Frame: 22 Bytes}

\begin{longtable}[]{@{}
  >{\raggedright\arraybackslash}p{(\columnwidth - 8\tabcolsep) * \real{0.2000}}
  >{\raggedright\arraybackslash}p{(\columnwidth - 8\tabcolsep) * \real{0.2000}}
  >{\raggedright\arraybackslash}p{(\columnwidth - 8\tabcolsep) * \real{0.2000}}
  >{\raggedright\arraybackslash}p{(\columnwidth - 8\tabcolsep) * \real{0.2000}}
  >{\raggedright\arraybackslash}p{(\columnwidth - 8\tabcolsep) * \real{0.2000}}@{}}
\toprule\noalign{}
\begin{minipage}[b]{\linewidth}\raggedright
Byte Index
\end{minipage} & \begin{minipage}[b]{\linewidth}\raggedright
ASCII Tag Prefix
\end{minipage} & \begin{minipage}[b]{\linewidth}\raggedright
Data Payload Field
\end{minipage} & \begin{minipage}[b]{\linewidth}\raggedright
Byte Type
\end{minipage} & \begin{minipage}[b]{\linewidth}\raggedright
Description
\end{minipage} \\
\midrule\noalign{}
\endhead
\bottomrule\noalign{}
\endlastfoot
\texttt{0} & \texttt{\#} (0x23) & Frame Start Tag & \texttt{char} & Marks start
of incoming telemetry frame \\
\texttt{1} & \texttt{A} (0x41) & Longitudinal Speed prefix & \texttt{char} &
Identifies following speed float payload \\
\texttt{2..5} & N/A & \texttt{xLongSpeed} & \texttt{float} (4 Bytes) & Simulated
longitudinal vehicle velocity {[}m/s{]} \\
\texttt{6} & \texttt{B} (0x42) & Lateral Speed prefix & \texttt{char} &
Identifies following lateral speed float \\
\texttt{7..10} & N/A & \texttt{xLatSpeed} & \texttt{float} (4 Bytes) & Simulated
lateral vehicle velocity {[}m/s{]} \\
\texttt{11} & \texttt{C} (0x43) & Yaw Rate prefix & \texttt{char} & Identifies
following heading yaw rate float \\
\texttt{12..15} & N/A & \texttt{xHeadingRate} & \texttt{float} (4 Bytes) &
Simulated yaw rate feedback {[}rad/s{]} \\
\texttt{16} & \texttt{D} (0x44) & Steering Angle feedback & \texttt{char} &
Identifies steering wheel angle status float \\
\texttt{17..20} & N/A & \texttt{xSteeringStatus} & \texttt{float} (4 Bytes) &
Simulated actual steering wheel angle {[}rad{]} \\
\texttt{21} & \texttt{\$} (0x24) & Frame Stop Tag & \texttt{char} & Marks end of
telemetry frame \\
\end{longtable}

\begin{center}\rule{0.5\linewidth}{0.5pt}\end{center}

\section{6. Safety Limits \& Command Validation
Rules}\label{safety-limits-command-validation-rules}

\begin{quote}
\textbf{Source Locations:} *
\texttt{vehicle\_interface/src/vehicle\_interface\_node.cpp} (C++ Node) *
\texttt{vehicle\_interface/microAutoware/src/HIL/taksControl.c} (STM32 Firmware)
*
\texttt{vehicle\_interface/microAutoware/src/nucleo-H753ZI\_microAutoware\_Serial\_HIL/Core/Src/taksControl.c}
(STM32 Firmware Project Copy)
\end{quote}

The interface layers enforce strict validation bounds, watchdogs, and override
logic to guarantee system stability and passenger safety:

\begin{longtable}[]{@{}
  >{\raggedright\arraybackslash}p{(\columnwidth - 8\tabcolsep) * \real{0.2000}}
  >{\raggedright\arraybackslash}p{(\columnwidth - 8\tabcolsep) * \real{0.2000}}
  >{\raggedright\arraybackslash}p{(\columnwidth - 8\tabcolsep) * \real{0.2000}}
  >{\raggedright\arraybackslash}p{(\columnwidth - 8\tabcolsep) * \real{0.2000}}
  >{\raggedright\arraybackslash}p{(\columnwidth - 8\tabcolsep) * \real{0.2000}}@{}}
\toprule\noalign{}
\begin{minipage}[b]{\linewidth}\raggedright
Rule Category
\end{minipage} & \begin{minipage}[b]{\linewidth}\raggedright
Applied Value / Range
\end{minipage} & \begin{minipage}[b]{\linewidth}\raggedright
Check Description
\end{minipage} & \begin{minipage}[b]{\linewidth}\raggedright
Action Taken on Violation
\end{minipage} & \begin{minipage}[b]{\linewidth}\raggedright
Source Node / Task
\end{minipage} \\
\midrule\noalign{}
\endhead
\bottomrule\noalign{}
\endlastfoot
\textbf{Throttle Clamping} & \texttt{{[}0.0,\ max\_throttle{]}} (Default max:
\texttt{0.4}) & Clamps input target throttle within mechanical bounds & Values
exceeding are capped at \texttt{max\_throttle} & \texttt{vehicle\_interface}
(C++) \\
\textbf{Brake Clamping} & \texttt{{[}0.0,\ max\_brake{]}} (Default max:
\texttt{0.8}) & Clamps input target brake within system pressure bounds & Values
exceeding are capped at \texttt{max\_brake} & \texttt{vehicle\_interface}
(C++) \\
\textbf{Steering Clamping} &
\texttt{{[}-max\_steering\_angle,\ max\_steering\_angle{]}} (Default max:
\texttt{1.0} rad) & Clamps input target tire angle to prevent physical lock
limits & Angles exceeding are capped at limits & \texttt{vehicle\_interface}
(C++) \\
\textbf{Steering Rate Limit} & \texttt{max\_steering\_rate} (Default max:
\texttt{5.0} rad/s) & Limits maximum target steering rotation rate & Transmitted
as target rate field in CAN command; actual rate-limiting is handled by Tier 3
ECU & \texttt{vehicle\_interface} (C++) \\
\textbf{Gear Anti-Chatter} & \texttt{margin\_time\_for\_gear\_change} (Default:
\texttt{2.0} sec) & Evaluates time elapsed since last gear state switch & Gear
requests are ignored if time delta is less & \texttt{vehicle\_interface}
(C++) \\
\textbf{Command Timeout} & \texttt{command\_timeout\_ms} (Default: \texttt{1000}
ms) & Checks time elapsed since last incoming ROS 2 command & Applies emergency
brake (\texttt{emergency\_brake}, default \texttt{0.7}) and zero throttle &
\texttt{vehicle\_interface} (C++) \\
\textbf{Manual Override} & \texttt{override\_active} flag (\texttt{0x340} CAN
Status bit 1) & Monitors driver manual steering/pedal override inputs &
Disengages autonomy (\texttt{is\_engaged\_} = false) and turns off CAN enable
flags & \texttt{vehicle\_interface} (C++) \\
\textbf{System Fault} & \texttt{system\_fault} flag (\texttt{0x340} CAN Status
bit 2) & Monitors hardware faults flagged by actuator ECUs & Disengages autonomy
immediately & \texttt{vehicle\_interface} (C++) \\
\textbf{HIL Autoware Command Timeout} & \texttt{TIMEOUT\_GET\_CONTROL\_ACTION}
(\texttt{110} ms) & Monitors incoming time difference between micro-ROS commands
& Increments data loss; triggers MANUAL mode fallback after \texttt{10}
consecutive drops & \texttt{microAutoware} (STM32 task) \\
\textbf{HIL Telemetry Status Timeout} & \texttt{TIMEOUT\_GET\_CARLA\_RX}
(\texttt{100} ms) & Monitors time difference between serial CARLA feedback
frames & Increments loss; triggers EMERGENCY stop mode after \texttt{10}
consecutive drops & \texttt{microAutoware} (STM32 task) \\
\textbf{Joystick Debounce} & \texttt{DEBOUNCE\_TICKS} (\texttt{1000} ticks) &
Avoids chatter when pressing the physical JoySW button & Ignores transition
signals sent within this interval & \texttt{microAutoware} (STM32 task) \\
\end{longtable}


\chapter{Brake System --- End-to-End
Function}\label{brake-system-end-to-end-function}

The E-Trike uses a \textbf{SYNTREE SEB} electro-hydraulic brake actuator
commanded via CAN. There is no mechanical cable from lever to master cylinder
--- the SEB is the sole hydraulic path. This document describes how braking
works in each operating mode, how the two control modes (Stroke and Pressure)
differ, and how multiple brake sources arbitrate to a single actuator.

\begin{center}\rule{0.5\linewidth}{0.5pt}\end{center}

\section{1. Physical path}\label{physical-path}

\begin{verbatim}
Brake lever (GPIO2, active-low) ──► SYS ESP32-S3
                                         │
                                    CAN 0x7B9 (50 Hz) ──► SYNTREE SEB ──► hydraulic master cylinder
                                                                              │
                                                                         calipers (front + rear)
\end{verbatim}

\begin{center}\rule{0.5\linewidth}{0.5pt}\end{center}

\section{2. Two control modes --- same frame, different
fields}\label{two-control-modes-same-frame-different-fields}

\texttt{0x7B9\ VCU\_SEB\_REQ} carries both stroke and pressure targets in every
frame. The \texttt{VCU\_SEB\_Control\_Mode} bit (byte 0, bit 2) selects which
one the SEB acts on:

\begin{longtable}[]{@{}
  >{\raggedright\arraybackslash}p{(\columnwidth - 8\tabcolsep) * \real{0.1277}}
  >{\raggedright\arraybackslash}p{(\columnwidth - 8\tabcolsep) * \real{0.1489}}
  >{\raggedright\arraybackslash}p{(\columnwidth - 8\tabcolsep) * \real{0.1277}}
  >{\raggedright\arraybackslash}p{(\columnwidth - 8\tabcolsep) * \real{0.3617}}
  >{\raggedright\arraybackslash}p{(\columnwidth - 8\tabcolsep) * \real{0.2340}}@{}}
\toprule\noalign{}
\begin{minipage}[b]{\linewidth}\raggedright
Mode
\end{minipage} & \begin{minipage}[b]{\linewidth}\raggedright
Field
\end{minipage} & \begin{minipage}[b]{\linewidth}\raggedright
Type
\end{minipage} & \begin{minipage}[b]{\linewidth}\raggedright
What it commands
\end{minipage} & \begin{minipage}[b]{\linewidth}\raggedright
Used when
\end{minipage} \\
\midrule\noalign{}
\endhead
\bottomrule\noalign{}
\endlastfoot
\textbf{Stroke (1)} & \texttt{VCU\_SEB\_Stroke\_Value\_Req} & u16, 0.05 mm/bit,
offset -30 & Pushrod moves to an exact physical position & MANUAL lever, ESTOP,
AUTO fallback \\
\textbf{Pressure (2)} & \texttt{VCU\_SEB\_Pre\_Value\_Req} & u8, 0.05 MPa/bit,
0--5 MPa & SEB's internal PID holds a target hydraulic pressure & AUTO modulated
braking via \texttt{0x205} \\
\end{longtable}

Both fields are present in every frame. The mode bit tells the SEB which one to
obey. SYS can switch modes mid-operation by changing the bit --- the SEB
transitions on the next received frame.

\textbf{Why Pressure Mode is better for autonomous braking:} Stroke Mode
commands a position. Pad wear, temperature, and fluid viscosity change the
relationship between stroke and braking force. At 15 mm stroke with cold thick
pads you get \textasciitilde2 MPa; same 15 mm with hot thin pads you get
\textasciitilde0.5 MPa. Pressure Mode commands force directly --- the SEB's
internal PID moves the pushrod until the pressure sensor reads the target,
compensating automatically.

\begin{center}\rule{0.5\linewidth}{0.5pt}\end{center}

\section{3. MANUAL mode --- binary braking}\label{manual-mode-binary-braking}

\begin{verbatim}
Rider squeezes lever
       │
       ▼  GPIO2 reads LOW (active-low)
       │
  brake_task (20 Hz): lever_pressed
       │
       ▼  send 0x7B9 {mode=Stroke, stroke=15mm} at 50 Hz
       │
  SEB pushrod → 15mm → hydraulic pressure builds → calipers clamp
       │
       ▼  brake_light GPIO21 = ON (OR logic)
\end{verbatim}

When the rider releases, SYS sends \texttt{\{mode=Stroke,\ stroke=0mm\}}. The
SEB releases pressure. The master cylinder spring-returns the pushrod.

Manual braking is \textbf{binary} --- 15 mm or 0 mm. No modulation. The rider
controls stopping distance by pulsing the lever, not by varying pressure. This
is acceptable because in MANUAL mode the rider also controls speed via the
throttle.

\begin{center}\rule{0.5\linewidth}{0.5pt}\end{center}

\section{4. AUTO mode --- three sources, one
actuator}\label{auto-mode-three-sources-one-actuator}

AUTO has three independent brake triggers converging on the same \texttt{0x7B9}
frame. Priority order:

\subsection{Priority: ESTOP \textgreater{} lever \textgreater{} modulated
\textgreater{} released}\label{priority-estop-lever-modulated-released}

\begin{verbatim}
brake_task @ 20 Hz:

  if ESTOP:
      send 0x7B9 {mode=Stroke, stroke=27mm}   ← full lock

  elif lever_pressed:
      send 0x7B9 {mode=Stroke, stroke=15mm}   ← driver always wins

  elif g_brake_pressure_kpa > 0:              ← from 0x205 RT_BRAKE_CMD
      raw = (uint8_t)(kpa * 0.02)             ← 5000 kPa → raw=100 (5 MPa)
      raw = clamp(raw, 0, 100)
      send 0x7B9 {mode=Pressure, pressure=raw} ← modulated braking

  else:
      send 0x7B9 {mode=Stroke, stroke=0mm}    ← released
\end{verbatim}

\textbf{The driver override is absolute.} If the rider squeezes the lever, SYS
switches back to Stroke Mode at 15 mm regardless of what \texttt{0x205}
commands. The human always wins.

\subsection{4.1 Source 1: Brake lever (driver
override)}\label{source-1-brake-lever-driver-override}

Identical to MANUAL mode. The lever is a binary switch --- pressed = 15 mm,
released = 0 mm. Available in all modes.

\subsection{\texorpdfstring{4.2 Source 2: \texttt{0x205\ RT\_BRAKE\_CMD}
(modulated autonomous
braking)}{4.2 Source 2: 0x205 RT\_BRAKE\_CMD (modulated autonomous braking)}}\label{source-2-0x205-rt_brake_cmd-modulated-autonomous-braking}

\begin{verbatim}
Jetson perception stack
       │
       ▼  "obstacle at 15m, decelerate at 300 kPa"
  0x301 HOST_BRAKE_REQ {300 kPa}  (high CAN, on demand)
       │
       ▼
RT dispatch_task:
   obstacle_distance = read HC-SR04
   rt_obstacle_kpa   = obstacle_to_brake(obstacle_distance)
   brake_kpa = max(rt_obstacle_kpa, jetson_301_kpa)
       │
       ▼
RT can_tx_low_task:
   send 0x205 {brake_kpa} at 50 Hz  (low CAN)
       │
       ▼
SYS dispatch_task:
   g_brake_pressure_kpa = 0x205 value  (atomic)
       │
       ▼
SYS brake_task:
   raw = (uint8_t)(kpa × 0.02)
   send 0x7B9 {mode=Pressure, pressure=raw}
       │
       ▼
SEB internal PID: holds target pressure
\end{verbatim}

RT uses max-select arbitration: the worse (higher) pressure between Jetson's
deceleration request and the obstacle sensor's emergency brake wins.

\subsection{4.3 Source 3: ESTOP}\label{source-3-estop}

Always Stroke Mode at 27 mm (maximum physical stroke). This path bypasses every
software decision --- it doesn't check \texttt{0x205}, doesn't check the lever,
doesn't care about Pressure Mode. It's the emergency stop.

\begin{center}\rule{0.5\linewidth}{0.5pt}\end{center}

\section{5. ESTOP mode --- maximum brake, no
questions}\label{estop-mode-maximum-brake-no-questions}

\begin{verbatim}
ESTOP trigger (button / CAN 0x001 / heartbeat timeout)
       │
       ▼
SYS brake_task:
   send 0x7B9 {mode=Stroke, stroke=27mm}   ← maximum physical stroke
       │
       ▼
SEB: full hydraulic pressure → maximum braking force
       │
       ▼
brake_light GPIO21 = ON (forced)
\end{verbatim}

27 mm is the SEB's physical maximum (raw 1140). This is independent of all other
brake sources.

\begin{center}\rule{0.5\linewidth}{0.5pt}\end{center}

\section{6. Mode-switching protocol}\label{mode-switching-protocol}

\subsection{Stroke → Pressure (0x205 transitions 0 →
positive)}\label{stroke-pressure-0x205-transitions-0-positive}

\begin{verbatim}
1. SYS is in Stroke Mode, stroke=0mm (released)
2. 0x205 arrives with positive kPa value
3. SYS: hold current stroke position (already released)
4. SYS: switch 0x7B9 mode bit from 1 (Stroke) to 2 (Pressure)
5. SYS: set VCU_SEB_Pre_Value_Req to converted raw value
6. SEB: receives frame, sees mode=Pressure, starts PID loop from 0 MPa → target
\end{verbatim}

\subsection{Pressure → Stroke (0x205 drops to
0)}\label{pressure-stroke-0x205-drops-to-0}

\begin{verbatim}
1. SYS: switch mode bit from 2 (Pressure) back to 1 (Stroke)
2. SYS: set stroke=0mm (released)
3. SEB: releases pressure
\end{verbatim}

The hold-then-switch pattern prevents pressure transients --- the SEB doesn't
jump between mode interpretations of the same command bytes.

\begin{center}\rule{0.5\linewidth}{0.5pt}\end{center}

\section{7. Conversion formula}\label{conversion-formula}

\begin{verbatim}
kPa → SEB raw:  raw = kPa × 0.02

Constants:
  kSebPressureScale  = 0.02f       (1 bit = 0.05 MPa; 1 MPa = 1000 kPa)
  kSebPressureOffset = 0.0f
  kSebMaxPressureRaw = 100         (100 × 0.05 = 5.0 MPa hardware limit)

Example conversions:
  1000 kPa → 1000 × 0.02 = 20   → 1.0 MPa (gentle deceleration)
  3000 kPa → 3000 × 0.02 = 60   → 3.0 MPa (firm stop)
  5000 kPa → 5000 × 0.02 = 100  → 5.0 MPa (maximum, emergency)
\end{verbatim}

Verified against SYNTREE SEB CAN protocol specification.
\texttt{VCU\_SEB\_Pre\_Value\_Req} is u8 at bit 32 of \texttt{0x7B9}, scale 0.05
MPa/bit, range 0--5 MPa.

\begin{center}\rule{0.5\linewidth}{0.5pt}\end{center}

\section{8. Stroke value reference}\label{stroke-value-reference}

\begin{verbatim}
raw = (physical_mm + 30.0) / 0.05

| Physical | Raw  | Use case              |
|----------|------|-----------------------|
| -5 mm    | 500  | Minimum               |
| 0 mm     | 600  | Released              |
| 15 mm    | 900  | Manual lever pressed  |
| 27 mm    | 1140 | ESTOP full brake      |
\end{verbatim}

\begin{center}\rule{0.5\linewidth}{0.5pt}\end{center}

\section{9. Brake light logic}\label{brake-light-logic}

\begin{verbatim}
brake_light_on = lever_pressed()           // GPIO2 — physical switch
              OR (mode == ESTOP)           // emergency
              OR g_light_state.brake_light // Jetson 0x302 (predictive/hazard)
\end{verbatim}

All three sources are local to SYS --- no CAN round-trip needed for the lever or
ESTOP. The Jetson bit is supplemental (predictive illumination before pressure
builds) and can never be the only trigger. Physical braking state always wins.

\begin{center}\rule{0.5\linewidth}{0.5pt}\end{center}

\section{10. Boot sequence}\label{boot-sequence}

\begin{verbatim}
Power-on
    │
    ▼
BRAKE_BOOT_WAIT (500ms) — do NOT transmit 0x7B9
    │
    ▼
BRAKE_LISTEN_SYNC:
    Wait for 0x721 SEB_STATUS from SEB
    Read current stroke position (SEB_Stroke_Value)
    Set initial command = current position (hold — don't move)
    Wait for SEB_Alignment_Status == 1
    Timeout: 2s → BRAKE_FAULT
    │
    ▼
BRAKE_ACTIVE: transmit 0x7B9 at 50 Hz continuously
\end{verbatim}

If the SEB never responds → \texttt{BRAKE\_FAULT}. SYS stops transmitting
\texttt{0x7B9}. The SEB enters its own timeout-fault. The rider has no hydraulic
braking. This is a known single-point failure --- there is currently no
mechanical brake fallback (tracked as issue M2).

\begin{center}\rule{0.5\linewidth}{0.5pt}\end{center}

\section{11. Complete brake source summary}\label{complete-brake-source-summary}

\begin{longtable}[]{@{}
  >{\raggedright\arraybackslash}p{(\columnwidth - 8\tabcolsep) * \real{0.1667}}
  >{\raggedright\arraybackslash}p{(\columnwidth - 8\tabcolsep) * \real{0.2500}}
  >{\raggedright\arraybackslash}p{(\columnwidth - 8\tabcolsep) * \real{0.1667}}
  >{\raggedright\arraybackslash}p{(\columnwidth - 8\tabcolsep) * \real{0.2500}}
  >{\raggedright\arraybackslash}p{(\columnwidth - 8\tabcolsep) * \real{0.1667}}@{}}
\toprule\noalign{}
\begin{minipage}[b]{\linewidth}\raggedright
Path
\end{minipage} & \begin{minipage}[b]{\linewidth}\raggedright
Trigger
\end{minipage} & \begin{minipage}[b]{\linewidth}\raggedright
Mode
\end{minipage} & \begin{minipage}[b]{\linewidth}\raggedright
Command
\end{minipage} & \begin{minipage}[b]{\linewidth}\raggedright
When
\end{minipage} \\
\midrule\noalign{}
\endhead
\bottomrule\noalign{}
\endlastfoot
Lever & Rider squeezes GPIO2 & Stroke (1) & 15 mm & MANUAL, AUTO (always
overrides) \\
ESTOP & Button / CAN \texttt{0x001} / HB loss & Stroke (1) & 27 mm (max) & Any
mode \\
Jetson decel & Perception → \texttt{0x301} → RT → \texttt{0x205} & Pressure (2)
& 0--5 MPa & AUTO only, no lever \\
Obstacle & HC-SR04 → RT → \texttt{0x205} & Pressure (2) & 0--5 MPa & AUTO only,
no lever \\
Released & No lever, no \texttt{0x205}, no ESTOP & Stroke (1) & 0 mm & AUTO
default \\
\end{longtable}

\begin{center}\rule{0.5\linewidth}{0.5pt}\end{center}

\emph{See also: \href{../architecture.md}{\texttt{architecture.md}} §8.6 for
brake control mechanisms,
\href{../can-dictionary.md}{\texttt{can-dictionary.md}} §0x7B9 and §0x721 for
bit-level frame layouts (also §0x731 SEB\_ErrInfo, §0x741 SEB\_Version, §0x6FB
SEB\_Test), \href{brake-unit.md}{\texttt{docs/brake-unit.md}} for SYNTREE SEB
protocol reference, \href{../issues.md}{\texttt{issues.md}} M2 for the
mechanical fallback gap.}


\chapter{SYNTREE SEB --- Electro-Hydraulic Brake
Unit}\label{syntree-seb-electro-hydraulic-brake-unit}

CAN-controlled brake actuator. Factory-programmed IDs (not reconfigurable).

\begin{center}\rule{0.5\linewidth}{0.5pt}\end{center}

\section{1. CAN Interface}\label{can-interface}

\begin{longtable}[]{@{}ll@{}}
\toprule\noalign{}
Parameter & Value \\
\midrule\noalign{}
\endhead
\bottomrule\noalign{}
\endlastfoot
Bus & Low-level CAN (500 kbit/s) \\
Command ID & \texttt{0x7B9} VCU\_SEB\_REQ (customized from factory default
\texttt{0x720}) \\
Status ID & \texttt{0x721} SEB\_STATUS \\
Command rate & 20 ms (50 Hz) --- \textbf{continuous transmission required} \\
Status rate & 10 ms (100 Hz) \\
Endianness & Motorola LSB (little-endian) \\
DLC (both) & 8 \\
\end{longtable}

\begin{center}\rule{0.5\linewidth}{0.5pt}\end{center}

\section{2. Command Frame --- 0x7B9
VCU\_SEB\_REQ}\label{command-frame-0x7b9-vcu_seb_req}

\begin{longtable}[]{@{}
  >{\raggedright\arraybackslash}p{(\columnwidth - 18\tabcolsep) * \real{0.1081}}
  >{\raggedright\arraybackslash}p{(\columnwidth - 18\tabcolsep) * \real{0.1486}}
  >{\raggedright\arraybackslash}p{(\columnwidth - 18\tabcolsep) * \real{0.0676}}
  >{\raggedright\arraybackslash}p{(\columnwidth - 18\tabcolsep) * \real{0.0811}}
  >{\raggedright\arraybackslash}p{(\columnwidth - 18\tabcolsep) * \real{0.0946}}
  >{\raggedright\arraybackslash}p{(\columnwidth - 18\tabcolsep) * \real{0.1081}}
  >{\raggedright\arraybackslash}p{(\columnwidth - 18\tabcolsep) * \real{0.0676}}
  >{\raggedright\arraybackslash}p{(\columnwidth - 18\tabcolsep) * \real{0.0676}}
  >{\raggedright\arraybackslash}p{(\columnwidth - 18\tabcolsep) * \real{0.0811}}
  >{\raggedright\arraybackslash}p{(\columnwidth - 18\tabcolsep) * \real{0.1757}}@{}}
\toprule\noalign{}
\begin{minipage}[b]{\linewidth}\raggedright
Signal
\end{minipage} & \begin{minipage}[b]{\linewidth}\raggedright
Start bit
\end{minipage} & \begin{minipage}[b]{\linewidth}\raggedright
Len
\end{minipage} & \begin{minipage}[b]{\linewidth}\raggedright
Type
\end{minipage} & \begin{minipage}[b]{\linewidth}\raggedright
Scale
\end{minipage} & \begin{minipage}[b]{\linewidth}\raggedright
Offset
\end{minipage} & \begin{minipage}[b]{\linewidth}\raggedright
Min
\end{minipage} & \begin{minipage}[b]{\linewidth}\raggedright
Max
\end{minipage} & \begin{minipage}[b]{\linewidth}\raggedright
Unit
\end{minipage} & \begin{minipage}[b]{\linewidth}\raggedright
Description
\end{minipage} \\
\midrule\noalign{}
\endhead
\bottomrule\noalign{}
\endlastfoot
\texttt{VCU\_SEB\_Alignment\_Enable} & 0 & 1 & bool & 1 & 0 & 0 & 1 & --- &
Calibration enable \\
\texttt{VCU\_SEB\_Control\_Enable} & 1 & 1 & bool & 1 & 0 & 0 & 1 & --- & Active
control enable \\
\texttt{VCU\_SEB\_Control\_Mode} & 2 & 1 & bool & 1 & 0 & 0 & 1 & --- &
0=Stroke, 1=Pressure \\
\texttt{VCU\_SEB\_AutoBrake} & 3 & 1 & bool & 1 & 0 & 0 & 1 & --- & Auto-brake
trigger \\
(reserved) & 4 & 4 & --- & --- & --- & --- & --- & --- & \\
(reserved) & --- & 8 & --- & --- & --- & --- & --- & --- & Byte 1 \\
\texttt{VCU\_SEB\_Stroke\_Value\_Req} & 16 & 16 & u16 & 0.05 & -30 & -5 & 27 &
mm & Requested stroke position \\
\texttt{VCU\_SEB\_Pre\_Value\_Req} & 32 & 8 & u8 & 0.05 & 0 & 0 & 5 & MPa &
Requested pressure. Raw = kPa × 0.02. ⚠️ Was 16-bit; CSV and can-dictionary
confirm 8-bit. \\
(reserved) & 40 & 8 & --- & --- & --- & --- & --- & --- & Byte 5 \\
\texttt{VCU\_SEB\_RollCnt\_Enable} & 48 & 1 & bool & 1 & 0 & 0 & 1 & --- &
\textbf{Must be 1} \\
\texttt{VCU\_SEB\_CheckSum\_Enable} & 49 & 1 & bool & 1 & 0 & 0 & 1 & --- &
\textbf{Must be 1} \\
(reserved) & 50 & 2 & --- & --- & --- & --- & --- & --- & \\
\texttt{VCU\_SEB\_RollCnt} & 52 & 4 & u8 & 1 & 0 & 0 & 15 & --- & Rolling
counter. Increment every frame. \\
\texttt{VCU\_SEB\_CheckSum} & 56 & 8 & u8 & 1 & 0 & 0 & 255 & --- & XOR of bytes
0--6, then \texttt{\^{}\ 0xFF} \\
\end{longtable}

\subsection{Byte layout}\label{byte-layout}

\begin{longtable}[]{@{}
  >{\raggedright\arraybackslash}p{(\columnwidth - 16\tabcolsep) * \real{0.2000}}
  >{\raggedright\arraybackslash}p{(\columnwidth - 16\tabcolsep) * \real{0.1000}}
  >{\raggedright\arraybackslash}p{(\columnwidth - 16\tabcolsep) * \real{0.1000}}
  >{\raggedright\arraybackslash}p{(\columnwidth - 16\tabcolsep) * \real{0.1000}}
  >{\raggedright\arraybackslash}p{(\columnwidth - 16\tabcolsep) * \real{0.1000}}
  >{\raggedright\arraybackslash}p{(\columnwidth - 16\tabcolsep) * \real{0.1000}}
  >{\raggedright\arraybackslash}p{(\columnwidth - 16\tabcolsep) * \real{0.1000}}
  >{\raggedright\arraybackslash}p{(\columnwidth - 16\tabcolsep) * \real{0.1000}}
  >{\raggedright\arraybackslash}p{(\columnwidth - 16\tabcolsep) * \real{0.1000}}@{}}
\toprule\noalign{}
\begin{minipage}[b]{\linewidth}\raggedright
Byte
\end{minipage} & \begin{minipage}[b]{\linewidth}\raggedright
0
\end{minipage} & \begin{minipage}[b]{\linewidth}\raggedright
1
\end{minipage} & \begin{minipage}[b]{\linewidth}\raggedright
2
\end{minipage} & \begin{minipage}[b]{\linewidth}\raggedright
3
\end{minipage} & \begin{minipage}[b]{\linewidth}\raggedright
4
\end{minipage} & \begin{minipage}[b]{\linewidth}\raggedright
5
\end{minipage} & \begin{minipage}[b]{\linewidth}\raggedright
6
\end{minipage} & \begin{minipage}[b]{\linewidth}\raggedright
7
\end{minipage} \\
\midrule\noalign{}
\endhead
\bottomrule\noalign{}
\endlastfoot
Content & Align{[}0{]}+CtrlEn{[}1{]}+Mode{[}2{]}+AutoBrk{[}3{]}+(gap){[}4:7{]} &
rsvd & stroke {[}7:0{]} & stroke {[}15:8{]} & press {[}7:0{]} & rsvd &
rsvd(2)+RollCntEn(1)+CksEn(1)+RollCnt(4) & checksum \\
\end{longtable}

\subsection{Stroke conversion}\label{stroke-conversion}

\begin{verbatim}
raw_value = (physical_mm + 30.0) / 0.05
physical_mm = raw_value × 0.05 − 30.0
\end{verbatim}

\begin{longtable}[]{@{}lll@{}}
\toprule\noalign{}
Physical & Raw & Use case \\
\midrule\noalign{}
\endhead
\bottomrule\noalign{}
\endlastfoot
−5 mm & 500 & Minimum \\
0 mm & 600 & Released (no brake) \\
15 mm & 900 & Manual lever pressed \\
27 mm & 1140 & ESTOP --- full brake \\
\end{longtable}

\begin{center}\rule{0.5\linewidth}{0.5pt}\end{center}

\section{3. Status Frame --- 0x721
SEB\_STATUS}\label{status-frame-0x721-seb_status}

\begin{longtable}[]{@{}
  >{\raggedright\arraybackslash}p{(\columnwidth - 18\tabcolsep) * \real{0.1081}}
  >{\raggedright\arraybackslash}p{(\columnwidth - 18\tabcolsep) * \real{0.1486}}
  >{\raggedright\arraybackslash}p{(\columnwidth - 18\tabcolsep) * \real{0.0676}}
  >{\raggedright\arraybackslash}p{(\columnwidth - 18\tabcolsep) * \real{0.0811}}
  >{\raggedright\arraybackslash}p{(\columnwidth - 18\tabcolsep) * \real{0.0946}}
  >{\raggedright\arraybackslash}p{(\columnwidth - 18\tabcolsep) * \real{0.1081}}
  >{\raggedright\arraybackslash}p{(\columnwidth - 18\tabcolsep) * \real{0.0676}}
  >{\raggedright\arraybackslash}p{(\columnwidth - 18\tabcolsep) * \real{0.0676}}
  >{\raggedright\arraybackslash}p{(\columnwidth - 18\tabcolsep) * \real{0.0811}}
  >{\raggedright\arraybackslash}p{(\columnwidth - 18\tabcolsep) * \real{0.1757}}@{}}
\toprule\noalign{}
\begin{minipage}[b]{\linewidth}\raggedright
Signal
\end{minipage} & \begin{minipage}[b]{\linewidth}\raggedright
Start bit
\end{minipage} & \begin{minipage}[b]{\linewidth}\raggedright
Len
\end{minipage} & \begin{minipage}[b]{\linewidth}\raggedright
Type
\end{minipage} & \begin{minipage}[b]{\linewidth}\raggedright
Scale
\end{minipage} & \begin{minipage}[b]{\linewidth}\raggedright
Offset
\end{minipage} & \begin{minipage}[b]{\linewidth}\raggedright
Min
\end{minipage} & \begin{minipage}[b]{\linewidth}\raggedright
Max
\end{minipage} & \begin{minipage}[b]{\linewidth}\raggedright
Unit
\end{minipage} & \begin{minipage}[b]{\linewidth}\raggedright
Description
\end{minipage} \\
\midrule\noalign{}
\endhead
\bottomrule\noalign{}
\endlastfoot
\texttt{SEB\_Alignment\_Status} & 0 & 1 & bool & 1 & 0 & 0 & 1 & --- & Alignment
Info Feedback. 1 = aligned. \\
\texttt{SEB\_Control\_Enable\_Status} & 1 & 1 & bool & 1 & 0 & 0 & 1 & --- &
Control Enable Feedback \\
\texttt{SEB\_Control\_Mode\_Status} & 2 & 2 & u8 & 1 & 0 & 0 & 3 & enum &
Control Mode Feedback \\
\texttt{SEB\_AutoBrake\_Status} & 4 & 1 & bool & 1 & 0 & 0 & 1 & --- & Auto
Brake Status Feedback \\
(unaccounted) & 5 & 1 & --- & --- & --- & --- & --- & --- & \\
\texttt{SEB\_Error\_Status} & 6 & 2 & u8 & 1 & 0 & 0 & 3 & enum & 0=No fault,
1=L1 minor, 2=L2 general, 3=L3 severe \\
(unaccounted) & --- & 8 & --- & --- & --- & --- & --- & --- & Byte 1 \\
\texttt{SEB\_Stroke\_Value} & 16 & 16 & u16 & 0.05 & -30 & -5 & 27 & mm & Stroke
Value Feedback \\
\texttt{SEB\_Pressure\_Value} & 24 & 8 & u8 & 0.05 & 0 & 0 & 5 & MPa & Pressure
Value Feedback. Overlaps Stroke{[}15:8{]} --- mode-dependent. \\
\texttt{SEB\_Angle\_Value} & 40 & 16 & i16 & 0.5 & 0 & -150 & 840 & --- & Angle
Feedback. Overlaps security echo bits at byte 6. \\
\texttt{SEB\_RollCnt\_Enable\_Status} & 48 & 1 & bool & 1 & 0 & 0 & 1 & --- &
Life Signal Status Feedback. Overlaps Angle{[}15:8{]}. \\
\texttt{SEB\_CheckSum\_Enable\_Status} & 49 & 1 & bool & 1 & 0 & 0 & 1 & --- &
Checksum Status Feedback. Overlaps Angle{[}15:8{]}. \\
(unaccounted) & 50 & 2 & --- & --- & --- & --- & --- & --- & \\
\texttt{SEB\_RollCnt\_Status} & 52 & 4 & u8 & 1 & 0 & 0 & 15 & --- & Life Signal
Feedback --- echoes rolling counter \\
\texttt{SEB\_CheckSum\_Status} & 56 & 8 & u8 & 1 & 0 & 0 & 255 & --- & Checksum
Feedback \\
\end{longtable}

\subsection{Byte layout}\label{byte-layout-1}

\begin{longtable}[]{@{}
  >{\raggedright\arraybackslash}p{(\columnwidth - 16\tabcolsep) * \real{0.2000}}
  >{\raggedright\arraybackslash}p{(\columnwidth - 16\tabcolsep) * \real{0.1000}}
  >{\raggedright\arraybackslash}p{(\columnwidth - 16\tabcolsep) * \real{0.1000}}
  >{\raggedright\arraybackslash}p{(\columnwidth - 16\tabcolsep) * \real{0.1000}}
  >{\raggedright\arraybackslash}p{(\columnwidth - 16\tabcolsep) * \real{0.1000}}
  >{\raggedright\arraybackslash}p{(\columnwidth - 16\tabcolsep) * \real{0.1000}}
  >{\raggedright\arraybackslash}p{(\columnwidth - 16\tabcolsep) * \real{0.1000}}
  >{\raggedright\arraybackslash}p{(\columnwidth - 16\tabcolsep) * \real{0.1000}}
  >{\raggedright\arraybackslash}p{(\columnwidth - 16\tabcolsep) * \real{0.1000}}@{}}
\toprule\noalign{}
\begin{minipage}[b]{\linewidth}\raggedright
Byte
\end{minipage} & \begin{minipage}[b]{\linewidth}\raggedright
0
\end{minipage} & \begin{minipage}[b]{\linewidth}\raggedright
1
\end{minipage} & \begin{minipage}[b]{\linewidth}\raggedright
2
\end{minipage} & \begin{minipage}[b]{\linewidth}\raggedright
3
\end{minipage} & \begin{minipage}[b]{\linewidth}\raggedright
4
\end{minipage} & \begin{minipage}[b]{\linewidth}\raggedright
5
\end{minipage} & \begin{minipage}[b]{\linewidth}\raggedright
6
\end{minipage} & \begin{minipage}[b]{\linewidth}\raggedright
7
\end{minipage} \\
\midrule\noalign{}
\endhead
\bottomrule\noalign{}
\endlastfoot
Content &
Align{[}0{]}+CtrlEn{[}1{]}+Mode{[}2:3{]}+AutoBrk{[}4{]}+(gap)+Error{[}6:7{]} &
(unacc.) & Stroke {[}7:0{]} & Stroke {[}15:8{]} / Pressure {[}7:0{]}
(mode-muxed) & (unacc.) & Angle {[}7:0{]} & Angle{[}15:8{]} /
RollCntEn{[}0{]}+CksEn{[}1{]}+(gap)+RollCnt{[}4:7{]} (overlap) & CksSum\_Stat \\
\end{longtable}

\begin{center}\rule{0.5\linewidth}{0.5pt}\end{center}

\section{4. Control Modes}\label{control-modes}

\subsection{Mode 0 --- Stroke Control
(Position-Based)}\label{mode-0-stroke-control-position-based}

Command the actuator's internal pushrod to move a specific distance in
millimeters. Best for mimicking human pedal travel or simple proportional
control.

\begin{itemize}
\tightlist
\item
  0 mm = no brake
\item
  15 mm = moderate braking
\item
  27 mm = full lock
\end{itemize}

\textbf{Current usage in E-Trike:} SYS uses Stroke Mode exclusively. The only
brake triggers are the physical lever (binary) and ESTOP (binary) --- both map
to fixed stroke positions. No closed-loop pressure control needed.

\subsection{Mode 1 --- Pressure Control
(Force-Based)}\label{mode-1-pressure-control-force-based}

Command a specific hydraulic pressure in the brake fluid lines (MPa). The unit's
internal PID moves the motor to maintain that exact pressure.

\textbf{Planned for AUTO mode} once the brake arbitration gap is closed.
Benefits: - Compensates for brake pad wear - Compensates for fluid temperature
expansion - Safer for autonomous deceleration --- maintains commanded
deceleration regardless of mechanical variance

\begin{center}\rule{0.5\linewidth}{0.5pt}\end{center}

\section{5. Security Bytes}\label{security-bytes}

\subsection{Rolling Counter}\label{rolling-counter}

4-bit value (0--15) in Byte 6, bits 4--7. \textbf{Must increment every frame.}
If the SEB receives two consecutive frames with the same counter value, it
assumes the controller is frozen and rejects the command.

\begin{verbatim}
roll_cnt = (roll_cnt + 1) & 0x0F;  // increment, wrap at 15
\end{verbatim}

\subsection{Checksum}\label{checksum}

8-bit XOR of bytes 0--6, then XOR with \texttt{0xFF}. Placed in Byte 7.

\begin{verbatim}
uint8_t checksum = 0;
for (int i = 0; i < 7; i++) checksum ^= raw_bytes[i];
checksum ^= 0xFF;
\end{verbatim}

\begin{quote}
\textbf{Verify against SYNTREE spec} --- some units use a different constant or
a shift operation.
\end{quote}

\begin{center}\rule{0.5\linewidth}{0.5pt}\end{center}

\section{6. Boot Sequence --- ``Listen Before
Speaking''}\label{boot-sequence-listen-before-speaking}

The SEB requires a strict startup handshake. Violating this causes sensor
detection errors.

\begin{verbatim}
State machine:

BRAKE_BOOT_WAIT:
  - 500 ms delay after power-on
  - DO NOT transmit any 0x7B9 frames
  - → BRAKE_LISTEN_SYNC

BRAKE_LISTEN_SYNC:
  - Wait for 0x721 SEB_STATUS frame
  - Read SEB_Stroke_Value (current physical position)
  - Set initial command target = current stroke (hold position — do not move)
  - Wait for SEB_Alignment_Status == 1
  - → BRAKE_ACTIVE

BRAKE_ACTIVE:
  - Transmit 0x7B9 at 50 Hz continuously
  - First frame commands exactly the current position (no movement)
  - Rolling counter + checksum on every frame
  - → BRAKE_FAULT on error

BRAKE_FAULT:
  - Stop transmitting
  - Log error
\end{verbatim}

\textbf{Critical rules:} 1. Never apply power while VCU is already sending
commands --- causes angle sensor detection error. 2. Never command a position
until the current position is known. 3. First frame after sync must command the
current position (no movement).

\begin{center}\rule{0.5\linewidth}{0.5pt}\end{center}

\section{7. C++ Implementation Guide}\label{c-implementation-guide}

\begin{verbatim}
#include <stdint.h>
#include <string.h>

#define CAN_ID_VCU_SEB_REQ 0x7B9  // customized from factory default 0x720

typedef struct {
    uint8_t align_enable   : 1;
    uint8_t control_enable : 1;
    uint8_t control_mode   : 1;   // 0=Stroke, 1=Pressure
    uint8_t auto_brake     : 1;
    uint8_t reserved_0     : 4;

    uint8_t reserved_1;

    // Physical = (raw × 0.05) − 30
    // Example: 0 mm → raw = (0 + 30) / 0.05 = 600
    uint16_t stroke_req;

    uint8_t pressure_req;     // u8: 0.05 MPa/bit, 0–5 MPa. Raw = kPa × 0.02

    uint8_t reserved_2;        // byte 5

    uint8_t roll_cnt_enable  : 1;  // MUST be 1
    uint8_t checksum_enable  : 1;  // MUST be 1
    uint8_t reserved_3       : 2;
    uint8_t rolling_counter  : 4;  // Increment 0–15 every frame

    uint8_t checksum;              // XOR(bytes[0..6]) ^ 0xFF
} __attribute__((packed)) VCU_SEB_Req_t;


void seb_send_command(float target_stroke_mm) {
    static uint8_t roll_cnt = 0;
    VCU_SEB_Req_t payload;
    memset(&payload, 0, sizeof(payload));

    // 1. Enable + Mode
    payload.align_enable  = 1;
    payload.control_enable = 1;
    payload.control_mode   = 0;   // Stroke Mode

    // 2. Stroke target (apply scale + offset)
    payload.stroke_req = (uint16_t)((target_stroke_mm + 30.0f) / 0.05f);

    // 3. Rolling counter
    payload.rolling_counter = roll_cnt;
    roll_cnt = (roll_cnt + 1) & 0x0F;

    // 4. Checksum
    uint8_t* raw = (uint8_t*)&payload;
    uint8_t cksum = 0;
    for (int i = 0; i < 7; i++) cksum ^= raw[i];
    payload.checksum = cksum ^ 0xFF;

    // 5. Transmit
    // CAN_Write(CAN_ID_VCU_SEB_REQ, 8, raw);
}
\end{verbatim}

\begin{center}\rule{0.5\linewidth}{0.5pt}\end{center}

\section{8. Error Handling}\label{error-handling}

\begin{longtable}[]{@{}
  >{\raggedright\arraybackslash}p{(\columnwidth - 4\tabcolsep) * \real{0.2558}}
  >{\raggedright\arraybackslash}p{(\columnwidth - 4\tabcolsep) * \real{0.3023}}
  >{\raggedright\arraybackslash}p{(\columnwidth - 4\tabcolsep) * \real{0.4419}}@{}}
\toprule\noalign{}
\begin{minipage}[b]{\linewidth}\raggedright
Condition
\end{minipage} & \begin{minipage}[b]{\linewidth}\raggedright
SEB Response
\end{minipage} & \begin{minipage}[b]{\linewidth}\raggedright
Controller Action
\end{minipage} \\
\midrule\noalign{}
\endhead
\bottomrule\noalign{}
\endlastfoot
Rolling counter repeats & Reject frame, trigger fault & Reset counter,
re-sync \\
Checksum mismatch & Discard frame & Re-transmit with correct checksum \\
Comm timeout (\textgreater20 ms no frame) & Internal fault & Restart boot
sequence \\
\texttt{SEB\_Error\_Status\ \textgreater{}\ 0} & Unit enters degraded mode &
Log, report via \texttt{0x011} \\
Sync timeout (no \texttt{0x721} for 2s) & --- & Remain in BRAKE\_FAULT \\
\end{longtable}

\begin{center}\rule{0.5\linewidth}{0.5pt}\end{center}

\section{9. Integration Notes}\label{integration-notes}

\begin{itemize}
\tightlist
\item
  \textbf{Controlled by:} SYS ESP32-S3 (low-level CAN)
\item
  \textbf{Current mode:} Stroke (1) --- lever press → 15 mm, ESTOP → 27 mm,
  released → 0 mm
\item
  \textbf{Planned mode:} Pressure (Mode 1) --- when RT brake arbitration path to
  SYS is implemented (§12.1 in architecture.md)
\item
  \textbf{No daily homing required.} One-time zero-calibration when first mated
  to mechanical assembly. Calibration retained across power cycles.
\item
  \textbf{Not reconfigurable.} CAN IDs are factory-programmed.
  \texttt{RT\_DRIVE\_CMD} placed at \texttt{0x204} (\texttt{0x202} is now
  SES\_ErrInfo).
\end{itemize}


\chapter{CAN Gateway Bridging}\label{can-gateway-bridging}

In a multi-bus CAN architecture, the \textbf{gateway} is the single node
connected to both buses. It selectively forwards messages between them, acting
as a controlled bridge.

On the E-Trike, the \textbf{RT ESP32-S3} is the gateway between the High-Level
CAN (Jetson + RT) and Low-Level CAN (RT + SYS + SYNTREE actuators).

\begin{center}\rule{0.5\linewidth}{0.5pt}\end{center}

\section{Why two buses?}\label{why-two-buses}

Separating CAN buses serves three purposes:

\begin{enumerate}
\def\labelenumi{\arabic{enumi}.}
\item
  \textbf{Security isolation.} Jetson runs a full Linux+ROS 2 stack --- complex,
  harder to harden. The actuators (steering, brake) are on a physically separate
  bus. Even if Jetson is compromised, it cannot directly command steering or
  brakes. It must go through RT, which validates and clamps every command.
\item
  \textbf{Bus load partitioning.} High-level bus carries variable-rate ROS
  traffic (up to 100 Hz drive commands, telemetry). Low-level bus carries fixed
  50 Hz actuator commands and safety status. Segregation prevents actuator
  commands from being delayed by bursty telemetry.
\item
  \textbf{Fault containment.} A bus-off on one bus doesn't take down the other.
  RT can continue operating actuators even if the Jetson link fails.
\end{enumerate}

\begin{center}\rule{0.5\linewidth}{0.5pt}\end{center}

\section{Gateway topology}\label{gateway-topology}

\begin{verbatim}
 High-Level CAN (500 kbit/s)          Low-Level CAN (500 kbit/s)
 ┌─────────────────────────┐          ┌─────────────────────────────┐
 │  Jetson                 │          │  SYS ESP32-S3               │
 │  (ROS 2, planning)      │          │  (safety, motor, brake)     │
 │                         │          │                             │
 │  TX: 0x300,0x301,0x302  │          │  TX: 0x011,0x012,0x110,    │
 │      0x001,0x7FC         │          │      0x120,0x600,0x7B9,    │
 │  RX: 0x011,0x120,0x210  │          │      0x001,0x7FE            │
 │      0x220,0x400,0x600  │          │  RX: 0x001,0x204,0x302,    │
 │      0x001,0x7FD         │          │      0x721,0x7FD            │
 └───────────┬─────────────┘          └──────────────┬──────────────┘
             │                                       │
             │          ┌──────────────┐             │
             └──────────┤ RT ESP32-S3  ├─────────────┘
                        │  (gateway)   │
                        │              │
                        │ TWAI +       │
                        │ MCP2515 SPI  │
                        └──────────────┘
\end{verbatim}

RT is the \textbf{only} node on both buses. No direct Jetson ↔ SYS path exists.

\begin{center}\rule{0.5\linewidth}{0.5pt}\end{center}

\section{Forwarding rules}\label{forwarding-rules}

The gateway applies explicit, static forwarding rules. Not all messages cross
the bridge.

\subsection{Low → High (upstream telemetry)}\label{low-high-upstream-telemetry}

\begin{longtable}[]{@{}lll@{}}
\toprule\noalign{}
ID & Name & Action \\
\midrule\noalign{}
\endhead
\bottomrule\noalign{}
\endlastfoot
\texttt{0x001} & SAFETY\_ESTOP & \textbf{Forward transparently} (same ID, same
payload) \\
\texttt{0x011} & SYS\_SAFETY\_STATUS & \textbf{Forward transparently} \\
\texttt{0x120} & SYS\_THROTTLE\_POS & \textbf{Forward transparently} \\
\texttt{0x600} & SYS\_DIAG & \textbf{Forward transparently} \\
\end{longtable}

\subsection{High → Low (downstream
commands)}\label{high-low-downstream-commands}

\begin{longtable}[]{@{}
  >{\raggedright\arraybackslash}p{(\columnwidth - 4\tabcolsep) * \real{0.2222}}
  >{\raggedright\arraybackslash}p{(\columnwidth - 4\tabcolsep) * \real{0.3333}}
  >{\raggedright\arraybackslash}p{(\columnwidth - 4\tabcolsep) * \real{0.4444}}@{}}
\toprule\noalign{}
\begin{minipage}[b]{\linewidth}\raggedright
ID
\end{minipage} & \begin{minipage}[b]{\linewidth}\raggedright
Name
\end{minipage} & \begin{minipage}[b]{\linewidth}\raggedright
Action
\end{minipage} \\
\midrule\noalign{}
\endhead
\bottomrule\noalign{}
\endlastfoot
\texttt{0x001} & SAFETY\_ESTOP & \textbf{Forward transparently} (ESTOP from
Jetson reaches SYS + actuators) \\
\texttt{0x302} & HOST\_LIGHT\_CMD & \textbf{Forward transparently} \\
\end{longtable}

\subsection{Consumed locally (not
forwarded)}\label{consumed-locally-not-forwarded}

\begin{longtable}[]{@{}
  >{\raggedright\arraybackslash}p{(\columnwidth - 4\tabcolsep) * \real{0.2667}}
  >{\raggedright\arraybackslash}p{(\columnwidth - 4\tabcolsep) * \real{0.4000}}
  >{\raggedright\arraybackslash}p{(\columnwidth - 4\tabcolsep) * \real{0.3333}}@{}}
\toprule\noalign{}
\begin{minipage}[b]{\linewidth}\raggedright
ID
\end{minipage} & \begin{minipage}[b]{\linewidth}\raggedright
Name
\end{minipage} & \begin{minipage}[b]{\linewidth}\raggedright
Why
\end{minipage} \\
\midrule\noalign{}
\endhead
\bottomrule\noalign{}
\endlastfoot
\texttt{0x300} & HOST\_DRIVE\_CMD & Consumed by RT --- kinematics resolve it to
\texttt{0x204} + \texttt{0x169} \\
\texttt{0x301} & HOST\_BRAKE\_REQUEST & Consumed by RT --- brake arbitration
(max-select) \\
\end{longtable}

\subsection{Generated locally (not forwarded, originated by
gateway)}\label{generated-locally-not-forwarded-originated-by-gateway}

\begin{longtable}[]{@{}llll@{}}
\toprule\noalign{}
ID & Name & Bus & Why \\
\midrule\noalign{}
\endhead
\bottomrule\noalign{}
\endlastfoot
\texttt{0x204} & RT\_DRIVE\_CMD & Low & RT output --- speed+gear → MTR \\
\texttt{0x169} & VCU\_SES\_REQ & Low & RT output --- steering angle → EPS-C \\
\texttt{0x210} & RT\_STATE\_REPORT & High & RT telemetry → Jetson \\
\texttt{0x220} & RT\_PID\_FEEDBACK & High & RT telemetry → Jetson \\
\texttt{0x400} & RT\_OBSTACLE\_DIST & High & RT telemetry → Jetson \\
\end{longtable}

\subsection{Bus-local only (never cross the
bridge)}\label{bus-local-only-never-cross-the-bridge}

\begin{longtable}[]{@{}
  >{\raggedright\arraybackslash}p{(\columnwidth - 6\tabcolsep) * \real{0.2000}}
  >{\raggedright\arraybackslash}p{(\columnwidth - 6\tabcolsep) * \real{0.3000}}
  >{\raggedright\arraybackslash}p{(\columnwidth - 6\tabcolsep) * \real{0.2500}}
  >{\raggedright\arraybackslash}p{(\columnwidth - 6\tabcolsep) * \real{0.2500}}@{}}
\toprule\noalign{}
\begin{minipage}[b]{\linewidth}\raggedright
ID
\end{minipage} & \begin{minipage}[b]{\linewidth}\raggedright
Name
\end{minipage} & \begin{minipage}[b]{\linewidth}\raggedright
Bus
\end{minipage} & \begin{minipage}[b]{\linewidth}\raggedright
Why
\end{minipage} \\
\midrule\noalign{}
\endhead
\bottomrule\noalign{}
\endlastfoot
\texttt{0x012} & SYS\_DCDC\_CMD & Low & DC-DC only on low bus \\
\texttt{0x110} & SYS\_MODE\_CMD & Low & Mode switch local to low bus \\
\texttt{0x201} & SES\_STATUS & Low & EPS-C feedback --- consumed by RT on low \\
\texttt{0x7B9} & VCU\_SEB\_REQ & Low & Brake command from SYS to SEB \\
\texttt{0x721} & SEB\_STATUS & Low & Brake feedback to SYS \\
\texttt{0x7FC}, \texttt{0x7FD}, \texttt{0x7FE} & HEARTBEAT (Jetson, RT, SYS) &
Both & Independent per bus (each bus has its own heartbeat domain) \\
\end{longtable}

\begin{center}\rule{0.5\linewidth}{0.5pt}\end{center}

\section{Transparent forwarding}\label{transparent-forwarding}

``Transparent'' means the gateway copies the frame verbatim: same CAN ID, same
DLC, same payload bytes. The receiver cannot tell whether the frame originated
from the original sender or passed through the gateway.

This is only safe because of the \textbf{one sender per ID} rule (except
heartbeat). Since each CAN ID has exactly one originator on each bus, there's no
ambiguity about which node sent it.

\begin{center}\rule{0.5\linewidth}{0.5pt}\end{center}

\section{Gateway implementation}\label{gateway-implementation}

On the RT ESP32-S3, the gateway runs in the \texttt{dispatch} task (priority 4):

\begin{verbatim}
dispatch_task:
  Block on can_rx_low_queue OR can_rx_high_queue (whichever fires first)

  Low bus frame received:
    0x011 → enqueue to gw_tx_high_queue    (forward SYS safety → Jetson)
    0x120 → enqueue to gw_tx_high_queue    (forward throttle pos → Jetson)
    0x600 → enqueue to gw_tx_high_queue    (forward diagnostics → Jetson)
    0x001 → enqueue to gw_tx_high_queue    (forward ESTOP → Jetson)
           + mode_set(Estop) locally
    0x201 → steer_feedback (sync angle, check following error)

  High bus frame received:
    0x300 → enqueue to cmd_queue            (drive command → control loop)
    0x301 → atomic store brake_request      (brake → arbitration)
    0x302 → enqueue to gw_tx_low_queue     (forward lights → SYS)
    0x001 → enqueue to gw_tx_low_queue     (forward ESTOP → SYS)
           + mode_set(Estop) locally
\end{verbatim}

The \texttt{can\_tx\_low} and \texttt{can\_tx\_high} tasks drain the gateway TX
queues and send the forwarded frames.

\begin{center}\rule{0.5\linewidth}{0.5pt}\end{center}

\section{Gateway queue full --- drop
policy}\label{gateway-queue-full-drop-policy}

If a gateway TX queue is full, frames are dropped. The exception is
\texttt{0x001} (SAFETY\_ESTOP), which is sent directly (not queued) to guarantee
delivery.

This is acceptable because: - \texttt{0x011}, \texttt{0x120}, \texttt{0x600} are
periodic --- a dropped frame is replaced by the next one. - \texttt{0x302} is
change-triggered --- if dropped, the next change retransmits. - \texttt{0x001}
is the only frame where one missed message could delay ESTOP, so it bypasses the
queue.

\begin{center}\rule{0.5\linewidth}{0.5pt}\end{center}

\section{Key principles}\label{key-principles}

\begin{enumerate}
\def\labelenumi{\arabic{enumi}.}
\tightlist
\item
  \textbf{RT is the only dual-bus node.} No alternative path exists. All
  cross-bus traffic goes through RT.
\item
  \textbf{Transparent forwarding} keeps the same CAN ID --- the receiver doesn't
  know or care about the gateway.
\item
  \textbf{Explicit allowlist, not a promiscuous bridge.} Only specific IDs
  cross. Everything else stays local to its bus.
\item
  \textbf{ESTOP bypasses queue.} The one frame that must never be dropped takes
  the fast path.
\item
  \textbf{Gateway is stateless for forwarding.} No buffering, no protocol
  translation, no sequence numbers --- just copy and send.
\end{enumerate}

\begin{center}\rule{0.5\linewidth}{0.5pt}\end{center}

\emph{See also: {[}{[}can-addressing-for-etrike{]}{]} for the CAN ID scheme,
{[}{[}distributed-architecture{]}{]} for three-node rationale,
{[}{[}architecture{]}{]} §2.3 for forwarding rules, §7.7 for gateway task
layout.}


\chapter{Defense-in-Depth Safety
Patterns}\label{defense-in-depth-safety-patterns}

The E-Trike safety architecture doesn't rely on any single mechanism. Multiple
independent layers check for faults --- if one misses, the next catches it. This
is \textbf{defense in depth}: overlapping, independent safety barriers so that
no single-point failure causes a hazardous event.

\begin{center}\rule{0.5\linewidth}{0.5pt}\end{center}

\section{Layer 1: Physical ESTOP button}\label{layer-1-physical-estop-button}

A physical NC (normally-closed) pushbutton wired to SYS ESP32-S3 GPIO1.

\begin{itemize}
\tightlist
\item
  Pressed → GPIO LOW → \texttt{safety\_task} (Core 0, priority 5) detects it
  within one scheduler tick (1 ms) → \texttt{mode\_set(Estop)}.
\item
  The button is wired fail-safe: a broken wire or disconnected plug reads as LOW
  = ESTOP.
\item
  This is the \textbf{fastest} path to ESTOP and works regardless of CAN bus
  state or firmware health (as long as the safety task runs).
\end{itemize}

\textbf{Checked in:} SYS \texttt{safety\_task} (direct GPIO read, 20 Hz poll).

\begin{center}\rule{0.5\linewidth}{0.5pt}\end{center}

\section{\texorpdfstring{Layer 2: CAN \texttt{0x001}
SAFETY\_ESTOP}{Layer 2: CAN 0x001 SAFETY\_ESTOP}}\label{layer-2-can-0x001-safety_estop}

Any node (Jetson, RT, SYS) can broadcast \texttt{0x001} on either CAN bus.

\begin{itemize}
\tightlist
\item
  An empty frame (DLC=0) with the highest-priority CAN ID (\texttt{0x001} wins
  every arbitration).
\item
  RT forwards it between buses (transparently, bypassing the gateway queue).
\item
  Each node's dispatch task processes it immediately.
\end{itemize}

\textbf{Why Jetson might trigger ESTOP:} Perception stack detects an imminent
collision that the obstacle sensor missed. \textbf{Why RT might trigger:}
Steering following error exceeds threshold. \textbf{Why SYS might trigger:}
Physical ESTOP button (redundant CAN path).

\textbf{Checked in:} Every node's CAN dispatch task.

\begin{center}\rule{0.5\linewidth}{0.5pt}\end{center}

\section{Layer 3: Heartbeat timeout}\label{layer-3-heartbeat-timeout}

Each node sends its own heartbeat (0x7FC/0x7FD/0x7FE) at 2 Hz on its bus. The
safety monitor checks:

\begin{itemize}
\tightlist
\item
  On RT: if Jetson heartbeat (\texttt{0x7FC} on high CAN) stops for
  \textgreater1500 ms → ESTOP in AUTO.
\item
  On SYS: if RT heartbeat (\texttt{0x7FD} on low CAN) stops for \textgreater1500
  ms → ESTOP in AUTO.
\end{itemize}

A missing heartbeat means the controller node has crashed, rebooted, or lost CAN
connectivity. The vehicle must stop.

\textbf{Checked in:} RT watchdog task, SYS safety task.

\begin{center}\rule{0.5\linewidth}{0.5pt}\end{center}

\section{Layer 4: Command staleness
(watchdog)}\label{layer-4-command-staleness-watchdog}

RT's watchdog task runs at 10 Hz and checks the timestamp of the last received
\texttt{0x300\ HOST\_DRIVE\_CMD}:

\begin{verbatim}
if (time_since_last_drive_cmd > 500 ms) {
    zero 0x204 RT_DRIVE_CMD       →  MTR drives motor at 0 V
    stop 0x169 VCU_SES_REQ        →  EPS-C timeout-faults
}
\end{verbatim}

If Jetson's planning stack hangs or the ROS→CAN bridge crashes, the vehicle
coasts to a stop (motor zero) and steering goes to its comm-fault safe state.
This is a \textbf{passive safety} response --- no active braking, but
acceleration stops and steering becomes inert.

\textbf{Checked in:} RT \texttt{watchdog} task (10 Hz).

\begin{center}\rule{0.5\linewidth}{0.5pt}\end{center}

\section{Layer 5: Steering following
error}\label{layer-5-steering-following-error}

While in AUTO mode, RT compares the commanded steering angle (\texttt{0x169})
against the actual angle reported by EPS-C (\texttt{0x201\ SES\_StrAngle}):

\begin{verbatim}
error = abs(cmd_angle_mdeg - actual_angle_mdeg);
if (error > 5000_mdeg && error_duration_ms > 300) {
    mode_set(Estop);
}
\end{verbatim}

This catches: - \textbf{Mechanical faults:} stuck linkage, rock jam, bent tie
rod. - \textbf{Actuator faults:} EPS-C motor failure, encoder failure, internal
control loop fault. - \textbf{CAN faults:} corrupted \texttt{0x169} frames
(caught because checksum failure → EPS-C rejects → angle doesn't move).

The 300 ms persistence requirement prevents false triggers from momentary CAN
glitches or sensor noise.

\textbf{Checked in:} RT \texttt{control} task (100 Hz), using
\texttt{SES\_StrAngle} from \texttt{0x201}.

\begin{center}\rule{0.5\linewidth}{0.5pt}\end{center}

\section{Layer 6: Dynamic angle clamp (speed-dependent steering
limit)}\label{layer-6-dynamic-angle-clamp-speed-dependent-steering-limit}

The maximum allowable steering angle is inversely proportional to vehicle speed:

\begin{longtable}[]{@{}lll@{}}
\toprule\noalign{}
Speed & Max angle & Rationale \\
\midrule\noalign{}
\endhead
\bottomrule\noalign{}
\endlastfoot
2 km/h & \textasciitilde40° & Low-speed maneuvering needs full range \\
10 km/h & \textasciitilde18° & Moderate cornering \\
25 km/h & \textasciitilde5° & High speed --- tight turns cause rollover \\
\end{longtable}

The clamp is applied in RT's \texttt{physics\_resolve()} function,
\textbf{after} the kinematic solve but \textbf{before} the steering angle is
sent to EPS-C. This means even if Jetson commands a dangerous turn (bug, bad
planning, or attack), RT clamps it to the safe envelope.

The formula:

\begin{verbatim}
max_angle = lerp(kSteerMaxAngleAtHighSpeed, kSteerMaxAngleAtLowSpeed,
                 clamp((speed - kLowSpeedThresh) / (kHighSpeedThresh - kLowSpeedThresh), 0, 1))
\end{verbatim}

This is distinct from the physics model's rollover threshold (§8 of
physics-model.md). The rollover threshold is a \emph{warning boundary} for the
planner. The dynamic clamp is an \emph{enforced hard limit} applied at the
actuator level.

\textbf{Checked in:} RT \texttt{control} task (100 Hz).

\begin{center}\rule{0.5\linewidth}{0.5pt}\end{center}

\section{Layer 7: Software hard-stops}\label{layer-7-software-hard-stops}

EPS-C mechanically accepts ±78° (its internal limit). The physical steering rack
on the tricycle has end-stops at approximately ±40°. If the firmware commands
beyond ±40°, the motor pushes against an immovable mechanical stop --- this can
strip gears, burn out the motor, or bend the linkage.

RT clamps ALL commanded angles to ±40° in firmware, regardless of source:

\begin{verbatim}
cmd_angle = clamp(cmd_angle, -40000_mdeg, +40000_mdeg);
\end{verbatim}

This is the innermost safety layer for steering --- it runs on every control
cycle, before any other clamp. Even if dynamic angle clamp misbehaves,
hard-stops prevent mechanical damage.

\textbf{Checked in:} RT \texttt{control} task (100 Hz).

\begin{center}\rule{0.5\linewidth}{0.5pt}\end{center}

\section{Layer 8: Brake light OR logic}\label{layer-8-brake-light-or-logic}

The brake light illuminates if ANY braking source is active:

\begin{verbatim}
brake_light_on = brake_lever_pressed()      // GPIO2 — physical lever
              OR (mode == Estop)            // ESTOP — full brake
              OR g_light_state.brake_light; // Jetson CAN 0x302 — predictive/hazard
\end{verbatim}

All sources are \textbf{local to SYS} --- no CAN round-trip needed for the
physical lever or ESTOP. The Jetson CAN bit is a \emph{supplemental} trigger
(predictive illumination) but can never be the \emph{only} trigger. The physical
braking state always wins.

This is a \textbf{fail-visible} pattern: if any brake reason exists, the light
is on. There's no way to be braking without the light illuminating.

\textbf{Checked in:} SYS \texttt{lights} task (20 Hz).

\begin{center}\rule{0.5\linewidth}{0.5pt}\end{center}

\section{The threshold + duration pattern}\label{the-threshold-duration-pattern}

Several safety checks use the same pattern:

\begin{longtable}[]{@{}llll@{}}
\toprule\noalign{}
Check & Threshold & Duration & Action \\
\midrule\noalign{}
\endhead
\bottomrule\noalign{}
\endlastfoot
Steering following error & \textgreater5° & \textgreater300 ms & ESTOP \\
Command staleness & no \texttt{0x300} & \textgreater500 ms & Zero outputs \\
Heartbeat timeout & no peer heartbeat & \textgreater1500 ms & ESTOP in AUTO \\
\end{longtable}

\textbf{Why duration matters:} CAN frames can be delayed by arbitration
(low-priority messages wait for high-priority ones) or by bus errors. A single
missed frame should not trigger ESTOP. The duration acts as a debounce --- it
confirms the condition is persistent, not transient.

\textbf{Why different durations:} Following error (300 ms) is the fastest --- a
stuck linkage is immediately dangerous. Command staleness (500 ms) allows for
brief Jetson hiccups. Heartbeat timeout (1500 ms) is conservative --- a node
crash takes longer to confirm than a CAN glitch.

\begin{center}\rule{0.5\linewidth}{0.5pt}\end{center}

\section{ESTOP as the convergence point}\label{estop-as-the-convergence-point}

All safety layers ultimately call \texttt{mode\_set(Estop)}. The ESTOP handler:

\begin{enumerate}
\def\labelenumi{\arabic{enumi}.}
\tightlist
\item
  Sets \texttt{g\_mode\ =\ Estop} (atomic, visible to all tasks).
\item
  Motor task: MCP4725 → 0 V, all gear relays → OFF.
\item
  Brake task: \texttt{0x7B9} stroke = max (full brake, \textasciitilde27 mm).
\item
  Steering: \texttt{0x169} stops transmitting → EPS-C timeout-faults.
\item
  DC-DC: \texttt{0x012} enable = 0 → 12V rail off.
\item
  Lights: brake ON, all others OFF.
\item
  Exit requires power-cycle (cannot leave ESTOP via mode switch).
\end{enumerate}

ESTOP is an \textbf{absorbing state} --- once entered, the only way out is a
full power-cycle. This prevents accidental re-engagement while a fault condition
persists.

\begin{center}\rule{0.5\linewidth}{0.5pt}\end{center}

\section{What defense-in-depth protects
against}\label{what-defense-in-depth-protects-against}

\begin{longtable}[]{@{}
  >{\raggedright\arraybackslash}p{(\columnwidth - 2\tabcolsep) * \real{0.4333}}
  >{\raggedright\arraybackslash}p{(\columnwidth - 2\tabcolsep) * \real{0.5667}}@{}}
\toprule\noalign{}
\begin{minipage}[b]{\linewidth}\raggedright
Failure mode
\end{minipage} & \begin{minipage}[b]{\linewidth}\raggedright
Caught by layer
\end{minipage} \\
\midrule\noalign{}
\endhead
\bottomrule\noalign{}
\endlastfoot
Jetson crashes & Layer 4 (staleness) → zero outputs \\
Jetson sends dangerous steering & Layer 6 (dynamic clamp) → capped to safe
envelope \\
Steering linkage jams & Layer 5 (following error) → ESTOP \\
EPS-C motor fails & Layer 5 (following error) → ESTOP \\
CAN bus dies & Layer 3 (heartbeat) + Layer 4 (staleness) \\
Rider hits ESTOP button & Layer 1 (GPIO) + Layer 2 (CAN redundancy) \\
Firmware hangs & External watchdog IC (see {[}{[}external-watchdog{]}{]}) \\
72V shorts to chassis & {[}{[}high-voltage-isolation{]}{]} \\
\end{longtable}

\begin{center}\rule{0.5\linewidth}{0.5pt}\end{center}

\emph{Primary reference: {[}{[}emergency-system{]}{]} for the complete ESTOP
system, trigger paths, emergency response matrix, rider's guide, and testing
procedures.}

\emph{See also: {[}{[}listen-before-speaking{]}{]} for actuator boot safety,
{[}{[}external-watchdog{]}{]} for hardware watchdog,
{[}{[}high-voltage-isolation{]}{]} for galvanic isolation,
{[}{[}physics-model{]}{]} §8 for rollover threshold, {[}{[}architecture{]}{]}
§7.6 for safety mechanisms, §3 for mode state machine.}


\chapter{Three-Node Distributed
Architecture}\label{three-node-distributed-architecture}

The E-Trike uses three physically separate compute nodes, each with a distinct
role:

\begin{longtable}[]{@{}
  >{\raggedright\arraybackslash}p{(\columnwidth - 6\tabcolsep) * \real{0.1622}}
  >{\raggedright\arraybackslash}p{(\columnwidth - 6\tabcolsep) * \real{0.2703}}
  >{\raggedright\arraybackslash}p{(\columnwidth - 6\tabcolsep) * \real{0.4054}}
  >{\raggedright\arraybackslash}p{(\columnwidth - 6\tabcolsep) * \real{0.1622}}@{}}
\toprule\noalign{}
\begin{minipage}[b]{\linewidth}\raggedright
Node
\end{minipage} & \begin{minipage}[b]{\linewidth}\raggedright
Hardware
\end{minipage} & \begin{minipage}[b]{\linewidth}\raggedright
OS / Framework
\end{minipage} & \begin{minipage}[b]{\linewidth}\raggedright
Role
\end{minipage} \\
\midrule\noalign{}
\endhead
\bottomrule\noalign{}
\endlastfoot
\textbf{Jetson} & Orin NX & Linux + ROS 2 & Perception, planning, high-level
control \\
\textbf{RT} & ESP32-S3 @ 240 MHz & FreeRTOS (ESP-IDF) & Realtime physics,
steering, CAN gateway \\
\textbf{SYS} & ESP32-S3 @ 240 MHz & FreeRTOS (ESP-IDF) & Safety, motor
actuation, body control \\
\end{longtable}

This is not an arbitrary split --- each node runs a fundamentally different
class of software.

\begin{center}\rule{0.5\linewidth}{0.5pt}\end{center}

\section{Why three nodes?}\label{why-three-nodes}

\subsection{1. Different failure
characteristics}\label{different-failure-characteristics}

\begin{longtable}[]{@{}
  >{\raggedright\arraybackslash}p{(\columnwidth - 4\tabcolsep) * \real{0.1875}}
  >{\raggedright\arraybackslash}p{(\columnwidth - 4\tabcolsep) * \real{0.4062}}
  >{\raggedright\arraybackslash}p{(\columnwidth - 4\tabcolsep) * \real{0.4062}}@{}}
\toprule\noalign{}
\begin{minipage}[b]{\linewidth}\raggedright
Node
\end{minipage} & \begin{minipage}[b]{\linewidth}\raggedright
Failure mode
\end{minipage} & \begin{minipage}[b]{\linewidth}\raggedright
Consequence
\end{minipage} \\
\midrule\noalign{}
\endhead
\bottomrule\noalign{}
\endlastfoot
Jetson & Kernel panic, ROS node crash, OOM kill & Loses planning --- vehicle
must coast/stop safely \\
RT & Watchdog reset, CAN bus-off & Loses physics/steering --- SYS heartbeats
detect this \\
SYS & Watchdog reset, CAN bus-off & Loses motor/brake actuation. MTR maintains
motor kill independently. Brake: SEB enters comm-fault on CAN loss (behavior
unverified --- hold or release). RT can still command EPS-C. Physical brake
lever is a GPIO input to SYS --- it cannot actuate SEB while SYS is rebooting
(by-wire system, no mechanical hydraulic path). \\
\end{longtable}

Each node has independent failure modes. A Jetson crash should not affect RT or
SYS. An RT crash should not prevent SYS from keeping the vehicle stopped.

\subsection{2. Different realtime
guarantees}\label{different-realtime-guarantees}

\begin{itemize}
\tightlist
\item
  \textbf{Jetson:} No hard realtime. Linux scheduler, memory pressure, disk I/O
  can introduce 10--100 ms jitter. Good enough for planning at 10--100 Hz.
\item
  \textbf{RT ESP32:} Hard realtime. FreeRTOS with 1000 Hz tick, task priorities
  enforced by preemptive scheduler. Kinematics + PID at 100 Hz with \textless50
  µs jitter. Steering CAN at 50 Hz with strict periodicity.
\item
  \textbf{SYS ESP32:} Hard realtime. Safety GPIO poll at 20 Hz, motor DAC at 100
  Hz, brake CAN at 50 Hz.
\end{itemize}

\subsection{3. Different safety criticality}\label{different-safety-criticality}

\begin{itemize}
\tightlist
\item
  \textbf{Jetson:} QM (quality managed). A bug here should be uncomfortable
  (jerky driving) but not dangerous. Safety envelope enforced by RT.
\item
  \textbf{RT:} ASIL-like (safety critical). Must guarantee steering angle is
  safe, speed is limited, ESTOP propagates. Failure can cause injury.
\item
  \textbf{SYS:} ASIL-like (safety critical). Must guarantee motor stops on
  ESTOP, brake engages, gear disengages. Failure can cause runaway.
\end{itemize}

\subsection{4. Physical separation enables independent power
domains}\label{physical-separation-enables-independent-power-domains}

\begin{itemize}
\tightlist
\item
  Jetson runs on 12 V (from DC-DC).
\item
  ESP32s run on 3.3 V (from 12 V via LDO).
\item
  The 72 V traction domain is galvanically isolated from all compute.
\end{itemize}

\begin{center}\rule{0.5\linewidth}{0.5pt}\end{center}

\section{Responsibility split (complete)}\label{responsibility-split-complete}

\begin{longtable}[]{@{}lccc@{}}
\toprule\noalign{}
Concern & Jetson & RT & SYS \\
\midrule\noalign{}
\endhead
\bottomrule\noalign{}
\endlastfoot
Perception / planning & ✓ & & \\
ROS 2 → CAN bridge & ✓ & & \\
CAN gateway (low ↔ high) & & ✓ & \\
Tricycle kinematics & & ✓ & \\
Speed PID & & ✓ & \\
Steering angle compute + CAN TX (\texttt{0x169}) & & ✓ & \\
Steering boot sync (LBS) & & ✓ & \\
Steering safety (dynamic clamp, hard-stops, following error) & & ✓ & \\
Obstacle speed limit & & ✓ & \\
Command staleness watchdog & & ✓ & \\
E-stop GPIO + button & & & ✓ \\
Brake lever → CAN (\texttt{0x7B9}, 50 Hz) & & & ✓ \\
Brake boot sync (LBS) & & & ✓ \\
Brake rolling counter + checksum & & & ✓ \\
DC-DC converter CAN control (\texttt{0x012}) & & & ✓ \\
Heartbeat monitoring & ✓ (RT, high) & ✓ (Jetson, high) & ✓ (RT, low) \\
Mode switch reading & & & ✓ \\
Throttle ADC read (0--5V) & & & ✓ \\
Throttle MCP4725 DAC output (0--5V) & & & ✓ \\
Gear 72V read (TLP281 opto) & & & ✓ \\
Gear 72V output (relay module) & & & ✓ \\
12V accessory power relay & & & ✓ \\
Mode indicator lights & & & ✓ \\
Signal lights (turn, brake, head) & & & ✓ \\
System diagnostics & & & ✓ \\
\end{longtable}

\begin{center}\rule{0.5\linewidth}{0.5pt}\end{center}

\section{Communication paths}\label{communication-paths}

\begin{verbatim}
Jetson ── High CAN ── RT ── Low CAN ── SYS ── Low CAN ── Actuators
  │                    │                   │
  │  0x300 drive cmd   │                   │
  │  0x301 brake req   │  0x204 drive sp   │
  │  0x302 lights      │  0x169 steer cmd  │
  │                    │                   │
  │  0x011 safety ◄────┤                   │
  │  0x120 throttle ◄──┤                   │
  │  0x210 state ◄─────┤                   │
  │  0x220 PID fb ◄────┤                   │
  │  0x400 obstacle ◄──┤                   │
  │  0x600 diag ◄──────┤                   │
  │                    │                   │  0x7B9 brake cmd → SEB
  │                    │                   │  0x012 dcdc cmd  → DC-DC
  │                    │                   │  0x721 SEB status ◄──
  │                    │  0x201 EPS-C ◄─────┤
\end{verbatim}

\textbf{Key property:} Jetson never talks directly to SYS or to actuators. All
commands pass through RT, which validates and clamps them.

\begin{center}\rule{0.5\linewidth}{0.5pt}\end{center}

\section{Dual-CAN on RT ESP32-S3}\label{dual-can-on-rt-esp32-s3}

RT is the only node with two CAN interfaces. It uses two different controllers:

\begin{longtable}[]{@{}
  >{\raggedright\arraybackslash}p{(\columnwidth - 8\tabcolsep) * \real{0.1087}}
  >{\raggedright\arraybackslash}p{(\columnwidth - 8\tabcolsep) * \real{0.2391}}
  >{\raggedright\arraybackslash}p{(\columnwidth - 8\tabcolsep) * \real{0.2391}}
  >{\raggedright\arraybackslash}p{(\columnwidth - 8\tabcolsep) * \real{0.1304}}
  >{\raggedright\arraybackslash}p{(\columnwidth - 8\tabcolsep) * \real{0.2826}}@{}}
\toprule\noalign{}
\begin{minipage}[b]{\linewidth}\raggedright
Bus
\end{minipage} & \begin{minipage}[b]{\linewidth}\raggedright
Controller
\end{minipage} & \begin{minipage}[b]{\linewidth}\raggedright
Interface
\end{minipage} & \begin{minipage}[b]{\linewidth}\raggedright
Pins
\end{minipage} & \begin{minipage}[b]{\linewidth}\raggedright
Transceiver
\end{minipage} \\
\midrule\noalign{}
\endhead
\bottomrule\noalign{}
\endlastfoot
Low-level & Built-in TWAI & Direct GPIO & TX=5, RX=4 & SN65HVD230 \\
High-level & MCP2515 & SPI & SCK=36, MOSI=37, MISO=38, CS=39, INT=40 &
SN65HVD230 \\
\end{longtable}

\subsection{Why two different controllers?}\label{why-two-different-controllers}

The ESP32-S3 has only one built-in TWAI controller (the S3 removed the second
one that the original ESP32 had). To get two CAN buses on one MCU, we add an
external MCP2515 via SPI.

\begin{itemize}
\tightlist
\item
  \textbf{TWAI (built-in):} Low-latency, hardware ACK, hardware filters,
  hardware TX mailbox. Used for the \textbf{low-level} bus (actuators) because
  steering commands need the lowest possible jitter (50 Hz fixed-rate).
\item
  \textbf{MCP2515 (SPI):} Additional latency from SPI transaction
  (\textasciitilde10 µs per frame at 10 MHz SPI). Used for the
  \textbf{high-level} bus (Jetson) where 100 Hz telemetry jitter of 1--2 ms is
  acceptable.
\end{itemize}

\subsection{SPI bandwidth check}\label{spi-bandwidth-check}

At 500 kbit/s CAN, worst-case frame is \textasciitilde130 bits (extended ID,
8-byte DLC, stuff bits). That's \textasciitilde260 µs on the bus. At 10 MHz SPI,
reading a 13-byte MCP2515 RX buffer takes \textasciitilde10 µs. SPI bandwidth is
not a bottleneck.

\begin{center}\rule{0.5\linewidth}{0.5pt}\end{center}

\section{Why not a single ESP32?}\label{why-not-a-single-esp32}

An earlier design note (\texttt{rtos-architecture.md}) describes a single
ESP32-S3 running all tasks on two cores with two TWAI controllers. That design
has been superseded for these reasons:

\begin{enumerate}
\def\labelenumi{\arabic{enumi}.}
\tightlist
\item
  \textbf{ESP32-S3 has only one TWAI.} The dual-TWAI design requires the
  original ESP32 (not S3), which has less SRAM and slower flash.
\item
  \textbf{No physical separation of safety.} If all software runs on one MCU, a
  single firmware bug (stack overflow, wild pointer, priority inversion) can
  take down both safety and actuation. With separate MCUs, SYS continues to
  function even if RT crashes.
\item
  \textbf{Independent power sequencing.} SYS powers up first and brings up the
  12 V rail. RT powers up second. Jetson boots last (longest boot time). This
  ordering is harder to enforce with a single MCU.
\item
  \textbf{CAN bus isolation is physical, not just logical.} With two TWAI
  controllers on one chip, both buses share a silicon die --- a latch-up on one
  CAN transceiver can propagate. Separate MCUs provide true galvanic isolation
  between bus domains (only connected through the CAN transceivers).
\end{enumerate}

\begin{center}\rule{0.5\linewidth}{0.5pt}\end{center}

\emph{See also: {[}{[}can-gateway-bridging{]}{]} for the forwarding rules
between buses, {[}{[}defense-in-depth-safety{]}{]} for layered safety across
nodes, {[}{[}architecture{]}{]} §1 for topology, §5 for responsibility split,
§7.2 for dual CAN hardware.}


\chapter{Emergency System Safety Analysis --- Three Critical
Issues}\label{emergency-system-safety-analysis-three-critical-issues}

This document analyzes three safety issues discovered during review of the ESTOP
and emergency handling design. Each issue is traced through the architecture
with its causal chain, quantified for risk, and assigned a recommended fix.

\begin{center}\rule{0.5\linewidth}{0.5pt}\end{center}

\section{Issue 1: ESTOP Exit Race Condition --- Steering Ramp Interrupted by
START
Button}\label{issue-1-estop-exit-race-condition-steering-ramp-interrupted-by-start-button}

\subsection{1.1 The scenario}\label{the-scenario}

\begin{enumerate}
\def\labelenumi{\arabic{enumi}.}
\tightlist
\item
  Non-obstacle ESTOP is triggered (button press, heartbeat loss, following
  error, CAN \texttt{0x001}).
\item
  RT steer state machine begins centering ramp: transmits
  \texttt{0x169\ VCU\_SES\_REQ} at 50 Hz with decreasing angle targets at 20°/s.
\item
  \textbf{Rider presses START button before the ramp completes.}
\item
  SYS \texttt{mode\_task} detects GPIO32 falling edge →
  \texttt{mode\_set(Manual)}.
\item
  SYS transmits CAN \texttt{0x110\ SYS\_MODE\_CMD} with \texttt{mode=0}
  (MANUAL).
\item
  RT \texttt{dispatch} receives \texttt{0x110} → calls
  \texttt{mode\_set(Manual)}.
\item
  RT steer SM: \textbf{MANUAL mode → RT does NOT send \texttt{0x169}.} Per
  architecture §7.6 mode behavior table: ``MANUAL: RT does NOT send 0x169. EPS-C
  standalone.''
\item
  \texttt{0x169} transmission stops \textbf{immediately} at whatever angle the
  ramp had reached.
\item
  EPS-C: 20 ms without \texttt{0x169} → internal comm-fault timeout.
\item
  EPS-C comm-fault behavior: unit-specific and \textbf{unknown}. May lock at
  current angle, go limp, or freewheel.
\end{enumerate}

\subsection{1.2 Causal trace}\label{causal-trace}

\begin{verbatim}
START btn press
  │
  ├── SYS mode_task @ 10 Hz
  │     └── g_mode = Manual
  │     └── CAN 0x110 {mode=MANUAL}
  │
  ├── RT dispatch receives 0x110
  │     └── mode_set(Manual)
  │           └── steer SM: ESTOP+STEER_ACTIVE → MANUAL
  │                 └── Stop transmitting 0x169   ← RACE: ramp may be mid-flight
  │
  └── EPS-C: 0x169 stops
        └── T+20ms: comm-fault timeout
              └── Behavior: UNKNOWN (lock / limp / freewheel at current angle)
\end{verbatim}

\textbf{Key architectural detail:} The START button's \texttt{mode\_set(Manual)}
is unconditional. There is no check on whether the steering ramp is complete.
The mode transition is instant. The steer SM obediently stops transmitting
\texttt{0x169} per MANUAL-mode rules --- it has no awareness that a ramp was in
progress.

\subsection{1.3 Quantified risk}\label{quantified-risk}

\begin{longtable}[]{@{}
  >{\raggedright\arraybackslash}p{(\columnwidth - 2\tabcolsep) * \real{0.6111}}
  >{\raggedright\arraybackslash}p{(\columnwidth - 2\tabcolsep) * \real{0.3889}}@{}}
\toprule\noalign{}
\begin{minipage}[b]{\linewidth}\raggedright
Parameter
\end{minipage} & \begin{minipage}[b]{\linewidth}\raggedright
Value
\end{minipage} \\
\midrule\noalign{}
\endhead
\bottomrule\noalign{}
\endlastfoot
Worst-case starting angle & ±40° (software hard-stop limit) \\
Ramp rate & 20°/s \\
Maximum ramp duration & 2.0 seconds (40° ÷ 20°/s) \\
Window of vulnerability & 0 to 2.0 seconds after ESTOP entry \\
Angle if interrupted at t=0.5s & 40° − (0.5 × 20) = \textbf{30° off-center} \\
Angle if interrupted at t=1.0s & \textbf{20° off-center} \\
EPS-C behavior after 0x169 stops & Unknown --- may \textbf{lock at that angle}
or go limp \\
\end{longtable}

At 30° steering lock in MANUAL mode, the vehicle cannot be ridden straight. The
rider must either: - Power-cycle and hope EPS-C releases on power-loss (unknown
behavior). - Wrestle the locked steering mechanically (risk of damaging EPS-C
gear train). - Press ESTOP again (which won't help --- already in a bad state).

\textbf{Risk assessment:} High severity, medium probability. The START button is
the only ESTOP exit path --- a rider who entered ESTOP on a curved road will
naturally want to exit as soon as they feel safe, likely before the 2-second
ramp completes if they're not aware of the ramp duration. Combined with unknown
EPS-C timeout behavior, this creates a non-deterministic safety outcome.

\subsection{1.4 Why it wasn't caught earlier}\label{why-it-wasnt-caught-earlier}

The two-tier ESTOP steering design (gap \#3 resolution in architecture §12)
focused on \emph{entering} ESTOP correctly --- hold vs ramp based on trigger ---
but never considered the \emph{exit} transition. The ramp was designed as an
ESTOP-mode behavior, and the MANUAL-mode behavior (``RT does NOT send 0x169'')
was designed for normal MANUAL operation. The boundary between them ---
ESTOP→MANUAL transition mid-ramp --- fell through the gap.

\subsection{1.5 Recommended fix: Deferred mode transition with ramp
completion}\label{recommended-fix-deferred-mode-transition-with-ramp-completion}

\textbf{Approach:} RT defers the ESTOP→MANUAL steering mode transition until the
centering ramp completes. No new CAN messages needed --- this is a local change
to RT's steer state machine.

\textbf{Implementation:}

In RT's \texttt{dispatch} task, when processing \texttt{0x110\ SYS\_MODE\_CMD}
with \texttt{mode=MANUAL} and the current mode is ESTOP:

\begin{verbatim}
// In dispatch task, on receiving 0x110 mode=MANUAL while in ESTOP:
if (g_steer_ramp_in_progress) {
    // Defer the mode transition for steering purposes.
    // SYS has already transitioned to MANUAL (brake, motor, lights all MANUAL).
    // RT continues the steering centering ramp.
    // Once ramp completes, steer SM transitions to MANUAL (stops 0x169).
    g_steer_deferred_exit_to_manual = true;
    // Do NOT call mode_set(Manual) yet — keep RT in ESTOP mode locally.
    // All other subsystems (motor setpoint, brake) follow SYS immediately.
} else {
    mode_set(Manual);  // ramp already complete — safe to transition
}
\end{verbatim}

In the steer state machine's \texttt{STEER\_ACTIVE} state:

\begin{verbatim}
case STEER_ACTIVE:
    if (g_mode == Mode::Estop && g_steer_ramp_in_progress) {
        // Continue centering ramp
        ramp_step();
        send_0x169(ramp_target);
        if (ramp_complete()) {
            g_steer_ramp_in_progress = false;
            if (g_steer_deferred_exit_to_manual) {
                // Ramp completed — safe to go silent
                stop_0x169();
                steer_state = STEER_IDLE_MANUAL;  // or equivalent
            }
        }
    } else if (g_mode == Mode::Manual) {
        // Normal MANUAL: silent
        stop_0x169();
    }
    break;
\end{verbatim}

\textbf{Key properties of this fix:} - No new CAN messages. No protocol change.
- No mode mismatch between SYS and RT: SYS transitions to MANUAL immediately
(brake lever responsive, motor pass-through, lights MANUAL behavior). RT defers
only the \emph{steering} aspect. - The brake transitions from ESTOP (max stroke)
to MANUAL (lever-based) immediately on START --- \textbf{this is correct}
because the rider pressed START specifically to regain control. - The steering
ramp continues uninterrupted. Once centered, RT goes silent per MANUAL rules. -
If the ramp cannot complete (mechanical jam → following error persists), the
existing fallback applies: silent-stop after \textgreater5° error for
\textgreater1s. The START button press doesn't change this. - Timeout: if the
ramp doesn't complete within 5 seconds (e.g., EPS-C not responding), RT
transitions to STEER\_FAULT. SYS is already in MANUAL, so rider has motor and
brake control but no steering assist.

\textbf{Edge case --- rider pressing START during obstacle-triggered ESTOP:} In
obstacle-triggered ESTOP, the steering holds current angle then silent-stops
after 500ms. During the hold phase,
\texttt{g\_steer\_ramp\_in\_progress\ =\ false} (it's a hold, not a ramp). So
START button exits immediately --- which is correct: the steering is already at
a fixed angle, and EPS-C will timeout-fault on 0x169 stop as it would in any
MANUAL transition. During the silent-stop phase (after 500ms), there's no 0x169
transmission at all --- START exits normally.

\begin{center}\rule{0.5\linewidth}{0.5pt}\end{center}

\section{Issue 2: SEB Boot Sync Timeout Contradiction --- ``Brake Lever Always
Works'' is
False}\label{issue-2-seb-boot-sync-timeout-contradiction-brake-lever-always-works-is-false}

\subsection{2.1 The contradiction}\label{the-contradiction}

Three statements in the architecture contradict each other:

\textbf{Statement A} (architecture §8.10, error handling table): \textgreater{}
SEB sync timeout: No 0x721 within 2s of boot → \texttt{BRAKE\_FAULT}, lever inop

\textbf{Statement B} (architecture §4.3, SYS failure section): \textgreater{}
Physical fallbacks during SYS failure: Brake lever (GPIO2) still works ---
physical hydraulic brake independent of CAN.

\textbf{Statement C} (architecture §8.6, BRAKE\_FAULT state): \textgreater{}
BRAKE\_FAULT: Stop transmitting

Statements B and C cannot both be true. If BRAKE\_FAULT means ``stop
transmitting 0x7B9'', then the brake lever --- which is just a GPIO2 input
switch read by SYS firmware --- cannot result in braking because the command
never reaches SEB.

Statement B is \textbf{factually incorrect}. The E-Trike has a
\textbf{brake-by-wire} system: the brake lever is an electrical switch, not a
hydraulic master cylinder. The actuation path is always:

\begin{verbatim}
Lever (GPIO2) → SYS firmware → CAN 0x7B9 → SEB → hydraulic pressure to calipers
\end{verbatim}

There is no mechanical or hydraulic bypass. ``Independent of CAN'' is false ---
the brake always depends on CAN 0x7B9 reaching SEB.

\subsection{2.2 How SEB boot sync failure
occurs}\label{how-seb-boot-sync-failure-occurs}

The Listen Before Speaking (LBS) sequence for the brake requires:

\begin{enumerate}
\def\labelenumi{\arabic{enumi}.}
\tightlist
\item
  \textbf{BRAKE\_BOOT\_WAIT} (500ms): Power-on delay. Do not transmit.
\item
  \textbf{BRAKE\_LISTEN\_SYNC}: Wait for \texttt{0x721\ SEB\_STATUS}. Extract
  current stroke. Wait for alignment.
\item
  \textbf{BRAKE\_ACTIVE}: Transmit \texttt{0x7B9} at 50 Hz continuously.
\end{enumerate}

If step 2 times out (no \texttt{0x721} within 2 seconds), the SM transitions to
BRAKE\_FAULT and \textbf{stops transmitting permanently}. There is no recovery
path defined.

This can happen from: - SEB not powered (separate power supply issue) - SEB CAN
interface not initialized yet (boots slower than SYS) - CAN bus wiring fault
between SYS and SEB - SEB internal fault preventing status transmission

\subsection{2.3 Causal trace --- worst case}\label{causal-trace-worst-case}

\begin{verbatim}
Power-on
  │
  ├── SYS boots in MANUAL mode (safe default)
  ├── brake_init(): BRAKE_BOOT_WAIT (500ms) → BRAKE_LISTEN_SYNC
  ├── SEB is powered but slower to boot (internal self-test)
  │
  ├── T+2500ms (500ms wait + 2000ms listen): no 0x721 received
  │     └── BRAKE_FAULT
  │           └── Stop transmitting 0x7B9
  │
  ├── Rider squeezes brake lever
  │     └── GPIO2 reads LOW
  │     └── brake_task: state == BRAKE_FAULT → does NOT transmit 0x7B9
  │     └── SEB receives no command
  │     └── **No braking occurs**
  │
  └── Rider has NO brake. Vehicle is in MANUAL mode with throttle and steering functional.
      The only way to stop: ESTOP button (but ESTOP also sends brake via 0x7B9 — same dead path).
\end{verbatim}

\textbf{Wait --- does ESTOP also fail?} Yes. In ESTOP, SYS sends \texttt{0x7B9}
with stroke=max from the \texttt{brake\_task}. But if the brake SM is in
BRAKE\_FAULT, it's not transmitting. The ESTOP brake command also goes through
the same dead CAN path.

However, the ESTOP button also: - Kills motor (MCP4725 = 0V via
\texttt{motor\_task} --- different task, still works) - Kills gear (relays OFF
--- different task, still works) - Turns off DC-DC (CAN 0x012 via
\texttt{dcdc\_task} --- different task, still works)

So ESTOP stops the motor but doesn't engage the brake in this scenario. The
vehicle coasts to a stop rather than braking actively.

\subsection{2.4 Quantified risk}\label{quantified-risk-1}

\begin{longtable}[]{@{}
  >{\raggedright\arraybackslash}p{(\columnwidth - 2\tabcolsep) * \real{0.6111}}
  >{\raggedright\arraybackslash}p{(\columnwidth - 2\tabcolsep) * \real{0.3889}}@{}}
\toprule\noalign{}
\begin{minipage}[b]{\linewidth}\raggedright
Parameter
\end{minipage} & \begin{minipage}[b]{\linewidth}\raggedright
Value
\end{minipage} \\
\midrule\noalign{}
\endhead
\bottomrule\noalign{}
\endlastfoot
Trigger condition & SEB CAN status not received within 2s of boot \\
Probability & Low (requires SEB power/CAN fault) but non-zero --- any CAN wiring
issue, SEB slow boot, or SEB internal fault \\
Consequence & \textbf{Complete brake loss} --- lever inoperative, ESTOP braking
inoperative \\
Mitigation in current design & None. BRAKE\_FAULT is terminal with no
recovery. \\
Vehicle state & MANUAL mode, motor and steering functional, no brake \\
Stop distance without brake & At 25 km/h on flat ground, coast-to-stop is
\textgreater50m \\
\end{longtable}

This is a \textbf{single-point failure} that takes out the brake: one CAN ID not
arriving at boot time disables the entire braking system for the remainder of
the power cycle.

\subsection{2.5 Why the architecture makes this
mistake}\label{why-the-architecture-makes-this-mistake}

The brake LBS was modeled on the steering LBS, where FAULT is a recoverable
state because EPS-C goes \textbf{standalone} in MANUAL mode --- the steering
wheel mechanically turns the rack regardless of CAN. The brake has no equivalent
``standalone'' mode. SEB requires continuous CAN commands to maintain any brake
pressure. The LBS pattern doesn't transfer cleanly to a by-wire actuator with no
mechanical fallback.

The statement ``physical brake lever independent of CAN'' appears to assume a
traditional hydraulic brake where the lever directly pushes fluid. But the SEB
is electro-hydraulic --- the lever is just a switch input to the MCU. This is a
domain confusion between traditional and by-wire braking.

\subsection{2.6 Recommended fix: Degraded transmit mode instead of terminal
FAULT}\label{recommended-fix-degraded-transmit-mode-instead-of-terminal-fault}

\textbf{Approach:} Replace BRAKE\_FAULT (terminal stop) with BRAKE\_DEGRADED
(transmit with safe defaults despite no sync). The brake is too critical to go
silent.

\textbf{Implementation:}

Replace the BRAKE\_FAULT state with BRAKE\_DEGRADED:

\begin{verbatim}
BRAKE_DEGRADED:
  - Transmit 0x7B9 at 50 Hz with conservative defaults
  - Stroke = 0 (no brake) when lever released
  - Stroke = kBrakeManualStroke (~15 mm) when lever pressed (GPIO2 LOW)
  - Rolling counter still increments
  - Checksum still computed
  - CAN 0x600 SYS_DIAG_RPT includes brake_degraded flag
  - Recovery: when first valid 0x721 arrives, sync current stroke
    and transition to BRAKE_ACTIVE
\end{verbatim}

This means: - The brake lever \textbf{always works} --- even without boot sync.
The statement ``brake lever always works'' becomes true. - The only difference
vs normal operation: without sync, SYS doesn't know the exact current stroke, so
it commands absolute positions rather than relative changes. - SEB may reject
some frames if the initial stroke value doesn't match its internal state. But a
frame with valid checksum + rolling counter at a plausible stroke value (0--27
mm) is likely to be accepted. - Once the first \texttt{0x721} arrives, normal
synced operation resumes.

\textbf{Why this is safe without sync:} - Stroke=0mm (released) and stroke=15mm
(pressed) are always within SEB's valid range (0--27mm). - The SEB's internal
PID will move to the commanded stroke regardless of starting position. - The
rolling counter increments → SEB sees liveness. - The checksum is correct → SEB
accepts the frame. - Worst case: SEB was at 27mm and receives stroke=0 → brake
releases. This is the correct MANUAL behavior (lever released = no brake). SEB
was at 0mm and receives stroke=15mm → brake engages. Also correct.

\textbf{Alternative consideration --- hardware fix:} If SYNTREE's protocol
requires exact sync (rejects frames when stroke doesn't match internal state), a
hardware solution is needed: wire the brake lever as a \textbf{redundant direct
input to SEB} (if SEB has a discrete brake input pin). This bypasses the MCU
entirely. But this depends on SEB hardware capability and adds wiring
complexity.

\begin{center}\rule{0.5\linewidth}{0.5pt}\end{center}

\section{Issue 3: Watchdog Reset Unbraked Window --- No Deterministic Brake
During MCU
Reset}\label{issue-3-watchdog-reset-unbraked-window-no-deterministic-brake-during-mcu-reset}

\subsection{3.1 The scenario}\label{the-scenario-1}

\begin{enumerate}
\def\labelenumi{\arabic{enumi}.}
\tightlist
\item
  SYS ESP32-S3 hangs (stack overflow, deadlock, crystal failure, latch-up).
\item
  \texttt{safety\_task} stops toggling GPIO23.
\item
  \textbf{T+100ms:} TPS3850 external watchdog IC asserts RST → ESP32 EN pin
  pulled LOW.
\item
  \textbf{All SYS GPIOs} go to reset state (high-impedance input).
\item
  Motor controller: MCP4725 loses I2C communication → output defaults to 0V (or
  last held value, depending on MCP4725 configuration). \textbf{Likely 0V}
  (MCP4725 has power-on reset to 0V, and I2C loss may not trigger a reset --- it
  holds last value until VCC drops).
\item
  Gear relays: GPIOs float → NPN transistor base goes LOW → relay coil
  de-energizes → contacts open → all gear lines = Neutral. \textbf{Confirmed
  safe.}
\item
  \textbf{T+120ms:} SEB has received no \texttt{0x7B9} for 20ms → internal
  comm-fault timeout.
\item
  SEB comm-fault behavior: \textbf{Unknown / TBD.} Options:

  \begin{itemize}
  \tightlist
  \item
    \textbf{Hold} last commanded pressure/stroke → brake maintained.
    \textbf{Safe.}
  \item
    \textbf{Release} pressure → brake lost. \textbf{Unsafe.}
  \item
    \textbf{Ramp release} over N seconds → partial brake during window.
    \textbf{Partially safe.}
  \end{itemize}
\item
  \textbf{T+\textasciitilde200ms:} ESP32 bootloader starts. SYS begins firmware
  init.
\item
  \textbf{T+\textasciitilde250ms:} FreeRTOS scheduler starts.
  \texttt{safety\_task} begins WDT toggle (re-arms watchdog).
\item
  \textbf{T+\textasciitilde700ms:} \texttt{brake\_init()} completes BOOT\_WAIT
  (500ms). Enters BRAKE\_LISTEN\_SYNC.
\item
  \textbf{T+\textasciitilde2700ms:} If SEB responds to \texttt{0x721}, brake
  enters ACTIVE. First \texttt{0x7B9} command transmitted.
\item
  \textbf{Total unbraked window if SEB releases on comm-fault:} T+120ms to
  T+2700ms = \textbf{\textasciitilde2.58 seconds.}
\end{enumerate}

\subsection{3.2 The MCP4725 ambiguity}\label{the-mcp4725-ambiguity}

The MCP4725 DAC is powered from the 5V rail (12V→5V LDO), not directly from the
ESP32's 3.3V. When the ESP32 resets:

\begin{itemize}
\tightlist
\item
  The I2C bus (SDA=GPIO15, SCL=GPIO16) goes high-impedance.
\item
  The MCP4725 retains its last programmed output voltage \textbf{as long as VCC
  (5V) is maintained.} \textbf{Motor kill during SYS reset --- covered by MTR.}
  MTR STM32 is the EGAS Level 1 Function Controller and runs on a separate MCU.
  When SYS watchdog resets, MTR continues running independently. MTR drives its
  own MCP4725 DAC and gear relays. MTR has a direct ESTOP GPIO input shared with
  SYS GPIO1. MTR also monitors SYS heartbeat (0x7FE) and broadcasts CAN 0x001 on
  timeout. During SYS reset, MTR kills motor (MCP4725=0V) and cuts gear (relays
  OFF) locally --- zero CAN dependency.
\end{itemize}

\textbf{Brake during SYS reset --- NOT covered.} SEB is commanded exclusively by
SYS via CAN 0x7B9. MTR has no SEB command path. RT could command SEB (it already
sends 0x7B9 in AUTO mode) but does not take over on SYS heartbeat loss in the
current design --- see Issue 7 for the fix. During the SYS reboot window, SEB
enters comm-fault after 20ms. Behavior is unverified (hold or release).

\subsection{3.3 Quantified risk (revised)}\label{quantified-risk-revised}

\begin{longtable}[]{@{}
  >{\raggedright\arraybackslash}p{(\columnwidth - 2\tabcolsep) * \real{0.6111}}
  >{\raggedright\arraybackslash}p{(\columnwidth - 2\tabcolsep) * \real{0.3889}}@{}}
\toprule\noalign{}
\begin{minipage}[b]{\linewidth}\raggedright
Parameter
\end{minipage} & \begin{minipage}[b]{\linewidth}\raggedright
Value
\end{minipage} \\
\midrule\noalign{}
\endhead
\bottomrule\noalign{}
\endlastfoot
Trigger & SYS firmware hang → external WDT fires \\
Motor response & Safe --- MTR kills motor locally (0V throttle, N gear) \\
Brake response & \textbf{Non-deterministic} --- depends on SEB comm-fault
behavior \\
Unbraked window (if SEB releases) & \textasciitilde2.58 seconds (20ms timeout +
\textasciitilde2.5s SYS reboot + brake LBS) \\
Distance at 25 km/h & \textasciitilde18 meters \\
Probability & Low (requires firmware hang + SEB releases on comm-fault) \\
\end{longtable}

\subsection{3.4 Why the SEB comm-fault behavior matters
critically}\label{why-the-seb-comm-fault-behavior-matters-critically}

The entire safety of the watchdog reset window hinges on one unknown:
\textbf{does SEB hold or release on CAN timeout?}

\begin{itemize}
\tightlist
\item
  If SEB \textbf{holds}: the brake stays engaged during the reset window. The
  vehicle decelerates under existing brake pressure. Safe.
\item
  If SEB \textbf{releases}: the vehicle coasts with no brake for
  \textasciitilde2.5s. At 25 km/h downhill, this is dangerous.
\end{itemize}

This is an \textbf{uncontrolled assumption} in the safety architecture. We are
depending on a behavior we haven't verified and can't control.

\subsection{3.5 Recommended fix: Hardware brake-hold on watchdog
reset}\label{recommended-fix-hardware-brake-hold-on-watchdog-reset}

\textbf{Approach:} Add a dedicated hardware signal from the TPS3850 watchdog
IC's RST output to a relay that controls SEB power or a discrete brake-engage
input on SEB (if available).

\textbf{Option A --- SEB power-cut relay (preferred):}

\begin{verbatim}
TPS3850 RST (active-low) ──► P-channel MOSFET gate
                                    │
                           SEB 12V power ──► MOSFET (normally ON)
                                    │
                           RST asserted → MOSFET OFF → SEB loses power
\end{verbatim}

If SEB loses power: - Its internal failsafe should engage (electro-hydraulic
units typically default to pressure-hold or full-release on power loss --- need
to verify with SYNTREE spec). - If SEB defaults to \textbf{hold} on power loss:
this is ideal. Power-cycle SEB, brake holds. - If SEB defaults to
\textbf{release} on power loss: this makes things worse.

\textbf{Option B --- Dedicated brake relay (more robust):}

\begin{verbatim}
TPS3850 RST ──► NC relay coil (energized during normal operation)
                      │
                Relay contacts: SEB enable line or brake pressure dump valve
                RST asserted → relay de-energizes → contacts close → brake engage
\end{verbatim}

This uses the same NC wiring principle as the ESTOP button. A watchdog reset (or
any power loss) de-energizes the relay, which mechanically engages a brake-hold
signal. The relay can be wired to: - A discrete ``emergency brake'' input on SEB
(if it has one --- check SYNTREE pinout) - A solenoid valve that locks hydraulic
pressure in the brake line

\textbf{Option C --- MTR as brake proxy (software, less robust):}

Add SEB command capability to MTR. When MTR detects SYS heartbeat loss
(\texttt{0x7FE} timeout), MTR starts transmitting \texttt{0x7B9} with
stroke=max. This requires: - MTR to have its own CAN ID for SEB command (can't
use 0x7B9 --- already used by RT/SYS in mode-gated scheme) - Or MTR to use a new
CAN ID that SEB accepts as a secondary command - SYNTREE SEB must support
multiple command sources (unlikely --- it's preprogrammed for specific IDs)

This is architecturally complex and depends on SEB capabilities we may not have.

\textbf{Option D --- Conservative: verify SEB comm-fault behavior and document
the assumption}

If SEB is verified to \textbf{hold pressure on CAN timeout}, the current design
is acceptable --- the unbraked window is only 20ms (time to comm-fault
detection), not 2.5s. This should be tested empirically. If SEB holds, document
it as a verified safety property. If SEB releases, Options A or B become
mandatory.

\subsection{3.6 Immediate action}\label{immediate-action}

\textbf{Test SEB comm-fault behavior:} 1. Command SEB to stroke=27mm (max brake)
via CAN 0x7B9 at 50 Hz. 2. After 1 second of sustained pressure, stop CAN
transmission. 3. Measure hydraulic pressure over the next 5 seconds. 4. If
pressure holds (\textgreater80\% after 1s): comm-fault = hold → acceptable. 5.
If pressure drops to zero within 500ms: comm-fault = release → \textbf{hardware
fix required before road testing.}

\begin{center}\rule{0.5\linewidth}{0.5pt}\end{center}

\section{Summary}\label{summary}

\begin{longtable}[]{@{}
  >{\raggedright\arraybackslash}p{(\columnwidth - 8\tabcolsep) * \real{0.1061}}
  >{\raggedright\arraybackslash}p{(\columnwidth - 8\tabcolsep) * \real{0.1515}}
  >{\raggedright\arraybackslash}p{(\columnwidth - 8\tabcolsep) * \real{0.1970}}
  >{\raggedright\arraybackslash}p{(\columnwidth - 8\tabcolsep) * \real{0.2424}}
  >{\raggedright\arraybackslash}p{(\columnwidth - 8\tabcolsep) * \real{0.3030}}@{}}
\toprule\noalign{}
\begin{minipage}[b]{\linewidth}\raggedright
Issue
\end{minipage} & \begin{minipage}[b]{\linewidth}\raggedright
Severity
\end{minipage} & \begin{minipage}[b]{\linewidth}\raggedright
Probability
\end{minipage} & \begin{minipage}[b]{\linewidth}\raggedright
Fix complexity
\end{minipage} & \begin{minipage}[b]{\linewidth}\raggedright
Recommended action
\end{minipage} \\
\midrule\noalign{}
\endhead
\bottomrule\noalign{}
\endlastfoot
\textbf{1. ESTOP exit race} & High & Medium & Low (software only, RT-local) &
Defer mode transition until ramp complete \\
\textbf{2. SEB brake lever contradiction} & High & Low & Low (software only,
SYS-local) & Replace BRAKE\_FAULT with BRAKE\_DEGRADED \\
\textbf{3. Watchdog unbraked window} & High & Low (if SEB holds) / High (if SEB
releases) & Medium (hardware) & \textbf{Test SEB comm-fault behavior first.} If
release: add NC brake relay on WDT RST line. \\
\end{longtable}

\textbf{Issue 1 and 2 are software-only fixes} that can be implemented
immediately in the RT and SYS firmware respectively. \textbf{Issue 3 requires
empirical testing} of SEB behavior before the fix path is known --- and may
require a hardware change.

\begin{center}\rule{0.5\linewidth}{0.5pt}\end{center}

\emph{See also: {[}{[}emergency-system{]}{]} §1-4 for the documented emergency
behaviors that these issues affect, {[}{[}external-watchdog{]}{]} for the
watchdog reset sequence, \href{../architecture.md}{architecture.md} §3 for mode
state machine, §8.6 for brake LBS, §8.10 for error handling table.}

\begin{center}\rule{0.5\linewidth}{0.5pt}\end{center}

\section{Issue 4: Obstacle-Triggered ESTOP ``Hold Angle'' Creates Rollover Risk
While
Cornering}\label{issue-4-obstacle-triggered-estop-hold-angle-creates-rollover-risk-while-cornering}

\subsection{4.1 The scenario}\label{the-scenario-2}

\begin{enumerate}
\def\labelenumi{\arabic{enumi}.}
\tightlist
\item
  Vehicle is cornering at 25 km/h with steering at 15° (within the 18° dynamic
  clamp limit at 10 km/h --- but at 25 km/h the dynamic clamp would limit to
  \textasciitilde5°. So for this scenario, let's assume the vehicle is at 15
  km/h, steering at 12°, which is within the dynamic clamp limit at that speed).
\item
  Jetson perception detects an obstacle → broadcasts CAN \texttt{0x001} on high
  bus.
\item
  RT receives \texttt{0x001}, classifies it as \textbf{obstacle-triggered}.
\item
  RT steer SM: ``Obstacle-triggered: hold current angle, silent-stop after
  500ms.''
\item
  Steering is held at 12°.
\item
  Brake engages at max stroke (\textasciitilde27 mm) --- full emergency braking.
\item
  \textbf{Combined cornering + hard braking on a three-wheeled vehicle} creates
  a lateral load transfer that approaches the rollover threshold.
\end{enumerate}

\subsection{4.2 Causal trace}\label{causal-trace-1}

The physics model defines the rollover threshold as:

\[a_y = \frac{v^2}{L}\tan(\delta) > \frac{g w}{2h}\]

The obstacle-triggered ESTOP design rationale is: ``don't steer --- you might
steer into the obstacle or adjacent lane.'' This is correct for straight-line
travel (δ ≈ 0°). It is \textbf{incorrect} when the vehicle is already cornering,
because:

\begin{enumerate}
\def\labelenumi{\arabic{enumi}.}
\tightlist
\item
  Holding a non-zero steering angle while emergency braking sustains lateral
  acceleration.
\item
  Hard braking transfers weight forward, reducing rear tire lateral grip.
\item
  The tricycle's single front wheel carries both steering and braking forces ---
  combined loading reduces the cornering force available before slip.
\item
  The dynamic angle clamp (architecture §7.6, Layer 6) is only applied to
  \textbf{commanded angles} in \texttt{physics\_resolve()}. During ESTOP, the
  hold behavior bypasses the dynamic clamp entirely --- there is no envelope
  check on the held angle.
\end{enumerate}

\textbf{The dynamic angle clamp says 12° is safe at 15 km/h during normal
driving.} But during emergency braking, the safe envelope shrinks because
braking reduces lateral grip margins. The hold behavior applies the
normal-driving envelope to an emergency-braking scenario --- these are different
physical conditions.

\subsection{4.3 Why the classification is too
coarse}\label{why-the-classification-is-too-coarse}

The current design classifies ESTOP triggers into exactly two categories:

\begin{longtable}[]{@{}
  >{\raggedright\arraybackslash}p{(\columnwidth - 4\tabcolsep) * \real{0.2093}}
  >{\raggedright\arraybackslash}p{(\columnwidth - 4\tabcolsep) * \real{0.3488}}
  >{\raggedright\arraybackslash}p{(\columnwidth - 4\tabcolsep) * \real{0.4419}}@{}}
\toprule\noalign{}
\begin{minipage}[b]{\linewidth}\raggedright
Trigger
\end{minipage} & \begin{minipage}[b]{\linewidth}\raggedright
Classification
\end{minipage} & \begin{minipage}[b]{\linewidth}\raggedright
Steering behavior
\end{minipage} \\
\midrule\noalign{}
\endhead
\bottomrule\noalign{}
\endlastfoot
Jetson perception, RT obstacle sensor & Obstacle & Hold angle, silent-stop after
500ms \\
Button, heartbeat, following error, stale command & Non-obstacle & Ramp to 0° at
20°/s \\
\end{longtable}

This classification assumes: - Obstacle = going straight, don't veer. -
Non-obstacle = any direction is acceptable to center.

But obstacle detection can occur while cornering. The ``hold angle'' response to
a cornering obstacle does not account for whether the held angle is safe during
emergency braking. The classification should consider vehicle state, not just
trigger source.

\subsection{4.4 Quantified risk}\label{quantified-risk-2}

\begin{longtable}[]{@{}
  >{\raggedright\arraybackslash}p{(\columnwidth - 2\tabcolsep) * \real{0.6111}}
  >{\raggedright\arraybackslash}p{(\columnwidth - 2\tabcolsep) * \real{0.3889}}@{}}
\toprule\noalign{}
\begin{minipage}[b]{\linewidth}\raggedright
Parameter
\end{minipage} & \begin{minipage}[b]{\linewidth}\raggedright
Value
\end{minipage} \\
\midrule\noalign{}
\endhead
\bottomrule\noalign{}
\endlastfoot
Scenario & Cornering at 15 km/h, δ = 12°, obstacle detected \\
Rollover threshold & Depends on CG height and track width --- not yet measured
for the physical trike \\
Braking deceleration & \textasciitilde0.6--0.8g (emergency brake on dry
pavement) \\
Weight transfer forward & \textasciitilde30--40\% of rear load shifted to
front \\
Lateral grip remaining & Significantly reduced during combined
braking+cornering \\
Severity & \textbf{High} --- rollover is a catastrophic failure mode \\
Probability & \textbf{Low} --- requires obstacle detection while cornering at
speed, which is a narrow window in urban operation \\
\end{longtable}

\subsection{4.5 Recommended fix: Dynamic-clamp the hold angle during obstacle
ESTOP}\label{recommended-fix-dynamic-clamp-the-hold-angle-during-obstacle-estop}

\textbf{Approach:} During obstacle-triggered ESTOP, the held angle is clamped to
the \textbf{dynamic angle clamp limit for the current speed}, using the SAME
envelope as normal AUTO steering commands. If the current angle exceeds the
limit, the steering ramps down to the limit (not to 0° --- to the speed-safe
limit).

\textbf{Behavior by scenario:}

\begin{longtable}[]{@{}
  >{\raggedright\arraybackslash}p{(\columnwidth - 6\tabcolsep) * \real{0.1296}}
  >{\raggedright\arraybackslash}p{(\columnwidth - 6\tabcolsep) * \real{0.2037}}
  >{\raggedright\arraybackslash}p{(\columnwidth - 6\tabcolsep) * \real{0.2778}}
  >{\raggedright\arraybackslash}p{(\columnwidth - 6\tabcolsep) * \real{0.3889}}@{}}
\toprule\noalign{}
\begin{minipage}[b]{\linewidth}\raggedright
Speed
\end{minipage} & \begin{minipage}[b]{\linewidth}\raggedright
Current δ
\end{minipage} & \begin{minipage}[b]{\linewidth}\raggedright
Dynamic limit
\end{minipage} & \begin{minipage}[b]{\linewidth}\raggedright
ESTOP hold behavior
\end{minipage} \\
\midrule\noalign{}
\endhead
\bottomrule\noalign{}
\endlastfoot
25 km/h & 2° (straight) & 5° & Hold at 2° (within limit, safe) \\
25 km/h & 8° (cornering, already past limit) & 5° & Ramp from 8° → 5° at 20°/s,
then hold at 5° \\
10 km/h & 12° (cornering) & \textasciitilde18° & Hold at 12° (within limit,
safe) \\
2 km/h & 40° (U-turn) & 40° & Hold at 40° (within limit, safe) \\
\end{longtable}

\textbf{Key properties:} - Straight-line obstacle detection: behavior is
unchanged (δ ≈ 0° is always within the dynamic limit). - Cornering at speed:
steering is reduced to the speed-safe envelope, preventing rollover. - The
vehicle still doesn't suddenly center (would risk lane departure) --- it ramps
to the safe limit, which at high speed is already near-straight
(\textasciitilde5°). - After 500ms, silent-stop as before --- EPS-C handles the
timeout.

\textbf{Implementation:} In RT's obstacle-triggered ESTOP path, before entering
the hold:

\begin{verbatim}
if (estop_trigger == ESTOP_TRIGGER_OBSTACLE) {
    // Clamp the hold angle to the dynamic limit for current speed
    float dynamic_limit = compute_dynamic_angle_limit(current_speed_mmps);
    float hold_angle = current_steering_angle_deg;
    if (fabs(hold_angle) > dynamic_limit) {
        // Ramp toward the safe limit
        hold_angle = ramp_toward(hold_angle, 
                      copysign(dynamic_limit, hold_angle), 
                      kSteerRampRateDegPerSec);
    }
    // Hold at the clamped angle, silent-stop after 500ms
}
\end{verbatim}

This reuses the existing dynamic angle clamp function that already runs in
\texttt{physics\_resolve()} --- no new math, just applying it in one more place.

\begin{center}\rule{0.5\linewidth}{0.5pt}\end{center}

\section{Issue 5: Jetson Heartbeat Loss → Controlled Stop Is Insufficient for
Perception
Failure}\label{issue-5-jetson-heartbeat-loss-controlled-stop-is-insufficient-for-perception-failure}

\subsection{5.1 The asymmetry}\label{the-asymmetry}

\begin{longtable}[]{@{}
  >{\raggedright\arraybackslash}p{(\columnwidth - 6\tabcolsep) * \real{0.3659}}
  >{\raggedright\arraybackslash}p{(\columnwidth - 6\tabcolsep) * \real{0.2195}}
  >{\raggedright\arraybackslash}p{(\columnwidth - 6\tabcolsep) * \real{0.2195}}
  >{\raggedright\arraybackslash}p{(\columnwidth - 6\tabcolsep) * \real{0.1951}}@{}}
\toprule\noalign{}
\begin{minipage}[b]{\linewidth}\raggedright
Heartbeat lost
\end{minipage} & \begin{minipage}[b]{\linewidth}\raggedright
Watcher
\end{minipage} & \begin{minipage}[b]{\linewidth}\raggedright
Timeout
\end{minipage} & \begin{minipage}[b]{\linewidth}\raggedright
Action
\end{minipage} \\
\midrule\noalign{}
\endhead
\bottomrule\noalign{}
\endlastfoot
SYS (0x7FE) & RT & 1000ms & CAN 0x001 → \textbf{ESTOP} (full brake, motor kill,
DCDC off) \\
Jetson (0x7FC) & RT & 1500ms & Zero setpoints → \textbf{controlled stop} (motor
0V, gear N, no active brake, steering silent) \\
\end{longtable}

The justification given (architecture §8.6): ``Jetson is QM, not
safety-critical. Its death triggers a controlled stop, not ESTOP. Three missed
frames protect against false triggers.''

But Jetson runs the \textbf{perception stack} --- obstacle detection, path
planning, collision avoidance. If Jetson dies while navigating traffic, the
vehicle continues coasting with: - No obstacle detection - No path planning - No
active braking (coast only) - No active steering (EPS-C timeout-faults)

A ``controlled stop'' in traffic is just coasting --- at 25 km/h,
\textgreater50m to stop on flat ground with no perception. The rider must
recognize the failure and manually brake.

\subsection{5.2 What the architecture actually
requires}\label{what-the-architecture-actually-requires}

The architecture's §6.2 Option D comparison table claims: ``SYS failure → brake
lost? No (RT takes over).'' This states RT should take over brake on any MCU
failure. But the heartbeat action table only triggers CAN 0x001 --- which SEB
ignores. RT does NOT take over 0x7B9 transmission on SYS heartbeat loss.

The option D architecture promises RT brake takeover on SYS failure but doesn't
implement it. And for Jetson failure, there's no brake takeover at all --- just
a coast.

\subsection{5.3 False-positive risk
assessment}\label{false-positive-risk-assessment}

The reason Jetson death isn't ESTOP is concern about false triggers: Jetson is a
Linux machine with non-realtime CAN stack. A momentary CAN buffer overflow or
scheduling delay could miss heartbeat frames.

But at 1500ms (3 missed frames at 2 Hz), a Linux machine that can't send a
single CAN frame in 1.5 seconds is genuinely dead --- not ``jittery.'' This is a
robust detection threshold. The false-positive risk from non-realtime Linux CAN
jitter at 1500ms is very low.

\subsection{5.4 Recommended fix: ``Assisted stop'' --- moderate brake on Jetson
loss}\label{recommended-fix-assisted-stop-moderate-brake-on-jetson-loss}

\textbf{Approach:} When RT detects Jetson heartbeat loss, instead of pure coast,
RT commands a moderate brake pressure via 0x205 (RT→SYS) and transitions SYS to
MANUAL mode. This is an intermediate response between ``controlled stop''
(coast) and ``ESTOP'' (full emergency brake with DCDC off).

\textbf{Behavior:}

\begin{longtable}[]{@{}
  >{\raggedright\arraybackslash}p{(\columnwidth - 4\tabcolsep) * \real{0.3667}}
  >{\raggedright\arraybackslash}p{(\columnwidth - 4\tabcolsep) * \real{0.3000}}
  >{\raggedright\arraybackslash}p{(\columnwidth - 4\tabcolsep) * \real{0.3333}}@{}}
\toprule\noalign{}
\begin{minipage}[b]{\linewidth}\raggedright
Condition
\end{minipage} & \begin{minipage}[b]{\linewidth}\raggedright
Current
\end{minipage} & \begin{minipage}[b]{\linewidth}\raggedright
Proposed
\end{minipage} \\
\midrule\noalign{}
\endhead
\bottomrule\noalign{}
\endlastfoot
Jetson heartbeat lost (1500ms) & Zero 0x204 + stop 0x169 → coast & Zero 0x204 +
stop 0x169 + \textbf{0x205 brake\_pressure\_kpa = 2000 (2 MPa,
\textasciitilde40\% of max)} + \textbf{SYS → MANUAL mode} \\
\end{longtable}

\textbf{Why 2000 kPa (2 MPa):} - SYNTREE SEB max pressure is 5 MPa (5000 kPa). -
2 MPa provides noticeable deceleration (\textasciitilde0.3--0.4g) without wheel
lockup. - This is roughly equivalent to moderate engine braking in a car --- the
vehicle slows decisively but not violently. - The rider can always override by
pressing the brake lever harder (lever pressure wins in the max-select brake
arbitration).

\textbf{Implementation path:} 1. RT detects Jetson heartbeat timeout. 2. RT sets
\texttt{0x205\ RT\_BRAKE\_CMD} = \texttt{kBrakeAssistedStopKpa} (2000 kPa). 3.
RT sends \texttt{0x110} with mode=MANUAL → SYS transitions to MANUAL. 4. SYS
receives \texttt{0x205\ \textgreater{}\ 0} → SEB Pressure Mode, target 2 MPa. 5.
Brake light ON (OR logic: mode change or 0x302 bit). 6. Motor throttle = 0V
(from 0x204 speed=0). 7. Gear = N (from 0x204 gear=N). 8. Rider takes over ---
lever overrides, MODE button available to re-enter AUTO after Jetson recovers.

\textbf{Why this is better than full ESTOP:} - DC-DC stays on → 12V rail stays
up → signal lights, brake light, indicators all functional. - Rider maintains
steering (EPS-C standalone in MANUAL) and can override brake. - No power-cycle
needed to recover --- just wait for Jetson to reboot and press MODE. - Safer in
traffic than a dark, unbraked vehicle (ESTOP kills all lights).

\textbf{Why this is better than pure coast:} - Active deceleration from the
moment of detection. - Brake light warns following vehicles. - Transition to
MANUAL gives the rider immediate control. - The 1500ms detection window +
\textasciitilde50ms to first brake pressure = \textasciitilde1.55s from Jetson
death to active braking. At 25 km/h, 10.8m of coast before brake engages ---
much better than \textgreater50m of pure coast.

\begin{center}\rule{0.5\linewidth}{0.5pt}\end{center}

\section{Issue 6: START Button Failure --- No Monitored Fallback for the Only
ESTOP
Exit}\label{issue-6-start-button-failure-no-monitored-fallback-for-the-only-estop-exit}

\subsection{6.1 The single point of failure}\label{the-single-point-of-failure}

The START button (GPIO32) is the \textbf{only} exit from ESTOP. The architecture
states:

\begin{quote}
``ESTOP exit: START button → MANUAL mode. (NOT MODE button, NOT CAN command, NOT
timeout)''
\end{quote}

There is no monitoring of GPIO32 health. Failure modes: - \textbf{Stuck HIGH:}
Button disconnected or internal pull-up failed → can never exit ESTOP. -
\textbf{Stuck LOW:} Button shorted to ground → could accidentally exit ESTOP
(though a falling edge transition is required, so a persistent LOW at power-on
wouldn't trigger exit). - \textbf{Debounce failure:} Noisy contact → multiple
rapid mode transitions → undefined behavior.

The only backup is power-cycle, which: 1. Restarts all four nodes (Jetson, RT,
SYS, MTR). 2. Requires SEB LBS boot sync (\textasciitilde2.5s until brake
available). 3. Requires EPS-C LBS boot sync (\textasciitilde2.5s until steering
sync). 4. May fail SEB sync (→ BRAKE\_DEGRADED, brake lever still works with the
fix from Issue 2). 5. Is impractical at roadside in traffic.

\subsection{6.2 Causal trace --- stuck START
button}\label{causal-trace-stuck-start-button}

\begin{verbatim}
Rider presses ESTOP (legitimate, e.g., obstacle)
  → Vehicle enters ESTOP: motor kill, full brake, DCDC off
  → Rider assesses: obstacle cleared, safe to proceed
  → Rider presses START button
  → GPIO32 reads HIGH (stuck — broken wire, bad solder joint, corroded contact)
  → mode_task: no falling edge detected
  → Mode remains ESTOP
  → Rider presses START repeatedly — no response
  → Rider's only option: power-cycle
  → All nodes reboot → 3+ seconds before drivable
  → If SEB sync fails during boot → BRAKE_DEGRADED
  → Roadside recovery takes 10+ seconds
\end{verbatim}

\subsection{6.3 Recommended fix: MODE button long-press as secondary ESTOP
exit}\label{recommended-fix-mode-button-long-press-as-secondary-estop-exit}

\textbf{Approach:} Add a secondary ESTOP exit path via the MODE button (GPIO11).
A long-press (3 seconds) of the MODE button exits ESTOP to MANUAL mode,
identical to a START button press. The MODE button is on a different GPIO ---
two independent physical buttons must fail for the rider to be completely locked
out.

\textbf{Implementation in SYS \texttt{mode\_task}:}

\begin{verbatim}
// Existing: START button (GPIO32) — exit ESTOP
if (falling_edge(gpio32) && g_mode == Mode::Estop) {
    mode_set(MANUAL);  // immediate exit
}

// NEW: MODE button (GPIO11) long-press — secondary ESTOP exit
if (gpio11_is_low) {
    gpio11_hold_ms += 10;  // mode_task runs at 10 Hz = 100ms per tick
    if (gpio11_hold_ms >= 3000 && g_mode == Mode::Estop) {
        mode_set(MANUAL);  // long-press exit
        gpio11_hold_ms = 0;
    }
} else {
    gpio11_hold_ms = 0;
}
\end{verbatim}

\textbf{START button health monitoring (additional):}

\begin{verbatim}
// In diag_task @ 1 Hz:
if (g_mode == Mode::Estop && estop_duration_ms > 30000) {
    // Been in ESTOP for >30s — check if START button is responsive
    // Log diagnostic. If button is stuck HIGH for the entire duration,
    // set a flag in 0x600 SYS_DIAG_RPT.
}
\end{verbatim}

\textbf{Why MODE button long-press, not short-press:} - In MANUAL/AUTO, MODE
short-press toggles modes --- muscle memory. - In ESTOP, a long-press prevents
accidental exit from bumping the MODE button. - A 3-second hold is a deliberate
action --- the rider is consciously choosing to exit ESTOP. - Different duration
from steering retry (START short-press for STEER\_FAULT vs START long-press 3s
for force-activate).

\begin{center}\rule{0.5\linewidth}{0.5pt}\end{center}

\section{Issue 7: SYS Crash Has a 1000ms Brake Gap --- RT Should Take Over 0x7B9
Immediately}\label{issue-7-sys-crash-has-a-1000ms-brake-gap-rt-should-take-over-0x7b9-immediately}

\subsection{7.1 The gap}\label{the-gap}

When SYS crashes:

\begin{longtable}[]{@{}
  >{\raggedright\arraybackslash}p{(\columnwidth - 4\tabcolsep) * \real{0.2308}}
  >{\raggedright\arraybackslash}p{(\columnwidth - 4\tabcolsep) * \real{0.2692}}
  >{\raggedright\arraybackslash}p{(\columnwidth - 4\tabcolsep) * \real{0.5000}}@{}}
\toprule\noalign{}
\begin{minipage}[b]{\linewidth}\raggedright
Time
\end{minipage} & \begin{minipage}[b]{\linewidth}\raggedright
Event
\end{minipage} & \begin{minipage}[b]{\linewidth}\raggedright
Brake state
\end{minipage} \\
\midrule\noalign{}
\endhead
\bottomrule\noalign{}
\endlastfoot
T+0 & SYS hangs. Last 0x7B9 frame sent at T−20ms (50 Hz). & Last commanded
stroke active \\
T+20ms & SEB: no 0x7B9 for 20ms → \textbf{comm-fault timeout.} Behavior unknown.
& \textbf{Brake lost} (or held, unverified) \\
T+1000ms & RT: SYS heartbeat (0x7FE) frozen for 2 missed frames → timeout. RT
broadcasts CAN 0x001. & SEB ignores CAN 0x001 (only responds to 0x7B9) \\
T+1000ms+ & MTR receives CAN 0x001 → kills motor + gear. & Motor safe, brake
still lost \\
\end{longtable}

\textbf{Brake gap: 20ms to 1000ms = \textasciitilde980ms.} At 25 km/h, this is
\textasciitilde6.8m of travel with no brake.

\subsection{7.2 The architecture already promises this
fix}\label{the-architecture-already-promises-this-fix}

Architecture §6.2 Option D explicitly claims:

\begin{quote}
``SYS failure → brake lost? \textbf{No (RT takes over)}''
\end{quote}

And:

\begin{quote}
``D preserves independent failure modes --- RT failure loses AUTO steering/brake
but SYS can still brake in MANUAL; \textbf{SYS failure loses MANUAL brake but RT
can still brake in AUTO.}''
\end{quote}

But the current heartbeat action table does NOT implement RT brake takeover on
SYS heartbeat loss. It only broadcasts CAN 0x001. RT already knows how to send
0x7B9 (it does so in AUTO mode at 50 Hz with pressure control). It simply
doesn't do so on SYS failure.

\subsection{7.3 Why the 1000ms timeout is wrong for the brake
path}\label{why-the-1000ms-timeout-is-wrong-for-the-brake-path}

Architecture §8.6 justifies the 1000ms inter-MCU heartbeat timeout:

\begin{quote}
``The fast detection path for actuator faults is the steering following-error
check (300ms → ESTOP) and 0x204 staleness (200ms, see below). The heartbeat
catches failures those checks miss.''
\end{quote}

This lists TWO fast detection paths: 1. Steering following error (300ms) ---
covers EPS-C 2. 0x204 staleness (200ms) --- covers motor

There is \textbf{no fast detection path for the SEB brake actuator.} The brake
relies entirely on the 1000ms heartbeat timeout. But the brake is the most
critical actuator --- motor kill without brake just means coasting. Brake
failure means no way to stop.

The FTTI (Fault Tolerant Time Interval) for brake loss is shorter than for motor
or steering loss. The vehicle must be able to stop within a bounded distance. At
25 km/h: - 200ms → 1.4m (acceptable) - 1000ms → 6.9m (marginal) - 2500ms → 17.4m
(unacceptable --- Issue 3's watchdog window)

\subsection{7.4 Recommended fix: RT brake takeover on SYS heartbeat loss, with
shortened
timeout}\label{recommended-fix-rt-brake-takeover-on-sys-heartbeat-loss-with-shortened-timeout}

\textbf{Approach:} Two changes:

\begin{enumerate}
\def\labelenumi{\arabic{enumi}.}
\tightlist
\item
  \textbf{Shorten RT's SYS heartbeat timeout from 1000ms to 200ms.} This is the
  FTTI-bound value (1.4m at 25 km/h). At 2 Hz heartbeat, this means detecting
  after 1 missed frame (200ms is \textgreater500ms period --- wait, at 2 Hz
  period=500ms, 1 missed frame = 500ms worst case, 2 missed = 1000ms). We need a
  faster heartbeat or a different detection mechanism.
\end{enumerate}

\textbf{Revised heartbeat rate for SYS:} Increase SYS heartbeat from 2 Hz to
\textbf{10 Hz} (100ms period). Then a timeout of 200ms = 2 missed frames. This
requires: - SYS \texttt{hb} task: change from 2 Hz to 10 Hz (trivial, 10 Hz is
well within capability). - Bandwidth: 1 byte DLC at 10 Hz = negligible
(\textless0.02\% of 500 kbit/s bus).

\begin{enumerate}
\def\labelenumi{\arabic{enumi}.}
\setcounter{enumi}{1}
\tightlist
\item
  \textbf{On SYS heartbeat timeout, RT immediately takes over 0x7B9 transmission
  with full brake.} RT already sends 0x7B9 in AUTO mode. On SYS heartbeat loss,
  RT sends 0x7B9 with stroke=max regardless of mode (mode gate opens on
  emergency). This is the dual-sender exception documented in §6.2 --- both RT
  and SYS may briefly transmit if SYS isn't truly dead, but SEB accepts
  whichever has valid checksum + rolling counter.
\end{enumerate}

\textbf{New behavior timeline:}

\begin{longtable}[]{@{}
  >{\raggedright\arraybackslash}p{(\columnwidth - 4\tabcolsep) * \real{0.2308}}
  >{\raggedright\arraybackslash}p{(\columnwidth - 4\tabcolsep) * \real{0.2692}}
  >{\raggedright\arraybackslash}p{(\columnwidth - 4\tabcolsep) * \real{0.5000}}@{}}
\toprule\noalign{}
\begin{minipage}[b]{\linewidth}\raggedright
Time
\end{minipage} & \begin{minipage}[b]{\linewidth}\raggedright
Event
\end{minipage} & \begin{minipage}[b]{\linewidth}\raggedright
Brake state
\end{minipage} \\
\midrule\noalign{}
\endhead
\bottomrule\noalign{}
\endlastfoot
T+0 & SYS hangs. Last 0x7B9 sent at T−20ms. & Last commanded stroke active \\
T+20ms & SEB comm-fault timeout (if SYS was the last sender). & Behavior
unknown \\
T+200ms & RT: SYS heartbeat (0x7FE, now at 10 Hz) frozen for 2 missed frames →
timeout. & --- \\
T+200ms & RT immediately starts sending 0x7B9 @ 50 Hz with stroke=max. & --- \\
T+220ms & SEB receives first RT 0x7B9 frame. & \textbf{Brake engages.} \\
\textbf{Brake gap} & \textbf{20ms to 220ms = 200ms worst case} & \textbf{1.4m at
25 km/h} \\
\end{longtable}

\textbf{Why RT can safely take over 0x7B9:}

RT already sends 0x7B9 in AUTO mode with the full SYNTREE protocol (rolling
counter, checksum, stroke/pressure control). Taking over in MANUAL/ESTOP on
emergency is just continuing to send with different parameters (stroke=max
instead of computed pressure). No new protocol logic needed. The SYNTREE
security bytes (rolling counter + checksum) prove liveness --- SEB accepts the
frame regardless of which node sent it.

\textbf{Collision risk during takeover (both RT and SYS sending 0x7B9):}

If SYS isn't fully dead (heartbeat was delayed but SYS is still transmitting
0x7B9), both RT and SYS temporarily send 0x7B9 on the low bus. Each sends at 50
Hz (20ms period). Both have valid checksums and rolling counters. SEB receives
two alternating streams with different counters --- this is within the CAN spec
(arbitration is per-frame) and SEB's acceptance criteria (valid checksum +
counter). SYS's frames have lever-based stroke (0--15mm); RT's frames have max
stroke (27mm). The max-select arbitration at the SEB level (it follows whichever
command has the higher stroke) would pick RT's 27mm --- which is the safe choice
(full brake).

\textbf{Return to normal operation:} When SYS heartbeat resumes (SYS reboots,
sends 0x7FE), RT stops sending 0x7B9 and SYS resumes normal mode-gated
operation. The transition is: RT stops → brief gap (≤20ms) → SYS starts. SEB
sees a 20ms gap from RT's last frame before SYS's first frame, which is within
the 20ms comm-fault threshold (marginal --- may need a brief overlap where both
send to avoid SEB timeout).

\textbf{Interaction with Issue 3 (watchdog brake window):}

With this fix, the watchdog brake window during SYS reset shrinks from
\textasciitilde2.5s to \textasciitilde200ms: RT detects SYS heartbeat loss
within 200ms and fills the brake gap with 0x7B9. SEB receives RT's brake command
while SYS reboots. When SYS comes back online and resumes 0x7B9 transmission, RT
detects SYS heartbeat return and yields. Total unbraked window: 200ms heartbeat
detection + 20ms first frame = 220ms worst case.

The SEB comm-fault behavior test (Issue 3.6) is still needed to confirm the 20ms
gap behavior, but the 2.5s window is eliminated.

\begin{center}\rule{0.5\linewidth}{0.5pt}\end{center}

\section{Updated Summary}\label{updated-summary}

\begin{longtable}[]{@{}
  >{\raggedright\arraybackslash}p{(\columnwidth - 8\tabcolsep) * \real{0.1061}}
  >{\raggedright\arraybackslash}p{(\columnwidth - 8\tabcolsep) * \real{0.1515}}
  >{\raggedright\arraybackslash}p{(\columnwidth - 8\tabcolsep) * \real{0.1970}}
  >{\raggedright\arraybackslash}p{(\columnwidth - 8\tabcolsep) * \real{0.2424}}
  >{\raggedright\arraybackslash}p{(\columnwidth - 8\tabcolsep) * \real{0.3030}}@{}}
\toprule\noalign{}
\begin{minipage}[b]{\linewidth}\raggedright
Issue
\end{minipage} & \begin{minipage}[b]{\linewidth}\raggedright
Severity
\end{minipage} & \begin{minipage}[b]{\linewidth}\raggedright
Probability
\end{minipage} & \begin{minipage}[b]{\linewidth}\raggedright
Fix complexity
\end{minipage} & \begin{minipage}[b]{\linewidth}\raggedright
Recommended action
\end{minipage} \\
\midrule\noalign{}
\endhead
\bottomrule\noalign{}
\endlastfoot
\textbf{1. ESTOP exit race} & High & Medium & Low (software, RT-local) & Defer
mode transition until ramp complete \\
\textbf{2. SEB brake lever contradiction} & High & Low & Low (software,
SYS-local) & Replace BRAKE\_FAULT with BRAKE\_DEGRADED \\
\textbf{3. Watchdog unbraked window} & High & Low→Very Low (with Issue 7 fix) &
Low (Issue 7 fix covers this) & Issue 7's RT brake takeover closes most of this
gap. Verify SEB comm-fault behavior for the 20ms edge. \\
\textbf{4. Obstacle ESTOP cornering rollover} & High & Low & Low (software,
RT-local) & Dynamic-clamp the hold angle during obstacle ESTOP \\
\textbf{5. Jetson death → coast is insufficient} & Medium & Low & Low (software,
RT-local) & Assisted stop: 2 MPa brake + MANUAL transition on Jetson heartbeat
loss \\
\textbf{6. START button failure --- no fallback} & Medium & Low & Low (software,
SYS-local) & MODE button long-press (3s) as secondary ESTOP exit \\
\textbf{7. SYS crash → 1000ms brake gap} & High & Low & Medium (SYS heartbeat
rate change + RT takeover logic) & Shorten SYS heartbeat to 10 Hz/200ms timeout;
RT takes over 0x7B9 on loss \\
\end{longtable}

\textbf{Issues 4, 5, and 6 are software-only.} Issue 7 requires a heartbeat rate
change on SYS and brake takeover logic on RT --- still software-only, but
touches both MCUs. \textbf{Issue 7 also significantly mitigates Issue 3} by
reducing the watchdog brake window from \textasciitilde2.5s to
\textasciitilde200ms.

\begin{center}\rule{0.5\linewidth}{0.5pt}\end{center}

\section{Issue 8: No Independent Brake Monitor --- ESTOP Could Silently Have No
Brakes}\label{issue-8-no-independent-brake-monitor-estop-could-silently-have-no-brakes}

\subsection{8.1 The gap}\label{the-gap-1}

All 8 safety layers protect motor speed (MCP4725=0V, gear=N, 0x204 staleness
check) and steering (following error, dynamic clamp, hard-stops), but
\textbf{none independently verify that the SEB actually applied braking force.}
During ESTOP, SYS (or RT, with the Issue 7 fix) transmits \texttt{0x7B9} with
stroke=max at 50 Hz. If SEB's CAN receiver is faulted, the system has no way to
know the brake command was never received or acted upon.

A stale \texttt{0x7B9} TX with no acknowledgment monitoring means ESTOP could
silently fail to brake --- the motor stops, but the vehicle coasts rather than
decelerating under brake force.

\subsection{8.2 What the SEB already provides (no new sensors
needed)}\label{what-the-seb-already-provides-no-new-sensors-needed}

Per \texttt{docs/by-wire\ -\ brake.csv}, the SEB\_STATUS frame (\texttt{0x721},
100 Hz = 10ms cycle) already contains:

\begin{longtable}[]{@{}
  >{\raggedright\arraybackslash}p{(\columnwidth - 8\tabcolsep) * \real{0.1951}}
  >{\raggedright\arraybackslash}p{(\columnwidth - 8\tabcolsep) * \real{0.1707}}
  >{\raggedright\arraybackslash}p{(\columnwidth - 8\tabcolsep) * \real{0.1463}}
  >{\raggedright\arraybackslash}p{(\columnwidth - 8\tabcolsep) * \real{0.1707}}
  >{\raggedright\arraybackslash}p{(\columnwidth - 8\tabcolsep) * \real{0.3171}}@{}}
\toprule\noalign{}
\begin{minipage}[b]{\linewidth}\raggedright
Signal
\end{minipage} & \begin{minipage}[b]{\linewidth}\raggedright
Field
\end{minipage} & \begin{minipage}[b]{\linewidth}\raggedright
Type
\end{minipage} & \begin{minipage}[b]{\linewidth}\raggedright
Range
\end{minipage} & \begin{minipage}[b]{\linewidth}\raggedright
Description
\end{minipage} \\
\midrule\noalign{}
\endhead
\bottomrule\noalign{}
\endlastfoot
\texttt{SEB\_Stroke\_Value} & Bytes 2-3 & uint16, scale 0.05, offset -30 & -5 to
27 mm & \textbf{Actual stroke position} \\
\texttt{SEB\_Pressure\_Value} & Byte 4 & uint8, scale 0.05 & 0 to 5 MPa &
\textbf{Actual hydraulic pressure} \\
\texttt{SEB\_Error\_Status} & Byte 0, bits 6-7 & uint2 & 0-3 & 0=no fault,
1=minor, 2=general, 3=\textbf{severe (request shutdown)} \\
\texttt{SEB\_RollCnt\_Status} & Byte 6, bits 4-7 & uint4 & 0-15 & Rolling
counter --- proves SEB is alive \\
\end{longtable}

Additionally, \texttt{0x731\ SEB\_ErrInfo} (10 Hz) provides detailed fault flags
including \texttt{SEB\_CanCom\_Err} (CAN Communication Fault, Level 3) --- the
SEB itself reports when it has lost CAN communication.

\textbf{This is sufficient to implement a brake following-error monitor entirely
in software, with zero new sensors.}

\subsection{8.3 Causal trace --- what a brake monitor
catches}\label{causal-trace-what-a-brake-monitor-catches}

\begin{verbatim}
ESTOP triggered (any source)
  → SYS/RT commands 0x7B9 {Stroke_Value_Req = 1140 (27mm max)}
  → SEB receives frame, validates checksum + rolling counter
     Case A: SEB healthy → stroke actuates to 27mm → 0x721 {Stroke_Value ≈ 1140}
             Monitor: |1140 - 1140| < threshold → brake OK
     Case B: SEB CAN receiver faulted → never receives 0x7B9
             → 0x721 {Stroke_Value = 0 (or last held value)}
             Monitor: |0 - 1140| >> threshold → BRAKE FAULT DETECTED
     Case C: SEB internal fault (motor stall, oil pressure loss)
             → 0x721 {SEB_Error_Status = 3 (severe), Stroke_Value stuck}
             → or 0x731 {SEB_Mtr_Stall_Err = 1, SEB_Oil_Err = 1}
             Monitor: error status or stroke mismatch → FAULT DETECTED
     Case D: SEB mechanically jammed (caliper stuck, hydraulic blockage)
             → 0x721 {Stroke_Value frozen, Pressure_Value rising abnormally}
             Monitor: stroke not tracking command → FAULT DETECTED
\end{verbatim}

\subsection{8.4 Recommended fix: Brake following-error monitor in
SYS}\label{recommended-fix-brake-following-error-monitor-in-sys}

\textbf{Approach:} Add a brake following-error check to SYS's \texttt{dispatch}
task (processing \texttt{0x721}) or \texttt{safety} task, analogous to the
steering following-error check on RT. Since we cannot add physical sensors, this
uses the existing SEB CAN feedback --- which is already being received by SYS.

\textbf{Implementation:}

\begin{verbatim}
// In SYS dispatch task, on receiving 0x721 SEB_STATUS:
void check_brake_following_error() {
    float cmd_stroke = last_commanded_stroke_mm;  // from brake_task's 0x7B9 TX
    float actual_stroke = seb_status.stroke_value_mm();  // from 0x721
    float stroke_error = fabs(cmd_stroke - actual_stroke);
    uint8_t seb_error_level = seb_status.error_status();  // bits 6-7

    // Check 1: Severe SEB internal fault
    if (seb_error_level >= 3) {
        log_brake_fault("SEB Level 3 fault — request shutdown");
        g_brake_monitor_fault = true;
    }

    // Check 2: Stroke following error (debounced)
    if (stroke_error > kBrakeStrokeErrorThresholdMm && cmd_stroke > kBrakeMinCmdForCheck) {
        brake_error_duration_ms += 10;  // 0x721 arrives at 100 Hz = 10ms period
        if (brake_error_duration_ms > kBrakeErrorDebounceMs) {
            log_brake_fault("Brake stroke following error: cmd=%.1f act=%.1f",
                            cmd_stroke, actual_stroke);
            g_brake_monitor_fault = true;
        }
    } else {
        brake_error_duration_ms = 0;
    }

    // Check 3: 0x721 staleness (no feedback at all)
    // If no 0x721 received for >100ms (10 missed frames at 100 Hz):
    // SEB is not communicating → brake may be dead
}
\end{verbatim}

\textbf{Constants:}

\begin{longtable}[]{@{}
  >{\raggedright\arraybackslash}p{(\columnwidth - 4\tabcolsep) * \real{0.3571}}
  >{\raggedright\arraybackslash}p{(\columnwidth - 4\tabcolsep) * \real{0.2500}}
  >{\raggedright\arraybackslash}p{(\columnwidth - 4\tabcolsep) * \real{0.3929}}@{}}
\toprule\noalign{}
\begin{minipage}[b]{\linewidth}\raggedright
Constant
\end{minipage} & \begin{minipage}[b]{\linewidth}\raggedright
Value
\end{minipage} & \begin{minipage}[b]{\linewidth}\raggedright
Rationale
\end{minipage} \\
\midrule\noalign{}
\endhead
\bottomrule\noalign{}
\endlastfoot
\texttt{kBrakeStrokeErrorThresholdMm} & 3.0 mm & \textasciitilde11\% of max
stroke (27mm). Allows for mechanical compliance but catches gross failures. \\
\texttt{kBrakeMinCmdForCheck} & 1.0 mm & Don't check when brake is released
(stroke=0) --- noise dominates. \\
\texttt{kBrakeErrorDebounceMs} & 100 ms & 10 consecutive frames at 100 Hz.
Prevents false triggers from hydraulic transients. \\
\texttt{0x721\ staleness\ timeout} & 100 ms & 10 missed frames at 100 Hz = SEB
CAN dead. \\
\end{longtable}

\textbf{What the monitor DOES when a fault is detected:}

Since ESTOP is already the maximum response (motor killed, brake commanded to
max), the brake monitor cannot trigger ``more ESTOP.'' But it CAN:

\begin{enumerate}
\def\labelenumi{\arabic{enumi}.}
\tightlist
\item
  \textbf{Set a persistent diagnostic flag} in CAN
  \texttt{0x600\ SYS\_DIAG\_RPT} --- records that ESTOP braking failed for
  post-incident analysis.
\item
  \textbf{Log the fault} to persistent storage (if available).
\item
  \textbf{Flash the brake light} in a distinct pattern (if 12V is available ---
  see Issue 12) to alert the rider that braking may be compromised.
\item
  \textbf{If the fault occurs during non-ESTOP operation} (lever braking in
  MANUAL): log the fault so it's caught early.
\end{enumerate}

\textbf{Why this is not identical to the steering following-error check:}

The steering following error triggers ESTOP --- it escalates the response. The
brake following error CANNOT escalate (ESTOP is already max brake). It serves a
different purpose: \textbf{fault recording and rider alerting.} Knowing that the
brake failed during an incident is critical for accident investigation and
system improvement.

\subsection{8.5 SEB\_Error\_Status as a primary
signal}\label{seb_error_status-as-a-primary-signal}

The \texttt{SEB\_Error\_Status} field in \texttt{0x721} is particularly valuable
because the SEB itself reports when it cannot function:

\begin{longtable}[]{@{}
  >{\raggedright\arraybackslash}p{(\columnwidth - 4\tabcolsep) * \real{0.2917}}
  >{\raggedright\arraybackslash}p{(\columnwidth - 4\tabcolsep) * \real{0.3750}}
  >{\raggedright\arraybackslash}p{(\columnwidth - 4\tabcolsep) * \real{0.3333}}@{}}
\toprule\noalign{}
\begin{minipage}[b]{\linewidth}\raggedright
Value
\end{minipage} & \begin{minipage}[b]{\linewidth}\raggedright
Meaning
\end{minipage} & \begin{minipage}[b]{\linewidth}\raggedright
Action
\end{minipage} \\
\midrule\noalign{}
\endhead
\bottomrule\noalign{}
\endlastfoot
0 & No fault & Normal operation \\
1 & Level 1 (minor, warning) & Log, continue \\
2 & Level 2 (general, warning) & Log, set diag flag \\
3 & \textbf{Level 3 (severe, request shutdown)} & \textbf{Immediate: set brake
fault flag, log, alert rider} \\
\end{longtable}

A Level 3 fault from SEB means the SEB itself has determined it cannot safely
operate. This is a manufacturer-defined condition and carries more weight than a
stroke comparison.

\begin{center}\rule{0.5\linewidth}{0.5pt}\end{center}

\section{Issue 9: CAN 0x001 Is Spoofable --- No Authentication,
DoS-Vulnerable}\label{issue-9-can-0x001-is-spoofable-no-authentication-dos-vulnerable}

\subsection{9.1 The gap}\label{the-gap-2}

CAN \texttt{0x001\ SAFETY\_ESTOP} is an empty frame (DLC=0) with the highest
arbitration priority. Any node on either bus can broadcast it. There is no
authentication --- no checksum, no sender identifier, no rolling counter. A
corrupted node that enters an error loop could flood \texttt{0x001}, causing:

\begin{enumerate}
\def\labelenumi{\arabic{enumi}.}
\tightlist
\item
  \textbf{Persistent ESTOP} --- every \texttt{0x001} frame re-triggers ESTOP on
  all receivers.
\item
  \textbf{CAN bus saturation} --- at 500 kbit/s, a minimum CAN frame is
  \textasciitilde47 bits (standard ID + DLC=0 + stuff bits + EOF). A babbling
  node could transmit thousands per second, consuming all bus bandwidth.
\item
  \textbf{No recovery except power-cycle} --- the flooding node must be
  physically removed from the bus.
\end{enumerate}

\subsection{9.2 CAN error containment
analysis}\label{can-error-containment-analysis}

CAN controllers have built-in error containment: a node that fails repeatedly
increments its Transmit Error Counter (TEC). At TEC \textgreater{} 255, the node
enters \textbf{bus-off} state and disconnects from the bus. This means a
flooding node WILL eventually be isolated --- but:

\begin{itemize}
\tightlist
\item
  The error recovery mechanism requires the node to successfully transmit before
  TEC decreases. A genuinely corrupted node in a tight loop may never recover.
\item
  During the flood period (before bus-off), all other traffic is blocked by
  \texttt{0x001} winning every arbitration.
\item
  If the flooding node is on the high bus (Jetson), RT bridges \texttt{0x001} to
  the low bus --- doubling the disruption.
\end{itemize}

\subsection{9.3 Recommended fix: Rate-limit ESTOP processing + sender
identification}\label{recommended-fix-rate-limit-estop-processing-sender-identification}

\textbf{Approach:} Two independent protections:

\textbf{Protection 1 --- Rate limit on ESTOP processing (software, all nodes):}

\begin{verbatim}
// In each node's CAN dispatch, on receiving 0x001:
static int64_t last_estop_frame_us = 0;
static int estop_count_in_window = 0;
constexpr int64_t kEstopRateWindowUs = 500'000;  // 500ms window
constexpr int kMaxEstopPerWindow = 2;             // max 2 per window

int64_t now = esp_timer_get_time();
if (now - last_estop_frame_us > kEstopRateWindowUs) {
    estop_count_in_window = 0;
    last_estop_frame_us = now;
}
estop_count_in_window++;

if (estop_count_in_window <= kMaxEstopPerWindow) {
    mode_set(Estop);  // process normally
} else {
    // Rate limit exceeded — ignore this frame
    // The first frame already triggered ESTOP; subsequent floods are noise
    log_warning("ESTOP flood detected (%d frames in %lld us)",
                estop_count_in_window, now - last_estop_frame_us);
}
\end{verbatim}

This means: the FIRST \texttt{0x001} frame triggers ESTOP (as intended).
Subsequent floods within the same 500ms window are ignored. A genuine second
ESTOP (e.g., physical button pressed after recovering from the first) would
arrive after the window expires.

\textbf{Protection 2 --- Add sender identification to ESTOP frames (protocol
change):}

Change \texttt{0x001\ SAFETY\_ESTOP} from DLC=0 to DLC=1 with a single byte
identifying the sender:

\begin{longtable}[]{@{}lll@{}}
\toprule\noalign{}
Byte & Field & Values \\
\midrule\noalign{}
\endhead
\bottomrule\noalign{}
\endlastfoot
0 & Sender ID & 0x00 = SYS, 0x01 = RT, 0x02 = Jetson, 0x03 = MTR \\
\end{longtable}

This allows receivers to: - Log which node triggered ESTOP (valuable for
diagnostics). - Detect duplicate frames from the same sender (ignore if the same
sender floods). - Apply per-sender rate limits (one node flooding doesn't block
other nodes' legitimate ESTOP).

The DLC=1 frame is 55 bits (vs 47 for DLC=0) --- the arbitration time difference
is negligible at 500 kbit/s (\textasciitilde16 µs). The safety benefit of
knowing who triggered ESTOP justifies the 8 extra bits.

\textbf{Recovery from ESTOP flood:} The rate limiter means the rider can press
START to exit ESTOP even during a flood --- because \texttt{0x001} frames
arriving during the rate-limit window are ignored. The flood source must still
be resolved (the corrupted node eventually bus-offs or is physically
disconnected), but the vehicle can be recovered to MANUAL mode.

\begin{center}\rule{0.5\linewidth}{0.5pt}\end{center}

\section{Issue 10: No ESTOP Acknowledgment from MTR
STM32}\label{issue-10-no-estop-acknowledgment-from-mtr-stm32}

\subsection{10.1 The gap}\label{the-gap-3}

When ESTOP is triggered (button, CAN 0x001, or heartbeat), SYS and RT have no
confirmation that MTR actually received and acted on it. MTR cuts throttle and
gear locally --- but this action is invisible to SYS/RT. The SYS function
monitor (EGAS Level 2) compares \texttt{0x204} setpoint vs \texttt{0x206}
feedback. But during ESTOP, \texttt{0x204} is forced to zero, so:

\begin{itemize}
\tightlist
\item
  If MTR was already at zero speed (vehicle stopped), a frozen MTR outputting
  zero would not generate a mismatch --- the monitor sees
  \texttt{0x204=0,\ 0x206=0} and assumes MTR is responding correctly.
\item
  If MTR freezes entirely (no \texttt{0x206} transmission), SYS has no
  \texttt{0x206} staleness check --- the absence of feedback is not monitored.
\item
  MTR has no dedicated heartbeat CAN ID. It monitors RT heartbeat
  (\texttt{0x7FD}) but doesn't send one of its own.
\end{itemize}

\subsection{10.2 What MTR already sends}\label{what-mtr-already-sends}

Per MTR config (\texttt{mtr-stm32/src/config.h}):

\begin{verbatim}
kIdSysThrottleSts = 0x120;  // MTR→RT/SYS, 100 Hz — throttle speed
kIdMotorFbk       = 0x206;  // MTR→SYS/RT, 50 Hz  — speed, gear state, fault flags
\end{verbatim}

\texttt{0x206\ MTR\_MOTOR\_FBK} at 50 Hz contains
\texttt{i16\ actual\_speed,\ u8\ gear\_state,\ u8\ fault\_flags}. The
\texttt{fault\_flags} field is the mechanism for MTR to report its internal
state.

\subsection{10.3 Recommended fix: ESTOP acknowledgment via 0x206
fault\_flags}\label{recommended-fix-estop-acknowledgment-via-0x206-fault_flags}

\textbf{Approach:} MTR sets a dedicated \texttt{ESTOP\_ACTIVE} bit in the
\texttt{0x206\ fault\_flags} byte when it has locally acted on ESTOP (throttle
DAC=0, all gear relays OFF). This is a software change on MTR.

\textbf{Implementation:}

\begin{verbatim}
// In MTR fault_flags bit definitions:
constexpr uint8_t kMtrFaultEstopActive = 0x01;  // bit 0: ESTOP active locally
constexpr uint8_t kMtrFaultDacFault    = 0x02;  // bit 1: MCP4725 I2C fault
constexpr uint8_t kMtrFaultGearFault   = 0x04;  // bit 2: gear relay fault
// ... other fault bits ...

// In MTR main loop, after ESTOP action:
if (estop_asserted) {
    dac_write(0);           // MCP4725 = 0V
    gear_all_off();         // all relays OFF
    fault_flags |= kMtrFaultEstopActive;  // ACKNOWLEDGE
} else {
    fault_flags &= ~kMtrFaultEstopActive;
}
// fault_flags is transmitted in 0x206 at 50 Hz
\end{verbatim}

\textbf{On SYS side (EGAS Level 2 monitor):}

\begin{verbatim}
// In SYS dispatch, processing 0x206:
if (g_mode == Mode::Estop) {
    if (!(mtr_fbk.fault_flags & kMtrFaultEstopActive)) {
        // MTR has NOT acknowledged ESTOP
        // This could mean: MTR didn't receive ESTOP, MTR CAN is dead, or MTR is frozen
        mtr_estop_ack_timeout_ms += 20;  // 0x206 at 50 Hz = 20ms period
        if (mtr_estop_ack_timeout_ms > 100) {  // 5 missed frames
            log_error("MTR failed to acknowledge ESTOP");
            g_mtr_estop_fault = true;
            // MTR is already supposed to kill motor — if it hasn't acknowledged,
            // something is wrong. Set diag flag for post-incident analysis.
        }
    } else {
        mtr_estop_ack_timeout_ms = 0;
        g_mtr_estop_fault = false;
    }
}

// Also: 0x206 staleness check (new)
if (time_since_last_0x206 > 200) {  // 10 missed frames at 50 Hz
    log_error("MTR feedback lost — 0x206 absent for >200ms");
    g_mtr_comms_lost = true;
}
\end{verbatim}

\textbf{What this gives us:} 1. Positive confirmation that MTR received and
acted on ESTOP (via \texttt{ESTOP\_ACTIVE} bit). 2. Detection of MTR
communication loss (via \texttt{0x206} staleness check). 3. The EGAS Level 2
monitor can now distinguish: ``MTR is healthy and braking'' vs ``MTR is silent
and possibly still powering the motor.''

\begin{center}\rule{0.5\linewidth}{0.5pt}\end{center}

\section{Issue 11: Startup Grace Period Masks Heartbeat but Not 0x204
Staleness}\label{issue-11-startup-grace-period-masks-heartbeat-but-not-0x204-staleness}

\subsection{11.1 The gap}\label{the-gap-4}

The startup grace period (3 seconds) masks heartbeat checks to prevent false
ESTOP at boot. But the \texttt{0x204} staleness check on SYS is NOT masked. On
cold boot:

\begin{enumerate}
\def\labelenumi{\arabic{enumi}.}
\tightlist
\item
  SYS boots first (fastest boot time).
\item
  RT is still booting --- \texttt{0x204} is absent.
\item
  At T+200ms after SYS \texttt{motor\_task} starts: no \texttt{0x204} received →
  staleness triggers: \texttt{setpoint.speed\ =\ 0,\ setpoint.gear\ =\ N}.
\end{enumerate}

This is a false trigger --- the vehicle isn't even running, so forcing speed=0
and gear=N is harmless (already the default). But it's logically inconsistent:
if the system is in a startup grace period where ``things aren't ready yet,''
ALL safety checks that depend on external data should be masked, not just
heartbeats.

The risk is low (the false trigger produces a safe outcome), but the design
principle is wrong --- it indicates the startup grace concept was applied
inconsistently.

\subsection{11.2 Recommended fix: Gate 0x204 staleness with the same startup
grace
period}\label{recommended-fix-gate-0x204-staleness-with-the-same-startup-grace-period}

\textbf{Approach:} Apply the same 3-second startup grace to the \texttt{0x204}
staleness check:

\begin{verbatim}
// In SYS motor_task:
bool startup_grace_active = (esp_timer_get_time() < kStartupGracePeriodUs);  // 3s

if (!startup_grace_active && time_since_last_0x204 > kSetpointStaleMs) {
    setpoint.speed = 0;
    setpoint.gear = N;
}
\end{verbatim}

This aligns the staleness check with the heartbeat check --- both are masked
during the 3-second startup window. After 3 seconds, the vehicle should have all
nodes online; if \texttt{0x204} is still absent, it's a genuine fault.

\begin{center}\rule{0.5\linewidth}{0.5pt}\end{center}

\section{Issue 12: ``Both Bulbs OFF = ESTOP'' Is Ambiguous at
Power-Off}\label{issue-12-both-bulbs-off-estop-is-ambiguous-at-power-off}

\subsection{12.1 The gap}\label{the-gap-5}

The rider's guide (§7.1 of emergency-system.md) states:

\begin{longtable}[]{@{}ll@{}}
\toprule\noalign{}
Indicators & Meaning \\
\midrule\noalign{}
\endhead
\bottomrule\noalign{}
\endlastfoot
Both bulbs OFF & ESTOP --- vehicle is emergency-stopped \\
\end{longtable}

But a powered-down vehicle also shows both bulbs OFF. The rider cannot
distinguish ESTOP-persisted from normal power-off without pressing START.
Pressing START on a powered-off vehicle does nothing (SYS isn't running to
detect the button). The rider might mistake power-off for ESTOP and try to
power-cycle --- which would actually start the vehicle in MANUAL mode,
potentially while they think they're troubleshooting an ESTOP condition.

Furthermore, during ESTOP the DC-DC converter is OFF (CAN
\texttt{0x012\ enable=0}), which means the \textbf{12V rail is dead.} Both the
AUTO bulb (GPIO25) and MANUAL bulb (GPIO26) are powered from the 12V accessory
rail via relay. If 12V is dead, the bulbs physically cannot illuminate --- even
if SYS drives the GPIOs. The brake light (GPIO21) is also on 12V and also cannot
illuminate.

\textbf{This means the OR logic ``brake light ON during ESTOP'' is physically
impossible with the current power architecture.} The 12V rail must be alive for
any lamp to work.

\subsection{12.2 Recommended fix: Keep DC-DC enabled during ESTOP; cut only the
accessory
relay}\label{recommended-fix-keep-dc-dc-enabled-during-estop-cut-only-the-accessory-relay}

\textbf{Approach:} During ESTOP, keep the DC-DC converter ON
(\texttt{0x012\ enable=1}) to maintain 12V for CAN transceivers, MCU power, and
the brake light. Cut only the 12V accessory relay (GPIO27) to kill non-safety
lights (headlight, turn signals, mode indicators).

\textbf{Revised ESTOP power behavior:}

\begin{longtable}[]{@{}
  >{\raggedright\arraybackslash}p{(\columnwidth - 4\tabcolsep) * \real{0.2750}}
  >{\raggedright\arraybackslash}p{(\columnwidth - 4\tabcolsep) * \real{0.3500}}
  >{\raggedright\arraybackslash}p{(\columnwidth - 4\tabcolsep) * \real{0.3750}}@{}}
\toprule\noalign{}
\begin{minipage}[b]{\linewidth}\raggedright
Component
\end{minipage} & \begin{minipage}[b]{\linewidth}\raggedright
Current ESTOP
\end{minipage} & \begin{minipage}[b]{\linewidth}\raggedright
Revised ESTOP
\end{minipage} \\
\midrule\noalign{}
\endhead
\bottomrule\noalign{}
\endlastfoot
DC-DC converter (0x012) & OFF & \textbf{ON} (maintains 12V for safety-critical
loads) \\
12V accessory relay (GPIO27) & OFF & OFF (cuts headlight, turn signals, mode
bulbs) \\
Brake light (GPIO21) & DEAD (no 12V) & \textbf{ON} (powered from DC-DC, not
through accessory relay) \\
Mode bulbs (GPIO25/26) & DEAD & OFF (cut by accessory relay --- intentional
ambiguity eliminated by brake light) \\
CAN transceivers & DEAD (no 5V) & \textbf{ALIVE} (5V from DC-DC→LDO) \\
MCUs & Variable & \textbf{ALIVE} (3.3V from DC-DC→LDO) \\
\end{longtable}

\textbf{This requires one wiring change:} The brake light (GPIO21 relay) must be
powered from the always-on DC-DC output, not from the accessory relay output.
This is a wiring change, not a new sensor. The MCUs are already on the always-on
rail (otherwise SYS couldn't run during ESTOP).

\textbf{With this change, ESTOP is visually distinct:} - Brake light:
\textbf{ON} (always, during ESTOP --- powered from DC-DC directly) - Mode bulbs:
\textbf{OFF} (cut by accessory relay) - Headlight, turn signals: \textbf{OFF}
(cut by accessory relay)

A powered-off vehicle: \textbf{everything OFF including brake light.} The rider
can now distinguish.

\textbf{Without the wiring change (immediate documentation fix):} Document that
DC-DC stays ON during ESTOP (to keep MCUs and CAN alive --- it always did, the
architecture was ambiguous). The brake light illuminates if it's on the
always-on rail; if it's not, document it as a known wiring limitation and
recommend the rewire.

\begin{center}\rule{0.5\linewidth}{0.5pt}\end{center}

\section{Issue 13: EPS-C Mechanical Jam Silent-Stop Recovery Path
Unclear}\label{issue-13-eps-c-mechanical-jam-silent-stop-recovery-path-unclear}

\subsection{13.1 The gap}\label{the-gap-6}

Emergency-system.md §3.3 describes the mechanical jam fallback: if following
error persists during ESTOP centering ramp (\textgreater5° for \textgreater1s),
fall back to silent-stop (stop transmitting \texttt{0x169}). EPS-C enters
internal fault-lock.

§4.6 says the boot sync timeout for EPS-C requires ``rider retry via START
short-press'' to reset to STEER\_LISTEN\_SYNC.

The question: does the mechanical-jam silent-stop during ESTOP enter the same
STEER\_FAULT state, making it recoverable via START short-press? Or does it
require a full power-cycle?

\subsection{13.2 State machine analysis}\label{state-machine-analysis}

During ESTOP centering ramp: - Steer SM is in \texttt{STEER\_ACTIVE}
(transmitting \texttt{0x169} with ramping targets). - Mechanical jam triggers:
following error \textgreater5° for \textgreater1s. - RT stops transmitting
\texttt{0x169} → EPS-C enters comm-fault.

At this point, the steer SM state should transition from \texttt{STEER\_ACTIVE}
to \texttt{STEER\_FAULT}. But the architecture doesn't specify this transition
--- it says ``fall back to silent-stop'' without naming the state.

If the steer SM remains in \texttt{STEER\_ACTIVE} (just silently stopped), then
when the rider exits ESTOP to MANUAL, the SM would still be
\texttt{STEER\_ACTIVE} but in MANUAL mode, which says ``do NOT send 0x169.''
This works --- no transmission. But the EPS-C is still fault-locked from the
silent-stop, so the rider has no steering assist in MANUAL.

If the steer SM transitions to \texttt{STEER\_FAULT}, the documented recovery
path applies: - Short-press START → reset to \texttt{STEER\_LISTEN\_SYNC} (EPS-C
must respond to \texttt{0x201}). - Long-press START (3s) + throttle at zero →
force-activate with target=0° (MANUAL only, AUTO locked out).

\subsection{13.3 Recommended fix: Explicit transition to STEER\_FAULT on
mechanical
jam}\label{recommended-fix-explicit-transition-to-steer_fault-on-mechanical-jam}

\textbf{Approach:} When the mechanical jam fallback triggers during ESTOP,
explicitly transition the steer SM to \texttt{STEER\_FAULT}. The existing
STEER\_FAULT recovery paths then apply:

\begin{verbatim}
// In steer SM, during ESTOP centering ramp:
if (following_error_deg > 5.0f && error_duration_ms > 1000) {
    // Mechanical jam — stop transmitting
    stop_0x169();
    steer_state = STEER_FAULT;  // explicit transition
    g_steer_ramp_in_progress = false;
    log_error("Steering mechanical jam during ESTOP centering");
}
\end{verbatim}

\textbf{Recovery sequence for the rider:} 1. Mechanical jam occurs during ESTOP
→ steering silent-stops. 2. Rider presses START → vehicle enters MANUAL mode. 3.
Steer SM is in \texttt{STEER\_FAULT} → no 0x169 transmission. 4. Rider
short-presses START → steer SM resets to \texttt{STEER\_LISTEN\_SYNC}. 5. If
EPS-C responds to \texttt{0x201} with \texttt{SES\_INF\_Angle\_Status\ ==\ 1}: →
\texttt{STEER\_ACTIVE} in MANUAL (EPS-C standalone --- steering wheel works
normally). 6. If EPS-C doesn't respond (still fault-locked): remains in
\texttt{STEER\_FAULT}. Rider long-presses START (3s) + throttle at zero →
force-activate with target=0°. 7. If force-activate also fails: power-cycle is
the final recovery.

\begin{center}\rule{0.5\linewidth}{0.5pt}\end{center}

\section{Issue 14: Fixed 5° Following Error Threshold at High
Speed}\label{issue-14-fixed-5-following-error-threshold-at-high-speed}

\subsection{14.1 The gap}\label{the-gap-7}

The steering following error threshold is a fixed 5° regardless of vehicle
speed. But at high speed, the dynamic clamp limits commanded steering to
\textasciitilde5°, making the 5° threshold represent a \textbf{100\% error}
(commanded 5°, actual 0°). At low speed, the dynamic clamp allows
\textasciitilde40°, making 5° only a \textbf{12.5\% error.}

The sensitivity is inverted: at high speed where precise steering matters most,
the threshold is coarsest relative to the command range. A 4° error at 25 km/h
(80\% of the allowed steering range) would NOT trigger ESTOP --- the vehicle
continues at speed with severely degraded steering authority.

\subsection{14.2 Quantified}\label{quantified}

\begin{longtable}[]{@{}lll@{}}
\toprule\noalign{}
Speed & Dynamic limit & 5° fixed threshold as \% of range \\
\midrule\noalign{}
\endhead
\bottomrule\noalign{}
\endlastfoot
2 km/h & 40° & 12.5\% --- appropriate \\
10 km/h & 18° & 27.8\% --- loose \\
25 km/h & 5° & \textbf{100\%} --- must have ZERO response to trigger \\
\end{longtable}

At 25 km/h, a 4° error (commanded 5°, actual 1°) represents the EPS-C delivering
only 20\% of the commanded steering angle. This is a degraded steering condition
at speed that the fixed 5° threshold would miss entirely.

\subsection{14.3 Recommended fix: Speed-scaled following error
threshold}\label{recommended-fix-speed-scaled-following-error-threshold}

\textbf{Approach:} Make the following error threshold proportional to the
dynamic angle clamp limit for the current speed, with a floor:

\begin{verbatim}
threshold = max(kMinFollowingErrorDeg, kFollowingErrorRatio * dynamic_limit)
\end{verbatim}

\begin{longtable}[]{@{}
  >{\raggedright\arraybackslash}p{(\columnwidth - 4\tabcolsep) * \real{0.3571}}
  >{\raggedright\arraybackslash}p{(\columnwidth - 4\tabcolsep) * \real{0.2500}}
  >{\raggedright\arraybackslash}p{(\columnwidth - 4\tabcolsep) * \real{0.3929}}@{}}
\toprule\noalign{}
\begin{minipage}[b]{\linewidth}\raggedright
Constant
\end{minipage} & \begin{minipage}[b]{\linewidth}\raggedright
Value
\end{minipage} & \begin{minipage}[b]{\linewidth}\raggedright
Rationale
\end{minipage} \\
\midrule\noalign{}
\endhead
\bottomrule\noalign{}
\endlastfoot
\texttt{kMinFollowingErrorDeg} & 2.0° & Floor --- prevents noise triggers at
very low speed \\
\texttt{kFollowingErrorRatio} & 0.25 & 25\% of the speed-safe range \\
\end{longtable}

\textbf{Behavior:}

\begin{longtable}[]{@{}llll@{}}
\toprule\noalign{}
Speed & Dynamic limit & New threshold & Old threshold \\
\midrule\noalign{}
\endhead
\bottomrule\noalign{}
\endlastfoot
2 km/h & 40° & max(2°, 10°) = \textbf{10°} & 5° (was too sensitive for
parking) \\
10 km/h & 18° & max(2°, 4.5°) = \textbf{4.5°} & 5° (similar) \\
25 km/h & 5° & max(2°, 1.25°) = \textbf{2°} & 5° (was 2.5× too loose) \\
\end{longtable}

At low speed: the threshold loosens (10° vs 5°), appropriate for parking
maneuvers where mechanical compliance and tire scrub produce larger following
errors that aren't safety-relevant.

At high speed: the threshold tightens (2° vs 5°), catching degraded steering
authority before it becomes dangerous. A 4° error at 25 km/h would now trigger
ESTOP --- which it should, because the EPS-C has lost \textasciitilde80\% of its
commanded authority.

\textbf{Implementation:} This is a one-line change in RT's following error check
--- replace the hardcoded \texttt{5.0f} with a call to compute the speed-scaled
threshold using the existing dynamic clamp value.

\begin{center}\rule{0.5\linewidth}{0.5pt}\end{center}

\section{Updated Summary (Issues 8-14)}\label{updated-summary-issues-8-14}

\begin{longtable}[]{@{}
  >{\raggedright\arraybackslash}p{(\columnwidth - 8\tabcolsep) * \real{0.1029}}
  >{\raggedright\arraybackslash}p{(\columnwidth - 8\tabcolsep) * \real{0.1471}}
  >{\raggedright\arraybackslash}p{(\columnwidth - 8\tabcolsep) * \real{0.2353}}
  >{\raggedright\arraybackslash}p{(\columnwidth - 8\tabcolsep) * \real{0.2206}}
  >{\raggedright\arraybackslash}p{(\columnwidth - 8\tabcolsep) * \real{0.2941}}@{}}
\toprule\noalign{}
\begin{minipage}[b]{\linewidth}\raggedright
Issue
\end{minipage} & \begin{minipage}[b]{\linewidth}\raggedright
Severity
\end{minipage} & \begin{minipage}[b]{\linewidth}\raggedright
Fix complexity
\end{minipage} & \begin{minipage}[b]{\linewidth}\raggedright
New hardware?
\end{minipage} & \begin{minipage}[b]{\linewidth}\raggedright
Recommended action
\end{minipage} \\
\midrule\noalign{}
\endhead
\bottomrule\noalign{}
\endlastfoot
\textbf{8. No brake monitor} & High & Low (software, SYS-local) & \textbf{None}
--- SEB provides 0x721 stroke/pressure/error at 100 Hz & Add brake
following-error check comparing 0x7B9 cmd vs 0x721 feedback. Monitor
SEB\_Error\_Status for Level 3 faults. Add 0x721 staleness check. \\
\textbf{9. 0x001 spoofable} & Medium & Low (software, all nodes) & None &
Rate-limit ESTOP processing (max 2 frames per 500ms window). Add DLC=1 sender-ID
byte to 0x001. \\
\textbf{10. No MTR ESTOP ACK} & Medium & Low (software, MTR+SYS) & None & MTR
sets ESTOP\_ACTIVE bit in 0x206 fault\_flags. SYS checks ACK within 100ms of
ESTOP. Add 0x206 staleness check. \\
\textbf{11. Startup grace inconsistent} & Low & Trivial (software, SYS) & None &
Gate 0x204 staleness check with same 3s startup grace as heartbeat. \\
\textbf{12. ESTOP HMI ambiguous} & Medium & Low (software + wiring) &
\textbf{Rewire only} --- brake light to always-on DC-DC rail, not accessory
relay & Keep DC-DC ON during ESTOP; cut only accessory relay (GPIO27). Brake
light illuminates during ESTOP. Mode bulbs OFF during ESTOP (distinct from
power-off where all lights OFF). \\
\textbf{13. EPS-C jam recovery unclear} & Low & Trivial (documentation +
software) & None & Explicitly transition to STEER\_FAULT on mechanical jam.
Document recovery via START short-press → LISTEN\_SYNC retry. \\
\textbf{14. Fixed 5° following error} & Medium & Trivial (software, RT-local) &
None & Speed-scaled threshold: \texttt{max(2°,\ 0.25\ ×\ dynamic\_limit)}. 2° at
25 km/h, 10° at 2 km/h. \\
\end{longtable}

\textbf{All seven issues are software-fixable.} Issue 12 requires a brake light
wiring change (to the always-on DC-DC rail) but no new sensors or components.


\chapter{Emergency Stop \& Emergency
Handling}\label{emergency-stop-emergency-handling}

The E-Trike emergency system uses \textbf{defense in depth} --- multiple
independent, overlapping safety layers. No single component failure defeats all
protections. This document describes how the emergency stop (ESTOP) works, how
the system handles other emergency conditions, and what the rider should expect.

\begin{center}\rule{0.5\linewidth}{0.5pt}\end{center}

\section{1. ESTOP Overview}\label{estop-overview}

ESTOP is the system's \textbf{absorbing safety state}. When entered, all
actuators return to their safe positions and the vehicle stops. ESTOP persists
until deliberately cleared by the rider --- it never self-clears.

\subsection{1.1 What happens in ESTOP}\label{what-happens-in-estop}

\begin{longtable}[]{@{}
  >{\raggedright\arraybackslash}p{(\columnwidth - 4\tabcolsep) * \real{0.4231}}
  >{\raggedright\arraybackslash}p{(\columnwidth - 4\tabcolsep) * \real{0.3077}}
  >{\raggedright\arraybackslash}p{(\columnwidth - 4\tabcolsep) * \real{0.2692}}@{}}
\toprule\noalign{}
\begin{minipage}[b]{\linewidth}\raggedright
Subsystem
\end{minipage} & \begin{minipage}[b]{\linewidth}\raggedright
Action
\end{minipage} & \begin{minipage}[b]{\linewidth}\raggedright
Notes
\end{minipage} \\
\midrule\noalign{}
\endhead
\bottomrule\noalign{}
\endlastfoot
\textbf{Motor throttle} & MCP4725 DAC = 0 V & Instant motor kill --- zero
voltage to controller \\
\textbf{Gear selection} & All relays OFF → Neutral & 72V gear lines
de-energized \\
\textbf{Brake (SEB)} & Stroke = max (\textasciitilde27 mm) & Full hydraulic
brake pressure, 50 Hz CAN 0x7B9 \\
\textbf{Steering (EPS-C)} & Obstacle trigger: hold angle, silent-stop after
500ms. Non-obstacle: ramp to 0° at 20°/s via active 0x169. Fallback to
silent-stop if following error persists \textgreater1s (mechanical jam). &
Two-tier response --- see §2.5 \\
\textbf{DC-DC converter} & CAN 0x012 enable = \textbf{1} (maintains 12V for
MCUs, CAN transceivers, and brake light --- MCUs need power to run safety tasks)
& \\
\textbf{12V accessory relay} & GPIO27 OFF (cuts headlight, turn signals, mode
bulbs --- non-safety loads) & \\
\textbf{Signal lights} & Brake light \textbf{ON} (powered from always-on DC-DC
rail, not through accessory relay), all others OFF & \\
\textbf{Mode indicators} & Both AUTO and MANUAL bulbs OFF & Dark dashboard =
ESTOP \\
\textbf{Throttle pass-through} & Ignored & ADC reads ignored, DAC forced to 0 \\
\textbf{Gear pass-through} & Ignored & All gear relays forced OFF \\
\end{longtable}

\subsection{1.2 ESTOP exit --- START button with deferred steering ramp
completion}\label{estop-exit-start-button-with-deferred-steering-ramp-completion}

ESTOP can only be exited by pressing the \textbf{START button} (green, GPIO32 on
SYS) or via \textbf{MODE button long-press} (GPIO11, 3-second hold). Both exit
to \textbf{MANUAL mode} --- the rider is always in direct control after an
ESTOP, never in AUTO.

\begin{itemize}
\tightlist
\item
  \textbf{START button (GPIO32):} Short press exits ESTOP → MANUAL immediately
  (subject to steering ramp completion, see below).
\item
  \textbf{MODE button (GPIO11):} Long-press (3 seconds) exits ESTOP → MANUAL.
  This is the secondary exit path --- two independent GPIOs on separate physical
  buttons ensure no single button failure locks the rider in ESTOP.
\item
  MODE button short-press is ignored in ESTOP (prevents accidental exit).
\item
  CAN commands cannot exit ESTOP.
\item
  No automatic timeout recovery.
\item
  Power-cycle always exits ESTOP (reboot starts in MANUAL by default) ---
  ultimate fallback.
\end{itemize}

\begin{verbatim}
ESTOP exit: START button → MANUAL mode
           MODE button long-press (3s) → MANUAL mode (secondary)
           Power-cycle → MANUAL mode (ultimate fallback)
           (NOT CAN command, NOT timeout)
\end{verbatim}

\textbf{START button health monitoring:} The \texttt{diag\_task} monitors ESTOP
duration. If the vehicle remains in ESTOP for \textgreater30 seconds with no
START button activity, a diagnostic flag is set in CAN
\texttt{0x600\ SYS\_DIAG\_RPT}. This catches a stuck or disconnected START
button before the rider discovers it at roadside.

\textbf{Critical: steering ramp completion before silent handoff.} During
non-obstacle ESTOP, RT is actively transmitting \texttt{0x169} to ramp steering
to 0° at 20°/s (see §3.2). If the rider presses START before the ramp completes,
\textbf{the mode transition is deferred for steering only} --- RT continues the
centering ramp to completion, then hands off silently to MANUAL mode. All other
subsystems (brake, motor, lights) transition to MANUAL immediately.

\textbf{Why this matters:} Without the deferral, pressing START mid-ramp would
stop \texttt{0x169} transmission immediately, leaving EPS-C to comm-fault at
whatever off-center angle the ramp had reached. A 30° steering lock in MANUAL
mode (where EPS-C is standalone and provides no active centering) would make the
vehicle unrideable. See {[}{[}emergency-safety-analysis{]}{]} §1 for the full
causal trace.

\textbf{The rider's experience:} 1. Press START button. 2. Brake transitions
from max to lever-controlled immediately (rider can release brake). 3. Motor and
gear transition to MANUAL pass-through immediately. 4. Steering continues to
center itself (up to 2 seconds from full lock). 5. Once centered, steering
becomes fully manual (EPS-C standalone). 6. The MANUAL indicator bulb
illuminates when the full transition is complete.

\begin{center}\rule{0.5\linewidth}{0.5pt}\end{center}

\section{2. ESTOP Trigger Paths --- The 8 Safety
Layers}\label{estop-trigger-paths-the-8-safety-layers}

\subsection{Layer 1: Physical ESTOP Button
(fastest)}\label{layer-1-physical-estop-button-fastest}

A \textbf{big red mushroom button} wired \textbf{normally-closed (NC)} to SYS
GPIO1 and MTR STM32 kEstopGpio.

\begin{longtable}[]{@{}ll@{}}
\toprule\noalign{}
Property & Value \\
\midrule\noalign{}
\endhead
\bottomrule\noalign{}
\endlastfoot
GPIO & SYS: GPIO1, MTR: kEstopGpio (TBD STM32 pin) \\
Type & NC (normally-closed), active-low \\
Detection & SYS \texttt{safety\_task} @ 20 Hz, MTR direct ISR \\
CAN redundancy & Also broadcasts CAN \texttt{0x001} on low bus \\
\end{longtable}

\textbf{Fail-safe by construction:} A cut wire, disconnected plug, or power loss
pulls the GPIO LOW --- all of which trigger ESTOP. There is no failure mode
where the button is broken but the system thinks it's fine.

\begin{verbatim}
Normal operation: GPIO reads HIGH (10k pull-up to 3.3V) → system runs
Button pressed:   GPIO reads LOW (switch connects to GND) → ESTOP
Wire cut/broken:  GPIO reads LOW (pull-up defeated) → ESTOP
Power loss:       GPIO reads LOW → ESTOP
\end{verbatim}

\textbf{MTR STM32 direct path:} The ESTOP button is also wired directly to the
MTR STM32 board. On ESTOP, MTR cuts throttle (MCP4725 = 0V) and gear (all relays
OFF) locally --- zero CAN dependency, zero software stack dependency. This is
the EGAS Level 3 hardware path.

\subsection{\texorpdfstring{Layer 2: CAN \texttt{0x001}
SAFETY\_ESTOP}{Layer 2: CAN 0x001 SAFETY\_ESTOP}}\label{layer-2-can-0x001-safety_estop}

Any node can broadcast \texttt{0x001} on either CAN bus to trigger ESTOP. This
is the redundant communication path.

\begin{longtable}[]{@{}
  >{\raggedright\arraybackslash}p{(\columnwidth - 2\tabcolsep) * \real{0.5882}}
  >{\raggedright\arraybackslash}p{(\columnwidth - 2\tabcolsep) * \real{0.4118}}@{}}
\toprule\noalign{}
\begin{minipage}[b]{\linewidth}\raggedright
Property
\end{minipage} & \begin{minipage}[b]{\linewidth}\raggedright
Value
\end{minipage} \\
\midrule\noalign{}
\endhead
\bottomrule\noalign{}
\endlastfoot
CAN ID & \texttt{0x001} --- highest priority on both buses (wins every
arbitration) \\
DLC & 1 (sender-ID byte: 0=SYS, 1=RT, 2=Jetson, 3=MTR) \\
Senders & Jetson, RT, SYS, MTR \\
Forwarding & RT transparently forwards between buses (bypasses gateway queue) \\
\end{longtable}

\textbf{Rate limiting:} Each node processes at most \textbf{2 ESTOP frames per
500ms window}. The first frame triggers ESTOP; subsequent floods are ignored.
This prevents a corrupted node from flooding the bus with persistent ESTOP
(denial-of-service). The sender-ID byte enables per-sender rate limiting and
diagnostic logging of which node triggered ESTOP.

\textbf{Why each node might trigger:} - \textbf{Jetson:} Perception stack
detects imminent collision that the obstacle sensor missed. - \textbf{RT:}
Steering following error exceeds threshold, or SYS heartbeat lost. -
\textbf{SYS:} Physical ESTOP button pressed (redundant CAN path), or RT
heartbeat lost. - \textbf{MTR:} Local ESTOP GPIO triggered, or motor fault
detected.

\subsection{Layer 3: Heartbeat Timeout --- Node Crash
Detection}\label{layer-3-heartbeat-timeout-node-crash-detection}

Each MCU sends a 1-byte alive counter on its own CAN ID at 2 Hz. A frozen
counter = a frozen node.

\begin{longtable}[]{@{}
  >{\raggedright\arraybackslash}p{(\columnwidth - 10\tabcolsep) * \real{0.1875}}
  >{\raggedright\arraybackslash}p{(\columnwidth - 10\tabcolsep) * \real{0.1875}}
  >{\raggedright\arraybackslash}p{(\columnwidth - 10\tabcolsep) * \real{0.1667}}
  >{\raggedright\arraybackslash}p{(\columnwidth - 10\tabcolsep) * \real{0.1042}}
  >{\raggedright\arraybackslash}p{(\columnwidth - 10\tabcolsep) * \real{0.1875}}
  >{\raggedright\arraybackslash}p{(\columnwidth - 10\tabcolsep) * \real{0.1667}}@{}}
\toprule\noalign{}
\begin{minipage}[b]{\linewidth}\raggedright
Watcher
\end{minipage} & \begin{minipage}[b]{\linewidth}\raggedright
Watchee
\end{minipage} & \begin{minipage}[b]{\linewidth}\raggedright
CAN ID
\end{minipage} & \begin{minipage}[b]{\linewidth}\raggedright
Bus
\end{minipage} & \begin{minipage}[b]{\linewidth}\raggedright
Timeout
\end{minipage} & \begin{minipage}[b]{\linewidth}\raggedright
Action
\end{minipage} \\
\midrule\noalign{}
\endhead
\bottomrule\noalign{}
\endlastfoot
SYS & RT & \texttt{0x7FD} & Low & 1000ms (2 missed) & ESTOP (AUTO only) \\
RT & SYS & \texttt{0x7FE} & Low & \textbf{200ms} (2 missed at 10 Hz) &
\textbf{RT brake takeover:} RT immediately starts transmitting 0x7B9 with
stroke=max (full brake) regardless of current mode. Also broadcasts CAN 0x001.
MTR kills motor + gear locally. Brake gap ≤220ms. \\
RT & Jetson & \texttt{0x7FC} & High & 1500ms (3 missed) & \textbf{Assisted
stop:} zero 0x204 speed, stop 0x169, 0x205 brake=2000 kPa (\textasciitilde2 MPa
moderate brake), transition SYS to MANUAL. Brake light ON. Rider can override
with lever. DC-DC stays on. \\
Jetson & RT & \texttt{0x7FD} & High & 1500ms (3 missed) & Stop publishing
\texttt{/cmd\_vel} \\
\end{longtable}

\textbf{Alive counter validation:} A frame with the same counter as the previous
frame = stuck CAN controller (MCU hung, controller DMA-ing from buffer). Treated
as a missed heartbeat.

\begin{verbatim}
bool heartbeat_is_fresh(uint8_t new_ctr) {
    if (new_ctr != last_alive_ctr) { last_alive_ctr = new_ctr; return true; }
    return false;  // frozen — same counter twice
}
\end{verbatim}

\textbf{Startup grace period:} Heartbeat checks are masked for the first
\textbf{3 seconds} after boot to prevent false ESTOP before the first frames
arrive.

\subsection{Layer 4: Command Staleness --- Jetson Hang
Detection}\label{layer-4-command-staleness-jetson-hang-detection}

RT monitors the last received \texttt{0x300\ HOST\_DRIVE\_CMD} timestamp. If
Jetson's planning stack or ROS→CAN bridge hangs:

\begin{longtable}[]{@{}llll@{}}
\toprule\noalign{}
Check & Watcher & Timeout & Action \\
\midrule\noalign{}
\endhead
\bottomrule\noalign{}
\endlastfoot
\texttt{0x300} staleness & RT \texttt{watchdog} task & 500ms & Zero
\texttt{0x204} + stop \texttt{0x169} \\
\texttt{0x204} staleness & SYS \texttt{motor} task & 200ms & Zero speed +
Neutral gear \\
\end{longtable}

The 500ms RT timeout allows for brief Jetson hiccups. The 200ms SYS timeout is
tighter because stale actuator data is immediately dangerous.

\subsection{Layer 5: Steering Following Error --- Mechanical Fault
Detection}\label{layer-5-steering-following-error-mechanical-fault-detection}

RT compares commanded steering angle (\texttt{0x169}) against actual angle
reported by EPS-C (\texttt{0x201\ SES\_StrAngle}) at 100 Hz.

\begin{longtable}[]{@{}
  >{\raggedright\arraybackslash}p{(\columnwidth - 2\tabcolsep) * \real{0.6111}}
  >{\raggedright\arraybackslash}p{(\columnwidth - 2\tabcolsep) * \real{0.3889}}@{}}
\toprule\noalign{}
\begin{minipage}[b]{\linewidth}\raggedright
Parameter
\end{minipage} & \begin{minipage}[b]{\linewidth}\raggedright
Value
\end{minipage} \\
\midrule\noalign{}
\endhead
\bottomrule\noalign{}
\endlastfoot
Error threshold & \textbf{Speed-scaled:}
\texttt{max(2°,\ 0.25\ ×\ dynamic\_limit)} --- 2° at 25 km/h, 4.5° at 10 km/h,
10° at 2 km/h \\
Persistence duration & \textgreater300ms \\
Action & ESTOP \\
\end{longtable}

\textbf{Why speed-scaled:} At high speed, the dynamic clamp limits steering to
\textasciitilde5°. A fixed 5° threshold would require a 100\% error (zero
response) to trigger --- missing an 80\% authority loss. At low speed (40°
limit), 5° is only 12.5\% --- too sensitive for parking maneuvers where tire
scrub produces normal following errors. The speed-scaled threshold is tight at
speed (where precision matters) and tolerant at low speed (where it doesn't).

This catches stuck linkage, rock jams, bent tie rods, EPS-C motor failure,
encoder failure, and corrupted CAN frames.

\subsection{Layer 6: Dynamic Angle Clamp --- Rollover
Prevention}\label{layer-6-dynamic-angle-clamp-rollover-prevention}

The maximum allowable steering angle is inversely proportional to vehicle speed,
enforced in RT's \texttt{physics\_resolve()} before the command reaches the
actuator:

\begin{longtable}[]{@{}ll@{}}
\toprule\noalign{}
Speed & Max angle \\
\midrule\noalign{}
\endhead
\bottomrule\noalign{}
\endlastfoot
2 km/h & \textasciitilde40° \\
10 km/h & \textasciitilde18° \\
25 km/h & \textasciitilde5° \\
\end{longtable}

Even if Jetson commands a dangerous turn (bug, bad planning, or adversarial
input), RT clamps it to the safe envelope.

\subsection{Layer 7: Software Hard-Stops --- Mechanical
Protection}\label{layer-7-software-hard-stops-mechanical-protection}

All commanded steering angles are clamped to ±40° regardless of source. This
prevents the motor from pushing against the physical end-stops (which could
strip gears, burn the motor, or bend the linkage).

\subsection{Layer 8: External Watchdog --- Last-Resort Hardware
Reset}\label{layer-8-external-watchdog-last-resort-hardware-reset}

Each ESP32-S3 has a dedicated external watchdog IC (TPS3850) on a separate chip.

\begin{longtable}[]{@{}lllll@{}}
\toprule\noalign{}
Node & Toggle GPIO & Toggled by & Period & Timeout \\
\midrule\noalign{}
\endhead
\bottomrule\noalign{}
\endlastfoot
RT & GPIO21 & \texttt{control\_task} & 100 Hz (every 10ms) &
\textasciitilde100ms \\
SYS & GPIO23 & \texttt{safety\_task} & 20 Hz (every 50ms) &
\textasciitilde100ms \\
\end{longtable}

If the MCU hangs and toggling stops, the watchdog IC asserts a hardware reset.
On reset: 1. All GPIOs go high-impedance → motor controller sees 0V throttle,
all gear relays open. 2. ESP32 reboots → starts in MANUAL mode → re-runs Listen
Before Speaking sequences. 3. Recovery time from watchdog fire to first valid
CAN frame: \textless500ms.

\textbf{The external watchdog catches failures that internal watchdogs miss:}
crystal oscillator failure, voltage brownout, silicon latch-up from ESD.

\subsection{Layer 9: Brake Following-Error Monitor --- Actuator Feedback
Verification}\label{layer-9-brake-following-error-monitor-actuator-feedback-verification}

The SEB provides rich feedback via CAN \texttt{0x721\ SEB\_STATUS} at 100 Hz:
\texttt{SEB\_Stroke\_Value} (actual stroke position, 0.05mm resolution),
\texttt{SEB\_Pressure\_Value} (actual hydraulic pressure, 0.05 MPa resolution),
and \texttt{SEB\_Error\_Status} (fault level 0-3 including ``severe, request
shutdown''). This data is already on the CAN bus --- no new sensors needed.

SYS compares the commanded brake stroke/pressure (from \texttt{0x7B9}) against
the actual feedback (from \texttt{0x721}):

\begin{longtable}[]{@{}
  >{\raggedright\arraybackslash}p{(\columnwidth - 6\tabcolsep) * \real{0.1944}}
  >{\raggedright\arraybackslash}p{(\columnwidth - 6\tabcolsep) * \real{0.3056}}
  >{\raggedright\arraybackslash}p{(\columnwidth - 6\tabcolsep) * \real{0.2778}}
  >{\raggedright\arraybackslash}p{(\columnwidth - 6\tabcolsep) * \real{0.2222}}@{}}
\toprule\noalign{}
\begin{minipage}[b]{\linewidth}\raggedright
Check
\end{minipage} & \begin{minipage}[b]{\linewidth}\raggedright
Threshold
\end{minipage} & \begin{minipage}[b]{\linewidth}\raggedright
Debounce
\end{minipage} & \begin{minipage}[b]{\linewidth}\raggedright
Action
\end{minipage} \\
\midrule\noalign{}
\endhead
\bottomrule\noalign{}
\endlastfoot
Stroke following error & abs(cmd − actual) \textgreater{} 3mm & 100ms (10
frames) & Log fault, set diag flag --- cannot escalate beyond ESTOP (already max
brake) \\
SEB Level 3 fault & \texttt{SEB\_Error\_Status\ ==\ 3} & Immediate & Log SEB
severe fault --- SEB itself reports it cannot function \\
0x721 staleness & No frame for \textgreater100ms & --- & SEB comms lost ---
brake system offline \\
\end{longtable}

\textbf{Why this matters even though ESTOP is already the maximum response:} If
\texttt{0x7B9\ stroke=max} is transmitted but SEB never receives or acts on it,
the vehicle coasts without brakes. The motor is killed (safe) but there's no
active deceleration. The brake monitor detects this condition and records it ---
essential for post-incident analysis and system improvement. During non-ESTOP
operation, the monitor catches brake degradation before it becomes critical.

\begin{center}\rule{0.5\linewidth}{0.5pt}\end{center}

\section{3. ESTOP Steering --- Two-Tier
Response}\label{estop-steering-two-tier-response}

The steering behavior during ESTOP depends on what triggered it:

\subsection{3.1 Obstacle-Triggered ESTOP}\label{obstacle-triggered-estop}

If ESTOP is triggered by an obstacle detection (Jetson perception, RT obstacle
sensor): 1. \textbf{Hold current angle, clamped to the dynamic angle limit} for
the current speed. If the current angle exceeds the speed-safe envelope (e.g.,
cornering at 15° when the dynamic limit at this speed is 5°), ramp the angle
down to the limit at 20°/s, then hold. 2. After \textbf{500ms} of hold at the
clamped angle: \textbf{silent-stop} (stop transmitting \texttt{0x169}). 3. EPS-C
enters internal comm-fault timeout.

\textbf{Why the dynamic clamp applies to the hold angle:} During emergency
braking while cornering, lateral load transfer reduces rear wheel grip. Holding
a large steering angle under hard braking on a three-wheeled vehicle approaches
the rollover threshold (see physics model §8:
\texttt{a\_y\ =\ v²/L·tan(δ)\ \textgreater{}\ g·w/(2h)}). The dynamic angle
clamp normally prevents this for commanded angles --- it must also apply to the
ESTOP hold angle. At speed, the clamp reduces the hold angle to near-straight
(\textasciitilde5° at 25 km/h), which is still compatible with ``don't veer into
adjacent lanes.'' At low speed (\textless2 km/h), the full steering range is
permitted.

\textbf{Straight-line case is unchanged:} If the vehicle is traveling straight
(δ ≈ 0°) when the obstacle is detected, the hold angle is 0° --- well within the
dynamic clamp at all speeds. The behavior is identical to the previous design.

See {[}{[}emergency-safety-analysis{]}{]} §4 for the full causal trace of the
cornering rollover scenario.

\subsection{3.2 Non-Obstacle ESTOP}\label{non-obstacle-estop}

For all other triggers (button press, heartbeat loss, following error, command
stale): 1. \textbf{Ramp to 0° at 20°/s} via active \texttt{0x169} transmission.
2. Continue transmitting during the ramp so EPS-C tracks the centering. 3. Hold
at 0° once centered. 4. \textbf{The ramp is non-interruptible by START button.}
See §1.2 --- the mode transition to MANUAL defers until the ramp completes. This
guarantees the vehicle is centered before the rider takes over.

\textbf{Rationale:} Straight wheels are the safest default position for a
stopped vehicle.

\subsection{3.3 Mechanical Jam Fallback}\label{mechanical-jam-fallback}

If the centering ramp encounters persistent following error (speed-scaled
threshold for \textgreater1s during the ramp): 1. \textbf{Fall back to
silent-stop} --- linkage is likely mechanically jammed. 2. Stop transmitting
\texttt{0x169}. 3. \textbf{Steer SM transitions to \texttt{STEER\_FAULT}}
explicitly. 4. EPS-C timeout-faults internally.

\textbf{Recovery:} After exiting ESTOP to MANUAL, the steer SM is in
\texttt{STEER\_FAULT}. The rider has two recovery options (same as EPS-C boot
sync failure): - \textbf{START short-press} → reset to
\texttt{STEER\_LISTEN\_SYNC} (waits for \texttt{0x201}, syncs current angle). -
\textbf{START long-press (3s) + throttle at zero} → force-activate with
target=0° (MANUAL mode only; AUTO locked out with blinking AUTO bulb). - If
neither works (EPS-C remains fault-locked): \textbf{power-cycle} as final
recovery.

\textbf{Rationale:} Don't burn out the steering motor fighting a mechanical jam.
Once the jam is cleared (manually by the rider), the standard STEER\_FAULT
recovery path restores steering without requiring a full power-cycle.

\begin{center}\rule{0.5\linewidth}{0.5pt}\end{center}

\section{4. Other Emergency Conditions}\label{other-emergency-conditions}

Beyond ESTOP, the system handles several emergency scenarios with graduated
responses.

\subsection{4.1 Jetson Failure}\label{jetson-failure}

\begin{longtable}[]{@{}
  >{\raggedright\arraybackslash}p{(\columnwidth - 4\tabcolsep) * \real{0.3000}}
  >{\raggedright\arraybackslash}p{(\columnwidth - 4\tabcolsep) * \real{0.3667}}
  >{\raggedright\arraybackslash}p{(\columnwidth - 4\tabcolsep) * \real{0.3333}}@{}}
\toprule\noalign{}
\begin{minipage}[b]{\linewidth}\raggedright
Symptom
\end{minipage} & \begin{minipage}[b]{\linewidth}\raggedright
Detection
\end{minipage} & \begin{minipage}[b]{\linewidth}\raggedright
Response
\end{minipage} \\
\midrule\noalign{}
\endhead
\bottomrule\noalign{}
\endlastfoot
Planning stack crash & No \texttt{0x300} for \textgreater500ms (RT watchdog) &
Zero \texttt{0x204} speed + stop \texttt{0x169} steering \\
ROS→CAN bridge crash & No \texttt{0x300} for \textgreater500ms & Same as
above \\
Jetson kernel panic & Heartbeat \texttt{0x7FC} frozen \textgreater1500ms & RT
zeros setpoints \\
Jetson sends dangerous steering & Dynamic angle clamp (Layer 6) & Capped to safe
envelope \\
\end{longtable}

\textbf{Response is controlled stop, not ESTOP.} The vehicle coasts to a stop
(motor zero, steering stops). The rider maintains manual brake and steering
control. MANUAL mode remains available.

\subsection{4.2 RT ESP32-S3 Failure}\label{rt-esp32-s3-failure}

\begin{longtable}[]{@{}
  >{\raggedright\arraybackslash}p{(\columnwidth - 4\tabcolsep) * \real{0.3000}}
  >{\raggedright\arraybackslash}p{(\columnwidth - 4\tabcolsep) * \real{0.3667}}
  >{\raggedright\arraybackslash}p{(\columnwidth - 4\tabcolsep) * \real{0.3333}}@{}}
\toprule\noalign{}
\begin{minipage}[b]{\linewidth}\raggedright
Symptom
\end{minipage} & \begin{minipage}[b]{\linewidth}\raggedright
Detection
\end{minipage} & \begin{minipage}[b]{\linewidth}\raggedright
Response
\end{minipage} \\
\midrule\noalign{}
\endhead
\bottomrule\noalign{}
\endlastfoot
RT crashes / hangs & SYS heartbeat monitor: \texttt{0x7FD} frozen
\textgreater1000ms & ESTOP (AUTO only) --- full brake, motor kill \\
RT control loop dies & SYS \texttt{0x204} staleness \textgreater200ms & Zero
speed + Neutral (faster than heartbeat) \\
RT CAN TX fails & TWAI TX errors & Log, auto-recover; persistent → EPS-C
timeout-fault \\
External watchdog fires & TPS3850 on GPIO21 & Hardware reset → safe state →
reboot → MANUAL \\
\end{longtable}

\textbf{Two independent checks on SYS:} 1. \texttt{0x204} staleness at 200ms →
zero speed (fast path, data quality). 2. \texttt{0x7FD} heartbeat at 1000ms →
ESTOP (slow path, node liveness).

\subsection{4.3 SYS ESP32-S3 Failure}\label{sys-esp32-s3-failure}

\begin{longtable}[]{@{}
  >{\raggedright\arraybackslash}p{(\columnwidth - 4\tabcolsep) * \real{0.3000}}
  >{\raggedright\arraybackslash}p{(\columnwidth - 4\tabcolsep) * \real{0.3667}}
  >{\raggedright\arraybackslash}p{(\columnwidth - 4\tabcolsep) * \real{0.3333}}@{}}
\toprule\noalign{}
\begin{minipage}[b]{\linewidth}\raggedright
Symptom
\end{minipage} & \begin{minipage}[b]{\linewidth}\raggedright
Detection
\end{minipage} & \begin{minipage}[b]{\linewidth}\raggedright
Response
\end{minipage} \\
\midrule\noalign{}
\endhead
\bottomrule\noalign{}
\endlastfoot
SYS crashes / hangs & RT heartbeat monitor: \texttt{0x7FE} frozen
\textgreater1000ms & RT broadcasts CAN \texttt{0x001} ESTOP \\
SYS throttle DAC fails & MCP4725 I2C NACK & Throttle = 0V (I2C bus reset) \\
SYS gear relay fails & Relay coil open & Gear = N (relay de-energizes ---
fail-safe) \\
External watchdog fires & TPS3850 on GPIO23 & Hardware reset → safe state →
reboot → MANUAL \\
\end{longtable}

\textbf{Physical fallbacks during SYS failure:} - Brake lever: SYS reads GPIO2
and transmits CAN 0x7B9 to SEB. This is a \textbf{by-wire} path --- the SEB is
electro-hydraulic with no mechanical linkage from lever to caliper. If SYS fails
completely (CAN TX dead), the brake lever cannot actuate the SEB. However, SYS
uses BRAKE\_DEGRADED (not terminal FAULT) on boot sync failure, so the lever
works as long as SYS can transmit CAN frames. If SYS is fully dead (watchdog
reset in progress), see §4.7 Electrical Faults for the watchdog brake window. -
Steering still responds to handlebars (EPS-C standalone in MANUAL). - ESTOP
button still works on RT (CAN \texttt{0x001}) and MTR (direct GPIO).

\subsection{4.4 MTR STM32 Failure}\label{mtr-stm32-failure}

\begin{longtable}[]{@{}
  >{\raggedright\arraybackslash}p{(\columnwidth - 4\tabcolsep) * \real{0.3000}}
  >{\raggedright\arraybackslash}p{(\columnwidth - 4\tabcolsep) * \real{0.3667}}
  >{\raggedright\arraybackslash}p{(\columnwidth - 4\tabcolsep) * \real{0.3333}}@{}}
\toprule\noalign{}
\begin{minipage}[b]{\linewidth}\raggedright
Symptom
\end{minipage} & \begin{minipage}[b]{\linewidth}\raggedright
Detection
\end{minipage} & \begin{minipage}[b]{\linewidth}\raggedright
Response
\end{minipage} \\
\midrule\noalign{}
\endhead
\bottomrule\noalign{}
\endlastfoot
MTR crashes / hangs & \texttt{0x206\ MTR\_MOTOR\_FBK} stops, \texttt{0x120}
stops & SYS detects mismatch between 0x204 setpoint and 0x206 feedback → CAN
\texttt{0x001} ESTOP \\
MTR ESTOP GPIO triggered & Direct GPIO read & Local throttle cut + gear OFF, CAN
\texttt{0x001} broadcast \\
MTR I2C bus fault & MCP4725 NACK & Motor controller sees 0V (DAC output defaults
to 0) \\
\end{longtable}

The MTR STM32 is the EGAS Level 1 Function Controller for motor actuation. SYS
ESP32-S3 is the Level 2 Function Monitor --- it compares commanded vs actual
speed and triggers ESTOP on mismatch.

\subsection{4.5 CAN Bus Failure}\label{can-bus-failure}

\begin{longtable}[]{@{}
  >{\raggedright\arraybackslash}p{(\columnwidth - 4\tabcolsep) * \real{0.3000}}
  >{\raggedright\arraybackslash}p{(\columnwidth - 4\tabcolsep) * \real{0.3667}}
  >{\raggedright\arraybackslash}p{(\columnwidth - 4\tabcolsep) * \real{0.3333}}@{}}
\toprule\noalign{}
\begin{minipage}[b]{\linewidth}\raggedright
Failure
\end{minipage} & \begin{minipage}[b]{\linewidth}\raggedright
Detection
\end{minipage} & \begin{minipage}[b]{\linewidth}\raggedright
Response
\end{minipage} \\
\midrule\noalign{}
\endhead
\bottomrule\noalign{}
\endlastfoot
Low CAN bus-off & TWAI TEC \textgreater{} 255 & Log, auto-recover. ESTOP if
persistent. \\
High CAN bus-off & MCP2515 error flags & Log, auto-recover. Zero setpoints until
restored. \\
CAN frame corruption & SYNTREE checksum fail & Frame silently discarded. Rolling
counter NOT incremented. \\
CAN transceiver dead & No frames from that node & Heartbeat timeout → ESTOP or
controlled stop. \\
\end{longtable}

\textbf{Bus-off recovery:} The CAN controller automatically attempts to recover
after 128 occurrences of 11 recessive bits (standard CAN spec). During bus-off,
the node cannot transmit. If recovery fails, the heartbeat timeout on the peer
node triggers.

\subsection{4.6 Actuator Faults}\label{actuator-faults}

\textbf{EPS-C (Steering):} \textbar{} Fault \textbar{} Detection \textbar{}
Response \textbar{}
\textbar-------\textbar-----------\textbar----------\textbar{} \textbar{} Motor
failure \textbar{} Following error \textgreater5° for \textgreater300ms
\textbar{} ESTOP \textbar{} \textbar{} Encoder failure \textbar{}
\texttt{SES\_INF\_Angle\_Status\ !=\ 1} \textbar{} STEER\_FAULT, stop
transmitting \textbar{} \textbar{} Comm timeout \textbar{} No \texttt{0x169} for
\textgreater20ms \textbar{} EPS-C internal fault lock \textbar{} \textbar{} Boot
sync timeout \textbar{} No \texttt{0x201} within 5s \textbar{} STEER\_FAULT;
rider retry via START short-press \textbar{}

\textbf{SEB (Brake):} \textbar{} Fault \textbar{} Detection \textbar{} Response
\textbar{} \textbar-------\textbar-----------\textbar----------\textbar{}
\textbar{} Comm timeout \textbar{} No \texttt{0x7B9} for \textgreater20ms
\textbar{} SEB internal fault lock \textbar{} \textbar{} Boot sync timeout
\textbar{} No \texttt{0x721} within 2s \textbar{} \textbf{BRAKE\_DEGRADED} ---
transmits 0x7B9 with lever-based stroke defaults (0mm released, 15mm pressed)
without sync. Recovers to BRAKE\_ACTIVE when first valid 0x721 arrives. Brake
lever remains functional at all times. \textbar{} \textbar{} Checksum failure
\textbar{} SEB rejects frame \textbar{} Frame dropped; counter still increments
\textbar{} \textbar{} Physical brake lever \textbar{} Always available --- GPIO2
→ SYS → CAN 0x7B9. The brake is \textbf{by-wire} (electro-hydraulic SEB, no
mechanical master cylinder). Lever availability depends on SYS continuing to
transmit 0x7B9. BRAKE\_DEGRADED guarantees transmission continues even without
boot sync. \textbar{}

\subsection{4.7 Electrical Faults}\label{electrical-faults}

\begin{longtable}[]{@{}
  >{\raggedright\arraybackslash}p{(\columnwidth - 4\tabcolsep) * \real{0.2500}}
  >{\raggedright\arraybackslash}p{(\columnwidth - 4\tabcolsep) * \real{0.3929}}
  >{\raggedright\arraybackslash}p{(\columnwidth - 4\tabcolsep) * \real{0.3571}}@{}}
\toprule\noalign{}
\begin{minipage}[b]{\linewidth}\raggedright
Fault
\end{minipage} & \begin{minipage}[b]{\linewidth}\raggedright
Protection
\end{minipage} & \begin{minipage}[b]{\linewidth}\raggedright
Response
\end{minipage} \\
\midrule\noalign{}
\endhead
\bottomrule\noalign{}
\endlastfoot
72V short to chassis & TLP281 galvanic isolation (4kV), TVS diodes, 1A fuses &
Fuse blows (\textless1ms), TVS clamps transient \\
72V on CAN bus & SN65HVD230 ±25V common-mode rating exceeded & Transceiver
destruction. Both CAN buses are physically separate --- a fault on one bus
cannot propagate. \\
12V rail short & DC-DC overcurrent protection + 12V relay cut & DC-DC shuts
down. 12V relay opens. \\
Load dump (motor regen) & SMCJ90CA TVS on all 72V lines & Clamps to 146V \\
Reverse polarity & Series Schottky on ESP32 12V input & Blocks reverse
current \\
ESD strike & TVS diodes + TLP281 isolation & Clamped to safe voltage \\
\end{longtable}

\subsection{4.8 Firmware Faults}\label{firmware-faults}

\begin{longtable}[]{@{}
  >{\raggedright\arraybackslash}p{(\columnwidth - 4\tabcolsep) * \real{0.2500}}
  >{\raggedright\arraybackslash}p{(\columnwidth - 4\tabcolsep) * \real{0.3929}}
  >{\raggedright\arraybackslash}p{(\columnwidth - 4\tabcolsep) * \real{0.3571}}@{}}
\toprule\noalign{}
\begin{minipage}[b]{\linewidth}\raggedright
Fault
\end{minipage} & \begin{minipage}[b]{\linewidth}\raggedright
Detection
\end{minipage} & \begin{minipage}[b]{\linewidth}\raggedright
Response
\end{minipage} \\
\midrule\noalign{}
\endhead
\bottomrule\noalign{}
\endlastfoot
Stack overflow & FreeRTOS stack canary & Task self-deletes, log error \\
Priority inversion & Preemptive scheduler + priority inheritance & Handled by
FreeRTOS \\
Deadlock & External watchdog timeout & Hardware reset within 100ms \\
Infinite ISR & External watchdog (independent oscillator) & Hardware reset ---
internal WDT would also be stuck \\
Flash corruption & Bootloader checksum & Boot into safe recovery mode \\
\end{longtable}

\begin{center}\rule{0.5\linewidth}{0.5pt}\end{center}

\section{5. EGAS 3-Level Motor Safety
Architecture}\label{egas-3-level-motor-safety-architecture}

The separation of MTR (STM32) from SYS (ESP32-S3) follows the EGAS 3-level
electronic throttle monitoring concept (ISO 26262 ASIL-C):

\begin{verbatim}
Level 3: Hardware — ESTOP button wired direct to both MTR and SYS
         TPS3850 external watchdog on each MCU. No software, no CAN.
         ESTOP press → MTR cuts throttle + gear instantly (local).

Level 2: Function Monitor — SYS ESP32-S3
         Monitors MTR via CAN: compares 0x204 setpoint vs 0x206 feedback.
         Mismatch > threshold → CAN 0x001 ESTOP.
         Also handles QM body functions (lights, DCDC, indicators).

Level 1: Function Controller — MTR STM32
         Normal actuation: reads sensors, drives MCP4725 DAC + gear relays.
         MANUAL: pass-through from grip/gear. AUTO: follows CAN 0x204.
         No wireless, no OS, minimal attack surface.
\end{verbatim}

\begin{longtable}[]{@{}
  >{\raggedright\arraybackslash}p{(\columnwidth - 2\tabcolsep) * \real{0.4231}}
  >{\raggedright\arraybackslash}p{(\columnwidth - 2\tabcolsep) * \real{0.5769}}@{}}
\toprule\noalign{}
\begin{minipage}[b]{\linewidth}\raggedright
Principle
\end{minipage} & \begin{minipage}[b]{\linewidth}\raggedright
Implementation
\end{minipage} \\
\midrule\noalign{}
\endhead
\bottomrule\noalign{}
\endlastfoot
\textbf{Freedom from interference} & Separate MCUs --- a SYS crash/hang cannot
block motor kill \\
\textbf{ASIL decomposition} & MTR (QM level sensor reads) + SYS monitor (ASIL-B
comparison) → combined ASIL-C \\
\textbf{Independent safe state path} & ESTOP button wired to both MCUs --- MTR
cuts throttle locally, zero CAN delay \\
\textbf{Diverse monitoring} & SYS compares commanded speed vs actual from 0x206
--- mismatch → ESTOP \\
\end{longtable}

\begin{center}\rule{0.5\linewidth}{0.5pt}\end{center}

\section{6. Emergency Response Matrix}\label{emergency-response-matrix}

\begin{longtable}[]{@{}
  >{\raggedright\arraybackslash}p{(\columnwidth - 16\tabcolsep) * \real{0.1406}}
  >{\raggedright\arraybackslash}p{(\columnwidth - 16\tabcolsep) * \real{0.1094}}
  >{\raggedright\arraybackslash}p{(\columnwidth - 16\tabcolsep) * \real{0.0938}}
  >{\raggedright\arraybackslash}p{(\columnwidth - 16\tabcolsep) * \real{0.1094}}
  >{\raggedright\arraybackslash}p{(\columnwidth - 16\tabcolsep) * \real{0.1562}}
  >{\raggedright\arraybackslash}p{(\columnwidth - 16\tabcolsep) * \real{0.0938}}
  >{\raggedright\arraybackslash}p{(\columnwidth - 16\tabcolsep) * \real{0.0781}}
  >{\raggedright\arraybackslash}p{(\columnwidth - 16\tabcolsep) * \real{0.1250}}
  >{\raggedright\arraybackslash}p{(\columnwidth - 16\tabcolsep) * \real{0.0938}}@{}}
\toprule\noalign{}
\begin{minipage}[b]{\linewidth}\raggedright
Trigger
\end{minipage} & \begin{minipage}[b]{\linewidth}\raggedright
Speed
\end{minipage} & \begin{minipage}[b]{\linewidth}\raggedright
Gear
\end{minipage} & \begin{minipage}[b]{\linewidth}\raggedright
Brake
\end{minipage} & \begin{minipage}[b]{\linewidth}\raggedright
Steering
\end{minipage} & \begin{minipage}[b]{\linewidth}\raggedright
DCDC
\end{minipage} & \begin{minipage}[b]{\linewidth}\raggedright
12V
\end{minipage} & \begin{minipage}[b]{\linewidth}\raggedright
Lights
\end{minipage} & \begin{minipage}[b]{\linewidth}\raggedright
Exit
\end{minipage} \\
\midrule\noalign{}
\endhead
\bottomrule\noalign{}
\endlastfoot
\textbf{ESTOP button} & 0V & N & Max & Obstacle: hold (dyn-clamped)→silent.
Non-obstacle: ramp→0°. & ON (MCU power) & OFF (accessories) & Brake ON & START
btn or MODE long-press → MANUAL \\
\textbf{CAN 0x001} & 0V & N & Max & Same as button & ON & OFF & Brake ON & START
btn or MODE long-press → MANUAL \\
\textbf{RT heartbeat lost} & 0V & N & Max & Same as button & ON & OFF & Brake ON
& START btn or MODE long-press → MANUAL \\
\textbf{SYS heartbeat lost} & 0V (MTR kills locally) & N (MTR cuts locally) &
Max (RT takes over 0x7B9) & EPS-C timeout* (RT still alive in AUTO) & ON (RT
alive) & OFF & Brake ON & SYS reboot → re-sync; MANUAL mode \\
\textbf{Steering follow err} & 0V & N & Max & Obstacle: hold
(dyn-clamped)→silent. Non-obstacle: ramp→0°. & ON & OFF & Brake ON & START btn
or MODE long-press → MANUAL \\
\textbf{Jetson heartbeat lost} & 0 mm/s & N & Moderate brake (2 MPa via 0x205) &
Stop 0x169 & ON & ON & Brake ON & Auto-recover on resume; MANUAL mode \\
\textbf{0x300 stale (500ms)} & 0 mm/s (coast) & N & Lever only & Stop 0x169 & ON
& ON & Normal & Auto-recover on resume \\
\textbf{0x204 stale (200ms)} & 0 mm/s & N & Lever only & No change & ON & ON &
Normal & Auto-recover on resume \\
\textbf{Watchdog reset} & 0V (MTR kills locally) & N (MTR cuts locally) & SEB
timeout* & EPS-C timeout* & OFF & OFF & All OFF & Reboot → MANUAL \\
\textbf{CAN bus-off} & Depends & Depends & Depends & Depends & Depends & Depends
& Depends & Auto-recover or ESTOP if persistent \\
\end{longtable}

\begin{quote}
* \textbf{Watchdog reset brake window (critical):} When SYS resets, MTR
maintains motor kill + gear neutral independently (separate MCU). However, SEB
is commanded exclusively by SYS via CAN 0x7B9. After 20ms without a command
frame, SEB enters internal comm-fault --- behavior is \textbf{empirically
unverified} (hold or release). If SEB holds pressure: \textasciitilde20ms
window, safe. If SEB releases: \textasciitilde2.5s window until SYS reboots and
brake LBS completes. \textbf{Must be tested before road operation.} A hardware
brake-hold relay gated by the TPS3850 RST line is recommended regardless, to
make brake behavior during MCU reset deterministic. See
{[}{[}emergency-safety-analysis{]}{]} §3 for full analysis.
\end{quote}

\begin{center}\rule{0.5\linewidth}{0.5pt}\end{center}

\section{7. Rider's Guide to Emergency
Situations}\label{riders-guide-to-emergency-situations}

\subsection{7.1 Understanding the dashboard}\label{understanding-the-dashboard}

\begin{longtable}[]{@{}
  >{\raggedright\arraybackslash}p{(\columnwidth - 2\tabcolsep) * \real{0.5714}}
  >{\raggedright\arraybackslash}p{(\columnwidth - 2\tabcolsep) * \real{0.4286}}@{}}
\toprule\noalign{}
\begin{minipage}[b]{\linewidth}\raggedright
Indicators
\end{minipage} & \begin{minipage}[b]{\linewidth}\raggedright
Meaning
\end{minipage} \\
\midrule\noalign{}
\endhead
\bottomrule\noalign{}
\endlastfoot
AUTO bulb ON, MANUAL bulb OFF & AUTO mode --- Jetson is driving \\
MANUAL bulb ON, AUTO bulb OFF & MANUAL mode --- rider is driving \\
Both bulbs OFF + \textbf{brake light ON} & ESTOP --- vehicle is
emergency-stopped, full brake engaged. Brake light is powered from the always-on
DC-DC rail and works even though the accessory relay is cut. \\
Everything OFF (brake light OFF, mode bulbs OFF) & Vehicle is \textbf{powered
off} --- normal state when key is off \\
AUTO bulb blinking 2 Hz & Degraded steering --- MANUAL mode only, AUTO locked
out \\
\end{longtable}

\textbf{How to tell ESTOP from power-off:} During ESTOP, the brake light is
\textbf{ON} (powered independently). When the vehicle is powered off, all lights
including the brake light are OFF. This is the primary visual distinction.

\subsection{7.2 If the vehicle enters ESTOP}\label{if-the-vehicle-enters-estop}

\begin{enumerate}
\def\labelenumi{\arabic{enumi}.}
\tightlist
\item
  \textbf{Stay seated and hold the handlebars.} The steering will either center
  itself or hold position.
\item
  \textbf{The brake will engage automatically} --- the vehicle will decelerate.
\item
  \textbf{Do NOT press the MODE button} --- it is ignored in ESTOP.
\item
  \textbf{To recover:} Press the \textbf{green START button} (GPIO32). The
  vehicle will enter MANUAL mode.
\item
  \textbf{Check surroundings} before releasing the brake lever and riding.
\item
  \textbf{If START button doesn't work:} Power-cycle the vehicle (key switch).
  This always starts in MANUAL mode.
\end{enumerate}

\subsection{7.3 If the vehicle behaves unexpectedly in
AUTO}\label{if-the-vehicle-behaves-unexpectedly-in-auto}

\begin{enumerate}
\def\labelenumi{\arabic{enumi}.}
\tightlist
\item
  \textbf{Press the ESTOP button} (big red mushroom). This is always the fastest
  path to stop.
\item
  \textbf{Or squeeze the brake lever.} This engages the brake directly (GPIO2 →
  SYS → SEB).
\item
  \textbf{Or switch to MANUAL} via the MODE button (if system is still
  responsive).
\item
  The steering has multiple safety clamps --- it cannot command dangerous angles
  even if Jetson malfunctions.
\end{enumerate}

\subsection{7.4 After a watchdog reset}\label{after-a-watchdog-reset}

If the system reboots unexpectedly (watchdog fired): 1. The vehicle starts in
\textbf{MANUAL mode}. 2. Steering and brake will re-sync via Listen Before
Speaking. 3. The reset reason is logged (\texttt{esp\_reset\_reason()}). 4. If
resets recur, stop riding and investigate.

\begin{center}\rule{0.5\linewidth}{0.5pt}\end{center}

\section{8. Testing the Emergency System}\label{testing-the-emergency-system}

\subsection{8.1 ESTOP Button Test}\label{estop-button-test}

\begin{enumerate}
\def\labelenumi{\arabic{enumi}.}
\tightlist
\item
  Vehicle stationary, MANUAL mode.
\item
  Press ESTOP button.
\item
  Verify: brake light ON, mode indicators OFF, throttle grip produces no
  response, gear selector produces no response.
\item
  Press START button → verify transition to MANUAL mode.
\item
  Verify: throttle and gear respond again.
\end{enumerate}

\subsection{8.2 CAN ESTOP Test}\label{can-estop-test}

\begin{enumerate}
\def\labelenumi{\arabic{enumi}.}
\tightlist
\item
  Vehicle stationary, on stands (wheels off ground).
\item
  Trigger CAN \texttt{0x001} from a CAN analyzer tool.
\item
  Verify same ESTOP behavior as physical button.
\end{enumerate}

\subsection{8.3 Heartbeat Timeout Test}\label{heartbeat-timeout-test}

\begin{enumerate}
\def\labelenumi{\arabic{enumi}.}
\tightlist
\item
  Vehicle stationary, on stands, AUTO mode.
\item
  Disconnect SYS from CAN bus (unplug SN65HVD230).
\item
  Verify: within 1000ms, RT detects SYS heartbeat loss and broadcasts CAN
  \texttt{0x001}.
\item
  Verify: ESTOP behavior activates.
\item
  Reconnect and restart.
\end{enumerate}

\subsection{8.4 Steering Following Error
Test}\label{steering-following-error-test}

\begin{enumerate}
\def\labelenumi{\arabic{enumi}.}
\tightlist
\item
  Vehicle stationary, on stands, AUTO mode.
\item
  Mechanically block the steering linkage.
\item
  Command a steering angle change via Jetson (or simulated CAN).
\item
  Verify: within 300ms of exceeding 5° error, ESTOP triggers.
\end{enumerate}

\subsection{8.5 Command Staleness Test}\label{command-staleness-test}

\begin{enumerate}
\def\labelenumi{\arabic{enumi}.}
\tightlist
\item
  Vehicle stationary, on stands, AUTO mode.
\item
  Stop Jetson ROS 2 bridge (or disconnect high CAN).
\item
  Verify: within 500ms, RT zeros \texttt{0x204} speed and stops \texttt{0x169}.
\item
  Verify: SYS detects \texttt{0x204} staleness within 200ms, forces speed=0 and
  gear=N.
\end{enumerate}

\subsection{8.6 External Watchdog Test}\label{external-watchdog-test}

\begin{enumerate}
\def\labelenumi{\arabic{enumi}.}
\tightlist
\item
  Comment out the WDT toggle in \texttt{safety\_task} (SYS) or
  \texttt{control\_task} (RT).
\item
  Flash and run.
\item
  Verify: MCU resets within \textasciitilde100ms.
\item
  Verify: during reset window, motor controller sees 0V throttle.
\item
  Restore toggle and re-flash.
\end{enumerate}

\subsection{8.7 Dynamic Angle Clamp Test}\label{dynamic-angle-clamp-test}

\begin{enumerate}
\def\labelenumi{\arabic{enumi}.}
\tightlist
\item
  Vehicle stationary, on stands, AUTO mode.
\item
  Send a CAN \texttt{0x300} command with speed=25 km/h and yaw\_rate requesting
  40° steering.
\item
  Verify: RT clamps steering to \textasciitilde5° (not 40°).
\item
  Verify: \texttt{0x169} transmitted angle is ≤5°.
\end{enumerate}

\begin{center}\rule{0.5\linewidth}{0.5pt}\end{center}

\section{9. Design Principles}\label{design-principles}

\begin{enumerate}
\def\labelenumi{\arabic{enumi}.}
\tightlist
\item
  \textbf{ESTOP bypasses queues.} The safety task preempts and writes directly
  to actuators --- no queue delay.
\item
  \textbf{ESTOP is an absorbing state.} Once entered, only deliberate human
  action (START button or power-cycle) exits.
\item
  \textbf{ESTOP exit goes to MANUAL, never AUTO.} The rider resumes direct
  control after any emergency.
\item
  \textbf{NC wiring for all safety inputs.} Cut wires and disconnected plugs
  read as ESTOP, not ``everything fine.''
\item
  \textbf{Threshold + duration for all fault checks.} Single-sample glitches
  don't trigger ESTOP --- faults must persist.
\item
  \textbf{Fail-safe by default.} De-energized relays = Neutral gear. Zero
  throttle voltage = motor stop. NC ESTOP button = pressed when disconnected.
\item
  \textbf{Fail-visible where possible.} Brake light always illuminates when any
  braking source is active. Both mode bulbs OFF = ESTOP.
\item
  \textbf{Independent safety layers.} Hardware (ESTOP button, watchdog IC) → CAN
  (0x001, heartbeats) → software (staleness, following error, clamps). No single
  failure defeats them all.
\end{enumerate}

\begin{center}\rule{0.5\linewidth}{0.5pt}\end{center}

\section{10. Related Documents}\label{related-documents}

\begin{longtable}[]{@{}
  >{\raggedright\arraybackslash}p{(\columnwidth - 2\tabcolsep) * \real{0.5263}}
  >{\raggedright\arraybackslash}p{(\columnwidth - 2\tabcolsep) * \real{0.4737}}@{}}
\toprule\noalign{}
\begin{minipage}[b]{\linewidth}\raggedright
Document
\end{minipage} & \begin{minipage}[b]{\linewidth}\raggedright
Content
\end{minipage} \\
\midrule\noalign{}
\endhead
\bottomrule\noalign{}
\endlastfoot
{[}{[}defense-in-depth-safety{]}{]} & Detailed breakdown of each safety layer \\
{[}{[}external-watchdog{]}{]} & TPS3850 watchdog IC --- timeout, safe state,
testing \\
{[}{[}high-voltage-isolation{]}{]} & 72V galvanic isolation --- TLP281, relays,
fuses, TVS \\
{[}{[}distributed-architecture{]}{]} & Three-node rationale ---
Jetson/RT/SYS/MTR split \\
{[}{[}listen-before-speaking{]}{]} & CAN actuator safe bootstrapping after
ESTOP/watchdog reset \\
{[}{[}syntree-security-protocol{]}{]} & Rolling counter + XOR checksum for
SYNTREE actuators \\
{[}{[}wiring{]}{]} & Pin-to-pin wiring including ESTOP button, brake lever,
START button \\
\href{../architecture.md}{architecture.md} §3 & Mode state machine (MANUAL →
AUTO → ESTOP) \\
\href{../architecture.md}{architecture.md} §6.1 & EGAS 3-level motor safety
architecture \\
\href{../architecture.md}{architecture.md} §7.6 & Steering safety mechanisms and
ESTOP behavior \\
\href{../architecture.md}{architecture.md} §8.6 & Brake control, heartbeats,
external watchdog, physical controls \\
\href{../can-dictionary.md}{can-dictionary.md} §0x001 & SAFETY\_ESTOP frame
definition \\
\end{longtable}

\begin{center}\rule{0.5\linewidth}{0.5pt}\end{center}

\emph{See also: {[}{[}defense-in-depth-safety{]}{]} for layered safety approach,
{[}{[}external-watchdog{]}{]} for hardware watchdog,
{[}{[}high-voltage-isolation{]}{]} for electrical fault protection,
\texttt{architecture.md} §6 for EGAS 3-level safety, §3 for mode state machine.}


\chapter{External Watchdog IC}\label{external-watchdog-ic}

The ESP32-S3 has a built-in watchdog timer, but it runs on the same silicon as
the application firmware. If the chip enters a latch-up state, browns out, or
the crystal oscillator fails, the internal watchdog may not fire.

An \textbf{external watchdog IC} is an independent hardware timer on a separate
chip. The firmware must toggle a GPIO periodically to ``pet'' it. If the
toggling stops, the watchdog asserts a reset signal to the MCU.

\begin{center}\rule{0.5\linewidth}{0.5pt}\end{center}

\section{Why an external watchdog?}\label{why-an-external-watchdog}

\begin{longtable}[]{@{}
  >{\raggedright\arraybackslash}p{(\columnwidth - 4\tabcolsep) * \real{0.1857}}
  >{\centering\arraybackslash}p{(\columnwidth - 4\tabcolsep) * \real{0.4000}}
  >{\centering\arraybackslash}p{(\columnwidth - 4\tabcolsep) * \real{0.4143}}@{}}
\toprule\noalign{}
\begin{minipage}[b]{\linewidth}\raggedright
Failure mode
\end{minipage} & \begin{minipage}[b]{\linewidth}\centering
Internal watchdog catches?
\end{minipage} & \begin{minipage}[b]{\linewidth}\centering
External watchdog catches?
\end{minipage} \\
\midrule\noalign{}
\endhead
\bottomrule\noalign{}
\endlastfoot
Task deadlock (priority inversion) & ✓ (if idle task starved) & ✓ (safety task
stops toggling) \\
Infinite loop in ISR & ✗ (ISR preempts watchdog) & ✓ (safety task never runs) \\
Crystal oscillator failure & ✗ (no clock = no watchdog) & ✓ (independent RC
oscillator) \\
Flash corruption (bad code) & ✗ (watchdog is code too) & ✓ (independent
silicon) \\
Voltage brownout & ✗ (watchdog browns out too) & ✓ (separate power-on-reset) \\
Latch-up (ESD, overvoltage) & ✗ (silicon locked) & ✓ (external IC unaffected) \\
\end{longtable}

The external watchdog is the \textbf{last-resort} safety layer. If everything
else fails --- ESTOP layers, heartbeat monitoring, command staleness --- the
watchdog ensures the MCU reboots into a safe state within a bounded time.

\begin{center}\rule{0.5\linewidth}{0.5pt}\end{center}

\section{Implementation on the E-Trike}\label{implementation-on-the-e-trike}

\subsection{Hardware}\label{hardware}

The SYS ESP32-S3 uses an external watchdog IC connected to GPIO21.

\begin{verbatim}
ESP32-S3 GPIO21 ──► WDI (Watchdog Input) pin on external watchdog IC
                         │
Watchdog IC RST ────────► EN (Enable) pin on ESP32-S3
\end{verbatim}

The \texttt{safety\_task} (Core 0, priority 5, 20 Hz) toggles GPIO21 every
iteration:

\begin{verbatim}
void safety_task() {
    TickType_t last_wake = xTaskGetTickCount();
    while (true) {
        // Toggle watchdog GPIO every iteration
        gpio_set_level(WATCHDOG_GPIO, 1);
        vTaskDelay(1);  // brief pulse
        gpio_set_level(WATCHDOG_GPIO, 0);

        // ... ESTOP GPIO check, heartbeat check ...

        vTaskDelayUntil(&last_wake, pdMS_TO_TICKS(50));  // 20 Hz = 50 ms
    }
}
\end{verbatim}

\subsection{Timeout selection}\label{timeout-selection}

The watchdog timeout is configured in hardware (typically via a resistor or pin
strapping on the watchdog IC). The timeout must be:

\begin{itemize}
\tightlist
\item
  \textbf{Long enough} that a transient spike in FreeRTOS tick latency (e.g.,
  flash erase, WiFi calibration burst) doesn't trigger it.
\item
  \textbf{Short enough} that a genuine firmware hang cuts power to the motor
  before the vehicle can accelerate dangerously.
\end{itemize}

For the E-Trike: \textbf{timeout = 100 ms}. This means: - The
\texttt{safety\_task} toggles every 50 ms (20 Hz). - Two missed toggles (100 ms)
triggers reset. - At 25 km/h (7 m/s), the vehicle moves \textasciitilde0.7 m
during the timeout window.

\subsection{Startup sequencing}\label{startup-sequencing}

On reset, the watchdog IC holds the ESP32 in reset for a startup delay
(\textasciitilde200 ms) to allow the power supply to stabilize. After release,
the ESP32 boots and the \texttt{safety\_task} begins toggling within
\textasciitilde50 ms (FreeRTOS scheduler start + task creation).

The watchdog \textbf{does not arm until the first toggle} --- most watchdog ICs
start with their timer disarmed and arm on the first WDI edge. This prevents a
reset loop during boot.

\begin{center}\rule{0.5\linewidth}{0.5pt}\end{center}

\section{Safe state on watchdog reset}\label{safe-state-on-watchdog-reset}

When the watchdog fires:

\begin{enumerate}
\def\labelenumi{\arabic{enumi}.}
\tightlist
\item
  ESP32-S3 EN pin is pulled LOW → MCU resets.
\item
  All GPIOs go to their reset state (high-impedance input).
\item
  Motor controller sees 0 V on throttle (MCP4725 loses I2C, output goes to 0 V).
\item
  All gear relays de-energize (GPIOs float → transistor off → relay open).
\item
  Brake: SEB continues its internal control loop with the last received
  \texttt{0x7B9} command. After 20 ms of no new frame, SEB enters comm-fault.
  \textbf{This is a known gap --- behavior is empirically unverified.} If SEB
  holds pressure on comm-fault, the window is \textasciitilde20ms (acceptable).
  If SEB releases, the vehicle coasts without brake for \textasciitilde2.5s (SYS
  reboot + brake LBS). A hardware brake-hold relay gated by the TPS3850 RST line
  is recommended to make brake behavior deterministic during MCU reset. See
  {[}{[}emergency-safety-analysis{]}{]} §3 for full causal trace and risk
  quantification.
\item
  After \textasciitilde200 ms, the ESP32 reboots and restores control.
\end{enumerate}

\subsection{Post-reset behavior}\label{post-reset-behavior}

After a watchdog reset, the firmware must: 1. Not assume any previous state is
valid. 2. Start in MANUAL mode (safest default). 3. Re-run the full LBS sequence
for all actuators (see {[}{[}listen-before-speaking{]}{]}). 4. Log the reset
cause (ESP32 provides \texttt{esp\_reset\_reason()}).

\begin{verbatim}
void app_main() {
    esp_reset_reason_t reason = esp_reset_reason();
    if (reason == ESP_RST_PANIC || reason == ESP_RST_WDT) {
        // External watchdog likely triggered
        // Log to persistent storage if available
    }

    // ... normal startup (MANUAL mode, LBS sequences) ...
}
\end{verbatim}

\begin{center}\rule{0.5\linewidth}{0.5pt}\end{center}

\section{Common external watchdog ICs}\label{common-external-watchdog-ics}

\begin{longtable}[]{@{}lll@{}}
\toprule\noalign{}
Part & Timeout & Features \\
\midrule\noalign{}
\endhead
\bottomrule\noalign{}
\endlastfoot
TPS3823-33 & Fixed 1.6 s & 3.3 V VCC, push-pull RST, WDI \\
MAX823 & 1.6 s typical & 5-pin SOT23, very common \\
STM6823 & 1.6 s & Low power, small SOT23-5 \\
TPS3431 & Programmable & Capacitor-programmable timeout, windowed mode \\
ADM8323 & 1.6 s & Ultra-low 6 µA supply current \\
\end{longtable}

For the E-Trike, a \textbf{programmable-timeout part} (e.g., TPS3431) is
preferred --- the 100 ms target is shorter than most fixed-timeout off-the-shelf
watchdogs (which typically start at 200--400 ms minimum).

\begin{center}\rule{0.5\linewidth}{0.5pt}\end{center}

\section{Testing the watchdog}\label{testing-the-watchdog}

\begin{enumerate}
\def\labelenumi{\arabic{enumi}.}
\tightlist
\item
  \textbf{Force a hang:} Comment out the GPIO toggle in \texttt{safety\_task}.
  Confirm the MCU resets within 100 ms.
\item
  \textbf{Verify safe state during reset:} Probe the MCP4725 output with an
  oscilloscope --- it should fall to 0 V within microseconds of the reset line
  asserting.
\item
  \textbf{Verify gear relays open:} Probe the relay coil driver --- it should
  de-energize immediately.
\item
  \textbf{Measure reset-boot-recovery time:} From watchdog fire to first valid
  CAN transmission. Should be \textless500 ms.
\end{enumerate}

\begin{center}\rule{0.5\linewidth}{0.5pt}\end{center}

\section{Architecture gap note}\label{architecture-gap-note}

The current \texttt{architecture.md} does \textbf{not} describe the external
watchdog. This note fills that gap. The watchdog should be added to the
architecture's SYS ESP32-S3 hardware section (§8.8) and the RTOS task layout
table (§8.7).

\begin{center}\rule{0.5\linewidth}{0.5pt}\end{center}

\emph{Primary reference: {[}{[}emergency-system{]}{]} for the complete ESTOP
system including watchdog's role in the 8-layer defense and post-reset recovery
procedures.}

\emph{See also: {[}{[}defense-in-depth-safety{]}{]} for the layered safety
approach that the watchdog is part of, {[}{[}listen-before-speaking{]}{]} for
the LBS sequence that runs after every watchdog reset, {[}{[}architecture{]}{]}
§8 for SYS ESP32-S3 details.}


\chapter{Syntree EPS-C \& SEB --- Technical
Reference}\label{syntree-eps-c-seb-technical-reference}

Extracted and cleaned from the original Syntree product specification PDFs
(\texttt{given-documentation.pdf}, \texttt{given-documentation-ocr.pdf}). The
source PDFs are bilingual (Chinese/English) scanned documents; the OCR quality
varies. This document distills the \textbf{technically actionable} information
needed for firmware integration. Always verify byte layouts against the original
PDFs before enabling actuator output.

\begin{center}\rule{0.5\linewidth}{0.5pt}\end{center}

\section{1. EPS-C --- Electric Power Steering
Column}\label{eps-c-electric-power-steering-column}

\textbf{Product:} Syntree steer-by-wire intelligent steering system (SES)
\textbf{Version:} A/10 (2024-11-22) \textbf{Category:} EPS-C

\subsection{1.1 Technical Parameters}\label{technical-parameters}

\begin{longtable}[]{@{}lll@{}}
\toprule\noalign{}
Parameter & Value & Unit \\
\midrule\noalign{}
\endhead
\bottomrule\noalign{}
\endlastfoot
Assembly weight & 3.5--4 & kg \\
ECU weight & 0.4 & kg \\
Rated motor power & 360 & W \\
Rated motor speed & 1480 & rpm \\
Rated motor torque & 2.36 & Nm \\
Rated motor current & 52 & A \\
Rated working voltage & 12 & V \\
Max working current & 30 & A \\
ECU operating voltage & 9--15 & V \\
Debugging voltage & 13.5 & V \\
Spline teeth & 30/36 & z \\
Spline module & 0.467/0.47 & m \\
Spline pressure angle & 30.5/45 & ° \\
Max steering angle (wheel end) & ±700 & ° \\
Reduction ratio & 27 / 20.5 / 13.5 & --- \\
Max steering torque & 55 / 41 / 27 & Nm \\
Steering accuracy & ±1 & ° \\
Max steering speed & 400 & °/s \\
Steering response time & \textless100 & ms \\
Max assist torque & (see datasheet) & Nm \\
\end{longtable}

\subsection{1.2 Communication}\label{communication}

\begin{itemize}
\tightlist
\item
  \textbf{Baud rate:} 500K (default), 250K optional
\item
  \textbf{Protocol:} CAN\_SES (Syntree proprietary steer-by-wire)
\item
  \textbf{Terminal resistance:} Not populated by default on ECU
\end{itemize}

\subsection{1.3 Operating Principle}\label{operating-principle}

The VCU transmits steering angle commands via CAN. The ECU calculates and drives
the motor, which provides steering torque through a worm-gear reduction
mechanism. An angle sensor monitors steering angle and feeds it back to the ECU
via PWM. The ECU makes real-time adjustments and reports actual angle back to
the VCU via CAN.

\subsection{1.4 CAN Frames (SES Protocol)}\label{can-frames-ses-protocol}

\subsubsection{VCU → EPS-C Command (ID: 0x169, cycle: 20 ms, DLC:
8)}\label{vcu-eps-c-command-id-0x169-cycle-20-ms-dlc-8}

\begin{longtable}[]{@{}ll@{}}
\toprule\noalign{}
Signal & Description \\
\midrule\noalign{}
\endhead
\bottomrule\noalign{}
\endlastfoot
\texttt{VCU\_SES\_Alignment\_Enable} & Alignment/zero-calibration enable \\
\texttt{VCU\_SES\_Control\_Enable} & Steering control enable \\
\texttt{VCU\_SES\_Tgt\_StrAngle} & Target steering angle \\
\texttt{VCU\_SES\_Tgt\_StrAngleSpd} & Target steering angle speed \\
\texttt{VCU\_SES\_RollCnt\_Enable} & Rolling counter enable \\
\texttt{VCU\_SES\_CheckSum\_Enable} & Checksum enable \\
\texttt{VCU\_SES\_RollCnt} & Rolling counter (liveliness signal) \\
\texttt{VCU\_Veh\_Spd\_Value} & Vehicle speed value \\
\texttt{VCU\_SES\_CheckSum} & Frame checksum \\
\end{longtable}

\begin{quote}
\textbf{⚠️ The exact byte layout, bit positions, scaling factors, and checksum
algorithm must be verified against the project-specific protocol document before
transmitting on hardware.} The \texttt{kSyntreeCanOutputEnabled} flag in
\texttt{config.h} defaults to \texttt{false} for this reason.
\end{quote}

\subsubsection{EPS-C → VCU Status (ID: 0x201, cycle: 10 ms, DLC:
8)}\label{eps-c-vcu-status-id-0x201-cycle-10-ms-dlc-8}

\begin{longtable}[]{@{}
  >{\raggedright\arraybackslash}p{(\columnwidth - 2\tabcolsep) * \real{0.3810}}
  >{\raggedright\arraybackslash}p{(\columnwidth - 2\tabcolsep) * \real{0.6190}}@{}}
\toprule\noalign{}
\begin{minipage}[b]{\linewidth}\raggedright
Signal
\end{minipage} & \begin{minipage}[b]{\linewidth}\raggedright
Description
\end{minipage} \\
\midrule\noalign{}
\endhead
\bottomrule\noalign{}
\endlastfoot
\texttt{SES\_INF\_Angle\_Status} & Actual angle feedback status \\
\texttt{SES\_Control\_Mode\_Status} & Control mode status (0x0 = normal, 0x1 =
fault level 1, 0x2 = fault level 2) \\
\texttt{SES\_Error\_Status} & Error/fault flags \\
\texttt{SES\_StrAngle} & Actual steering angle \\
\texttt{SES\_Tgt\_StrAngleSpd} & Actual steering angle speed \\
\texttt{SES\_SteeringWheel\_Torq} & Steering wheel torque \\
\texttt{SES\_RollCnt\_Enable\_Status} & Rolling counter valid \\
\texttt{SES\_CheckSum\_Enable\_Status} & Checksum valid \\
\texttt{SES\_RollCnt\_Status} & Rolling counter echo \\
\texttt{SES\_CheckSum\_Status} & Checksum value \\
\texttt{SES\_SW\_Version} & Software version (separate frame, 1000 ms cycle) \\
\texttt{SES\_HW\_Version} & Hardware version (separate frame) \\
\texttt{SES\_MtrCurt} & Motor current (10 ms cycle) \\
\texttt{SES\_ECU\_Temp} & ECU temperature \\
\texttt{SES\_Anl\_Volt} & Analog voltage \\
\end{longtable}

\subsection{1.5 EPS-C Fault Codes}\label{eps-c-fault-codes}

\begin{longtable}[]{@{}ll@{}}
\toprule\noalign{}
Fault & Level \\
\midrule\noalign{}
\endhead
\bottomrule\noalign{}
\endlastfoot
ECU Under-voltage & L2 \\
ECU Over-voltage & L2 \\
CAN Communication Error & L1 \\
ECU Temperature & L3 \\
Domain drive SC (short circuit) & L3 \\
Domain drive V (voltage) & L3 \\
Domain drive T (temperature) & L3 \\
Temperature Sensor Error & L3 \\
Angle Sensor Channel O/C (open circuit) & L3 \\
Angle Sensor Channel AF & L3 \\
Angle Sensor S O/C & L3 \\
Angle Sensor S AF & L3 \\
Sensor Power Error & L3 \\
Alignment Error & L1 \\
Over-angle Error & L2 \\
Steering Stall Error & L3 \\
Motor Current Error & L3 \\
Sensor O/C Error & L3 \\
Torque Sensor T1 O/C & L3 \\
Torque Sensor T2 O/C & L3 \\
Torque Sensor T1 AF & L3 \\
Torque Sensor T2 AF & L3 \\
Sent Angle Error & L1 \\
Steering Ralling Error & L3 \\
EPPROM Error & L2 \\
\end{longtable}

\begin{quote}
Fault levels: L1 = highest severity, L4 = lowest severity
\end{quote}

\subsection{1.6 Installation Notes}\label{installation-notes}

\begin{itemize}
\tightlist
\item
  ECU must be zero-calibrated when matched to a new assembly, then re-powered
\item
  After zero-calibration, the ECU + assembly pair is unique and not
  interchangeable
\item
  ECU housing must be connected to chassis ground
\item
  Low-voltage system shares common ground with the vehicle
\end{itemize}

\begin{center}\rule{0.5\linewidth}{0.5pt}\end{center}

\section{2. SEB --- Electronic Brake System}\label{seb-electronic-brake-system}

\textbf{Product:} Syntree wire-controlled braking system (SEB) \textbf{Version:}
A/19 (2024-08-02) \textbf{Category:} SEB

\subsection{2.1 Technical Parameters}\label{technical-parameters-1}

\begin{longtable}[]{@{}lll@{}}
\toprule\noalign{}
Parameter & Value & Unit \\
\midrule\noalign{}
\endhead
\bottomrule\noalign{}
\endlastfoot
Max stroke & 27 & mm \\
Stroke accuracy & ±0.5 & mm \\
Stroke response time & \textasciitilde50 & ms \\
Pressure-building response time & \textasciitilde150 & ms \\
Max oil pressure & 5 & MPa \\
Brake fluid type & DOT4 & --- \\
Operating voltage & 12 & V \\
\end{longtable}

\subsection{2.2 Operating Principle}\label{operating-principle-1}

The SEB is an electro-hydraulic brake system based on traditional hydraulic
brake architecture. It consists of a drive motor, reduction mechanism, master
cylinder, angle sensor, oil pressure sensor, and control unit. The VCU sends
brake commands via CAN. The ECU drives the motor to build hydraulic pressure in
the master cylinder, which actuates the calipers. Pressure and stroke sensors
provide feedback.

\subsection{2.3 CAN Frames (SEB Protocol)}\label{can-frames-seb-protocol}

\subsubsection{VCU → SEB Command (ID: 0x7B0, cycle: 20 ms, DLC:
8)}\label{vcu-seb-command-id-0x7b0-cycle-20-ms-dlc-8}

\begin{longtable}[]{@{}ll@{}}
\toprule\noalign{}
Signal & Description \\
\midrule\noalign{}
\endhead
\bottomrule\noalign{}
\endlastfoot
\texttt{VCU\_SEB\_Alignment\_Enable} & Alignment enable \\
\texttt{VCU\_SEB\_Control\_Enable} & Brake control enable \\
\texttt{VCU\_SEB\_Control\_Mode} & Control mode (stroke or pressure) \\
\texttt{VCU\_SEB\_AutoBrake} & Automatic brake trigger \\
\texttt{VCU\_SEB\_Stroke\_Value\_Req} & Target stroke value \\
\texttt{VCU\_SEB\_Prs\_Value\_Req} & Target oil pressure value \\
\texttt{VCU\_SEB\_RollCnt\_Enable} & Rolling counter enable \\
\texttt{VCU\_SEB\_CheckSum\_Enable} & Checksum enable \\
\texttt{VCU\_SEB\_RollCnt} & Rolling counter \\
\texttt{VCU\_SEB\_CheckSum} & Frame checksum \\
\end{longtable}

\begin{quote}
\textbf{⚠️ Same warning as EPS-C:} verify byte layout and checksum algorithm
before transmitting.
\end{quote}

\subsubsection{SEB → VCU Status (ID: 0x721, cycle: 10 ms, DLC:
8)}\label{seb-vcu-status-id-0x721-cycle-10-ms-dlc-8}

\begin{longtable}[]{@{}
  >{\raggedright\arraybackslash}p{(\columnwidth - 2\tabcolsep) * \real{0.3810}}
  >{\raggedright\arraybackslash}p{(\columnwidth - 2\tabcolsep) * \real{0.6190}}@{}}
\toprule\noalign{}
\begin{minipage}[b]{\linewidth}\raggedright
Signal
\end{minipage} & \begin{minipage}[b]{\linewidth}\raggedright
Description
\end{minipage} \\
\midrule\noalign{}
\endhead
\bottomrule\noalign{}
\endlastfoot
\texttt{SEB\_INF\_Alignment\_Status} & Alignment status feedback \\
\texttt{SEB\_Control\_Enable\_Status} & Control enable status \\
\texttt{SEB\_Control\_Mode\_Status} & Control mode feedback \\
\texttt{SEB\_AutoBrake\_Status} & Auto-brake status \\
\texttt{SEB\_Error\_Status} & Error flags (0x0 = normal, 0x1 = fault level 1,
0x2 = fault level 2) \\
\texttt{SEB\_Stroke\_Value} & Actual stroke \\
\texttt{SEB\_Prs\_Value} & Actual oil pressure \\
\texttt{SEB\_AutoBrake\_Value} & Auto-brake target value \\
\texttt{SEB\_RollCnt\_Enable\_Status} & Rolling counter valid \\
\texttt{SEB\_CheckSum\_Enable\_Status} & Checksum valid \\
\texttt{SEB\_RollCnt\_Status} & Rolling counter echo \\
\texttt{SEB\_CheckSum\_Status} & Checksum echo \\
\texttt{SEB\_SW\_Version} & Software version (1000 ms cycle) \\
\texttt{SEB\_HW\_Version} & Hardware version (1000 ms cycle) \\
\end{longtable}

\subsection{2.4 SEB Fault Codes}\label{seb-fault-codes}

\begin{longtable}[]{@{}ll@{}}
\toprule\noalign{}
Fault & Description \\
\midrule\noalign{}
\endhead
\bottomrule\noalign{}
\endlastfoot
\texttt{SEB\_ECU\_UnderVolt\_Err} & ECU under-voltage \\
\texttt{SEB\_ECU\_OverVolt\_Err} & ECU over-voltage \\
\texttt{SEB\_CanCom\_Err} & CAN communication timeout \\
\texttt{SEB\_ECU\_Temp\_Err} & ECU over-temperature \\
\texttt{SEB\_Domain\_drive\_SC\_Err} & Motor driver short circuit \\
\texttt{SEB\_Domain\_drive\_V\_Err} & Motor driver voltage fault \\
\texttt{SEB\_Domain\_drive\_T\_Err} & Motor driver over-temperature \\
\texttt{SEB\_AngleSensor\_P\_O/C\_Err} & Angle sensor primary open circuit \\
\texttt{SEB\_AngleSensor\_P\_AF\_Err} & Angle sensor primary range fault \\
\texttt{SEB\_AngleSensor\_S\_O/C\_Err} & Angle sensor secondary open circuit \\
\texttt{SEB\_AngleSensor\_S\_AF\_Err} & Angle sensor secondary range fault \\
\texttt{SEB\_NoPrsSensor\_Err} & No pressure sensor detected \\
\texttt{SEB\_Sensor\_VCL\_Err} & Sensor supply voltage fault \\
\texttt{SEB\_Alignment\_Err} & Alignment/zero-calibration error \\
\texttt{SEB\_AngleOver\_Err} & Angle exceeded limit \\
\texttt{SEB\_Str\_Stall\_Err} & Stroke stall \\
\texttt{SEB\_Mtr\_O/C\_Err} & Motor open circuit \\
\texttt{SEB\_Oil\_Err} & Hydraulic oil fault \\
\texttt{SEB\_HotOil\_Err} & Oil over-temperature \\
\texttt{SEB\_Snsr\_Value\_Err} & Sensor value out of range \\
\texttt{SEB\_Mtr\_NoBrake\_Err} & Motor fails to brake \\
\texttt{SEB\_PrsSensor\_Over\_Err} & Pressure sensor over-range \\
\texttt{SEB\_Low\_Volt\_Charging} & Low voltage charging fault \\
\end{longtable}

\subsection{2.5 Debug/Test Commands}\label{debugtest-commands}

The SEB supports test-mode commands for bleeding and calibration. These use CAN
IDs \texttt{0x789} and \texttt{0x7B9} with a 9-byte payload (extended frame,
DLC=8 with length byte preceding).

\textbf{Stroke mode test} (Byte0 = 0x02): - Byte1--Byte2: Target stroke =
\texttt{(physical\_mm\ +\ 30)\ /\ 0.05} - Max stroke: 27 mm - Example 8mm
stroke: \texttt{{[}02\ 58\ 00\ 00\ 00\ 00\ 00\ 00{]}} (with length=9 prefix)

\textbf{Oil pressure mode test} (Byte0 = 0x06): - Byte3: Target oil pressure =
\texttt{physical\_MPa\ /\ 0.05} - Max pressure: 5 MPa

\textbf{Cyclic venting} (for brake bleeding): - Alternate between extend and
retract stroke commands - 2S interval between commands, 4S per cycle - Repeat
until no bubbles at caliper bleed screw

\subsection{2.6 Installation Notes}\label{installation-notes-1}

\begin{itemize}
\tightlist
\item
  After installation, the oil circuit must be bled (non-air procedure using ECU
  stroke commands)
\item
  Oil circuit is established when stroke command produces ≥2 MPa feedback
  pressure
\item
  Use DOT4 brake fluid only
\item
  Power-on pressure maintenance: before power-off, send 4 MPa pressure command;
  mechanical self-locking maintains \textasciitilde4 MPa after power-off
\item
  Avoid sustained \textgreater5 MPa commands --- may cause motor overheating
\end{itemize}

\begin{center}\rule{0.5\linewidth}{0.5pt}\end{center}

\section{3. DCDC Converter --- CAN Protocol}\label{dcdc-converter-can-protocol}

\textbf{Document:} CAN BUS COMMUNICATION SPECIFICATION Version 1.2/1.3
\textbf{Baud rate:} 250 Kbps \textbf{Frame format:} CAN 2.0B extended (29-bit
ID), J1939-based

\subsection{3.1 Node Addresses}\label{node-addresses}

\begin{longtable}[]{@{}ll@{}}
\toprule\noalign{}
Node & Address (SA) \\
\midrule\noalign{}
\endhead
\bottomrule\noalign{}
\endlastfoot
VCU & 0x27 \\
DCDC & 0x2B \\
\end{longtable}

\subsection{3.2 Message 1: VCU → DCDC (ID: 0x10262B27, 100 ms
cycle)}\label{message-1-vcu-dcdc-id-0x10262b27-100-ms-cycle}

\begin{longtable}[]{@{}
  >{\raggedright\arraybackslash}p{(\columnwidth - 6\tabcolsep) * \real{0.2000}}
  >{\raggedright\arraybackslash}p{(\columnwidth - 6\tabcolsep) * \real{0.2667}}
  >{\raggedright\arraybackslash}p{(\columnwidth - 6\tabcolsep) * \real{0.3000}}
  >{\raggedright\arraybackslash}p{(\columnwidth - 6\tabcolsep) * \real{0.2333}}@{}}
\toprule\noalign{}
\begin{minipage}[b]{\linewidth}\raggedright
Byte
\end{minipage} & \begin{minipage}[b]{\linewidth}\raggedright
Signal
\end{minipage} & \begin{minipage}[b]{\linewidth}\raggedright
Scaling
\end{minipage} & \begin{minipage}[b]{\linewidth}\raggedright
Notes
\end{minipage} \\
\midrule\noalign{}
\endhead
\bottomrule\noalign{}
\endlastfoot
Byte0 & Control & Bit0--1: Working mode (00=disable, 01=enable). Bit2--7:
Reserved & \\
Byte1 & Max Charging Voltage Low & 0.1 V/bit, offset 0 & e.g.~Vset=140 → 14V \\
Byte2 & Max Charging Voltage High & 0.1 V/byte & \\
Byte3 & Max Charging Current Low & 0.1 A/bit, offset 0 & e.g.~Aset=200 → 20A \\
Byte4 & Max Charging Current High & 0.1 A/byte & \\
Byte5 & Reserved & 0xFF & \\
Byte6 & Reserved & 0xFF & \\
Byte7 & Reset control & Bit0--1: Reset (00=no reset, 01=reset). Bit2--7:
Reserved & \\
\end{longtable}

\begin{itemize}
\tightlist
\item
  If VCU sends enable + zero voltage/current → DCDC outputs system default
\item
  If no message received within 6s → DCDC enters communication error state and
  stops output
\end{itemize}

\subsection{3.3 Message 2: DCDC → VCU (ID: 0x18F8622B, 100 ms
cycle)}\label{message-2-dcdc-vcu-id-0x18f8622b-100-ms-cycle}

\begin{longtable}[]{@{}llll@{}}
\toprule\noalign{}
Byte & Signal & Scaling & Notes \\
\midrule\noalign{}
\endhead
\bottomrule\noalign{}
\endlastfoot
Byte0 & System state & Table 1-1 (see §3.4) & \\
Byte1 & Temperature & 1°C/bit, offset 40°C & e.g.~value 120 → 80°C \\
Byte2 & Output Voltage Low & 0.05 V/bit, offset 0 & \\
Byte3 & Output Voltage High & 0.05 V/byte & \\
Byte4 & Output Current Low & 0.05 A/bit, offset 0 & \\
Byte5 & Output Current High & 0.05 A/byte & \\
Byte6 & Error flags & Table 1-2 (see §3.5) & \\
Byte7 & Version info & Software version & \\
\end{longtable}

\subsection{3.4 DCDC System State (Byte0)}\label{dcdc-system-state-byte0}

\begin{longtable}[]{@{}lll@{}}
\toprule\noalign{}
Bits & Value & State \\
\midrule\noalign{}
\endhead
\bottomrule\noalign{}
\endlastfoot
Bit0--2 & 000 & Stop \\
& 001 & Charging \\
& 010 & Charging complete \\
& 011 & Reserved \\
Bit3--4 & 00 & Fault Level 1 (highest) \\
& 01 & Fault Level 2 \\
& 10 & Fault Level 3 \\
& 11 & Fault Level 4 (lowest) \\
Bit5--7 & 000 & Ready \\
& 100 & Power Up \\
& 101 & Error \\
& 111 & Diagnostic/Calibration \\
\end{longtable}

\subsection{3.5 DCDC Fault Codes}\label{dcdc-fault-codes}

\begin{longtable}[]{@{}ll@{}}
\toprule\noalign{}
Code & Fault \\
\midrule\noalign{}
\endhead
\bottomrule\noalign{}
\endlastfoot
11 & Output over-voltage \\
12 & Output over-current \\
101 & Input over-voltage \\
102 & Input under-voltage \\
103 & Over-temperature protection \\
151 & CAN communication receive timeout \\
152 & Output under-voltage \\
\end{longtable}

\begin{center}\rule{0.5\linewidth}{0.5pt}\end{center}

\section{4. Key Integration Notes}\label{key-integration-notes}

\begin{enumerate}
\def\labelenumi{\arabic{enumi}.}
\item
  \textbf{Both EPS-C and SEB use proprietary CAN protocols with checksums and
  rolling counters.} Do not transmit actuator commands until the byte layout,
  scaling, and checksum algorithm are confirmed against the project-specific
  protocol document from Syntree.
\item
  \textbf{EPS-C requires zero-calibration} when paired with a new assembly. This
  is a one-time procedure, likely triggered via the \texttt{Alignment\_Enable}
  flag in the command frame.
\item
  \textbf{SEB requires brake bleeding} after installation. The ECU supports
  automated bleeding via stroke-mode test commands.
\item
  \textbf{Both ECUs have CAN termination disabled by default.} External 120Ω
  resistors must be installed at the physical ends of each CAN bus.
\item
  \textbf{EPS-C and SEB are on the private CAN bus} (TWAI1 on the ESP32-S3).
  They must never be exposed to the public Jetson CAN bus.
\item
  \textbf{VCU address 0x27} appears in the DCDC spec for J1939 addressing. The
  EPS-C and SEB protocols may use different addressing --- verify with Syntree.
\item
  \textbf{Fault handling:} Both ECUs report faults via status frames. Level-1
  faults (L1) are most severe. The ESP32-S3 should monitor
  \texttt{Error\_Status} and trigger ESTOP on critical faults.
\end{enumerate}

\begin{center}\rule{0.5\linewidth}{0.5pt}\end{center}

\emph{Source PDFs: \texttt{given-documentation.pdf} (58 pages, Syntree EPS-C +
DCDC CAN spec), \texttt{given-documentation-ocr.pdf} (58 pages, alternate scan
with higher-quality text layer). Both preserved in the project root for
byte-layout verification.}


\chapter{High-Voltage Isolation \&
Protection}\label{high-voltage-isolation-protection}

The E-Trike has a \textbf{72 V traction battery} and a \textbf{12 V accessory
rail}. The ESP32-S3 MCUs run at \textbf{3.3 V}. These voltage domains must never
meet --- a single fault connecting 72 V to an ESP32 GPIO destroys the MCU and
potentially everything on the CAN bus.

We use three strategies: \textbf{galvanic isolation} for input sensing,
\textbf{mechanical relay isolation} for output switching, and \textbf{protection
circuits} on every line that carries high voltage.

\begin{center}\rule{0.5\linewidth}{0.5pt}\end{center}

\section{Voltage domains}\label{voltage-domains}

\begin{verbatim}
72 V Traction Battery
  ├── Motor controller (72 V direct)
  ├── Gear selector lines (72 V signaling)
  ├── DC-DC converter (72V → 12V)
  │
  └─12 V Accessory Rail
      ├── Signal lights (12 V)
      ├── Headlight (12 V)
      ├── Mode indicator LEDs (12 V)
      ├── Brake light (12 V)
      └── ESP32-S3 (via 3.3 V LDO regulator)
          └── CAN bus (3.3 V logic, 5 V transceiver)
\end{verbatim}

\begin{center}\rule{0.5\linewidth}{0.5pt}\end{center}

\section{Galvanic isolation --- TLP281 optoisolators
(input)}\label{galvanic-isolation-tlp281-optoisolators-input}

\textbf{Problem:} The gear selector uses 72 V lines. When a gear position is
selected, 72 V appears on the corresponding wire. The ESP32 GPIOs are 3.3 V
tolerant only.

\textbf{Solution:} TLP281 optoisolators provide galvanic (optical) isolation
between the 72 V domain and the 3.3 V domain.

\begin{verbatim}
72 V gear wire ──┬── R1 (current limit) ──┬── TLP281 LED (anode)
                 │                         │
                 └── R2 (voltage divider) ─┴── TLP281 LED (cathode) ── GND_72V

                                                  │  Optical barrier
                                                  │  (4 kV isolation)

ESP32 GPIO ◀── R3 (pull-up to 3.3V) ──┬── TLP281 phototransistor (collector)
                                       │
                                      GND_3V3
\end{verbatim}

\textbf{How it works:} 1. 72 V present → LED emits light → phototransistor
conducts → GPIO reads LOW. 2. 72 V absent → LED off → phototransistor open →
GPIO reads HIGH (pull-up).

Note the logic inversion: HIGH on GPIO means ``no gear signal.'' This is handled
in the \texttt{gear} task's FSM.

\textbf{Key properties of TLP281:} - Isolation voltage: 2500 Vrms minimum (far
exceeds 72 VDC requirement). - CTR (Current Transfer Ratio): 100--600\% at 5 mA
LED current. - The input resistor network is sized so that 72 V produces
\textasciitilde5 mA through the LED (within the device's absolute maximum rating
and optimal CTR range).

\begin{center}\rule{0.5\linewidth}{0.5pt}\end{center}

\section{Mechanical relay isolation
(output)}\label{mechanical-relay-isolation-output}

\textbf{Problem:} SYS ESP32 must switch 72 V to the motor controller's gear
inputs (D, S, R lines). These lines carry full traction voltage.

\textbf{Solution:} Mechanical relays driven by ESP32 GPIOs via a transistor
driver.

\begin{verbatim}
ESP32 GPIO ── R1 ── NPN transistor base
                     │
                     ├── Collector ── Relay coil ── 12V
                     │
                     └── Emitter ── GND

Relay contacts (NO):
  72V bus ──┬── Common ── NO contact ──┬── Gear input on motor controller
             │                          │
             └── 1A fuse ───────────────┘
\end{verbatim}

\textbf{Why relays instead of MOSFETs:} - Galvanic isolation: relay coil and
contacts are physically separated. - Fail-safe: when coil is de-energized (ESP32
off, crashed, or ESTOP), contacts open → no gear signal. - 72 VDC switching
doesn't require complex gate drive or bootstrap circuits. - Relay coil runs on
12 V (via transistor), not 72 V.

\textbf{Why not SSRs (solid-state relays):} - Mechanical relays are simpler to
debug (audible click). - Lower cost for low-frequency switching (gear changes
are rare, not PWM). - No leakage current in OFF state.

\begin{center}\rule{0.5\linewidth}{0.5pt}\end{center}

\section{Protection circuits --- 1A fuse + SMCJ90CA
TVS}\label{protection-circuits-1a-fuse-smcj90ca-tvs}

Every line that carries or connects to 72 V has a protection circuit:

\begin{verbatim}
72V source ──┬── 1A fast-blow fuse ──┬── Load (relay contact, gear wire, etc.)
              │                       │
              │                       ├── SMCJ90CA TVS diode ── GND_72V
              │                       │
              └───────────────────────┘
\end{verbatim}

\subsection{1A fast-blow fuse}\label{a-fast-blow-fuse}

\begin{itemize}
\tightlist
\item
  \textbf{Purpose:} Overcurrent protection. If a gear line shorts to chassis or
  another wire, the fuse blows before the wire melts or the battery is damaged.
\item
  \textbf{Rating:} 1A fast-blow. Gear signaling draws microamps (into the motor
  controller's high-impedance input) --- any current above 1A indicates a short.
\item
  \textbf{One per line:} Each gear line (D, S, R) has its own fuse. A short on
  one gear doesn't take down the others.
\end{itemize}

\subsection{SMCJ90CA bidirectional TVS
diode}\label{smcj90ca-bidirectional-tvs-diode}

\begin{itemize}
\tightlist
\item
  \textbf{Purpose:} Transient voltage suppression. Protects against:

  \begin{itemize}
  \tightlist
  \item
    \textbf{Load dump:} When the motor controller suddenly disconnects from the
    battery under load, the wiring inductance generates a high-voltage spike.
  \item
    \textbf{ESD:} Electrostatic discharge during handling or from road static.
  \item
    \textbf{Switching transients:} Relay contact bounce and motor controller PWM
    noise.
  \end{itemize}
\item
  \textbf{Rating:} SMCJ90CA --- 90 V standoff, 146 V clamping, 1500 W peak pulse
  power.

  \begin{itemize}
  \tightlist
  \item
    Standoff voltage (90 V) \textgreater{} nominal bus voltage (72 V) --- diode
    doesn't conduct in normal operation.
  \item
    Clamping voltage (146 V) \textless{} breakdown voltage of adjacent
    components.
  \end{itemize}
\item
  \textbf{Bidirectional:} Protects against both positive and negative
  transients. The `C' suffix = bidirectional. The `A' suffix = 5\% tolerance on
  breakdown voltage.
\end{itemize}

\subsection{Why both?}\label{why-both}

\begin{itemize}
\tightlist
\item
  The \textbf{fuse} protects against sustained overcurrent (shorts, wiring
  faults).
\item
  The \textbf{TVS} protects against transient overvoltage (spikes, load dump,
  ESD).
\item
  Together they cover the two major electrical fault categories. Neither alone
  is sufficient.
\end{itemize}

\begin{center}\rule{0.5\linewidth}{0.5pt}\end{center}

\section{12 V accessory protection}\label{v-accessory-protection}

The 12 V rail powers signal lights, headlights, and the ESP32's LDO regulator.
It's downstream of the DC-DC converter (72V→12V).

Protection on the 12 V rail is less stringent (lower voltage, lower energy) but
still present: - \textbf{Reverse polarity protection:} A series Schottky diode
on the ESP32 power input prevents damage if the 12 V wiring is accidentally
reversed. - \textbf{Bulk capacitance:} 470 µF electrolytic on the 12 V rail at
the ESP32 input smooths voltage dips when lights switch on/off.

\begin{center}\rule{0.5\linewidth}{0.5pt}\end{center}

\section{Physical wiring rules}\label{physical-wiring-rules}

\begin{enumerate}
\def\labelenumi{\arabic{enumi}.}
\tightlist
\item
  \textbf{72 V wiring is orange.} 12 V wiring is red. Ground is black. CAN is
  yellow/green twisted pair. Color coding prevents accidental cross-connection.
\item
  \textbf{72 V terminals are insulated boot-style.} No exposed metal. All
  connectors are latching (not friction-fit).
\item
  \textbf{Fuses are accessible without disassembly.} Mounted on a terminal block
  with a clear cover.
\item
  \textbf{The 72 V battery has a service disconnect} (Anderson connector or
  equivalent) that physically breaks both positive and negative. Pull it before
  any wiring work.
\end{enumerate}

\begin{center}\rule{0.5\linewidth}{0.5pt}\end{center}

\section{Why this matters for the CAN
bus}\label{why-this-matters-for-the-can-bus}

The CAN bus shares a common ground reference. If 72 V leaks onto the CAN ground,
every transceiver on the bus sees 72 V common-mode and fails --- potentially all
at once. Galvanic isolation of the high-voltage signals prevents this. The
SN65HVD230 transceivers are rated for ±25 V common-mode; 72 V would destroy them
instantly.

\begin{center}\rule{0.5\linewidth}{0.5pt}\end{center}

\emph{Primary reference: {[}{[}emergency-system{]}{]} for ESTOP electrical fault
handling, emergency response matrix, and testing procedures.}

\emph{See also: {[}{[}defense-in-depth-safety{]}{]} for ESTOP layers including
motor/gear shutdown, {[}{[}architecture{]}{]} §5 for responsibility split, §8.6
for gear output details.}


\chapter{I/O Data --- Complete Variable
Inventory}\label{io-data-complete-variable-inventory}

This document catalogs every input and output variable across all four system
components: - \textbf{Host} (Jetson Orin) --- high-level autonomy - \textbf{RT}
(ESP32-S3) --- real-time physics, steering, CAN gateway - \textbf{SYS}
(ESP32-S3) --- safety (EGAS Level 2), brake control, body control - \textbf{MTR}
(STM32) --- dedicated motor controller (EGAS Level 1)

All CAN IDs and frame layouts are defined in
\texttt{shared/can/can\_protocol.h}. Physical pin mappings are in
\texttt{docs/wiring.md}.

\subsection{Type Notation}\label{type-notation}

\begin{longtable}[]{@{}
  >{\raggedright\arraybackslash}p{(\columnwidth - 6\tabcolsep) * \real{0.2778}}
  >{\raggedright\arraybackslash}p{(\columnwidth - 6\tabcolsep) * \real{0.2222}}
  >{\raggedright\arraybackslash}p{(\columnwidth - 6\tabcolsep) * \real{0.2500}}
  >{\raggedright\arraybackslash}p{(\columnwidth - 6\tabcolsep) * \real{0.2500}}@{}}
\toprule\noalign{}
\begin{minipage}[b]{\linewidth}\raggedright
Notation
\end{minipage} & \begin{minipage}[b]{\linewidth}\raggedright
C Type
\end{minipage} & \begin{minipage}[b]{\linewidth}\raggedright
Meaning
\end{minipage} & \begin{minipage}[b]{\linewidth}\raggedright
Example
\end{minipage} \\
\midrule\noalign{}
\endhead
\bottomrule\noalign{}
\endlastfoot
\texttt{i8} & \texttt{int8\_t} & Signed 8-bit integer & --- \\
\texttt{i16} & \texttt{int16\_t} & Signed 16-bit integer & \texttt{i16\ (mm/s)}
--- speed \\
\texttt{i24} & 24-bit signed (packed in CAN) & Signed 24-bit, non-standard width
& \texttt{i24\ (mrad/s)} --- yaw rate \\
\texttt{i32} & \texttt{int32\_t} & Signed 32-bit integer & \texttt{i32\ (kPa)}
--- brake pressure \\
\texttt{u8} & \texttt{uint8\_t} & Unsigned 8-bit integer & \texttt{u8\ enum} ---
mode, gear \\
\texttt{u16} & \texttt{uint16\_t} & Unsigned 16-bit integer &
\texttt{u16\ (0.05\ mm/bit)} --- stroke \\
\texttt{u32} & \texttt{uint32\_t} & Unsigned 32-bit integer & \texttt{u32\ (mm)}
--- distance \\
\texttt{u8\ bool} & \texttt{uint8\_t} & Boolean packed in a byte (0 or 1) &
\texttt{SYS\_EstopActive} \\
\texttt{u8\ enum} & \texttt{uint8\_t} & Enumeration packed in a byte &
\texttt{\{0=Manual,\ 1=Auto,\ 2=ESTOP\}} \\
\texttt{u8\ bitmask} & \texttt{uint8\_t} & Bitfield, each bit is a flag &
\texttt{fault\_flags} \\
\texttt{DLC=0} & --- & Zero-length CAN frame (no payload, event signal) &
\texttt{SAFETY\_ESTOP} \\
\end{longtable}

\subsection{Physical Units}\label{physical-units}

\begin{longtable}[]{@{}
  >{\raggedright\arraybackslash}p{(\columnwidth - 6\tabcolsep) * \real{0.1500}}
  >{\raggedright\arraybackslash}p{(\columnwidth - 6\tabcolsep) * \real{0.2250}}
  >{\raggedright\arraybackslash}p{(\columnwidth - 6\tabcolsep) * \real{0.2500}}
  >{\raggedright\arraybackslash}p{(\columnwidth - 6\tabcolsep) * \real{0.3750}}@{}}
\toprule\noalign{}
\begin{minipage}[b]{\linewidth}\raggedright
Unit
\end{minipage} & \begin{minipage}[b]{\linewidth}\raggedright
Meaning
\end{minipage} & \begin{minipage}[b]{\linewidth}\raggedright
Used For
\end{minipage} & \begin{minipage}[b]{\linewidth}\raggedright
Concrete range
\end{minipage} \\
\midrule\noalign{}
\endhead
\bottomrule\noalign{}
\endlastfoot
\texttt{mm/s} & Millimeters per second & Speed & {[}-500, 3000{]} = -0.5 to 3.0
m/s \\
\texttt{mrad/s} & Milliradians per second & Yaw rate & {[}-3000, 3000{]} \\
\texttt{kPa} & Kilopascals & Brake pressure & 0--12000 (\textasciitilde12 MPa
max) \\
\texttt{mm} & Millimeters & Distance & Obstacle: 0--4000 \\
\texttt{0.1°/bit} & Tenths of a degree per LSB & Steering angle & 455 = 45.5° \\
\texttt{0.05\ mm/bit} & 0.05 mm per LSB & Brake stroke & -30 mm offset \\
\texttt{0.05\ MPa/bit} & 0.05 MPa per LSB & Brake pressure & 1 bar = 0.1 MPa ≈ 2
bits \\
\texttt{°/s} & Degrees per second & Steering angle speed & EPS-C ramp rate \\
\texttt{Nm} & Newton-meters & Torque & Steering wheel resistance \\
\end{longtable}

\begin{center}\rule{0.5\linewidth}{0.5pt}\end{center}

\section{1. Host (Jetson Orin)}\label{host-jetson-orin}

\subsection{1.1 CAN Inputs --- Host receives}\label{can-inputs-host-receives}

\begin{longtable}[]{@{}
  >{\raggedright\arraybackslash}p{(\columnwidth - 10\tabcolsep) * \real{0.1481}}
  >{\raggedright\arraybackslash}p{(\columnwidth - 10\tabcolsep) * \real{0.1667}}
  >{\raggedright\arraybackslash}p{(\columnwidth - 10\tabcolsep) * \real{0.1852}}
  >{\raggedright\arraybackslash}p{(\columnwidth - 10\tabcolsep) * \real{0.1111}}
  >{\raggedright\arraybackslash}p{(\columnwidth - 10\tabcolsep) * \real{0.2407}}
  >{\raggedright\arraybackslash}p{(\columnwidth - 10\tabcolsep) * \real{0.1481}}@{}}
\toprule\noalign{}
\begin{minipage}[b]{\linewidth}\raggedright
CAN ID
\end{minipage} & \begin{minipage}[b]{\linewidth}\raggedright
Message
\end{minipage} & \begin{minipage}[b]{\linewidth}\raggedright
Variable
\end{minipage} & \begin{minipage}[b]{\linewidth}\raggedright
Type
\end{minipage} & \begin{minipage}[b]{\linewidth}\raggedright
Range / Enum
\end{minipage} & \begin{minipage}[b]{\linewidth}\raggedright
Source
\end{minipage} \\
\midrule\noalign{}
\endhead
\bottomrule\noalign{}
\endlastfoot
\texttt{0x001} & SAFETY\_ESTOP & (no payload) & DLC=0 & --- & any node \\
\texttt{0x011} & SYS\_SAFETY\_STS & \texttt{SYS\_EstopActive} & u8 bool & 0/1 &
SYS → RT → Host \\
\texttt{0x011} & SYS\_SAFETY\_STS & \texttt{SYS\_HeartbeatOk} & u8 bool & 0/1 &
SYS → RT → Host \\
\texttt{0x120} & SYS\_THROTTLE\_STS & \texttt{SYS\_ThrottleSpeed} & i16 (mm/s) &
--- & MTR → RT → Host \\
\texttt{0x210} & RT\_STATE\_RPT & \texttt{RT\_Mode} & u8 enum & \{0=Manual,
1=Auto, 2=ESTOP\} & RT \\
\texttt{0x210} & RT\_STATE\_RPT & \texttt{RT\_SteerValid} & u8 bool & 0/1 &
RT \\
\texttt{0x210} & RT\_STATE\_RPT & \texttt{RT\_Reversing} & u8 bool & 0/1 & RT \\
\texttt{0x220} & RT\_PID\_RPT & \texttt{RT\_PidSetpoint} & i16 (mm/s) & --- & RT
(reserved) \\
\texttt{0x220} & RT\_PID\_RPT & \texttt{RT\_PidMeasured} & i16 (mm/s) & --- & RT
(reserved) \\
\texttt{0x220} & RT\_PID\_RPT & \texttt{RT\_PidOutput} & i16 & --- & RT
(reserved) \\
\texttt{0x400} & RT\_OBSTACLE\_RPT & \texttt{RT\_ObstacleDistance} & u32 (mm) &
--- & RT \\
\texttt{0x600} & SYS\_DIAG\_RPT & \texttt{SYS\_DiagMode} & u8 & --- & SYS → RT →
Host \\
\texttt{0x600} & SYS\_DIAG\_RPT & \texttt{SYS\_DiagBrakeEngaged} & u8 bool & 0/1
& SYS → RT → Host \\
\texttt{0x600} & SYS\_DIAG\_RPT & \texttt{SYS\_DiagHeartbeatOk} & u8 bool & 0/1
& SYS → RT → Host \\
\texttt{0x600} & SYS\_DIAG\_RPT & \texttt{SYS\_DiagEstopActive} & u8 bool & 0/1
& SYS → RT → Host \\
\texttt{0x600} & SYS\_DIAG\_RPT & \texttt{SYS\_DiagFreeHeapKb} & u16 & --- & SYS
→ RT → Host \\
\texttt{0x600} & SYS\_DIAG\_RPT & \texttt{SYS\_DiagTec} & u8 & --- & SYS → RT →
Host \\
\texttt{0x600} & SYS\_DIAG\_RPT & \texttt{SYS\_DiagRec} & u8 & --- & SYS → RT →
Host \\
\texttt{0x7FD} & RT\_HEARTBEAT & \texttt{alive\_ctr} & u8 & --- & RT \\
\end{longtable}

\subsection{1.2 CAN Outputs --- Host sends}\label{can-outputs-host-sends}

\begin{longtable}[]{@{}
  >{\raggedright\arraybackslash}p{(\columnwidth - 10\tabcolsep) * \real{0.1481}}
  >{\raggedright\arraybackslash}p{(\columnwidth - 10\tabcolsep) * \real{0.1667}}
  >{\raggedright\arraybackslash}p{(\columnwidth - 10\tabcolsep) * \real{0.1852}}
  >{\raggedright\arraybackslash}p{(\columnwidth - 10\tabcolsep) * \real{0.1111}}
  >{\raggedright\arraybackslash}p{(\columnwidth - 10\tabcolsep) * \real{0.2407}}
  >{\raggedright\arraybackslash}p{(\columnwidth - 10\tabcolsep) * \real{0.1481}}@{}}
\toprule\noalign{}
\begin{minipage}[b]{\linewidth}\raggedright
CAN ID
\end{minipage} & \begin{minipage}[b]{\linewidth}\raggedright
Message
\end{minipage} & \begin{minipage}[b]{\linewidth}\raggedright
Variable
\end{minipage} & \begin{minipage}[b]{\linewidth}\raggedright
Type
\end{minipage} & \begin{minipage}[b]{\linewidth}\raggedright
Range / Enum
\end{minipage} & \begin{minipage}[b]{\linewidth}\raggedright
Target
\end{minipage} \\
\midrule\noalign{}
\endhead
\bottomrule\noalign{}
\endlastfoot
\texttt{0x001} & SAFETY\_ESTOP & (no payload) & DLC=0 & --- & all \\
\texttt{0x300} & HOST\_DRIVE\_CMD & \texttt{HOST\_DriveSpeed} & i32 (mm/s) &
{[}-500, 3000{]} & RT \\
\texttt{0x300} & HOST\_DRIVE\_CMD & \texttt{HOST\_YawRate} & i24 (mrad/s) &
{[}-3000, 3000{]} & RT \\
\texttt{0x300} & HOST\_DRIVE\_CMD & \texttt{HOST\_Gear} & u8 enum & \{N=0, D=1,
S=2, R=3\} & RT \\
\texttt{0x301} & HOST\_BRAKE\_REQ & \texttt{HOST\_BrakePressure} & i32 (kPa) &
--- & RT \\
\texttt{0x302} & HOST\_LIGHT\_CMD & \texttt{HOST\_LeftTurn} & bool & 0/1 & RT \\
\texttt{0x302} & HOST\_LIGHT\_CMD & \texttt{HOST\_RightTurn} & bool & 0/1 &
RT \\
\texttt{0x302} & HOST\_LIGHT\_CMD & \texttt{HOST\_BrakeLight} & bool & 0/1 &
RT \\
\texttt{0x302} & HOST\_LIGHT\_CMD & \texttt{HOST\_Headlight} & bool & 0/1 &
RT \\
\texttt{0x7FC} & JETSON\_HEARTBEAT & \texttt{alive\_ctr} & u8 & --- & RT \\
\end{longtable}

\begin{center}\rule{0.5\linewidth}{0.5pt}\end{center}

\section{2. RT (ESP32-S3)}\label{rt-esp32-s3}

Role: real-time physics model, steering control (EPS-C via CAN), CAN gateway
between high-side and low-side buses, obstacle braking, heartbeat/watchdog
monitoring.

\subsection{2.1 CAN Inputs --- RT receives (low-side
bus)}\label{can-inputs-rt-receives-low-side-bus}

\begin{longtable}[]{@{}
  >{\raggedright\arraybackslash}p{(\columnwidth - 10\tabcolsep) * \real{0.1481}}
  >{\raggedright\arraybackslash}p{(\columnwidth - 10\tabcolsep) * \real{0.1667}}
  >{\raggedright\arraybackslash}p{(\columnwidth - 10\tabcolsep) * \real{0.1852}}
  >{\raggedright\arraybackslash}p{(\columnwidth - 10\tabcolsep) * \real{0.1111}}
  >{\raggedright\arraybackslash}p{(\columnwidth - 10\tabcolsep) * \real{0.2407}}
  >{\raggedright\arraybackslash}p{(\columnwidth - 10\tabcolsep) * \real{0.1481}}@{}}
\toprule\noalign{}
\begin{minipage}[b]{\linewidth}\raggedright
CAN ID
\end{minipage} & \begin{minipage}[b]{\linewidth}\raggedright
Message
\end{minipage} & \begin{minipage}[b]{\linewidth}\raggedright
Variable
\end{minipage} & \begin{minipage}[b]{\linewidth}\raggedright
Type
\end{minipage} & \begin{minipage}[b]{\linewidth}\raggedright
Range / Enum
\end{minipage} & \begin{minipage}[b]{\linewidth}\raggedright
Source
\end{minipage} \\
\midrule\noalign{}
\endhead
\bottomrule\noalign{}
\endlastfoot
\texttt{0x001} & SAFETY\_ESTOP & (no payload) & DLC=0 & --- & any \\
\texttt{0x011} & SYS\_SAFETY\_STS & \texttt{SYS\_EstopActive} & u8 bool & 0/1 &
SYS \\
\texttt{0x011} & SYS\_SAFETY\_STS & \texttt{SYS\_HeartbeatOk} & u8 bool & 0/1 &
SYS \\
\texttt{0x110} & SYS\_MODE\_CMD & \texttt{SYS\_Mode} & u8 enum & \{0=Manual,
1=Auto, 2=ESTOP\} & SYS \\
\texttt{0x120} & SYS\_THROTTLE\_STS & \texttt{SYS\_ThrottleSpeed} & i16 (mm/s) &
--- & MTR \\
\texttt{0x201} & SES\_STATUS & \texttt{SES\_INF\_Angle\_Status} & bool & 0/1 &
EPS-C \\
\texttt{0x201} & SES\_STATUS & \texttt{SES\_Control\_Mode\_Status} & u8 enum &
--- & EPS-C \\
\texttt{0x201} & SES\_STATUS & \texttt{SES\_Error\_Status} & u8 enum & \{0=N,
1=L1, 2=L2, 3=L3\} & EPS-C \\
\texttt{0x201} & SES\_STATUS & \texttt{SES\_StrAngle} & i16 (0.1°/bit) &
{[}-780, 780{]} & EPS-C \\
\texttt{0x201} & SES\_STATUS & \texttt{EPS\_SteeringWheel\_Torq} & u8 (Nm) & ---
& EPS-C \\
\texttt{0x201} & SES\_STATUS & \texttt{SES\_Tgt\_StrAngleSpd} & i16 (0.5°/s/bit)
& --- & EPS-C \\
\texttt{0x201} & SES\_STATUS & \texttt{SES\_RollCnt\_Enable\_Status} & bool &
0/1 & EPS-C \\
\texttt{0x201} & SES\_STATUS & \texttt{SES\_CheckSum\_Enable\_Status} & bool &
0/1 & EPS-C \\
\texttt{0x201} & SES\_STATUS & \texttt{SES\_RollCnt\_Status} & u8 (low 4 bits) &
0--15 & EPS-C \\
\texttt{0x201} & SES\_STATUS & \texttt{SES\_CheckSum\_Status} & u8 & XOR over
bytes 0--6 & EPS-C \\
\texttt{0x7FE} & SYS\_HEARTBEAT & \texttt{alive\_ctr} & u8 & --- & SYS \\
\end{longtable}

\subsection{2.2 CAN Inputs --- RT receives (high-side bus, via
MCP2515)}\label{can-inputs-rt-receives-high-side-bus-via-mcp2515}

\begin{longtable}[]{@{}
  >{\raggedright\arraybackslash}p{(\columnwidth - 10\tabcolsep) * \real{0.1481}}
  >{\raggedright\arraybackslash}p{(\columnwidth - 10\tabcolsep) * \real{0.1667}}
  >{\raggedright\arraybackslash}p{(\columnwidth - 10\tabcolsep) * \real{0.1852}}
  >{\raggedright\arraybackslash}p{(\columnwidth - 10\tabcolsep) * \real{0.1111}}
  >{\raggedright\arraybackslash}p{(\columnwidth - 10\tabcolsep) * \real{0.2407}}
  >{\raggedright\arraybackslash}p{(\columnwidth - 10\tabcolsep) * \real{0.1481}}@{}}
\toprule\noalign{}
\begin{minipage}[b]{\linewidth}\raggedright
CAN ID
\end{minipage} & \begin{minipage}[b]{\linewidth}\raggedright
Message
\end{minipage} & \begin{minipage}[b]{\linewidth}\raggedright
Variable
\end{minipage} & \begin{minipage}[b]{\linewidth}\raggedright
Type
\end{minipage} & \begin{minipage}[b]{\linewidth}\raggedright
Range / Enum
\end{minipage} & \begin{minipage}[b]{\linewidth}\raggedright
Source
\end{minipage} \\
\midrule\noalign{}
\endhead
\bottomrule\noalign{}
\endlastfoot
\texttt{0x300} & HOST\_DRIVE\_CMD & \texttt{HOST\_DriveSpeed} & i32 (mm/s) &
{[}-500, 3000{]} & Host \\
\texttt{0x300} & HOST\_DRIVE\_CMD & \texttt{HOST\_YawRate} & i24 (mrad/s) &
{[}-3000, 3000{]} & Host \\
\texttt{0x300} & HOST\_DRIVE\_CMD & \texttt{HOST\_Gear} & u8 enum & \{N,D,S,R\}
& Host \\
\texttt{0x301} & HOST\_BRAKE\_REQ & \texttt{HOST\_BrakePressure} & i32 (kPa) &
--- & Host \\
\texttt{0x302} & HOST\_LIGHT\_CMD & light bits (4× bool) & u8 bitmask & --- &
Host \\
\texttt{0x7FC} & JETSON\_HEARTBEAT & \texttt{alive\_ctr} & u8 & --- & Host \\
\end{longtable}

\subsection{2.3 Physical Inputs --- RT}\label{physical-inputs-rt}

\begin{longtable}[]{@{}lllll@{}}
\toprule\noalign{}
Variable & GPIO & Type & Connected To & Notes \\
\midrule\noalign{}
\endhead
\bottomrule\noalign{}
\endlastfoot
\texttt{CAN\_RX\_LOW} & 4 & TWAI & SN65HVD230 RXD & Low-side CAN bus \\
\texttt{MCP2515\_MISO} & 38 & SPI & MCP2515 SO & High-side CAN RX data \\
\texttt{MCP2515\_INT} & 40 & Digital & MCP2515 INT & High-side CAN RX
interrupt \\
\texttt{ENC\_REAR\_MOTOR\_A} & 1 & PCNT & Motor encoder A & Speed feedback ---
active \\
\texttt{ENC\_REAR\_MOTOR\_B} & 2 & PCNT & Motor encoder B & Quadrature phase
B \\
\texttt{ENC\_FRONT\_WHEEL\_A} & 3 & PCNT & Front wheel & Sensor not yet
fitted \\
\texttt{ENC\_FRONT\_WHEEL\_B} & 6 & PCNT & Front wheel & Sensor not yet
fitted \\
\texttt{ENC\_REAR\_LEFT\_A} & 9 & PCNT & Rear left wheel & Sensor not yet
fitted \\
\texttt{ENC\_REAR\_LEFT\_B} & 12 & PCNT & Rear left wheel & Sensor not yet
fitted \\
\texttt{ENC\_REAR\_RIGHT\_A} & 13 & PCNT & Rear right wheel & Sensor not yet
fitted \\
\texttt{ENC\_REAR\_RIGHT\_B} & 14 & PCNT & Rear right wheel & Sensor not yet
fitted \\
\texttt{I2C\_SDA} (IMU) & 10 & I2C & Optional IMU & --- \\
\end{longtable}

\subsection{2.4 CAN Outputs --- RT sends (low-side
bus)}\label{can-outputs-rt-sends-low-side-bus}

\begin{longtable}[]{@{}
  >{\raggedright\arraybackslash}p{(\columnwidth - 12\tabcolsep) * \real{0.1333}}
  >{\raggedright\arraybackslash}p{(\columnwidth - 12\tabcolsep) * \real{0.1500}}
  >{\raggedright\arraybackslash}p{(\columnwidth - 12\tabcolsep) * \real{0.1667}}
  >{\raggedright\arraybackslash}p{(\columnwidth - 12\tabcolsep) * \real{0.1000}}
  >{\raggedright\arraybackslash}p{(\columnwidth - 12\tabcolsep) * \real{0.2167}}
  >{\raggedright\arraybackslash}p{(\columnwidth - 12\tabcolsep) * \real{0.1000}}
  >{\raggedright\arraybackslash}p{(\columnwidth - 12\tabcolsep) * \real{0.1333}}@{}}
\toprule\noalign{}
\begin{minipage}[b]{\linewidth}\raggedright
CAN ID
\end{minipage} & \begin{minipage}[b]{\linewidth}\raggedright
Message
\end{minipage} & \begin{minipage}[b]{\linewidth}\raggedright
Variable
\end{minipage} & \begin{minipage}[b]{\linewidth}\raggedright
Type
\end{minipage} & \begin{minipage}[b]{\linewidth}\raggedright
Range / Enum
\end{minipage} & \begin{minipage}[b]{\linewidth}\raggedright
Rate
\end{minipage} & \begin{minipage}[b]{\linewidth}\raggedright
Target
\end{minipage} \\
\midrule\noalign{}
\endhead
\bottomrule\noalign{}
\endlastfoot
\texttt{0x001} & SAFETY\_ESTOP & (no payload) & DLC=0 & --- & event & all \\
\texttt{0x204} & RT\_DRIVE\_CMD & \texttt{RT\_MotorSpeed} & i32 (mm/s) &
{[}-500, 3000{]} & 100 Hz & MTR \\
\texttt{0x204} & RT\_DRIVE\_CMD & \texttt{RT\_Gear} & u8 enum & \{N,D,S,R\} &
100 Hz & MTR \\
\texttt{0x205} & RT\_BRAKE\_CMD & \texttt{RT\_BrakePressure} & i32 (kPa) & --- &
50 Hz & SYS \\
\texttt{0x169} & VCU\_SES\_REQ & \texttt{VCU\_SES\_Alignment\_Enable} & bool &
0/1 & 50 Hz & EPS-C \\
\texttt{0x169} & VCU\_SES\_REQ & \texttt{VCU\_SES\_Control\_Enable} & bool & 0/1
& 50 Hz & EPS-C \\
\texttt{0x169} & VCU\_SES\_REQ & \texttt{VCU\_SES\_Control\_Mode} & u8 & --- &
50 Hz & EPS-C \\
\texttt{0x169} & VCU\_SES\_REQ & \texttt{VCU\_SES\_Tgt\_StrAngle} & i16
(0.1°/bit) & {[}-780, 780{]} & 50 Hz & EPS-C \\
\texttt{0x169} & VCU\_SES\_REQ & \texttt{VCU\_SES\_Tgt\_StrAngleSpd} & u8 (°/s)
& --- & 50 Hz & EPS-C \\
\texttt{0x169} & VCU\_SES\_REQ & \texttt{roll\_cnt\_enable} & bool & must be 1 &
50 Hz & EPS-C \\
\texttt{0x169} & VCU\_SES\_REQ & \texttt{checksum\_enable} & bool & must be 1 &
50 Hz & EPS-C \\
\texttt{0x169} & VCU\_SES\_REQ & \texttt{VCU\_SES\_RollCnt} & u8 (low 4 bits) &
0--15 rolling & 50 Hz & EPS-C \\
\texttt{0x169} & VCU\_SES\_REQ & \texttt{VCU\_SES\_CheckSum} & u8 & XOR over
bytes 0--6 & 50 Hz & EPS-C \\
\texttt{0x169} & VCU\_SES\_REQ & \texttt{VCU\_Veh\_Spd\_Value} & u8 (km/h) &
0--255 & 50 Hz & EPS-C \\
\texttt{0x302} & HOST\_LIGHT\_CMD (fwd) & light bits & u8 bitmask & --- & on
change & SYS \\
\texttt{0x7FD} & RT\_HEARTBEAT & \texttt{alive\_ctr} & u8 & --- & 2 Hz & SYS \\
\end{longtable}

\subsection{2.5 CAN Outputs --- RT sends (high-side bus, via
MCP2515)}\label{can-outputs-rt-sends-high-side-bus-via-mcp2515}

\begin{longtable}[]{@{}
  >{\raggedright\arraybackslash}p{(\columnwidth - 12\tabcolsep) * \real{0.1333}}
  >{\raggedright\arraybackslash}p{(\columnwidth - 12\tabcolsep) * \real{0.1500}}
  >{\raggedright\arraybackslash}p{(\columnwidth - 12\tabcolsep) * \real{0.1667}}
  >{\raggedright\arraybackslash}p{(\columnwidth - 12\tabcolsep) * \real{0.1000}}
  >{\raggedright\arraybackslash}p{(\columnwidth - 12\tabcolsep) * \real{0.2167}}
  >{\raggedright\arraybackslash}p{(\columnwidth - 12\tabcolsep) * \real{0.1000}}
  >{\raggedright\arraybackslash}p{(\columnwidth - 12\tabcolsep) * \real{0.1333}}@{}}
\toprule\noalign{}
\begin{minipage}[b]{\linewidth}\raggedright
CAN ID
\end{minipage} & \begin{minipage}[b]{\linewidth}\raggedright
Message
\end{minipage} & \begin{minipage}[b]{\linewidth}\raggedright
Variable
\end{minipage} & \begin{minipage}[b]{\linewidth}\raggedright
Type
\end{minipage} & \begin{minipage}[b]{\linewidth}\raggedright
Range / Enum
\end{minipage} & \begin{minipage}[b]{\linewidth}\raggedright
Rate
\end{minipage} & \begin{minipage}[b]{\linewidth}\raggedright
Target
\end{minipage} \\
\midrule\noalign{}
\endhead
\bottomrule\noalign{}
\endlastfoot
\texttt{0x011} & SYS\_SAFETY\_STS (fwd) & \texttt{SYS\_EstopActive} & u8 bool &
0/1 & 5 Hz & Host \\
\texttt{0x011} & SYS\_SAFETY\_STS (fwd) & \texttt{SYS\_HeartbeatOk} & u8 bool &
0/1 & 5 Hz & Host \\
\texttt{0x120} & SYS\_THROTTLE\_STS (fwd) & \texttt{SYS\_ThrottleSpeed} & i16
(mm/s) & --- & 100 Hz & Host \\
\texttt{0x210} & RT\_STATE\_RPT & \texttt{RT\_Mode} & u8 enum & \{M,A,ESTOP\} &
10 Hz & Host \\
\texttt{0x210} & RT\_STATE\_RPT & \texttt{RT\_SteerValid} & u8 bool & 0/1 & 10
Hz & Host \\
\texttt{0x210} & RT\_STATE\_RPT & \texttt{RT\_Reversing} & u8 bool & 0/1 & 10 Hz
& Host \\
\texttt{0x220} & RT\_PID\_RPT & \texttt{RT\_PidSetpoint} & i16 (mm/s) & --- &
reserved & Host \\
\texttt{0x220} & RT\_PID\_RPT & \texttt{RT\_PidMeasured} & i16 (mm/s) & --- &
reserved & Host \\
\texttt{0x220} & RT\_PID\_RPT & \texttt{RT\_PidOutput} & i16 & --- & reserved &
Host \\
\texttt{0x400} & RT\_OBSTACLE\_RPT & \texttt{RT\_ObstacleDistance} & u32 (mm) &
--- & 10 Hz & Host \\
\texttt{0x600} & SYS\_DIAG\_RPT (fwd) & 8-byte diagnostic struct & struct & ---
& 1 Hz & Host \\
\texttt{0x7FD} & RT\_HEARTBEAT & \texttt{alive\_ctr} & u8 & --- & 2 Hz & Host \\
\end{longtable}

\subsection{2.6 Physical Outputs --- RT}\label{physical-outputs-rt}

\begin{longtable}[]{@{}
  >{\raggedright\arraybackslash}p{(\columnwidth - 8\tabcolsep) * \real{0.2381}}
  >{\raggedright\arraybackslash}p{(\columnwidth - 8\tabcolsep) * \real{0.1429}}
  >{\raggedright\arraybackslash}p{(\columnwidth - 8\tabcolsep) * \real{0.1429}}
  >{\raggedright\arraybackslash}p{(\columnwidth - 8\tabcolsep) * \real{0.3095}}
  >{\raggedright\arraybackslash}p{(\columnwidth - 8\tabcolsep) * \real{0.1667}}@{}}
\toprule\noalign{}
\begin{minipage}[b]{\linewidth}\raggedright
Variable
\end{minipage} & \begin{minipage}[b]{\linewidth}\raggedright
GPIO
\end{minipage} & \begin{minipage}[b]{\linewidth}\raggedright
Type
\end{minipage} & \begin{minipage}[b]{\linewidth}\raggedright
Connected To
\end{minipage} & \begin{minipage}[b]{\linewidth}\raggedright
Notes
\end{minipage} \\
\midrule\noalign{}
\endhead
\bottomrule\noalign{}
\endlastfoot
\texttt{CAN\_TX\_LOW} & 5 & TWAI & SN65HVD230 TXD & Low-side CAN bus \\
\texttt{SPI\_SCK} & 36 & SPI & MCP2515 SCK & 10 MHz clock \\
\texttt{SPI\_MOSI} & 37 & SPI & MCP2515 SI & High-side CAN TX data \\
\texttt{SPI\_CS} & 39 & Digital & MCP2515 CS & Active-low chip select \\
\texttt{I2C\_SCL} (IMU) & 11 & I2C & Optional IMU & Clock \\
\texttt{WDT\_TOGGLE} & 21 & Digital & TPS3850 WDI & Toggled at 100 Hz by control
task \\
\end{longtable}

\begin{center}\rule{0.5\linewidth}{0.5pt}\end{center}

\section{3. SYS (ESP32-S3)}\label{sys-esp32-s3}

Role: safety monitoring (EGAS Level 2), ESTOP handling, brake control (SEB via
CAN), body control (lights, DCDC, mode management), watchdog. Motor actuation is
on MTR only --- SYS monitors via CAN 0x206 feedback.

\subsection{3.1 CAN Inputs --- SYS receives}\label{can-inputs-sys-receives}

\begin{longtable}[]{@{}
  >{\raggedright\arraybackslash}p{(\columnwidth - 10\tabcolsep) * \real{0.1481}}
  >{\raggedright\arraybackslash}p{(\columnwidth - 10\tabcolsep) * \real{0.1667}}
  >{\raggedright\arraybackslash}p{(\columnwidth - 10\tabcolsep) * \real{0.1852}}
  >{\raggedright\arraybackslash}p{(\columnwidth - 10\tabcolsep) * \real{0.1111}}
  >{\raggedright\arraybackslash}p{(\columnwidth - 10\tabcolsep) * \real{0.2407}}
  >{\raggedright\arraybackslash}p{(\columnwidth - 10\tabcolsep) * \real{0.1481}}@{}}
\toprule\noalign{}
\begin{minipage}[b]{\linewidth}\raggedright
CAN ID
\end{minipage} & \begin{minipage}[b]{\linewidth}\raggedright
Message
\end{minipage} & \begin{minipage}[b]{\linewidth}\raggedright
Variable
\end{minipage} & \begin{minipage}[b]{\linewidth}\raggedright
Type
\end{minipage} & \begin{minipage}[b]{\linewidth}\raggedright
Range / Enum
\end{minipage} & \begin{minipage}[b]{\linewidth}\raggedright
Source
\end{minipage} \\
\midrule\noalign{}
\endhead
\bottomrule\noalign{}
\endlastfoot
\texttt{0x001} & SAFETY\_ESTOP & (no payload) & DLC=0 & --- & any \\
\texttt{0x205} & RT\_BRAKE\_CMD & \texttt{RT\_BrakePressure} & i32 (kPa) & --- &
RT \\
\texttt{0x206} & MTR\_MOTOR\_FBK & \texttt{actual\_speed\_mmps} & i16 (mm/s) &
--- & MTR \\
\texttt{0x206} & MTR\_MOTOR\_FBK & \texttt{gear\_state} & u8 enum & \{N,D,S,R\}
& MTR \\
\texttt{0x206} & MTR\_MOTOR\_FBK & \texttt{fault\_flags} & u8 bitmask & --- &
MTR \\
\texttt{0x302} & HOST\_LIGHT\_CMD & light bits (4× bool) & u8 bitmask & --- & RT
(fwd) \\
\texttt{0x721} & SEB\_STATUS & \texttt{SEB\_Alignment\_Status} & bool & 0/1 &
SEB \\
\texttt{0x721} & SEB\_STATUS & \texttt{SEB\_Control\_Enable\_Status} & bool &
0/1 & SEB \\
\texttt{0x721} & SEB\_STATUS & \texttt{SEB\_Control\_Mode\_Status} & u8 enum &
--- & SEB \\
\texttt{0x721} & SEB\_STATUS & \texttt{SEB\_AutoBrake\_Status} & bool & 0/1 &
SEB \\
\texttt{0x721} & SEB\_STATUS & \texttt{SEB\_Error\_Status} & u8 enum & \{0=N,
1=L1, 2=L2, 3=L3\} & SEB \\
\texttt{0x721} & SEB\_STATUS & \texttt{SEB\_Stroke\_Value} & u16 (0.05 mm/bit) &
offset -30 mm & SEB \\
\texttt{0x721} & SEB\_STATUS & \texttt{SEB\_Pressure\_Value} & u8 (0.05 MPa/bit)
& --- & SEB \\
\texttt{0x721} & SEB\_STATUS & \texttt{SEB\_Angle\_Value} & i16 (0.5°/bit) &
{[}-150, 840{]} & SEB \\
\texttt{0x721} & SEB\_STATUS & \texttt{SEB\_RollCnt\_Enable\_Status} & bool &
0/1 & SEB \\
\texttt{0x721} & SEB\_STATUS & \texttt{SEB\_CheckSum\_Enable\_Status} & bool &
0/1 & SEB \\
\texttt{0x721} & SEB\_STATUS & \texttt{SEB\_RollCnt\_Status} & u8 (low 4 bits) &
--- & SEB \\
\texttt{0x721} & SEB\_STATUS & \texttt{SEB\_CheckSum\_Status} & u8 & --- &
SEB \\
\texttt{0x7FD} & RT\_HEARTBEAT & \texttt{alive\_ctr} & u8 & --- & RT \\
\end{longtable}

\subsection{3.2 Physical Inputs --- SYS}\label{physical-inputs-sys}

\begin{longtable}[]{@{}
  >{\raggedright\arraybackslash}p{(\columnwidth - 8\tabcolsep) * \real{0.2500}}
  >{\raggedright\arraybackslash}p{(\columnwidth - 8\tabcolsep) * \real{0.1500}}
  >{\raggedright\arraybackslash}p{(\columnwidth - 8\tabcolsep) * \real{0.1500}}
  >{\raggedright\arraybackslash}p{(\columnwidth - 8\tabcolsep) * \real{0.2750}}
  >{\raggedright\arraybackslash}p{(\columnwidth - 8\tabcolsep) * \real{0.1750}}@{}}
\toprule\noalign{}
\begin{minipage}[b]{\linewidth}\raggedright
Variable
\end{minipage} & \begin{minipage}[b]{\linewidth}\raggedright
GPIO
\end{minipage} & \begin{minipage}[b]{\linewidth}\raggedright
Type
\end{minipage} & \begin{minipage}[b]{\linewidth}\raggedright
Electrical
\end{minipage} & \begin{minipage}[b]{\linewidth}\raggedright
Notes
\end{minipage} \\
\midrule\noalign{}
\endhead
\bottomrule\noalign{}
\endlastfoot
\texttt{CAN\_RX} & 4 & TWAI & SN65HVD230 RXD & Low-side CAN bus \\
\texttt{ESTOP\_BTN} & 1 & Digital NC & 10k pull-up to 3.3V, LOW = estop & Red
mushroom button. Shared with MTR --- only signal wired to both MCUs. \\
\texttt{BRAKE\_LEVER} & 2 & Digital & Internal pull-up, LOW = pressed & Physical
brake lever → SEB via CAN 0x7B9 \\
\texttt{START\_BTN} & 32 & Digital & Internal pull-up, LOW = pressed & Green
button, exits ESTOP \\
\texttt{MODE\_BTN} & 11 & Digital & Internal pull-up, LOW = pressed &
Manual/Auto toggle → publishes CAN 0x110 \\
\texttt{SW\_LEFT\_TURN} & 3 & Digital & Internal pull-up, LOW = pressed &
Handlebar switch \\
\texttt{SW\_RIGHT\_TURN} & 6 & Digital & Internal pull-up, LOW = pressed &
Handlebar switch \\
\texttt{SW\_HEADLIGHT} & 7 & Digital & Internal pull-up, LOW = pressed &
Handlebar switch \\
\end{longtable}

\subsection{3.3 CAN Outputs --- SYS sends}\label{can-outputs-sys-sends}

\begin{longtable}[]{@{}
  >{\raggedright\arraybackslash}p{(\columnwidth - 12\tabcolsep) * \real{0.1333}}
  >{\raggedright\arraybackslash}p{(\columnwidth - 12\tabcolsep) * \real{0.1500}}
  >{\raggedright\arraybackslash}p{(\columnwidth - 12\tabcolsep) * \real{0.1667}}
  >{\raggedright\arraybackslash}p{(\columnwidth - 12\tabcolsep) * \real{0.1000}}
  >{\raggedright\arraybackslash}p{(\columnwidth - 12\tabcolsep) * \real{0.2167}}
  >{\raggedright\arraybackslash}p{(\columnwidth - 12\tabcolsep) * \real{0.1000}}
  >{\raggedright\arraybackslash}p{(\columnwidth - 12\tabcolsep) * \real{0.1333}}@{}}
\toprule\noalign{}
\begin{minipage}[b]{\linewidth}\raggedright
CAN ID
\end{minipage} & \begin{minipage}[b]{\linewidth}\raggedright
Message
\end{minipage} & \begin{minipage}[b]{\linewidth}\raggedright
Variable
\end{minipage} & \begin{minipage}[b]{\linewidth}\raggedright
Type
\end{minipage} & \begin{minipage}[b]{\linewidth}\raggedright
Range / Enum
\end{minipage} & \begin{minipage}[b]{\linewidth}\raggedright
Rate
\end{minipage} & \begin{minipage}[b]{\linewidth}\raggedright
Target
\end{minipage} \\
\midrule\noalign{}
\endhead
\bottomrule\noalign{}
\endlastfoot
\texttt{0x001} & SAFETY\_ESTOP & (no payload) & DLC=0 & --- & event & all \\
\texttt{0x011} & SYS\_SAFETY\_STS & \texttt{SYS\_EstopActive} & u8 bool & 0/1 &
5 Hz & RT → Host \\
\texttt{0x011} & SYS\_SAFETY\_STS & \texttt{SYS\_HeartbeatOk} & u8 bool & 0/1 &
5 Hz & RT → Host \\
\texttt{0x012} & SYS\_DCDC\_CMD & \texttt{SYS\_DcdcEnable} & u8 bool & 0/1 & on
change & DC-DC converter \\
\texttt{0x110} & SYS\_MODE\_CMD & \texttt{SYS\_Mode} & u8 enum & \{0=Manual,
1=Auto, 2=ESTOP\} & on change & RT, MTR \\
\texttt{0x600} & SYS\_DIAG\_RPT & \texttt{SYS\_DiagMode} & u8 & --- & 1 Hz & RT
→ Host \\
\texttt{0x600} & SYS\_DIAG\_RPT & \texttt{SYS\_DiagBrakeEngaged} & u8 bool & 0/1
& 1 Hz & RT → Host \\
\texttt{0x600} & SYS\_DIAG\_RPT & \texttt{SYS\_DiagHeartbeatOk} & u8 bool & 0/1
& 1 Hz & RT → Host \\
\texttt{0x600} & SYS\_DIAG\_RPT & \texttt{SYS\_DiagEstopActive} & u8 bool & 0/1
& 1 Hz & RT → Host \\
\texttt{0x600} & SYS\_DIAG\_RPT & \texttt{SYS\_DiagFreeHeapKb} & u16 & --- & 1
Hz & RT → Host \\
\texttt{0x600} & SYS\_DIAG\_RPT & \texttt{SYS\_DiagTec} & u8 & --- & 1 Hz & RT →
Host \\
\texttt{0x600} & SYS\_DIAG\_RPT & \texttt{SYS\_DiagRec} & u8 & --- & 1 Hz & RT →
Host \\
\texttt{0x7B9} & VCU\_SEB\_REQ & (see §5.2) & --- & --- & 50 Hz & SEB
(MANUAL/ESTOP only) \\
\texttt{0x7FE} & SYS\_HEARTBEAT & \texttt{alive\_ctr} & u8 & --- & 2 Hz & RT \\
\end{longtable}

\subsection{3.4 Physical Outputs --- SYS}\label{physical-outputs-sys}

\begin{longtable}[]{@{}
  >{\raggedright\arraybackslash}p{(\columnwidth - 8\tabcolsep) * \real{0.2381}}
  >{\raggedright\arraybackslash}p{(\columnwidth - 8\tabcolsep) * \real{0.1429}}
  >{\raggedright\arraybackslash}p{(\columnwidth - 8\tabcolsep) * \real{0.1429}}
  >{\raggedright\arraybackslash}p{(\columnwidth - 8\tabcolsep) * \real{0.3095}}
  >{\raggedright\arraybackslash}p{(\columnwidth - 8\tabcolsep) * \real{0.1667}}@{}}
\toprule\noalign{}
\begin{minipage}[b]{\linewidth}\raggedright
Variable
\end{minipage} & \begin{minipage}[b]{\linewidth}\raggedright
GPIO
\end{minipage} & \begin{minipage}[b]{\linewidth}\raggedright
Type
\end{minipage} & \begin{minipage}[b]{\linewidth}\raggedright
Connected To
\end{minipage} & \begin{minipage}[b]{\linewidth}\raggedright
Notes
\end{minipage} \\
\midrule\noalign{}
\endhead
\bottomrule\noalign{}
\endlastfoot
\texttt{CAN\_TX} & 5 & TWAI & SN65HVD230 TXD & Low-side CAN bus \\
\texttt{LIGHT\_LEFT\_TURN} & 18 & Digital & Relay → 12V lamp & Left turn
signal \\
\texttt{LIGHT\_RIGHT\_TURN} & 19 & Digital & Relay → 12V lamp & Right turn
signal \\
\texttt{BRAKE\_LIGHT} & 21 & Digital & Relay → 12V lamp & Brake light \\
\texttt{HEADLIGHT} & 22 & Digital & Relay → 12V lamp & Low-beam headlight \\
\texttt{BULB\_AUTO} & 25 & Digital & Relay → 12V indicator & Auto mode lamp \\
\texttt{BULB\_MANUAL} & 26 & Digital & Relay → 12V indicator & Manual mode
lamp \\
\texttt{RELAY\_12V} & 27 & Digital & Relay → accessory bus & Cuts accessory
power in ESTOP \\
\texttt{WDT\_TOGGLE} & 23 & Digital & TPS3850 WDI & Toggled at 20 Hz by safety
task \\
\end{longtable}

\begin{center}\rule{0.5\linewidth}{0.5pt}\end{center}

\section{4. MTR (STM32)}\label{mtr-stm32}

Role: dedicated motor controller (EGAS Level 1 --- Function Controller). Sole
owner of all motor-related I/O. Active in all modes --- behavior gated by
\texttt{0x110\ SYS\_MODE\_CMD} from SYS.

\subsection{4.0 Mode-Gated Behavior}\label{mode-gated-behavior}

\begin{longtable}[]{@{}
  >{\raggedright\arraybackslash}p{(\columnwidth - 8\tabcolsep) * \real{0.1154}}
  >{\centering\arraybackslash}p{(\columnwidth - 8\tabcolsep) * \real{0.1731}}
  >{\raggedright\arraybackslash}p{(\columnwidth - 8\tabcolsep) * \real{0.3269}}
  >{\raggedright\arraybackslash}p{(\columnwidth - 8\tabcolsep) * \real{0.2500}}
  >{\raggedright\arraybackslash}p{(\columnwidth - 8\tabcolsep) * \real{0.1346}}@{}}
\toprule\noalign{}
\begin{minipage}[b]{\linewidth}\raggedright
Mode
\end{minipage} & \begin{minipage}[b]{\linewidth}\centering
\texttt{0x110}
\end{minipage} & \begin{minipage}[b]{\linewidth}\raggedright
Throttle source
\end{minipage} & \begin{minipage}[b]{\linewidth}\raggedright
Gear source
\end{minipage} & \begin{minipage}[b]{\linewidth}\raggedright
Notes
\end{minipage} \\
\midrule\noalign{}
\endhead
\bottomrule\noalign{}
\endlastfoot
Manual & 0 & Read \texttt{THROTTLE\_ADC} directly & Pass-through
\texttt{GEAR\_x\_SENSE} → relays & Rider directly controls motor \\
Auto & 1 & Follow CAN \texttt{0x204} \texttt{RT\_MotorSpeed} & Follow CAN
\texttt{0x204} \texttt{RT\_Gear} & RT kinematics drives motor \\
ESTOP & 2 & DAC = 0 (cut throttle) & All gear relays off & Hardware kill + CAN
kill \\
\end{longtable}

\subsection{4.1 CAN Inputs --- MTR receives}\label{can-inputs-mtr-receives}

\begin{longtable}[]{@{}llllll@{}}
\toprule\noalign{}
CAN ID & Message & Variable & Type & Range / Enum & Source \\
\midrule\noalign{}
\endhead
\bottomrule\noalign{}
\endlastfoot
\texttt{0x001} & SAFETY\_ESTOP & (no payload) & DLC=0 & --- & any \\
\texttt{0x110} & SYS\_MODE\_CMD & \texttt{SYS\_Mode} & u8 enum & \{0=M, 1=A,
2=E\} & SYS \\
\texttt{0x204} & RT\_DRIVE\_CMD & \texttt{RT\_MotorSpeed} & i32 (mm/s) &
{[}-500, 3000{]} & RT \\
\texttt{0x204} & RT\_DRIVE\_CMD & \texttt{RT\_Gear} & u8 enum & \{N,D,S,R\} &
RT \\
\texttt{0x7FD} & RT\_HEARTBEAT & \texttt{alive\_ctr} & u8 & --- & RT \\
\end{longtable}

\subsection{4.2 Physical Inputs --- MTR}\label{physical-inputs-mtr}

\begin{longtable}[]{@{}
  >{\raggedright\arraybackslash}p{(\columnwidth - 8\tabcolsep) * \real{0.2564}}
  >{\raggedright\arraybackslash}p{(\columnwidth - 8\tabcolsep) * \real{0.1282}}
  >{\raggedright\arraybackslash}p{(\columnwidth - 8\tabcolsep) * \real{0.1538}}
  >{\raggedright\arraybackslash}p{(\columnwidth - 8\tabcolsep) * \real{0.2821}}
  >{\raggedright\arraybackslash}p{(\columnwidth - 8\tabcolsep) * \real{0.1795}}@{}}
\toprule\noalign{}
\begin{minipage}[b]{\linewidth}\raggedright
Variable
\end{minipage} & \begin{minipage}[b]{\linewidth}\raggedright
Pin
\end{minipage} & \begin{minipage}[b]{\linewidth}\raggedright
Type
\end{minipage} & \begin{minipage}[b]{\linewidth}\raggedright
Electrical
\end{minipage} & \begin{minipage}[b]{\linewidth}\raggedright
Notes
\end{minipage} \\
\midrule\noalign{}
\endhead
\bottomrule\noalign{}
\endlastfoot
\texttt{CAN\_RX} & TBD & CAN & CAN transceiver RXD & Low-side CAN bus \\
\texttt{ESTOP\_BTN} & TBD & Digital NC & Shared with SYS GPIO1 & Hardwired Level
3 kill --- zero CAN delay \\
\texttt{THROTTLE\_ADC} & TBD & ADC & Voltage divider 5V→3.3V & Throttle grip
0--5V \\
\texttt{GEAR\_D\_SENSE} & TBD & Digital & TLP281 optocoupler ch1 & Gear D 72V
sense \\
\texttt{GEAR\_S\_SENSE} & TBD & Digital & TLP281 optocoupler ch2 & Gear S 72V
sense \\
\texttt{GEAR\_R\_SENSE} & TBD & Digital & TLP281 optocoupler ch3 & Gear R 72V
sense \\
\end{longtable}

\subsection{4.3 CAN Outputs --- MTR sends}\label{can-outputs-mtr-sends}

\begin{longtable}[]{@{}
  >{\raggedright\arraybackslash}p{(\columnwidth - 12\tabcolsep) * \real{0.1333}}
  >{\raggedright\arraybackslash}p{(\columnwidth - 12\tabcolsep) * \real{0.1500}}
  >{\raggedright\arraybackslash}p{(\columnwidth - 12\tabcolsep) * \real{0.1667}}
  >{\raggedright\arraybackslash}p{(\columnwidth - 12\tabcolsep) * \real{0.1000}}
  >{\raggedright\arraybackslash}p{(\columnwidth - 12\tabcolsep) * \real{0.2167}}
  >{\raggedright\arraybackslash}p{(\columnwidth - 12\tabcolsep) * \real{0.1000}}
  >{\raggedright\arraybackslash}p{(\columnwidth - 12\tabcolsep) * \real{0.1333}}@{}}
\toprule\noalign{}
\begin{minipage}[b]{\linewidth}\raggedright
CAN ID
\end{minipage} & \begin{minipage}[b]{\linewidth}\raggedright
Message
\end{minipage} & \begin{minipage}[b]{\linewidth}\raggedright
Variable
\end{minipage} & \begin{minipage}[b]{\linewidth}\raggedright
Type
\end{minipage} & \begin{minipage}[b]{\linewidth}\raggedright
Range / Enum
\end{minipage} & \begin{minipage}[b]{\linewidth}\raggedright
Rate
\end{minipage} & \begin{minipage}[b]{\linewidth}\raggedright
Target
\end{minipage} \\
\midrule\noalign{}
\endhead
\bottomrule\noalign{}
\endlastfoot
\texttt{0x120} & SYS\_THROTTLE\_STS & \texttt{actual\_speed\_mmps} & i16 (mm/s)
& --- & 100 Hz & RT (→ Host) \\
\texttt{0x206} & MTR\_MOTOR\_FBK & \texttt{actual\_speed\_mmps} & i16 (mm/s) &
--- & 50 Hz & SYS, RT \\
\texttt{0x206} & MTR\_MOTOR\_FBK & \texttt{gear\_state} & u8 enum & \{N,D,S,R\}
& 50 Hz & SYS, RT \\
\texttt{0x206} & MTR\_MOTOR\_FBK & \texttt{fault\_flags} & u8 bitmask &
bit0=ESTOP\_ACTIVE & 50 Hz & SYS, RT \\
\end{longtable}

\subsection{4.4 Physical Outputs --- MTR}\label{physical-outputs-mtr}

\begin{longtable}[]{@{}lllll@{}}
\toprule\noalign{}
Variable & Pin & Type & Connected To & Notes \\
\midrule\noalign{}
\endhead
\bottomrule\noalign{}
\endlastfoot
\texttt{CAN\_TX} & TBD & CAN & CAN transceiver TXD & Low-side CAN bus \\
\texttt{MCP4725\_SDA} & TBD & I2C & MCP4725 DAC (addr 0x60) & Throttle 0--5V
output \\
\texttt{MCP4725\_SCL} & TBD & I2C & MCP4725 DAC & I2C clock \\
\texttt{GEAR\_D\_OUT} & TBD & Digital & Relay channel 1 & 72V to ECU gear D \\
\texttt{GEAR\_S\_OUT} & TBD & Digital & Relay channel 2 & 72V to ECU gear S \\
\texttt{GEAR\_R\_OUT} & TBD & Digital & Relay channel 3 & 72V to ECU gear R \\
\end{longtable}

\begin{center}\rule{0.5\linewidth}{0.5pt}\end{center}

\section{5. External Actuators (non-host CAN
nodes)}\label{external-actuators-non-host-can-nodes}

\subsection{5.1 SYNTREE EPS-C (Steering
Actuator)}\label{syntree-eps-c-steering-actuator}

\textbf{CAN Inputs} (receives \texttt{0x169\ VCU\_SES\_REQ} from RT):

\begin{longtable}[]{@{}lll@{}}
\toprule\noalign{}
Variable & Type & Range \\
\midrule\noalign{}
\endhead
\bottomrule\noalign{}
\endlastfoot
\texttt{VCU\_SES\_Alignment\_Enable} & bool & 0/1 \\
\texttt{VCU\_SES\_Control\_Enable} & bool & 0/1 \\
\texttt{VCU\_SES\_Control\_Mode} & u8 & --- \\
\texttt{VCU\_SES\_Tgt\_StrAngle} & i16 (0.1°/bit) & {[}-780, 780{]} \\
\texttt{VCU\_SES\_Tgt\_StrAngleSpd} & u8 (°/s) & --- \\
\texttt{VCU\_Veh\_Spd\_Value} & u8 (km/h) & 0--255 \\
\texttt{roll\_cnt\_enable} & bool & must be 1 \\
\texttt{checksum\_enable} & bool & must be 1 \\
\texttt{VCU\_SES\_RollCnt} & u8 (low 4 bits) & 0--15 \\
\texttt{VCU\_SES\_CheckSum} & u8 & XOR(0..6) \\
\end{longtable}

\textbf{CAN Outputs} (sends \texttt{0x201\ SES\_STATUS} to RT):

\begin{longtable}[]{@{}lll@{}}
\toprule\noalign{}
Variable & Type & Range \\
\midrule\noalign{}
\endhead
\bottomrule\noalign{}
\endlastfoot
\texttt{SES\_INF\_Angle\_Status} & bool & 0/1 \\
\texttt{SES\_Control\_Mode\_Status} & u8 enum & --- \\
\texttt{SES\_Error\_Status} & u8 enum & \{0=N, 1=L1, 2=L2, 3=L3\} \\
\texttt{SES\_StrAngle} & i16 (0.1°/bit) & --- \\
\texttt{SES\_Tgt\_StrAngleSpd} & i16 (0.5°/s/bit) & --- \\
\texttt{EPS\_SteeringWheel\_Torq} & u8 (Nm) & --- \\
\texttt{SES\_RollCnt\_Enable\_Status} & bool & 0/1 \\
\texttt{SES\_CheckSum\_Enable\_Status} & bool & 0/1 \\
\texttt{SES\_RollCnt\_Status} & u8 (low 4 bits) & 0--15 \\
\texttt{SES\_CheckSum\_Status} & u8 & XOR over bytes 0--6 \\
\end{longtable}

\subsection{5.2 SYNTREE SEB (Brake Actuator)}\label{syntree-seb-brake-actuator}

\textbf{CAN Inputs} (receives \texttt{0x7B9\ VCU\_SEB\_REQ} --- mode-gated
sender: SYS in MANUAL/ESTOP, RT in AUTO):

\begin{longtable}[]{@{}lll@{}}
\toprule\noalign{}
Variable & Type & Range \\
\midrule\noalign{}
\endhead
\bottomrule\noalign{}
\endlastfoot
\texttt{VCU\_SEB\_Alignment\_Enable} & bool & 0/1 \\
\texttt{VCU\_SEB\_Control\_Enable} & bool & 0/1 \\
\texttt{VCU\_SEB\_Control\_Mode} & bool & 0=Stroke, 1=Pressure \\
\texttt{VCU\_SEB\_AutoBrake} & bool & 0/1 \\
\texttt{VCU\_SEB\_Stroke\_Value\_Req} & u16 (0.05 mm/bit) & offset -30 mm \\
\texttt{VCU\_SEB\_Pre\_Value\_Req} & u8 (0.05 MPa/bit) & --- \\
\texttt{VCU\_SEB\_RollCnt\_Enable} & bool & must be 1 \\
\texttt{VCU\_SEB\_CheckSum\_Enable} & bool & must be 1 \\
\texttt{VCU\_SEB\_RollCnt} & u8 (low 4 bits) & 0--15 \\
\texttt{VCU\_SEB\_CheckSum} & u8 & XOR(0..6) \\
\end{longtable}

\textbf{CAN Outputs} (sends \texttt{0x721\ SEB\_STATUS} to SYS):

\begin{longtable}[]{@{}lll@{}}
\toprule\noalign{}
Variable & Type & Range \\
\midrule\noalign{}
\endhead
\bottomrule\noalign{}
\endlastfoot
\texttt{SEB\_Alignment\_Status} & bool & 0/1 \\
\texttt{SEB\_Control\_Enable\_Status} & bool & 0/1 \\
\texttt{SEB\_Control\_Mode\_Status} & u8 enum & --- \\
\texttt{SEB\_AutoBrake\_Status} & bool & 0/1 \\
\texttt{SEB\_Error\_Status} & u8 enum & \{0=N, 1=L1, 2=L2, 3=L3\} \\
\texttt{SEB\_Stroke\_Value} & u16 (0.05 mm/bit) & --- \\
\texttt{SEB\_Pressure\_Value} & u8 (0.05 MPa/bit) & --- \\
\texttt{SEB\_Angle\_Value} & i16 (0.5°/bit) & {[}-150, 840{]} \\
\texttt{SEB\_RollCnt\_Enable\_Status} & bool & 0/1 \\
\texttt{SEB\_CheckSum\_Enable\_Status} & bool & 0/1 \\
\texttt{SEB\_RollCnt\_Status} & u8 (low 4 bits) & --- \\
\texttt{SEB\_CheckSum\_Status} & u8 & --- \\
\end{longtable}

\subsection{5.3 DC-DC Converter}\label{dc-dc-converter}

\textbf{CAN Inputs} (receives \texttt{0x012\ SYS\_DCDC\_CMD} from SYS):

\begin{longtable}[]{@{}lll@{}}
\toprule\noalign{}
Variable & Type & Range \\
\midrule\noalign{}
\endhead
\bottomrule\noalign{}
\endlastfoot
\texttt{SYS\_DcdcEnable} & u8 bool & 0/1 \\
\end{longtable}

\begin{center}\rule{0.5\linewidth}{0.5pt}\end{center}

\section{6. Data Flow Diagram}\label{data-flow-diagram}

\begin{verbatim}
                         HIGH-SIDE CAN BUS (500 kbps)
 ┌──────────┐                                            ┌──────────┐
 │   HOST   │                                            │    RT    │
 │ (Jetson) │                                            │ (ESP32)  │
 │          │──0x300 DriveCmd───────────────────────────→│          │
 │          │──0x301 BrakeReq───────────────────────────→│          │
 │          │──0x302 LightCmd───────────────────────────→│          │
 │          │──0x7FC JetsonHB───────────────────────────→│          │
 │          │                                            │          │
 │          │←─0x011 SafetySts───────────────────────────│          │
 │          │←─0x120 ThrottleSts─────────────────────────│          │
 │          │←─0x210 StateRpt────────────────────────────│          │
 │          │←─0x220 PidRpt──────────────────────────────│          │
 │          │←─0x400 ObstacleRpt─────────────────────────│          │
 │          │←─0x600 DiagRpt─────────────────────────────│          │
 │          │←─0x7FD RTHB────────────────────────────────│          │
 └──────────┘                                            └─────┬────┘
                                                              │
                        LOW-SIDE CAN BUS (500 kbps)           │
       ┌──────────────────────────────────────────────────────┤
       │                                                      │
 ┌─────┴──────┐    ┌──────────┐    ┌──────────┐    ┌─────────┴──┐
 │    SYS     │    │   MTR    │    │  EPS-C   │    │    SEB     │
 │  (ESP32)   │    │ (STM32)  │    │(Steering)│    │  (Brake)  │
 │            │    │          │    │          │    │           │
 │→0x110 Mode─┼───→│ 0x110    │    │          │    │           │
 │→0x7B9 REQ──┼────┼──────────┼────┼──────────┼───→│←0x7B9────│
 │←0x721 STS──┼────┼──────────┼────┼──────────┼────│→0x721────│
 │            │    │          │    │←0x169────│    │           │
 │            │    │Motor DAC │    │→0x201────│    │           │
 │            │    │Gear rel. │    │          │    │           │
 │→Lights     │    │→0x120───→│    │          │    │           │
 │→DCDC 0x012 │    │→0x206───→│    │          │    │           │
 │→0x7FE HB──→│    │          │    │          │    │           │
 │←0x206 FBK──│    │          │    │          │    │           │
 └────────────┘    └──────────┘    └──────────┘    └───────────┘
\end{verbatim}

\subsection{Key paths}\label{key-paths}

\begin{longtable}[]{@{}
  >{\raggedright\arraybackslash}p{(\columnwidth - 2\tabcolsep) * \real{0.3158}}
  >{\raggedright\arraybackslash}p{(\columnwidth - 2\tabcolsep) * \real{0.6842}}@{}}
\toprule\noalign{}
\begin{minipage}[b]{\linewidth}\raggedright
Path
\end{minipage} & \begin{minipage}[b]{\linewidth}\raggedright
Description
\end{minipage} \\
\midrule\noalign{}
\endhead
\bottomrule\noalign{}
\endlastfoot
SYS → MTR (0x110) & Mode gating: Manual → MTR reads ADC/gear; Auto → MTR follows
CAN 0x204; ESTOP → cut all \\
Host → RT → MTR (0x204) & Autonomous drive: speed + gear from RT kinematics →
MTR DAC/relays \\
Host → RT → SEB (0x7B9) & AUTO brake: 1-hop from RT kinematics to SEB \\
SYS → SEB (0x7B9) & MANUAL/ESTOP brake: brake lever → SYS GPIO2 → SEB CAN \\
RT → EPS-C (0x169) & Steering angle commands from RT to steering actuator \\
MTR → RT → Host (0x120) & Speed telemetry upstream for autonomy stack \\
MTR → SYS (0x206) & EGAS Level 2 monitoring: SYS compares 0x204 cmd vs 0x206
actual \\
SYS → RT → Host (0x011, 0x600) & Safety status and diagnostics forwarded through
RT gateway \\
\end{longtable}


\chapter{Latency Issues}\label{latency-issues}

\section{Summary}\label{summary}

End-to-end latency from Jetson 0x300 (HOST\_DRIVE\_CMD) to SYS MCP4725 DAC
output, analysed by tracing every line of code in \texttt{rt-esp32/src/main.cpp}
and \texttt{sys-esp32/src/main.cpp}.

\begin{longtable}[]{@{}
  >{\raggedright\arraybackslash}p{(\columnwidth - 4\tabcolsep) * \real{0.2778}}
  >{\raggedright\arraybackslash}p{(\columnwidth - 4\tabcolsep) * \real{0.2500}}
  >{\raggedright\arraybackslash}p{(\columnwidth - 4\tabcolsep) * \real{0.4722}}@{}}
\toprule\noalign{}
\begin{minipage}[b]{\linewidth}\raggedright
Scenario
\end{minipage} & \begin{minipage}[b]{\linewidth}\raggedright
Latency
\end{minipage} & \begin{minipage}[b]{\linewidth}\raggedright
Dominant factor
\end{minipage} \\
\midrule\noalign{}
\endhead
\bottomrule\noalign{}
\endlastfoot
Best case (all phases aligned) & \textasciitilde200 µs & CAN wire + interrupt
wake \\
Typical (normal traffic, avg alignment) & \textasciitilde23 ms & Three
\textasciitilde5 ms avg scheduling delays + 7.5 ms TX polling \\
Worst case, normal & \textasciitilde56 ms & All 10 ms periods aligned against us
+ TX jitter \\
Worst case, degraded (boot / CAN fault) & \textasciitilde540 ms & Dispatch
blocks 500 ms waiting for low-bus heartbeat \\
Stop command (zero-speed drop bug) & +500 ms & Watchdog timeout before zero
written to cmd\_q \\
\end{longtable}

CAN bus utilisation is \textasciitilde15 \% (low bus) and \textasciitilde6 \%
(high bus) --- bus contention is not a factor.

\begin{center}\rule{0.5\linewidth}{0.5pt}\end{center}

\section{1. Physical layer (all
sub-millisecond)}\label{physical-layer-all-sub-millisecond}

\begin{longtable}[]{@{}
  >{\raggedright\arraybackslash}p{(\columnwidth - 4\tabcolsep) * \real{0.2917}}
  >{\raggedright\arraybackslash}p{(\columnwidth - 4\tabcolsep) * \real{0.4583}}
  >{\raggedright\arraybackslash}p{(\columnwidth - 4\tabcolsep) * \real{0.2500}}@{}}
\toprule\noalign{}
\begin{minipage}[b]{\linewidth}\raggedright
Layer
\end{minipage} & \begin{minipage}[b]{\linewidth}\raggedright
Mechanism
\end{minipage} & \begin{minipage}[b]{\linewidth}\raggedright
Time
\end{minipage} \\
\midrule\noalign{}
\endhead
\bottomrule\noalign{}
\endlastfoot
CAN frame on wire, 5-byte DLC @ 500 kbit/s & Bit time 2 µs, \textasciitilde64
bit times + stuffing & \textasciitilde130 µs \\
MCP2515 SPI register read/write @ 10 MHz & 3-byte transaction & \textasciitilde5
µs \\
MCP2515 full-frame receive & \textasciitilde10 SPI transactions (SIDH, SIDL,
DLC, data, clear IF) & \textasciitilde70 µs \\
MCP2515 full-frame send & \textasciitilde12 SPI transactions (TXB0CTRL, SIDH,
SIDL, DLC, data, RTS) & \textasciitilde70 µs \\
TWAI transmit & Hardware TX buffer, typically free & \textasciitilde50 µs \\
MCP4725 I2C DAC write & 100--400 kHz, 3 bytes (not yet implemented --- stub
only) & \textasciitilde100--300 µs \\
\end{longtable}

\begin{center}\rule{0.5\linewidth}{0.5pt}\end{center}

\section{2. Queue / task pipeline (8 hops)}\label{queue-task-pipeline-8-hops}

\begin{verbatim}
Jetson TX 0x300
  │  High CAN bus ~130 µs
  ▼
[A] t_can_rx_high (prio 5) — MCP2515 polling
  │  xQueueSend(g_can_rx_high_q, timeout=0, depth 16)
  ▼
[B] t_dispatch (prio 4) — dual-queue router
  │  xQueueOverwrite(g_cmd_q, depth 4)
  ▼
[C] t_control (prio 4, 100 Hz) — kinematics
  │  xQueueOverwrite(g_setpoint_q, depth 4)
  ▼
[D] t_can_tx_low (prio 3) — periodic TX
  │  twai_transmit (10 ms timeout)
  │  Low CAN bus ~130 µs
  ▼
[E] task_can_rx (SYS, prio 5) — TWAI interrupt-driven
  │  xQueueSend(g_can_rx_queue, timeout=0, depth 16)
  ▼
[F] task_dispatch (SYS, prio 4)
  │  atomic store → g_setpoint_speed_mmps
  ▼
[G] task_motor (SYS, prio 4, 100 Hz)
  │  MCP4725 DAC
  ▼
Motor voltage change
\end{verbatim}

\begin{center}\rule{0.5\linewidth}{0.5pt}\end{center}

\section{3. Issue A --- MCP2515 polling, not
interrupt-driven}\label{issue-a-mcp2515-polling-not-interrupt-driven}

\textbf{File:} \texttt{rt-esp32/src/can\_driver\_mcp2515.cpp:230-277}

GPIO40 (MCP\_INT) is configured as a pull-up input but \textbf{never used as an
interrupt source}. The receive loop polls \texttt{read\_status()} via SPI,
sleeping 1 ms between checks:

\begin{verbatim}
while (true) {
    uint8_t status = read_status();        // ~5 µs SPI
    if (rx0 || rx1) { /* read frame */ }
    else {
        if (deadline exceeded) return false;
        vTaskDelay(pdMS_TO_TICKS(1));     // ← 1 ms sleep
    }
}
\end{verbatim}

\begin{longtable}[]{@{}ll@{}}
\toprule\noalign{}
Scenario & Latency \\
\midrule\noalign{}
\endhead
\bottomrule\noalign{}
\endlastfoot
Frame arrives just before \texttt{read\_status()} & \textasciitilde70 µs \\
Frame arrives just after \texttt{read\_status()} returned empty &
\textasciitilde1 ms \\
Average & \textasciitilde500 µs \\
\end{longtable}

\textbf{Fix:} Configure GPIO40 as edge-triggered ISR → FreeRTOS task
notification. Eliminates polling, reduces RX latency to interrupt + scheduling
time (\textasciitilde100 µs).

\begin{center}\rule{0.5\linewidth}{0.5pt}\end{center}

\section{4. Issue B --- Dispatch priority inversion
(CRITICAL)}\label{issue-b-dispatch-priority-inversion-critical}

\textbf{File:} \texttt{rt-esp32/src/main.cpp:70-72}

\begin{verbatim}
if (xQueueReceive(g_can_rx_low_q, &fr, 0) != pdTRUE &&
    xQueueReceive(g_can_rx_high_q, &fr, 0) != pdTRUE) {
    xQueueReceive(g_can_rx_low_q, &fr, portMAX_DELAY);  // ← blocks on low ONLY
}
\end{verbatim}

The dispatch task blocks on the \textbf{low} queue with infinite timeout. A
0x300 frame sitting in \texttt{g\_can\_rx\_high\_q} \textbf{cannot be processed}
until a low-bus frame wakes dispatch. The high bus is gated by low-bus traffic.

Low-bus periodic sources that keep dispatch alive:

\begin{longtable}[]{@{}lll@{}}
\toprule\noalign{}
Source & Rate & Max gap \\
\midrule\noalign{}
\endhead
\bottomrule\noalign{}
\endlastfoot
0x201 (EPS-C steering status) & 100 Hz & 10 ms \\
0x721 (SEB brake status) & 100 Hz & 10 ms \\
0x120 (SYS throttle status) & 100 Hz & 10 ms \\
\end{longtable}

In degraded states: - During boot (before SYNTREE modules online): \textbf{500
ms} (2 Hz heartbeats only) - If SYNTREE modules fault: \textbf{500 ms} -
Complete low-bus silence: \textbf{∞ (deadlock)}

\textbf{Fix:} Use \texttt{xQueueCreateSet} to block on both queues
simultaneously, or use task notifications from \texttt{t\_can\_rx\_low} and
\texttt{t\_can\_rx\_high} to unblock dispatch when either queue has data.

\begin{center}\rule{0.5\linewidth}{0.5pt}\end{center}

\section{5. Issue C --- Zero-speed commands silently
dropped}\label{issue-c-zero-speed-commands-silently-dropped}

\textbf{File:} \texttt{rt-esp32/src/main.cpp:96}

\begin{verbatim}
if (cmd_buf.speed_mmps) {
    xQueueOverwrite(g_cmd_q, &cmd_buf);
    g_watchdog.feed(esp_timer_get_time());
}
\end{verbatim}

A Jetson command with \texttt{speed\_mmps\ ==\ 0} (commanded stop) is silently
discarded. The previous non-zero command persists in \texttt{g\_cmd\_q}. The
watchdog at 10 Hz eventually overwrites with zero after 500 ms of staleness, but
a commanded stop can take \textbf{up to 500 ms} to reach the motor.

This also interacts with the watchdog: if Jetson sends a zero-speed command
(meaning ``I'm alive but stopped''), the watchdog is \textbf{not fed}, so it
will incorrectly flag staleness and the system will enter fault behaviour even
though the host is healthy.

\textbf{Fix:} Remove the \texttt{if\ (cmd\_buf.speed\_mmps)} guard. Always write
the command and feed the watchdog. If zero-speed must be handled specially, use
a separate flag.

\begin{center}\rule{0.5\linewidth}{0.5pt}\end{center}

\section{6. Issue D --- CAN TX polling
jitter}\label{issue-d-can-tx-polling-jitter}

\textbf{File:} \texttt{rt-esp32/src/main.cpp:128-149}

\begin{verbatim}
if (xTaskGetTickCount() - t100 >= pdMS_TO_TICKS(10)) {
    t100 = xTaskGetTickCount();
    // ... send 0x204 ...
}
vTaskDelay(pdMS_TO_TICKS(5));  // free-running 5 ms sleep
\end{verbatim}

The TX loop uses \texttt{vTaskDelay(5ms)} --- NOT \texttt{vTaskDelayUntil}. The
100 Hz window is checked at 5 ms granularity, producing 10--15 ms periods
instead of steady 10 ms:

\begin{verbatim}
Timeline (ms):  0    5   10   15   20   25   30
Poll wake:      X    X    X    X    X    X    X
100Hz check:   [=]       [===]      [===]
TX 0x204:       Y            Y            Y
                ^10ms^      ^15ms^      ^10ms^
\end{verbatim}

±5 ms jitter on 0x204 output. This cascades into the SYS 200 ms staleness check
and motor control smoothness. Also applies to the 50 Hz steering path (0x169)
where a 15 ms + 25 ms = 40 ms gap could trigger the SYNTREE EPS-C 20 ms
comm-fault timeout.

Additionally, \texttt{t\_can\_tx\_low} is priority 3 while \texttt{t\_control}
is priority 4. CAN transmission always yields to computation --- backwards for a
realtime gateway.

\textbf{Fix:} Use
\texttt{vTaskDelayUntil(\&last\_tx\_100,\ pdMS\_TO\_TICKS(10))} for the 100 Hz
loop body. Separate the 100 Hz and 50 Hz transmissions into their own
\texttt{vTaskDelayUntil}-scheduled blocks. Consider raising TX task priority to
4.

\begin{center}\rule{0.5\linewidth}{0.5pt}\end{center}

\section{7. Issue E --- SYNTREE 20 ms timeout with 50 Hz TX
(future)}\label{issue-e-syntree-20-ms-timeout-with-50-hz-tx-future}

\textbf{Files:} \texttt{rt-esp32/src/config.h:18},
\texttt{rt-esp32/src/main.cpp:142-146}

\begin{verbatim}
constexpr int kSteerCmdRateHz = 50;  // SYNTREE 20 ms period
\end{verbatim}

The SYNTREE EPS-C triggers an internal comm fault if 0x169 frames stop for
\textgreater20 ms. Transmission at exactly 50 Hz (20 ms period) has \textbf{zero
margin}. Any jitter --- from the 5 ms TX polling granularity, control task
preemption, or CAN arbitration --- pushes the inter-frame gap past 20 ms.

When combined with Issue D (±5 ms jitter), the effective period is 15--25 ms,
guaranteeing periodic comm faults on the EPS-C.

\textbf{Note:} The steering path is not yet wired (see §9). This issue will
manifest when steering is connected.

\textbf{Fix:} Set \texttt{kSteerCmdRateHz\ =\ 100} (10 ms period, 10 ms margin).
At minimum, use 67 Hz (15 ms period, 5 ms margin).

\begin{center}\rule{0.5\linewidth}{0.5pt}\end{center}

\section{8. Issue F --- No WCET analysis or timing
instrumentation}\label{issue-f-no-wcet-analysis-or-timing-instrumentation}

No task execution time is measured anywhere in the codebase. There are no
\texttt{esp\_timer\_get\_time()} calls for profiling, no WCET comments, and no
period overrun detection.

Specifically unknown: - Kinematics compute time (atan2, clamp, multiply ---
likely \textless10 µs but unmeasured) - MCP2515 full-frame send time under
contention - Any blocking in I2C DAC writes (when implemented) - Whether any
task overruns its period under worst-case load

\textbf{Fix:} Add \texttt{esp\_timer\_get\_time()} instrumentation at key
pipeline points (MCP2515 RX → dispatch → control → TX → SYS RX → motor DAC). Log
max/min latency once per second on the diag task. Add period-overrun detection
in every \texttt{vTaskDelayUntil} loop.

\begin{center}\rule{0.5\linewidth}{0.5pt}\end{center}

\section{9. Known unwired paths}\label{known-unwired-paths}

Two latency-critical paths exist in the architecture but are not yet implemented
in code:

\begin{longtable}[]{@{}
  >{\raggedright\arraybackslash}p{(\columnwidth - 2\tabcolsep) * \real{0.4286}}
  >{\raggedright\arraybackslash}p{(\columnwidth - 2\tabcolsep) * \real{0.5714}}@{}}
\toprule\noalign{}
\begin{minipage}[b]{\linewidth}\raggedright
Path
\end{minipage} & \begin{minipage}[b]{\linewidth}\raggedright
Status
\end{minipage} \\
\midrule\noalign{}
\endhead
\bottomrule\noalign{}
\endlastfoot
0x300 → steering (0x169 via SYNTREE EPS-C) & \texttt{g\_steering.set\_target()}
never called. \texttt{steer\_angle\_mdeg} from physics model computed but
discarded. Steering state machine stuck in LISTEN\_SYNC (always receives
\texttt{INT16\_MIN}). 0x201 feedback angle parsed but not forwarded to steering
control. \\
0x301 → brake (0x205 via SYS → SYNTREE SEB) & \texttt{brake\_arbitrate()} result
cast to \texttt{(void)}. No 0x205 frame constructed or sent by RT. \\
\end{longtable}

Only the drive command path (0x300 → 0x204 → MCP4725 DAC) is functional
end-to-end. Both steering and brake paths require wiring before latency testing
can begin.

\begin{center}\rule{0.5\linewidth}{0.5pt}\end{center}

\section{10. Jitter model}\label{jitter-model}

The output timing (SYS motor DAC update) is jittered by three independent 100 Hz
oscillators that are not phase-locked:

\begin{verbatim}
RT Control:     |⏸ 10ms ⏸|⏸ 10ms ⏸|⏸ 10ms ⏸|
RT TX (0x204):  |⏸ 10-15ms ⏸|⏸ 10-15ms ⏸|
SYS Motor:      |⏸ 10ms ⏸|⏸ 10ms ⏸|⏸ 10ms ⏸|
\end{verbatim}

Variance sources: 1. RT control: \texttt{vTaskDelayUntil} --- ±0 (fixed phase,
no drift) --- \textbf{0 jitter} 2. RT TX: free-running poll --- \textbf{±5 ms
jitter} 3. RT dispatch: gated by low-bus traffic --- \textbf{0--10 ms (normal)
to 500 ms (degraded)} 4. SYS motor: \texttt{vTaskDelayUntil} --- ±0 (fixed
phase) --- \textbf{0 jitter}

Total peak-to-peak end-to-end jitter under normal conditions:
\textbf{\textasciitilde30 ms}.

\begin{center}\rule{0.5\linewidth}{0.5pt}\end{center}

\section{11. Recommendation priority}\label{recommendation-priority}

\begin{longtable}[]{@{}
  >{\raggedright\arraybackslash}p{(\columnwidth - 6\tabcolsep) * \real{0.1071}}
  >{\raggedright\arraybackslash}p{(\columnwidth - 6\tabcolsep) * \real{0.2500}}
  >{\raggedright\arraybackslash}p{(\columnwidth - 6\tabcolsep) * \real{0.3571}}
  >{\raggedright\arraybackslash}p{(\columnwidth - 6\tabcolsep) * \real{0.2857}}@{}}
\toprule\noalign{}
\begin{minipage}[b]{\linewidth}\raggedright
\#
\end{minipage} & \begin{minipage}[b]{\linewidth}\raggedright
Issue
\end{minipage} & \begin{minipage}[b]{\linewidth}\raggedright
Severity
\end{minipage} & \begin{minipage}[b]{\linewidth}\raggedright
Effort
\end{minipage} \\
\midrule\noalign{}
\endhead
\bottomrule\noalign{}
\endlastfoot
1 & Dispatch priority inversion (§4) & Critical --- up to 500 ms, potential
deadlock & Small \\
2 & Zero-speed command drop (§5) & Critical --- safety impact on stop latency &
Trivial \\
3 & Wire steering and brake paths (§9) & High --- blocks latency testing on
those paths & Medium \\
4 & MCP2515 interrupt (§3) & Medium --- unnecessary 1 ms poll delay & Small \\
5 & TX polling jitter (§4) & Medium --- ±5 ms cascading jitter & Small \\
6 & SYNTREE 20 ms margin (§7) & Medium --- will cause comm faults when steering
wired & Trivial (config) \\
7 & Timing instrumentation (§8) & Medium --- blocks validation of all other
fixes & Medium \\
\end{longtable}


\chapter{Light System Design --- E-Trike}\label{light-system-design-e-trike}

All lights controlled by SYS ESP32-S3. Design covers all modes: MANUAL, AUTO,
ESTOP. No other node directly drives lamps --- SYS is the sole output stage.

\begin{center}\rule{0.5\linewidth}{0.5pt}\end{center}

\section{1. Light inventory}\label{light-inventory}

\begin{longtable}[]{@{}lllll@{}}
\toprule\noalign{}
\# & Light & GPIO & Type & Priority \\
\midrule\noalign{}
\endhead
\bottomrule\noalign{}
\endlastfoot
1 & Brake light & 21 & OUT --- relay → 12V lamp & Safety-critical \\
2 & Left turn lamp & 18 & OUT --- relay → 12V lamp & Safety-relevant \\
3 & Right turn lamp & 19 & OUT --- relay → 12V lamp & Safety-relevant \\
4 & Headlight & 22 & OUT --- relay → 12V lamp & Non-safety \\
5 & Reverse light & \textbf{TBD} & OUT --- relay → 12V lamp & Non-safety \\
6 & Position/running lights & \textbf{TBD} & OUT --- relay → 12V lamp &
Non-safety \\
7 & AUTO mode bulb & 25 & OUT --- relay → 12V bulb & Non-safety \\
8 & MANUAL mode bulb & 26 & OUT --- relay → 12V bulb & Non-safety \\
\end{longtable}

GPIOs 8 and 9 are available for reverse and position lights. Need relay modules
(same type as turn lamps).

\begin{center}\rule{0.5\linewidth}{0.5pt}\end{center}

\section{2. Brake light --- GPIO21}\label{brake-light-gpio21}

\textbf{Safety-critical. Must work in all modes without CAN dependency.}

\begin{verbatim}
brake_light = (brake_lever_pressed)        // GPIO2 — physical lever
           OR (mode == ESTOP)              // ESTOP — full brake
           OR g_light_state.brake_light;   // CAN 0x302 bit2 — Jetson supplemental
\end{verbatim}

All three sources are local to SYS. The Jetson CAN bit is supplemental only ---
it can illuminate the brake light predictively (obstacle detected before
pressure builds) but can never be the only reason the brake light is OFF when
physical braking is active.

ESTOP forces the brake light ON regardless of all other inputs.

\begin{center}\rule{0.5\linewidth}{0.5pt}\end{center}

\section{3. Turn signals --- GPIO18/19}\label{turn-signals-gpio1819}

\subsection{3.1 MANUAL mode --- handlebar
switches}\label{manual-mode-handlebar-switches}

\begin{longtable}[]{@{}
  >{\raggedright\arraybackslash}p{(\columnwidth - 6\tabcolsep) * \real{0.2857}}
  >{\raggedright\arraybackslash}p{(\columnwidth - 6\tabcolsep) * \real{0.2143}}
  >{\raggedright\arraybackslash}p{(\columnwidth - 6\tabcolsep) * \real{0.2143}}
  >{\raggedright\arraybackslash}p{(\columnwidth - 6\tabcolsep) * \real{0.2857}}@{}}
\toprule\noalign{}
\begin{minipage}[b]{\linewidth}\raggedright
Switch
\end{minipage} & \begin{minipage}[b]{\linewidth}\raggedright
GPIO
\end{minipage} & \begin{minipage}[b]{\linewidth}\raggedright
Type
\end{minipage} & \begin{minipage}[b]{\linewidth}\raggedright
Action
\end{minipage} \\
\midrule\noalign{}
\endhead
\bottomrule\noalign{}
\endlastfoot
Left turn & 3 & Momentary, active-low, pull-up & Press → toggles left blink.
Press again → stops. \\
Right turn & 6 & Momentary, active-low, pull-up & Press → toggles right blink.
Press again → stops. \\
\end{longtable}

Toggle behavior: - Press once: blinking starts (500ms on/off) - Press same side
again: blinking stops - Press opposite side: cancels current side, starts new
side - Both pressed simultaneously: hazard flashers (sync blink)

No CAN involvement. Works even if both CAN buses are dead.

\subsection{3.2 AUTO mode --- two-tier
control}\label{auto-mode-two-tier-control}

\begin{verbatim}
                    ┌──────────────────────┐
                    │  0x302 HOST_LIGHT_CMD │
                    │  (Jetson → RT → SYS) │
                    │  bit0: left_turn     │
                    │  bit1: right_turn    │
                    │  bit4: auto_turn_en  │
                    └──────────┬───────────┘
                               │
            ┌──────────────────┼──────────────────┐
            ▼                  ▼                  ▼
    auto_turn_en = 0    auto_turn_en = 1    ESTOP
    Jetson direct       Low-level auto      all OFF
\end{verbatim}

\textbf{Tier 1 --- Low-level auto-logic} (\texttt{auto\_turn\_en\ =\ 1}):

RT monitors \texttt{SES\_StrAngle} from EPS-C status (\texttt{0x201}) at 100 Hz.
Logic runs on RT:

\begin{verbatim}
If abs(angle) > 15° for > 500ms → set turn signal direction
If abs(angle) < 5° for > 1000ms → cancel turn signal
\end{verbatim}

RT communicates the result to SYS by originating \texttt{0x302} directly on the
low bus --- same CAN ID it normally forwards from Jetson on the high bus. The
turn bits (0-1) are set based on steering angle. RT merges its auto-turn bits
with forwarded Jetson bits via bitwise OR --- Jetson can always add a turn
signal even when auto-logic hasn't triggered.

This logic handles lane changes (short angle spikes don't trigger), gentle
curves (angle stays elevated), and return-to-center after a turn. The thresholds
(15°, 5°) and timings (500ms, 1000ms) are configurable via RT config.h.

\textbf{Tier 2 --- Jetson direct control} (\texttt{auto\_turn\_en\ =\ 0}):

Jetson sets bits 0-1 of \texttt{0x302}. SYS uses them directly. No auto-logic.
Jetson's planning stack knows intended trajectories and can signal BEFORE the
turn starts (predictive signaling). This is the default when auto\_turn\_en is
0.

\textbf{Priority when both active:} \texttt{auto\_turn\_en\ =\ 1} uses OR logic
--- auto-turn bits OR Jetson bits. Jetson can always add a signal; it just can't
suppress one that auto-logic has triggered (intentional --- better to signal
unnecessarily than miss a signal).

\subsection{3.3 Blink pattern (all modes)}\label{blink-pattern-all-modes}

500ms ON, 500ms OFF. Managed by \texttt{lights\_task} at 20 Hz. One timer for
both sides --- they blink in phase. Both active = hazard flashers (synchronized
blink, not alternating).

\subsection{3.4 ESTOP override}\label{estop-override}

All turn signals OFF regardless of mode, switches, or CAN state.

\begin{center}\rule{0.5\linewidth}{0.5pt}\end{center}

\section{4. Headlight --- GPIO22}\label{headlight-gpio22}

\begin{longtable}[]{@{}
  >{\raggedright\arraybackslash}p{(\columnwidth - 2\tabcolsep) * \real{0.4000}}
  >{\raggedright\arraybackslash}p{(\columnwidth - 2\tabcolsep) * \real{0.6000}}@{}}
\toprule\noalign{}
\begin{minipage}[b]{\linewidth}\raggedright
Mode
\end{minipage} & \begin{minipage}[b]{\linewidth}\raggedright
Control
\end{minipage} \\
\midrule\noalign{}
\endhead
\bottomrule\noalign{}
\endlastfoot
MANUAL & Handlebar toggle switch (GPIO7). Each press toggles on/off. State
persists across mode changes. \\
AUTO & \texttt{g\_light\_state.headlight} from CAN \texttt{0x302} bit3 overrides
manual state. \\
ESTOP & OFF \\
\end{longtable}

No auto-logic needed --- headlight is purely manual or Jetson-commanded.

\begin{center}\rule{0.5\linewidth}{0.5pt}\end{center}

\section{5. Reverse light --- TBD GPIO}\label{reverse-light-tbd-gpio}

\begin{verbatim}
reverse_light = (current_gear == R)
\end{verbatim}

Gear state comes from MTR via \texttt{0x206\ MTR\_MOTOR\_FBK} (gear\_state
field, updated at 50 Hz). Works identically in MANUAL and AUTO --- whenever the
gear selector or CAN command puts the vehicle in reverse, the light illuminates.
No CAN dependency for the decision (gear state is always local to SYS after
receiving \texttt{0x206}).

Fails safe: if \texttt{0x206} stops arriving, last known gear state decays to N
after timeout → reverse light turns off.

\begin{center}\rule{0.5\linewidth}{0.5pt}\end{center}

\section{6. Position/running lights --- TBD
GPIO}\label{positionrunning-lights-tbd-gpio}

\begin{verbatim}
position_lights = (mode != ESTOP)
\end{verbatim}

ON whenever the vehicle is powered. OFF on ESTOP (12V rail may be dead anyway).
Simple state-derived --- no CAN, no mode switching. These are low-intensity
markers visible from the side, required by most vehicle regulations.

Optional CAN override: if \texttt{g\_light\_state.position\_lights} (new bit in
\texttt{0x302}) is explicitly set to 0, the lights turn off for energy saving or
stealth mode. But the default (no CAN signal present) is ON.

\begin{center}\rule{0.5\linewidth}{0.5pt}\end{center}

\section{7. Mode indicator bulbs ---
GPIO25/26}\label{mode-indicator-bulbs-gpio2526}

\begin{longtable}[]{@{}lll@{}}
\toprule\noalign{}
Mode & AUTO bulb (GPIO25) & MANUAL bulb (GPIO26) \\
\midrule\noalign{}
\endhead
\bottomrule\noalign{}
\endlastfoot
MANUAL & OFF & ON \\
AUTO & ON & OFF \\
ESTOP & OFF & OFF \\
\end{longtable}

Both OFF = ESTOP (visually distinct from both MANUAL and AUTO). Bulbs powered
from 12V accessory rail via relays --- they go dark on ESTOP regardless of MCU
state.

\begin{center}\rule{0.5\linewidth}{0.5pt}\end{center}

\section{8. CAN interface --- 0x302
HOST\_LIGHT\_CMD}\label{can-interface-0x302-host_light_cmd}

\subsection{Current (DLC=1)}\label{current-dlc1}

\begin{longtable}[]{@{}lll@{}}
\toprule\noalign{}
Byte & Bit & Signal \\
\midrule\noalign{}
\endhead
\bottomrule\noalign{}
\endlastfoot
0 & 0 & left\_turn \\
0 & 1 & right\_turn \\
0 & 2 & brake\_light \\
0 & 3 & headlight \\
0 & 4-7 & (reserved) \\
\end{longtable}

\subsection{Proposed expansion (DLC=2)}\label{proposed-expansion-dlc2}

\begin{longtable}[]{@{}llll@{}}
\toprule\noalign{}
Byte & Bit & Signal & Description \\
\midrule\noalign{}
\endhead
\bottomrule\noalign{}
\endlastfoot
0 & 0 & left\_turn & Turn signal ON \\
0 & 1 & right\_turn & Turn signal ON \\
0 & 2 & brake\_light & Supplemental brake (OR with lever + ESTOP) \\
0 & 3 & headlight & Headlight ON \\
0 & 4 & auto\_turn\_enable & Enable Tier 1 auto-turn logic (AUTO mode only) \\
0 & 5 & position\_lights & Position lights ON (0=force off, 1=ON) \\
0 & 6-7 & (reserved) & Set to 0 \\
1 & 0-7 & (reserved) & Future use \\
\end{longtable}

Byte 1 added for future needs without changing the ID. Backward compatible ---
receivers that expect DLC=1 just ignore the extra byte.

\begin{center}\rule{0.5\linewidth}{0.5pt}\end{center}

\section{9. RT auto-turn signal path}\label{rt-auto-turn-signal-path}

RT monitors \texttt{0x201\ SES\_StrAngle} at 100 Hz. When
\texttt{auto\_turn\_enable} is active, RT originates \texttt{0x302} on the low
bus:

\begin{verbatim}
RT receives 0x302 from Jetson on high bus (forward to low)
  → Stores Jetson's bits for forwarding
  → If auto_turn_en = 1:
      Run steering angle logic
      merged = Jetson_bits OR auto_turn_bits
      Send merged as 0x302 on low bus
  → Else:
      Send Jetson_bits as-is on low bus
\end{verbatim}

SYS receives \texttt{0x302} on low bus. It cannot tell whether the bits came
from Jetson (forwarded) or RT (originated) --- same ID, same payload. This is
intentional --- SYS doesn't need to know.

\begin{center}\rule{0.5\linewidth}{0.5pt}\end{center}

\section{10. ESTOP behavior}\label{estop-behavior}

\begin{longtable}[]{@{}ll@{}}
\toprule\noalign{}
Light & ESTOP state \\
\midrule\noalign{}
\endhead
\bottomrule\noalign{}
\endlastfoot
Brake & \textbf{ON} (forced) \\
Left turn & OFF \\
Right turn & OFF \\
Headlight & OFF \\
Reverse & OFF \\
Position & OFF \\
AUTO bulb & OFF \\
MANUAL bulb & OFF \\
\end{longtable}

All lamps except brake go dark. The 12V accessory relay (GPIO27) is cut on
ESTOP, so all relay-driven lights lose power at the source regardless of MCU
GPIO state.


\chapter{Listen Before Speaking (LBS) --- Safe CAN Actuator
Bootstrapping}\label{listen-before-speaking-lbs-safe-can-actuator-bootstrapping}

When a CAN actuator powers on, it has no idea what state the vehicle is in. If
the controller immediately sends a position command, the actuator will jerk to
that target --- potentially causing a dangerous steering or brake transient.

\textbf{Listen Before Speaking (LBS)} is the pattern that prevents this. The
controller waits for the actuator's first status frame, reads the \emph{current
physical position}, and sets its initial command to match. Only then does it
begin transmitting commands.

This is used for both SYNTREE actuators on our tricycle.

\begin{center}\rule{0.5\linewidth}{0.5pt}\end{center}

\section{Why LBS is necessary}\label{why-lbs-is-necessary}

CAN actuators are stateful devices that maintain their own internal control
loops:

\begin{itemize}
\tightlist
\item
  \textbf{EPS-C} (steering): runs an internal position loop. If you send
  \texttt{0x169} with angle=0° and the wheel is physically at 15°, it will slew
  to 0° at maximum motor speed. This can yank the handlebars out of the rider's
  hands.
\item
  \textbf{SEB} (brake): runs an internal pressure/stroke loop. If you command
  stroke=0 while the calipers are partially engaged, the brake releases
  abruptly.
\end{itemize}

Both actuators also \textbf{require continuous 50 Hz command frames} once
active. If frames stop for \textgreater20 ms, the actuator enters a comm-fault
timeout (lock, limp, or freewheel --- behavior is unit-specific).

\begin{center}\rule{0.5\linewidth}{0.5pt}\end{center}

\section{The LBS state machine}\label{the-lbs-state-machine}

\begin{verbatim}
Power-on
    │
    ▼
┌──────────────┐
│  BOOT_WAIT   │  500 ms delay — do NOT transmit any command frames
│              │  Actuator is powering up, may not be ready
└──────┬───────┘
       │ 500 ms elapsed
       ▼
┌──────────────┐
│ LISTEN_SYNC  │  Wait for first status frame (0x201 or 0x721)
│              │  Extract current physical position (angle or stroke)
│              │  Set active_target = current_physical_position
│              │  Wait for alignment_status bit == 1
└──────┬───────┘
       │ Status received + aligned
       ▼
┌──────────────┐
│   ACTIVE     │  Transmit command at 50 Hz continuously
│              │  First frame commands "stay where you are"
│              │  Then follow controller targets with rate limiting
└──────┬───────┘
       │ Timeout (no status for >2s) or ESTOP
       ▼
┌──────────────┐
│    FAULT     │  Stop transmitting
│              │  Actuator will timeout-fault internally
└──────────────┘
\end{verbatim}

\begin{center}\rule{0.5\linewidth}{0.5pt}\end{center}

\section{EPS-C (Steering) specifics}\label{eps-c-steering-specifics}

\begin{longtable}[]{@{}
  >{\raggedright\arraybackslash}p{(\columnwidth - 2\tabcolsep) * \real{0.6111}}
  >{\raggedright\arraybackslash}p{(\columnwidth - 2\tabcolsep) * \real{0.3889}}@{}}
\toprule\noalign{}
\begin{minipage}[b]{\linewidth}\raggedright
Parameter
\end{minipage} & \begin{minipage}[b]{\linewidth}\raggedright
Value
\end{minipage} \\
\midrule\noalign{}
\endhead
\bottomrule\noalign{}
\endlastfoot
Status ID & \texttt{0x201} SES\_STATUS (100 Hz from EPS-C) \\
Position field & \texttt{SES\_StrAngle} (int16, 0.1°/bit) \\
Alignment field & \texttt{SES\_INF\_Angle\_Status} (bit 0 of status byte) \\
Command ID & \texttt{0x169} VCU\_SES\_REQ (50 Hz to EPS-C) \\
Boot wait & 500 ms \\
Sync timeout & 2 seconds → STEER\_FAULT, remain in MANUAL \\
Mode behavior & MANUAL: do NOT send \texttt{0x169}, EPS-C standalone. AUTO: send
at 50 Hz. ESTOP: stop sending. \\
\end{longtable}

\subsection{Steering LBS in C++}\label{steering-lbs-in-c}

\begin{verbatim}
enum class SteerState {
    STEER_BOOT_WAIT,    // 500 ms power-on delay
    STEER_LISTEN_SYNC,  // Wait for 0x201, read angle, wait for aligned
    STEER_ACTIVE,       // Normal operation, transmit 0x169 at 50 Hz
    STEER_FAULT         // Timeout or ESTOP, stop transmitting
};

void steer_state_machine() {
    switch (state) {
    case STEER_BOOT_WAIT:
        if (millis_since_boot() > 500) {
            state = STEER_LISTEN_SYNC;
        }
        break;

    case STEER_LISTEN_SYNC:
        if (received_0x201) {
            // CRITICAL: match current physical position
            int16_t current_angle_raw = ses_status.angle_raw;
            active_target_angle = raw_to_mdeg(current_angle_raw);
            if (ses_status.alignment_ok) {
                state = STEER_ACTIVE;
            }
        } else if (time_in_state > 2000) {
            state = STEER_FAULT;
        }
        break;

    case STEER_ACTIVE:
        send_0x169_at_50hz();
        // Monitor following error
        if (abs(cmd_angle - actual_angle) > 5000_mdeg &&
            error_duration_ms > 300) {
            mode_set(Estop);  // stuck linkage
        }
        break;

    case STEER_FAULT:
        // Silent — let EPS-C timeout-fault internally
        break;
    }
}
\end{verbatim}

\begin{center}\rule{0.5\linewidth}{0.5pt}\end{center}

\section{SEB (Brake) specifics}\label{seb-brake-specifics}

\begin{longtable}[]{@{}
  >{\raggedright\arraybackslash}p{(\columnwidth - 2\tabcolsep) * \real{0.6111}}
  >{\raggedright\arraybackslash}p{(\columnwidth - 2\tabcolsep) * \real{0.3889}}@{}}
\toprule\noalign{}
\begin{minipage}[b]{\linewidth}\raggedright
Parameter
\end{minipage} & \begin{minipage}[b]{\linewidth}\raggedright
Value
\end{minipage} \\
\midrule\noalign{}
\endhead
\bottomrule\noalign{}
\endlastfoot
Status ID & \texttt{0x721} SEB\_STATUS (100 Hz from SEB) \\
Position field & \texttt{SEB\_Stroke\_Value} (uint16, scale 0.05, offset -30) \\
Alignment field & \texttt{SEB\_Alignment\_Status} \\
Command ID & \texttt{0x7B9} VCU\_SEB\_REQ (50 Hz to SEB) \\
Boot wait & 500 ms \\
Sync timeout & 2 seconds → BRAKE\_FAULT, brake lever inoperative until
resolved \\
\end{longtable}

The brake LBS is identical in structure to steering but runs on the SYS ESP32-S3
instead of RT.

\begin{center}\rule{0.5\linewidth}{0.5pt}\end{center}

\section{Why 500 ms boot wait?}\label{why-500-ms-boot-wait}

SYNTREE actuators run an internal bootloader and self-test on power-up. During
this window:

\begin{itemize}
\tightlist
\item
  The CAN interface may not be ready (no ACK, no status frames).
\item
  The internal position sensor may not be calibrated.
\item
  Sending commands during boot can trigger a fault latch.
\end{itemize}

500 ms is conservative for these units. Shorter times risk startup faults;
longer times delay readiness.

\begin{center}\rule{0.5\linewidth}{0.5pt}\end{center}

\section{Why this matters for safety}\label{why-this-matters-for-safety}

Without LBS, the worst case on power-up:

\begin{enumerate}
\def\labelenumi{\arabic{enumi}.}
\tightlist
\item
  Vehicle was parked with steering at full right lock (\textasciitilde40°).
\item
  Controller boots faster than EPS-C and immediately commands 0°.
\item
  EPS-C slews 40° at max motor speed --- handlebars whip to center.
\item
  If a rider is mounting the vehicle, this can cause injury.
\end{enumerate}

LBS guarantees the first command is always ``stay exactly where you are.''

\begin{center}\rule{0.5\linewidth}{0.5pt}\end{center}

\section{Comm-fault timeout (the other side of
LBS)}\label{comm-fault-timeout-the-other-side-of-lbs}

The continuous 50 Hz requirement is the \textbf{speaking} obligation once LBS
completes:

\begin{itemize}
\tightlist
\item
  EPS-C: if \texttt{0x169} stops for \textgreater20 ms → internal comm fault →
  locks or goes limp (TBD by unit spec).
\item
  SEB: if \texttt{0x7B9} stops → similar timeout.
\end{itemize}

This means the controller has an ongoing duty to transmit --- silence is
interpreted as failure. The actuator's internal safety logic handles the timeout
independently; the controller doesn't need to send an explicit ``stop'' command.

\begin{center}\rule{0.5\linewidth}{0.5pt}\end{center}

\emph{Primary reference: {[}{[}emergency-system{]}{]} for ESTOP behavior,
watchdog recovery sequence (which re-runs LBS), and testing procedures that
exercise the LBS state machine.}

\emph{See also: {[}{[}defense-in-depth-safety{]}{]} for following error and
dynamic angle clamp, {[}{[}syntree-security-protocol{]}{]} for rolling counter +
checksum, {[}{[}architecture{]}{]} §7.6 for the steering boot sequence, §8.6 for
the brake boot sequence.}


\chapter{Physics Model for the Tricycle
Platform}\label{physics-model-for-the-tricycle-platform}

This document defines the vehicle model used by
\texttt{rt-esp32/src/physics\_model.cpp}. It is the tricycle kinematic model,
with an implementation focused on converting ROS-style motion commands into
steering and rear-motor setpoints.

\section{1. Scope}\label{scope}

The current firmware uses a rigid-body, planar, non-holonomic model:

\begin{itemize}
\tightlist
\item
  rear axle midpoint is the reference point,
\item
  the front wheel provides steering,
\item
  the rear axle provides propulsion,
\item
  wheel slip is ignored in the base model.
\end{itemize}

This is sufficient for low-speed motion control, path tracking, and odometry
conversion. It does not model tire deformation, load transfer, or rollover
dynamics unless those are added separately.

\section{2. State, inputs, and geometry}\label{state-inputs-and-geometry}

State vector:

\[
\mathbf{q} = \begin{bmatrix} x \\ y \\ \theta \end{bmatrix}
\]

where:

\begin{itemize}
\tightlist
\item
  \texttt{x}, \texttt{y} are the coordinates of the rear-axle midpoint in the
  global frame,
\item
  \texttt{theta} is the vehicle heading.
\end{itemize}

Control inputs:

\[
\mathbf{u} = \begin{bmatrix} v \\ \delta \end{bmatrix}
\]

where:

\begin{itemize}
\tightlist
\item
  \texttt{v} is the forward linear velocity,
\item
  \texttt{delta} is the front-wheel steering angle.
\end{itemize}

Geometric parameters:

\begin{itemize}
\tightlist
\item
  \texttt{L} = wheelbase,
\item
  \texttt{w} = rear track width,
\item
  \texttt{r\_f} = front wheel radius,
\item
  \texttt{r\_r} = rear wheel radius.
\end{itemize}

Current RT configuration uses:

\begin{itemize}
\tightlist
\item
  \texttt{L\ =\ 1500\ mm}
\item
  \texttt{w\ =\ 800\ mm}
\item
  \texttt{r\_r\ =\ 200\ mm}
\item
  steering limit = \texttt{45\ deg}
\item
  low-speed threshold = \texttt{50\ mm/s}
\end{itemize}

\section{3. Base tricycle kinematics}\label{base-tricycle-kinematics}

For pure rolling and no lateral slip, the forward kinematics are:

\[
\dot{x} = v \cos(\theta)
\]

\[
\dot{y} = v \sin(\theta)
\]

\[
\dot{\theta} = \frac{v}{L} \tan(\delta)
\]

The non-holonomic constraint is:

\[
\dot{x} \sin(\theta) - \dot{y} \cos(\theta) = 0
\]

This states that the vehicle cannot move sideways instantaneously.

\section{4. Inverse steering solve used by the
firmware}\label{inverse-steering-solve-used-by-the-firmware}

The RT ESP32 does not command wheel angles directly from \texttt{cmd\_vel}.
Instead, it resolves the desired motion into:

\begin{itemize}
\tightlist
\item
  rear motor speed,
\item
  front steering angle,
\item
  validity flags for control logic.
\end{itemize}

The expected input type is the \texttt{DriveCmd} struct:

\begin{verbatim}
struct DriveCmd {
    int32_t speed_mmps = 0;      // linear.x  [mm/s]
    int32_t yaw_rate_mrad_s = 0; // angular.z [millirad/s]
};
\end{verbatim}

For normal motion, steering is computed from the requested yaw rate and speed:

\[
\delta = \arctan\left(\frac{L \omega}{|v|}\right)
\]

where \texttt{omega} is the commanded yaw rate.

The implementation uses \texttt{atan2(L\ *\ omega,\ abs(v))} so the sign comes
from the yaw-rate command and the zero-speed case stays numerically stable.

\subsection{Low-speed behavior}\label{low-speed-behavior}

When \texttt{\textbar{}v\textbar{}} falls below the low-speed threshold, the
steering angle is not re-estimated from the command. The firmware holds the last
valid steering angle and decays it toward zero. This avoids noisy steering
changes near standstill.

\subsection{Output limits}\label{output-limits}

The solver clamps:

\begin{itemize}
\tightlist
\item
  forward speed to \texttt{kMaxSpeedFwdMmps},
\item
  reverse speed to \texttt{kMaxSpeedRevMmps},
\item
  steering to \texttt{kSteerLimitDeg}.
\end{itemize}

The resolved output is the \texttt{ResolvedSetpoint} struct:

\begin{verbatim}
struct ResolvedSetpoint {
    int32_t motor_speed_mmps = 0; // rear motor target [mm/s]
    int32_t steer_angle_mdeg = 0; // front steer angle [millideg]
    bool steer_valid = false;
    bool reversing = false;
};
\end{verbatim}

\section{5. Wheel-level kinematics}\label{wheel-level-kinematics}

If the rear wheels are treated separately, the platform behaves like a
differential-drive odometry model on the rear axle.

Rear wheel linear speeds:

\[
v_l = \omega_{rl} r_r
\]

\[
v_r = \omega_{rr} r_r
\]

Vehicle speed and yaw rate:

\[
v = \frac{v_r + v_l}{2}
\]

\[
\dot{\theta} = \frac{v_r - v_l}{w}
\]

Forward odometry in the global frame:

\[
\dot{x} = \frac{r_r(\omega_{rr} + \omega_{rl})}{2}\cos(\theta)
\]

\[
\dot{y} = \frac{r_r(\omega_{rr} + \omega_{rl})}{2}\sin(\theta)
\]

\[
\dot{\theta} = \frac{r_r(\omega_{rr} - \omega_{rl})}{w}
\]

\section{6. Rear-wheel inverse kinematics}\label{rear-wheel-inverse-kinematics}

For a commanded \texttt{v} and \texttt{delta}, the turn radius is:

\[
R = \frac{L}{\tan(\delta)}
\]

That gives the rear-wheel path speeds:

\[
v_l = v \left(1 - \frac{w}{2L}\tan(\delta)\right)
\]

\[
v_r = v \left(1 + \frac{w}{2L}\tan(\delta)\right)
\]

Converted to wheel angular velocity:

\[
\omega_{rl} = \frac{v_l}{r_r}
\]

\[
\omega_{rr} = \frac{v_r}{r_r}
\]

This matters when the rear axle is driven by independent motors. With a solid
axle and differential, the hardware resolves the speed difference mechanically.

\section{7. Front-wheel driven form}\label{front-wheel-driven-form}

If the front wheel is driven and measured directly, its linear speed is:

\[
v_f = \omega_f r_f
\]

The rear-axle forward speed becomes:

\[
v = v_f \cos(\delta)
\]

and the state derivatives are:

\[
\dot{x} = (\omega_f r_f)\cos(\delta)\cos(\theta)
\]

\[
\dot{y} = (\omega_f r_f)\cos(\delta)\sin(\theta)
\]

\[
\dot{\theta} = \frac{\omega_f r_f}{L}\sin(\delta)
\]

\section{8. Dynamic limits and rollover}\label{dynamic-limits-and-rollover}

The kinematic model is not enough at higher speeds. A tricycle has a high center
of gravity and a narrow track, so lateral acceleration can trigger rollover.

Lateral acceleration during a turn:

\[
a_y = v\dot{\theta} = \frac{v^2}{L}\tan(\delta)
\]

The rollover threshold is approximately:

\[
\frac{v^2}{L}\tan(\delta) > \frac{g w}{2h}
\]

where:

\begin{itemize}
\tightlist
\item
  \texttt{m} = vehicle mass,
\item
  \texttt{h} = center-of-gravity height,
\item
  \texttt{w} = rear track width,
\item
  \texttt{g} = gravity.
\end{itemize}

For an advanced controller, this inequality should be enforced as a control
constraint so the command never enters an unstable region.

\section{9. Slip-angle model}\label{slip-angle-model}

At higher speed, tire slip becomes non-negligible. A linearized tire model can
use slip angles and cornering stiffness:

\[
\alpha_f = \delta - \arctan\left(\frac{\dot{y} + L_f \dot{\theta}}{\dot{x}}\right)
\]

\[
\alpha_r = -\arctan\left(\frac{\dot{y} - L_r \dot{\theta}}{\dot{x}}\right)
\]

Lateral forces:

\[
F_{yf} = C_{\alpha f}\alpha_f
\]

\[
F_{yr} = C_{\alpha r}\alpha_r
\]

Vehicle dynamics:

\[
m(\ddot{y} + \dot{x}\dot{\theta}) = F_{yf}\cos(\delta) + F_{yr}
\]

\[
I_z \ddot{\theta} = L_fF_{yf}\cos(\delta) - L_rF_{yr}
\]

This model is useful for higher-fidelity estimation or model-predictive control,
but it is not required for the current RT ESP32 solver.

\section{10. Obstacle speed limiting}\label{obstacle-speed-limiting}

The physics module also provides a simple obstacle-based speed limiter:

\begin{itemize}
\tightlist
\item
  stop completely at or below \texttt{kObstacleStopDistMM},
\item
  pass through target speed at or beyond \texttt{kObstacleClearDistMM},
\item
  interpolate linearly in between.
\end{itemize}

This is a separate safety layer from steering geometry.

\section{11. Practical summary}\label{practical-summary}

Use the base tricycle model for:

\begin{itemize}
\tightlist
\item
  command resolution,
\item
  low-speed trajectory generation,
\item
  basic odometry,
\item
  steering actuation.
\end{itemize}

Add the dynamic model only when the controller must reason about:

\begin{itemize}
\tightlist
\item
  tire slip,
\item
  high-speed cornering,
\item
  load transfer,
\item
  rollover risk.
\end{itemize}


\chapter{PID Speed Control}\label{pid-speed-control}

The RT ESP32-S3 runs a \textbf{PID controller} at 100 Hz to track the speed
setpoint computed by the tricycle kinematics solver. This closes the loop
between the desired speed (from Jetson or physics model) and the actual speed
(measured from the rear motor encoder).

\begin{center}\rule{0.5\linewidth}{0.5pt}\end{center}

\section{Why PID?}\label{why-pid}

A PID (Proportional-Integral-Derivative) controller is the standard choice for
embedded motor speed control because:

\begin{enumerate}
\def\labelenumi{\arabic{enumi}.}
\tightlist
\item
  \textbf{Simple to implement:} Three terms, multiply-accumulate, no matrix
  math.
\item
  \textbf{Well understood:} Tuning procedures are mature, behavior is
  predictable.
\item
  \textbf{Sufficient for the problem:} Motor speed control on a tricycle doesn't
  require model-predictive control or adaptive tuning --- the load (rider +
  vehicle mass) is fairly constant, and the motor controller's internal current
  loop handles the fast dynamics.
\end{enumerate}

\begin{center}\rule{0.5\linewidth}{0.5pt}\end{center}

\section{PID loop structure}\label{pid-loop-structure}

\begin{verbatim}
                    ┌──────────┐
setpoint_speed ────►│    +     │────► P × error ────────────────┐
                    │   error  │                                │
measured_speed ────►│    -     │────► I × ∫error ───────────────┤───► output ──► MCP4725 DAC
                    └──────────┘                                │
                               └────► D × d(error)/dt ──────────┘
\end{verbatim}

\subsection{Proportional term (P)}\label{proportional-term-p}

\begin{verbatim}
P_out = Kp × error
\end{verbatim}

\begin{itemize}
\tightlist
\item
  \textbf{What it does:} Produces an output proportional to the current error.
  The further from target, the harder it pushes.
\item
  \textbf{Too high:} Oscillation. The controller overshoots, corrects,
  overshoots the other way.
\item
  \textbf{Too low:} Sluggish response. The vehicle takes too long to reach
  target speed.
\end{itemize}

\subsection{Integral term (I)}\label{integral-term-i}

\begin{verbatim}
integral += error × dt
integral = clamp(integral, -max_integral, +max_integral)
I_out = Ki × integral
\end{verbatim}

\begin{itemize}
\tightlist
\item
  \textbf{What it does:} Accumulates past error to eliminate steady-state
  offset. If the vehicle is cruising 50 mm/s below target, the I term slowly
  ramps up until the error is zero.
\item
  \textbf{Integral windup:} If the motor can't physically reach the target
  (e.g., uphill against max power), the integral grows without bound. When the
  load reduces, the inflated integral causes a massive overshoot.
  \textbf{Clamping} the integral prevents this.
\item
  \textbf{Resetting on mode change:} When switching from ESTOP back to AUTO
  (after power-cycle), the integral must be zeroed. Otherwise, the accumulated
  error from the ESTOP period would cause a throttle surge.
\end{itemize}

\subsection{Derivative term (D)}\label{derivative-term-d}

\begin{verbatim}
derivative = (error - previous_error) / dt
D_out = Kd × derivative
\end{verbatim}

\begin{itemize}
\tightlist
\item
  \textbf{What it does:} Reacts to the rate of change of error. If the vehicle
  is approaching the target quickly, the D term backs off to prevent overshoot.
\item
  \textbf{Noise sensitivity:} The derivative amplifies measurement noise
  (encoder jitter). In practice, it's often computed on the measurement directly
  (\texttt{-Kd\ ×\ d(measurement)/dt}) rather than on the error --- this is
  called ``derivative on measurement'' and avoids derivative kicks on setpoint
  changes.
\end{itemize}

\begin{center}\rule{0.5\linewidth}{0.5pt}\end{center}

\section{Current tuning parameters}\label{current-tuning-parameters}

\begin{verbatim}
constexpr float kPidKp = 1.0f;
constexpr float kPidKi = 0.1f;
constexpr float kPidKd = 0.05f;
constexpr float kPidMaxIntegral = 500.0f;
constexpr int   kControlLoopHz = 100;
\end{verbatim}

\begin{longtable}[]{@{}lll@{}}
\toprule\noalign{}
Parameter & Value & Unit \\
\midrule\noalign{}
\endhead
\bottomrule\noalign{}
\endlastfoot
Kp & 1.0 & DAC counts per mm/s error \\
Ki & 0.1 & DAC counts per mm/s·s accumulated error \\
Kd & 0.05 & DAC counts per mm/s² error rate \\
Max integral & ±500 & DAC counts \\
Loop rate & 100 & Hz \\
\end{longtable}

\begin{quote}
\textbf{Note:} These gains are initial values. They assume the MCP4725 DAC
output range (0--4095) maps linearly to the motor controller's speed range.
Actual tuning requires closed-loop testing with the real motor and load.
\end{quote}

\begin{center}\rule{0.5\linewidth}{0.5pt}\end{center}

\section{Implementation}\label{implementation}

\begin{verbatim}
class SpeedPID {
public:
    void reset() {
        integral_ = 0.0f;
        prev_error_ = 0.0f;
        prev_measurement_ = 0.0f;
    }

    float update(float setpoint, float measurement, float dt) {
        // Compute error
        float error = setpoint - measurement;

        // Proportional
        float p_term = Kp_ * error;

        // Integral (with anti-windup clamping)
        integral_ += error * dt;
        integral_ = std::clamp(integral_, -max_integral_, max_integral_);
        float i_term = Ki_ * integral_;

        // Derivative on measurement (reduces derivative kick)
        float derivative = -(measurement - prev_measurement_) / dt;
        float d_term = Kd_ * derivative;

        // Sum and clamp output
        float output = p_term + i_term + d_term;
        output = std::clamp(output, 0.0f, 4095.0f);  // DAC range

        // Store state for next iteration
        prev_error_ = error;
        prev_measurement_ = measurement;

        return output;
    }

private:
    float Kp_ = 1.0f;
    float Ki_ = 0.1f;
    float Kd_ = 0.05f;
    float max_integral_ = 500.0f;
    float integral_ = 0.0f;
    float prev_error_ = 0.0f;
    float prev_measurement_ = 0.0f;
};
\end{verbatim}

The PID runs at 100 Hz in the RT \texttt{control} task:

\begin{verbatim}
void control_task() {
    TickType_t last_wake = xTaskGetTickCount();
    SpeedPID pid;
    float dt = 1.0f / 100.0f;  // 100 Hz = 10 ms

    while (true) {
        DriveCmd cmd;
        if (xQueueReceive(cmd_queue, &cmd, 0) == pdTRUE) {
            // Physics resolve: cmd → setpoint (steer angle, motor speed)
            ResolvedSetpoint sp = physics_resolve(cmd);
            setpoint = sp;
        }

        // Read current speed from encoder
        float measured_speed = encoder_read_speed_mmps();

        // PID update
        float dac_value = pid.update(setpoint.motor_speed_mmps,
                                      measured_speed, dt);

        // Push to downstream (0x204 for MTR + 0x169 for EPS-C)
        ActuatorOutput out = {
            .dac_value = static_cast<uint16_t>(dac_value),
            .gear = setpoint.gear,
            .steer_angle_mdeg = setpoint.steer_angle_mdeg,
        };
        xQueueOverwrite(setpoint_queue, &out);

        vTaskDelayUntil(&last_wake, pdMS_TO_TICKS(10));
    }
}
\end{verbatim}

\begin{center}\rule{0.5\linewidth}{0.5pt}\end{center}

\section{Tuning procedure}\label{tuning-procedure}

For the tricycle's rear motor (brushless DC, \textasciitilde1--3 kW, controlled
by an external motor controller that accepts 0--5 V analog input):

\begin{enumerate}
\def\labelenumi{\arabic{enumi}.}
\item
  \textbf{Start with P only.} Set Ki = Kd = 0. Increase Kp until the speed
  oscillates with constant amplitude. Set Kp to \textasciitilde60\% of that
  value.
\item
  \textbf{Add I.} Increase Ki until the steady-state error is eliminated within
  \textasciitilde2 seconds of a step change. Watch for overshoot --- if the
  vehicle overshoots and then slowly settles, Ki is too high.
\item
  \textbf{Add D (optional).} Increase Kd to dampen overshoot from step changes.
  On a vehicle with significant inertia, D can help. On a low-inertia system, D
  often amplifies encoder noise and is counterproductive.
\item
  \textbf{Test load rejection.} On a slight incline, command a constant speed.
  The I term should compensate for the additional load without oscillation.
\item
  \textbf{Test step response.} Command 0 → 1500 mm/s. Measure rise time
  (10--90\%), overshoot (\%), and settling time (±5\%).
\end{enumerate}

\begin{center}\rule{0.5\linewidth}{0.5pt}\end{center}

\section{What the PID does NOT control}\label{what-the-pid-does-not-control}

\begin{itemize}
\tightlist
\item
  \textbf{Steering angle:} EPS-C runs its own internal position loop. RT
  commands an angle target; EPS-C controls the motor to reach it.
\item
  \textbf{Brake pressure:} SEB runs its own internal pressure/stroke loop. SYS
  commands a stroke target; SEB controls the pump.
\item
  \textbf{Motor commutation:} The motor controller handles FOC (Field-Oriented
  Control) internally. RT only sends a 0--5 V speed reference.
\end{itemize}

The PID is the outermost speed loop. The motor controller's internal
current/torque loop is faster (typically 10--20 kHz) and handles the actual
motor phase control.

\begin{center}\rule{0.5\linewidth}{0.5pt}\end{center}

\emph{See also: {[}{[}physics-model{]}{]} for the tricycle kinematics that
produce the speed setpoint, {[}{[}actuator-interfacing{]}{]} for the MCP4725 DAC
that converts PID output to voltage, {[}{[}architecture{]}{]} §7.6 for the
control loop integration.}


\chapter{SYNTREE EPS-C --- Steer-by-Wire
Unit}\label{syntree-eps-c-steer-by-wire-unit}

CAN-controlled steering actuator. Factory-programmed IDs (not reconfigurable).

\begin{center}\rule{0.5\linewidth}{0.5pt}\end{center}

\section{1. CAN Interface}\label{can-interface}

\begin{longtable}[]{@{}ll@{}}
\toprule\noalign{}
Parameter & Value \\
\midrule\noalign{}
\endhead
\bottomrule\noalign{}
\endlastfoot
Bus & Low-level CAN (500 kbit/s) \\
Command ID & \texttt{0x169} VCU\_SES\_REQ (customized from factory default
\texttt{0x200}) \\
Status ID & \texttt{0x201} SES\_STATUS \\
Command rate & 20 ms (50 Hz) --- \textbf{continuous transmission required} \\
Status rate & 10 ms (100 Hz) \\
Endianness & Motorola LSB (little-endian) \\
DLC (both) & 8 \\
\end{longtable}

\begin{quote}
\textbf{ID note}: Factory command ID \texttt{0x200} is the factory default,
customized to \texttt{0x169}. \texttt{RT\_DRIVE\_CMD} is placed at
\texttt{0x204} to avoid collision.
\end{quote}

\begin{center}\rule{0.5\linewidth}{0.5pt}\end{center}

\section{2. Command Frame --- 0x169
VCU\_SES\_REQ}\label{command-frame-0x169-vcu_ses_req}

\begin{longtable}[]{@{}
  >{\raggedright\arraybackslash}p{(\columnwidth - 18\tabcolsep) * \real{0.1081}}
  >{\raggedright\arraybackslash}p{(\columnwidth - 18\tabcolsep) * \real{0.1486}}
  >{\raggedright\arraybackslash}p{(\columnwidth - 18\tabcolsep) * \real{0.0676}}
  >{\raggedright\arraybackslash}p{(\columnwidth - 18\tabcolsep) * \real{0.0811}}
  >{\raggedright\arraybackslash}p{(\columnwidth - 18\tabcolsep) * \real{0.0946}}
  >{\raggedright\arraybackslash}p{(\columnwidth - 18\tabcolsep) * \real{0.1081}}
  >{\raggedright\arraybackslash}p{(\columnwidth - 18\tabcolsep) * \real{0.0676}}
  >{\raggedright\arraybackslash}p{(\columnwidth - 18\tabcolsep) * \real{0.0676}}
  >{\raggedright\arraybackslash}p{(\columnwidth - 18\tabcolsep) * \real{0.0811}}
  >{\raggedright\arraybackslash}p{(\columnwidth - 18\tabcolsep) * \real{0.1757}}@{}}
\toprule\noalign{}
\begin{minipage}[b]{\linewidth}\raggedright
Signal
\end{minipage} & \begin{minipage}[b]{\linewidth}\raggedright
Start bit
\end{minipage} & \begin{minipage}[b]{\linewidth}\raggedright
Len
\end{minipage} & \begin{minipage}[b]{\linewidth}\raggedright
Type
\end{minipage} & \begin{minipage}[b]{\linewidth}\raggedright
Scale
\end{minipage} & \begin{minipage}[b]{\linewidth}\raggedright
Offset
\end{minipage} & \begin{minipage}[b]{\linewidth}\raggedright
Min
\end{minipage} & \begin{minipage}[b]{\linewidth}\raggedright
Max
\end{minipage} & \begin{minipage}[b]{\linewidth}\raggedright
Unit
\end{minipage} & \begin{minipage}[b]{\linewidth}\raggedright
Description
\end{minipage} \\
\midrule\noalign{}
\endhead
\bottomrule\noalign{}
\endlastfoot
\texttt{VCU\_SES\_Alignment\_Enable} & 0 & 1 & bool & 1 & 0 & 0 & 1 & --- & 1 =
enable calibration \\
\texttt{VCU\_SES\_Control\_Enable} & 1 & 1 & bool & 1 & 0 & 0 & 1 & --- & 1 =
enable active control (rising-edge trigger, no separate Control\_Mode signal) \\
(reserved) & 2 & 6 & --- & --- & --- & --- & --- & --- & Bits 2--7 --- not
enumerated in CSV \\
(reserved) & --- & 8 & --- & --- & --- & --- & --- & --- & Byte 1 --- not
enumerated in CSV \\
\texttt{VCU\_SES\_Tgt\_StrAngle} & 16 & 16 & i16 & 0.1 & -3000 & -700 & 700 &
deg & Target angle. Negative = left, positive = right. Raw 30000→0°. Note:
offset -3000 is unusual --- physical scale is 0.1°/bit centered at 30000 raw for
0°. \\
\texttt{VCU\_SES\_Tgt\_StrAngleSpd} & 32 & 16 & u16 & 1 & 0 & 125 & 525 & deg/s
& Maximum turning speed / slew rate \\
\texttt{VCU\_SES\_RollCnt\_Enable} & 40 & 1 & bool & 1 & 0 & 0 & 1 & --- &
\textbf{Must be 1} \\
\texttt{VCU\_SES\_CheckSum\_Enable} & 41 & 1 & bool & 1 & 0 & 0 & 1 & --- &
\textbf{Must be 1} \\
(reserved) & 42 & 2 & --- & --- & --- & --- & --- & --- & \\
\texttt{VCU\_SES\_RollCnt} & 44 & 4 & u8 & 1 & 0 & 0 & 15 & --- & Rolling
counter. Increment every frame. \\
\texttt{VCU\_Veh\_Spd\_Value} & 48 & 8 & u8 & 1 & 0 & 0 & 255 & km/h & Vehicle
speed \\
\texttt{VCU\_SES\_CheckSum} & 56 & 8 & u8 & 1 & 0 & 0 & 255 & --- & Sum of bytes
0--6, low byte only \\
\end{longtable}

\subsection{Byte layout}\label{byte-layout}

\begin{longtable}[]{@{}
  >{\raggedright\arraybackslash}p{(\columnwidth - 16\tabcolsep) * \real{0.2000}}
  >{\raggedright\arraybackslash}p{(\columnwidth - 16\tabcolsep) * \real{0.1000}}
  >{\raggedright\arraybackslash}p{(\columnwidth - 16\tabcolsep) * \real{0.1000}}
  >{\raggedright\arraybackslash}p{(\columnwidth - 16\tabcolsep) * \real{0.1000}}
  >{\raggedright\arraybackslash}p{(\columnwidth - 16\tabcolsep) * \real{0.1000}}
  >{\raggedright\arraybackslash}p{(\columnwidth - 16\tabcolsep) * \real{0.1000}}
  >{\raggedright\arraybackslash}p{(\columnwidth - 16\tabcolsep) * \real{0.1000}}
  >{\raggedright\arraybackslash}p{(\columnwidth - 16\tabcolsep) * \real{0.1000}}
  >{\raggedright\arraybackslash}p{(\columnwidth - 16\tabcolsep) * \real{0.1000}}@{}}
\toprule\noalign{}
\begin{minipage}[b]{\linewidth}\raggedright
Byte
\end{minipage} & \begin{minipage}[b]{\linewidth}\raggedright
0
\end{minipage} & \begin{minipage}[b]{\linewidth}\raggedright
1
\end{minipage} & \begin{minipage}[b]{\linewidth}\raggedright
2
\end{minipage} & \begin{minipage}[b]{\linewidth}\raggedright
3
\end{minipage} & \begin{minipage}[b]{\linewidth}\raggedright
4
\end{minipage} & \begin{minipage}[b]{\linewidth}\raggedright
5
\end{minipage} & \begin{minipage}[b]{\linewidth}\raggedright
6
\end{minipage} & \begin{minipage}[b]{\linewidth}\raggedright
7
\end{minipage} \\
\midrule\noalign{}
\endhead
\bottomrule\noalign{}
\endlastfoot
Content & Align{[}0{]}+CtrlEn{[}1{]}+rsvd{[}2-7{]} & rsvd & angle {[}7:0{]} &
angle {[}15:8{]} & speed {[}7:0{]} &
speed\href{lower\%20nibble}{11:8}+security(upper nibble) & veh\_spd &
checksum \\
\end{longtable}

\subsection{Unit conversion (internal mdeg ↔ SYNTREE raw, offset
-3000)}\label{unit-conversion-internal-mdeg-syntree-raw-offset--3000}

The CSV defines an offset of -3000 for \texttt{VCU\_SES\_Tgt\_StrAngle}, meaning
0° corresponds to a raw value of 30000.

\begin{verbatim}
SYNTREE raw = (internal_angle_mdeg / 100) + 30000    (45500 mdeg → 455 + 30000 = 30455 raw → 45.5°)
internal_mdeg = (SYNTREE raw - 30000) × 100          (30455 raw → 455 × 100 = 45500 mdeg → 45.5°)
\end{verbatim}

\begin{center}\rule{0.5\linewidth}{0.5pt}\end{center}

\section{3. Status Frame --- 0x201
SES\_STATUS}\label{status-frame-0x201-ses_status}

\begin{longtable}[]{@{}
  >{\raggedright\arraybackslash}p{(\columnwidth - 18\tabcolsep) * \real{0.1081}}
  >{\raggedright\arraybackslash}p{(\columnwidth - 18\tabcolsep) * \real{0.1486}}
  >{\raggedright\arraybackslash}p{(\columnwidth - 18\tabcolsep) * \real{0.0676}}
  >{\raggedright\arraybackslash}p{(\columnwidth - 18\tabcolsep) * \real{0.0811}}
  >{\raggedright\arraybackslash}p{(\columnwidth - 18\tabcolsep) * \real{0.0946}}
  >{\raggedright\arraybackslash}p{(\columnwidth - 18\tabcolsep) * \real{0.1081}}
  >{\raggedright\arraybackslash}p{(\columnwidth - 18\tabcolsep) * \real{0.0676}}
  >{\raggedright\arraybackslash}p{(\columnwidth - 18\tabcolsep) * \real{0.0676}}
  >{\raggedright\arraybackslash}p{(\columnwidth - 18\tabcolsep) * \real{0.0811}}
  >{\raggedright\arraybackslash}p{(\columnwidth - 18\tabcolsep) * \real{0.1757}}@{}}
\toprule\noalign{}
\begin{minipage}[b]{\linewidth}\raggedright
Signal
\end{minipage} & \begin{minipage}[b]{\linewidth}\raggedright
Start bit
\end{minipage} & \begin{minipage}[b]{\linewidth}\raggedright
Len
\end{minipage} & \begin{minipage}[b]{\linewidth}\raggedright
Type
\end{minipage} & \begin{minipage}[b]{\linewidth}\raggedright
Scale
\end{minipage} & \begin{minipage}[b]{\linewidth}\raggedright
Offset
\end{minipage} & \begin{minipage}[b]{\linewidth}\raggedright
Min
\end{minipage} & \begin{minipage}[b]{\linewidth}\raggedright
Max
\end{minipage} & \begin{minipage}[b]{\linewidth}\raggedright
Unit
\end{minipage} & \begin{minipage}[b]{\linewidth}\raggedright
Description
\end{minipage} \\
\midrule\noalign{}
\endhead
\bottomrule\noalign{}
\endlastfoot
\texttt{SES\_INF\_Angle\_Status} & 0 & 1 & bool & 1 & 0 & 0 & 1 & --- & Center
Finding Status. 0=Center Finding, 1=Found. \\
\texttt{SES\_Control\_Mode\_Status} & 1 & 2 & u8 & 1 & 0 & 0 & 3 & enum &
Control Mode Feedback. 0=Manual, 1=Automatic. \\
(unaccounted) & 3 & 3 & --- & --- & --- & --- & --- & --- & Bits 3--5 --- not
enumerated in CSV \\
\texttt{SES\_Error\_Status} & 6 & 2 & u8 & 1 & 0 & 0 & 3 & enum & Error Status.
0=Normal, 1=L1, 2=L2, 3=L3. \\
(unaccounted) & --- & 8 & --- & --- & --- & --- & --- & --- & Byte 1 --- not
enumerated in CSV \\
\texttt{SES\_StrAngle} & 16 & 16 & u16 & 0.1 & -3000 & -700 & 700 & ° & Steering
Angle. Unsigned per CSV. Raw 30000→0°. \\
\texttt{SES\_Tgt\_StrAngleSpd} & 32 & 16 & i16 & 0.5 & 0 & 0 & 1480 & °/s &
Target Angle Speed feedback. Overlaps Torq at byte 5. \\
\texttt{EPS\_SteeringWheel\_Torq} & 40 & 8 & u8 & 0.1 & -12.1 & -12 & 12 & Nm &
Steering Wheel Torque. Init 0x79 (121=0 Nm). Overlaps StrAngleSpd{[}15:8{]}. \\
\texttt{SES\_RollCnt\_Enable\_Status} & 48 & 1 & bool & 1 & 0 & 0 & 1 & --- &
Life Signal Enable Feedback. 0=Invalid, 1=Valid. \\
\texttt{SES\_CheckSum\_Enable\_Status} & 49 & 1 & bool & 1 & 0 & 0 & 1 & --- &
Checksum Enable Feedback. \\
(unaccounted) & 50 & 2 & --- & --- & --- & --- & --- & --- & \\
\texttt{SES\_RollCnt\_Status} & 52 & 4 & u8 & 1 & 0 & 0 & 15 & --- & Life Signal
Feedback. Rolling counter 0--15. \\
\texttt{SES\_CheckSum\_Status} & 56 & 8 & u8 & 1 & 0 & 0 & 255 & --- & Checksum
Feedback = XOR(bytes 0--6) \^{} 0xFF. \\
\end{longtable}

\subsection{Byte layout}\label{byte-layout-1}

\begin{longtable}[]{@{}
  >{\raggedright\arraybackslash}p{(\columnwidth - 16\tabcolsep) * \real{0.2000}}
  >{\raggedright\arraybackslash}p{(\columnwidth - 16\tabcolsep) * \real{0.1000}}
  >{\raggedright\arraybackslash}p{(\columnwidth - 16\tabcolsep) * \real{0.1000}}
  >{\raggedright\arraybackslash}p{(\columnwidth - 16\tabcolsep) * \real{0.1000}}
  >{\raggedright\arraybackslash}p{(\columnwidth - 16\tabcolsep) * \real{0.1000}}
  >{\raggedright\arraybackslash}p{(\columnwidth - 16\tabcolsep) * \real{0.1000}}
  >{\raggedright\arraybackslash}p{(\columnwidth - 16\tabcolsep) * \real{0.1000}}
  >{\raggedright\arraybackslash}p{(\columnwidth - 16\tabcolsep) * \real{0.1000}}
  >{\raggedright\arraybackslash}p{(\columnwidth - 16\tabcolsep) * \real{0.1000}}@{}}
\toprule\noalign{}
\begin{minipage}[b]{\linewidth}\raggedright
Byte
\end{minipage} & \begin{minipage}[b]{\linewidth}\raggedright
0
\end{minipage} & \begin{minipage}[b]{\linewidth}\raggedright
1
\end{minipage} & \begin{minipage}[b]{\linewidth}\raggedright
2
\end{minipage} & \begin{minipage}[b]{\linewidth}\raggedright
3
\end{minipage} & \begin{minipage}[b]{\linewidth}\raggedright
4
\end{minipage} & \begin{minipage}[b]{\linewidth}\raggedright
5
\end{minipage} & \begin{minipage}[b]{\linewidth}\raggedright
6
\end{minipage} & \begin{minipage}[b]{\linewidth}\raggedright
7
\end{minipage} \\
\midrule\noalign{}
\endhead
\bottomrule\noalign{}
\endlastfoot
Content & AngleSts{[}0{]}+ModeSts{[}1:2{]}+(gap)+Error{[}6:7{]} & (unacc.) &
StrAngle {[}7:0{]} & StrAngle {[}15:8{]} & Speed {[}7:0{]} & Speed{[}15:8{]} /
Torq {[}7:0{]} (overlap) & RollCntEn{[}0{]}+CksEn{[}1{]}+(gap)+RollCnt{[}4:7{]}
& CksSum\_Stat \\
\end{longtable}

\begin{center}\rule{0.5\linewidth}{0.5pt}\end{center}

\section{4. Control Mode --- Angle Control}\label{control-mode-angle-control}

Command the steering rack to a specific angle. The EPS-C's internal PID drives
the motor until the built-in angle sensor matches the target.

\textbf{Slew rate (\texttt{VCU\_SES\_Tgt\_StrAngleSpd})} is a required second
dimension --- unlike the brake, every steering movement must specify a maximum
turning speed in degrees/second. Low slew rate = smooth, controlled turns. High
slew rate = aggressive snap (risk of traction loss or rack damage).

\textbf{RT implementation:} Slew rate is speed-dependent. At low vehicle speed →
low slew rate for comfort. At high speed → fast slew rate for responsiveness,
but within the dynamic angle clamp (§5.2).

\begin{center}\rule{0.5\linewidth}{0.5pt}\end{center}

\section{5. Safety Mechanisms (RT
ESP32-S3)}\label{safety-mechanisms-rt-esp32-s3}

These are implemented in the RT firmware, not in the EPS-C unit itself.

\subsection{5.1 Software Hard-Stops}\label{software-hard-stops}

The EPS-C accepts commands up to ±780° (raw ±7800). The physical trike steering
rack maxes out around ±35--45°. Commanding beyond the mechanical limit will
stall the motor, potentially burning out the driver or snapping a tie-rod.

\begin{longtable}[]{@{}
  >{\raggedright\arraybackslash}p{(\columnwidth - 4\tabcolsep) * \real{0.3548}}
  >{\raggedright\arraybackslash}p{(\columnwidth - 4\tabcolsep) * \real{0.2258}}
  >{\raggedright\arraybackslash}p{(\columnwidth - 4\tabcolsep) * \real{0.4194}}@{}}
\toprule\noalign{}
\begin{minipage}[b]{\linewidth}\raggedright
Parameter
\end{minipage} & \begin{minipage}[b]{\linewidth}\raggedright
Value
\end{minipage} & \begin{minipage}[b]{\linewidth}\raggedright
Description
\end{minipage} \\
\midrule\noalign{}
\endhead
\bottomrule\noalign{}
\endlastfoot
\texttt{kSteerHardLimitDeg} & 40.0 & Software clamp, inside physical
end-stops \\
Enforcement & RT \texttt{control\_task} & Any Jetson command exceeding ±40° is
clamped before \texttt{0x169} TX \\
\end{longtable}

\subsection{5.2 Dynamic Angle Clamp (Rollover
Prevention)}\label{dynamic-angle-clamp-rollover-prevention}

A 40° turn at 2 km/h is safe. A 40° turn at 25 km/h will flip the trike. Maximum
allowable steering angle is inversely proportional to vehicle speed.

\begin{longtable}[]{@{}ll@{}}
\toprule\noalign{}
Speed & Max angle \\
\midrule\noalign{}
\endhead
\bottomrule\noalign{}
\endlastfoot
2 km/h (\textasciitilde555 mm/s) & ±40° \\
25 km/h (\textasciitilde6944 mm/s) & ±5° \\
Interpolation & Linear between these points \\
\end{longtable}

Enforcement: RT \texttt{control\_task} clamps resolved angle to
\texttt{max\_angle\ =\ f(RT\_PidMeasured)} before CAN TX. This overrides Jetson
regardless of what \texttt{/cmd\_vel} requests.

\subsection{5.3 Following Error Detection}\label{following-error-detection}

If the physical wheel gets stuck (rock jam, linkage failure) while the EPS-C
tries to turn:

\begin{longtable}[]{@{}lll@{}}
\toprule\noalign{}
Parameter & Value & Description \\
\midrule\noalign{}
\endhead
\bottomrule\noalign{}
\endlastfoot
Threshold & 5° & \texttt{\textbar{}commanded\ −\ SES\_StrAngle\textbar{}} \\
Duration & 300 ms & Must persist this long before trigger \\
Action & \texttt{mode\_set(ESTOP)} & System-wide emergency stop \\
\end{longtable}

RT compares the last commanded angle (from \texttt{0x169}) against the feedback
angle (from \texttt{0x201\ SES\_StrAngle}) every control tick.

\begin{center}\rule{0.5\linewidth}{0.5pt}\end{center}

\section{6. Security Bytes}\label{security-bytes}

\subsection{Rolling Counter}\label{rolling-counter}

Same as SEB: 4-bit value (0--15), increment every frame. Two consecutive frames
with same counter → EPS-C rejects, triggers fault.

\subsection{Checksum}\label{checksum}

Same algorithm as SEB: \texttt{XOR(bytes{[}0..6{]})\ \^{}\ 0xFF}. Placed in Byte
7.

\subsection{Security Enable Bits (Byte 5) --- Unique to
EPS-C}\label{security-enable-bits-byte-5-unique-to-eps-c}

Unlike the SEB, the EPS-C requires explicit enable bits to activate its security
checks:

\begin{itemize}
\tightlist
\item
  \texttt{roll\_cnt\_enable} (Byte 5, bit 0): \textbf{Must be 1}
\item
  \texttt{checksum\_enable} (Byte 5, bit 1): \textbf{Must be 1}
\end{itemize}

If either is 0, the unit may ignore perfectly valid rolling counter and checksum
values.

\begin{center}\rule{0.5\linewidth}{0.5pt}\end{center}

\section{7. Boot Sequence --- ``Listen Before
Speaking''}\label{boot-sequence-listen-before-speaking}

The EPS-C has an absolute encoder that retains calibration across power cycles.
No physical homing sweep is required on every boot. However, a strict software
handshake is mandatory.

\begin{verbatim}
State machine:

STEER_BOOT_WAIT:
  - 500 ms delay after power-on
  - DO NOT transmit any 0x169 frames
  - → STEER_LISTEN_SYNC

STEER_LISTEN_SYNC:
  - Wait for 0x201 SES_STATUS frame
  - Read SES_StrAngle (current physical angle)
  - CRITICAL: Set active_target_angle = current_physical_angle
    (If trike was parked with wheels at 15°, first command must be 15° — not 0°)
  - Wait for SES_INF_Angle_Status == 1 (aligned)
  - → STEER_ACTIVE

STEER_ACTIVE:
  - Transmit 0x169 at 50 Hz continuously
  - First frame commands exactly the current angle (no movement)
  - Then follow Jetson targets with dynamic clamp applied
  - Monitor following error → ESTOP on fault

STEER_FAULT:
  - Stop transmitting 0x169
  - EPS-C will timeout-fault (lock or limp — verify behavior with SYNTREE spec)
\end{verbatim}

\textbf{Critical rules:} 1. Never apply power while VCU is already sending
commands --- causes angle sensor detection error. 2. Never command 0° on boot if
the wheels are physically turned --- the resulting snap can damage the rack or
cause injury. 3. First frame after sync must command the current angle (stay
where you are). 4. \texttt{SES\_INF\_Angle\_Status} must be 1 before AUTO mode
engages. Drive motor locked out until aligned.

\begin{center}\rule{0.5\linewidth}{0.5pt}\end{center}

\section{8. One-Time Calibration}\label{one-time-calibration}

\begin{itemize}
\tightlist
\item
  \textbf{Required:} Only when the EPS-C is first mated to a new mechanical
  assembly.
\item
  \textbf{Procedure:} Trike stationary,
  \texttt{VCU\_SES\_Alignment\_Enable\ =\ 1}. Motor physically finds end-stops
  or index pulse to establish center.
\item
  \textbf{After calibration:} ECU is permanently paired to that hardware. Not
  interchangeable between vehicles.
\item
  \textbf{Power cycles:} Calibration state retained. No daily homing needed.
\item
  \textbf{Boot-time:} Only the software sync handshake (§7) is required --- no
  mechanical movement.
\end{itemize}

\begin{center}\rule{0.5\linewidth}{0.5pt}\end{center}

\section{9. C++ Implementation Guide}\label{c-implementation-guide}

\begin{verbatim}
#include <stdint.h>
#include <string.h>

// Factory default — NOT reconfigurable
#define CAN_ID_VCU_SES_REQ 0x169  // customized from factory default 0x200

typedef struct {
    uint8_t align_enable    : 1;   // Byte0,b0
    uint8_t control_enable  : 1;   // Byte0,b1 — rising-edge trigger, no separate Control_Mode
    uint8_t reserved_0      : 6;   // Byte0,b2-7

    uint8_t reserved_1;            // Byte1 — not enumerated in CSV

    // Signed i16, scale 0.1 deg/bit, offset -3000, min -700, max 700.
    // Example: 0° → 30000 (raw), 45.5° right → 30455 (raw)
    int16_t target_angle;

    uint16_t target_speed;         // Bytes4-5, u16, scale 1 deg/s, 125-525

    // Byte 5 lower nibble: target_speed[11:8]
    // Byte 5 upper nibble: security
    uint8_t roll_cnt_enable  : 1;  // Byte5,b40 — MUST be 1
    uint8_t checksum_enable  : 1;  // Byte5,b41 — MUST be 1
    uint8_t reserved_2       : 2;  // Byte5,b42-43
    uint8_t rolling_counter  : 4;  // Byte5,b44-47, 0-15

    uint8_t vehicle_speed;         // Byte6, u8, 0-255 km/h
    uint8_t checksum;              // Byte7, sum(bytes[0..6]) & 0xFF
} __attribute__((packed)) VCU_SES_Req_t;


void ses_send_command(float target_angle_deg, uint16_t speed_limit_deg_s) {
    static uint8_t roll_cnt = 0;
    VCU_SES_Req_t payload;
    memset(&payload, 0, sizeof(payload));

    // 1. Enables (no separate Control_Mode — control_enable rising-edge triggers angle control)
    payload.align_enable   = 1;
    payload.control_enable = 1;

    // 2. Security enables — MUST be 1
    payload.roll_cnt_enable = 1;
    payload.checksum_enable = 1;

    // 3. Angle target (internal mdeg → SYNTREE raw with offset -3000)
    //    internal_angle_mdeg / 100 = target_angle_deg × 10, then add 30000 offset
    payload.target_angle = (int16_t)(target_angle_deg * 10.0f + 30000);

    // 4. Slew rate (16-bit, 125-525)
    payload.target_speed = speed_limit_deg_s;

    // 5. Rolling counter
    payload.rolling_counter = roll_cnt;
    roll_cnt = (roll_cnt + 1) & 0x0F;

    // 6. Vehicle speed
    payload.vehicle_speed = 0;  // set from vehicle state

    // 7. Checksum (sum of bytes 0-6, low byte only)
    uint8_t* raw = (uint8_t*)&payload;
    uint16_t sum = 0;
    for (int i = 0; i < 7; i++) sum += raw[i];
    payload.checksum = sum & 0xFF;

    // 8. Transmit
    // CAN_Write(CAN_ID_VCU_SES_REQ, 8, raw);
}
\end{verbatim}

\begin{center}\rule{0.5\linewidth}{0.5pt}\end{center}

\section{10. Error Handling}\label{error-handling}

\begin{longtable}[]{@{}
  >{\raggedright\arraybackslash}p{(\columnwidth - 4\tabcolsep) * \real{0.2444}}
  >{\raggedright\arraybackslash}p{(\columnwidth - 4\tabcolsep) * \real{0.3333}}
  >{\raggedright\arraybackslash}p{(\columnwidth - 4\tabcolsep) * \real{0.4222}}@{}}
\toprule\noalign{}
\begin{minipage}[b]{\linewidth}\raggedright
Condition
\end{minipage} & \begin{minipage}[b]{\linewidth}\raggedright
EPS-C Response
\end{minipage} & \begin{minipage}[b]{\linewidth}\raggedright
Controller Action
\end{minipage} \\
\midrule\noalign{}
\endhead
\bottomrule\noalign{}
\endlastfoot
Rolling counter repeats & Reject frame, fault & Reset counter, re-sync \\
Checksum mismatch & Discard frame & Re-transmit with correct checksum \\
Comm timeout (\textgreater20 ms no frame) & Internal comm fault & Restart boot
sequence \\
Security enables = 0 & May ignore frame unpredictably & Firmware bug --- enforce
in struct \\
\texttt{SES\_Error\_Status\ \textgreater{}\ 0} & Degraded mode & Log, report,
ESTOP if L2/L3 \\
Following error \textgreater5° for 300 ms & --- & RT triggers system ESTOP \\
Sync timeout (no \texttt{0x201} for 2s) & --- & Remain in STEER\_FAULT, MANUAL
only \\
EPS-C timeout-fault behavior & TBD --- lock, center, or freewheel & Verify
against SYNTREE spec \\
\end{longtable}

\begin{center}\rule{0.5\linewidth}{0.5pt}\end{center}

\section{11. Integration Notes}\label{integration-notes}

\begin{itemize}
\tightlist
\item
  \textbf{Controlled by:} RT ESP32-S3 (low-level CAN)
\item
  \textbf{Mode:} Angle (1) --- only supported mode
\item
  \textbf{No daily homing.} Absolute encoder retains calibration across power
  cycles.
\item
  \textbf{Not reconfigurable.} Factory CAN IDs are fixed. EPS-C command
  \texttt{0x169} (customized from factory default \texttt{0x200}).
  \texttt{RT\_DRIVE\_CMD} is at \texttt{0x204} to avoid collision.
\item
  \textbf{MANUAL mode:} RT does not send \texttt{0x169}. EPS-C operates
  standalone from rider steering wheel input.
\item
  \textbf{AUTO mode:} RT sends \texttt{0x169} at 50 Hz with resolved angle +
  dynamic clamp + slew rate.
\item
  \textbf{ESTOP:} RT stops sending \texttt{0x169}. EPS-C timeout-faults
  (behavior TBD).
\end{itemize}


\chapter{SYNTREE CAN Security Protocol --- Rolling Counter +
Checksum}\label{syntree-can-security-protocol-rolling-counter-checksum}

SYNTREE EPS-C and SEB actuators require two security bytes in every command
frame: a \textbf{rolling counter} and an \textbf{XOR checksum}. If either is
wrong, the actuator silently discards the frame --- no error response, no fault
flag, just ignored.

This is a \textbf{liveness check}, not encryption. It proves the controller
hasn't crashed and is still computing fresh frames (not stuck in a loop
replaying the same buffer).

\begin{center}\rule{0.5\linewidth}{0.5pt}\end{center}

\section{Why SYNTREE requires this}\label{why-syntree-requires-this}

CAN is a broadcast bus. Without security bytes, a crashed controller could:

\begin{enumerate}
\def\labelenumi{\arabic{enumi}.}
\tightlist
\item
  \textbf{Replay the last valid frame forever.} The CAN peripheral's TX mailbox
  auto-retransmits on error. If the CPU hangs but the peripheral keeps running,
  the actuator sees a valid-looking 50 Hz stream with plausible angle values ---
  and never detects the fault.
\item
  \textbf{Send stale data after a watchdog reset.} On reboot, the firmware might
  reinitialize and immediately resume transmitting the last command from before
  the crash, before it has reacquired sensor data or mode state.
\end{enumerate}

Rolling counter + checksum breaks both scenarios: a crashed CPU stops
incrementing the counter (actuator detects liveness loss), and a reboot resets
the counter (actuator detects discontinuity).

\begin{center}\rule{0.5\linewidth}{0.5pt}\end{center}

\section{Frame layout (both EPS-C and
SEB)}\label{frame-layout-both-eps-c-and-seb}

Both actuators use the same 8-byte security scheme in bytes 5--7:

\begin{longtable}[]{@{}
  >{\raggedright\arraybackslash}p{(\columnwidth - 4\tabcolsep) * \real{0.2308}}
  >{\raggedright\arraybackslash}p{(\columnwidth - 4\tabcolsep) * \real{0.2692}}
  >{\raggedright\arraybackslash}p{(\columnwidth - 4\tabcolsep) * \real{0.5000}}@{}}
\toprule\noalign{}
\begin{minipage}[b]{\linewidth}\raggedright
Byte
\end{minipage} & \begin{minipage}[b]{\linewidth}\raggedright
Field
\end{minipage} & \begin{minipage}[b]{\linewidth}\raggedright
Description
\end{minipage} \\
\midrule\noalign{}
\endhead
\bottomrule\noalign{}
\endlastfoot
5 & Security enables & Bit 0 = rolling counter enable, Bit 1 = checksum enable.
Always \texttt{0x03} (both enabled). \\
6 & Rolling counter (low nibble) + reserved (high nibble) & 4-bit counter,
increments 0→15→0 every frame. High nibble = 0. \\
7 & Checksum & \texttt{XOR(bytes\ 0..6)\ \^{}\ 0xFF} \\
\end{longtable}

\subsection{Byte 5 --- security enables}\label{byte-5-security-enables}

\begin{verbatim}
Bit 0: rolling_counter_enable   (1 = enabled)
Bit 1: checksum_enable          (1 = enabled)
Bits 2-7: reserved (0)

Value: 0x03
\end{verbatim}

Both must be set to 1. If either is 0, the frame is rejected regardless of
content.

\subsection{Byte 6 --- rolling counter}\label{byte-6-rolling-counter}

\begin{verbatim}
Bits 0-3: rolling_counter (0–15)
Bits 4-7: reserved (0)

After each transmission: counter = (counter + 1) & 0x0F
\end{verbatim}

The actuator tracks the last received counter. It expects
\texttt{counter\ =\ (previous\ +\ 1)\ \&\ 0x0F}. If the counter doesn't
increment (stuck) or jumps backward (reboot), the actuator may fault --- exact
behavior is unit-specific.

\subsection{Byte 7 --- checksum}\label{byte-7-checksum}

\begin{verbatim}
uint8_t checksum = 0;
for (int i = 0; i < 7; i++) {
    checksum ^= frame_bytes[i];
}
checksum ^= 0xFF;
\end{verbatim}

Or equivalently: XOR all 7 bytes together, then invert all bits.

\begin{center}\rule{0.5\linewidth}{0.5pt}\end{center}

\section{C++ implementation}\label{c-implementation}

\begin{verbatim}
struct SyntreeSecurityFrame {
    uint8_t bytes[8];

    void set_security_enables() {
        bytes[5] = 0x03;  // both rolling counter and checksum enabled
    }

    void set_rolling_counter(uint8_t counter) {
        bytes[6] = (counter & 0x0F);  // low nibble only
    }

    void compute_checksum() {
        uint8_t cks = 0;
        for (int i = 0; i < 7; i++) {
            cks ^= bytes[i];
        }
        bytes[7] = cks ^ 0xFF;
    }

    void finalize(uint8_t counter) {
        set_security_enables();
        set_rolling_counter(counter);
        compute_checksum();
    }
};

// Usage in 50 Hz transmit loop
static uint8_t rolling_counter = 0;
SyntreeSecurityFrame cmd;

// EPS-C steering command
cmd.bytes[0] = 0x01;           // control mode = Angle
cmd.bytes[1] = angle_raw & 0xFF;      // angle low byte
cmd.bytes[2] = (angle_raw >> 8) & 0xFF; // angle high byte
cmd.bytes[3] = angle_speed;     // slew rate
cmd.bytes[4] = 0x00;            // reserved
cmd.finalize(rolling_counter);
rolling_counter = (rolling_counter + 1) & 0x0F;

twai_transmit(&cmd, portMAX_DELAY);
\end{verbatim}

\begin{center}\rule{0.5\linewidth}{0.5pt}\end{center}

\section{Verification by the actuator}\label{verification-by-the-actuator}

When EPS-C or SEB receives a frame:

\begin{enumerate}
\def\labelenumi{\arabic{enumi}.}
\tightlist
\item
  Check byte 5 bits 0 and 1 --- both must be \texttt{1}. If not → discard
  silently.
\item
  Compute \texttt{XOR(bytes\ 0..6)\ \^{}\ 0xFF} --- must match byte 7. If not →
  discard silently.
\item
  Check byte 6 low nibble --- must be \texttt{(previous\ +\ 1)\ \&\ 0x0F}. If
  stuck or wrong → may fault (unit-specific).
\item
  If all checks pass → accept the command and actuate.
\end{enumerate}

\begin{center}\rule{0.5\linewidth}{0.5pt}\end{center}

\section{Failure modes}\label{failure-modes}

\begin{longtable}[]{@{}
  >{\raggedright\arraybackslash}p{(\columnwidth - 4\tabcolsep) * \real{0.2571}}
  >{\raggedright\arraybackslash}p{(\columnwidth - 4\tabcolsep) * \real{0.2000}}
  >{\raggedright\arraybackslash}p{(\columnwidth - 4\tabcolsep) * \real{0.5429}}@{}}
\toprule\noalign{}
\begin{minipage}[b]{\linewidth}\raggedright
Failure
\end{minipage} & \begin{minipage}[b]{\linewidth}\raggedright
Cause
\end{minipage} & \begin{minipage}[b]{\linewidth}\raggedright
Actuator behavior
\end{minipage} \\
\midrule\noalign{}
\endhead
\bottomrule\noalign{}
\endlastfoot
Checksum mismatch & Corrupted frame (bus noise, bad termination) or software bug
(wrong XOR) & Frame silently discarded. Rolling counter NOT incremented. \\
Counter stuck & Controller CPU hung, CAN peripheral replaying TX mailbox &
Actuator detects no counter change. May trigger comm-fault after N repeated
frames (unit-specific timeout, typically 5--10 frames = 100--200 ms). \\
Counter jump & Controller rebooted --- counter reset to 0 while actuator
expected next value & Actuator may fault or re-sync depending on unit firmware.
Boot LBS ensures sync. \\
Security enables = 0 & Software bug (forgot to set byte 5) or endianness error &
All frames silently discarded. Actuator appears unresponsive. \\
\end{longtable}

\begin{center}\rule{0.5\linewidth}{0.5pt}\end{center}

\section{Debugging silent discards}\label{debugging-silent-discards}

If the actuator seems unresponsive despite valid-looking CAN traffic:

\begin{enumerate}
\def\labelenumi{\arabic{enumi}.}
\tightlist
\item
  Verify byte 5 = \texttt{0x03} on every frame.
\item
  Verify checksum with a CAN analyzer that can compute XOR across bytes.
\item
  Verify the rolling counter increments monotonically (0, 1, 2, \ldots, 15, 0,
  1, \ldots).
\item
  Check that byte ordering matches SYNTREE's little-endian expectation.
\end{enumerate}

The most common bug: forgetting to call \texttt{finalize()} or calling it before
setting all payload bytes (so bytes 0--4 are zero when checksum is computed,
then payload is written afterward, invalidating the checksum).

\begin{center}\rule{0.5\linewidth}{0.5pt}\end{center}

\section{Why not CRC?}\label{why-not-crc}

SYNTREE chose XOR + inversion over a proper CRC:

\begin{itemize}
\tightlist
\item
  XOR is computationally trivial (no lookup table, no polynomial division).
\item
  Combined with the rolling counter, it covers both integrity (corruption
  detection) and liveness (stuck-controller detection).
\item
  A standalone CRC only proves the frame wasn't corrupted in transit, not that
  the controller is still alive.
\end{itemize}

The downside: XOR catches single-bit errors but is weaker against multi-bit
errors than CRC-8 or CRC-16. SYNTREE's design trades stronger error detection
for simplicity and the addition of liveness checking via the counter.

\begin{center}\rule{0.5\linewidth}{0.5pt}\end{center}

\emph{Primary reference: {[}{[}emergency-system{]}{]} for how rolling counter
failure modes fit into the 8-layer safety system and trigger ESTOP.}

\emph{See also: {[}{[}listen-before-speaking{]}{]} for the boot sequence that
must complete before security bytes matter, {[}{[}architecture{]}{]} §7.6 for
EPS-C protocol, §8.6 for SEB protocol, {[}{[}can-dictionary{]}{]} for bit-level
frame layouts.}


\chapter{Throttle DAC Design --- MCP4725 I2C
DAC}\label{throttle-dac-design-mcp4725-i2c-dac}

Why MCP4725, not a digital potentiometer, not PWM.

\begin{center}\rule{0.5\linewidth}{0.5pt}\end{center}

\section{1. What the motor controller
expects}\label{what-the-motor-controller-expects}

The motor controller's throttle input is a \textbf{0--5V analog DC voltage},
same as what the physical throttle grip produces. The throttle grip is a
potentiometer wired as a voltage divider --- twist produces a steady voltage
proportional to angle.

MCP4725 mimics this by outputting a programmable DC voltage:

\begin{verbatim}
SYS ADC reads grip voltage → SYS DAC recreates voltage → Motor controller
\end{verbatim}

\begin{center}\rule{0.5\linewidth}{0.5pt}\end{center}

\section{2. Why not X9C103S (digital
potentiometer)}\label{why-not-x9c103s-digital-potentiometer}

X9C103S is a digital potentiometer --- it varies a \textbf{resistance} (0 to
10kΩ), not a voltage.

\begin{longtable}[]{@{}
  >{\raggedright\arraybackslash}p{(\columnwidth - 4\tabcolsep) * \real{0.3571}}
  >{\raggedright\arraybackslash}p{(\columnwidth - 4\tabcolsep) * \real{0.3214}}
  >{\raggedright\arraybackslash}p{(\columnwidth - 4\tabcolsep) * \real{0.3214}}@{}}
\toprule\noalign{}
\begin{minipage}[b]{\linewidth}\raggedright
Property
\end{minipage} & \begin{minipage}[b]{\linewidth}\raggedright
X9C103S
\end{minipage} & \begin{minipage}[b]{\linewidth}\raggedright
MCP4725
\end{minipage} \\
\midrule\noalign{}
\endhead
\bottomrule\noalign{}
\endlastfoot
Output type & Variable resistor & Voltage source \\
Resolution & 100 steps (\textasciitilde7-bit) & 4096 steps (12-bit) \\
Interface & Up/down pulses (relative) & I2C (set absolute value) \\
Power-on state & Unknown wiper position & Known: outputs 0V until configured \\
Set absolute value & No --- must cycle to endpoint first & Yes --- write 12-bit
value \\
External parts needed & Voltage reference, buffer op-amp & None (direct voltage
output) \\
Throttle feel & Steppy (100 discrete levels) & Smooth (4096 levels) \\
\end{longtable}

To use X9C103S, you'd need to: 1. Wire it as voltage divider with external
reference 2. On boot, pulse 100 steps down to calibrate to zero (slow, wears the
wiper) 3. Accept only 100 throttle positions --- noticeable steps to the rider

To use MCP4725: 1. Write \texttt{code\ =\ speed/3000\ ×\ 4095} via I2C 2. Output
pin produces that voltage 3. Done.

\begin{center}\rule{0.5\linewidth}{0.5pt}\end{center}

\section{3. Why not PWM from ESP32}\label{why-not-pwm-from-esp32}

ESP32's LEDC peripheral can produce PWM. With an RC low-pass filter, PWM
approximates a DC voltage. But:

\begin{longtable}[]{@{}
  >{\raggedright\arraybackslash}p{(\columnwidth - 4\tabcolsep) * \real{0.2500}}
  >{\raggedright\arraybackslash}p{(\columnwidth - 4\tabcolsep) * \real{0.4250}}
  >{\raggedright\arraybackslash}p{(\columnwidth - 4\tabcolsep) * \real{0.3250}}@{}}
\toprule\noalign{}
\begin{minipage}[b]{\linewidth}\raggedright
Property
\end{minipage} & \begin{minipage}[b]{\linewidth}\raggedright
PWM + RC filter
\end{minipage} & \begin{minipage}[b]{\linewidth}\raggedright
MCP4725 DAC
\end{minipage} \\
\midrule\noalign{}
\endhead
\bottomrule\noalign{}
\endlastfoot
Ripple & Inherent --- filter cutoff trades ripple vs response time & Near-zero
(true DC) \\
Resolution & ESP32 LEDC: up to 16-bit, but effective bits drop with ripple &
True 12-bit \\
Response time & Filter delay (e.g., 1 kHz PWM needs \textasciitilde10ms to
settle) & I2C write latency only (\textasciitilde100µs) \\
Output impedance & High (from RC filter) & Low (op-amp buffered) \\
Failure mode & PWM stuck HIGH → RC charges to 3.3V → motor sees voltage & I2C
timeout → DAC holds last value or drops to 0 depending on config \\
\end{longtable}

PWM+filter \emph{can} work. But the MCP4725 is purpose-built for this ---
cleaner signal, simpler failure analysis.

\begin{center}\rule{0.5\linewidth}{0.5pt}\end{center}

\section{4. Voltage domain bridging (3.3V vs
5V)}\label{voltage-domain-bridging-3.3v-vs-5v}

ESP32 I2C operates at 3.3V. MCP4725 needs to output 0--5V to the motor
controller.

\begin{verbatim}
ESP32-S3 (3.3V)              MCP4725 (5V domain)           Motor Controller
┌──────────────┐            ┌──────────────────┐          ┌───────────────┐
│              │            │                  │          │               │
│  I2C SDA ────┼────────────►  SDA             │          │               │
│  I2C SCL ────┼────────────►  SCL          VOUT├──────────► 0-5V throttle │
│              │            │                  │          │               │
│  3.3V        │            │  VDD = 5V        │          │               │
│              │            │  GND             │          │               │
└──────────────┘            └──────────────────┘          └───────────────┘
\end{verbatim}

\textbf{Problem:} MCP4725 at VDD=5V has VIH ≈ 0.7×5V = 3.5V. ESP32's 3.3V I2C
HIGH may not reliably cross this threshold.

\textbf{Option A: Power MCP4725 at 5V, level-shift I2C}

\begin{verbatim}
ESP32 ──► BSS138 MOSFET level shifter ──► MCP4725 VDD=5V ──► VOUT=0-5V
\end{verbatim}

Clean 0-5V output. Needs shifter module (2× BSS138 + 4× 10k pull-ups) --- cheap,
standard.

\textbf{Option B: Power MCP4725 at 3.3V, amplify output}

\begin{verbatim}
ESP32 ──► MCP4725 VDD=3.3V ──► VOUT=0-3.3V ──► op-amp ×1.52 gain ──► 0-5V
\end{verbatim}

Safe I2C levels. Adds op-amp. MCP4725 datasheet allows VDD down to 2.7V.

\textbf{Option C: Motor controller accepts 3.3V?} Some controllers treat
anything above \textasciitilde3V as ``full throttle.'' If so, no level shifting
needed. Only way to know: test it. See parallel task WO-14 (motor controller
analog input characterization).

\textbf{Recommended until proven otherwise:} Option A (5V MCP4725 + I2C level
shifter). Most reliable, standard pattern, well-documented.

\begin{center}\rule{0.5\linewidth}{0.5pt}\end{center}

\section{5. I2C details}\label{i2c-details}

\begin{longtable}[]{@{}ll@{}}
\toprule\noalign{}
Parameter & Value \\
\midrule\noalign{}
\endhead
\bottomrule\noalign{}
\endlastfoot
Device & MCP4725A0T-E/CH \\
Address & 0x60 (7-bit) \\
SDA & GPIO15 \\
SCL & GPIO16 \\
VDD & 5V (via Option A level shifter) or 3.3V (via Option B) \\
Resolution & 12-bit (0--4095) \\
Output range & 0V to VDD \\
Write command & I2C fast mode (400 kHz): 3 bytes (cmd + dataH + dataL) \\
\end{longtable}

Mapping:
\texttt{DAC\_code\ =\ abs(speed\_mmps)\ /\ kThrottleMaxSpeedMmps\ ×\ 4095}

\begin{center}\rule{0.5\linewidth}{0.5pt}\end{center}

\section{6. Failure modes}\label{failure-modes}

\begin{longtable}[]{@{}
  >{\raggedright\arraybackslash}p{(\columnwidth - 4\tabcolsep) * \real{0.2195}}
  >{\raggedright\arraybackslash}p{(\columnwidth - 4\tabcolsep) * \real{0.4146}}
  >{\raggedright\arraybackslash}p{(\columnwidth - 4\tabcolsep) * \real{0.3659}}@{}}
\toprule\noalign{}
\begin{minipage}[b]{\linewidth}\raggedright
Failure
\end{minipage} & \begin{minipage}[b]{\linewidth}\raggedright
MCP4725 behavior
\end{minipage} & \begin{minipage}[b]{\linewidth}\raggedright
System effect
\end{minipage} \\
\midrule\noalign{}
\endhead
\bottomrule\noalign{}
\endlastfoot
I2C bus hung & Holds last written value & Motor stays at last speed --- stale
setpoint check catches this (\textgreater200ms → zero) \\
MCP4725 power lost & VOUT drops to 0V & Motor stops --- safe \\
ESP32 hung, WDT resets & I2C pins go Hi-Z, pull-ups bring SDA/SCL HIGH & MCP4725
holds last value until next I2C write. WDT reset clears ESP32 → ESP32 re-inits
DAC to 0 \\
ESTOP triggered & Software writes 0 to DAC & 0V output → motor stops \\
\end{longtable}


\chapter{E-Trike Wiring Reference}\label{e-trike-wiring-reference}

Pin-to-pin wiring for both ESP32-S3 boards, CAN buses, and all peripherals.
Derived from \texttt{architecture.md} and \texttt{config.h}.

\begin{center}\rule{0.5\linewidth}{0.5pt}\end{center}

\section{1. CAN Bus Topology}\label{can-bus-topology}

Two physically separate CAN buses at 500 kbit/s. Each requires 120Ω termination
at both ends.

\subsection{Low-Level CAN Bus}\label{low-level-can-bus}

\begin{verbatim}
  RT-ESP32 (TWAI)          SYS-ESP32 (TWAI)       EPS-C          SEB         DC-DC
  TX=5  RX=4              TX=5  RX=4           (steering)      (brake)    72V→12V
   │     │                  │     │                │              │           │
   ├──┬──┘                  ├──┬──┘                │              │           │
   │  │                     │  │                   │              │           │
 CAN_H CAN_L              CAN_H CAN_L           CAN_H CAN_L   CAN_H CAN_L  CAN_H CAN_L
   │  │                     │  │                   │  │           │  │        │  │
   └──┼─────────────────────┼──┼───────────────────┼──┼───────────┼──┼────────┼──┤
      └─────────────────────┴──┴───────────────────┴──┴───────────┴──┴────────┴──┘
                                   120Ω terminator at each physical end
\end{verbatim}

\begin{longtable}[]{@{}
  >{\raggedright\arraybackslash}p{(\columnwidth - 6\tabcolsep) * \real{0.2353}}
  >{\raggedright\arraybackslash}p{(\columnwidth - 6\tabcolsep) * \real{0.2647}}
  >{\raggedright\arraybackslash}p{(\columnwidth - 6\tabcolsep) * \real{0.2941}}
  >{\raggedright\arraybackslash}p{(\columnwidth - 6\tabcolsep) * \real{0.2059}}@{}}
\toprule\noalign{}
\begin{minipage}[b]{\linewidth}\raggedright
Signal
\end{minipage} & \begin{minipage}[b]{\linewidth}\raggedright
RT GPIO
\end{minipage} & \begin{minipage}[b]{\linewidth}\raggedright
SYS GPIO
\end{minipage} & \begin{minipage}[b]{\linewidth}\raggedright
Notes
\end{minipage} \\
\midrule\noalign{}
\endhead
\bottomrule\noalign{}
\endlastfoot
CAN TX & 5 & 5 & To SN65HVD230 TXD \\
CAN RX & 4 & 4 & From SN65HVD230 RXD \\
CAN\_GND & --- & --- & \textbf{Must} connect grounds if nodes on separate power
supplies \\
\end{longtable}

\subsection{High-Level CAN Bus}\label{high-level-can-bus}

\begin{verbatim}
  Jetson Orin NX           RT-ESP32 (MCP2515 SPI)
  CAN interface           SCK=36 MOSI=37 MISO=38 CS=39 INT=40
      │                        │
   CAN_H CAN_L              CAN_H CAN_L
      │  │                     │  │
      └──┼─────────────────────┼──┘
         └─────────────────────┘
           120Ω terminator at each end
\end{verbatim}

\begin{center}\rule{0.5\linewidth}{0.5pt}\end{center}

\section{2. RT ESP32-S3 --- Input/Output
Tables}\label{rt-esp32-s3-inputoutput-tables}

\subsection{2.1 INPUTS (signals RT reads from external
hardware)}\label{inputs-signals-rt-reads-from-external-hardware}

\begin{longtable}[]{@{}
  >{\raggedright\arraybackslash}p{(\columnwidth - 10\tabcolsep) * \real{0.0698}}
  >{\raggedright\arraybackslash}p{(\columnwidth - 10\tabcolsep) * \real{0.1860}}
  >{\raggedright\arraybackslash}p{(\columnwidth - 10\tabcolsep) * \real{0.1395}}
  >{\raggedright\arraybackslash}p{(\columnwidth - 10\tabcolsep) * \real{0.1395}}
  >{\raggedright\arraybackslash}p{(\columnwidth - 10\tabcolsep) * \real{0.3023}}
  >{\raggedright\arraybackslash}p{(\columnwidth - 10\tabcolsep) * \real{0.1628}}@{}}
\toprule\noalign{}
\begin{minipage}[b]{\linewidth}\raggedright
\#
\end{minipage} & \begin{minipage}[b]{\linewidth}\raggedright
Signal
\end{minipage} & \begin{minipage}[b]{\linewidth}\raggedright
GPIO
\end{minipage} & \begin{minipage}[b]{\linewidth}\raggedright
Type
\end{minipage} & \begin{minipage}[b]{\linewidth}\raggedright
Connected To
\end{minipage} & \begin{minipage}[b]{\linewidth}\raggedright
Notes
\end{minipage} \\
\midrule\noalign{}
\endhead
\bottomrule\noalign{}
\endlastfoot
1 & CAN RX (low bus) & 4 & TWAI RX & SN65HVD230 RXD pin & Low-level CAN: SYS,
EPS-C, SEB, DC-DC \\
2 & MCP2515 MISO & 38 & SPI & MCP2515 SO pin & High-level CAN RX data \\
3 & MCP2515 INT & 40 & Digital & MCP2515 INT pin & Interrupt: RX buffer ready \\
4 & Encoder A (rear motor) & 1 & PCNT & Encoder sensor & Speed feedback ---
\textbf{active} \\
5 & Encoder B (rear motor) & 2 & PCNT & Encoder sensor & Quadrature phase B \\
6 & Encoder A (front wheel) & 3 & PCNT & Encoder sensor & \textbf{Sensor TBD} \\
7 & Encoder B (front wheel) & 6 & PCNT & Encoder sensor & \textbf{Sensor TBD} \\
8 & Encoder A (rear left) & 9 & PCNT & Encoder sensor & \textbf{Sensor TBD} \\
9 & Encoder B (rear left) & 12 & PCNT & Encoder sensor & \textbf{Sensor TBD} \\
10 & Encoder A (rear right) & 13 & PCNT & Encoder sensor & \textbf{Sensor
TBD} \\
11 & Encoder B (rear right) & 14 & PCNT & Encoder sensor & \textbf{Sensor
TBD} \\
12 & I2C SDA & 10 & I2C & IMU (optional) & Pull-up to 3.3V required \\
\end{longtable}

\subsection{2.2 OUTPUTS (signals RT drives to external
hardware)}\label{outputs-signals-rt-drives-to-external-hardware}

\begin{longtable}[]{@{}
  >{\raggedright\arraybackslash}p{(\columnwidth - 10\tabcolsep) * \real{0.0698}}
  >{\raggedright\arraybackslash}p{(\columnwidth - 10\tabcolsep) * \real{0.1860}}
  >{\raggedright\arraybackslash}p{(\columnwidth - 10\tabcolsep) * \real{0.1395}}
  >{\raggedright\arraybackslash}p{(\columnwidth - 10\tabcolsep) * \real{0.1395}}
  >{\raggedright\arraybackslash}p{(\columnwidth - 10\tabcolsep) * \real{0.3023}}
  >{\raggedright\arraybackslash}p{(\columnwidth - 10\tabcolsep) * \real{0.1628}}@{}}
\toprule\noalign{}
\begin{minipage}[b]{\linewidth}\raggedright
\#
\end{minipage} & \begin{minipage}[b]{\linewidth}\raggedright
Signal
\end{minipage} & \begin{minipage}[b]{\linewidth}\raggedright
GPIO
\end{minipage} & \begin{minipage}[b]{\linewidth}\raggedright
Type
\end{minipage} & \begin{minipage}[b]{\linewidth}\raggedright
Connected To
\end{minipage} & \begin{minipage}[b]{\linewidth}\raggedright
Notes
\end{minipage} \\
\midrule\noalign{}
\endhead
\bottomrule\noalign{}
\endlastfoot
1 & CAN TX (low bus) & 5 & TWAI TX & SN65HVD230 TXD pin & Commands to SYS +
EPS-C \\
2 & SPI SCK & 36 & SPI & MCP2515 SCK pin & 10 MHz clock \\
3 & SPI MOSI & 37 & SPI & MCP2515 SI pin & High-level CAN TX data \\
4 & SPI CS & 39 & Digital & MCP2515 CS pin & Chip select, active low \\
5 & I2C SCL & 11 & I2C & IMU (optional) & Clock \\
6 & WDT Toggle & 21 & Digital & TPS3850 WDI pin & Toggled at 100 Hz by
control\_task \\
\end{longtable}

\begin{center}\rule{0.5\linewidth}{0.5pt}\end{center}

\section{3. SYS ESP32-S3 --- Input/Output
Tables}\label{sys-esp32-s3-inputoutput-tables}

\subsection{3.1 INPUTS (signals SYS reads from external
hardware)}\label{inputs-signals-sys-reads-from-external-hardware}

\begin{longtable}[]{@{}
  >{\raggedright\arraybackslash}p{(\columnwidth - 10\tabcolsep) * \real{0.0698}}
  >{\raggedright\arraybackslash}p{(\columnwidth - 10\tabcolsep) * \real{0.1860}}
  >{\raggedright\arraybackslash}p{(\columnwidth - 10\tabcolsep) * \real{0.1395}}
  >{\raggedright\arraybackslash}p{(\columnwidth - 10\tabcolsep) * \real{0.1395}}
  >{\raggedright\arraybackslash}p{(\columnwidth - 10\tabcolsep) * \real{0.3023}}
  >{\raggedright\arraybackslash}p{(\columnwidth - 10\tabcolsep) * \real{0.1628}}@{}}
\toprule\noalign{}
\begin{minipage}[b]{\linewidth}\raggedright
\#
\end{minipage} & \begin{minipage}[b]{\linewidth}\raggedright
Signal
\end{minipage} & \begin{minipage}[b]{\linewidth}\raggedright
GPIO
\end{minipage} & \begin{minipage}[b]{\linewidth}\raggedright
Type
\end{minipage} & \begin{minipage}[b]{\linewidth}\raggedright
Connected To
\end{minipage} & \begin{minipage}[b]{\linewidth}\raggedright
Notes
\end{minipage} \\
\midrule\noalign{}
\endhead
\bottomrule\noalign{}
\endlastfoot
1 & CAN RX (low bus) & 4 & TWAI RX & SN65HVD230 RXD & RT commands, MTR feedback
(0x206), SEB feedback \\
2 & ESTOP Button & 1 & Digital, NC & Red mushroom → GND & 10k pull-up to 3.3V.
LOW = ESTOP. Shared with MTR. \\
3 & Brake Lever & 2 & Digital, NO & Lever switch → GND & Internal pull-up. LOW =
pressed. → SEB via CAN 0x7B9. \\
4 & START Button & 32 & Digital, NO & Green momentary → GND & Internal pull-up.
LOW = pressed. \\
5 & MODE Button & 11 & Digital, NO & Momentary → GND & Internal pull-up. LOW =
pressed. → publishes CAN 0x110. \\
6 & Left Turn Switch & 3 & Digital, NO & Handlebar momentary → GND & Internal
pull-up \\
7 & Right Turn Switch & 6 & Digital, NO & Handlebar momentary → GND & Internal
pull-up \\
8 & Headlight Switch & 7 & Digital, NO & Handlebar toggle → GND & Internal
pull-up \\
\end{longtable}

\subsection{3.2 OUTPUTS (signals SYS drives to external
hardware)}\label{outputs-signals-sys-drives-to-external-hardware}

\begin{longtable}[]{@{}
  >{\raggedright\arraybackslash}p{(\columnwidth - 10\tabcolsep) * \real{0.0698}}
  >{\raggedright\arraybackslash}p{(\columnwidth - 10\tabcolsep) * \real{0.1860}}
  >{\raggedright\arraybackslash}p{(\columnwidth - 10\tabcolsep) * \real{0.1395}}
  >{\raggedright\arraybackslash}p{(\columnwidth - 10\tabcolsep) * \real{0.1395}}
  >{\raggedright\arraybackslash}p{(\columnwidth - 10\tabcolsep) * \real{0.3023}}
  >{\raggedright\arraybackslash}p{(\columnwidth - 10\tabcolsep) * \real{0.1628}}@{}}
\toprule\noalign{}
\begin{minipage}[b]{\linewidth}\raggedright
\#
\end{minipage} & \begin{minipage}[b]{\linewidth}\raggedright
Signal
\end{minipage} & \begin{minipage}[b]{\linewidth}\raggedright
GPIO
\end{minipage} & \begin{minipage}[b]{\linewidth}\raggedright
Type
\end{minipage} & \begin{minipage}[b]{\linewidth}\raggedright
Connected To
\end{minipage} & \begin{minipage}[b]{\linewidth}\raggedright
Notes
\end{minipage} \\
\midrule\noalign{}
\endhead
\bottomrule\noalign{}
\endlastfoot
1 & CAN TX (low bus) & 5 & TWAI TX & SN65HVD230 TXD & Telemetry to RT, mode cmd
(0x110), brake cmd (0x7B9) to SEB \\
2 & Left Turn Lamp & 18 & Digital & Relay coil → 12V & Relay COM=12V, NO=lamp \\
3 & Right Turn Lamp & 19 & Digital & Relay coil → 12V & Relay COM=12V,
NO=lamp \\
4 & Brake Light & 21 & Digital & Relay coil → 12V & Relay COM=12V, NO=lamp \\
5 & Headlight & 22 & Digital & Relay coil → 12V & Relay COM=12V, NO=lamp \\
6 & AUTO Bulb & 25 & Digital & Relay coil → 12V & Relay COM=12V, NO=bulb \\
7 & MANUAL Bulb & 26 & Digital & Relay coil → 12V & Relay COM=12V, NO=bulb \\
8 & 12V Power Relay & 27 & Digital & Relay coil → 12V & Relay COM=12V bus,
NO=accessories \\
9 & WDT Toggle & 23 & Digital & TPS3850 WDI pin & Toggled at 20 Hz by
safety\_task \\
\end{longtable}

\subsection{3.3 Safety --- Dashboard Buttons}\label{safety-dashboard-buttons}

\begin{longtable}[]{@{}
  >{\raggedright\arraybackslash}p{(\columnwidth - 6\tabcolsep) * \real{0.2857}}
  >{\raggedright\arraybackslash}p{(\columnwidth - 6\tabcolsep) * \real{0.2143}}
  >{\raggedright\arraybackslash}p{(\columnwidth - 6\tabcolsep) * \real{0.2143}}
  >{\raggedright\arraybackslash}p{(\columnwidth - 6\tabcolsep) * \real{0.2857}}@{}}
\toprule\noalign{}
\begin{minipage}[b]{\linewidth}\raggedright
Button
\end{minipage} & \begin{minipage}[b]{\linewidth}\raggedright
GPIO
\end{minipage} & \begin{minipage}[b]{\linewidth}\raggedright
Type
\end{minipage} & \begin{minipage}[b]{\linewidth}\raggedright
Wiring
\end{minipage} \\
\midrule\noalign{}
\endhead
\bottomrule\noalign{}
\endlastfoot
ESTOP & 1 & NC, active-low & Big red mushroom. GPIO → button → GND. 10k pull-up
to 3.3V. \\
START & 32 & NO, active-low & Green momentary. GPIO → button → GND. Internal
pull-up. \\
MODE & 11 & NO, active-low & Momentary. GPIO → button → GND. Internal
pull-up. \\
Brake Lever & 2 & Active-low & GPIO → lever switch → GND. Internal pull-up. \\
\end{longtable}

\subsection{3.4 Signal Lights}\label{signal-lights}

\begin{longtable}[]{@{}
  >{\raggedright\arraybackslash}p{(\columnwidth - 4\tabcolsep) * \real{0.3478}}
  >{\raggedright\arraybackslash}p{(\columnwidth - 4\tabcolsep) * \real{0.2609}}
  >{\raggedright\arraybackslash}p{(\columnwidth - 4\tabcolsep) * \real{0.3913}}@{}}
\toprule\noalign{}
\begin{minipage}[b]{\linewidth}\raggedright
Signal
\end{minipage} & \begin{minipage}[b]{\linewidth}\raggedright
GPIO
\end{minipage} & \begin{minipage}[b]{\linewidth}\raggedright
Circuit
\end{minipage} \\
\midrule\noalign{}
\endhead
\bottomrule\noalign{}
\endlastfoot
Left Turn Switch & 3 & Handlebar momentary, active-low. GPIO → switch → GND.
Pull-up. \\
Right Turn Switch & 6 & Handlebar momentary, active-low \\
Headlight Switch & 7 & Handlebar toggle, active-low \\
Left Turn Lamp & 18 & GPIO → relay coil → 12V. Relay COM → 12V, NO → lamp. \\
Right Turn Lamp & 19 & GPIO → relay → 12V → lamp \\
Brake Light & 21 & GPIO → relay → 12V → lamp \\
Headlight & 22 & GPIO → relay → 12V → lamp \\
\end{longtable}

\subsection{3.5 Indicators \& Power}\label{indicators-power}

\begin{longtable}[]{@{}
  >{\raggedright\arraybackslash}p{(\columnwidth - 4\tabcolsep) * \real{0.3478}}
  >{\raggedright\arraybackslash}p{(\columnwidth - 4\tabcolsep) * \real{0.2609}}
  >{\raggedright\arraybackslash}p{(\columnwidth - 4\tabcolsep) * \real{0.3913}}@{}}
\toprule\noalign{}
\begin{minipage}[b]{\linewidth}\raggedright
Signal
\end{minipage} & \begin{minipage}[b]{\linewidth}\raggedright
GPIO
\end{minipage} & \begin{minipage}[b]{\linewidth}\raggedright
Circuit
\end{minipage} \\
\midrule\noalign{}
\endhead
\bottomrule\noalign{}
\endlastfoot
AUTO Bulb & 25 & GPIO → relay → 12V → bulb \\
MANUAL Bulb & 26 & GPIO → relay → 12V → bulb \\
12V Relay & 27 & GPIO → relay coil → 12V. Relay COM → 12V bus, NO →
accessories. \\
WDT Toggle & 23 & GPIO → TPS3850 WDI pin. Toggled at 20 Hz by safety\_task. \\
\end{longtable}

\begin{center}\rule{0.5\linewidth}{0.5pt}\end{center}

\section{4. MTR STM32 --- Input/Output
Tables}\label{mtr-stm32-inputoutput-tables}

Motor controller (EGAS Level 1). Sole owner of all throttle/gear I/O. Mode-gated
via CAN \texttt{0x110} from SYS.

\subsection{4.1 INPUTS (signals MTR reads from external
hardware)}\label{inputs-signals-mtr-reads-from-external-hardware}

\begin{longtable}[]{@{}
  >{\raggedright\arraybackslash}p{(\columnwidth - 10\tabcolsep) * \real{0.0714}}
  >{\raggedright\arraybackslash}p{(\columnwidth - 10\tabcolsep) * \real{0.1905}}
  >{\raggedright\arraybackslash}p{(\columnwidth - 10\tabcolsep) * \real{0.1190}}
  >{\raggedright\arraybackslash}p{(\columnwidth - 10\tabcolsep) * \real{0.1429}}
  >{\raggedright\arraybackslash}p{(\columnwidth - 10\tabcolsep) * \real{0.3095}}
  >{\raggedright\arraybackslash}p{(\columnwidth - 10\tabcolsep) * \real{0.1667}}@{}}
\toprule\noalign{}
\begin{minipage}[b]{\linewidth}\raggedright
\#
\end{minipage} & \begin{minipage}[b]{\linewidth}\raggedright
Signal
\end{minipage} & \begin{minipage}[b]{\linewidth}\raggedright
Pin
\end{minipage} & \begin{minipage}[b]{\linewidth}\raggedright
Type
\end{minipage} & \begin{minipage}[b]{\linewidth}\raggedright
Connected To
\end{minipage} & \begin{minipage}[b]{\linewidth}\raggedright
Notes
\end{minipage} \\
\midrule\noalign{}
\endhead
\bottomrule\noalign{}
\endlastfoot
1 & CAN RX (low bus) & TBD & CAN RX & CAN transceiver RXD & RT commands (0x204),
SYS mode (0x110) \\
2 & ESTOP Button & TBD & Digital, NC & Red mushroom → GND & Shared with SYS
GPIO1. Hardwired Level 3 kill --- zero CAN delay. \\
3 & Throttle ADC & TBD & ADC & Voltage divider from 0--5V grip & 5V→3.3V via
resistive divider \\
4 & Gear D Sense & TBD & Digital & TLP281 ch1 output & 72V gear line →
optoisolator → GPIO \\
5 & Gear S Sense & TBD & Digital & TLP281 ch2 output & 72V gear line →
optoisolator → GPIO \\
6 & Gear R Sense & TBD & Digital & TLP281 ch3 output & 72V gear line →
optoisolator → GPIO \\
\end{longtable}

\subsection{4.2 OUTPUTS (signals MTR drives to external
hardware)}\label{outputs-signals-mtr-drives-to-external-hardware}

\begin{longtable}[]{@{}
  >{\raggedright\arraybackslash}p{(\columnwidth - 10\tabcolsep) * \real{0.0714}}
  >{\raggedright\arraybackslash}p{(\columnwidth - 10\tabcolsep) * \real{0.1905}}
  >{\raggedright\arraybackslash}p{(\columnwidth - 10\tabcolsep) * \real{0.1190}}
  >{\raggedright\arraybackslash}p{(\columnwidth - 10\tabcolsep) * \real{0.1429}}
  >{\raggedright\arraybackslash}p{(\columnwidth - 10\tabcolsep) * \real{0.3095}}
  >{\raggedright\arraybackslash}p{(\columnwidth - 10\tabcolsep) * \real{0.1667}}@{}}
\toprule\noalign{}
\begin{minipage}[b]{\linewidth}\raggedright
\#
\end{minipage} & \begin{minipage}[b]{\linewidth}\raggedright
Signal
\end{minipage} & \begin{minipage}[b]{\linewidth}\raggedright
Pin
\end{minipage} & \begin{minipage}[b]{\linewidth}\raggedright
Type
\end{minipage} & \begin{minipage}[b]{\linewidth}\raggedright
Connected To
\end{minipage} & \begin{minipage}[b]{\linewidth}\raggedright
Notes
\end{minipage} \\
\midrule\noalign{}
\endhead
\bottomrule\noalign{}
\endlastfoot
1 & CAN TX (low bus) & TBD & CAN TX & CAN transceiver TXD & Speed telemetry
(0x120), motor feedback (0x206) \\
2 & MCP4725 SDA & TBD & I2C & MCP4725 DAC (addr 0x60) & Throttle 0--5V output.
Pull-up to 3.3V. \\
3 & MCP4725 SCL & TBD & I2C & MCP4725 DAC & I2C clock \\
4 & Gear D Out & TBD & Digital & 4-ch relay module IN1 & 72V to ECU gear D
wire \\
5 & Gear S Out & TBD & Digital & 4-ch relay module IN2 & 72V to ECU gear S
wire \\
6 & Gear R Out & TBD & Digital & 4-ch relay module IN3 & 72V to ECU gear R
wire \\
\end{longtable}

\subsection{4.3 Throttle --- 0--5V (MTR-owned)}\label{throttle-05v-mtr-owned}

\begin{longtable}[]{@{}
  >{\raggedright\arraybackslash}p{(\columnwidth - 4\tabcolsep) * \real{0.3636}}
  >{\raggedright\arraybackslash}p{(\columnwidth - 4\tabcolsep) * \real{0.2273}}
  >{\raggedright\arraybackslash}p{(\columnwidth - 4\tabcolsep) * \real{0.4091}}@{}}
\toprule\noalign{}
\begin{minipage}[b]{\linewidth}\raggedright
Signal
\end{minipage} & \begin{minipage}[b]{\linewidth}\raggedright
Pin
\end{minipage} & \begin{minipage}[b]{\linewidth}\raggedright
Circuit
\end{minipage} \\
\midrule\noalign{}
\endhead
\bottomrule\noalign{}
\endlastfoot
Throttle Read & TBD (ADC) & 0--5V grip → \textbf{voltage divider} (2-resistor,
\textasciitilde5V→3.3V) → ADC pin \\
Throttle Output & I2C (TBD) & \textbf{MCP4725 DAC} (addr 0x60, VCC=5V). Output
0--5V → motor controller \\
\end{longtable}

\begin{verbatim}
Throttle Grip (0-5V) ──┬── R1 (e.g. 1.8kΩ) ──┬── ADC pin
                        │                      │
                        └── R2 (e.g. 3.3kΩ) ──┴── GND

MCP4725: VCC=5V, GND, SDA, SCL, VOUT→Motor Controller
\end{verbatim}

\subsection{4.4 Gear --- Bidirectional 72V
(MTR-owned)}\label{gear-bidirectional-72v-mtr-owned}

\textbf{Input} (read gear selector, galvanic isolation):

\begin{longtable}[]{@{}
  >{\raggedright\arraybackslash}p{(\columnwidth - 4\tabcolsep) * \real{0.3636}}
  >{\raggedright\arraybackslash}p{(\columnwidth - 4\tabcolsep) * \real{0.2273}}
  >{\raggedright\arraybackslash}p{(\columnwidth - 4\tabcolsep) * \real{0.4091}}@{}}
\toprule\noalign{}
\begin{minipage}[b]{\linewidth}\raggedright
Signal
\end{minipage} & \begin{minipage}[b]{\linewidth}\raggedright
Pin
\end{minipage} & \begin{minipage}[b]{\linewidth}\raggedright
Circuit
\end{minipage} \\
\midrule\noalign{}
\endhead
\bottomrule\noalign{}
\endlastfoot
Gear D Sense & TBD & 72V → \textbf{TLP281 optoisolator ch1} (input:
72V+resistor, output: GPIO + pull-up) \\
Gear S Sense & TBD & TLP281 optoisolator ch2 \\
Gear R Sense & TBD & TLP281 optoisolator ch3 \\
\end{longtable}

\begin{verbatim}
72V Gear D ──┬── R (current limit, ~33kΩ/2W) ── TLP281 pin1 (anode)
             └── TLP281 pin2 (cathode) ── GND_72V
TLP281 pin4 (collector) ── GPIO ── 10k pull-up ── 3.3V
TLP281 pin3 (emitter) ── GND
\end{verbatim}

\textbf{Output} (mimic 72V to ECU, mechanical relay isolation):

\begin{longtable}[]{@{}
  >{\raggedright\arraybackslash}p{(\columnwidth - 4\tabcolsep) * \real{0.3636}}
  >{\raggedright\arraybackslash}p{(\columnwidth - 4\tabcolsep) * \real{0.2273}}
  >{\raggedright\arraybackslash}p{(\columnwidth - 4\tabcolsep) * \real{0.4091}}@{}}
\toprule\noalign{}
\begin{minipage}[b]{\linewidth}\raggedright
Signal
\end{minipage} & \begin{minipage}[b]{\linewidth}\raggedright
Pin
\end{minipage} & \begin{minipage}[b]{\linewidth}\raggedright
Circuit
\end{minipage} \\
\midrule\noalign{}
\endhead
\bottomrule\noalign{}
\endlastfoot
Gear D Out & TBD & GPIO → 4-ch 5V relay module IN1 → \textbf{COM=72V} (via 1A
fuse) → NO → ECU Gear D \\
Gear S Out & TBD & GPIO → relay IN2 → COM=72V → NO → ECU Gear S \\
Gear R Out & TBD & GPIO → relay IN3 → COM=72V → NO → ECU Gear R \\
\end{longtable}

\textbf{Protection}:

\begin{verbatim}
72V Battery ──┬──[1A fast-blow fuse]──┬── Relay COM (D) ── NO ── ECU Gear D ──┬── [TVS SMCJ90CA] ── GND
              │                       ├── Relay COM (S) ── NO ── ECU Gear S ──┼── [TVS SMCJ90CA] ── GND
              │                       └── Relay COM (R) ── NO ── ECU Gear R ──┴── [TVS SMCJ90CA] ── GND
\end{verbatim}

\begin{longtable}[]{@{}lll@{}}
\toprule\noalign{}
Protection & Part & Spec \\
\midrule\noalign{}
\endhead
\bottomrule\noalign{}
\endlastfoot
Fuse & 1A fast-blow & 72V, 1A \\
TVS ×3 & SMCJ90CA bidirectional & 90--100V standoff, 1500W peak \\
\end{longtable}

\subsection{4.5 Mode-Gated Behavior}\label{mode-gated-behavior}

\begin{longtable}[]{@{}
  >{\raggedright\arraybackslash}p{(\columnwidth - 8\tabcolsep) * \real{0.1500}}
  >{\centering\arraybackslash}p{(\columnwidth - 8\tabcolsep) * \real{0.2750}}
  >{\raggedright\arraybackslash}p{(\columnwidth - 8\tabcolsep) * \real{0.2500}}
  >{\raggedright\arraybackslash}p{(\columnwidth - 8\tabcolsep) * \real{0.1500}}
  >{\raggedright\arraybackslash}p{(\columnwidth - 8\tabcolsep) * \real{0.1750}}@{}}
\toprule\noalign{}
\begin{minipage}[b]{\linewidth}\raggedright
Mode
\end{minipage} & \begin{minipage}[b]{\linewidth}\centering
CAN 0x110
\end{minipage} & \begin{minipage}[b]{\linewidth}\raggedright
Throttle
\end{minipage} & \begin{minipage}[b]{\linewidth}\raggedright
Gear
\end{minipage} & \begin{minipage}[b]{\linewidth}\raggedright
Notes
\end{minipage} \\
\midrule\noalign{}
\endhead
\bottomrule\noalign{}
\endlastfoot
Manual & 0 & ADC → DAC pass-through & Sense → relay pass-through & Rider direct
control \\
Auto & 1 & Follow CAN 0x204 \texttt{RT\_MotorSpeed} & Follow CAN 0x204
\texttt{RT\_Gear} & RT kinematics \\
ESTOP & 2 & DAC = 0 (cut throttle) & All relays off & Hardware + CAN kill \\
\end{longtable}

\begin{center}\rule{0.5\linewidth}{0.5pt}\end{center}

\section{5. External Watchdog}\label{external-watchdog}

Each ESP32 has a TPS3850 (or equivalent) window watchdog IC.

\begin{longtable}[]{@{}lll@{}}
\toprule\noalign{}
Pin & RT & SYS \\
\midrule\noalign{}
\endhead
\bottomrule\noalign{}
\endlastfoot
WDI (input) & GPIO21 & GPIO23 \\
MR (output) & ESP32 EN pin & ESP32 EN pin \\
Window & 100ms & 100ms \\
\end{longtable}

\textbf{Wiring}: GPIO → WDI. TPS3850 MR → ESP32 EN (active-low reset). If GPIO
stops toggling for \textgreater100ms, TPS3850 pulls MR LOW → ESP32 hardware
reset.

\begin{center}\rule{0.5\linewidth}{0.5pt}\end{center}

\section{6. Power Distribution}\label{power-distribution}

\begin{verbatim}
72V Traction Battery
  │
  ├── Motor Controller (72V) ── Motor (MCP4725 0–5V from MTR drives throttle input)
  │
  ├── DC-DC Converter (72V→12V, CAN: 0x012)
  │     │
  │     ├── 12V Relay (SYS GPIO27) ── 12V Accessory Bus
  │     │     ├── Signal lamps (via relays)
  │     │     ├── Mode bulbs (via relays)
  │     │     └── Headlight
  │     │
  │     ├── ESP32-S3 Dev Boards (5V regulated)
  │     └── STM32 Board (5V regulated)
  │           └── MCP4725 VCC (5V) — on MTR, provides 0–5V to motor controller
  │
  ├── SYNTREE EPS-C (steering) — internal power
  ├── SYNTREE SEB (brake) — internal power
  │
  └── 72V Gear Lines (via 1A fuse → relay COM terminals on MTR)

CAN_GND: Connect ground references if any CAN node uses isolated power.
\end{verbatim}

\begin{center}\rule{0.5\linewidth}{0.5pt}\end{center}

\section{7. Quick Reference --- GPIO
Summary}\label{quick-reference-gpio-summary}

\subsection{RT ESP32-S3}\label{rt-esp32-s3}

\begin{verbatim}
GPIO 1  : Encoder A (rear motor)
GPIO 2  : Encoder B (rear motor)
GPIO 3  : Encoder A (front wheel, TBD)
GPIO 4  : CAN RX (low-level)
GPIO 5  : CAN TX (low-level)
GPIO 6  : Encoder B (front wheel, TBD)
GPIO 7  : (unused)
GPIO 8  : (unused)
GPIO 9  : Encoder A (rear left, TBD)
GPIO 10 : I2C SDA (IMU optional)
GPIO 11 : I2C SCL (IMU optional)
GPIO 12 : Encoder B (rear left, TBD)
GPIO 13 : Encoder A (rear right, TBD)
GPIO 14 : Encoder B (rear right, TBD)
GPIO 21 : WDT toggle
GPIO 36 : SPI SCK (MCP2515)
GPIO 37 : SPI MOSI (MCP2515)
GPIO 38 : SPI MISO (MCP2515)
GPIO 39 : SPI CS (MCP2515)
GPIO 40 : MCP2515 INT
\end{verbatim}

\subsection{SYS ESP32-S3}\label{sys-esp32-s3}

\begin{verbatim}
GPIO 1  : ESTOP button (NC, active-low) — shared with MTR
GPIO 2  : Brake lever (active-low) → SEB via CAN 0x7B9
GPIO 3  : Left turn switch
GPIO 4  : CAN RX (low-level)
GPIO 5  : CAN TX (low-level)
GPIO 6  : Right turn switch
GPIO 7  : Headlight switch
GPIO 11 : MODE button → publishes CAN 0x110
GPIO 18 : Left turn lamp (relay)
GPIO 19 : Right turn lamp (relay)
GPIO 21 : Brake light (relay)
GPIO 22 : Headlight (relay)
GPIO 23 : WDT toggle (20 Hz)
GPIO 25 : AUTO bulb (relay)
GPIO 26 : MANUAL bulb (relay)
GPIO 27 : 12V power relay
GPIO 32 : START button
\end{verbatim}

\subsection{MTR STM32}\label{mtr-stm32}

\begin{verbatim}
ESTOP  : ESTOP button (NC, active-low) — shared with SYS GPIO1
ADC    : Throttle grip 0–5V (via voltage divider)
GPIO_x : Gear D sense (TLP281 ch1)
GPIO_x : Gear S sense (TLP281 ch2)
GPIO_x : Gear R sense (TLP281 ch3)
GPIO_x : Gear D output (relay ch1)
GPIO_x : Gear S output (relay ch2)
GPIO_x : Gear R output (relay ch3)
I2C    : MCP4725 DAC SDA (addr 0x60)
I2C    : MCP4725 DAC SCL
CAN_RX : CAN transceiver RXD
CAN_TX : CAN transceiver TXD
\end{verbatim}

\begin{center}\rule{0.5\linewidth}{0.5pt}\end{center}

\section{8. CAN Termination}\label{can-termination}

Each bus requires a 120Ω resistor between CAN\_H and CAN\_L at both physical
ends. Place one at each farthest node.

\begin{longtable}[]{@{}ll@{}}
\toprule\noalign{}
Bus & Termination points \\
\midrule\noalign{}
\endhead
\bottomrule\noalign{}
\endlastfoot
Low-Level & RT ESP32 + DC-DC converter (or farthest actuator) \\
High-Level & Jetson + RT ESP32 \\
\end{longtable}

Without proper termination, signal reflections cause bit errors and bus-off
states.


% ═══════════════════════════════════════════════════════════════════════
%  PART III -- THEORY & CONCEPTS (notes/)
% ═══════════════════════════════════════════════════════════════════════

\part{Theory \& Concepts}
\label{part:notes}

\textit{Background theory needed to understand the design decisions in Parts I--II. These are the notes --- learning material covering the fundamental concepts.}

\chapter{Analog Interfacing Basics}\label{analog-interfacing-basics}

Most embedded systems bridge digital logic (MCU) and analog reality (voltages,
currents, sensors, actuators). The E-Trike does this in several places: reading
a 0--5 V throttle grip via ADC, generating a 0--5 V motor speed signal via DAC,
and sensing 72 V gear lines through optoisolators.

This note covers the fundamental concepts behind each interface.

\begin{center}\rule{0.5\linewidth}{0.5pt}\end{center}

\section{1. ADC --- Analog-to-Digital
Conversion}\label{adc-analog-to-digital-conversion}

An \textbf{ADC} converts a continuous voltage into a discrete digital number.
The ESP32-S3 has a built-in 12-bit SAR ADC with multiple channels.

\subsection{Resolution}\label{resolution}

A 12-bit ADC divides the input range into 2¹² = 4096 steps. Each step
represents:

\begin{verbatim}
step_size = V_ref / 2^bits = 3.3V / 4096 ≈ 0.8 mV
\end{verbatim}

\begin{longtable}[]{@{}llll@{}}
\toprule\noalign{}
Bits & Steps & Step size at 3.3V & Step size at 5.0V \\
\midrule\noalign{}
\endhead
\bottomrule\noalign{}
\endlastfoot
8 & 256 & 12.9 mV & 19.5 mV \\
10 & 1024 & 3.2 mV & 4.9 mV \\
\textbf{12} & \textbf{4096} & \textbf{0.8 mV} & \textbf{1.2 mV} \\
16 & 65536 & 0.05 mV & 0.08 mV \\
\end{longtable}

For a throttle grip that moves from 0--5 V, 12 bits gives \textasciitilde1.2 mV
resolution --- far more than a human hand can position (and more than the grip
potentiometer's own noise floor).

\subsection{Voltage dividers}\label{voltage-dividers}

The ESP32 ADC is 3.3 V max. The throttle signal is 0--5 V. A \textbf{voltage
divider} scales it down:

\begin{verbatim}
5V input ──┬── R1 (e.g., 10kΩ) ──┬── ADC pin (0–3.3V)
           │                      │
           │                      ├── R2 (e.g., 20kΩ)
           │                      │
           └──────────────────────┴── GND

V_out = V_in × R2 / (R1 + R2) = 5V × 20k / (10k + 20k) = 3.33V
\end{verbatim}

Choose R1, R2 so that max input voltage → just under V\_ref. Keep total
resistance (R1 + R2) above 10 kΩ to limit current draw, but below 100 kΩ to keep
the ADC's sampling capacitor charging fast.

\subsection{Sampling and noise}\label{sampling-and-noise}

SAR ADCs work by charging an internal sampling capacitor, then comparing it
against a reference. If the source impedance is too high, the capacitor doesn't
fully charge before the comparison starts → reading is low.

\begin{itemize}
\tightlist
\item
  Add a small capacitor (100 nF) from ADC pin to GND --- acts as a charge
  reservoir.
\item
  Average multiple readings in software:
  \texttt{value\ =\ (value\ *\ 3\ +\ new\_reading)\ /\ 4} (exponential moving
  average).
\item
  A reading of 0 or 4095 that persists is likely a fault (shorted to GND or
  VCC).
\end{itemize}

\begin{center}\rule{0.5\linewidth}{0.5pt}\end{center}

\section{2. DAC --- Digital-to-Analog
Conversion}\label{dac-digital-to-analog-conversion}

A \textbf{DAC} does the reverse: converts a digital number into a continuous
voltage. The E-Trike uses the MCP4725, a 12-bit I2C DAC, to generate the 0--5 V
signal that controls motor speed.

\subsection{Why a dedicated DAC instead of
PWM?}\label{why-a-dedicated-dac-instead-of-pwm}

\begin{longtable}[]{@{}
  >{\raggedright\arraybackslash}p{(\columnwidth - 6\tabcolsep) * \real{0.2857}}
  >{\raggedright\arraybackslash}p{(\columnwidth - 6\tabcolsep) * \real{0.3714}}
  >{\raggedright\arraybackslash}p{(\columnwidth - 6\tabcolsep) * \real{0.1714}}
  >{\raggedright\arraybackslash}p{(\columnwidth - 6\tabcolsep) * \real{0.1714}}@{}}
\toprule\noalign{}
\begin{minipage}[b]{\linewidth}\raggedright
Approach
\end{minipage} & \begin{minipage}[b]{\linewidth}\raggedright
How it works
\end{minipage} & \begin{minipage}[b]{\linewidth}\raggedright
Pros
\end{minipage} & \begin{minipage}[b]{\linewidth}\raggedright
Cons
\end{minipage} \\
\midrule\noalign{}
\endhead
\bottomrule\noalign{}
\endlastfoot
\textbf{PWM + RC filter} & MCU outputs a square wave; low-pass filter smooths it
to an average voltage & Free (built-in peripheral) & Ripple, slow settling,
sensitive to component tolerance \\
\textbf{I2C/SPI DAC} & Dedicated IC outputs true analog voltage & Clean output,
\textless10 µs settling, no ripple & Additional BOM cost (\textasciitilde\$1) \\
\end{longtable}

For a safety-critical throttle signal, the DAC's clean output is worth the cost.
PWM ripple on a throttle line can cause the motor controller to oscillate
between drive and coast, producing surging.

\subsection{MCP4725 basics}\label{mcp4725-basics}

\begin{verbatim}
ESP32 I2C (SDA=15, SCL=16) ──► MCP4725 (addr 0x60) ──► 0–5V analog ──► Motor controller
                                   │
                                  VCC = 5V (not 3.3V — the DAC's output range equals its supply)
\end{verbatim}

12-bit value → output voltage:

\begin{verbatim}
V_out = (digital_value / 4095) × VCC
\end{verbatim}

So \texttt{digital\_value\ =\ 2047} → \textasciitilde2.5 V at VCC=5V.

\textbf{I2C write sequence:}

\begin{verbatim}
START | ADDR+W | ACK | DATA_H | ACK | DATA_L | ACK | STOP
\end{verbatim}

The MCP4725 expects two data bytes: the upper 8 bits of the 12-bit value in the
first byte, the lower 4 bits in the upper nibble of the second byte.

\begin{verbatim}
void dac_write(uint16_t value_12bit) {
    uint8_t data[2];
    data[0] = (value_12bit >> 4) & 0xFF;    // upper 8 bits
    data[1] = (value_12bit & 0x0F) << 4;     // lower 4 bits → upper nibble
    i2c_write(I2C_ADDR, data, 2);
}
\end{verbatim}

\begin{center}\rule{0.5\linewidth}{0.5pt}\end{center}

\section{3. Optoisolators --- Galvanic Isolation for Digital
Signals}\label{optoisolators-galvanic-isolation-for-digital-signals}

An \textbf{optoisolator} (optocoupler) transfers a signal using light, providing
electrical isolation between two voltage domains. No current flows between input
and output --- the signal crosses an optical gap.

The E-Trike uses TLP281 optoisolators to sense 72 V gear selector lines from a
3.3 V ESP32 GPIO.

\subsection{How it works}\label{how-it-works}

\begin{verbatim}
72V domain                             3.3V domain
──────────                             ──────────

72V gear wire ──► R_limit ──► LED ──► GND_72V
                                │
                           light │  (optical barrier, 2.5 kV isolation)
                                │
                                ▼
                          Phototransistor
                                │
ESP32 GPIO ◄── R_pullup ────────┤
                                │
                              GND_3V3
\end{verbatim}

\begin{itemize}
\tightlist
\item
  \textbf{72 V present} → LED emits light → phototransistor conducts → GPIO
  reads LOW.
\item
  \textbf{72 V absent} → LED off → phototransistor open → pull-up resistor pulls
  GPIO HIGH.
\end{itemize}

Note the \textbf{logic inversion}: HIGH on GPIO = ``no gear signal.'' This is
normal for optoisolator circuits --- handle the inversion in software.

\subsection{Key optoisolator parameters}\label{key-optoisolator-parameters}

\begin{longtable}[]{@{}
  >{\raggedright\arraybackslash}p{(\columnwidth - 4\tabcolsep) * \real{0.3056}}
  >{\raggedright\arraybackslash}p{(\columnwidth - 4\tabcolsep) * \real{0.2500}}
  >{\raggedright\arraybackslash}p{(\columnwidth - 4\tabcolsep) * \real{0.4444}}@{}}
\toprule\noalign{}
\begin{minipage}[b]{\linewidth}\raggedright
Parameter
\end{minipage} & \begin{minipage}[b]{\linewidth}\raggedright
Meaning
\end{minipage} & \begin{minipage}[b]{\linewidth}\raggedright
TLP281 typical
\end{minipage} \\
\midrule\noalign{}
\endhead
\bottomrule\noalign{}
\endlastfoot
Isolation voltage & Max voltage between input and output & 2500 Vrms \\
CTR (Current Transfer Ratio) & Output current / input current, as \% &
100--600\% @ 5 mA \\
Input forward voltage & Voltage drop across LED & \textasciitilde1.15 V @ 5
mA \\
Rise/fall time & Speed of signal transition & 2--5 µs (fast enough for 50 Hz
gear sensing) \\
\end{longtable}

\subsection{Sizing the input resistor}\label{sizing-the-input-resistor}

The resistor on the 72 V side limits LED current:

\begin{verbatim}
R = (V_in - V_f) / I_f = (72V - 1.15V) / 0.005A ≈ 14.2 kΩ
\end{verbatim}

Use a resistor with adequate power rating: P = I²R = (0.005)² × 14200 = 0.36 W →
use a 0.5 W or 1 W resistor.

\begin{center}\rule{0.5\linewidth}{0.5pt}\end{center}

\section{4. Relays --- Mechanical Isolation for Power
Signals}\label{relays-mechanical-isolation-for-power-signals}

When you need to switch a high-voltage signal ON/OFF from a low-voltage GPIO, a
\textbf{mechanical relay} provides both isolation and power handling.

\begin{verbatim}
ESP32 GPIO ──► R_base ──► NPN transistor base
                            │
                            ├── Collector ──► Relay coil ──► 12V
                            │
                            └── Emitter ──► GND

Relay contacts (electrically isolated from coil):
  COM ──► NO (normally open) ──► Load (e.g., 72V gear input on motor controller)
\end{verbatim}

\textbf{Why relays for gear outputs?}

\begin{longtable}[]{@{}
  >{\raggedright\arraybackslash}p{(\columnwidth - 4\tabcolsep) * \real{0.4000}}
  >{\raggedright\arraybackslash}p{(\columnwidth - 4\tabcolsep) * \real{0.2800}}
  >{\raggedright\arraybackslash}p{(\columnwidth - 4\tabcolsep) * \real{0.3200}}@{}}
\toprule\noalign{}
\begin{minipage}[b]{\linewidth}\raggedright
Property
\end{minipage} & \begin{minipage}[b]{\linewidth}\raggedright
Relay
\end{minipage} & \begin{minipage}[b]{\linewidth}\raggedright
MOSFET
\end{minipage} \\
\midrule\noalign{}
\endhead
\bottomrule\noalign{}
\endlastfoot
Galvanic isolation & ✓ (coil and contacts are separate) & ✗ (single silicon
die) \\
Fail-safe on power loss & ✓ (contacts open) & ✗ (body diode may conduct) \\
72 VDC switching & Simple & Requires gate drive with bootstrap \\
Switching speed & \textasciitilde5--10 ms & \textless1 µs \\
Audible feedback & Click confirms operation & Silent \\
\end{longtable}

For gear selection (switching once every few seconds at most), relay speed is
irrelevant --- and the audible click is a useful diagnostic.

\begin{center}\rule{0.5\linewidth}{0.5pt}\end{center}

\section{5. Protection Circuits}\label{protection-circuits}

Any line that connects to a high-voltage domain needs protection.

\subsection{Fuse (overcurrent)}\label{fuse-overcurrent}

A fast-blow fuse in series with every 72 V line. If the line shorts to chassis,
the fuse blows before the wire melts or the battery is damaged. Size the fuse
slightly above the maximum normal load. For gear signaling (microamps into
high-impedance input), 1 A is conservative --- any current above 1 A is
definitely a short.

\subsection{TVS diode (overvoltage
transient)}\label{tvs-diode-overvoltage-transient}

A \textbf{Transient Voltage Suppression (TVS)} diode clamps voltage spikes.
Place between the signal line and GND:

\begin{verbatim}
72V line ──┬── Fuse ──┬── Load
           │           │
           │           ├── TVS diode ── GND
           │           │
           └───────────┘
\end{verbatim}

Choose a TVS with: - \textbf{Standoff voltage \textgreater{} normal operating
voltage} (diode doesn't conduct normally). - \textbf{Clamping voltage
\textless{} destruction voltage of protected components}. -
\textbf{Bidirectional} for lines that can see both positive and negative
transients.

The E-Trike uses SMCJ90CA: 90 V standoff (above 72 V nominal), 146 V clamping,
1500 W peak. The `C' suffix = bidirectional; `A' suffix = 5\% tolerance.

\begin{center}\rule{0.5\linewidth}{0.5pt}\end{center}

\section{6. Common Pitfalls}\label{common-pitfalls}

\begin{longtable}[]{@{}
  >{\raggedright\arraybackslash}p{(\columnwidth - 4\tabcolsep) * \real{0.3333}}
  >{\raggedright\arraybackslash}p{(\columnwidth - 4\tabcolsep) * \real{0.4815}}
  >{\raggedright\arraybackslash}p{(\columnwidth - 4\tabcolsep) * \real{0.1852}}@{}}
\toprule\noalign{}
\begin{minipage}[b]{\linewidth}\raggedright
Pitfall
\end{minipage} & \begin{minipage}[b]{\linewidth}\raggedright
What happens
\end{minipage} & \begin{minipage}[b]{\linewidth}\raggedright
Fix
\end{minipage} \\
\midrule\noalign{}
\endhead
\bottomrule\noalign{}
\endlastfoot
\textbf{Missing voltage divider} & 5 V signal fed directly to 3.3 V ADC pin →
pin damage & Always check max input voltage against V\_ref \\
\textbf{High source impedance} & ADC reading droops --- sampling capacitor
doesn't charge & Add parallel capacitor (100 nF) at ADC pin, or use a buffer
op-amp \\
\textbf{DAC powered by MCU's 3.3V} & Output max is 3.3 V, but motor controller
expects 0--5 V & DAC VCC = 5 V rail, not 3.3 V \\
\textbf{No TVS on 72 V lines} & Load dump or inductive kick destroys
transceiver/opto & TVS + fuse on every high-voltage line \\
\textbf{Forgetting logic inversion with optos} & GPIO HIGH when signal is
present --- opposite of expected & Account for inversion in software; document
in pin table \\
\textbf{Relay coil without flyback diode} & Inductive kick on turn-off destroys
transistor & Diode across relay coil (cathode to VCC) \\
\end{longtable}

\begin{center}\rule{0.5\linewidth}{0.5pt}\end{center}

\emph{See also: {[}{[}high-voltage-isolation{]}{]} for the E-Trike's specific
72V isolation design, {[}{[}external-watchdog{]}{]} for hardware watchdog,
\texttt{architecture.md} §8.6 for throttle/gear control, §8.8 for pin
assignments.}


\chapter{CAN Gateway Design}\label{can-gateway-design}

A \textbf{CAN gateway} is a node connected to two (or more) CAN buses that
selectively forwards messages between them. It acts as a controlled bridge ---
not every frame crosses, and some are transformed in transit.

The E-Trike's RT ESP32-S3 is the gateway between the High-Level CAN (Jetson +
RT) and Low-Level CAN (RT + SYS + actuators). It uses two different CAN
controllers (built-in TWAI for low, external MCP2515 via SPI for high) and
applies explicit forwarding rules.

\begin{center}\rule{0.5\linewidth}{0.5pt}\end{center}

\section{1. Why Multiple CAN Buses?}\label{why-multiple-can-buses}

A single CAN bus works for simple systems. Beyond \textasciitilde50\% bus load
or when security isolation matters, you split into multiple buses.

\subsection{Reason 1: Bus load
partitioning}\label{reason-1-bus-load-partitioning}

Each CAN bus has finite bandwidth (500 kbit/s on the E-Trike). A standard frame
with 8-byte payload takes \textasciitilde130 bits, so max theoretical throughput
is \textasciitilde3,800 frames/s. At 50\% target load: \textasciitilde1,900
frames/s.

\begin{verbatim}
High bus load (if everything were on one bus):
  Jetson drive cmd:      100 Hz × 130 bits = 13,000 bps
  RT steer cmd:           50 Hz × 130 bits =  6,500 bps
  SYS throttle status:   100 Hz × 130 bits = 13,000 bps
  RT PID telemetry:       10 Hz × 130 bits =  1,300 bps
  EPS-C status:          100 Hz × 130 bits = 13,000 bps
  SEB status:            100 Hz × 130 bits = 13,000 bps
  ...
  Total: easily exceeds 50% load → low-priority messages delayed
\end{verbatim}

Splitting into two buses keeps each bus under 30\% load --- plenty of headroom
for bursts.

\subsection{Reason 2: Security isolation}\label{reason-2-security-isolation}

Jetson runs a full Linux + ROS 2 stack --- complex, network-connected, harder to
harden. The actuators (steering, brake) are on a physically separate bus. Even
if Jetson is fully compromised, it cannot directly command steering or brakes.
Every command passes through RT, which validates and clamps.

\begin{verbatim}
Jetson (untrusted) ──► High CAN ──► RT (validates) ──► Low CAN ──► Actuators
\end{verbatim}

\subsection{Reason 3: Fault containment}\label{reason-3-fault-containment}

A bus-off on one bus doesn't take down the other. If the High CAN fails (e.g.,
Jetson CAN controller hangs), RT continues sending steering commands and
heartbeats on the Low bus. SYS is unaffected and continues monitoring safety.

\begin{center}\rule{0.5\linewidth}{0.5pt}\end{center}

\section{2. The Three Forwarding
Categories}\label{the-three-forwarding-categories}

Every CAN message falls into exactly one category. This is an explicit design
decision, not emergent behavior.

\subsection{Category 1: Transparent
Forward}\label{category-1-transparent-forward}

The gateway copies the frame to the other bus \textbf{unchanged} --- same CAN
ID, same payload, same DLC. The receiver cannot tell whether the original sender
or the gateway transmitted it.

\textbf{When to use:} The message is relevant to nodes on both buses, and the
format doesn't need translation.

\begin{verbatim}
SYS sends 0x011 on low bus ──► RT copies to high bus ──► Jetson receives 0x011
                             (same ID, same payload)
\end{verbatim}

\textbf{Rule:} Transparent forward is safe only when each forwarded CAN ID has
exactly \textbf{one sender per bus}. If SYS sends \texttt{0x011} and RT also
generates \texttt{0x011}, forwarding creates ambiguity. The E-Trike enforces
``one sender per ID'' for all forwarded IDs.

\subsection{Category 2: Consume and
Regenerate}\label{category-2-consume-and-regenerate}

The gateway receives a frame on one bus, processes it internally, and transmits
a \textbf{different} CAN ID with \textbf{different} payload on the other bus.
This is a translation, not a forward.

\textbf{When to use:} The message semantics change across buses, or the gateway
adds validation/transformation.

\begin{verbatim}
Jetson sends 0x300 {speed, yaw} on high bus
    → RT consumes it
    → RT runs kinematics + PID
    → RT generates 0x204 {speed, gear} + 0x169 {angle} on low bus

The low-bus messages are completely different from the high-bus message that triggered them.
\end{verbatim}

\subsection{Category 3: Bus-Local}\label{category-3-bus-local}

The message never crosses the bus boundary. It serves only nodes on its own bus.

\textbf{When to use:} The message is only meaningful to nodes on one bus, or
crossing would create a security/safety problem.

Examples from the E-Trike: - \texttt{0x201} SES\_STATUS (EPS-C → RT): Only RT
needs steering feedback. Jetson gets processed telemetry via
\texttt{0x210\ RT\_STATE\_RPT}. - \texttt{0x7FD} RT\_HEARTBEAT (low bus):
Heartbeats are per-bus liveness domains --- bridging them creates the ambiguity
described in {[}{[}heartbeat-monitoring{]}{]}. - \texttt{0x012} SYS\_DCDC\_CMD
(SYS → DC-DC): The DC-DC converter is only on the low bus. Jetson has no reason
to see or control it directly.

\begin{center}\rule{0.5\linewidth}{0.5pt}\end{center}

\section{3. Gateway Implementation --- The Dispatch
Task}\label{gateway-implementation-the-dispatch-task}

On the E-Trike, the gateway logic runs in RT's \texttt{dispatch} task (priority
4):

\begin{verbatim}
dispatch_task loop:
  1. Block on can_rx_low_queue OR can_rx_high_queue
  2. Dequeue frame
  3. Switch on CAN ID:
       Low bus IDs:
         0x001 → forward to high + mode_set(Estop) locally
         0x011 → forward to high
         0x120 → forward to high
         0x600 → forward to high
         0x201 → steer_feedback (consume locally, never forward)
         0x7FD → update RT heartbeat liveness (consume locally)

       High bus IDs:
         0x001 → forward to low + mode_set(Estop) locally
         0x300 → enqueue to cmd_queue (consume → generate 0x204 + 0x169)
         0x301 → atomic store brake_request (consume locally)
         0x302 → forward to low
         0x7FC → update Jetson heartbeat liveness (consume locally)

      All other IDs → ignore (bus-local)
  4. Enqueue forwarded frames to gw_tx_low_queue or gw_tx_high_queue
\end{verbatim}

The actual CAN transmission happens in the \texttt{can\_tx\_low} and
\texttt{can\_tx\_high} tasks --- the dispatch task only queues, never blocks on
TX.

\begin{center}\rule{0.5\linewidth}{0.5pt}\end{center}

\section{4. ESTOP Bypasses the Queue}\label{estop-bypasses-the-queue}

The gateway's TX queues have finite depth. If a queue is full, normal frames are
dropped (the next periodic frame replaces them). But \texttt{0x001}
(SAFETY\_ESTOP) is the one frame that must never be dropped.

\begin{verbatim}
if (can_id == 0x001) {
    // Direct TX, bypass queue — guarantee delivery
    twai_transmit(&frame, 0);
    mcp2515_transmit(&frame, 0);
} else {
    // Normal path: enqueue for TX task
    xQueueSend(gw_tx_queue, &frame, 0);
}
\end{verbatim}

This is a general pattern: the most critical safety signal takes the fast path.
Everything else can tolerate queueing.

\begin{center}\rule{0.5\linewidth}{0.5pt}\end{center}

\section{5. Gateway as Security Boundary}\label{gateway-as-security-boundary}

The gateway is a \textbf{policy enforcement point}. It doesn't just pass bits
--- it validates and clamps:

\begin{longtable}[]{@{}
  >{\raggedright\arraybackslash}p{(\columnwidth - 4\tabcolsep) * \real{0.2258}}
  >{\raggedright\arraybackslash}p{(\columnwidth - 4\tabcolsep) * \real{0.2258}}
  >{\raggedright\arraybackslash}p{(\columnwidth - 4\tabcolsep) * \real{0.5484}}@{}}
\toprule\noalign{}
\begin{minipage}[b]{\linewidth}\raggedright
Check
\end{minipage} & \begin{minipage}[b]{\linewidth}\raggedright
Where
\end{minipage} & \begin{minipage}[b]{\linewidth}\raggedright
What it prevents
\end{minipage} \\
\midrule\noalign{}
\endhead
\bottomrule\noalign{}
\endlastfoot
Steering angle clamp & RT control task, before \texttt{0x169} TX & Jetson
commanding physically impossible or dangerous angles \\
Speed limit clamp & RT control task, before \texttt{0x204} TX & Jetson
commanding overspeed \\
Obstacle speed limit & RT control task & Running into detected obstacles \\
Command staleness & RT watchdog task & Jetson hung → zero setpoints \\
\end{longtable}

A transparent CAN-to-CAN bridge (like a dumb repeater) would forward anything.
The gateway pattern gives you a point of control to enforce safety invariants.

\begin{center}\rule{0.5\linewidth}{0.5pt}\end{center}

\section{6. Common Pitfalls}\label{common-pitfalls}

\begin{longtable}[]{@{}
  >{\raggedright\arraybackslash}p{(\columnwidth - 4\tabcolsep) * \real{0.3333}}
  >{\raggedright\arraybackslash}p{(\columnwidth - 4\tabcolsep) * \real{0.4815}}
  >{\raggedright\arraybackslash}p{(\columnwidth - 4\tabcolsep) * \real{0.1852}}@{}}
\toprule\noalign{}
\begin{minipage}[b]{\linewidth}\raggedright
Pitfall
\end{minipage} & \begin{minipage}[b]{\linewidth}\raggedright
What happens
\end{minipage} & \begin{minipage}[b]{\linewidth}\raggedright
Fix
\end{minipage} \\
\midrule\noalign{}
\endhead
\bottomrule\noalign{}
\endlastfoot
\textbf{Promiscuous bridge} & Forwarding ALL frames between buses → bus load
doubles, ID collisions, security bypass & Explicit allowlist for each
direction \\
\textbf{Bridging heartbeats} & Two different alive counters, same CAN ID on one
bus --- ambiguous liveness & Never forward heartbeat frames. Each bus has its
own. \\
\textbf{Queue depth too small} & Forwarded frames dropped under load & Profile
peak bus load; size queues for 2× worst case \\
\textbf{Gateway task priority too low} & Frames pile up in RX queue while
lower-prio tasks run & Dispatch task at priority 4 --- just below CAN RX, above
control \\
\textbf{Bridging with translation but keeping same ID} & Receivers see the same
ID from two sources --- which one is authoritative? & Each CAN ID has exactly
one sender per bus. If translation is needed, use a different ID. \\
\end{longtable}

\begin{center}\rule{0.5\linewidth}{0.5pt}\end{center}

\emph{See also: {[}{[}can-protocol{]}{]} for arbitration and bus basics,
{[}{[}heartbeat-monitoring{]}{]} for why heartbeats stay local,
\texttt{architecture.md} §2.3 for the E-Trike's forwarding rules, §7.3--7.4 for
gateway message tables.}


\chapter{CAN Hardware --- Physical Implementation
Basics}\label{can-hardware-physical-implementation-basics}

When moving from theory to physical implementation, CAN is unforgiving if
hardware fundamentals are ignored. These are the non-negotiable basics --- the
things that cause 90\% of failures in embedded systems.

\begin{center}\rule{0.5\linewidth}{0.5pt}\end{center}

\section{1. The 120Ω Termination Rule}\label{the-120ux3c9-termination-rule}

CAN relies on voltage differentials, making it susceptible to signal reflections
bouncing from the ends of the wires.

\begin{itemize}
\tightlist
\item
  \textbf{Exactly two} 120Ω resistors --- one at each extreme physical end of
  the main bus line, across CAN\_H and CAN\_L.
\item
  Do \textbf{not} put terminators on branch nodes (stubs).
\item
  Do \textbf{not} use more or fewer than two terminators.
\item
  With power off, measuring across CAN\_H to CAN\_L should read
  \textbf{\textasciitilde60Ω} (two 120Ω in parallel).
\item
  Tolerance: 120Ω ±5\% or ±10\% is fine; 118Ω works (under 2\% difference).
\end{itemize}

\textbf{Symptom of missing/wrong termination:} Intermittent comms, corrupted
frames, or total failure --- especially worse at higher bitrates.

\begin{verbatim}
Correct termination layout:

  120Ω ---+-----+-----+-----+--- 120Ω
          |     |     |     |
        Node  Node  Node  Node
        (short stub, no termination)
\end{verbatim}

\begin{center}\rule{0.5\linewidth}{0.5pt}\end{center}

\section{2. The ``Two-Wire'' Myth --- Common
Ground}\label{the-two-wire-myth-common-ground}

CAN is often advertised as a ``two-wire'' protocol. \textbf{This is misleading.}

\begin{itemize}
\tightlist
\item
  While signaling is differential, the transceivers need a common reference
  frame to keep voltages within their common-mode range.
\item
  If nodes are powered by separate, isolated power supplies, you \textbf{must}
  run a third wire connecting their grounds (CAN\_GND).
\item
  Failing to tie grounds together is a primary cause of transceivers burning out
  or dropping packets in industrial environments.
\item
  \textbf{Alternative:} Use galvanic isolation for the CAN bus.
\end{itemize}

\begin{center}\rule{0.5\linewidth}{0.5pt}\end{center}

\section{3. Controller vs.~Transceiver}\label{controller-vs.-transceiver}

Microcontrollers (ESP32, STM32, etc.) often advertise a ``CAN Interface'' ---
but this means only a CAN \textbf{Controller} is built into the silicon.

\begin{longtable}[]{@{}
  >{\raggedright\arraybackslash}p{(\columnwidth - 4\tabcolsep) * \real{0.2683}}
  >{\raggedright\arraybackslash}p{(\columnwidth - 4\tabcolsep) * \real{0.3171}}
  >{\raggedright\arraybackslash}p{(\columnwidth - 4\tabcolsep) * \real{0.4146}}@{}}
\toprule\noalign{}
\begin{minipage}[b]{\linewidth}\raggedright
Component
\end{minipage} & \begin{minipage}[b]{\linewidth}\raggedright
What it does
\end{minipage} & \begin{minipage}[b]{\linewidth}\raggedright
Built into MCU?
\end{minipage} \\
\midrule\noalign{}
\endhead
\bottomrule\noalign{}
\endlastfoot
\textbf{CAN Controller} & Handles logic, framing, arbitration, CRC & Sometimes
(e.g., ESP32 TWAI, STM32 bxCAN/FDCAN) \\
\textbf{CAN Transceiver} & Converts TX/RX logic levels to differential
CAN\_H/CAN\_L voltages & \textbf{Never} --- always a separate IC \\
\end{longtable}

You \textbf{always} need an external CAN transceiver IC (e.g., TJA1050,
SN65HVD230, MCP2551) between the microcontroller's TX/RX pins and the physical
bus wires.

\textbf{A known-good beginner bench setup:}

\begin{verbatim}
MCU/SoC A                  Twisted pair bus                     MCU/SoC B
+----------+        +---------------------------+        +----------+
| CAN ctrl |  TXD   |                           |   TXD | CAN ctrl |
|         o--------| Transceiver A      B      |-------o         |
|         o--------| RXD             RXD       |-------o         |
+----------+        |     CANH ======== CANH    |        +----------+
                    |     CANL ======== CANL    |
                    |       |               |    |
                    |      120Ω           120Ω  |
                    |       |               |    |
                    |      GND-----------GND    |
                    +---------------------------+
\end{verbatim}

The important beginner rule: \textbf{termination belongs at the two physical
ends of the main cable}, not at every node. CAN\_H, CAN\_L, and a ground
reference must all be connected: CAN\_H→CAN\_H, CAN\_L→CAN\_L, GND→GND, firm
wiring, and matching bitrate across all nodes.

\begin{verbatim}
MCU (ESP32/STM32)                    Bus
+------------------+       +--------+==== CAN_H
| CAN Controller   |--TXD--|        |
|                  |--RXD--|  CAN   |
+------------------+       | Trans- |
                           | ceiver |
                           |        |==== CAN_L
                           +--------+
\end{verbatim}

\textbf{Symptom of missing transceiver:} TXD toggles but CAN\_H/CAN\_L stay flat
--- no frames seen externally.

\begin{quote}
Also verify the transceiver is in \textbf{normal mode} (not standby/silent). On
Microchip MCP2561/2, a low STBY pin selects normal mode; floating it can put the
part into standby. Check your transceiver's datasheet.
\end{quote}

\begin{center}\rule{0.5\linewidth}{0.5pt}\end{center}

\section{4. Wire Twisting is Not Optional}\label{wire-twisting-is-not-optional}

\begin{itemize}
\tightlist
\item
  CAN\_H and CAN\_L \textbf{must} be tightly twisted together (twisted pair).
\item
  Twisting ensures external electromagnetic interference (EMI) strikes both
  wires equally. Because the receiver only looks at the \emph{difference}
  between the two voltages, the common-mode interference is mathematically
  canceled out.
\item
  Loose jumper wires might work on a clean lab bench, but on a mobile chassis
  near high-current motor drivers or actuators, they will fail.
\end{itemize}

\begin{center}\rule{0.5\linewidth}{0.5pt}\end{center}

\section{5. Bus Topology}\label{bus-topology}

\textbf{Wrong for bring-up at speed (star topology):}

\begin{verbatim}
          Node A
            |
Node B ---- Star ---- Node C
            |
          Node D
\end{verbatim}

This feels neat on a bench, but T-connections and star-like layouts increase
reflections. The problem gets worse as bitrate rises.

\textbf{Correct: Linear trunk with short stubs.}

\begin{verbatim}
120Ω ---+-----+-----+-----+--- 120Ω
        |     |     |     |
      Node  Node  Node  Node
      <30cm  <30cm  <30cm  <30cm
       stub   stub   stub   stub
\end{verbatim}

\textbf{Wrong topologies to avoid:}

\begin{verbatim}
❌ Star topology — reflections at the junction kill signal integrity
❌ Long stubs (>30cm at 1 Mbps) — impedance discontinuity causes reflections
❌ Ring topology — not supported by CAN
\end{verbatim}

\subsection{Bus Length vs.~Bitrate}\label{bus-length-vs.-bitrate}

\begin{longtable}[]{@{}lll@{}}
\toprule\noalign{}
Bitrate & Max Bus Length & Max Stub Length \\
\midrule\noalign{}
\endhead
\bottomrule\noalign{}
\endlastfoot
1 Mbit/s & 40 m & 0.3 m \\
500 kbit/s & 100 m & 0.5 m \\
250 kbit/s & 200 m & 1 m \\
125 kbit/s & 500 m & 2 m \\
50 kbit/s & 1000 m & 5 m \\
\end{longtable}

Higher bitrates demand shorter buses and shorter stubs.

\begin{center}\rule{0.5\linewidth}{0.5pt}\end{center}

\section{6. Managing Bus Load}\label{managing-bus-load}

Because CAN is broadcast-based, dumping data indiscriminately will cripple the
network.

\begin{itemize}
\tightlist
\item
  \textbf{Never} put CAN transmissions in an unconstrained loop.
\item
  Keep average bus load \textbf{below 50\%}. Above 70\%, lower-priority messages
  are systematically starved of bus access, leading to timeouts and control lag.
\item
  Optimize by reducing message frequency, payload size, or segmenting the
  network.
\item
  A CAN analyzer can display real-time bus load.
\end{itemize}

\begin{center}\rule{0.5\linewidth}{0.5pt}\end{center}

\section{7. Minimum Two Nodes}\label{minimum-two-nodes}

CAN requires at least \textbf{2 active nodes} on the bus. A frame's ACK slot
must be driven dominant by a receiver --- without an ACK, the transmitter enters
Error Passive mode, then Bus Off.

\begin{itemize}
\tightlist
\item
  Testing with a single board and a passive analyzer won't work (analyzers in
  listen-only mode don't ACK).
\item
  For solo development: use internal loopback mode, or connect a second active
  node.
\end{itemize}

\begin{center}\rule{0.5\linewidth}{0.5pt}\end{center}

\emph{See also: {[}{[}can-protocol{]}{]} for protocol theory,
{[}{[}can-troubleshooting{]}{]} for debugging procedures,
{[}{[}can-addressing-for-etrike{]}{]} for our project's ID assignments.}


\chapter{How CAN Works --- Protocol \&
Theory}\label{how-can-works-protocol-theory}

\textbf{Controller Area Network (CAN)} is a robust, multi-master, message-based
serial communication protocol. Originally developed by Bosch for automotive
multiplex wiring, it is now a foundational standard in robotics, industrial
automation, and embedded systems.

\begin{center}\rule{0.5\linewidth}{0.5pt}\end{center}

\section{Physical Layer}\label{physical-layer}

CAN uses a differential two-wire bus: \textbf{CAN High (CAN\_H)} and \textbf{CAN
Low (CAN\_L)}. Nodes do not have a master-slave relationship; any node can
broadcast when the bus is idle.

\begin{quote}
⚠️ \textbf{The ``two-wire'' label is a trap.} While signaling is differential,
the transceivers need a common voltage reference. If your nodes run on separate,
isolated power supplies, you \textbf{must} run a third wire connecting their
grounds (CAN\_GND). Without it, voltages drift beyond the common-mode range ---
transceivers burn out or drop packets. See
{[}{[}can-hardware-basics\#2-the-two-wire-myth-common-ground{]}{]}.
\end{quote}

\subsection{Logic States}\label{logic-states}

Instead of standard high/low voltages, CAN uses two states based on the voltage
\emph{difference} between CAN\_H and CAN\_L:

\begin{longtable}[]{@{}lllll@{}}
\toprule\noalign{}
State & Meaning & CAN\_H & CAN\_L & Differential \\
\midrule\noalign{}
\endhead
\bottomrule\noalign{}
\endlastfoot
\textbf{Dominant} & Logic 0 & \textasciitilde3.5V & \textasciitilde1.5V &
\textasciitilde2V \\
\textbf{Recessive} & Logic 1 & \textasciitilde2.5V & \textasciitilde2.5V &
\textasciitilde0V (idle) \\
\end{longtable}

\textbf{The Golden Rule:} A Dominant bit (0) physically overwrites a Recessive
bit (1) on the bus. This is the physical mechanism that makes collision
resolution possible.

\begin{center}\rule{0.5\linewidth}{0.5pt}\end{center}

\section{Addressing \& Priority}\label{addressing-priority}

CAN is \textbf{content-centric}, not node-centric. A standard CAN frame does not
contain sender or receiver addresses. Instead, it broadcasts a \textbf{Message
Identifier (ID)}.

\begin{itemize}
\tightlist
\item
  \textbf{Message Identifier:} The ID says \emph{what} the data is (e.g., motor
  RPM, steering angle, sensor fault). Every node receives every message, checks
  the ID against its hardware acceptance filters, and decides whether to process
  or ignore it.
\item
  \textbf{Priority:} The ID also determines message priority. \textbf{Lower ID =
  higher priority.} A message with ID \texttt{0x001} will always beat
  \texttt{0x050} in arbitration.
\end{itemize}

\begin{center}\rule{0.5\linewidth}{0.5pt}\end{center}

\section{Bitwise Arbitration (CSMA/CD+AMP)}\label{bitwise-arbitration-csmacdamp}

Because any node can transmit, collisions are inevitable. CAN handles them
without data loss using \textbf{Carrier Sense Multiple Access / Collision
Detection + Arbitration on Message Priority}.

When two or more nodes begin transmitting simultaneously, they undergo
\textbf{Non-Destructive Bitwise Arbitration}:

\begin{enumerate}
\def\labelenumi{\arabic{enumi}.}
\tightlist
\item
  Both nodes transmit their Identifier bits starting with the most significant
  bit (MSB), while simultaneously monitoring the bus.
\item
  If Node A transmits a Recessive bit (1) but Node B transmits a Dominant bit
  (0), Node B's dominant bit pulls the bus to the Dominant state.
\item
  Node A reads the bus, sees a 0 despite transmitting a 1, and
  \textbf{immediately realizes it has lost arbitration}.
\item
  Node A instantly stops transmitting and switches to receiving mode. Node B
  continues its payload \textbf{uninterrupted} --- no data is lost.
\end{enumerate}

\begin{quote}
\textbf{Key insight:} The lower the ID number, the more leading zeros (dominant
bits), so lower IDs always win arbitration. This is why safety-critical messages
get the lowest IDs.
\end{quote}

\begin{center}\rule{0.5\linewidth}{0.5pt}\end{center}

\section{Frame Types}\label{frame-types}

There are four primary frame types:

\subsection{1. Data Frame}\label{data-frame}

The standard frame for transmitting payload data from a transmitter to
receivers.

\begin{longtable}[]{@{}ll@{}}
\toprule\noalign{}
Field & Purpose \\
\midrule\noalign{}
\endhead
\bottomrule\noalign{}
\endlastfoot
SOF & Start of Frame (single dominant bit) \\
ID & 11-bit or 29-bit identifier; lower = higher priority \\
RTR & Remote Transmission Request (dominant for data frames) \\
Control & IDE bit, reserved bit, DLC (Data Length Code: 0--8 bytes) \\
Data & Payload (0--8 bytes Classical, up to 64 bytes FD, up to 2048 bytes XL) \\
CRC & Cyclic Redundancy Check for error detection \\
ACK & Receiver acknowledges by overwriting the ACK slot with a dominant bit \\
EOF & End of Frame (7 recessive bits) \\
\end{longtable}

\subsection{2. Remote Frame}\label{remote-frame}

A node requests specific data by sending a frame with the RTR bit set to
recessive. The node holding that data responds with a Data Frame.
\emph{Deprecated in CAN FD and CAN XL.}

\subsection{3. Error Frame}\label{error-frame}

If any node detects a protocol or CRC error, it immediately broadcasts an Error
Frame (a deliberate sequence of 6 dominant bits). This violates the bit-stuffing
rules, causing all other nodes to discard the corrupted message and forcing
retransmission.

\subsection{4. Overload Frame}\label{overload-frame}

Used by a receiving node to request a delay between consecutive frames if
overwhelmed. Rarely used in modern fast microcontrollers.

\begin{center}\rule{0.5\linewidth}{0.5pt}\end{center}

\section{CAN Standards Evolution}\label{can-standards-evolution}

\subsection{Classical CAN (CAN 2.0, ISO
11898-1:2015)}\label{classical-can-can-2.0-iso-11898-12015}

\begin{longtable}[]{@{}lllll@{}}
\toprule\noalign{}
Variant & ID Length & Unique IDs & Payload & Max Baud \\
\midrule\noalign{}
\endhead
\bottomrule\noalign{}
\endlastfoot
CAN 2.0A (Standard) & 11-bit & 2,048 & 8 bytes & 1 Mbit/s \\
CAN 2.0B (Extended) & 29-bit & 536+ million & 8 bytes & 1 Mbit/s \\
\end{longtable}

\subsection{CAN FD (Flexible Data-Rate)}\label{can-fd-flexible-data-rate}

Introduced to solve the bandwidth bottleneck of Classical CAN.

\begin{itemize}
\tightlist
\item
  \textbf{Dual bitrates:} Standard slower speed (e.g., 500 kbit/s) during
  arbitration, switching to a faster speed (2--8 Mbit/s) during the data payload
  phase.
\item
  \textbf{Payload:} Up to \textbf{64 bytes}.
\item
  \textbf{Compatibility:} Fully backward compatible (an FD node understands
  Classical CAN; a Classical node will error on FD frames).
\end{itemize}

\subsection{CAN XL (ISO 11898-1:2024)}\label{can-xl-iso-11898-12024}

The third generation, bridging CAN FD and 100BASE-T1 Automotive Ethernet.

\begin{itemize}
\tightlist
\item
  \textbf{Data phase:} Up to \textbf{20 Mbit/s} (using push-pull signaling via
  CAN SIC XL transceivers).
\item
  \textbf{Payload:} Up to \textbf{2048 bytes} --- can encapsulate entire
  Ethernet (TCP/IP) frames.
\item
  \textbf{New features:}

  \begin{itemize}
  \tightlist
  \item
    \textbf{VCID (Virtual CAN Network ID):} Splits one physical bus into
    multiple logical networks (like VLANs).
  \item
    \textbf{SDT (Service Data Unit Type):} Upper-layer indicator describing the
    payload type.
  \item
    \textbf{AF (Acceptance Field):} 32-bit field for content- or node-based
    addressing, separate from the 11-bit priority ID.
  \end{itemize}
\end{itemize}

\begin{center}\rule{0.5\linewidth}{0.5pt}\end{center}

\section{Higher-Layer Protocols}\label{higher-layer-protocols}

Bare-metal CAN only provides OSI Layers 1 (Physical) and 2 (Data Link).
Higher-layer protocols add addressing, multi-packet transmission, and network
management:

\begin{itemize}
\tightlist
\item
  \textbf{SAE J1939} --- Standard for heavy-duty vehicles. Forces node-based
  addressing onto CAN 2.0B's 29-bit ID by partitioning it into a Parameter Group
  Number (PGN, what the data is) and a Source Address (who sent it).
\item
  \textbf{CANopen} --- Popular in industrial automation, robotics, and motor
  controllers. Uses an Object Dictionary model with Process Data Objects (PDOs,
  real-time) and Service Data Objects (SDOs, configuration).
\item
  \textbf{DeviceNet / UAVCAN (Cyphal)} --- Specialized frameworks for industrial
  and aerospace applications.
\end{itemize}

\begin{center}\rule{0.5\linewidth}{0.5pt}\end{center}

\emph{See also: {[}{[}can-hardware-basics{]}{]} for physical implementation,
{[}{[}can-troubleshooting{]}{]} for debugging,
{[}{[}can-addressing-for-etrike{]}{]} for our project's ID scheme.}


\chapter{CAN Troubleshooting --- Common Mistakes \&
Debugging}\label{can-troubleshooting-common-mistakes-debugging}

Most beginner CAN failures are not caused by mysterious protocol bugs --- they
come from a small set of repeat offenders. This note covers how to identify and
fix them, and the diagnostic tools to use.

\begin{quote}
⚠️ \textbf{Safety warning:} Don't connect an ESP32 (or any dev board) directly
to a road vehicle's CAN bus unless you're absolutely certain what you're doing
--- you could disable critical systems and create safety risks. Use an OBD2
adapter or test on a bench first.
\end{quote}

\begin{center}\rule{0.5\linewidth}{0.5pt}\end{center}

\section{Quick Debugging Checklist}\label{quick-debugging-checklist}

Before posting a forum question or concluding your controller is defective, work
through these:

\begin{enumerate}
\def\labelenumi{\arabic{enumi}.}
\tightlist
\item
  \textbf{\textasciitilde60Ω across CAN\_H/CAN\_L?} (power off) --- verifies two
  120Ω terminators at bus ends
\item
  \textbf{Voltages sane?} CAN\_H ≈ 2.5--3.5V, CAN\_L ≈ 1.5--2.5V (powered, idle
  bus)
\item
  \textbf{Real transceiver on every node?} --- not just a CAN-capable MCU pin
\item
  \textbf{Transceiver in normal mode?} --- not standby or silent
\item
  \textbf{≥2 active nodes on the bus?} --- unless in self-test/loopback
\item
  \textbf{CAN\_H→CAN\_H, CAN\_L→CAN\_L, shared GND?} --- polarity and ground
  reference
\item
  \textbf{Linear trunk topology?} --- no star, short stubs
\item
  \textbf{All nodes at same baud rate?} --- and same sample point configuration
\item
  \textbf{Classical CAN working first?} --- before enabling CAN FD
\item
  \textbf{Error counters checked?} --- read TEC/REC before retrying forever
\item
  \textbf{ISR/foreground buffer ownership clear?} --- no shared-buffer races
\end{enumerate}

\begin{center}\rule{0.5\linewidth}{0.5pt}\end{center}

\section{Common Mistakes --- Symptoms \&
Fixes}\label{common-mistakes-symptoms-fixes}

\subsection{Physical Layer}\label{physical-layer}

\begin{longtable}[]{@{}
  >{\raggedright\arraybackslash}p{(\columnwidth - 4\tabcolsep) * \real{0.3750}}
  >{\raggedright\arraybackslash}p{(\columnwidth - 4\tabcolsep) * \real{0.4167}}
  >{\raggedright\arraybackslash}p{(\columnwidth - 4\tabcolsep) * \real{0.2083}}@{}}
\toprule\noalign{}
\begin{minipage}[b]{\linewidth}\raggedright
Mistake
\end{minipage} & \begin{minipage}[b]{\linewidth}\raggedright
Symptoms
\end{minipage} & \begin{minipage}[b]{\linewidth}\raggedright
Fix
\end{minipage} \\
\midrule\noalign{}
\endhead
\bottomrule\noalign{}
\endlastfoot
\textbf{No transceiver} (MCU has controller only) & TXD toggles, CAN\_H/CAN\_L
stay flat & Add ISO 11898-compatible transceiver to each node \\
\textbf{Transceiver in standby/silent} & Internal loopback works, bus silent &
Drive STBY/S pin to normal mode \\
\textbf{Only one active node} & Tx FIFO fills, endless retransmissions, ACK
errors & Use ≥2 active nodes, or internal loopback for self-test \\
\textbf{Missing/wrong termination} & Intermittent comms, worse at higher
bitrates & Exactly two 120Ω at bus ends → measure \textasciitilde60Ω \\
\textbf{Star topology / long stubs} & Works on bench, fails at speed or with
more nodes & Linear trunk, stubs \textless30cm at 1 Mbps \\
\textbf{CAN\_H/CAN\_L swapped} & No comms, decode errors & Verify CAN\_H→CAN\_H,
CAN\_L→CAN\_L \\
\textbf{No common ground} & Erratic errors, transceiver burnout & Run dedicated
signal ground, or use galvanic isolation \\
\end{longtable}

\begin{quote}
⚠️ \textbf{The ``two-wire'' label is a trap.} CAN is often advertised as a
two-wire protocol (CAN\_H and CAN\_L). This is misleading. While signaling is
differential, the transceivers need a common reference frame --- otherwise
voltages drift beyond the common-mode range and transceivers burn out or drop
packets. If your nodes run on separate, isolated power supplies, you
\textbf{must} run a third wire connecting their grounds (CAN\_GND). Failing to
tie grounds together is a primary cause of inexplicable failures in real-world
setups. See {[}{[}can-hardware-basics\#2-the-two-wire-myth-common-ground{]}{]}.
\end{quote}

\subsection{Configuration \& Timing}\label{configuration-timing}

\begin{longtable}[]{@{}
  >{\raggedright\arraybackslash}p{(\columnwidth - 4\tabcolsep) * \real{0.3750}}
  >{\raggedright\arraybackslash}p{(\columnwidth - 4\tabcolsep) * \real{0.4167}}
  >{\raggedright\arraybackslash}p{(\columnwidth - 4\tabcolsep) * \real{0.2083}}@{}}
\toprule\noalign{}
\begin{minipage}[b]{\linewidth}\raggedright
Mistake
\end{minipage} & \begin{minipage}[b]{\linewidth}\raggedright
Symptoms
\end{minipage} & \begin{minipage}[b]{\linewidth}\raggedright
Fix
\end{minipage} \\
\midrule\noalign{}
\endhead
\bottomrule\noalign{}
\endlastfoot
\textbf{Mismatched baud rates} & Error storms, no valid frames & All nodes use
identical baud rate --- compute from oscillator, don't guess \\
\textbf{Wrong sample point} & Marginal comms, occasional errors & Set sample
point to \textasciitilde75--87.5\% of bit time (e.g., 14 of 16 time quanta) \\
\textbf{RC oscillator at ≥250 kbps} & Unreliable timing & Use quartz/crystal
oscillator --- RC accuracy may be insufficient \\
\textbf{CAN FD without data-phase config} & Classical works, FD fails &
Configure nominal bitrate, data bitrate, and TDC/SSP explicitly \\
\textbf{Classical node on FD bus} & Error frames, bus-off & Keep all nodes
Classical during bring-up; move to FD together \\
\textbf{Standard/Extended ID mismatch} & Sender ``works,'' receiver never
matches & Consistent 11-bit vs 29-bit choice across controller, filters,
analyzer \\
\end{longtable}

\subsection{Software}\label{software}

\begin{longtable}[]{@{}
  >{\raggedright\arraybackslash}p{(\columnwidth - 4\tabcolsep) * \real{0.3750}}
  >{\raggedright\arraybackslash}p{(\columnwidth - 4\tabcolsep) * \real{0.4167}}
  >{\raggedright\arraybackslash}p{(\columnwidth - 4\tabcolsep) * \real{0.2083}}@{}}
\toprule\noalign{}
\begin{minipage}[b]{\linewidth}\raggedright
Mistake
\end{minipage} & \begin{minipage}[b]{\linewidth}\raggedright
Symptoms
\end{minipage} & \begin{minipage}[b]{\linewidth}\raggedright
Fix
\end{minipage} \\
\midrule\noalign{}
\endhead
\bottomrule\noalign{}
\endlastfoot
\textbf{Treating CAN ID as node address} & Wrong priorities, awkward filter
design & Design IDs around message meaning \& priority; addressing is
application-layer \\
\textbf{Wrong filter mask polarity} & Bus active but app sees no frames & Start
accept-all, then narrow. Rule: mask bit 1 = compare, 0 = ignore \\
\textbf{Unsafe ISR/foreground buffer sharing} & Corrupted payload, zeros
appearing, lost messages & ISR writes, foreground consumes --- separate
ownership via queue or double buffer \\
\textbf{Duplicate listeners / missing shutdown} & Duplicate callbacks, memory
growth & \texttt{Notifier.stop()}, avoid adding same listener twice, use
\texttt{ThreadSafeBus} in Python \\
\end{longtable}

\begin{center}\rule{0.5\linewidth}{0.5pt}\end{center}

\section{Step-by-Step Fix Procedures}\label{step-by-step-fix-procedures}

\subsection{Physical layer first --- before touching
software}\label{physical-layer-first-before-touching-software}

\begin{enumerate}
\def\labelenumi{\arabic{enumi}.}
\tightlist
\item
  Power the bus down. Measure CAN\_H to CAN\_L resistance → should be
  \textbf{\textasciitilde60Ω}. If not, fix termination and topology before
  anything else.
\item
  Verify every node has a real transceiver and it's in \textbf{normal mode} (not
  standby/silent).
\item
  Verify CAN\_H→CAN\_H, CAN\_L→CAN\_L, and GND are connected correctly.
\item
  Reduce to \textbf{two nodes only} on a short linear cable. This single
  procedure fixes ``no waveform'', ``bus silent'', ``works in loopback only'',
  and many ``random'' overrun reports.
\end{enumerate}

\subsection{ACK errors / ``nothing is
transmitted''}\label{ack-errors-nothing-is-transmitted}

Real CAN requires an ACK from another node. If you have only one transmitting
node plus a passive analyzer, the controller keeps retrying forever and the Tx
FIFO appears to ``stick''.

\begin{itemize}
\tightlist
\item
  Use internal loopback to confirm the controller block works.
\item
  Connect a second \textbf{active} ACK-capable node, then retest in normal mode.
\item
  Listen-only mode won't ACK --- it can monitor but can't serve as the only
  partner in a two-device test.
\end{itemize}

\subsection{Baud rate \& bit timing}\label{baud-rate-bit-timing}

\begin{itemize}
\tightlist
\item
  Pick one conservative Classical CAN bitrate. Make every node use identical
  nominal timing.
\item
  Compute timing from oscillator and controller clock --- don't tune by
  folklore. Apply identically to every node.
\item
  Rule of thumb: \textbf{12--20 time quanta} with sample point at
  \textbf{75--87.5\%} of bit time.
\end{itemize}

\subsection{CAN FD specific}\label{can-fd-specific}

\begin{itemize}
\tightlist
\item
  Make Classical CAN work first. Enable FD only after specifying \textbf{both}
  nominal bitrate and data-phase bitrate.
\item
  If FD works only when slowed down or with BRS disabled: suspect wrong data
  timing or missing TDC (transceiver delay compensation) before suspecting
  payload format bugs.
\end{itemize}

\subsection{Classical vs FD mismatch}\label{classical-vs-fd-mismatch}

\begin{itemize}
\tightlist
\item
  Put every device on Classical CAN first. Once stable, move every live device
  to \textbf{ISO CAN FD} together. Only then tune data-phase bitrate.
\item
  A Classical node on an FD bus will send error frames --- it doesn't understand
  FD frames. Use a gateway if legacy Classical nodes must remain.
\end{itemize}

\subsection{Error frames / Error Passive / Bus
Off}\label{error-frames-error-passive-bus-off}

\begin{itemize}
\tightlist
\item
  An error frame is the \emph{mechanism} by which CAN deletes a bad frame ---
  not the ``bug'' itself.
\item
  Fix the physical or timing cause, inspect counters (TEC/REC) and error status,
  then run the controller's documented recovery sequence instead of blindly
  hammering retransmissions.
\end{itemize}

\subsection{Message ID, masks \& filters}\label{message-id-masks-filters}

\begin{itemize}
\tightlist
\item
  Separate protocol facts from project conventions. CAN ID = arbitration
  priority, not automatically a node address.
\item
  Start with \textbf{accept-all} or passive monitoring. After traffic is proven
  good, add filters.
\item
  Classic bit-mask rule: \textbf{mask bit 1 = compare this bit, mask bit 0 =
  ignore this bit.} All ones = exact match, all zeros = everything matches.
\item
  On MCAN-style controllers, filters are checked in order; evaluation stops at
  the \textbf{first} match --- ordering matters.
\end{itemize}

\subsection{Software: races, listeners,
overflows}\label{software-races-listeners-overflows}

\begin{itemize}
\tightlist
\item
  \textbf{Linux SocketCAN:} Uses local loopback semantics --- host-local
  visibility ≠ proof the bus is healthy. Supports \texttt{vcan} for testing
  without hardware.
\item
  \textbf{python-can:} Filtering is often in hardware/kernel space, not Python
  space. Avoid adding the same listener twice. Use \texttt{ThreadSafeBus} for
  multi-threaded access. \texttt{Notifier.stop()} matters --- listeners may
  flush state on shutdown.
\item
  \textbf{Bare-metal/RTOS:} Never let ISR and foreground code write the same
  frame buffer without ownership rules. ``CAN data became zero'' is usually a
  race condition.
\end{itemize}

\begin{center}\rule{0.5\linewidth}{0.5pt}\end{center}

\section{Understanding Error States}\label{understanding-error-states}

CAN controllers maintain two counters: \textbf{TEC} (Transmit Error Counter) and
\textbf{REC} (Receive Error Counter).

\begin{longtable}[]{@{}
  >{\raggedright\arraybackslash}p{(\columnwidth - 4\tabcolsep) * \real{0.2500}}
  >{\raggedright\arraybackslash}p{(\columnwidth - 4\tabcolsep) * \real{0.3929}}
  >{\raggedright\arraybackslash}p{(\columnwidth - 4\tabcolsep) * \real{0.3571}}@{}}
\toprule\noalign{}
\begin{minipage}[b]{\linewidth}\raggedright
State
\end{minipage} & \begin{minipage}[b]{\linewidth}\raggedright
Condition
\end{minipage} & \begin{minipage}[b]{\linewidth}\raggedright
Behavior
\end{minipage} \\
\midrule\noalign{}
\endhead
\bottomrule\noalign{}
\endlastfoot
\textbf{Error Active} & TEC ≤ 127 and REC ≤ 127 & Normal operation; can send
Active Error Flags \\
\textbf{Error Passive} & TEC \textgreater{} 127 or REC \textgreater{} 127 & Can
still communicate but sends Passive Error Flags; waits longer between
retransmissions \\
\textbf{Bus Off} & TEC \textgreater{} 255 & Node disconnects from bus; must be
manually recovered \\
\end{longtable}

\textbf{Key insight:} An Error Frame is the \emph{mechanism} by which CAN
deletes a bad frame --- it's not the ``bug'' itself. Fix the root cause
(physical, timing, or configuration), then recover the controller per its
datasheet.

\begin{center}\rule{0.5\linewidth}{0.5pt}\end{center}

\section{Diagnostic Tools \& Workflow}\label{diagnostic-tools-workflow}

Use tools in this order: \textbf{multimeter → oscilloscope → CAN analyzer →
logic analyzer → load tools.}

\subsection{1. Multimeter (always first)}\label{multimeter-always-first}

\begin{itemize}
\tightlist
\item
  Power off, measure CAN\_H to CAN\_L resistance → should be
  \textbf{\textasciitilde60Ω}
\item
  Check continuity: CAN\_H, CAN\_L, GND on each connector
\end{itemize}

\subsection{2. Oscilloscope}\label{oscilloscope}

Ask four questions in order: 1. Does TXD toggle when software sends a frame? 2.
Does RXD follow expected behavior? 3. Do CAN\_H and CAN\_L move differentially?
4. Do edges and idle levels look stable, or do they ring/distort?

\textbf{Key insight:} If TXD toggles but CAN\_H/CAN\_L stay flat → missing
transceiver, transceiver in standby, missing power, or wrong pin mux.

\subsection{3. CAN Analyzer (e.g., PCAN-View,
Kvaser)}\label{can-analyzer-e.g.-pcan-view-kvaser}

\begin{itemize}
\tightlist
\item
  Start in \textbf{listen-only} mode for passive observation --- don't disturb
  an existing bus
\item
  Configure the exact nominal bitrate before anything else
\item
  Only after passive observation is sound should you use it as an active node
\item
  Monitor bus load and error counters
\end{itemize}

\begin{quote}
\textbf{Common trap:} USB-to-CAN connected ≠ another node ACKing you.
Listen-only analyzers do not ACK.
\end{quote}

\subsection{4. Logic Analyzer}\label{logic-analyzer}

\begin{itemize}
\tightlist
\item
  Use only for the \textbf{controller/transceiver boundary} --- probe TXD/RXD
\item
  Do NOT use as a substitute for a proper differential CAN instrument
\item
  If probing CAN\_H/CAN\_L directly, ensure the analyzer's electrical limits are
  appropriate
\end{itemize}

\subsection{5. Bus Load \& Replay Tools}\label{bus-load-replay-tools}

\begin{itemize}
\tightlist
\item
  Linux: \texttt{cangen}, \texttt{candump}, \texttt{canplayer},
  \texttt{cansequence} from \texttt{can-utils}
\item
  Increase load gradually; watch for lost frames, overruns, sequencing issues
\item
  \texttt{vcan} (virtual CAN) is excellent for testing logic without hardware
\end{itemize}

\begin{center}\rule{0.5\linewidth}{0.5pt}\end{center}

\section{Troubleshooting Flowchart}\label{troubleshooting-flowchart}

\begin{verbatim}
No CAN communication
    │
    ├─ Power off: ~60Ω between CAN_H & CAN_L?
    │   └─ No → Fix termination & bus topology
    │
    ├─ Real transceiver on every node, in normal mode?
    │   └─ No → Add transceiver or drive STBY/S to normal
    │
    ├─ CAN_H→CAN_H, CAN_L→CAN_L, GND shared?
    │   └─ No → Fix wiring & polarity
    │
    ├─ ≥2 active nodes to provide ACK?
    │   └─ No → Add second active node or use self-test loopback
    │
    ├─ Classical CAN works at low nominal bitrate?
    │   └─ No → Recalculate nominal bit timing & sample point
    │
    ├─ Using CAN FD?
    │   ├─ No → Add IDs, filters, interrupts, then raise load gradually
    │   └─ Yes → Nominal + data phase bit timing + TDC configured?
    │       └─ No → Set dbitrate, data sample point, TDC/SSP
    │
    ├─ Error counters rising / bus-off?
    │   └─ Yes → Read TEC/REC, inspect waveform, fix root cause, recover bus
    │
    └─ Check filters, std/ext ID choice, software races & queue overruns
\end{verbatim}

\begin{center}\rule{0.5\linewidth}{0.5pt}\end{center}

\section{Minimal Debugging Code}\label{minimal-debugging-code}

\subsection{Linux SocketCAN (C) --- includes error frame
reporting}\label{linux-socketcan-c-includes-error-frame-reporting}

\begin{verbatim}
#include <stdio.h>
#include <string.h>
#include <unistd.h>
#include <sys/ioctl.h>
#include <sys/socket.h>
#include <net/if.h>
#include <linux/can.h>
#include <linux/can/raw.h>

int main(void) {
    int s = socket(PF_CAN, SOCK_RAW, CAN_RAW);
    if (s < 0) return 1;

    // Subscribe to error frames for debugging
    can_err_mask_t err_mask = CAN_ERR_TX_TIMEOUT | CAN_ERR_BUSOFF;
    setsockopt(s, SOL_CAN_RAW, CAN_RAW_ERR_FILTER, &err_mask, sizeof(err_mask));

    struct ifreq ifr = {0};
    strncpy(ifr.ifr_name, "can0", IFNAMSIZ - 1);
    if (ioctl(s, SIOCGIFINDEX, &ifr) < 0) return 2;

    struct sockaddr_can addr = {0};
    addr.can_family  = AF_CAN;
    addr.can_ifindex = ifr.ifr_ifindex;
    if (bind(s, (struct sockaddr *)&addr, sizeof(addr)) < 0) return 3;

    struct can_frame f = {0};
    f.can_id  = 0x123;
    f.can_dlc = 2;
    f.data[0] = 0xAA;
    f.data[1] = 0x55;

    if (write(s, &f, sizeof(f)) != sizeof(f)) return 4;
    close(s);
    return 0;
}
\end{verbatim}

\subsection{C++ --- ISR-safe double buffer
pattern}\label{c-isr-safe-double-buffer-pattern}

\begin{verbatim}
#include <array>
#include <atomic>
#include <cstdint>

struct CanFrame {
    uint32_t id{};
    uint8_t  dlc{};
    std::array<uint8_t, 8> data{};
};

static CanFrame isr_buffer;
static std::atomic<bool> frame_ready{false};

void can_rx_isr(const CanFrame& hw_frame) {
    isr_buffer = hw_frame;                          // ISR writes only
    frame_ready.store(true, std::memory_order_release);
}

bool try_fetch_frame(CanFrame& out) {
    if (!frame_ready.exchange(false, std::memory_order_acq_rel))
        return false;
    out = isr_buffer;                               // foreground consumes only
    return true;
}
\end{verbatim}

\begin{quote}
For sustained traffic, replace the single slot with a ring buffer or RTOS queue.
\end{quote}

\subsection{Python (python-can) --- with thread safety \&
filters}\label{python-python-can-with-thread-safety-filters}

\begin{verbatim}
import can

filters = [
    {"can_id": 0x123, "can_mask": 0x7FF, "extended": False},
]

bus = can.ThreadSafeBus(interface="socketcan", channel="vcan0", can_filters=filters)
notifier = can.Notifier(bus, listeners=[can.Printer()], timeout=0.2)

try:
    msg = can.Message(arbitration_id=0x123,
                      data=[0x01, 0x02, 0x03],
                      is_extended_id=False)
    bus.send(msg, timeout=0.5)
    rx = bus.recv(timeout=1.0)
    print("RX:", rx)
finally:
    notifier.stop()
    bus.shutdown()
\end{verbatim}

\begin{quote}
Use \texttt{vcan0} to test logic before involving real hardware.
\end{quote}

\begin{center}\rule{0.5\linewidth}{0.5pt}\end{center}

\section{Reference Materials}\label{reference-materials}

\subsection{Official specifications}\label{official-specifications}

\begin{itemize}
\tightlist
\item
  ISO 11898-1 (protocol) and ISO 11898-2 (physical layer) --- current editions:
  2024 and 2026
\item
  \href{https://www.bosch-semiconductors.com/ip-modules/can-protocols/}{Bosch
  CAN Protocols} --- primary source for CAN FD
\item
  \href{https://www.bosch-semiconductors.com/ip-modules/can-fd-ip-modules/can-fd/}{Bosch
  CAN FD overview}
\end{itemize}

\subsection{Vendor documentation}\label{vendor-documentation}

\begin{itemize}
\tightlist
\item
  Microchip AN754 --- bit timing for CAN
\item
  Microchip CAN mechanism overview, filtering, and error-handling pages
\item
  NXP bit-timing guidance and calculator
\item
  TI platform FAQ, MCAN debug/getting-started notes, and physical-layer debug
  article
\item
  \href{https://www.kernel.org/doc/html/latest/networking/can.html}{Linux
  SocketCAN documentation}
\item
  \href{https://python-can.readthedocs.io/}{python-can API documentation}
\end{itemize}

\subsection{Tools}\label{tools}

\begin{itemize}
\tightlist
\item
  \textbf{CAN analyzers:} PCAN-View (PEAK), Kvaser, KJElectronics
\item
  \textbf{Linux can-utils:} \texttt{cangen}, \texttt{candump},
  \texttt{canplayer}, \texttt{cansequence}
\item
  \textbf{Virtual CAN:} \texttt{vcan0} --- test logic without any hardware
\item
  \textbf{Oscilloscope} with CAN/CAN FD decoding (e.g., PicoScope)
\end{itemize}

\subsection{Books \& specs}\label{books-specs}

\begin{itemize}
\tightlist
\item
  Bosch CAN 2.0B specification
\item
  CANopen DS301/DS303 (for higher-layer protocol reference)
\end{itemize}

\begin{center}\rule{0.5\linewidth}{0.5pt}\end{center}

\emph{See also: {[}{[}can-protocol{]}{]} for protocol theory,
{[}{[}can-hardware-basics{]}{]} for physical setup,
{[}{[}can-addressing-for-etrike{]}{]} for our project's ID assignments.}


\chapter{Distributed Safety Patterns}\label{distributed-safety-patterns}

Safety-critical embedded systems don't rely on a single mechanism to prevent
hazards. They stack multiple independent layers so that if one fails, the next
catches it. This is \textbf{defense in depth} --- a pattern borrowed from
nuclear, aviation, and automotive safety engineering.

The E-Trike has 8 independent safety layers (§6 of architecture.md). Each layer
assumes the one above it might fail.

\begin{center}\rule{0.5\linewidth}{0.5pt}\end{center}

\section{1. Defense in Depth --- The Core
Principle}\label{defense-in-depth-the-core-principle}

\begin{verbatim}
┌──────────────────────────────────────┐
│  Layer 1: Physical ESTOP button      │  ← fastest, simplest, hardest to defeat
├──────────────────────────────────────┤
│  Layer 2: CAN 0x001 SAFETY_ESTOP     │  ← redundant comm path for ESTOP
├──────────────────────────────────────┤
│  Layer 3: Heartbeat timeout          │  ← detects crashed/frozen nodes
├──────────────────────────────────────┤
│  Layer 4: Command staleness          │  ← detects hung planning stack
├──────────────────────────────────────┤
│  Layer 5: Steering following error   │  ← detects stuck linkage / bad actuator
├──────────────────────────────────────┤
│  Layer 6: Dynamic steering clamp     │  ← prevents dangerous commands
├──────────────────────────────────────┤
│  Layer 7: Software hard-stops        │  ← prevents mechanical damage
├──────────────────────────────────────┤
│  Layer 8: Brake light OR logic       │  ← fail-visible: always shows braking
└──────────────────────────────────────┘
\end{verbatim}

\textbf{The rule for layers:} each layer must be \emph{independent} --- a
failure that defeats one layer must not defeat others. Independence comes from:
- Different hardware (GPIO vs CAN vs watchdog IC) - Different software (separate
tasks, separate nodes) - Different physical principles (electrical vs optical vs
mechanical)

\begin{center}\rule{0.5\linewidth}{0.5pt}\end{center}

\section{2. ESTOP as an Absorbing State}\label{estop-as-an-absorbing-state}

An emergency stop isn't just ``stop the motor.'' It must guarantee:

\begin{enumerate}
\def\labelenumi{\arabic{enumi}.}
\tightlist
\item
  \textbf{All actuators return to safe state.} Motor → 0V. Gear → neutral. Brake
  → max. Steering → stop commanding (actuator timeout-faults). DC-DC → off. 12V
  relay → off.
\item
  \textbf{The state persists.} No automatic recovery. No ``ESTOP cleared,
  resuming.'' The system stays stopped until a deliberate human action.
\item
  \textbf{The exit path is explicit and guarded.} The E-Trike requires pressing
  the START button (GPIO32) --- not the MODE button, not a CAN command, not a
  timeout. And ESTOP exits to MANUAL (rider in direct control), never to AUTO.
\end{enumerate}

\begin{verbatim}
ESTOP exit rule:
  START button → MANUAL mode
  (NOT: MODE button, CAN command, power-cycle timeout)
\end{verbatim}

\textbf{Why this matters:} If ESTOP could be cleared by the same buggy software
that triggered it, you'd have an oscillating system --- ESTOP → clear → ESTOP →
clear --- with the vehicle lurching between stop and go.

\begin{center}\rule{0.5\linewidth}{0.5pt}\end{center}

\section{3. Fail-Safe vs Fail-Operational}\label{fail-safe-vs-fail-operational}

\begin{longtable}[]{@{}
  >{\raggedright\arraybackslash}p{(\columnwidth - 4\tabcolsep) * \real{0.2632}}
  >{\raggedright\arraybackslash}p{(\columnwidth - 4\tabcolsep) * \real{0.5000}}
  >{\raggedright\arraybackslash}p{(\columnwidth - 4\tabcolsep) * \real{0.2368}}@{}}
\toprule\noalign{}
\begin{minipage}[b]{\linewidth}\raggedright
Strategy
\end{minipage} & \begin{minipage}[b]{\linewidth}\raggedright
Behavior on failure
\end{minipage} & \begin{minipage}[b]{\linewidth}\raggedright
Example
\end{minipage} \\
\midrule\noalign{}
\endhead
\bottomrule\noalign{}
\endlastfoot
\textbf{Fail-safe} & System enters safe state, stops functioning & ESTOP: motor
off, brake on, 12V off \\
\textbf{Fail-operational} & System continues with degraded capability & Jetson
fails → RT zeros setpoints → vehicle coasts (motor off but steering still
responds to rider) \\
\end{longtable}

The E-Trike uses \textbf{fail-safe} for the safety-critical path (ESTOP kills
everything) but \textbf{fail-operational} for non-critical failures (Jetson loss
= coast, not ESTOP). The distinction matters: you don't want to ESTOP on every
Jetson hiccup, but you do want to ESTOP on steering fault.

\begin{center}\rule{0.5\linewidth}{0.5pt}\end{center}

\section{4. Normally-Closed (NC) Wiring}\label{normally-closed-nc-wiring}

A \textbf{normally-closed} switch conducts current in its resting state.
Pressing it \emph{breaks} the circuit.

\begin{verbatim}
Normal operation:  GPIO reads HIGH (pull-up holds it) → system runs
Button pressed:    GPIO reads LOW (switch connects to GND) → ESTOP
Wire cut/broken:   GPIO reads LOW (pull-up defeated, floating→LOW) → ESTOP
Power loss:        GPIO reads LOW → ESTOP
\end{verbatim}

An NC switch is \textbf{fail-safe by construction.} A broken wire, disconnected
plug, or power loss all trigger ESTOP --- they can't produce a ``silent
failure'' where the button is broken but the system thinks it's fine.

The alternative --- normally-open (NO) --- is cheaper but unsafe for ESTOP: a
cut wire looks exactly like ``button not pressed.''

\begin{center}\rule{0.5\linewidth}{0.5pt}\end{center}

\section{5. Threshold + Duration Debounce}\label{threshold-duration-debounce}

Safety checks don't trigger on a single out-of-bounds sample. They wait for the
condition to \emph{persist}.

\begin{verbatim}
if (abs(cmd - actual) > THRESHOLD) {
    error_duration_ms += dt;
    if (error_duration_ms > DURATION) {
        mode_set(Estop);  // confirmed persistent fault
    }
} else {
    error_duration_ms = 0;  // transient — reset counter
}
\end{verbatim}

\textbf{Why duration matters:} CAN frames can be delayed by arbitration (a
low-priority message waits for a high-priority one to finish). A single missed
or delayed frame is a transient, not a failure. The duration confirms ``this is
real, not noise.''

\begin{longtable}[]{@{}
  >{\raggedright\arraybackslash}p{(\columnwidth - 6\tabcolsep) * \real{0.1489}}
  >{\raggedright\arraybackslash}p{(\columnwidth - 6\tabcolsep) * \real{0.2340}}
  >{\raggedright\arraybackslash}p{(\columnwidth - 6\tabcolsep) * \real{0.2128}}
  >{\raggedright\arraybackslash}p{(\columnwidth - 6\tabcolsep) * \real{0.4043}}@{}}
\toprule\noalign{}
\begin{minipage}[b]{\linewidth}\raggedright
Check
\end{minipage} & \begin{minipage}[b]{\linewidth}\raggedright
Threshold
\end{minipage} & \begin{minipage}[b]{\linewidth}\raggedright
Duration
\end{minipage} & \begin{minipage}[b]{\linewidth}\raggedright
Why this duration?
\end{minipage} \\
\midrule\noalign{}
\endhead
\bottomrule\noalign{}
\endlastfoot
Steering following error & \textgreater5° & 300 ms & Stuck linkage is
immediately dangerous. Short debounce. \\
Command staleness & No \texttt{0x300} & 500 ms & Jetson hiccups are common.
Longer tolerance. \\
Heartbeat timeout & No fresh counter & 200 ms (inter-MCU) & FTTI-bound. Must
catch within 1.4 m of travel. \\
\end{longtable}

\begin{center}\rule{0.5\linewidth}{0.5pt}\end{center}

\section{6. Independent Safety Layers --- Node
Separation}\label{independent-safety-layers-node-separation}

A single MCU running all safety + actuation code can fail entirely from one bug
(stack overflow, wild pointer, priority inversion). Physical separation means:

\begin{verbatim}
SYS ESP32-S3 (safety, motor, brake)          RT ESP32-S3 (physics, steering, gateway)
         │                                             │
         │  0x7FD/0x7FE heartbeats                     │
         │◄────────────────────────────────────────────►│
         │                                             │
         │  If RT crashes: SYS detects HB timeout       │
         │  → triggers CAN 0x001 ESTOP                  │
         │  → motor stops, brake engages                │
         │  → SYS continues running independently       │
\end{verbatim}

SYS doesn't need RT to be alive to stop the vehicle. The ESTOP button, brake
lever, and motor kill all work directly on SYS GPIOs. This is the architectural
decoupling that makes defense in depth real --- not just multiple \texttt{if}
statements on the same chip.

\begin{center}\rule{0.5\linewidth}{0.5pt}\end{center}

\section{7. Fail-Visible: The Brake Light OR
Logic}\label{fail-visible-the-brake-light-or-logic}

Some safety states must be \textbf{visible} to the outside world, not just
enforced internally:

\begin{verbatim}
brake_light_on = brake_lever_pressed()     // physical lever — local GPIO
              OR (mode == Estop)           // ESTOP — always brake
              OR g_light_state.brake;      // Jetson CAN 0x302 — predictive
\end{verbatim}

All sources are OR'd together. The brake light illuminates if ANY braking reason
exists. There's no code path that can produce ``braking but light off.'' This is
\textbf{fail-visible}: the safety state is externally observable, so a following
vehicle or rider can see it even if other systems have failed.

\begin{center}\rule{0.5\linewidth}{0.5pt}\end{center}

\section{8. Safety Pattern Checklist}\label{safety-pattern-checklist}

When designing a safety-critical embedded system, verify:

\begin{itemize}
\tightlist
\item[$\square$]
  Every safety function has at least two independent implementations.
\item[$\square$]
  No single point of failure defeats all safety layers.
\item[$\square$]
  ESTOP is an absorbing state --- deliberate exit only.
\item[$\square$]
  Safety-critical inputs are NC (fail to safe state on disconnect).
\item[$\square$]
  Threshold checks include a persistence duration (debounce).
\item[$\square$]
  Safety-critical outputs are fail-visible where possible (lights, indicators).
\item[$\square$]
  Nodes that must catch each other's failures exchange independent heartbeats.
\item[$\square$]
  The external watchdog is on a different chip from the MCU.
\end{itemize}

\begin{center}\rule{0.5\linewidth}{0.5pt}\end{center}

\emph{Primary reference: \texttt{docs/emergency-system.md} for the complete
ESTOP system, all trigger paths, emergency response matrix, rider's guide, and
testing procedures.}

\emph{See also: {[}{[}heartbeat-monitoring{]}{]} for liveness detection,
{[}{[}external-watchdog{]}{]} for hardware watchdog,
{[}{[}state-machine-design{]}{]} for ESTOP FSM, \texttt{architecture.md} §6 for
the E-Trike's safety layers.}


\chapter{Endianness \& Binary Protocols}\label{endianness-binary-protocols}

When a CAN frame carries a 16-bit or 32-bit value (speed, angle, pressure), the
sender and receiver must agree on \emph{byte order}. Getting it wrong silently
swaps bytes --- your 3000 mm/s becomes 768 mm/s, or worse, a negative number.

The E-Trike uses \textbf{big-endian (MSB first)} for multi-byte CAN fields ---
except SYNTREE actuators which use \textbf{little-endian (LSB first)} per their
factory protocol. The \texttt{can-dictionary.md} marks endianness on every
signal.

\begin{center}\rule{0.5\linewidth}{0.5pt}\end{center}

\section{1. Big-Endian vs Little-Endian}\label{big-endian-vs-little-endian}

\textbf{Big-endian (MSB first):} The most significant byte comes first in
memory.

\begin{verbatim}
Value: 0x12345678 (32-bit)
Memory:  [12] [34] [56] [78]
          ↑                 ↑
        low addr         high addr
\end{verbatim}

\textbf{Little-endian (LSB first):} The least significant byte comes first.

\begin{verbatim}
Value: 0x12345678 (32-bit)
Memory:  [78] [56] [34] [12]
          ↑                 ↑
        low addr         high addr
\end{verbatim}

\textbf{How to remember:} Big-endian is ``human readable'' --- the bytes appear
in the order you'd write the number. Little-endian is ``computer efficient'' ---
the low byte is at the low address, making addition and carry propagation faster
in hardware.

\subsection{Practical example}\label{practical-example}

\begin{verbatim}
uint32_t speed = 0x00000BB8;  // 3000 in decimal, 0x0BB8 in hex

// Big-endian on CAN bus: [00] [00] [0B] [B8]
// Little-endian on CAN bus: [B8] [0B] [00] [00]
\end{verbatim}

If the receiver expects big-endian but gets little-endian, it reads
\texttt{0xB80B0000} = 3,087,433,728. Completely wrong, no obvious error.

\begin{center}\rule{0.5\linewidth}{0.5pt}\end{center}

\section{2. Bit Numbering --- Motorola vs Intel
Format}\label{bit-numbering-motorola-vs-intel-format}

CAN signal descriptions specify a \textbf{start bit} and \textbf{length}. But
``start bit'' means different things to different conventions:

\begin{longtable}[]{@{}lll@{}}
\toprule\noalign{}
Convention & Start bit meaning & Used by \\
\midrule\noalign{}
\endhead
\bottomrule\noalign{}
\endlastfoot
\textbf{Motorola (big-endian)} & MSB of the signal & Classical CAN DBs,
AUTOSAR \\
\textbf{Intel (little-endian)} & LSB of the signal & Many modern tools \\
\end{longtable}

For example, a 16-bit signal at start bit 0, length 16:

\begin{verbatim}
Motorola: bit 0 = MSB of signal → signal bits are [0..15] = [MSB..LSB]
Intel:    bit 0 = LSB of signal → signal bits are [0..15] = [LSB..MSB]
\end{verbatim}

\textbf{The E-Trike convention:} All signals use big-endian (Motorola) unless
marked otherwise. SYNTREE protocols use ``Motorola LSB'' --- Motorola bit
numbering but little-endian byte order. See the signal tables in
\texttt{can-dictionary.md}.

\begin{center}\rule{0.5\linewidth}{0.5pt}\end{center}

\section{3. Why CAN Uses Big-Endian}\label{why-can-uses-big-endian}

CAN itself is big-endian at the bit level: the most significant bit of the CAN
ID is transmitted first during arbitration. This is why lower IDs (more leading
zeros in the MSB) win arbitration --- the first dominant bit (0) from a
competing node overrides a recessive bit (1) from another, and the lower-ID
frame continues.

Most automotive standards (J1939, CANopen) follow this convention and define
multi-byte signals as big-endian. The E-Trike does the same for consistency ---
except where SYNTREE factory protocols dictate little-endian.

\begin{center}\rule{0.5\linewidth}{0.5pt}\end{center}

\section{4. C/C++ Struct Packing}\label{cc-struct-packing}

Embedded firmware often maps a C struct directly onto CAN payload bytes:

\begin{verbatim}
// Big-endian: RT_DRIVE_CMD (0x204), DLC=5
struct __attribute__((packed)) RtDriveCmd {
    uint8_t  speed_byte3;  // MSB of speed
    uint8_t  speed_byte2;
    uint8_t  speed_byte1;
    uint8_t  speed_byte0;  // LSB of speed
    uint8_t  gear;
};

// Convert struct to uint32_t speed:
int32_t speed = (cmd.speed_byte3 << 24) |
                (cmd.speed_byte2 << 16) |
                (cmd.speed_byte1 << 8)  |
                cmd.speed_byte0;
\end{verbatim}

\textbf{Three rules for safe struct mapping:}

\begin{enumerate}
\def\labelenumi{\arabic{enumi}.}
\tightlist
\item
  \textbf{Always \texttt{\_\_attribute\_\_((packed))}} --- prevents the compiler
  from inserting padding bytes that shift everything.
\item
  \textbf{Always verify with \texttt{static\_assert(sizeof(Struct)\ ==\ DLC)}}
  --- catches compiler quirks.
\item
  \textbf{Never use bitfields for CAN protocols} --- the C standard doesn't
  specify bitfield ordering, so the same code on different compilers (GCC vs
  Clang vs IAR) produces different byte layouts.
\end{enumerate}

Bad (non-portable):

\begin{verbatim}
struct Bad {
    uint8_t mode : 2;     // ← ordering is compiler-defined!
    uint8_t flags : 6;
};
\end{verbatim}

Good (portable):

\begin{verbatim}
struct Good {
    uint8_t byte0;        // ← explicit bit operations extract fields
};
uint8_t mode = byte0 & 0x03;
uint8_t flags = (byte0 >> 2) & 0x3F;
\end{verbatim}

\begin{center}\rule{0.5\linewidth}{0.5pt}\end{center}

\section{5. Endianness Conversion}\label{endianness-conversion}

When your CPU's native endianness differs from the protocol:

\begin{verbatim}
// Host-to-network (big-endian) — for standard CAN signals
uint32_t htonl(uint32_t host) {
    return ((host & 0xFF) << 24) |
           ((host & 0xFF00) << 8) |
           ((host & 0xFF0000) >> 8) |
           ((host & 0xFF000000) >> 24);
}

// For SYNTREE (little-endian), just memcpy if the CPU is little-endian (ESP32 is).
// If your CPU is big-endian, reverse.
\end{verbatim}

The ESP32 (Xtensa LX7) is \textbf{little-endian}. So: - For E-Trike big-endian
CAN signals: reverse the bytes before sending. - For SYNTREE little-endian
signals: use native byte order --- no conversion needed.

\begin{center}\rule{0.5\linewidth}{0.5pt}\end{center}

\section{6. Common Gotchas}\label{common-gotchas}

\begin{longtable}[]{@{}
  >{\raggedright\arraybackslash}p{(\columnwidth - 4\tabcolsep) * \real{0.2353}}
  >{\raggedright\arraybackslash}p{(\columnwidth - 4\tabcolsep) * \real{0.3824}}
  >{\raggedright\arraybackslash}p{(\columnwidth - 4\tabcolsep) * \real{0.3824}}@{}}
\toprule\noalign{}
\begin{minipage}[b]{\linewidth}\raggedright
Gotcha
\end{minipage} & \begin{minipage}[b]{\linewidth}\raggedright
What happens
\end{minipage} & \begin{minipage}[b]{\linewidth}\raggedright
How to catch
\end{minipage} \\
\midrule\noalign{}
\endhead
\bottomrule\noalign{}
\endlastfoot
\textbf{Forgot \texttt{packed} attribute} & Compiler inserts padding, shifts
fields by 1--3 bytes &
\texttt{static\_assert(sizeof(MyStruct)\ ==\ expected\_dlc,\ "wrong\ size")} \\
\textbf{Bitfield across compilers} & GCC and Clang order bitfields differently &
Use explicit bit masks, never bitfields for protocols \\
\textbf{Mixed endianness on one bus} & Some signals correct, others silently
wrong & Every signal in can-dictionary.md has an endianness column \\
\textbf{Sign extension on narrow types} & Casting \texttt{uint8\_t} to
\texttt{int16\_t} sign-extends bit 7 & Use explicit masks:
\texttt{value\ =\ (int16\_t)((raw\_high\ \textless{}\textless{}\ 8)\ \textbar{}\ raw\_low)}
--- but watch for the \texttt{int16\_t} cast after OR, not before \\
\textbf{Endianness of the CAN ID itself} & Mismatch between \texttt{0x200} in
code and on the bus & The CAN controller handles ID endianness. Just use the
numeric value. \\
\end{longtable}

\begin{center}\rule{0.5\linewidth}{0.5pt}\end{center}

\emph{See also: {[}{[}can-protocol{]}{]} for CAN frame structure,
\texttt{can-dictionary.md} for bit-level layouts of every signal,
\texttt{docs/steering-unit.md} for SYNTREE little-endian layout.}


\chapter{Heartbeat \& Liveness Monitoring}\label{heartbeat-liveness-monitoring}

In a distributed system, a node can fail silently. Its CAN controller might keep
retransmitting the last frame from a DMA buffer while the CPU is frozen. Its
power supply might dip below brownout threshold. Its crystal might drift out of
tolerance. \textbf{Liveness monitoring} answers the question: \emph{is the node
on the other end of the bus still alive and thinking?}

The E-Trike uses heartbeat frames (\texttt{0x7FD}, \texttt{0x7FE},
\texttt{0x7FC}) on every CAN bus for this purpose.

\begin{center}\rule{0.5\linewidth}{0.5pt}\end{center}

\section{1. The Problem: Silent Failure}\label{the-problem-silent-failure}

Consider a CAN bus without heartbeats. Node A sends commands at 50 Hz. Node B
receives them. One day, Node B's CPU freezes but its CAN controller (a separate
hardware block) keeps DMA-ing the last valid frame from its TX mailbox. Node A
sees a perfect 50 Hz stream and assumes Node B is healthy.

Node B is dead. Node A doesn't know. The system continues operating with no
safety monitor.

\textbf{This is not theoretical.} It happens when: - A watchdog reset leaves the
CAN controller initialized but the application hung in a boot loop. - A power
glitch browns out the CPU core but not the CAN transceiver. - A stack overflow
corrupts the task but the CAN TX mailbox was already loaded.

\begin{center}\rule{0.5\linewidth}{0.5pt}\end{center}

\section{2. The Alive Counter --- Why Not Just ``Frame
Present''?}\label{the-alive-counter-why-not-just-frame-present}

A naive heartbeat sends an empty frame periodically. The receiver checks ``did I
get a frame recently?'' But with the DMA replay problem above, the CAN
controller answers ``yes'' --- forever.

The fix is an \textbf{alive counter}: a single byte that increments every frame.

\begin{verbatim}
Frame 1: alive_ctr = 0x01
Frame 2: alive_ctr = 0x02
Frame 3: alive_ctr = 0x03   ← CPU freezes here
Frame 4: alive_ctr = 0x03   ← CAN controller replays frame 3 forever
Frame 5: alive_ctr = 0x03   ← receiver detects: same counter twice → node is dead
\end{verbatim}

The receiver stores the last counter value. If a new frame arrives with the same
counter, the node is frozen --- even if the frame timing is perfect.

\begin{verbatim}
bool heartbeat_is_fresh(uint8_t new_ctr) {
    if (new_ctr != last_ctr) {
        last_ctr = new_ctr;
        return true;       // counter changed — node is alive
    }
    return false;          // frozen — same counter twice, CAN controller on autopilot
}
\end{verbatim}

\begin{center}\rule{0.5\linewidth}{0.5pt}\end{center}

\section{3. Timeout Selection --- FTTI}\label{timeout-selection-ftti}

The \textbf{Fault Tolerant Time Interval (FTTI)} is the maximum time a fault can
persist before it becomes hazardous. It comes from ISO 26262 (functional safety)
and drives every timeout in a safety-critical system.

\begin{longtable}[]{@{}
  >{\raggedright\arraybackslash}p{(\columnwidth - 4\tabcolsep) * \real{0.2000}}
  >{\raggedright\arraybackslash}p{(\columnwidth - 4\tabcolsep) * \real{0.3600}}
  >{\raggedright\arraybackslash}p{(\columnwidth - 4\tabcolsep) * \real{0.4400}}@{}}
\toprule\noalign{}
\begin{minipage}[b]{\linewidth}\raggedright
Bus
\end{minipage} & \begin{minipage}[b]{\linewidth}\raggedright
Timeout
\end{minipage} & \begin{minipage}[b]{\linewidth}\raggedright
Rationale
\end{minipage} \\
\midrule\noalign{}
\endhead
\bottomrule\noalign{}
\endlastfoot
Low CAN (RT↔SYS) & \textbf{200 ms} & At 25 km/h (\textasciitilde7 m/s), the
trike travels \textasciitilde1.4 m. A steering fault persisting longer than this
puts the vehicle in an adjacent lane or off the road. \\
High CAN (RT↔Jetson) & \textbf{500 ms} & Jetson failure results in a
\emph{controlled stop} (zero setpoints), not an immediate ESTOP. The rider can
still override. 500 ms is long enough to survive a Jetson hiccup, short enough
that the vehicle doesn't coast into an intersection. \\
\end{longtable}

\textbf{The rule:} faster vehicle × more critical system = shorter timeout. 200
ms for inter-MCU is an automotive-grade choice. A passenger car at highway speed
demands \textless100 ms FTTI for steering; at 25 km/h, 200 ms is conservative.

\begin{center}\rule{0.5\linewidth}{0.5pt}\end{center}

\section{4. Startup Grace Period}\label{startup-grace-period}

At boot, no heartbeats have been exchanged yet. If you check immediately, you'd
trigger ESTOP before any node has had a chance to send its first frame.

\textbf{Solution:} a startup grace period (typically 3 seconds) where
\texttt{heartbeat\_ok()} returns \texttt{true} regardless.

\begin{verbatim}
bool heartbeat_ok() {
    int64_t now = esp_timer_get_time();
    if (last_hb_us == 0) {
        // Never received a heartbeat yet — still in startup
        return (now < kStartupGraceUs);  // 3 seconds
    }
    return (now - last_hb_us) < kTimeoutUs;  // 200 ms
}
\end{verbatim}

After the grace period, the safety monitor arms and real heartbeat checking
begins. If a node never comes online, the grace period expires and ESTOP
triggers --- preventing operation with a dead safety node.

\begin{center}\rule{0.5\linewidth}{0.5pt}\end{center}

\section{5. Independent Buses, Independent
Heartbeats}\label{independent-buses-independent-heartbeats}

Heartbeat frames MUST NOT be bridged between CAN buses. Each bus is an
independent liveness domain.

\textbf{Why bridging heartbeats creates ambiguity:}

\begin{verbatim}
Low bus:  SYS sends 0x7FE {ctr=42}, RT sends 0x7FD {ctr=17}
          ↓ RT bridges SYS's 0x7FE to high bus
High bus: Jetson receives 0x7FE {ctr=42} — whose counter is this?
          SYS? RT? Is the counter monotonic? Which node died?
\end{verbatim}

With the same CAN ID on the same bus, you can't tell which sender is which. The
E-Trike now uses separate IDs for separate nodes on the low bus (\texttt{0x7FD}
for RT, \texttt{0x7FE} for SYS) and monitors each independently.

\textbf{The liveness matrix} (who monitors whom):

\begin{longtable}[]{@{}lllll@{}}
\toprule\noalign{}
Monitor & Watches & Bus & Timeout & Action on loss \\
\midrule\noalign{}
\endhead
\bottomrule\noalign{}
\endlastfoot
SYS & RT (\texttt{0x7FD}) & Low & 200 ms & ESTOP (AUTO only) \\
RT & SYS (\texttt{0x7FE}) & Low & 200 ms & CAN \texttt{0x001} ESTOP \\
RT & Jetson (\texttt{0x7FC}) & High & 500 ms & Zero setpoints \\
Jetson & RT (\texttt{0x7FD}) & High & 500 ms & Stop publishing \\
\end{longtable}

RT monitors both buses because it's the gateway --- it's the only node that
knows about failures on both sides.

\begin{center}\rule{0.5\linewidth}{0.5pt}\end{center}

\section{6. Heartbeat vs.~External
Watchdog}\label{heartbeat-vs.-external-watchdog}

A heartbeat catches a frozen node on the CAN bus. An \textbf{external watchdog}
catches a frozen MCU regardless of CAN state (see
{[}{[}external-watchdog{]}{]}). They're complementary:

\begin{longtable}[]{@{}
  >{\raggedright\arraybackslash}p{(\columnwidth - 4\tabcolsep) * \real{0.1579}}
  >{\centering\arraybackslash}p{(\columnwidth - 4\tabcolsep) * \real{0.3509}}
  >{\centering\arraybackslash}p{(\columnwidth - 4\tabcolsep) * \real{0.4912}}@{}}
\toprule\noalign{}
\begin{minipage}[b]{\linewidth}\raggedright
Failure
\end{minipage} & \begin{minipage}[b]{\linewidth}\centering
Heartbeat catches?
\end{minipage} & \begin{minipage}[b]{\linewidth}\centering
External watchdog catches?
\end{minipage} \\
\midrule\noalign{}
\endhead
\bottomrule\noalign{}
\endlastfoot
CPU hung, CAN controller alive & ✓ (same counter twice) & ✗ (watchdog GPIO may
still toggle from DMA) \\
CPU hung, CAN controller dead & ✓ (no frames) & ✓ \\
Crystal failure & ✓ (no frames) & ✓ (watchdog IC has independent oscillator) \\
CAN transceiver dead & ✓ (no frames from that node) & ✗ (MCU is fine) \\
Power brownout & ✓ (no frames) & ✓ \\
\end{longtable}

Neither alone is sufficient for safety-critical systems. Both together cover the
failure space.

\begin{center}\rule{0.5\linewidth}{0.5pt}\end{center}

\section{7. Common Pitfalls}\label{common-pitfalls}

\begin{longtable}[]{@{}
  >{\raggedright\arraybackslash}p{(\columnwidth - 4\tabcolsep) * \real{0.3333}}
  >{\raggedright\arraybackslash}p{(\columnwidth - 4\tabcolsep) * \real{0.4815}}
  >{\raggedright\arraybackslash}p{(\columnwidth - 4\tabcolsep) * \real{0.1852}}@{}}
\toprule\noalign{}
\begin{minipage}[b]{\linewidth}\raggedright
Pitfall
\end{minipage} & \begin{minipage}[b]{\linewidth}\raggedright
What happens
\end{minipage} & \begin{minipage}[b]{\linewidth}\raggedright
Fix
\end{minipage} \\
\midrule\noalign{}
\endhead
\bottomrule\noalign{}
\endlastfoot
\textbf{Empty heartbeat frame} & CAN controller replays it forever --- receiver
can't detect CPU freeze & Always include an alive counter byte \\
\textbf{Heartbeat on same ID from multiple nodes} & Ambiguous which node is dead
& Assign a unique CAN ID to each heartbeat sender, or use separate buses \\
\textbf{Bridging heartbeats} & Two different counters on the same bus with the
same ID --- receiver can't distinguish & Never forward heartbeat frames. Each
bus is its own liveness domain. \\
\textbf{Timeout too short} & Transient bus errors cause false ESTOP & Set
timeout ≥ 3× the heartbeat period. At 2 Hz (500 ms), timeout ≥ 150 ms. \\
\textbf{Timeout too long} & Vehicle travels dangerous distance before ESTOP &
Bound by FTTI. At 25 km/h with 200 ms timeout: 1.4 m travel. \\
\textbf{No startup grace period} & ESTOP at boot before any heartbeats are sent
& Mask liveness checks for the first 3 seconds \\
\end{longtable}

\begin{center}\rule{0.5\linewidth}{0.5pt}\end{center}

\emph{See also: {[}{[}external-watchdog{]}{]} for hardware watchdog,
{[}{[}can-gateway-bridging{]}{]} for why heartbeats stay local,
\texttt{architecture.md} §8.6 for heartbeat implementation,
\texttt{can-dictionary.md} §1 0x7FD/0x7FE/0x7FC for frame layout.}


\chapter{PID Control Fundamentals}\label{pid-control-fundamentals}

A \textbf{PID controller} closes the loop between a desired value (setpoint) and
a measured value (process variable), producing an output that drives the error
toward zero. It's the workhorse of embedded control --- motor speed,
temperature, position, voltage --- anywhere you need a system to track a target.

The E-Trike uses PID on the RT ESP32-S3 for rear-motor speed control. Jetson
says ``go 1500 mm/s'', RT measures the actual speed from the encoder, the PID
computes a DAC value for the motor controller, and the loop repeats at 100 Hz.

\begin{center}\rule{0.5\linewidth}{0.5pt}\end{center}

\section{1. The Closed-Loop Concept}\label{the-closed-loop-concept}

\begin{verbatim}
                    ┌──────────┐
setpoint ────►[+]──►│   PID    │──► output ──► [Plant / Motor] ──► actual speed
              ▲     └──────────┘                                    │
              │                                                     │
              └────────────────── measurement ◄─────────────────────┘
\end{verbatim}

Without feedback (open-loop): you command 50\% throttle, the vehicle
does\ldots{} something. Uphill it crawls, downhill it speeds. The controller has
no idea.

With feedback (closed-loop): the PID sees the error and adjusts. Uphill → error
grows → integral term ramps up → throttle increases until the error is zero.

\begin{center}\rule{0.5\linewidth}{0.5pt}\end{center}

\section{2. The Three Terms}\label{the-three-terms}

\subsection{Proportional (P)}\label{proportional-p}

\begin{verbatim}
P_out = Kp × error
\end{verbatim}

\begin{itemize}
\tightlist
\item
  \textbf{What it does:} Reacts to the \emph{present} error. The further from
  target, the harder it pushes.
\item
  \textbf{Intuition:} Like a spring --- the force is proportional to how far
  you've stretched it.
\item
  \textbf{Too high:} Oscillation. The controller overshoots, corrects,
  overshoots the other way.
\item
  \textbf{Too low:} Sluggish. The vehicle takes too long to reach target speed.
\item
  \textbf{Alone:} Never reaches the target exactly. There's always a small
  residual error (steady-state offset) because as error → 0, P\_out → 0.
\end{itemize}

\subsection{Integral (I)}\label{integral-i}

\begin{verbatim}
integral += error × dt
integral = clamp(integral, −max, +max)
I_out = Ki × integral
\end{verbatim}

\begin{itemize}
\tightlist
\item
  \textbf{What it does:} Accumulates \emph{past} error. If the vehicle cruises
  50 mm/s below target for 2 seconds, the integral term grows until the offset
  is eliminated.
\item
  \textbf{Intuition:} Like a water bucket under a dripping faucet --- the longer
  the error persists, the fuller the bucket, and the stronger the correction.
\item
  \textbf{Why clamping matters (anti-windup):} If the motor can't physically
  reach the target (uphill, headwind), the integral grows without bound. When
  the load reduces, the inflated integral causes a massive overshoot.
  \textbf{Clamp the integral} to prevent this.
\item
  \textbf{When to reset:} On mode change (ESTOP→MANUAL) or on large setpoint
  changes, zero the integral. Otherwise, accumulated error from the previous
  operating point causes a transient surge.
\end{itemize}

\subsection{Derivative (D)}\label{derivative-d}

\begin{verbatim}
derivative = (error − prev_error) / dt
D_out = Kd × derivative
\end{verbatim}

\begin{itemize}
\tightlist
\item
  \textbf{What it does:} Reacts to the \emph{rate of change} of error. If the
  vehicle is approaching the target quickly, the D term backs off to prevent
  overshoot.
\item
  \textbf{Intuition:} Like a shock absorber --- it resists rapid changes. The
  faster you approach, the harder it pushes back.
\item
  \textbf{Noise problem:} The derivative amplifies measurement noise. A 1-bit
  encoder jitter looks like a massive instantaneous speed change.
\end{itemize}

\textbf{Derivative on measurement (the fix):}

Instead of differentiating the error (which jumps on setpoint changes),
differentiate the measurement directly:

\begin{verbatim}
// Derivative on error — causes "derivative kick" on setpoint step
derivative = (error - prev_error) / dt;

// Derivative on measurement — smooth response to setpoint changes
derivative = -(measurement - prev_measurement) / dt;
\end{verbatim}

The setpoint only enters through the P and I terms. The D term only damps actual
system motion. This avoids a spike in output when the setpoint changes abruptly.

\begin{center}\rule{0.5\linewidth}{0.5pt}\end{center}

\section{3. The Complete PID Equation}\label{the-complete-pid-equation}

\begin{verbatim}
float pid_update(float setpoint, float measurement, float dt) {
    float error = setpoint - measurement;

    // Proportional — present error
    float p_term = Kp * error;

    // Integral — accumulated past error (with anti-windup)
    integral += error * dt;
    integral = clamp(integral, -max_integral, +max_integral);
    float i_term = Ki * integral;

    // Derivative on measurement — rate of change (no kick)
    float derivative = -(measurement - prev_measurement) / dt;
    float d_term = Kd * derivative;

    prev_measurement = measurement;

    return clamp(p_term + i_term + d_term, out_min, out_max);
}
\end{verbatim}

\begin{center}\rule{0.5\linewidth}{0.5pt}\end{center}

\section{4. Tuning --- Making It Work}\label{tuning-making-it-work}

Tuning is finding Kp, Ki, Kd values that give fast response without oscillation.

\subsection{Manual tuning procedure}\label{manual-tuning-procedure}

\begin{enumerate}
\def\labelenumi{\arabic{enumi}.}
\tightlist
\item
  \textbf{Set Ki = Kd = 0.} Start with P only.
\item
  \textbf{Increase Kp} until the system oscillates with \emph{constant
  amplitude} (neither growing nor decaying). This is the \textbf{ultimate gain}
  Ku.
\item
  \textbf{Set Kp = 0.6 × Ku.} This gives a stable proportional response.
\item
  \textbf{Add I.} Increase Ki until steady-state error is eliminated within
  \textasciitilde2 seconds of a step change. Watch for overshoot --- if the
  system overshoots and slowly settles, Ki is too high.
\item
  \textbf{Add D (if needed).} Increase Kd to dampen overshoot. On a high-inertia
  system (vehicle), D helps. On a low-inertia system, D often amplifies noise
  --- skip it.
\end{enumerate}

\subsection{What good tuning looks like}\label{what-good-tuning-looks-like}

\begin{verbatim}
Speed ▲
      │        ┌──────────────  setpoint
      │       ╱
      │      ╱                 ← slight overshoot (5-10%), quick settling
      │     ╱
      │    ╱
      │   ╱
      │  ╱
      │ ╱
      │╱
      └────────────────────────► time
\end{verbatim}

\begin{itemize}
\tightlist
\item
  \textbf{Rise time} (10\%→90\%): fast but not aggressive
\item
  \textbf{Overshoot}: \textless10\% of setpoint step
\item
  \textbf{Settling time} (±5\% of setpoint): \textless1 second
\item
  \textbf{No sustained oscillation} at steady state
\end{itemize}

\begin{center}\rule{0.5\linewidth}{0.5pt}\end{center}

\section{5. The E-Trike's PID in Context}\label{the-e-trikes-pid-in-context}

The PID on RT is the \textbf{outermost speed loop}:

\begin{verbatim}
Jetson cmd_vel ──► Kinematics ──► speed setpoint ──► PID ──► DAC value ──► Motor Controller ──► Motor
                         ▲                               │
                         │                               │
                         └── encoder feedback ◄──────────┘
\end{verbatim}

The motor controller (a separate unit) runs its own internal current/torque loop
at 10--20 kHz. The E-Trike's PID only needs to be fast enough to track
vehicle-speed dynamics --- 100 Hz is more than sufficient for a
\textasciitilde300 kg tricycle with seconds-scale acceleration.

\subsection{What the PID does NOT control}\label{what-the-pid-does-not-control}

\begin{itemize}
\tightlist
\item
  \textbf{Steering angle} --- EPS-C runs its own internal position loop. RT
  commands an angle; the unit's firmware handles the motor.
\item
  \textbf{Brake pressure} --- SEB runs its own internal pressure/stroke loop.
  SYS commands a stroke; the unit handles the pump.
\item
  \textbf{Motor commutation} --- The motor controller handles FOC
  (Field-Oriented Control) internally. RT only sends a 0--5 V speed reference.
\end{itemize}

\begin{center}\rule{0.5\linewidth}{0.5pt}\end{center}

\section{6. When PID Isn't Enough}\label{when-pid-isnt-enough}

PID is the right tool for single-input, single-output (SISO) systems with
linear-ish behavior. You need more when:

\begin{longtable}[]{@{}ll@{}}
\toprule\noalign{}
Situation & Better approach \\
\midrule\noalign{}
\endhead
\bottomrule\noalign{}
\endlastfoot
Multiple interacting variables (speed + steering) & MIMO control, MPC \\
Large dead time (seconds of delay) & Smith predictor \\
Highly nonlinear plant & Gain scheduling, adaptive control \\
Hard constraints (must never exceed limit) & MPC with constraints \\
\end{longtable}

For a tricycle speed controller, PID is the right choice --- simple,
well-understood, and sufficient.

\begin{center}\rule{0.5\linewidth}{0.5pt}\end{center}

\emph{See also: {[}{[}physics-model{]}{]} for the kinematics that produce the
speed setpoint, {[}{[}actuator-interfacing{]}{]} for the MCP4725 DAC,
\texttt{architecture.md} §7.6 for the control loop integration.}


\chapter{RTOS Task Design Fundamentals}\label{rtos-task-design-fundamentals}

A \textbf{Real-Time Operating System (RTOS)} lets you split firmware into
independent \emph{tasks} that the scheduler interleaves on one or more CPU
cores. Without an RTOS, embedded code runs in a single \texttt{while(1)} loop
--- add one blocking call and everything stalls. With an RTOS, if one task
blocks waiting for a CAN frame, the scheduler runs another task instead.

The E-Trike uses \textbf{FreeRTOS} on both ESP32-S3 nodes (9 tasks on RT, 15 on
SYS). Every design decision --- task count, priority assignments, queue usage
--- follows from concepts below.

\begin{center}\rule{0.5\linewidth}{0.5pt}\end{center}

\section{1. Preemptive Scheduling}\label{preemptive-scheduling}

A \textbf{preemptive scheduler} interrupts a running task when a higher-priority
task becomes ready. The higher-priority task runs immediately --- no cooperative
``yield'' needed.

\begin{verbatim}
Time ──────────────────────────────────────────────►

Task A (pri 2)  ████████████░░░░░░░░░░░░░░░░░░████████
Task B (pri 5)  ░░░░░░░░████████████░░░░░░░░░░░░░░░░░░
                 ▲                    ▲
                 │ B unblocks         │ B blocks again
                 │ (CAN frame arrives)│ A resumes where it left off
\end{verbatim}

The key property: \textbf{the highest-priority ready task always runs.} If two
tasks share a priority, they time-slice (round-robin).

\begin{center}\rule{0.5\linewidth}{0.5pt}\end{center}

\section{2. Task Priorities --- Why They
Matter}\label{task-priorities-why-they-matter}

FreeRTOS uses numeric priorities (0 = lowest, configMAX\_PRIORITIES−1 =
highest). Choosing them is a design decision, not an afterthought.

\begin{longtable}[]{@{}
  >{\raggedright\arraybackslash}p{(\columnwidth - 4\tabcolsep) * \real{0.3846}}
  >{\raggedright\arraybackslash}p{(\columnwidth - 4\tabcolsep) * \real{0.4872}}
  >{\raggedright\arraybackslash}p{(\columnwidth - 4\tabcolsep) * \real{0.1282}}@{}}
\toprule\noalign{}
\begin{minipage}[b]{\linewidth}\raggedright
Priority level
\end{minipage} & \begin{minipage}[b]{\linewidth}\raggedright
What belongs here
\end{minipage} & \begin{minipage}[b]{\linewidth}\raggedright
Why
\end{minipage} \\
\midrule\noalign{}
\endhead
\bottomrule\noalign{}
\endlastfoot
\textbf{Highest (5)} & CAN receive, safety GPIO poll & Must never miss a frame
or button press. If the CAN RX task is delayed, frames are lost. \\
\textbf{High (4)} & Dispatch, control loop, mode FSM & Business logic. Runs
frequently, must not be starved by telemetry. \\
\textbf{Medium (3)} & CAN transmit, actuator FSM, lights & Periodic output.
Jitter is acceptable (a few ms late won't break anything). \\
\textbf{Low (2)} & Obstacle sensor, indicators, power relay & Nice to have.
Dropped reads are tolerably stale. \\
\textbf{Lowest (1)} & Diagnostics, heartbeat & Background housekeeping. Can be
delayed indefinitely without safety impact. \\
\end{longtable}

\textbf{Rule of thumb from the E-Trike design:} - Input tasks (CAN RX, GPIO) go
at the top --- data must be captured. - Processing tasks (dispatch, control) go
one level down --- transform captured data. - Output tasks (CAN TX, actuator
write) go at medium --- push results out. - Telemetry and housekeeping go at the
bottom --- never block safety.

\begin{center}\rule{0.5\linewidth}{0.5pt}\end{center}

\section{3. Queues --- Why ``Queues Over Shared
State''}\label{queues-why-queues-over-shared-state}

A FreeRTOS queue is a thread-safe FIFO pipe. One task sends, another receives.
No mutexes, no semaphores.

\begin{verbatim}
Task A (CAN RX)          Queue (16 slots)          Task B (Dispatch)
     │                                                 │
     ├── xQueueSend(q, &frame, timeout) ──►            │
     │                                                 ├── xQueueReceive(q, &frame, timeout)
     │                                                 │
\end{verbatim}

\textbf{Why queues over shared state (global variables + mutexes)?}

\begin{longtable}[]{@{}
  >{\raggedright\arraybackslash}p{(\columnwidth - 2\tabcolsep) * \real{0.5263}}
  >{\raggedright\arraybackslash}p{(\columnwidth - 2\tabcolsep) * \real{0.4737}}@{}}
\toprule\noalign{}
\begin{minipage}[b]{\linewidth}\raggedright
Approach
\end{minipage} & \begin{minipage}[b]{\linewidth}\raggedright
Problem
\end{minipage} \\
\midrule\noalign{}
\endhead
\bottomrule\noalign{}
\endlastfoot
Global variable + mutex & Priority inversion: low-prio task holds mutex,
high-prio task blocked waiting --- effectively runs at low prio until mutex is
released. \\
Global variable + no mutex & Data races: ISR writes byte 3 while foreground
reads byte 1 → corrupted frame. \\
Queue & No locks, no races. ISR-safe \texttt{xQueueSendFromISR()}. Sender and
receiver have clear ownership. \\
\end{longtable}

The E-Trike's design principle \#1 captures this: \emph{Queues over shared
state. No mutexes, no semaphores.}

\subsection{Queue overflow}\label{queue-overflow}

Every queue has a fixed depth. If the sender produces faster than the receiver
consumes:

\begin{verbatim}
xQueueSend(q, &frame, 0)  →  returns errQUEUE_FULL
\end{verbatim}

The sender must decide: drop the oldest (overwrite), drop the newest (skip), or
block until space is available. The E-Trike uses \textbf{overwrite} for setpoint
queues (only the latest command matters) and \textbf{drop} for telemetry queues
(a missed periodic update is replaced by the next one).

\begin{center}\rule{0.5\linewidth}{0.5pt}\end{center}

\section{4. Tick Rate}\label{tick-rate}

FreeRTOS's \textbf{tick interrupt} fires at a fixed frequency (1000 Hz on the
E-Trike). Every tick, the scheduler checks if a higher-priority task is ready.

\begin{longtable}[]{@{}
  >{\raggedright\arraybackslash}p{(\columnwidth - 2\tabcolsep) * \real{0.4583}}
  >{\raggedright\arraybackslash}p{(\columnwidth - 2\tabcolsep) * \real{0.5417}}@{}}
\toprule\noalign{}
\begin{minipage}[b]{\linewidth}\raggedright
Tick rate
\end{minipage} & \begin{minipage}[b]{\linewidth}\raggedright
Implication
\end{minipage} \\
\midrule\noalign{}
\endhead
\bottomrule\noalign{}
\endlastfoot
100 Hz (10 ms) & Coarse. Tasks get 10 ms granularity. A 100 Hz control loop
needs to run every other tick. \\
\textbf{1000 Hz (1 ms)} & Standard for motor control. 100 Hz loop runs every 10
ticks. Jitter \textless1 ms. \\
10000 Hz (100 µs) & Aggressive. CPU spends significant time in tick ISR. Only
for sub-ms control. \\
\end{longtable}

The E-Trike uses 1000 Hz because the steering CAN command at 50 Hz (20 ms
period) needs tight periodicity --- a missed 20 ms window causes EPS-C comm
fault.

\begin{center}\rule{0.5\linewidth}{0.5pt}\end{center}

\section{5. Task Periodicity Patterns}\label{task-periodicity-patterns}

\subsection{Periodic task (fixed rate)}\label{periodic-task-fixed-rate}

\begin{verbatim}
void periodic_task() {
    TickType_t last_wake = xTaskGetTickCount();
    while (true) {
        // Do work
        vTaskDelayUntil(&last_wake, pdMS_TO_TICKS(10));  // 100 Hz = 10 ms
    }
}
\end{verbatim}

\texttt{vTaskDelayUntil} corrects for drift --- if the work takes 3 ms, it
sleeps 7 ms (not 10 ms). Over time, the period stays exact.

\subsection{Event-driven task (blocking)}\label{event-driven-task-blocking}

\begin{verbatim}
void event_task() {
    while (true) {
        CanFrame frame;
        if (xQueueReceive(rx_queue, &frame, portMAX_DELAY) == pdTRUE) {
            process(frame);
        }
    }
}
\end{verbatim}

The task blocks (consumes zero CPU) until a frame arrives. This is the pattern
for CAN RX and dispatch tasks.

\begin{center}\rule{0.5\linewidth}{0.5pt}\end{center}

\section{6. Stack Sizing}\label{stack-sizing}

Every FreeRTOS task has its own stack. Too small → stack overflow → corrupted
memory (hardest bug to debug). Too large → wasted RAM.

\begin{longtable}[]{@{}
  >{\raggedright\arraybackslash}p{(\columnwidth - 4\tabcolsep) * \real{0.3667}}
  >{\raggedright\arraybackslash}p{(\columnwidth - 4\tabcolsep) * \real{0.4667}}
  >{\raggedright\arraybackslash}p{(\columnwidth - 4\tabcolsep) * \real{0.1667}}@{}}
\toprule\noalign{}
\begin{minipage}[b]{\linewidth}\raggedright
Task type
\end{minipage} & \begin{minipage}[b]{\linewidth}\raggedright
Typical stack
\end{minipage} & \begin{minipage}[b]{\linewidth}\raggedright
Why
\end{minipage} \\
\midrule\noalign{}
\endhead
\bottomrule\noalign{}
\endlastfoot
CAN RX (event) & 4096 B & Needs buffer for incoming frame + function call
depth \\
Control (periodic) & 4096 B & Floating-point math (kinematics, PID) uses stack
heavily \\
CAN TX (event) & 3072 B & Simpler call chain, smaller structs \\
GPIO poll & 2048 B & Minimal logic, no deep calls \\
Heartbeat & 2048 B & Counter increment + CAN send \\
\end{longtable}

\textbf{How to verify:} Fill the stack with a known pattern (0xA5) at task
creation, then inspect the high-water mark after runtime. ESP-IDF provides
\texttt{uxTaskGetStackHighWaterMark()}.

\begin{center}\rule{0.5\linewidth}{0.5pt}\end{center}

\section{7. Common Pitfalls}\label{common-pitfalls}

\begin{longtable}[]{@{}
  >{\raggedright\arraybackslash}p{(\columnwidth - 4\tabcolsep) * \real{0.3333}}
  >{\raggedright\arraybackslash}p{(\columnwidth - 4\tabcolsep) * \real{0.4815}}
  >{\raggedright\arraybackslash}p{(\columnwidth - 4\tabcolsep) * \real{0.1852}}@{}}
\toprule\noalign{}
\begin{minipage}[b]{\linewidth}\raggedright
Pitfall
\end{minipage} & \begin{minipage}[b]{\linewidth}\raggedright
What happens
\end{minipage} & \begin{minipage}[b]{\linewidth}\raggedright
Fix
\end{minipage} \\
\midrule\noalign{}
\endhead
\bottomrule\noalign{}
\endlastfoot
\textbf{Priority inversion} & Low-prio task holds a resource, medium-prio task
starves it, high-prio task blocks on the resource & Use queues (lock-free) or
priority inheritance mutexes \\
\textbf{Task starvation} & A high-prio task never blocks → lower-prio tasks
never run & Every high-prio task must block (queue, delay, semaphore) \\
\textbf{Too many tasks} & Context-switch overhead dominates CPU time & One task
per independent rate or blocking boundary. Don't split logic into unnecessary
micro-tasks. \\
\textbf{Blocking in a high-prio task} & The blocking call waits forever → all
lower-prio tasks starve & Always use timeouts on blocking calls in high-prio
tasks \\
\textbf{printf in an ISR} & ISR blocks on UART → system hangs & ISRs must only
do queue sends and GPIO toggles. No I/O. \\
\end{longtable}

\begin{center}\rule{0.5\linewidth}{0.5pt}\end{center}

\section{8. Why FreeRTOS for Safety-Critical
Embedded?}\label{why-freertos-for-safety-critical-embedded}

\begin{itemize}
\tightlist
\item
  \textbf{Preemptive} --- a safety task at priority 5 runs within 1 tick (1 ms)
  of its condition becoming true.
\item
  \textbf{Memory isolation via tasks} --- a stack overflow in the diagnostics
  task doesn't corrupt the CAN RX task's stack.
\item
  \textbf{Mature and audited} --- deployed in millions of devices,
  MISRA-compliant variant available.
\item
  \textbf{Lightweight} --- runs on a \$2 ESP32-S3 with plenty of headroom.
\end{itemize}

The alternative --- bare-metal superloop --- works for simple devices. But when
you have 9+ concurrent concerns at different rates (100 Hz control, 50 Hz CAN, 2
Hz heartbeat, event-driven dispatch), an RTOS is the right tool.

\begin{center}\rule{0.5\linewidth}{0.5pt}\end{center}

\emph{See also: {[}{[}pid-speed-control{]}{]} for the 100 Hz control loop,
{[}{[}can-protocol{]}{]} for CAN RX/TX task design, \texttt{architecture.md}
§7.7 for RT task layout, §8.7 for SYS task layout.}


\chapter{Shared Code Architecture --- Monorepo Embedded
Pattern}\label{shared-code-architecture-monorepo-embedded-pattern}

When a system has \textbf{multiple MCUs} that must agree on protocol definitions
(CAN IDs, frame layouts, enums), you need a \textbf{single source of truth} for
those definitions --- but you can't link them at runtime because each MCU is an
independent physical device with its own flash.

The e-trike uses a \textbf{header-only shared library} consumed by each project
via compiler include paths. No symlinks, no copies, no separate build artifact.

\begin{center}\rule{0.5\linewidth}{0.5pt}\end{center}

\section{The Problem}\label{the-problem}

\begin{verbatim}
┌─────────────┐     CAN bus      ┌─────────────┐
│  RT ESP32   │◄═══════════════►│  SYS ESP32   │
└─────────────┘                  └─────────────┘

Both must agree on:
  • CAN message IDs (0x169 = steering, 0x205 = brake, …)
  • Frame layouts (which byte is speed, gear, mode)
  • Enum values (Mode::Auto = 1, Gear::D = 1)
  • Serialisation order (big-endian vs little-endian per field)
\end{verbatim}

If you duplicate these definitions, they \textbf{will} drift apart --- a field
gets reordered in one file but not the other, and suddenly the brake command
reads as a speed value on the other side.

\begin{center}\rule{0.5\linewidth}{0.5pt}\end{center}

\section{The Pattern: Header-Only Shared
Library}\label{the-pattern-header-only-shared-library}

\begin{verbatim}
project/
├── shared/               ← single source of truth
│   ├── can/
│   │   ├── can_protocol.h    # CAN IDs, Frame struct, all payload types
│   │   └── can_driver.h      # RAII TWAI driver (ESP-IDF dependent)
│   └── os/
│       ├── endian.h           # Big-endian read/write helpers
│       ├── queue.h            # Type-safe FreeRTOS queue wrapper
│       └── result.h           # Result<T,E> error handling
│
├── rt-esp32/             ← consumer A
│   ├── platformio.ini        # build_flags: -I ../shared
│   └── src/
│       └── main.cpp          # #include "can/can_protocol.h"
│
└── sys-esp32/            ← consumer B
    ├── platformio.ini        # build_flags: -I ../shared
    └── src/
        └── main.cpp          # #include "can/can_protocol.h"
\end{verbatim}

\subsection{Key properties}\label{key-properties}

\begin{enumerate}
\def\labelenumi{\arabic{enumi}.}
\item
  \textbf{No \texttt{.cpp} files in \texttt{shared/}.} Everything is
  \texttt{inline}, \texttt{constexpr}, or \texttt{template}. The shared code
  compiles directly into each consumer's translation units --- no shared object,
  no library to link.
\item
  \textbf{Include via path, not symlink.} Each \texttt{platformio.ini} adds
  \texttt{-I\ ../shared} to \texttt{build\_flags}. When the compiler sees
  \texttt{\#include\ "can/can\_protocol.h"}, it searches the include path and
  finds \texttt{../shared/can/can\_protocol.h}.
\item
  \textbf{Each MCU gets its own copy.} The shared code is compiled into
  \texttt{rt-esp32.bin} and into \texttt{sys-esp32.bin} independently. There's
  no runtime sharing --- the CAN bus is the only communication channel between
  them.
\end{enumerate}

\begin{center}\rule{0.5\linewidth}{0.5pt}\end{center}

\section{Why Header-Only?}\label{why-header-only}

\begin{longtable}[]{@{}ll@{}}
\toprule\noalign{}
Constraint & Why it rules out alternatives \\
\midrule\noalign{}
\endhead
\bottomrule\noalign{}
\endlastfoot
Embedded (no OS linker) & Can't load a \texttt{.so} or \texttt{.dll} at
runtime \\
Separate flash chips & Can't share a memory-mapped code region \\
No filesystem on MCU & Can't distribute definitions as data files \\
Compile-time is free & \texttt{inline}/\texttt{constexpr} has zero runtime
cost \\
\end{longtable}

Header-only is the \textbf{lightest-weight sharing mechanism} available: the
compiler does all the work, and the linker sees nothing new.

\begin{center}\rule{0.5\linewidth}{0.5pt}\end{center}

\section{The Include Path Mechanism}\label{the-include-path-mechanism}

In PlatformIO (ESP-IDF framework):

\begin{verbatim}
; rt-esp32/platformio.ini (identical line in sys-esp32/platformio.ini)
build_flags =
    -I ../shared
\end{verbatim}

This tells GCC/Clang: \emph{``when resolving \texttt{\#include} directives, also
look in \texttt{../shared} relative to the project directory.''}

The include chain works like this:

\begin{verbatim}
main.cpp                           shared/can/can_protocol.h
┌──────────────────────────┐       ┌──────────────────────────┐
│ #include "can/can_protocol.h"────►│ #include "os/endian.h"──┼──► shared/os/endian.h
│                          │       │                          │
│ // uses can::Frame       │       │ namespace can { … }     │
│ // uses can::Mode        │       └──────────────────────────┘
│ // uses can::VcuSesReq   │
└──────────────────────────┘
\end{verbatim}

Each \texttt{\#include} with quotes (\texttt{"…"}) checks the current directory
first, then the \texttt{-I} paths. Angle-bracket includes
(\texttt{\textless{}…\textgreater{}}) check only system paths --- FreeRTOS and
ESP-IDF headers use this.

\begin{center}\rule{0.5\linewidth}{0.5pt}\end{center}

\section{Optional: Per-Project Wrappers}\label{optional-per-project-wrappers}

When the shared code needs project-specific configuration, each consumer adds a
thin wrapper:

\begin{verbatim}
// sys-esp32/src/can_driver.h — NOT a copy, just configuration glue
#include "can/can_driver.h"   // the shared one
#include "config.h"           // sys-esp32's own pin definitions

namespace sys {
inline can::CanDriver::Config make_can_config() {
    can::CanDriver::Config cfg;
    cfg.tx_gpio    = sys::kCanTxGpio;    // 4 for SYS, 5 for RT
    cfg.rx_gpio    = sys::kCanRxGpio;
    cfg.bitrate_hz = sys::kCanBitrateHz;
    return cfg;
}
}
\end{verbatim}

The wrapper \textbf{does not duplicate} the shared driver --- it only fills in
local constants. The heavy logic stays in \texttt{shared/}.

\begin{center}\rule{0.5\linewidth}{0.5pt}\end{center}

\section{How It Lands on Hardware}\label{how-it-lands-on-hardware}

\begin{verbatim}
shared/                  rt-esp32/                  sys-esp32/
(headers only)           (consumer)                 (consumer)
                             │                          │
                    ┌────────▼────────┐      ┌─────────▼────────┐
                    │ pio run          │      │ pio run          │
                    │ g++ -I ../shared │      │ g++ -I ../shared │
                    │ … main.cpp       │      │ … main.cpp       │
                    └────────┬────────┘      └─────────┬────────┘
                             │                          │
                    ┌────────▼────────┐      ┌─────────▼────────┐
                    │ rt-esp32.elf    │      │ sys-esp32.elf    │
                    │ (shared code    │      │ (shared code     │
                    │  compiled in)   │      │  compiled in)    │
                    └────────┬────────┘      └─────────┬────────┘
                             │                          │
                    ┌────────▼────────┐      ┌─────────▼────────┐
                    │ pio run -t      │      │ pio run -t       │
                    │ upload          │      │ upload           │
                    └────────┬────────┘      └─────────┬────────┘
                             │                          │
                    ┌────────▼────────┐      ┌─────────▼────────┐
                    │ RT ESP32-S3     │      │ SYS ESP32-S3     │
                    │ (flash chip)    │      │ (flash chip)      │
                    └─────────────────┘      └──────────────────┘

                    Each MCU has its own independent
                    copy of the shared definitions.
                    They agree because they came from
                    the same source files.
\end{verbatim}

\begin{center}\rule{0.5\linewidth}{0.5pt}\end{center}

\section{When This Pattern Works}\label{when-this-pattern-works}

\begin{longtable}[]{@{}
  >{\raggedright\arraybackslash}p{(\columnwidth - 2\tabcolsep) * \real{0.5000}}
  >{\raggedright\arraybackslash}p{(\columnwidth - 2\tabcolsep) * \real{0.5000}}@{}}
\toprule\noalign{}
\begin{minipage}[b]{\linewidth}\raggedright
✅ Good fit
\end{minipage} & \begin{minipage}[b]{\linewidth}\raggedright
❌ Poor fit
\end{minipage} \\
\midrule\noalign{}
\endhead
\bottomrule\noalign{}
\endlastfoot
Protocol definitions (CAN IDs, enums, frame layouts) & Large function bodies
(bloats every consumer) \\
Lightweight inline helpers (serialisation, bit packing) & Stateful singletons
(each MCU gets its own copy --- that's usually wrong) \\
Template/constexpr utilities & Code that needs per-project conditional
compilation (\texttt{\#ifdef} hell) \\
2--10 consumers in a monorepo & 50+ consumers (consider a proper library with
versioning) \\
\end{longtable}

\begin{center}\rule{0.5\linewidth}{0.5pt}\end{center}

\section{Alternatives}\label{alternatives}

\begin{longtable}[]{@{}
  >{\raggedright\arraybackslash}p{(\columnwidth - 2\tabcolsep) * \real{0.4348}}
  >{\raggedright\arraybackslash}p{(\columnwidth - 2\tabcolsep) * \real{0.5652}}@{}}
\toprule\noalign{}
\begin{minipage}[b]{\linewidth}\raggedright
Approach
\end{minipage} & \begin{minipage}[b]{\linewidth}\raggedright
When to use
\end{minipage} \\
\midrule\noalign{}
\endhead
\bottomrule\noalign{}
\endlastfoot
\textbf{Git submodule} & Shared code lives in a separate repo; consumers pin a
specific commit. Good when the shared code has its own release cycle. \\
\textbf{Static library (.a)} & Shared code has \texttt{.cpp} files; you want to
compile once and link. Works on embedded if the toolchain supports it, but adds
build complexity. \\
\textbf{Code generation} & Protocol definitions are the source of truth (e.g., a
DBC file); each project runs a generator to produce its own header. Avoids the
include-path assumption entirely. \\
\textbf{Copy-paste} & Never. Don't. \\
\end{longtable}

\begin{center}\rule{0.5\linewidth}{0.5pt}\end{center}

\section{Related Notes}\label{related-notes}

\begin{itemize}
\tightlist
\item
  {[}{[}can-protocol{]}{]} --- what the shared definitions actually define
\item
  {[}{[}endianness-binary-protocols{]}{]} --- why the serialisation helpers
  exist
\item
  {[}{[}rtos-task-design{]}{]} --- how the shared queue wrapper is used in tasks
\item
  {[}{[}can-gateway-design{]}{]} --- how RT bridges between the two CAN buses
  using these shared IDs
\end{itemize}


\chapter{State Machine Design for Embedded
Systems}\label{state-machine-design-for-embedded-systems}

A \textbf{state machine} (finite state machine, FSM) models a system that is
always in exactly one \emph{state} and transitions between states in response to
\emph{events}. In embedded firmware, state machines replace tangled
\texttt{if/else} chains with explicit, testable behavior.

The E-Trike uses four state machines: - \textbf{Mode FSM} --- MANUAL → AUTO →
ESTOP (architecture §3) - \textbf{Steering boot FSM} --- BOOT\_WAIT →
LISTEN\_SYNC → ACTIVE → FAULT (§7.6) - \textbf{Brake boot FSM} --- BOOT\_WAIT →
LISTEN\_SYNC → ACTIVE → FAULT (§8.6) - \textbf{Gear FSM} --- N → D → S → R
transitions (§8.6)

\begin{center}\rule{0.5\linewidth}{0.5pt}\end{center}

\section{1. Why State Machines?}\label{why-state-machines}

Without a state machine, mode-dependent behavior looks like this:

\begin{verbatim}
// Anti-pattern: flag soup
void motor_task() {
    if (estop_active) {
        dac_write(0);
        gear_off();
    } else if (auto_mode && drive_cmd_valid) {
        dac_write(speed_to_dac(setpoint));
        gear_set(setpoint.gear);
    } else if (manual_mode) {
        dac_write(adc_read());
        gear_mirror_tlp281();
    }
    // What about transitioning? What about startup? What if two flags are true?
}
\end{verbatim}

Two problems: (1) the logic is scattered across many \texttt{if} blocks in many
tasks, and (2) contradictory states (what if \texttt{auto\_mode} AND
\texttt{estop\_active} are both true?) produce undefined behavior.

A state machine makes the behavior explicit:

\begin{verbatim}
enum class Mode { Manual, Auto, Estop };

void motor_task() {
    switch (g_mode) {
        case Mode::Manual:  motor_manual();  break;
        case Mode::Auto:    motor_auto();    break;
        case Mode::Estop:   motor_estop();   break;
    }
}
\end{verbatim}

\textbf{One variable, one switch, one behavior per state.} The FSM \emph{is} the
documentation.

\begin{center}\rule{0.5\linewidth}{0.5pt}\end{center}

\section{2. States, Events, Transitions}\label{states-events-transitions}

A state machine has three ingredients:

\begin{longtable}[]{@{}
  >{\raggedright\arraybackslash}p{(\columnwidth - 4\tabcolsep) * \real{0.3438}}
  >{\raggedright\arraybackslash}p{(\columnwidth - 4\tabcolsep) * \real{0.3750}}
  >{\raggedright\arraybackslash}p{(\columnwidth - 4\tabcolsep) * \real{0.2812}}@{}}
\toprule\noalign{}
\begin{minipage}[b]{\linewidth}\raggedright
Ingredient
\end{minipage} & \begin{minipage}[b]{\linewidth}\raggedright
What it is
\end{minipage} & \begin{minipage}[b]{\linewidth}\raggedright
Example
\end{minipage} \\
\midrule\noalign{}
\endhead
\bottomrule\noalign{}
\endlastfoot
\textbf{State} & A distinct mode of operation & \texttt{MANUAL}, \texttt{AUTO},
\texttt{ESTOP} \\
\textbf{Event} & Something that triggers a change & Button press, CAN frame,
timeout \\
\textbf{Transition} & Moving from one state to another &
\texttt{MANUAL\ ──{[}ESTOP\ btn{]}──►\ ESTOP} \\
\end{longtable}

\begin{verbatim}
         ┌──────────┐
    ┌───►│  MANUAL  │◀───┐
    │    └─────┬────┘    │
    │     MODE btn       MODE btn
    │          │          │
    │    ┌─────▼────┐    │
    │    │   AUTO   │    │
    │    └─────┬────┘    │
    │          │          │
    │  ESTOP btn / CAN 0x001 / HB timeout
    │          │          │
    │    ┌─────▼────┐    │
    │    │  ESTOP   │────┘
    │    └─────┬────┘
    │         │ START btn
    │         ▼
    └──────── MANUAL
\end{verbatim}

Key property: \textbf{ESTOP is an absorbing state} --- once entered, only a
deliberate START action exits. No automatic recovery.

\begin{center}\rule{0.5\linewidth}{0.5pt}\end{center}

\section{3. Implementation Patterns in
C/C++}\label{implementation-patterns-in-cc}

\subsection{Pattern 1: Switch-case (flat
FSM)}\label{pattern-1-switch-case-flat-fsm}

Best for simple machines with \textless8 states.

\begin{verbatim}
enum class State { BootWait, ListenSync, Active, Fault };
State g_state = State::BootWait;

void steer_fsm() {
    switch (g_state) {
    case State::BootWait:
        if (millis_since_boot() > 500) {
            g_state = State::ListenSync;
        }
        break;

    case State::ListenSync:
        if (received_ses_status) {
            active_target = ses_angle;
            if (ses_aligned) g_state = State::Active;
        } else if (time_in_state() > 2000) {
            g_state = State::Fault;
        }
        break;

    case State::Active:
        send_steer_cmd();
        if (following_error_too_large()) {
            mode_set(Estop);  // triggers system-level ESTOP
        }
        break;

    case State::Fault:
        // Silent — actuator timeout-faults internally
        break;
    }
}
\end{verbatim}

\subsection{Pattern 2: Function pointer
table}\label{pattern-2-function-pointer-table}

Best when each state has complex behavior and you want separate functions.

\begin{verbatim}
typedef void (*state_fn_t)();

void st_boot_wait()  { /* ... */ }
void st_listen_sync() { /* ... */ }
void st_active()     { /* ... */ }
void st_fault()      { /* ... */ }

state_fn_t state_table[] = { st_boot_wait, st_listen_sync, st_active, st_fault };

void steer_fsm() {
    state_table[g_state]();  // dispatch
}
\end{verbatim}

\subsection{Pattern 3: Transition table}\label{pattern-3-transition-table}

Best when transitions are complex and you want them in one place.

\begin{verbatim}
struct Transition {
    State from, to;
    bool (*guard)();      // condition function
    void (*action)();     // what to do on transition
};

Transition table[] = {
    { BootWait, ListenSync, timeout_500ms,  nullptr },
    { ListenSync, Active,  ses_aligned,    nullptr },
    { Active, Fault,       timeout_2s,     log_fault },
    // ...
};
\end{verbatim}

The E-Trike uses \textbf{Pattern 1 (switch-case)} --- simple, readable, and
sufficient for each FSM's complexity.

\begin{center}\rule{0.5\linewidth}{0.5pt}\end{center}

\section{4. Entry and Exit Actions}\label{entry-and-exit-actions}

When transitioning between states, you often need cleanup/initialization:

\begin{verbatim}
void mode_set(Mode new_mode) {
    Mode old = g_mode;           // snapshot for exit action
    g_mode = new_mode;           // atomic — all tasks see the new mode immediately

    // Exit actions
    if (old == Mode::Auto) {
        steer_stop_transmitting();   // EPS-C comm fault on timeout
    }

    // Entry actions
    if (new_mode == Mode::Estop) {
        dac_write(0);                // motor off
        gear_all_off();              // neutral
        brake_full();                // 0x7B9 stroke = max
        dcdc_disable();              // 12V rail off
        lights_brake_on();           // brake light forced ON
    }
}
\end{verbatim}

\textbf{Important:} \texttt{g\_mode} is written \emph{before} entry actions run.
This is intentional --- if a high-priority safety task reads \texttt{g\_mode}
mid-transition, it sees the new mode and acts accordingly, even if entry actions
haven't finished yet.

\begin{center}\rule{0.5\linewidth}{0.5pt}\end{center}

\section{5. Hierarchical State Machines}\label{hierarchical-state-machines}

A flat FSM works for 3--5 states. When you have mode logic AND sub-mode logic,
you need hierarchy:

\begin{verbatim}
MANUAL
  ├── Throttle: pass-through
  ├── Gear: pass-through
  └── Lights: handlebar switches

AUTO
  ├── Throttle: CAN setpoint → DAC
  ├── Gear: CAN setpoint → relays
  ├── Lights: CAN 0x302 → GPIOs
  └── Steering: active control with dynamic clamp

ESTOP
  ├── Throttle: 0V forced
  ├── Gear: all OFF
  ├── Brake: max stroke
  └── Lights: brake ON, all others OFF
\end{verbatim}

The top-level mode FSM selects which sub-behavior runs. In the E-Trike, this is
implemented by having each task check \texttt{g\_mode} and branch accordingly,
rather than a formal hierarchical statechart.

\begin{center}\rule{0.5\linewidth}{0.5pt}\end{center}

\section{6. Common Pitfalls}\label{common-pitfalls}

\begin{longtable}[]{@{}
  >{\raggedright\arraybackslash}p{(\columnwidth - 4\tabcolsep) * \real{0.3333}}
  >{\raggedright\arraybackslash}p{(\columnwidth - 4\tabcolsep) * \real{0.4815}}
  >{\raggedright\arraybackslash}p{(\columnwidth - 4\tabcolsep) * \real{0.1852}}@{}}
\toprule\noalign{}
\begin{minipage}[b]{\linewidth}\raggedright
Pitfall
\end{minipage} & \begin{minipage}[b]{\linewidth}\raggedright
What happens
\end{minipage} & \begin{minipage}[b]{\linewidth}\raggedright
Fix
\end{minipage} \\
\midrule\noalign{}
\endhead
\bottomrule\noalign{}
\endlastfoot
\textbf{Ghost state} & Code falls through to unintended behavior & Always have a
\texttt{default:} case in switch, even if it just logs an error \\
\textbf{Hidden transitions} & A state change happens in a low-level function,
not visible in the FSM & All \texttt{g\_mode\ =\ X} assignments go through
\texttt{mode\_set()}. Grep for \texttt{g\_mode\ =} to catch violations. \\
\textbf{Re-entrant transitions} & Entry action triggers another state change →
infinite loop & Guard: \texttt{if\ (g\_mode\ ==\ new\_mode)\ return;} at the top
of \texttt{mode\_set()} \\
\textbf{Non-atomic state write} & A 32-bit write on an 8-bit MCU is non-atomic &
Use \texttt{uint8\_t} or \texttt{enum\ class\ :\ uint8\_t} for state variables
--- single-instruction write on ARM/XTensa \\
\textbf{Race between FSM and ISR} & ISR sets state while FSM is mid-entry-action
& ISR only sets flags/queues. FSM checks them synchronously. \\
\end{longtable}

\begin{center}\rule{0.5\linewidth}{0.5pt}\end{center}

\section{7. The Boot→Listen→Active Pattern}\label{the-bootlistenactive-pattern}

A recurring FSM shape in the E-Trike, used for both steering and brake actuator
startup:

\begin{verbatim}
BOOT_WAIT ──► LISTEN_SYNC ──► ACTIVE ──► FAULT
  500ms        await status    50 Hz TX    timeout
  no TX        read position   monitor     silent
\end{verbatim}

This is a \textbf{safe bootstrapping pattern} for CAN actuators. Transmitting
before the actuator is ready can latch a fault; transmitting the wrong position
(e.g., 0° on a turned wheel) can cause dangerous transients. See
{[}{[}listen-before-speaking{]}{]} for the full rationale.

\begin{center}\rule{0.5\linewidth}{0.5pt}\end{center}

\emph{See also: {[}{[}listen-before-speaking{]}{]} for the boot sequence
pattern, {[}{[}defense-in-depth-safety{]}{]} for the ESTOP absorbing state,
\texttt{architecture.md} §3 for the mode FSM, §7.6 for steering FSM, §8.6 for
brake FSM.}


% ═══════════════════════════════════════════════════════════════════════
%  PART IV -- STANDARDS (standards/)
% ═══════════════════════════════════════════════════════════════════════

\part{Standards}
\label{part:standards}

\textit{Industry standards that inform the E-Trike architecture. These are the pillars the design rests on --- not just reference material but active constraints on every design decision.}

\chapter{AUTOSAR E2E --- End-to-End Communication
Protection}\label{autosar-e2e-end-to-end-communication-protection}

AUTOSAR (AUTomotive Open System ARchitecture) defines patterns for safe
communication over unreliable transport. We do not implement the full AUTOSAR
stack, but we adopt its core E2E Protection principles for CAN frames.

\section{What we apply}\label{what-we-apply}

\subsection{Alive counter}\label{alive-counter}

Every safety-relevant periodic frame carries a monotonically incrementing
counter. A frozen counter = a frozen sender (MCU hung but CAN controller still
DMA-ing from buffer).

\begin{longtable}[]{@{}ll@{}}
\toprule\noalign{}
Frame & Counter field \\
\midrule\noalign{}
\endhead
\bottomrule\noalign{}
\endlastfoot
\texttt{0x7FF} HEARTBEAT (all) & \texttt{u8\ alive\_ctr} --- increments every
500ms \\
\texttt{0x200} VCU\_SES\_REQ & \texttt{VCU\_SES\_RollCnt} --- 4-bit, 0--15,
every 20ms \\
\texttt{0x720} VCU\_SEB\_REQ & \texttt{VCU\_SEB\_RollCnt} --- 4-bit, 0--15,
every 20ms \\
\end{longtable}

\subsection{Checksum}\label{checksum}

The SYNTREE units require XOR-based checksums on every command frame. Same
frames with different data produce different checksums --- corrupted frames are
detected.

\begin{longtable}[]{@{}ll@{}}
\toprule\noalign{}
Frame & Algorithm \\
\midrule\noalign{}
\endhead
\bottomrule\noalign{}
\endlastfoot
\texttt{0x200} VCU\_SES\_REQ & \texttt{XOR(bytes{[}0..6{]})\ \^{}\ 0xFF} \\
\texttt{0x720} VCU\_SEB\_REQ & \texttt{XOR(bytes{[}0..6{]})\ \^{}\ 0xFF} \\
\end{longtable}

\subsection{Timeout supervision}\label{timeout-supervision}

Receivers independently track the arrival time of expected frames. If no valid
frame arrives within the timeout, the receiver takes autonomous safety action.

\begin{longtable}[]{@{}llll@{}}
\toprule\noalign{}
Monitor & Expects & Timeout & Action \\
\midrule\noalign{}
\endhead
\bottomrule\noalign{}
\endlastfoot
SYS & RT \texttt{0x7FF} & 200ms & ESTOP (AUTO) \\
RT & SYS \texttt{0x7FF} & 200ms & CAN \texttt{0x001} ESTOP \\
RT & Jetson \texttt{0x7FF} / \texttt{0x300} & 500ms & Zero setpoints \\
EPS-C & \texttt{0x200} every 20ms & \textasciitilde20ms & Internal comm fault \\
SEB & \texttt{0x720} every 20ms & \textasciitilde20ms & Internal comm fault \\
RT & \texttt{SES\_StrAngle} following error & abs(error) \textgreater5° for
300ms & ESTOP \\
\end{longtable}

\subsection{E2E Profile 1 (simplified)}\label{e2e-profile-1-simplified}

AUTOSAR defines E2E Profile 1 as: CRC + Alive Counter + Data ID. We implement a
simplified version:

\begin{longtable}[]{@{}ll@{}}
\toprule\noalign{}
AUTOSAR element & Our equivalent \\
\midrule\noalign{}
\endhead
\bottomrule\noalign{}
\endlastfoot
CRC & SYNTREE checksum (\texttt{XOR\ \^{}\ 0xFF}) \\
Alive Counter & Rolling counter / \texttt{alive\_ctr} \\
Data ID & Implicit --- CAN ID uniquely identifies the message type \\
Timeout & Per-receiver, configurable \\
\end{longtable}

\section{What we do NOT use}\label{what-we-do-not-use}

\begin{itemize}
\tightlist
\item
  Full AUTOSAR RTE (Runtime Environment)
\item
  AUTOSAR COM stack
\item
  E2E Profiles 2--6 (CRC-8/CRC-16 variants)
\item
  Configuration via ARXML
\end{itemize}

\section{References}\label{references}

\begin{itemize}
\tightlist
\item
  AUTOSAR SWS E2E Library (R20-11)
\item
  AUTOSAR RS Safety Extensions (R20-11)
\end{itemize}


\chapter{FreeRTOS + ESP-IDF --- RTOS \&
Framework}\label{freertos-esp-idf-rtos-framework}

Both ESP32-S3 nodes run \textbf{ESP-IDF} (Espressif IoT Development Framework)
with its built-in \textbf{FreeRTOS} kernel. Preemptive, tickless idle, dual-core
Xtensa LX7 @ 240 MHz, 1000 Hz scheduler tick.

\section{What we use}\label{what-we-use}

\subsection{Task isolation}\label{task-isolation}

Tasks communicate through \textbf{queues, not shared state}. Queues are
thread-safe by design --- one task writes, another reads. No mutexes, no
semaphores. This keeps the system deterministic and deadlock-free.

\begin{longtable}[]{@{}
  >{\raggedright\arraybackslash}p{(\columnwidth - 4\tabcolsep) * \real{0.3750}}
  >{\raggedright\arraybackslash}p{(\columnwidth - 4\tabcolsep) * \real{0.2083}}
  >{\raggedright\arraybackslash}p{(\columnwidth - 4\tabcolsep) * \real{0.4167}}@{}}
\toprule\noalign{}
\begin{minipage}[b]{\linewidth}\raggedright
Pattern
\end{minipage} & \begin{minipage}[b]{\linewidth}\raggedright
API
\end{minipage} & \begin{minipage}[b]{\linewidth}\raggedright
Use case
\end{minipage} \\
\midrule\noalign{}
\endhead
\bottomrule\noalign{}
\endlastfoot
\texttt{xQueueOverwrite} (depth 4) & \texttt{xQueueOverwrite(q,\ \&item)} &
\texttt{cmd\_queue}, \texttt{setpoint\_queue} --- only latest value matters \\
\texttt{xQueueSend} timeout=0 (depth 16) & \texttt{xQueueSend(q,\ \&item,\ 0)} &
CAN RX queues --- drop if consumer is slow (FIFO back-pressure) \\
\texttt{xQueueSend} timeout=0 (depth 8) & \texttt{xQueueSend(q,\ \&item,\ 0)} &
Gateway queues --- non-critical forwarded frames \\
\texttt{portMAX\_DELAY} block &
\texttt{xQueueReceive(q,\ \&item,\ portMAX\_DELAY)} & \texttt{dispatch\_task}
--- zero CPU when idle \\
\end{longtable}

\subsection{Fixed-rate loops}\label{fixed-rate-loops}

\texttt{vTaskDelayUntil()} corrects for execution time, keeping periods stable
even if one iteration runs long. Jitter \textless{} 50 µs on ESP32-S3.

\subsection{Atomic variables}\label{atomic-variables}

For single values where a queue is overkill:

\begin{longtable}[]{@{}
  >{\raggedright\arraybackslash}p{(\columnwidth - 6\tabcolsep) * \real{0.3226}}
  >{\raggedright\arraybackslash}p{(\columnwidth - 6\tabcolsep) * \real{0.1935}}
  >{\raggedright\arraybackslash}p{(\columnwidth - 6\tabcolsep) * \real{0.2581}}
  >{\raggedright\arraybackslash}p{(\columnwidth - 6\tabcolsep) * \real{0.2258}}@{}}
\toprule\noalign{}
\begin{minipage}[b]{\linewidth}\raggedright
Variable
\end{minipage} & \begin{minipage}[b]{\linewidth}\raggedright
Type
\end{minipage} & \begin{minipage}[b]{\linewidth}\raggedright
Writes
\end{minipage} & \begin{minipage}[b]{\linewidth}\raggedright
Reads
\end{minipage} \\
\midrule\noalign{}
\endhead
\bottomrule\noalign{}
\endlastfoot
\texttt{g\_mode} & \texttt{std::atomic\textless{}int\textgreater{}} & safety,
dispatch, mode & motor, brake, control, can\_tx, watchdog \\
\texttt{g\_brake\_request\_kpa} &
\texttt{std::atomic\textless{}int32\_t\textgreater{}} & dispatch & control \\
\texttt{g\_steer\_feedback\_angle\_mdeg} &
\texttt{std::atomic\textless{}int32\_t\textgreater{}} & dispatch & control \\
\texttt{g\_light\_state} &
\texttt{std::atomic\textless{}LightState\textgreater{}} & dispatch &
lights\_task \\
\texttt{g\_obstacle\_mm} &
\texttt{std::atomic\textless{}uint32\_t\textgreater{}} & obstacle\_task &
control \\
\texttt{g\_throttle\_speed\_mmps} &
\texttt{std::atomic\textless{}int16\_t\textgreater{}} & throttle\_task &
motor\_task, can\_tx \\
\end{longtable}

\texttt{std::atomic\textless{}int\textgreater{}} is lock-free on Xtensa LX7
(single word, atomic store/load).

\subsection{Priority-based preemption}\label{priority-based-preemption}

Higher priority tasks preempt lower ones within the same tick. The scheduler
always runs the highest-priority ready task. Priority 5 (safety, CAN RX) is
non-negotiable --- these tasks must never be delayed.

\subsection{Tickless idle}\label{tickless-idle}

FreeRTOS enters tickless idle when no task is ready. The ESP32-S3's deep sleep
features are not used (CAN bus must stay active), but the scheduler itself does
not waste cycles on an idle tick.

\section{What we do NOT use}\label{what-we-do-not-use}

\begin{itemize}
\tightlist
\item
  SMP (Symmetric Multi-Processing) task pinning --- both cores use the same
  priority scheme
\item
  FreeRTOS task notifications --- queues are preferred for auditability
\item
  Software timers --- dedicated tasks with \texttt{vTaskDelayUntil} instead
\item
  ESP-IDF's \texttt{esp\_timer} for control loops --- FreeRTOS primitives only
\end{itemize}

\section{References}\label{references}

\begin{itemize}
\tightlist
\item
  FreeRTOS Kernel Developer Docs (v10.5+)
\item
  ESP-IDF Programming Guide (v5.x)
\item
  ESP32-S3 Technical Reference Manual
\end{itemize}


\chapter{ISO 11898 --- CAN Bus}\label{iso-11898-can-bus}

The physical and data-link layer standard for Controller Area Network. Both CAN
buses in the E-Trike operate under ISO 11898-1:2015 (Classical CAN).

\section{What we use}\label{what-we-use}

\begin{longtable}[]{@{}
  >{\raggedright\arraybackslash}p{(\columnwidth - 2\tabcolsep) * \real{0.4500}}
  >{\raggedright\arraybackslash}p{(\columnwidth - 2\tabcolsep) * \real{0.5500}}@{}}
\toprule\noalign{}
\begin{minipage}[b]{\linewidth}\raggedright
Feature
\end{minipage} & \begin{minipage}[b]{\linewidth}\raggedright
Our config
\end{minipage} \\
\midrule\noalign{}
\endhead
\bottomrule\noalign{}
\endlastfoot
Physical layer & ISO 11898-2 (high-speed CAN, differential two-wire + GND) \\
Frame format & CAN 2.0B, 11-bit standard identifiers \\
Bitrate & 500 kbit/s (both buses) \\
Transceiver & SN65HVD230 (3.3V compatible) \\
Termination & 120Ω at each end of the bus \\
Payload (DLC) & 0--8 bytes \\
Endianness & Big-endian (MSB first) for our messages; Motorola LSB for SYNTREE
units \\
\end{longtable}

\section{Principles we rely on}\label{principles-we-rely-on}

\textbf{Bitwise arbitration.} Lower CAN ID = higher priority. No data lost on
collision --- the lower ID wins and continues transmitting. Safety-critical IDs
(\texttt{0x001} ESTOP, \texttt{0x01X} safety/dcdc) occupy the lowest ID block.

\textbf{Content-addressed, not node-addressed.} The CAN ID says \emph{what} the
data is, not \emph{who} sent it or \emph{who} should receive it. Every node
hears every frame and uses hardware acceptance filters to decide what to
process.

\textbf{One sender per ID per bus.} Two nodes sending the same ID on the same
bus causes ACK errors and bus degradation. Our IDs are unique per sender (except
heartbeat \texttt{0x7FF} which multiple nodes share with distinct alive
counters).

\textbf{Error confinement.} Each CAN controller tracks Transmit Error Counter
(TEC) and Receive Error Counter (REC). TEC \textgreater{} 255 → bus-off state
(node disconnects). TEC/REC reported by SYS via \texttt{0x600} diagnostics.

\section{Ground reference (the ``third
wire'')}\label{ground-reference-the-third-wire}

CAN is differential but transceivers need a common-mode voltage reference. If
nodes are on separate power supplies, the ground pins must be connected
(CAN\_GND). Without it, voltages drift beyond the common-mode range and
transceivers burn out or drop packets.

\section{References}\label{references}

\begin{itemize}
\tightlist
\item
  ISO 11898-1:2015 --- Data link layer and physical signaling
\item
  ISO 11898-2:2016 --- High-speed medium access unit
\item
  Bosch CAN Specification 2.0 (1991) --- Original protocol definition
\end{itemize}


\chapter{ISO 26262 --- Functional Safety (Road
Vehicles)}\label{iso-26262-functional-safety-road-vehicles}

The international standard for functional safety of electrical/electronic
systems in production road vehicles. Our architecture targets \textbf{ASIL-C}
for the steering subsystem.

\section{What we apply}\label{what-we-apply}

\subsection{ASIL decomposition}\label{asil-decomposition}

\begin{longtable}[]{@{}
  >{\raggedright\arraybackslash}p{(\columnwidth - 4\tabcolsep) * \real{0.3235}}
  >{\raggedright\arraybackslash}p{(\columnwidth - 4\tabcolsep) * \real{0.3529}}
  >{\raggedright\arraybackslash}p{(\columnwidth - 4\tabcolsep) * \real{0.3235}}@{}}
\toprule\noalign{}
\begin{minipage}[b]{\linewidth}\raggedright
Subsystem
\end{minipage} & \begin{minipage}[b]{\linewidth}\raggedright
ASIL target
\end{minipage} & \begin{minipage}[b]{\linewidth}\raggedright
Rationale
\end{minipage} \\
\midrule\noalign{}
\endhead
\bottomrule\noalign{}
\endlastfoot
Steering (EPS-C) & \textbf{ASIL-C} & Steer-by-wire --- loss of steering control
→ severe hazard \\
Braking (SEB) & ASIL-B & Electro-hydraulic --- mechanical backup via lever \\
Motor control & ASIL-A & ESC/Motor controller handles low-level safety \\
Body (lights, DCDC) & QM & Non-safety; failure does not cause harm \\
\end{longtable}

\subsection{Fault Tolerant Time Interval
(FTTI)}\label{fault-tolerant-time-interval-ftti}

The maximum time a fault can persist before the system must reach a safe state.

\begin{longtable}[]{@{}
  >{\raggedright\arraybackslash}p{(\columnwidth - 6\tabcolsep) * \real{0.1628}}
  >{\raggedright\arraybackslash}p{(\columnwidth - 6\tabcolsep) * \real{0.1395}}
  >{\raggedright\arraybackslash}p{(\columnwidth - 6\tabcolsep) * \real{0.4419}}
  >{\raggedright\arraybackslash}p{(\columnwidth - 6\tabcolsep) * \real{0.2558}}@{}}
\toprule\noalign{}
\begin{minipage}[b]{\linewidth}\raggedright
Fault
\end{minipage} & \begin{minipage}[b]{\linewidth}\raggedright
FTTI
\end{minipage} & \begin{minipage}[b]{\linewidth}\raggedright
Our response time
\end{minipage} & \begin{minipage}[b]{\linewidth}\raggedright
Mechanism
\end{minipage} \\
\midrule\noalign{}
\endhead
\bottomrule\noalign{}
\endlastfoot
RT→SYS comm loss & \textless200 ms & 200 ms & \texttt{0x7FF} alive counter
timeout → ESTOP \\
SYS→RT comm loss & \textless200 ms & 200 ms & \texttt{0x7FF} alive counter
timeout → RT sends \texttt{0x001} \\
EPS-C comm loss & \textless50 ms & 20 ms (1 frame gap) & EPS-C internal timeout
fault \\
ESTOP button press & \textless50 ms & 50 ms (20 Hz poll) & \texttt{safety\_task}
→ \texttt{mode\_set(Estop)} \\
Jetson→RT comm loss & \textless500 ms & 500 ms & Watchdog → zero setpoints
(controlled stop) \\
\end{longtable}

\subsection{Defense-in-depth}\label{defense-in-depth}

ESTOP is not a single check --- it's layered across multiple tasks and hardware
domains:

\begin{enumerate}
\def\labelenumi{\arabic{enumi}.}
\tightlist
\item
  \textbf{Hardware} --- ESTOP button (normally-closed, active-low) wired
  directly to MCU GPIO with pull-up
\item
  \textbf{Firmware} --- \texttt{safety\_task} (prio 5) polls ESTOP GPIO at 20 Hz
  and calls \texttt{mode\_set(Estop)}
\item
  \textbf{Firmware} --- \texttt{motor\_set\_speed()},
  \texttt{gear\_set\_outputs()}, \texttt{dcdc\_set\_output()} all independently
  check \texttt{g\_mode} before actuating
\item
  \textbf{CAN} --- Any node can trigger ESTOP via \texttt{0x001} (highest
  priority ID)
\item
  \textbf{External watchdog} --- Hardware IC (TPS3850) resets MCU if
  \texttt{safety\_task} stops toggling GPIO
\end{enumerate}

\subsection{Freedom from interference}\label{freedom-from-interference}

\begin{itemize}
\tightlist
\item
  RT and SYS are \textbf{physically separate MCUs} on different PCBs
\item
  Low-level CAN is \textbf{electrically isolated} from high-level CAN (RT uses
  separate TWAI + MCP2515 controllers)
\item
  SYS has \textbf{no direct connection to Jetson} --- all cross-bus traffic goes
  through RT gateway
\item
  Manual mode is \textbf{pass-through} --- rider has direct
  mechanical/electrical control independent of CAN or firmware
\end{itemize}

\subsection{Safe state}\label{safe-state}

On any ESTOP trigger, the system reaches safe state: - Motor: MCP4725 DAC = 0V,
all gear relays OFF - Brake: SEB commanded to max stroke (27mm) via
\texttt{0x720} - Steering: EPS-C commands stopped (timeout-faults internally) -
12V: DC-DC converter disabled via \texttt{0x012} + accessory relay OFF -
External watchdog: stops toggling → MCU resets to safe defaults

\section{What we do NOT apply}\label{what-we-do-not-apply}

\begin{itemize}
\tightlist
\item
  Full ISO 26262 lifecycle (HARA, FMEA, safety case documentation) --- this is a
  prototype
\item
  ASIL-D (highest) --- not required for a three-wheeled vehicle at our speed
  class
\item
  Redundant MCU lockstep --- single-core ESP32-S3 is sufficient with external
  watchdog
\item
  ISO 26262-6 software development process --- FreeRTOS + ESP-IDF are not
  certified
\end{itemize}

\section{References}\label{references}

\begin{itemize}
\tightlist
\item
  ISO 26262-1:2018 --- Vocabulary
\item
  ISO 26262-3:2018 --- Concept phase (HARA)
\item
  ISO 26262-5:2018 --- Product development at the hardware level
\item
  ISO 26262-6:2018 --- Product development at the software level
\end{itemize}


% ═══════════════════════════════════════════════════════════════════════
%  APPENDIX
% ═══════════════════════════════════════════════════════════════════════

\appendix
\part*{Appendix}
\addcontentsline{toc}{part}{Appendix}

\chapter{Reference Index}

\section{Core Documents}

\begin{itemize}[leftmargin=*,nosep]
  \item \texttt{architecture.md} --- System architecture (topology, CAN catalog, FSMs, RT/SYS design)
  \item \texttt{can-dictionary.md} --- Bit-level CAN signal layouts for all IDs on both buses
\end{itemize}

\section{Component Documentation (\texttt{docs/})}

\begin{itemize}[leftmargin=*,nosep]
  \item \texttt{steering-unit.md} --- SYNTREE EPS-C protocol reference and C++ implementation
  \item \texttt{brake-unit.md} --- SYNTREE SEB protocol reference and C++ implementation
  \item \texttt{physics-model.md} --- Tricycle kinematics (forward/inverse, rollover, slip angles)
  \item \texttt{listen-before-speaking.md} --- CAN actuator safe bootstrapping pattern
  \item \texttt{can-gateway-bridging.md} --- CAN gateway forwarding rules and implementation
  \item \texttt{defense-in-depth-safety.md} --- 8-layer safety architecture
  \item \texttt{syntree-security-protocol.md} --- Rolling counter + XOR checksum
  \item \texttt{high-voltage-isolation.md} --- 72V galvanic isolation (TLP281, relays, fuses, TVS)
  \item \texttt{distributed-architecture.md} --- Three-node rationale (Jetson/RT/SYS split)
  \item \texttt{actuator-interfacing.md} --- MCP4725 DAC throttle, gear pass-through, relay logic
  \item \texttt{external-watchdog.md} --- External watchdog IC design and testing
  \item \texttt{pid-speed-control.md} --- PID speed control design and tuning
  \item \texttt{given-documentation-ocr.md} --- Source PDF extraction (EPS-C/SEB/DCDC specs)
\end{itemize}

\section{Theory Notes (\texttt{notes/})}

\begin{itemize}[leftmargin=*,nosep]
  \item \texttt{can-protocol.md} --- CAN protocol theory (arbitration, frame types, standards)
  \item \texttt{can-hardware-basics.md} --- CAN physical layer (termination, topology, transceivers)
  \item \texttt{can-troubleshooting.md} --- CAN debugging (common mistakes, error states, tools)
  \item \texttt{rtos-task-design.md} --- RTOS task design (priorities, scheduling, queues, tick rate)
  \item \texttt{pid-control.md} --- PID control fundamentals (P/I/D terms, integral windup, tuning)
  \item \texttt{state-machine-design.md} --- State machine design (FSMs, hierarchical states, C++ patterns)
  \item \texttt{heartbeat-monitoring.md} --- Heartbeat \& liveness (alive counters, FTTI, timeout selection)
  \item \texttt{distributed-safety-patterns.md} --- Distributed safety (defense in depth, ESTOP, fail-safe)
  \item \texttt{endianness-binary-protocols.md} --- Endianness \& binary protocols (big/little endian, struct packing)
  \item \texttt{can-gateway-design.md} --- CAN gateway design (multi-bus architecture, forwarding categories)
  \item \texttt{analog-interfacing.md} --- Analog interfacing (ADC, DAC, optoisolators, relays, protection)
\end{itemize}

\section{Standards (\texttt{standards/})}

\begin{itemize}[leftmargin=*,nosep]
  \item \texttt{iso-11898-can.md} --- CAN bus physical and data-link layer (ISO 11898-1:2015)
  \item \texttt{iso-26262-functional-safety.md} --- Functional safety for road vehicles (targeting ASIL-C steering)
  \item \texttt{autosar-e2e-protection.md} --- End-to-End communication protection (alive counters, checksums, timeouts)
  \item \texttt{freertos-esp-idf.md} --- FreeRTOS + ESP-IDF framework (task isolation, queues, tick rate)
\end{itemize}

\section{Firmware \& Source}

\begin{itemize}[leftmargin=*,nosep]
  \item \texttt{rt-esp32/} --- RT ESP32-S3 firmware (PlatformIO, ESP-IDF + FreeRTOS)
  \item \texttt{sys-esp32/} --- SYS ESP32-S3 firmware (PlatformIO, ESP-IDF + FreeRTOS)
  \item \texttt{shared/} --- Shared CAN protocol headers between RT and SYS
  \item \texttt{legacy/} --- Deprecated design notes and historical references
  \item \texttt{given-documentation.pdf} --- Original Syntree specification PDF (58 pages)
  \item \texttt{given-documentation-ocr.pdf} --- Alternate OCR scan of original specification
  \item \texttt{work-orders/} --- Manufacturing and assembly work orders
\end{itemize}

\vfill

\begin{center}
  \rule{0.5\textwidth}{0.4pt}\\[4mm]
  {\small\textit{E-Trike System Documentation v0.0.3-alpha}}\\[2mm]
  {\small Generated 2026-06-24 --- Covers architecture.md, can-dictionary.md, 14 component docs, 12 theory notes, 4 standards}\\[2mm]
  {\small Compiled with \LaTeX\ --- PDF/A-2b compatible}
\end{center}

\end{document}
